% Options for packages loaded elsewhere
\PassOptionsToPackage{unicode}{hyperref}
\PassOptionsToPackage{hyphens}{url}
%
\documentclass[
  openright]{book}
\usepackage{amsmath,amssymb}
\usepackage{setspace}
\usepackage{iftex}
\ifPDFTeX
  \usepackage[T1]{fontenc}
  \usepackage[utf8]{inputenc}
  \usepackage{textcomp} % provide euro and other symbols
\else % if luatex or xetex
  \usepackage{unicode-math} % this also loads fontspec
  \defaultfontfeatures{Scale=MatchLowercase}
  \defaultfontfeatures[\rmfamily]{Ligatures=TeX,Scale=1}
\fi
\usepackage{lmodern}
\ifPDFTeX\else
  % xetex/luatex font selection
\fi
% Use upquote if available, for straight quotes in verbatim environments
\IfFileExists{upquote.sty}{\usepackage{upquote}}{}
\IfFileExists{microtype.sty}{% use microtype if available
  \usepackage[]{microtype}
  \UseMicrotypeSet[protrusion]{basicmath} % disable protrusion for tt fonts
}{}
\makeatletter
\@ifundefined{KOMAClassName}{% if non-KOMA class
  \IfFileExists{parskip.sty}{%
    \usepackage{parskip}
  }{% else
    \setlength{\parindent}{0pt}
    \setlength{\parskip}{6pt plus 2pt minus 1pt}}
}{% if KOMA class
  \KOMAoptions{parskip=half}}
\makeatother
\usepackage{xcolor}
\usepackage[paperwidth=6in,paperheight=9in,inner=0.875in,outer=0.625in,top=0.75in,bottom=0.75in,includehead,includefoot]{geometry}
\usepackage{longtable,booktabs,array}
\usepackage{calc} % for calculating minipage widths
% Correct order of tables after \paragraph or \subparagraph
\usepackage{etoolbox}
\makeatletter
\patchcmd\longtable{\par}{\if@noskipsec\mbox{}\fi\par}{}{}
\makeatother
% Allow footnotes in longtable head/foot
\IfFileExists{footnotehyper.sty}{\usepackage{footnotehyper}}{\usepackage{footnote}}
\makesavenoteenv{longtable}
\setlength{\emergencystretch}{3em} % prevent overfull lines
\providecommand{\tightlist}{%
  \setlength{\itemsep}{0pt}\setlength{\parskip}{0pt}}
\setcounter{secnumdepth}{-\maxdimen} % remove section numbering
\ifLuaTeX
\usepackage[bidi=basic]{babel}
\else
\usepackage[bidi=default]{babel}
\fi
\babelprovide[main,import]{american}
% get rid of language-specific shorthands (see #6817):
\let\LanguageShortHands\languageshorthands
\def\languageshorthands#1{}
% DFP Book — LaTeX preamble additions
% Loaded via --include-in-header (verbatim, after pandoc's default preamble)

% graphicx: for \includegraphics (chapter art pages)
\usepackage{graphicx}

% unicode-math: full Unicode math support with TeX Gyre Pagella Math
% Must be loaded before \setmainfont if fontspec not yet active.
% Pandoc only loads fontspec when mainfont= is set; we did not set it,
% so fontspec is not yet loaded here — unicode-math will load it.
\usepackage{unicode-math}

% Main text font: TeX Gyre Pagella (Palatino-compatible)
% Fallback to its companion math font for characters not in the text font:
% subscripts (₀–₉, ₋, ₙ), superscripts (⁰⁻ⁿ⁺⁶), letterlike (ℕ ℤ ℝ), etc.
% Fallback 1: TeX Gyre Pagella Math — letterlike symbols ℕ ℤ ℝ, math operators
% Fallback 2: FreeSerif — subscripts ₀₁…ₙ₋, superscripts ⁰⁻ⁿ⁺, geometric shapes □
\directlua{
  luaotfload.add_fallback("mainfallback", {
    "TeX Gyre Pagella Math:",
    "FreeSerif:"
  })
}
\setmainfont[RawFeature={fallback=mainfallback}]{TeX Gyre Pagella}
% Companion math font with full Unicode math symbol coverage
\setmathfont{TeX Gyre Pagella Math}
% Sans-serif
\setsansfont{TeX Gyre Heros}
% Monospace: set below with FreeSerif fallback for Mathematical Alphanumeric Symbols

% CJK support via luatexja — provides fallback for Japanese haiku text
\usepackage{luatexja-fontspec}
\setmainjfont{Noto Serif CJK JP}

% luatexja classifies many Unicode characters as JAchar (Japanese), routing
% them to NotoSerifCJKJP which lacks them → missing-glyph boxes in the PDF.
% Fix: declare each affected block as ALchar (non-Japanese) so the normal
% Latin/math fonts handle them.  Japanese text (hiragana U+3040–, katakana,
% kanji U+4E00–) is unaffected — none of those are in the ranges below.
%
% Range 9: Typographic quotes U+2000–U+206F (General Punctuation, incl. ' ')
%   Fixes "Author' s" → "Author's" (xkanjiskip after closing-quote glyph).
\ltjdefcharrange{9}{"2000-"206F}
\ltjsetparameter{jacharrange={-9}}
%
% Range 10: Superscripts/subscripts and Letterlike Symbols
%   U+2070–U+214F covers ⁰⁻ⁿ (U+2070–), ₀₁…₉ (U+2080–), ℕℤℝ (U+2100–).
\ltjdefcharrange{10}{"2070-"214F}
\ltjsetparameter{jacharrange={-10}}
%
% Range 11: Mathematical Operators U+2200–U+22FF
%   Covers ∀ ∈ ∧ ∨ ∘ ≃ ≡ ⊂ ⊎ ⊢ ⊤ ⊥ ⋮ and all other math operators.
\ltjdefcharrange{11}{"2200-"22FF}
\ltjsetparameter{jacharrange={-11}}
%
% Range 12: Miscellaneous Technical + Box Drawing + Geometric Shapes U+2300–U+25FF
%   Covers ─ (U+2500), box-drawing chars, and □ (U+25A1) used as QED symbol.
\ltjdefcharrange{12}{"2300-"25FF}
\ltjsetparameter{jacharrange={-12}}
%
% Range 13: Mathematical Alphanumeric Symbols U+1D400–U+1D7FF
%   Covers 𝟘 (U+1D7D8) and 𝟙 (U+1D7D9) used in Agda for empty/unit types.
%   These are not in DejaVu Sans Mono; add a monofont fallback to FreeSerif.
\ltjdefcharrange{13}{"1D400-"1D7FF}
\ltjsetparameter{jacharrange={-13}}

% Monospace fallback for Mathematical Alphanumeric Symbols (𝟘 𝟙 etc.)
% DejaVu Sans Mono lacks this block; FreeSerif covers it and is available.
\directlua{
  luaotfload.add_fallback("monofallback", {
    "FreeSerif:"
  })
}
\setmonofont[Scale=0.88, RawFeature={fallback=monofallback}]{DejaVu Sans Mono}

% Subtitle support — pandoc uses \providecommand{\subtitle}[1]{} (no-op).
% Define \subtitle to store the text; \posttitle renders it below the title.
\usepackage{titling}
\newcommand{\thesubtitle}{}
\newcommand{\subtitle}[1]{\renewcommand{\thesubtitle}{#1}}
\posttitle{\par\large\itshape\thesubtitle\par\end{center}\vskip 0.5em}

% \linethickness is a 1-arg command in LaTeX2e picture mode, but pandoc uses
% it as a bare dimension in \rule{0.5\linewidth}{\linethickness} for <hr>.
% Redefine as a fixed dimension so \rule works without picture mode.
\def\linethickness{0.4pt}

% Better tables: tabularx for stretchy columns, booktabs for rules
\usepackage{tabularx}
\usepackage{booktabs}
% Landscape pages (pdflscape rotates the page in the PDF viewer)
\usepackage{pdflscape}
% rotating: sidewaystable/sidewaysfigure floats for landscape tables
\usepackage{rotating}

% Commutative diagrams via tikz-cd
% Use \begin{tikzcd}...\end{tikzcd} inside $$...$$ for diagrams.
% Do NOT use \xymatrix — the xy package is not loaded.
\usepackage{tikz-cd}

% Dirac notation: \ket{..} \bra{..} \braket{..|..}
\usepackage{braket}

% ZX-calculus diagrams (Chapter 8 onward): \begin{ZX} ... \end{ZX}
% NOTE: do NOT \usepackage{zx-calculus} — with TeX Live 2023 its
% \ProvidesPackage line has a non-date version string ("v2.0 ...") that the
% LaTeX kernel chokes on ("Missing = inserted for \ifnum"). Loading the TikZ
% library directly is equivalent and works.
\usetikzlibrary{zx-calculus}

% String-diagram styles (Chapters 1 onward). Diagrams read LEFT TO RIGHT.
%   \begin{tikzpicture}[sd] ... \end{tikzpicture}
%   wire   — a quantum wire           sdbox  — a process box
%   sdstate / sdeffect — a state (no input) / an effect (no output)
%   sdnode — the unnamed three-legged node (Chapter 7 only): \node[sdnode] at (x,y) {};
%   sdcut  — the cut through a diagram (Chapter 11+): \draw[sdcut] ...;  not a wire
%   cup: \draw[wire] (x,y1) to[out=180,in=180,looseness=2] (x,y2);   (opens right: a STATE)
%   cap: \draw[wire] (x,y1) to[out=0,in=0,looseness=2] (x,y2);       (closes right: an EFFECT)
\tikzset{
  sd/.style={baseline=(current bounding box.center), line cap=round, x=1em, y=1em},
  wire/.style={line width=0.9pt},
  cwire/.style={line width=0.5pt, double, double distance=1.2pt},
  sdbox/.style={draw, line width=0.9pt, fill=white, minimum width=1.7em,
                minimum height=1.7em, inner sep=2pt, font=\small},
  sdstate/.style={sdbox, rounded corners=0.85em},
  sdeffect/.style={sdstate},
  sdlabel/.style={font=\footnotesize, inner sep=1pt},
  % the unnamed three-legged node of Chapter 7 (|000> + |111>, UNNORMALIZED); not a ZX spider
  sdnode/.style={circle, fill, inner sep=0pt, minimum size=0.56em},
  % the cut (Chapter 11+): a thin grey dotted curve through a diagram; never dashed
  % (dashed is Chapter 8's Hadamard edge). Not a wire: it counts no line.
  sdcut/.style={densely dotted, gray, line width=0.6pt},
}
% \sdcup{x}{y1}{y2} — a cup (STATE, opens to the right) joining rows y1 and y2 at column x
% \sdcap{x}{y1}{y2} — a cap (EFFECT, closes on the right) joining rows y1 and y2 at column x
% Both are UNNORMALIZED: cup = |00> + |11>.
\newcommand{\sdcup}[3]{\draw[wire] (#1,#2) to[out=180,in=180,looseness=1.7] (#1,#3);}
\newcommand{\sdcap}[3]{\draw[wire] (#1,#2) to[out=0,in=0,looseness=1.7] (#1,#3);}
% \sdground{x}{y} — the discard / ground symbol (Chapter 9 onward, Stage D): a wire ending
% in the electrical-ground glyph (short stem + shrinking bars). Connection point at (x,y),
% glyph to the right; draw the incoming wire up to (x,y), then \sdground{x}{y}.
\newcommand{\sdground}[2]{%
  \draw[wire] (#1,#2) -- (#1+0.5,#2);
  \draw[wire] (#1+0.5,#2-0.3) -- (#1+0.5,#2+0.3);
  \draw[wire] (#1+0.65,#2-0.18) -- (#1+0.65,#2+0.18);
  \draw[wire] (#1+0.8,#2-0.08) -- (#1+0.8,#2+0.08);
}

% Page headers and footers
\usepackage{fancyhdr}
\pagestyle{fancy}
\fancyhf{}
\fancyhead[LE]{\small\itshape Dirac, Feynman, Penrose}
\fancyhead[RO]{\small\itshape \leftmark}
\fancyfoot[C]{\thepage}
\renewcommand{\headrulewidth}{0.4pt}
% Strip the auto-counter from running chapter heads — show title text only.
\renewcommand{\chaptermark}[1]{\markboth{#1}{}}

% \cleartoleftpage — advance to the next even (left-hand) page.
% Used to place chapter art on the verso page before the chapter opens recto.
\newcommand{\cleartoleftpage}{%
  \clearpage
  \ifodd\value{page}\mbox{}\thispagestyle{empty}\clearpage\fi
}

% Chapter/section heading style
\usepackage{titlesec}
\titleformat{\chapter}[display]
  {\normalfont\large\bfseries\centering}
  {}{0pt}{\Large}
\titlespacing*{\chapter}{0pt}{40pt}{30pt}

% Fix titlesec + hyperref bookmark destination mismatch.
% titlesec rewrites \section/\chapter rendering; after a list environment
% \@currentHref still holds the last \item anchor (e.g. Item.18) instead
% of the new section anchor. Patching phantomsection into titlesec's
% display and block renderers resets \@currentHref before the bookmark
% is written, so PDF destinations land on the correct pages.
\usepackage{etoolbox}
\makeatletter
\patchcmd{\ttlh@display}{\parindent \z@}{\parindent \z@\phantomsection}{}{}
\patchcmd{\ttlh@block}{\parindent \z@}{\parindent \z@\phantomsection}{}{}
\patchcmd{\ttlh@hang}{\parindent \z@}{\parindent \z@\phantomsection}{}{}
\makeatother

% Part heading style — centred, generous spacing, no "Part N" label
\titleformat{\part}[display]
  {\normalfont\huge\bfseries\centering}
  {}{0pt}{\vspace{2em}\huge}
\titlespacing*{\part}{0pt}{80pt}{60pt}

% Suppress running headers on part-opener and empty pages
\fancypagestyle{plain}{\fancyhf{}\fancyfoot[C]{\thepage}\renewcommand{\headrulewidth}{0pt}}


% Index support — makeindex integration; intoc=true adds entry to ToC
\usepackage{imakeidx}
\makeindex[title=Index, intoc=true]

% ToC depth: show only chapters (depth 0) — keeps the ToC clean
\setcounter{tocdepth}{0}
\ifLuaTeX
  \usepackage{selnolig}  % disable illegal ligatures
\fi
\IfFileExists{bookmark.sty}{\usepackage{bookmark}}{\usepackage{hyperref}}
\IfFileExists{xurl.sty}{\usepackage{xurl}}{} % add URL line breaks if available
\urlstyle{same}
\hypersetup{
  pdftitle={Dirac, Feynman, Penrose},
  pdflang={en-US},
  pdfsubject={Quantum Computing, Diagrammatic Reasoning, Physics},
  pdfkeywords={quantum computing, ZX-calculus, string diagrams, quantum
information},
  hidelinks,
  pdfcreator={LaTeX via pandoc}}

\title{Dirac, Feynman, Penrose}
\usepackage{etoolbox}
\makeatletter
\providecommand{\subtitle}[1]{% add subtitle to \maketitle
  \apptocmd{\@title}{\par {\large #1 \par}}{}{}
}
\makeatother
\subtitle{An Entangled Golden Braid}
\author{}
\date{}

\begin{document}
\frontmatter
\maketitle

{
\setcounter{tocdepth}{2}
\tableofcontents
}
\setstretch{1.15}
\mainmatter
\hypertarget{overture-for-a-single-string}{%
\chapter{Overture for a Single
String}\label{overture-for-a-single-string}}

\begin{center}\rule{0.5\linewidth}{0.5pt}\end{center}

\emph{木枯や}\\
\emph{裏返りまた}\\
\emph{裏返る}

\emph{kogarashi ya}\\
\emph{uragaeri mata}\\
\emph{uragaeru}

\emph{The withering wind ---}\\
\emph{something turns over, and then}\\
\emph{turns over again.}

\begin{center}\rule{0.5\linewidth}{0.5pt}\end{center}

\emph{{[}Achilles's table, after supper. He is bent over a coin and a
much-scribbled sheet; the Tortoise has just let herself in.{]}}

\emph{{[}Outside, the first hard wind of the year is stripping the lane
of its last leaves and worrying at the window, which has gone black and
cold; a little of the draught finds its way in under the door. The fire
has not long been lit and has not yet made up its mind --- it catches,
gutters, catches again, and throws a light that will not sit still on
anything. On the wall above the hearth hangs a woodcut, Escher's
\textup{Day and Night}: the chequered fields between two towns, daylight
lying over the one side and night over the other. Supper is cleared; a
pot of tea stands cooling at the far corner of the table, and beside it
the coin, the pencil and the much-scribbled sheet are Achilles's.{]}}

\textbf{Tortoise:} You started without me.

\textbf{Achilles:} Only fidgeting. Come round to this side and sit ---
good, now the table reads the same way for us both. Was the road very
cold?

\textbf{Tortoise:} Cold enough that I'm grateful for the chair, and the
tea, and a table indoors. Don't mind me; you have the face of a man in
the middle of something.

\textbf{Achilles:} In the middle of it, and quite turned around in the
middle of it. I'm glad you've come.

\emph{{[}The Tortoise, settled beside him, holds her hands a moment
toward the grate, where the flames are still finding their feet.
Achilles reaches across for the cooling pot and pours. For a while there
is only the wind at the glass and the small business of getting
warm.{]}}

\hypertarget{first-challenge-a-coin-on-a-string}{%
\section{\texorpdfstring{First Challenge --- \emph{A Coin on a
String}}{First Challenge --- A Coin on a String}}\label{first-challenge-a-coin-on-a-string}}

\emph{{[}Achilles's table. A coin, a pencil, and a sheet of paper
crowded with his marks. The Tortoise sits beside him, on the same side
of the table.{]}}

\textbf{Achilles:} Over. Over again. Then I laid it heads up --- or was
that before the second turn? \emph{{[}He turns the coin over once more
and makes another mark.{]}} Tortoise, I have been turning this coin over
since supper, I have written down everything I did to it, and I can no
longer tell from my own notes which face ought to be showing.

\textbf{Tortoise:} Which face is showing?

\textbf{Achilles:} Tails. But is it \emph{right} that it's tails? That
is what the paper was for.

\textbf{Tortoise:} Before you cross out one more line --- may I? I keep
a length of string in my shell for exactly this kind of muddle. Laid
flat along the table, it holds things in the order they happen, which is
the one thing your paper won't do.

\textbf{Achilles:} By all means. My notes have not covered themselves in
glory tonight.

\textbf{Tortoise:} Your paper says what you did. It doesn't say in what
order, or to what. \emph{{[}From her shell she draws a length of string;
she lays it across the table in front of them, left to right, and pulls
it taut between her hands until the last of the slack is gone and it
lies dead straight. She pins down its two ends, and plucks it once: a
single clear note sounds and dies away, and the string settles into
stillness.{]}} Let us keep the record somewhere it can't get out of
order.

\textbf{Achilles:} A string. There's nothing on it.

\textbf{Tortoise:} Then so far nothing has been done. \emph{{[}She
writes ``heads up'' on a slip of paper and ties it to the string's left
end.{]}} The end tag: how the traveller starts. Nothing comes before it.
\emph{{[}She writes ``turn it over'' on another slip, folds it, and sets
it astride the string halfway along, like a saddle.{]}} And a tag: what
is done on the way. Set the coin as the end tag says, and walk it along.

\emph{{[}Achilles lays the coin heads up at the right-hand end of the
string, where his hand is.{]}}

\textbf{Tortoise:} Read it from the left. The end tag is where it
starts.

\textbf{Achilles:} \emph{{[}The coin sits heads up at the right-hand
end, where his hand is. He carries it round to the left end, under the
``heads up'' end tag, and slides it rightward along the table beside the
string --- the string and its one saddled tag holding still --- and
where the tag hangs he stops and turns the coin over in his fingers, so
that the face which was heads is now tails; then he carries it on to the
right-hand end.{]}} Tails up. Well, of course it's tails up.

\emph{{[}Under the door the draught comes again, harder than before, and
lays the hearth-flames flat over to one side. They burn that way a
moment, all of them leaning the one direction.{]}}

\textbf{Tortoise:} And if it had started the other way? \emph{{[}She
unties the end tag and ties on another: ``tails up''.{]}}

\textbf{Achilles:} \emph{{[}setting the coin tails up at the left end
and walking it along, turning it over at the tag{]}} Heads. So what
you're saying is, the string remembers what I did to the coin.

\textbf{Tortoise:} The string says what is \emph{to be} done, and in
what order, to whatever the end tag sends along. It would say so if you
never walked a coin past it at all.

\textbf{Achilles:} Then it can say more than I've done. \emph{{[}He
writes ``turn it over'' on a fresh slip, folds it, and sets it on the
string beside the first.{]}} And two of them?

\begin{center}\rule{0.5\linewidth}{0.5pt}\end{center}

\hypertarget{second-challenge-there-and-back}{%
\section{\texorpdfstring{Second Challenge --- \emph{There and
Back}}{Second Challenge --- There and Back}}\label{second-challenge-there-and-back}}

\textbf{Tortoise:} What do you expect?

\textbf{Achilles:} One tag turns the coin over. Two tags must do twice
as much.

\textbf{Tortoise:} Twice as much as turning over. And how will the coin
look, when it has been turned over twice as much?

\textbf{Achilles:} Very\ldots{} thoroughly tails. No --- I had better
walk it. \emph{{[}The Tortoise ties the ``heads up'' end tag on again.
He sets the coin heads up at the left end and walks it along: over at
the first tag, over at the second.{]}} Tails, then heads. It arrives as
it left.

\emph{{[}As he says it, the gust falls away outside. The flames right
themselves and stand up where they had stood before, unhurried, come
back to exactly where they began.{]}}

\textbf{Tortoise:} From that start.

\textbf{Achilles:} \emph{{[}The end tag is changed to ``tails up''; he
walks the coin again.{]}} Heads, then tails. As it left, again. And a
coin has only the two ways of starting, so I've walked them all.

\emph{{[}Two ``turn it over'' tags hang astride the string a little
apart, the end tag at the left. The Tortoise slides them along the
string toward each other until they meet, and lifts the touching pair
straight up and off. The string stays pinned and taut and the end tag
stays at its knot; only the pair is gone, and where they hung there is
now nothing to be done. The string is bare but for its end tag.{]}}

\textbf{Tortoise:} Walk it again.

\textbf{Achilles:} Past what? \emph{{[}He slides the coin from the left
end to the right without stopping.{]}} There was nothing to do. It
arrives as it left --- which is exactly what it did a moment ago, with
both tags on. \ldots Oh. I needn't have walked it. Two turn-overs side
by side are so much bare string: wherever I see the pair, I may lift it
off.

\textbf{Tortoise:} And wherever the string is bare, you may set a pair
on.

\textbf{Achilles:} Whyever would I want to?

\textbf{Tortoise:} Another evening.

\textbf{Achilles:} Then any tag, hung twice, comes off?

\begin{center}\rule{0.5\linewidth}{0.5pt}\end{center}

\hypertarget{third-challenge-which-tag-first}{%
\section{\texorpdfstring{Third Challenge --- \emph{Which Tag
First?}}{Third Challenge --- Which Tag First?}}\label{third-challenge-which-tag-first}}

\textbf{Tortoise:} \emph{{[}writing on a fresh slip{]}} ``Heads up,
whatever it was.'' \emph{{[}She writes a second the same, folds both,
and sets them on the string side by side.{]}} Any tag, you said.

\textbf{Achilles:} Twins. Off they come.

\textbf{Tortoise:} Walk it first. The end tag still says tails.

\textbf{Achilles:} \emph{{[}walking the coin{]}} Tails. At the first tag
I lay it heads up. At the second it is heads up already, and I leave it.
It arrives heads --- but along a bare string it would have arrived
tails! The pair can't come off.

\textbf{Tortoise:} \emph{{[}lifting one of the two off{]}} And now?

\textbf{Achilles:} \emph{{[}walking it again{]}} Heads, just the same.
So one of the twins comes off, and never both. Hung twice, this tag does
what it does once.

\textbf{Tortoise:} So not every tag.

\textbf{Achilles:} Very well: \emph{some} tags come off in pairs. But
the order along the string can't matter, at least. It didn't for the two
turn-overs.

\textbf{Tortoise:} The two turn-overs were the same tag. \emph{{[}She
sets a ``turn it over'' tag on the string to the left of the ``heads up,
whatever it was'' tag, and changes the end tag to ``heads up''.{]}} Walk
it.

\textbf{Achilles:} Heads. Over: tails. Heads up, whatever it was: heads.
It arrives heads.

\emph{{[}On the string, ``turn it over'' hangs at the left and ``heads
up, whatever it was'' at the right. The Tortoise lifts both clear of the
string and sets them down again with their places swapped --- ``heads
up, whatever it was'' now at the left, ``turn it over'' now at the
right. They are the very same two tags, and the coin's start is the very
same start; what has changed is only their order along the string.{]}}

\textbf{Tortoise:} Same tags. Walk it again.

\textbf{Achilles:} Heads. Heads up: it is already. Over: tails. It
arrives \emph{tails}. Same coin, same start, the same two tags, and the
other face --- so the order along the string is part of what the string
says.

\emph{{[}The tea has stood long enough to want reviving. Achilles pours
a little milk into his cup, where it blooms and greys the dark brew; the
Tortoise takes hers the other way about, the tea let down onto the milk
already in the cup, and the two settle to different colours.{]}}

\textbf{Tortoise:} Change the end tag, and walk it once more.

\textbf{Achilles:} \emph{{[}with ``tails up'' tied on{]}} Tails. Heads
up; over; tails again. Both starts arrive alike! Once the coin has
passed that ``heads up'' tag, nothing further along the string can tell
how it started, and nothing I could hang there would bring it back. The
turn-over tag never loses anything, and this one does. Is there a rule
for which tags lose things?

\emph{{[}He tops up both cups from the jug while he thinks, milk into
each until the two are of a colour. Whatever they had been --- the
strong cup and the weak --- they are the same pale tea now, and no
amount of looking into either would tell you which had been which.{]}}

\textbf{Tortoise:} Another evening. \emph{{[}She lifts both tags off.
The ``heads up, whatever it was'' tags go away into her shell; the two
``turn it over'' tags stay on the table.{]}} Every walk so far began
with a coin you could see. \emph{{[}She takes the coin and a die behind
her shell, where the die rattles once and settles; then, by the rule she
has set herself, she sets the coin down just short of the string's left
end, out of his sight, and lays her foot over it. It is one coin still,
with one true face --- she knows which, and will not say. She unties the
end tag; her foot does not lift again all evening.{]}} How does this one
start?

\begin{center}\rule{0.5\linewidth}{0.5pt}\end{center}

\hypertarget{fourth-challenge-the-covered-coin}{%
\section{\texorpdfstring{Fourth Challenge --- \emph{The Covered
Coin}}{Fourth Challenge --- The Covered Coin}}\label{fourth-challenge-the-covered-coin}}

\textbf{Achilles:} Heads up or tails up. It must be one of them; coins
don't lie on their edges under tortoises. What did the die say?

\textbf{Tortoise:} Heads, unless the die showed a six. I shan't tell you
what it showed.

\textbf{Achilles:} Then lift your foot.

\textbf{Tortoise:} That would tell us how the coin lies. I asked what we
may honestly tie to the string, now, with my foot where it is. ``Heads
up''?

\textbf{Achilles:} I couldn't swear to it. Nor to ``tails up''. Neither
tag is honest, and yet the coin \emph{is} one or the other. Let me take
a plainer case first: a coin tossed and covered. Evens.

\textbf{Tortoise:} By all means. \emph{{[}She cuts a strip of paper and
folds it in six.{]}} An end tag for a coin under a foot. The strip is
tied on by one end. From the knot out to where the shading changes is
heads' share; from there out to the free end is tails'. The shading only
marks the place: it is how far from the knot that counts, not which part
is dark.

\textbf{Achilles:} \emph{{[}shading three parts of the six from one end,
and tying the strip to the string's left end by its dark end{]}} From
the knot, three parts, and then the change: heads, three parts in six.
Tails, the other three. A tossed coin.

\textbf{Tortoise:} \emph{{[}setting one ``turn it over'' tag on the
string{]}} Walk it.

\textbf{Achilles:} The coin is under your foot.

\textbf{Tortoise:} Walk what you have.

\textbf{Achilles:} \emph{{[}untying the strip and carrying it along the
string as far as the tag{]}} And how does one turn over a wager?

\textbf{Tortoise:} If I turned the coin over where it lies, unseen, what
would become of heads' share?

\textbf{Achilles:} It would be tails' share, and tails' would be heads'.

\textbf{Tortoise:} Then turn the strip end for end, so that it will be
tied on by its other end.

\textbf{Achilles:} \emph{{[}The strip is in his hands at the tag,
untied; the end that had sat at the knot is the dark one, and from it
three parts run dark to the change of shade --- heads' share --- then
three parts white to the free end. He turns it end for end and holds its
white end to the string. What has moved is which end is held to the
string; what has not moved is the change of shade, still three parts in
from the held end, for on this strip it sits dead in the middle.{]}}
Tied on by this end: from the knot, three parts, and then the change ---
heads, three in six. \emph{{[}He turns it back, and round again: three
parts from either end, every time.{]}} It reads the same from either
end. Then for a coin nobody has seen, this tag does nothing! It needn't
wait for a twin; it may come off alone.

\textbf{Tortoise:} That is your tossed coin. Mine is under my foot.
\emph{{[}Achilles lays his strip aside. She folds a second strip in six,
shades five parts from one end, leaves the sixth white, and ties it to
the string's left end by its dark end.{]}}

\emph{{[}The fire has burned down while they worked, and sits now as one
low bed of embers: most of it still redly alight, and over in one corner
a patch already fallen to grey ash. It gives less light than it did, and
the room has drawn in a little closer around the table.{]}}

\textbf{Achilles:} From the knot, five parts, and then the change:
heads, five parts in six; tails, one. That's your die.

\textbf{Tortoise:} Walk it.

\emph{{[}The fire drops with a soft collapse of coals, and for a breath
the room holds still --- only the wind. In the lower light the woodcut
on the wall seems to come forward off its nail, nearer than it had hung
all evening.{]}}

\textbf{Achilles:} \emph{{[}This strip hangs by its dark end: from the
knot, five parts dark to the change --- heads' share, five in six ---
then one white part to the free end. He unties it, carries it along to
the tag, turns it end for end, and holds its white end to the string.
The strip is the same length it was; but the change of shade now sits
only one part in from the held end, where a moment ago it sat five parts
in, and the two shares have traded sizes.{]}} Tied on by this end: from
the knot, one part --- and then the change. Heads, one in six; tails,
five. It moved! \emph{{[}He looks up, the strip still in his hands, dark
end and white end, at the print on the wall --- Escher's
\textup{Day and Night}.{]}} Black birds flying one way, white birds the
other, over fields that are neither yet. And it is all one sheet of
paper.

\textbf{Tortoise:} And the tag?

\textbf{Achilles:} At work all along. My strip couldn't show it: turned
end for end, half-and-half is half-and-half.

\textbf{Tortoise:} One start was not enough to tell a tag from no tag.
\emph{{[}She sets the second ``turn it over'' tag on the string beside
the first.{]}} Go on.

\textbf{Achilles:} \emph{{[}carrying the strip on to the second tag and
turning it end for end again{]}} Tied on by its dark end again: from the
knot, five parts, and then the change --- heads, five in six, as it
began. Turned end for end, it is still one strip --- the same length
however it's held, and the two shares together always the whole of it.
\emph{{[}He carries it back and ties it to the string's left end by its
dark end.{]}}

\emph{{[}The Tortoise slides the two tags together and lifts them both
off. The strip hangs at the left end of a bare string, as she first tied
it.{]}}

\textbf{Achilles:} As before. The pair comes off for the strip as it did
for the coin, and the coin itself never stirred. So a start is either a
fact or a wager. There is nothing else it could be.

\emph{{[}The bed of embers has settled again while they talked: most of
the live red gathered at one end of the grate now, the barest coal at
the other, and more grey between them than there was. The wind takes up
and dies and takes up again beyond the glass.{]}}

\emph{{[}The Tortoise reaches into her shell and brings out a square of
tracing paper.{]}}

\begin{center}\rule{0.5\linewidth}{0.5pt}\end{center}

\hypertarget{fifth-challenge-the-arrow}{%
\section{\texorpdfstring{Fifth Challenge --- \emph{The
Arrow}}{Fifth Challenge --- The Arrow}}\label{fifth-challenge-the-arrow}}

\textbf{Tortoise:} Tracing paper this time, and a handful of
matchsticks. I choose tracing paper because a mark on it shows through
from behind, so we may turn it over and still follow the very same mark;
the matches are nothing but a small ruler --- set end to end, they say
how far a thing reaches.

\textbf{Achilles:} Matchsticks for a ruler. Very well --- show me a
start.

\textbf{Tortoise:} \emph{{[}She lays the square flat on the table. From
its near left-hand corner she lays three matchsticks end to end along
the bottom edge, and four up the side.{]}} Along the table, the way the
string runs: call that heads-ward. Away from us: tails-ward.
\emph{{[}With her pencil she draws an arrow from the corner to the spot
three matches along and four matches up.{]}} A start.

\textbf{Achilles:} A strip with a corner in it! Three parts heads-ward
and four tails-ward: seven parts in all. Your arrow is seven matches
long.

\textbf{Tortoise:} Measure it.

\textbf{Achilles:} \emph{{[}laying matchsticks end to end along the
arrow, from the corner to its tip{]}} Five. Not seven. But nine and
sixteen make twenty-five, and five fives are twenty-five: it's the
squares on the two reaches that make up the square on the arrow, not the
reaches themselves. Then these aren't shares --- three and four
overshoot the whole. As a wager it's nonsense.

\emph{{[}The teapot has sat cooling at the far corner of the table this
hour past. Achilles reaches for it now without rising --- straight
across the corner, the short way, a shorter reach than following the two
edges round --- finds it barely warm, and pours anyway.{]}}

\textbf{Tortoise:} As a wager, yes.

\textbf{Achilles:} And it's no fact either: it lies along neither edge.
What does it mean, to reach three heads-ward?

\textbf{Tortoise:} Another evening. Tonight, only what a tag does to it.
\emph{{[}The arrow reaches three matches heads-ward and four tails-ward,
and five along its own length. She clears the matches away and turns the
square face-down about the line running from the arrow's corner to the
far corner. The corner stays pinned where it was and the line she turns
it about stays put; what swings across is the rest of the square, so
that the heads-ward edge comes to lie where the tails-ward edge was, and
the reverse --- and through the paper the arrow shows, as long as
ever.{]}} Turned over.

\textbf{Achilles:} \emph{{[}laying matches again: along the bottom, up
the side, along the arrow{]}} Four heads-ward, three tails-ward. The
reaches have changed places, as the shares did. And along the arrow,
five still: turned over, it is still one arrow. \emph{{[}He clears the
matches and turns the square over again about the same line.{]}} Three
and four. Home.

\emph{{[}He unties the strip and ties the square to the string's left
end by the arrow's corner. He sets the two ``turn it over'' tags on the
string side by side. His hand goes to the knot, to untie the square for
its walk --- and stops. He slides the two tags together and lifts them
off.{]}}

\textbf{Achilles:} \ldots Oh. I needn't walk anything; I never need
have. The string and its tags never once asked what was travelling --- a
coin, a strip, an arrow. Each has its own way of being turned over, and
the tag doesn't care which. Two of them side by side come off for all
three.

\textbf{Tortoise:} They always did.

\textbf{Achilles:} But which ways may an arrow lie? And when it comes to
the right-hand end of the string --- how should I find out what had
arrived?

\textbf{Tortoise:} Another evening.

\textbf{Achilles:} You say that often.

\emph{{[}The wind has dropped to a long low sound at the glass, gone
from worrying at the house to something almost companionable. The fire
is a clear bed of coals now with hardly a flame left on it, warming
without lighting; the hour has got late without either of them marking
it.{]}}

\textbf{Tortoise:} It is a thrifty phrase. It puts nothing away for
good.

\textbf{Achilles:} \emph{{[}He gathers the loose tags into a small pile
and slides his crowded, useless sheet to the far edge of the table.{]}}
Then I'll leave off crossing things out. The string kept the better
record of the two.

\emph{{[}She unties the square and puts it away. The string lies across
the table, pinned and straight, with nothing on it, as it did at the
start. Her foot has not moved from the coin. Achilles plucks the string
once: the same note.{]}}

\textbf{Achilles:} But what does travel along it, then? A coin?

\textbf{Tortoise:} Not just a coin.

\emph{{[}The fire is banked low and steady, a clear bed of coals that
will keep till morning, and the wind has all but given over. Above the
hearth the woodcut has gone quiet in the dim --- the chequered fields
and the two towns settled back into the wall. Neither of them moves to
clear anything away just yet.{]}}

\hypertarget{chapter-1-not-just-a-bit}{%
\chapter{Chapter 1: Not Just a Bit}\label{chapter-1-not-just-a-bit}}

\emph{--- Overture for a Single String ---}

\emph{Part I: The Single Wire}

\emph{Escher: ``Day and Night'' --- white birds fly one way into the
night, black birds the other way into the day, and both grow out of the
fields between}

\begin{center}\rule{0.5\linewidth}{0.5pt}\end{center}

When something carries the answer to a yes-or-no question, what exactly
is it carrying?

A light switch is up or down. A coin on a table shows heads or tails. A
door is open or shut. Each of these settles one question, and it seems
that nothing could be simpler: there are two possibilities, the world
has picked one, and the thing in front of us is the record of the pick.
Every message ever sent, every photograph stored in a telephone, every
sum done by a machine is built from records of this kind, a great many
of them, and the theory of computing begins by taking one of them
seriously.

But put a hand over the coin before anybody has looked. There are still
two possibilities, and the world has still picked one, yet the honest
description of what is under the hand is no longer ``heads'' or
``tails''. It is something like ``heads, five times in six''. That
description is not a face of the coin. It is a pair of numbers, and it
has to be carried from place to place, and changed when the coin is
turned over, just as a face would be. So already there are two different
things that can travel down the same line: a fact, and a wager.

This chapter is about a third. It can be written in the same place as
the other two, it is changed by the same operations, and it is neither a
fact nor a wager. We shall define it carefully and say honestly how
little we yet know about it. Along the way we need a language for things
that travel and things that are done to them, and we shall build two:
one of columns and matrices, written on a blackboard, and one of lines
and boxes, drawn in a sketchbook. They will say the same things. The
interest of the book is in how differently they say them.

The smallest piece of all this is a single line with something
travelling along it. What travels? A coin, to begin with. Not just a
coin.

\begin{center}\rule{0.5\linewidth}{0.5pt}\end{center}

\hypertarget{a-coin-on-a-string}{%
\section{§1.1 A Coin on a String}\label{a-coin-on-a-string}}

Start with the plainest case. A \textbf{bit} is a thing with two
possible values, which we call \(0\) and \(1\). A coin lying on a table
is a bit: say heads is \(0\) and tails is \(1\). We want to describe two
matters: how a bit starts, and what is done to it.

\textbf{On the blackboard.} We could write the value of a bit as the
digit \(0\) or the digit \(1\) and be done. Instead we write it as a
column of two numbers, with a \(1\) in the place that belongs to the
value and a \(0\) in the other place:

\[\text{the state } 0 \;=\; \begin{pmatrix}1\\0\end{pmatrix}, \qquad\qquad \text{the state } 1 \;=\; \begin{pmatrix}0\\1\end{pmatrix}.\]

The convention, fixed for the whole book, is that \textbf{the first
(top) entry belongs to \(0\) and the second to \(1\)}. Read a column as
two answers, to the questions ``is it \(0\)?'' and ``is it \(1\)?''; the
state \(0\) answers \emph{yes, no}.

Something done to a bit is called a \textbf{process}, and on the
blackboard a process is a \(2\times2\) matrix, a square of four numbers.
A matrix acts on a column by the usual rule,

\[\begin{pmatrix}a&b\\c&d\end{pmatrix}\begin{pmatrix}x\\y\end{pmatrix} \;=\; \begin{pmatrix}ax+by\\cx+dy\end{pmatrix}.\]

Try the rule on our two columns. With \(x=1,y=0\) the result is
\(\begin{pmatrix}a\\c\end{pmatrix}\), the first column of the matrix;
with \(x=0,y=1\) it is \(\begin{pmatrix}b\\d\end{pmatrix}\), the second.
So a matrix wears its behaviour on its face: \textbf{its first column is
what it does to the state \(0\), and its second column is what it does
to the state \(1\).} We shall use this remark more than once.

The process we care about most in this chapter turns the coin over. It
is called \textbf{NOT}, and we write \(X\) for its matrix. It must send
the state \(0\) to the state \(1\) and the state \(1\) to the state
\(0\), so by the remark just made its columns are forced:

\[X \;=\; \begin{pmatrix}0&1\\1&0\end{pmatrix}.\]

Check it by the rule:

\[X\begin{pmatrix}1\\0\end{pmatrix}=\begin{pmatrix}0\cdot1+1\cdot0\\1\cdot1+0\cdot0\end{pmatrix}=\begin{pmatrix}0\\1\end{pmatrix},
\qquad
X\begin{pmatrix}0\\1\end{pmatrix}=\begin{pmatrix}0\cdot0+1\cdot1\\1\cdot0+0\cdot1\end{pmatrix}=\begin{pmatrix}1\\0\end{pmatrix}.\]

\textbf{In the sketchbook.} We now draw the same things, and since these
are the first pictures of the book, each element is described before it
appears.

A \textbf{wire} is a horizontal line. It stands for the system: the
thing that has a value, here our one coin. A wire is not a moment and
not a place; it is the coin's whole career, drawn from left to right.

A \textbf{box} is a rectangle sitting on the wire, with the wire
entering on its left and leaving on its right, and a name written
inside. It stands for a process: something done to the system. What
enters on the left is what the process is given; what leaves on the
right is what it hands on.

A \textbf{state} is a box with \emph{no wire coming in}: nothing enters
it, and one wire leaves it on the right. It stands for a way of
starting. To tell states from processes at a glance we draw a state with
rounded ends. The two states we have so far are labelled \(0\) and
\(1\).

\[\text{a wire:}\;\;
\begin{tikzpicture}[sd, baseline=-0.55ex]
  \draw[wire] (0,0) -- (3,0);
\end{tikzpicture}
\qquad
\text{a box:}\;\;
\begin{tikzpicture}[sd, baseline=-0.55ex]
  \draw[wire] (0,0) -- (4,0);
  \node[sdbox, text height=1.5ex, text depth=0.5ex] at (2,0) {$A$};
\end{tikzpicture}
\qquad
\text{a state:}\;\;
\begin{tikzpicture}[sd, baseline=-0.55ex]
  \draw[wire] (0,0) -- (3,0);
  \node[sdstate, text height=1.5ex, text depth=0.5ex] at (0,0) {$0$};
\end{tikzpicture}\]

One more element, and it is not a shape but a way of placing shapes.
Several things may sit \textbf{in a row} on the same wire, and a row is
read \textbf{from left to right: first this, then that.} A state at the
left end followed by a box means ``start like this, then do that''.
Whatever is further to the left happens earlier.

With these four elements the two blackboard calculations above become
two picture equations. Start in the state \(0\), then apply NOT: the
result is the state \(1\). And the other way about:

\[\begin{tikzpicture}[sd, baseline=-0.55ex]
  \draw[wire] (0,0) -- (6,0);
  \node[sdstate, text height=1.5ex, text depth=0.5ex] at (0,0) {$0$};
  \node[sdbox, text height=1.5ex, text depth=0.5ex] at (3.2,0) {$\mathrm{NOT}$};
\end{tikzpicture}
\;=\;
\begin{tikzpicture}[sd, baseline=-0.55ex]
  \draw[wire] (0,0) -- (3,0);
  \node[sdstate, text height=1.5ex, text depth=0.5ex] at (0,0) {$1$};
\end{tikzpicture}
\qquad\qquad
\begin{tikzpicture}[sd, baseline=-0.55ex]
  \draw[wire] (0,0) -- (6,0);
  \node[sdstate, text height=1.5ex, text depth=0.5ex] at (0,0) {$1$};
  \node[sdbox, text height=1.5ex, text depth=0.5ex] at (3.2,0) {$\mathrm{NOT}$};
\end{tikzpicture}
\;=\;
\begin{tikzpicture}[sd, baseline=-0.55ex]
  \draw[wire] (0,0) -- (3,0);
  \node[sdstate, text height=1.5ex, text depth=0.5ex] at (0,0) {$0$};
\end{tikzpicture}\]

We shall need to refer to these two equations by name. The left one is
\textbf{\emph{NOT on 0}}; the right one is \textbf{\emph{NOT on 1}}.
Each side of each equation is a state: no wire comes in, one wire goes
out. An equals sign between two pictures means that they stand for the
same column or the same matrix, entry for entry. In this chapter every
picture equation is exact in that sense. See \texttt{Code/Ch1/Flip.py}.

\textbf{\emph{Why this definition.}} Why write a bit as a column of two
numbers, when one digit would do? For two reasons, and both are
investments. First, once a start is a column, a process can be a matrix,
and everything that matrices know how to do --- above all, to be
multiplied --- becomes available. Second, a column has room in it. The
two columns above use only the numbers \(0\) and \(1\); the places where
they sit could hold other numbers. In §1.4 and §1.5 we shall change what
is written in the column \emph{without changing a single matrix or a
single picture}, and that is the whole plot of the chapter.

\begin{quote}
\textbf{Feynman's Notebook.}

\emph{What a bit is made of.} No machine has ever held a \(0\) or a
\(1\). What a machine holds is some physical thing that can be put into
either of two conditions and found, later, reliably in the one it was
put in: a voltage that is high or low, a patch of magnetic material
lying one way or the other, a hole punched or not punched in a card. The
\(0\) and the \(1\) are names we give to the two conditions. The coin on
the table is as good a bit as any of these, only slower.

This is why the theory can afford to be abstract. Nielsen and Chuang, in
the textbook on which the classical half of this book leans
(\emph{Quantum Computation and Quantum Information}, §1.2), introduce
the bit as the fundamental concept of classical computation and
classical information, and then describe its quantum counterpart
deliberately as a mathematical object, so that what is proved does not
depend on any particular way of building one. Our wire is drawn in the
same spirit. It does not say what the system is made of. It says only
that there is one.
\end{quote}

\begin{center}\rule{0.5\linewidth}{0.5pt}\end{center}

\hypertarget{there-and-back}{%
\section{§1.2 There and Back}\label{there-and-back}}

Turn the coin over, and then turn it over again. Anyone can say what
happens. The point of this section is to say it in our two languages and
to notice how much, or how little, each has really said.

\textbf{On the blackboard.} Start in the state \(0\) and apply \(X\)
twice, one step at a time, using the two calculations of §1.1:

\[\begin{pmatrix}1\\0\end{pmatrix}\;\xrightarrow{\;X\;}\;\begin{pmatrix}0\\1\end{pmatrix}\;\xrightarrow{\;X\;}\;\begin{pmatrix}1\\0\end{pmatrix},
\qquad\qquad
\begin{pmatrix}0\\1\end{pmatrix}\;\xrightarrow{\;X\;}\;\begin{pmatrix}1\\0\end{pmatrix}\;\xrightarrow{\;X\;}\;\begin{pmatrix}0\\1\end{pmatrix}.\]

Each start comes home.

\textbf{In the sketchbook.} Put two NOT boxes in a row after the state
\(0\). Reading from the left, the state \(0\) followed by the first NOT
is, by \emph{NOT on 0}, the state \(1\); and the state \(1\) followed by
the remaining NOT is, by \emph{NOT on 1}, the state \(0\):

\[\begin{tikzpicture}[sd, baseline=-0.55ex]
  \draw[wire] (0,0) -- (9.4,0);
  \node[sdstate, text height=1.5ex, text depth=0.5ex] at (0,0) {$0$};
  \node[sdbox, text height=1.5ex, text depth=0.5ex] at (3.2,0) {$\mathrm{NOT}$};
  \node[sdbox, text height=1.5ex, text depth=0.5ex] at (6.6,0) {$\mathrm{NOT}$};
\end{tikzpicture}
\;=\;
\begin{tikzpicture}[sd, baseline=-0.55ex]
  \draw[wire] (0,0) -- (6,0);
  \node[sdstate, text height=1.5ex, text depth=0.5ex] at (0,0) {$1$};
  \node[sdbox, text height=1.5ex, text depth=0.5ex] at (3.2,0) {$\mathrm{NOT}$};
\end{tikzpicture}
\;=\;
\begin{tikzpicture}[sd, baseline=-0.55ex]
  \draw[wire] (0,0) -- (3,0);
  \node[sdstate, text height=1.5ex, text depth=0.5ex] at (0,0) {$0$};
\end{tikzpicture}\]

The same two steps in the other order bring the state \(1\) home.

So far this is an \emph{example}: two particular starts, walked through
two particular boxes. It says that for a coin we can see, two turns are
as good as none. It would be more useful to have a statement about the
two boxes themselves, with no start mentioned --- a statement one could
apply in the middle of a long row without asking what had come in from
the left. Such a statement needs a picture of ``no boxes at all''. That
is the \textbf{bare wire}: a wire with nothing on it, standing for the
process of doing nothing. On the blackboard doing nothing is the
\textbf{identity matrix}

\[I \;=\; \begin{pmatrix}1&0\\0&1\end{pmatrix},\]

which by the rule of §1.1 returns every column unchanged.

\begin{quote}
\textbf{Theorem 1.1 (NOT twice is nothing).} \(X\,X = I\), exactly. In
pictures: two NOT boxes in a row are the bare wire.
\end{quote}

\[\begin{tikzpicture}[sd, baseline=-0.55ex]
  \draw[wire] (0,0) -- (8.4,0);
  \node[sdbox, text height=1.5ex, text depth=0.5ex] at (2.5,0) {$\mathrm{NOT}$};
  \node[sdbox, text height=1.5ex, text depth=0.5ex] at (5.9,0) {$\mathrm{NOT}$};
\end{tikzpicture}
\;=\;
\begin{tikzpicture}[sd, baseline=-0.55ex]
  \draw[wire] (0,0) -- (3,0);
\end{tikzpicture}\]

\emph{Proof idea.} A matrix is known once its two columns are known, and
its columns are what it does to the states \(0\) and \(1\). The example
has just shown that ``NOT, then NOT'' sends each of those two states
where doing nothing sends it. So the two processes have the same
columns.

\emph{Proof sketch.} The matrix of ``NOT, then NOT'' is the product
\(XX\) (why a row of boxes is a product is the business of §1.3). Its
first column is
\(X\bigl(X\begin{pmatrix}1\\0\end{pmatrix}\bigr)=\begin{pmatrix}1\\0\end{pmatrix}\)
and its second is
\(X\bigl(X\begin{pmatrix}0\\1\end{pmatrix}\bigr)=\begin{pmatrix}0\\1\end{pmatrix}\);
those are the columns of \(I\). \(\;\blacksquare\) See
\texttt{Code/Ch1/TwiceOver.py}.

This is the headline theorem of the chapter, and §1.7 proves it twice,
in full. It may seem a great deal of ceremony for a coin turned over
twice. The reason for the ceremony is that the theorem is about the
\emph{matrix}, not about the coin. Before the chapter is over the wire
will be carrying things that are not coins, and the theorem will hold
for them without a word being changed.

\begin{center}\rule{0.5\linewidth}{0.5pt}\end{center}

\hypertarget{boxes-in-a-row}{%
\section{§1.3 Boxes in a Row}\label{boxes-in-a-row}}

Two copies of the same box in a row cannot show us everything about
rows. We need a second kind of box.

\textbf{On the blackboard.} How many processes are there that take a bit
and return a bit? A process must say where \(0\) goes and where \(1\)
goes, and each has two places to go, so there are exactly four. Writing
each as the matrix whose columns are the two results:

\[I=\begin{pmatrix}1&0\\0&1\end{pmatrix},\qquad
X=\begin{pmatrix}0&1\\1&0\end{pmatrix},\qquad
S_0=\begin{pmatrix}1&1\\0&0\end{pmatrix},\qquad
S_1=\begin{pmatrix}0&0\\1&1\end{pmatrix}.\]

The first does nothing and the second is NOT. The third, \textbf{SET-0},
sends both states to the state \(0\) --- ``heads up, whatever it was''
--- and the fourth, \textbf{SET-1}, sends both to the state \(1\).

Now do one process and then another. If a column \(v\) is first given to
a matrix \(A\), the result is \(Av\); if that is then given to \(B\),
the result is \(B(Av)\). Matrix multiplication is defined precisely so
that this equals \((BA)v\): the product \(BA\) is the single matrix that
does the work of ``\(A\), then \(B\)''. Notice where the letters have
ended up.

\textbf{In the sketchbook.} Two boxes in a row, \(A\) on the left and
\(B\) on the right, are read as always from left to right: first \(A\),
then \(B\). Draw an outline around the two of them and what is inside
has one wire in and one wire out: \textbf{a row of boxes is itself a
box}. Its matrix is \(BA\).

\[\begin{tikzpicture}[sd, baseline=-0.55ex]
  \draw[wire] (0,0) -- (7,0);
  \node[sdbox, text height=1.5ex, text depth=0.5ex] at (2,0) {$A$};
  \node[sdbox, text height=1.5ex, text depth=0.5ex] at (5,0) {$B$};
\end{tikzpicture}
\;=\;
\begin{tikzpicture}[sd, baseline=-0.55ex]
  \draw[wire] (0,0) -- (4.4,0);
  \node[sdbox, text height=1.5ex, text depth=0.5ex] at (2.2,0) {$BA$};
\end{tikzpicture}\]

This is the one convention in the book most worth saying loudly, so here
it is, loudly. \textbf{Pictures are read from left to right. Matrix
products are read from right to left. The row ``\(A\), then \(B\)'' is
the product \(BA\).} The reason is nothing deeper than the habit of
writing a matrix to the \emph{left} of the column it acts on, so that
the matrix nearest the column acts first. Neither convention will
change; the reader who remembers that they run opposite ways will be
spared a great many sign errors and transposed answers later in the
book.

With NOT alone the convention was invisible, since \(XX\) reads the same
in both directions. With two \emph{different} boxes it shows. Take NOT
and SET-0, in both orders:

\[\text{NOT, then SET-0:}\quad S_0X=\begin{pmatrix}1&1\\0&0\end{pmatrix}\begin{pmatrix}0&1\\1&0\end{pmatrix}=\begin{pmatrix}1&1\\0&0\end{pmatrix}=S_0,\]

\[\text{SET-0, then NOT:}\quad XS_0=\begin{pmatrix}0&1\\1&0\end{pmatrix}\begin{pmatrix}1&1\\0&0\end{pmatrix}=\begin{pmatrix}0&0\\1&1\end{pmatrix}=S_1.\]

In words: turn the coin over and then lay it heads up, and it is heads
up; lay it heads up and then turn it over, and it is tails up. The same
two boxes, in the two orders, are two different processes:

\[\begin{tikzpicture}[sd, baseline=-0.55ex]
  \draw[wire] (0,0) -- (9.4,0);
  \node[sdbox, text height=1.5ex, text depth=0.5ex] at (2.5,0) {$\mathrm{NOT}$};
  \node[sdbox, text height=1.5ex, text depth=0.5ex] at (6.4,0) {$\mathrm{SET}_0$};
\end{tikzpicture}
\;=\;
\begin{tikzpicture}[sd, baseline=-0.55ex]
  \draw[wire] (0,0) -- (5,0);
  \node[sdbox, text height=1.5ex, text depth=0.5ex] at (2.5,0) {$\mathrm{SET}_0$};
\end{tikzpicture}\]

\[\begin{tikzpicture}[sd, baseline=-0.55ex]
  \draw[wire] (0,0) -- (9.4,0);
  \node[sdbox, text height=1.5ex, text depth=0.5ex] at (2.9,0) {$\mathrm{SET}_0$};
  \node[sdbox, text height=1.5ex, text depth=0.5ex] at (6.8,0) {$\mathrm{NOT}$};
\end{tikzpicture}
\;=\;
\begin{tikzpicture}[sd, baseline=-0.55ex]
  \draw[wire] (0,0) -- (5,0);
  \node[sdbox, text height=1.5ex, text depth=0.5ex] at (2.5,0) {$\mathrm{SET}_1$};
\end{tikzpicture}\]

\textbf{Order matters.} See \texttt{Code/Ch1/OrderMatters.py}, which
also checks the left-to-right, right-to-left convention on a pair of
matrices with no helpful symmetries at all.

The bare wire of §1.2 now has a second description: it is the
\emph{empty row}, the row with no boxes on it, and its matrix is \(I\).
Putting a bare stretch of wire before or after a box changes nothing,
which on the blackboard reads \(AI=A=IA\).

Two further lessons can be read from SET-0. First, Theorem 1.1 is a fact
about NOT and not a fact about rows: two SET-0 boxes in a row are
\emph{not} the bare wire, since \(S_0S_0=S_0\). Doing it twice is the
same as doing it once. Second, nothing placed after SET-0 can repair it.
For any matrix \(B\) whatever,

\[B\,S_0 \;=\; \begin{pmatrix}b_{11}&b_{12}\\b_{21}&b_{22}\end{pmatrix}\begin{pmatrix}1&1\\0&0\end{pmatrix} \;=\; \begin{pmatrix}b_{11}&b_{11}\\b_{21}&b_{21}\end{pmatrix},\]

a matrix with two equal columns, and \(I\) does not have two equal
columns. Once the coin has been laid heads up, nothing further along the
wire can tell how it started. We record this as a fact about this one
matrix. What distinguishes the boxes that can be undone from those that
cannot is a question for Chapter 3.

\begin{quote}
\textbf{Dirac's Blackboard.}

\emph{All sixteen rows of two.} Four boxes, taken two at a time with
order respected, make sixteen rows. Here are all of them. The entry in
the row labelled \(A\) and the column labelled \(B\) is the process
``\(A\), then \(B\)'', which is the matrix product \(BA\) --- the
\emph{column} label is written on the \emph{left} of the product.

\[\begin{array}{c|cccc}
\text{first}\backslash\text{then} & I & X & S_0 & S_1\\ \hline
I & I & X & S_0 & S_1\\
X & X & I & S_0 & S_1\\
S_0 & S_0 & S_1 & S_0 & S_1\\
S_1 & S_1 & S_0 & S_0 & S_1
\end{array}\]

Three things to read off. Every entry is again one of the four, as it
must be, since a bit-process followed by a bit-process is a bit-process.
The table is not symmetric about its diagonal: compare the entry for
\(X\), then \(S_0\) (which is \(S_0\)) with the entry for \(S_0\), then
\(X\) (which is \(S_1\)). And the symbol \(I\) appears on the diagonal
exactly twice: only doing nothing, and NOT, are undone by a second copy
of themselves.
\end{quote}

\textbf{\emph{Why this definition.}} Why let a row of boxes \emph{be} a
box, rather than keep rows and boxes as two kinds of thing? Because then
every rule about boxes applies to rows for nothing, and an equation such
as Theorem 1.1, which says that a certain row equals a certain simpler
row, may be used anywhere inside a longer one.

What we have drawn in this section are the first \textbf{circuits} of
the book: boxes on a wire, doing to a bit the things a classical
computer does. They are, admittedly, the circuits of a computer with a
single bit in it. Real circuits have many wires side by side and boxes
that join them, and those need one new idea --- how to draw, and how to
calculate, two systems at once --- which belongs to Chapter 4.

\begin{quote}
\textbf{Penrose's Sketchbook.}

\emph{Where the pictures come from, and what they do not say.} Drawing a
calculation is an old idea. In 1971 Roger Penrose published
``Applications of negative dimensional tensors''; as Peter Selinger's
survey of graphical languages tells it, Penrose proposed drawing each
map as a box, with a wire for each of its indices, and joining the wires
of indices that are paired. In 2004 Samson Abramsky and Bob Coecke, of
the Oxford University Computing Laboratory, set out to replace ad hoc
calculations with matrices and normalizing constants by conceptual
definitions and proofs; of their networks of boxes they observe (in
their §4) that the quantum information seems to flow along the line. Our
reading of a wire as a system and a box as a process descends from that
work.

Two cautions, to carry through the book. A picture does not say what its
wire carries; the next two sections lean on that. And at this stage a
picture has no rules of its own: every equation between pictures in this
chapter is a matrix equation underneath, and is believed for that reason
and no other.
\end{quote}

\begin{center}\rule{0.5\linewidth}{0.5pt}\end{center}

\hypertarget{the-covered-coin}{%
\section{§1.4 The Covered Coin}\label{the-covered-coin}}

Now roll a die out of sight. If it shows a six, lay the coin tails up;
otherwise lay it heads up. Cover the coin. How does it start?

Neither of our two state boxes is an honest answer. The coin under the
hand is heads or it is tails, but \emph{we} can say only that it is
heads in five cases out of six.

\textbf{On the blackboard.} A \textbf{probabilistic bit} is described by
a column

\[\begin{pmatrix}p\\1-p\end{pmatrix}, \qquad 0\le p\le 1,\]

in which the first entry is the probability that the value is \(0\) and
the second the probability that it is \(1\). Such a column --- entries
that are real and not negative, and that \textbf{sum to \(1\)} --- is
called a \textbf{stochastic vector}. The covered coin is
\(\begin{pmatrix}5/6\\1/6\end{pmatrix}\). The two columns of §1.1 are
stochastic vectors too, the ones in which nothing is in doubt.

What should a process do to such a column? Suppose the process is the
matrix \(A\). With probability \(p\) the coin is really in the state
\(0\), and then the result is the first column of \(A\); with
probability \(1-p\) it is in the state \(1\), and the result is the
second column. The honest description of the result is the two mixed in
those proportions:

\[p\cdot\bigl(\text{first column of }A\bigr)+(1-p)\cdot\bigl(\text{second column of }A\bigr) \;=\; A\begin{pmatrix}p\\1-p\end{pmatrix}.\]

The right-hand side is the rule of §1.1 read backwards. So \textbf{the
same matrices act, in the same way}; nothing needs to be redefined. For
NOT,

\[X\begin{pmatrix}p\\1-p\end{pmatrix}=\begin{pmatrix}1-p\\p\end{pmatrix}, \qquad\text{for instance}\qquad X\begin{pmatrix}5/6\\1/6\end{pmatrix}=\begin{pmatrix}1/6\\5/6\end{pmatrix}:\]

the two shares change places. Nobody has touched the coin's description
by hand; the matrix did it.

Does a process always return a legal description? Add the two entries of
\(Av\) using the rule of §1.1: the sum is \((a+c)\,x+(b+d)\,y\). If each
\emph{column} of \(A\) sums to \(1\), that is \(x+y\), and the total is
kept. All four matrices of §1.3 have entries that are not negative and
columns that sum to \(1\) (such matrices are called \textbf{stochastic
matrices}), and that is exactly why they turn stochastic vectors into
stochastic vectors.

And two NOTs? There is nothing to prove. Theorem 1.1 says \(XX=I\), and
\(I\) returns every column unchanged, whatever is written in it:

\[X\Bigl(X\begin{pmatrix}5/6\\1/6\end{pmatrix}\Bigr)=X\begin{pmatrix}1/6\\5/6\end{pmatrix}=\begin{pmatrix}5/6\\1/6\end{pmatrix}.\]

\textbf{In the sketchbook.} There is no new element. The wire is the
same wire and the boxes are the same boxes; only the label inside the
state box has changed. We write the two entries of the column in it,
first entry first:

\[\begin{tikzpicture}[sd, baseline=-0.55ex]
  \draw[wire] (0,0) -- (8.6,0);
  \node[sdstate, inner xsep=0.6em, text height=1.5ex, text depth=0.5ex] at (0,0) {$5/6,\;1/6$};
  \node[sdbox, text height=1.5ex, text depth=0.5ex] at (5.6,0) {$\mathrm{NOT}$};
\end{tikzpicture}
\;=\;
\begin{tikzpicture}[sd, baseline=-0.55ex]
  \draw[wire] (0,0) -- (4.8,0);
  \node[sdstate, inner xsep=0.6em, text height=1.5ex, text depth=0.5ex] at (0,0) {$1/6,\;5/6$};
\end{tikzpicture}\]

Every picture of §1.1--§1.3 stands exactly as it was drawn. None of them
ever said what the wire carries.

\hypertarget{a-trap}{%
\subsection{A trap}\label{a-trap}}

Here is a tempting experiment that goes wrong. Take the fairest covered
coin, heads and tails in equal shares, and apply NOT:

\[X\begin{pmatrix}1/2\\1/2\end{pmatrix}=\begin{pmatrix}1/2\\1/2\end{pmatrix}.\]

The column has not moved. One might conclude that for a coin nobody has
seen, NOT does nothing, and might as well be taken off the wire \emph{by
itself}. In pictures the tempting inference is from the first line below
to the denial of the second:

\[\begin{tikzpicture}[sd, baseline=-0.55ex]
  \draw[wire] (0,0) -- (8.6,0);
  \node[sdstate, inner xsep=0.6em, text height=1.5ex, text depth=0.5ex] at (0,0) {$1/2,\;1/2$};
  \node[sdbox, text height=1.5ex, text depth=0.5ex] at (5.6,0) {$\mathrm{NOT}$};
\end{tikzpicture}
\;=\;
\begin{tikzpicture}[sd, baseline=-0.55ex]
  \draw[wire] (0,0) -- (4.8,0);
  \node[sdstate, inner xsep=0.6em, text height=1.5ex, text depth=0.5ex] at (0,0) {$1/2,\;1/2$};
\end{tikzpicture}\]

\[\begin{tikzpicture}[sd, baseline=-0.55ex]
  \draw[wire] (0,0) -- (5,0);
  \node[sdbox, text height=1.5ex, text depth=0.5ex] at (2.5,0) {$\mathrm{NOT}$};
\end{tikzpicture}
\;\neq\;
\begin{tikzpicture}[sd, baseline=-0.55ex]
  \draw[wire] (0,0) -- (3,0);
\end{tikzpicture}\]

Both lines are true. NOT leaves the fair coin's description alone and is
nevertheless not the bare wire, as the five-in-six coin showed a moment
ago. The moral is worth a sentence of its own: \textbf{one state is not
enough to tell two boxes apart.} Agreement on a single start, however
natural a start it is, proves nothing about the boxes. How many starts
\emph{are} enough is a question we shall have to answer in §1.7, since
the sketchbook proof of Theorem 1.1 depends on it. See
\texttt{Code/Ch1/HiddenCoin.py}.

\textbf{\emph{Why this definition.}} Why describe ignorance by a column
of two numbers summing to \(1\), rather than by a single number \(p\)?
Because the column was already there. In §1.1 we chose to write a bit in
two places precisely so that something other than \(0\) and \(1\) could
later be written in them, and so that the matrices would not have to be
told. A single number \(p\) would have needed a new rule for each
process; the column needs none.

\begin{center}\rule{0.5\linewidth}{0.5pt}\end{center}

\hypertarget{the-arrow}{%
\section{§1.5 The Arrow}\label{the-arrow}}

So a start is either a fact or a wager? It is natural to think that
there is nothing else it could be. The columns say otherwise. A column
has two places, and we have filled them in two ways: with a \(1\) and a
\(0\); and with two shares of a whole. Here is a third way.

\textbf{On the blackboard.} A \textbf{qubit} is a system whose state is
a column of two \emph{complex} numbers

\[\begin{pmatrix}\alpha\\\beta\end{pmatrix}, \qquad\text{with}\qquad |\alpha|^2+|\beta|^2=1 .\]

Here \(|z|\) is the modulus of a complex number: for \(z=a+bi\) with
\(a,b\) real, \(|z|^2=a^2+b^2\). The set of all columns of two complex
numbers is written \(\mathbb{C}^2\), so a qubit state is a column in
\(\mathbb{C}^2\) whose \textbf{squared moduli sum to \(1\)}. The entries
need not be shares of anything: being complex numbers, they may be
negative, or not real at all.

Put the two rules side by side on one example:

\[\begin{pmatrix}3/5\\4/5\end{pmatrix}: \qquad \Bigl(\tfrac35\Bigr)^2+\Bigl(\tfrac45\Bigr)^2=\tfrac{9}{25}+\tfrac{16}{25}=1, \qquad\text{but}\qquad \tfrac35+\tfrac45=\tfrac75 .\]

This column is a perfectly good qubit state and is \emph{not} a
stochastic vector. Going the other way, the covered coin
\(\begin{pmatrix}5/6\\1/6\end{pmatrix}\) is a stochastic vector and not
a qubit state, since \(\tfrac{25}{36}+\tfrac1{36}=\tfrac{26}{36}\). The
two rules are different rules, and the two kinds of column are different
kinds of thing. Only the states \(0\) and \(1\) are legal under both ---
and as bits besides. Another legal qubit state, to show an entry that is
not real, is \(\begin{pmatrix}3/5\\4i/5\end{pmatrix}\), since
\(|4i/5|^2=\tfrac{16}{25}\).

When both entries are real, there is a picture of the rule that everyone
knows. Draw an arrow in the plane, starting at the origin, reaching
\(\alpha\) to the right and \(\beta\) upward. By Pythagoras its length
squared is \(\alpha^2+\beta^2\). A real qubit state is an arrow
\textbf{of length \(1\)}; our example reaches three fifths one way and
four fifths the other, and the old triangle with sides three, four and
five does the rest. The dialogue draws its arrow this way, on paper. It
is an honest picture of a slice of the truth: the full space
\(\mathbb{C}^2\) has complex entries and does not fit on a sheet of
paper.

The matrices act as before, by the rule of §1.1, which never asked what
kind of numbers it was multiplying. NOT exchanges the entries,

\[X\begin{pmatrix}\alpha\\\beta\end{pmatrix}=\begin{pmatrix}\beta\\\alpha\end{pmatrix},
\qquad\text{for instance}\qquad
X\begin{pmatrix}3/5\\4/5\end{pmatrix}=\begin{pmatrix}4/5\\3/5\end{pmatrix},\]

and since \(|\beta|^2+|\alpha|^2=|\alpha|^2+|\beta|^2\), NOT turns a
qubit state into a qubit state. Two NOTs bring it home, by Theorem 1.1,
with nothing new to prove. See \texttt{Code/Ch1/Arrow.py}.

Not every box of §1.3 is so well behaved with the new passenger. SET-0
keeps sums but not sums of squares:
\(S_0\begin{pmatrix}3/5\\4/5\end{pmatrix}=\begin{pmatrix}7/5\\0\end{pmatrix}\),
which is no qubit state. Which boxes may sit on a wire that carries a
qubit is, again, Chapter 3's question.

\textbf{In the sketchbook.} Once more, no new element. The wire now
stands for \(\mathbb{C}^2\); a general state box is labelled \(\psi\),
the Greek letter psi, standing for a column
\(\begin{pmatrix}\alpha\\\beta\end{pmatrix}\); and NOT acts on it as the
blackboard says:

\[\begin{tikzpicture}[sd, baseline=-0.55ex]
  \draw[wire] (0,0) -- (6,0);
  \node[sdstate, text height=1.5ex, text depth=0.5ex] at (0,0) {$\psi$};
  \node[sdbox, text height=1.5ex, text depth=0.5ex] at (3.2,0) {$\mathrm{NOT}$};
\end{tikzpicture}
\;=\;
\begin{tikzpicture}[sd, baseline=-0.55ex]
  \draw[wire] (0,0) -- (3.2,0);
  \node[sdstate, text height=1.5ex, text depth=0.5ex] at (0,0) {$X\psi$};
\end{tikzpicture}\]

The three passengers the wire has carried in this chapter, in one table:

\begin{longtable}[]{@{}
  >{\raggedright\arraybackslash}p{(\columnwidth - 4\tabcolsep) * \real{0.3333}}
  >{\raggedright\arraybackslash}p{(\columnwidth - 4\tabcolsep) * \real{0.3333}}
  >{\raggedright\arraybackslash}p{(\columnwidth - 4\tabcolsep) * \real{0.3333}}@{}}
\toprule\noalign{}
\begin{minipage}[b]{\linewidth}\raggedright
passenger
\end{minipage} & \begin{minipage}[b]{\linewidth}\raggedright
its column
\end{minipage} & \begin{minipage}[b]{\linewidth}\raggedright
the rule
\end{minipage} \\
\midrule\noalign{}
\endhead
\bottomrule\noalign{}
\endlastfoot
a bit & \(\begin{pmatrix}1\\0\end{pmatrix}\) or
\(\begin{pmatrix}0\\1\end{pmatrix}\) & one entry is \(1\), the other
\(0\) \\
a probabilistic bit & \(\begin{pmatrix}p\\1-p\end{pmatrix}\) & real, not
negative, \textbf{entries sum to \(1\)} \\
a qubit & \(\begin{pmatrix}\alpha\\\beta\end{pmatrix}\) & complex,
\textbf{squared moduli sum to \(1\)} \\
\end{longtable}

\textbf{\emph{Why this definition.}} Here we must be candid: this
chapter does not say. We have \emph{defined} the qubit's rule; we have
not explained it. Why squares, and not the entries themselves? What do
\(\alpha\) and \(\beta\) have to do with what anyone would see who
looked? If they are not shares, what are they? Those are the questions
of Chapter 2, and the possibility of negative entries, which we have
mentioned and then carefully not used, is where that chapter begins.
What this chapter can say is that the definition is \emph{not} a
relabelling of the covered coin: the column
\(\begin{pmatrix}3/5\\4/5\end{pmatrix}\) is legal under one rule and
illegal under the other, so whatever a qubit state is, it is not a wager
about a hidden bit written in unusual units.

\begin{quote}
\textbf{Feynman's Notebook.}

\emph{Two levels of anything.} What in the world is described by such a
column? Nielsen and Chuang give a short list (\emph{Quantum Computation
and Quantum Information}, §1.2): the two polarizations of a photon; the
alignment of a nuclear spin in a uniform magnetic field; two states of
an electron orbiting a single atom, the ground state and an excited one.
As with the bit, the two labels \(0\) and \(1\) name two conditions of a
physical thing that can be reliably told apart. What is new is what else
the thing can do.

The idea of computing with such systems is usually traced to Richard
Feynman's ``Simulating Physics with Computers'', published in the
\emph{International Journal of Theoretical Physics} in 1982. As Nielsen
and Chuang tell it (§1.1), Feynman pointed to the difficulty of
simulating quantum mechanical systems on ordinary computers and
suggested that computers built on quantum mechanical principles would
avoid it. The word ``qubit'' is younger: Nielsen and Chuang credit it to
Benjamin Schumacher, in 1995.
\end{quote}

\begin{center}\rule{0.5\linewidth}{0.5pt}\end{center}

\hypertarget{the-rosetta-table}{%
\section{§1.6 The Rosetta Table}\label{the-rosetta-table}}

This chapter seeds the book's running dictionary between the blackboard
and the sketchbook. Five of the seven rows introduce picture elements
--- the wire, the box, the state, the row, and the bare wire; one names
the two states we use most; and the last is the book's first rewrite
rule.

\begingroup\renewcommand{\arraystretch}{1.5}

\begin{longtable}[]{@{}
  >{\raggedright\arraybackslash}p{(\columnwidth - 4\tabcolsep) * \real{0.2727}}
  >{\raggedright\arraybackslash}p{(\columnwidth - 4\tabcolsep) * \real{0.3636}}
  >{\raggedright\arraybackslash}p{(\columnwidth - 4\tabcolsep) * \real{0.3636}}@{}}
\toprule\noalign{}
\begin{minipage}[b]{\linewidth}\raggedright
\#
\end{minipage} & \begin{minipage}[b]{\linewidth}\raggedright
Algebra (Dirac)
\end{minipage} & \begin{minipage}[b]{\linewidth}\raggedright
Picture (Penrose)
\end{minipage} \\
\midrule\noalign{}
\endhead
\bottomrule\noalign{}
\endlastfoot
1 & a two-valued system, and the space its columns live in: \(\{0,1\}\),
then stochastic vectors, then \(\mathbb{C}^2\) & a \textbf{wire} ---
\emph{new element} \\
2 & a \(2\times2\) matrix \(A\) & a \textbf{box} labelled \(A\) on the
wire, input on the left, output on the right --- \emph{new element} \\
3 & a column \(\begin{pmatrix}\alpha\\\beta\end{pmatrix}\) & a
\textbf{state}: a box with no input and one wire out --- \emph{new
element} \\
4 & \(\begin{pmatrix}1\\0\end{pmatrix}\) and
\(\begin{pmatrix}0\\1\end{pmatrix}\) & the state boxes labelled \(0\)
and \(1\) \\
5 & the product \(BA\) (first \(A\), then \(B\)) & \(A\) and \(B\)
\textbf{in a row} on the wire, \(A\) on the left: read the picture left
to right, the product right to left --- \emph{new element} \\
6 & the identity matrix \(I\) & the \textbf{bare wire} --- \emph{new
element} \\
7 & \(XX=I\) & two NOT boxes in a row \(=\) the bare wire --- the first
\textbf{rewrite rule} \\
\end{longtable}

\endgroup

All seven rows are exact: the same numbers on both sides, entry for
entry.

\begin{center}\rule{0.5\linewidth}{0.5pt}\end{center}

\hypertarget{the-theorem-twice}{%
\section{§1.7 The Theorem, Twice}\label{the-theorem-twice}}

\begin{quote}
\textbf{Theorem 1.1 (NOT twice is nothing).} Let \(X\) be the matrix of
NOT and \(I\) the identity matrix. Then

\[X\,X \;=\; I, \qquad\text{that is,}\qquad \begin{pmatrix}0&1\\1&0\end{pmatrix}\begin{pmatrix}0&1\\1&0\end{pmatrix}=\begin{pmatrix}1&0\\0&1\end{pmatrix},\]

exactly. In pictures: two NOT boxes in a row are the bare wire.
\end{quote}

The example of §1.2 walked two starts through the two boxes. Here are
two proofs of the statement about the boxes themselves.

\hypertarget{diracs-blackboard}{%
\subsection{Dirac's Blackboard}\label{diracs-blackboard}}

The entry of a product \(BA\) in a given row and column is found by
running along that row of \(B\) and down that column of \(A\),
multiplying in pairs and adding. For \(XX\) all four entries are:

\[\begin{aligned}
\text{row 1, column 1:}&\quad 0\cdot0+1\cdot1 \;=\; 1,\\
\text{row 1, column 2:}&\quad 0\cdot1+1\cdot0 \;=\; 0,\\
\text{row 2, column 1:}&\quad 1\cdot0+0\cdot1 \;=\; 0,\\
\text{row 2, column 2:}&\quad 1\cdot1+0\cdot0 \;=\; 1.
\end{aligned}\]

Therefore

\[X\,X \;=\; \begin{pmatrix}0&1\\1&0\end{pmatrix}\begin{pmatrix}0&1\\1&0\end{pmatrix} \;=\; \begin{pmatrix}1&0\\0&1\end{pmatrix} \;=\; I.\]

And so, for \emph{every} column \(v\) of two numbers --- be they \(0\)
and \(1\), shares of a whole, or complex ---

\[X(Xv) \;=\; (XX)\,v \;=\; I\,v \;=\; v,\]

where the first equality is the defining property of the matrix product
(§1.3) and the last is the defining property of \(I\). Four entries,
eight multiplications, and nothing assumed about \(v\).
\(\;\blacksquare\)

\hypertarget{penroses-sketchbook}{%
\subsection{Penrose's Sketchbook}\label{penroses-sketchbook}}

The sketchbook owns, at this stage, exactly two equations that mention
NOT: \emph{NOT on 0} and \emph{NOT on 1}, from §1.1. Both have a state
box in them. The theorem has none. To get from equations about states to
an equation about boxes we need a principle, and we state it, and its
source, before using it.

\textbf{\emph{Two states suffice.}} \emph{If two boxes (or two rows of
boxes) give equal pictures when the state \(0\) is placed before them,
and equal pictures when the state \(1\) is placed before them, then they
are equal.}

This is not a fact the sketchbook has found for itself. It is lent by
the blackboard, where it is the remark of §1.1: the first column of a
matrix is what it does to the state \(0\), the second is what it does to
the state \(1\), and a matrix is nothing more than its two columns. It
is the other side of the trap of §1.4: one state was not enough to tell
two boxes apart; these two are. We also use, without further comment,
that a piece of a picture may be replaced by an equal piece --- equal
factors give equal products.

\textbf{Step 1 --- \emph{NOT on 0}.} Place the state \(0\) before the
two NOT boxes. The state and the first box together are the left side of
\emph{NOT on 0}; replace them by its right side, the state \(1\).

\[\begin{tikzpicture}[sd, baseline=-0.55ex]
  \draw[wire] (0,0) -- (9.4,0);
  \node[sdstate, text height=1.5ex, text depth=0.5ex] at (0,0) {$0$};
  \node[sdbox, text height=1.5ex, text depth=0.5ex] at (3.2,0) {$\mathrm{NOT}$};
  \node[sdbox, text height=1.5ex, text depth=0.5ex] at (6.6,0) {$\mathrm{NOT}$};
\end{tikzpicture}
\;\;\overset{\raisebox{0pt}[\height][1.2ex]{\scriptsize NOT on 0}}{=}\;\;
\begin{tikzpicture}[sd, baseline=-0.55ex]
  \draw[wire] (0,0) -- (6,0);
  \node[sdstate, text height=1.5ex, text depth=0.5ex] at (0,0) {$1$};
  \node[sdbox, text height=1.5ex, text depth=0.5ex] at (3.2,0) {$\mathrm{NOT}$};
\end{tikzpicture}\]

\textbf{Step 2 --- \emph{NOT on 1}.} What is left is the left side of
\emph{NOT on 1}; replace it by the state \(0\). A state followed by a
bare stretch of wire is that state, so this is also the state \(0\)
placed before the bare wire.

\[\begin{tikzpicture}[sd, baseline=-0.55ex]
  \draw[wire] (0,0) -- (6,0);
  \node[sdstate, text height=1.5ex, text depth=0.5ex] at (0,0) {$1$};
  \node[sdbox, text height=1.5ex, text depth=0.5ex] at (3.2,0) {$\mathrm{NOT}$};
\end{tikzpicture}
\;\;\overset{\raisebox{0pt}[\height][1.2ex]{\scriptsize NOT on 1}}{=}\;\;
\begin{tikzpicture}[sd, baseline=-0.55ex]
  \draw[wire] (0,0) -- (3,0);
  \node[sdstate, text height=1.5ex, text depth=0.5ex] at (0,0) {$0$};
\end{tikzpicture}\]

\textbf{Steps 3 and 4 --- \emph{NOT on 1}, then \emph{NOT on 0}.} Place
the state \(1\) before the two NOT boxes and use the same two equations
in the other order.

\[\begin{tikzpicture}[sd, baseline=-0.55ex]
  \draw[wire] (0,0) -- (9.4,0);
  \node[sdstate, text height=1.5ex, text depth=0.5ex] at (0,0) {$1$};
  \node[sdbox, text height=1.5ex, text depth=0.5ex] at (3.2,0) {$\mathrm{NOT}$};
  \node[sdbox, text height=1.5ex, text depth=0.5ex] at (6.6,0) {$\mathrm{NOT}$};
\end{tikzpicture}
\;\;\overset{\raisebox{0pt}[\height][1.2ex]{\scriptsize NOT on 1}}{=}\;\;
\begin{tikzpicture}[sd, baseline=-0.55ex]
  \draw[wire] (0,0) -- (6,0);
  \node[sdstate, text height=1.5ex, text depth=0.5ex] at (0,0) {$0$};
  \node[sdbox, text height=1.5ex, text depth=0.5ex] at (3.2,0) {$\mathrm{NOT}$};
\end{tikzpicture}\]

\[\begin{tikzpicture}[sd, baseline=-0.55ex]
  \draw[wire] (0,0) -- (6,0);
  \node[sdstate, text height=1.5ex, text depth=0.5ex] at (0,0) {$0$};
  \node[sdbox, text height=1.5ex, text depth=0.5ex] at (3.2,0) {$\mathrm{NOT}$};
\end{tikzpicture}
\;\;\overset{\raisebox{0pt}[\height][1.2ex]{\scriptsize NOT on 0}}{=}\;\;
\begin{tikzpicture}[sd, baseline=-0.55ex]
  \draw[wire] (0,0) -- (3,0);
  \node[sdstate, text height=1.5ex, text depth=0.5ex] at (0,0) {$1$};
\end{tikzpicture}\]

\textbf{Step 5 --- \emph{two states suffice}.} The row of two NOT boxes
and the bare wire give equal pictures after the state \(0\) (Steps 1--2)
and equal pictures after the state \(1\) (Steps 3--4). By the principle,
they are equal:

\[\begin{tikzpicture}[sd, baseline=-0.55ex]
  \draw[wire] (0,0) -- (8.4,0);
  \node[sdbox, text height=1.5ex, text depth=0.5ex] at (2.5,0) {$\mathrm{NOT}$};
  \node[sdbox, text height=1.5ex, text depth=0.5ex] at (5.9,0) {$\mathrm{NOT}$};
\end{tikzpicture}
\;\;\overset{\raisebox{0pt}[\height][1.2ex]{\scriptsize two states suffice}}{=}\;\;
\begin{tikzpicture}[sd, baseline=-0.55ex]
  \draw[wire] (0,0) -- (3,0);
\end{tikzpicture}
\qquad\blacksquare\]

\hypertarget{the-verdict}{%
\subsection{The verdict}\label{the-verdict}}

The algebra wins, plainly. The blackboard proof is one matrix product,
four entries, and it owes nothing to anyone. The sketchbook proof is
five steps, and its last step is not a picture-move at all: \emph{two
states suffice} is a fact about matrices and their columns, borrowed
from the blackboard. At this stage the sketchbook has no rule of its own
that could make two boxes vanish, and it would be dishonest to pretend
otherwise.

What the picture gives in return is not a shorter proof but a different
kind of statement. The equation between the two pictures never mentions
what the wire carries; and from this page on it is a \textbf{rewrite
rule} --- wherever two NOT boxes stand side by side in a row, they may
be lifted off, and wherever a wire is bare, two may be hung on it ---
which is how every later picture proof in this book will work. That is a
promise about the future, not a victory in the present. It proves an
equality of processes, and nothing else.

\begin{center}\rule{0.5\linewidth}{0.5pt}\end{center}

\hypertarget{the-round-trip}{%
\section{§1.8 The Round Trip}\label{the-round-trip}}

Here is the chapter's master diagram. We call it \textbf{The Round
Trip}: a wire, open at both ends, running from left to right, with two
NOT boxes in a row on it. Out, and home again. Its evaluation is Theorem
1.1:

\[\begin{tikzpicture}[sd, baseline=-0.55ex]
  \draw[wire] (0,0) -- (8.4,0);
  \node[sdbox, text height=1.5ex, text depth=0.5ex] at (2.5,0) {$\mathrm{NOT}$};
  \node[sdbox, text height=1.5ex, text depth=0.5ex] at (5.9,0) {$\mathrm{NOT}$};
\end{tikzpicture}
\;\;=\;\;
\begin{tikzpicture}[sd, baseline=-0.55ex]
  \draw[wire] (0,0) -- (3,0);
\end{tikzpicture}\]

On the blackboard, \(XX=I\); the equation is exact. We do not prove it
again. Instead we \emph{read} it, three times, by placing a state box at
the open left end. The picture is the same each time, and only the label
in the state box changes.

\textbf{First reading: a bit.} The passenger is a coin we can see, say
the state \(0\).

\[\begin{tikzpicture}[sd, baseline=-0.55ex]
  \draw[wire] (0,0) -- (9.4,0);
  \node[sdstate, text height=1.5ex, text depth=0.5ex] at (0,0) {$0$};
  \node[sdbox, text height=1.5ex, text depth=0.5ex] at (3.2,0) {$\mathrm{NOT}$};
  \node[sdbox, text height=1.5ex, text depth=0.5ex] at (6.6,0) {$\mathrm{NOT}$};
\end{tikzpicture}
\;=\;
\begin{tikzpicture}[sd, baseline=-0.55ex]
  \draw[wire] (0,0) -- (3,0);
  \node[sdstate, text height=1.5ex, text depth=0.5ex] at (0,0) {$0$};
\end{tikzpicture}
\qquad\qquad
XX\begin{pmatrix}1\\0\end{pmatrix}=\begin{pmatrix}1\\0\end{pmatrix}\]

\textbf{Second reading: a probabilistic bit.} The passenger is a coin
under a hand; the state box is labelled \(p\), standing for the column
\(\begin{pmatrix}p\\1-p\end{pmatrix}\).

\[\begin{tikzpicture}[sd, baseline=-0.55ex]
  \draw[wire] (0,0) -- (9.4,0);
  \node[sdstate, text height=1.5ex, text depth=0.5ex] at (0,0) {$p$};
  \node[sdbox, text height=1.5ex, text depth=0.5ex] at (3.2,0) {$\mathrm{NOT}$};
  \node[sdbox, text height=1.5ex, text depth=0.5ex] at (6.6,0) {$\mathrm{NOT}$};
\end{tikzpicture}
\;=\;
\begin{tikzpicture}[sd, baseline=-0.55ex]
  \draw[wire] (0,0) -- (3,0);
  \node[sdstate, text height=1.5ex, text depth=0.5ex] at (0,0) {$p$};
\end{tikzpicture}
\qquad\qquad
XX\begin{pmatrix}p\\1-p\end{pmatrix}=\begin{pmatrix}p\\1-p\end{pmatrix}\]

\textbf{Third reading: a qubit.} The passenger is a column
\(\psi=\begin{pmatrix}\alpha\\\beta\end{pmatrix}\) of \(\mathbb{C}^2\)
with \(|\alpha|^2+|\beta|^2=1\).

\[\begin{tikzpicture}[sd, baseline=-0.55ex]
  \draw[wire] (0,0) -- (9.4,0);
  \node[sdstate, text height=1.5ex, text depth=0.5ex] at (0,0) {$\psi$};
  \node[sdbox, text height=1.5ex, text depth=0.5ex] at (3.2,0) {$\mathrm{NOT}$};
  \node[sdbox, text height=1.5ex, text depth=0.5ex] at (6.6,0) {$\mathrm{NOT}$};
\end{tikzpicture}
\;=\;
\begin{tikzpicture}[sd, baseline=-0.55ex]
  \draw[wire] (0,0) -- (3,0);
  \node[sdstate, text height=1.5ex, text depth=0.5ex] at (0,0) {$\psi$};
\end{tikzpicture}
\qquad\qquad
XX\begin{pmatrix}\alpha\\\beta\end{pmatrix}=\begin{pmatrix}\alpha\\\beta\end{pmatrix}\]

One picture, three passengers, one answer. In each reading the passenger
really does make the trip --- the state \(0\) becomes the state \(1\),
the shares change places, the entries \(\alpha\) and \(\beta\) are
exchanged --- and in each it comes home exactly as it left. The equation
did not need to be told which reading was intended, because it was never
about the passenger. See \texttt{Code/Ch1/TwiceOver.py} for the equation
and for the first two passengers, and \texttt{Code/Ch1/Arrow.py} for the
third.

Let us also use the equation once as what §1.7 said it had become, a
rewrite rule. Here is a row of five NOT boxes. Lift off one adjacent
pair, then another:

\[\begin{tikzpicture}[sd, baseline=-0.55ex]
  \draw[wire] (0,0) -- (18.2,0);
  \node[sdbox, text height=1.5ex, text depth=0.5ex] at (2.3,0) {$\mathrm{NOT}$};
  \node[sdbox, text height=1.5ex, text depth=0.5ex] at (5.7,0) {$\mathrm{NOT}$};
  \node[sdbox, text height=1.5ex, text depth=0.5ex] at (9.1,0) {$\mathrm{NOT}$};
  \node[sdbox, text height=1.5ex, text depth=0.5ex] at (12.5,0) {$\mathrm{NOT}$};
  \node[sdbox, text height=1.5ex, text depth=0.5ex] at (15.9,0) {$\mathrm{NOT}$};
\end{tikzpicture}\]

\[\overset{\raisebox{0pt}[\height][1.2ex]{\scriptsize round trip}}{=}\;\;
\begin{tikzpicture}[sd, baseline=-0.55ex]
  \draw[wire] (0,0) -- (11.4,0);
  \node[sdbox, text height=1.5ex, text depth=0.5ex] at (2.3,0) {$\mathrm{NOT}$};
  \node[sdbox, text height=1.5ex, text depth=0.5ex] at (5.7,0) {$\mathrm{NOT}$};
  \node[sdbox, text height=1.5ex, text depth=0.5ex] at (9.1,0) {$\mathrm{NOT}$};
\end{tikzpicture}
\;\;\overset{\raisebox{0pt}[\height][1.2ex]{\scriptsize round trip}}{=}\;\;
\begin{tikzpicture}[sd, baseline=-0.55ex]
  \draw[wire] (0,0) -- (4.6,0);
  \node[sdbox, text height=1.5ex, text depth=0.5ex] at (2.3,0) {$\mathrm{NOT}$};
\end{tikzpicture}\]

No state was plugged in and no matrix was multiplied: \(X^5=X\), by
lifting boxes off a wire. It is a small thing. It is also the first
calculation in this book done entirely in the sketchbook, on credit from
the blackboard, and the later chapters will run up a considerable
account.

\begin{center}\rule{0.5\linewidth}{0.5pt}\end{center}

\hypertarget{the-other-end-of-the-wire}{%
\section{§1.9 The Other End of the
Wire}\label{the-other-end-of-the-wire}}

Let us gather what we have. A system is a wire; a process is a box,
which is a \(2\times2\) matrix; a way of starting is a state, which is a
column (§1.1). Two NOTs in a row are the bare wire, which we proved
twice, and the algebra won (§1.2, §1.7). Boxes in a row are a matrix
product written backwards, and their order matters (§1.3). The column
can hold a fact; it can hold a wager, with entries that sum to \(1\)
(§1.4); and it can hold something that is neither, with complex entries
whose squared moduli sum to \(1\) (§1.5). Through all three the pictures
did not change, and \emph{The Round Trip} came out the same (§1.8).

Now look at the pictures once more and notice what none of them has.
Every wire in this chapter starts somewhere --- we have a box with no
input, and we have used it on every page --- but no wire \emph{ends}.
Each one runs off the right-hand side of the picture, still carrying its
passenger, and nobody has asked the passenger anything. For a coin that
hardly mattered: one looks, and sees a face. For a covered coin one
lifts the hand, and the wager is settled by a fact. But for the third
passenger we have not said what looking would be, or what one would see,
or how often. We have a column of two complex numbers and a rule about
their squared moduli, and no connection yet between those numbers and
anything that happens.

That is the business of the next chapter. It will draw the thing that
closes the right-hand end of a wire, and find that a wire closed at both
ends is no longer a process at all but a single number. It will say what
that number has to do with what is seen, and why the rule must use
squares. And it will finally use the fact we mentioned in §1.5 and then
set carefully aside: that the entries of a qubit's column, unlike any
wager's, are allowed to be negative.

\hypertarget{nocturne-in-negative-numbers}{%
\chapter{Nocturne in Negative
Numbers}\label{nocturne-in-negative-numbers}}

\begin{center}\rule{0.5\linewidth}{0.5pt}\end{center}

\emph{水底に}\\
\emph{もう一つある}\\
\emph{雪催}

\emph{minasoko ni}\\
\emph{mō hitotsu aru}\\
\emph{yukimoyoi}

\emph{Under the water}\\
\emph{there is another one, too ---}\\
\emph{sky heavy with snow.}

\begin{center}\rule{0.5\linewidth}{0.5pt}\end{center}

\emph{{[}The snow has not come yet but means to: all evening the sky has
hung low and close, the loaded yellow-grey that stands before a heavy
fall, and the shutters are drawn against it. Above the hearth, where the
firelight can reach it, hangs a print Achilles had framed in the spring
--- Escher's \textup{Sky and Water I}: dark birds ranked across the top,
pale fish across the bottom, and a broad middle band where the two come
to the same size and fit close, wing against fin. As the coals shift,
the light picks out the birds and loses the fish, then loses the birds
and finds the fish. The table stands bare and scrubbed, waiting.{]}}

\emph{{[}Evening; a low fire burning down. The Tortoise sits exactly
where she sat last evening, one foot planted just so --- she has not
shifted it since. Achilles comes back from the cold hall with a fresh
pot and two cups.{]}}

\textbf{Achilles:} Still there. You've not budged an inch since last
evening --- I do believe you've taken root in my chair.

\textbf{Tortoise:} I have a coin to mind. \emph{{[}She settles more
comfortably, without moving the foot.{]}} Some things keep better left
exactly as they are.

\textbf{Achilles:} \emph{{[}pouring{]}} Then keep it. It's a filthy
night for anything but the fire and a warm cup. \emph{{[}He hands her
one.{]}} That overture of last evening is still going round in my head,
mind --- I hummed it half the day and got the ending wrong every time.

\textbf{Tortoise:} The good ones do that. \emph{{[}warming her hands
round the cup{]}} There is nothing to be done on a night like this but
sit close to the fire and think slowly.

\textbf{Achilles:} Slowly suits me. Though I'll warn you --- there's a
question left over from last evening that has kept me from my sleep.

\textbf{Tortoise:} They always are. Let me get my strings out of my
shell, and we shall see what can be done about it.

\begin{center}\rule{0.5\linewidth}{0.5pt}\end{center}

\hypertarget{first-challenge-the-right-hand-end}{%
\section{\texorpdfstring{First Challenge --- \emph{The Right-Hand
End}}{First Challenge --- The Right-Hand End}}\label{first-challenge-the-right-hand-end}}

\emph{{[}Evening; Achilles's table. The Tortoise draws a pale string out
of her shell and a dark one after it. She pins the pale string taut and
straight, left to right, on the left half of the table; she pins the
dark string on the same line after it, a hand's breadth of bare table
between. She plucks the pale string once: one note. Under her foot is a
coin that was covered on an earlier evening; nobody has looked at it
since.{]}}

\textbf{Achilles:} You never answered me. An arrow sets out from the
left-hand end. When it comes to the right-hand end, how do I find out
what has arrived?

\textbf{Tortoise:} \emph{{[}She rests a fingertip on the pale string,
then on the dark one lying after it.{]}} Gently, now. This pale string
is the one we shall work along: your arrow sets out from its left end,
and your question waits at the right. The dark one lies after it, a good
hand's breadth of bare table between, so they never muddle --- it will
have its turn. And everything I carry about tonight I'll draw on one of
these little squares of tracing paper: an arrow inked in from a corner,
stiff enough to knot on by that corner. Tracing paper, mind, so that
when I want the same arrow in two places at once I can copy it rather
than move it.

\textbf{Tortoise:} \emph{{[}She brings out a square of tracing paper
with an arrow drawn on it from the near left-hand corner, and ties it by
that corner to the pale string's left end.{]}} Last evening's arrow:
three matches heads-ward, four tails-ward. I called this the end tag,
then. Tonight a string has an end to spare, so call it the start tag.

\textbf{Achilles:} Five matches long, I remember. And the dark string?

\textbf{Tortoise:} Later. \emph{{[}She draws an arrow on a second
square, from the near left-hand corner straight along the bottom edge,
five matches long.{]}} Heads-ward, and nothing else. It would do for a
start tag.

\textbf{Achilles:} A dull one. A coin lying heads up.

\textbf{Tortoise:} Dull or not, watch what I do with it. A string has an
end to spare tonight, and a tag can be tied on for either errand: to
send an arrow off from the near end, or to sit at the far end and wait
for whatever comes along the string to it.

\textbf{Tortoise:} \emph{{[}The start tag at the left end is set to send
its reach off down the string. She takes a fresh square drawn the very
same way and carries it along the pale string to the right-hand end as
it lies --- not turned round on the table, not turned over --- so the
arrow keeps pointing exactly as it did; only she ties it on now by the
far right-hand corner, the corner opposite, so that the knot, and not
the arrow, faces back along the string toward the start.{]}} The arrow
points as it pointed. Only the knot has changed its mind: this tag
doesn't send anything along the string. It waits for what comes. A
closing tag.

\textbf{Achilles:} Then hang something beyond it, and let us see what
falls off the end.

\textbf{Tortoise:} Where?

\textbf{Achilles:} \emph{{[}His hand goes to the right of the closing
tag, and finds bare table.{]}} There's no string left to hang it on. So
what you're saying is, a start at one end, a question at the other, and
the string hands over a coin. Heads, I expect.

\textbf{Tortoise:} Then let me show you how we put a question to the
string at all. I have another sheet, ruled with two lines that cross
square, an arrowhead on each --- picture the crossing as the question
itself, and each line as one of the answers it might be given, the one
way or the other. Laid over an arrow, it lets us read how far the arrow
reaches toward each answer. For measuring that reach I keep this
square-cornered card, whose true corner drops a straight line from the
arrow's tip down onto a ruled line, and matchsticks all cut to one
length, for a ruler --- so every reach we find is spoken in matches.

\textbf{Tortoise:} Let us find what it hands over. \emph{{[}She brings
out a sheet of tracing paper ruled with two lines crossing at right
angles, an arrowhead on each. She lays it over the closing tag with one
line along the tag's arrow, arrowheads agreeing; then she slides it back
along the pale string, as it lies, and lets it rest on the start tag
with the crossing on the arrow's tail.{]}} A question with two answers:
heads-ward along this line, tails-ward along that. \emph{{[}She hands
him a square-cornered card.{]}} From the arrow's tip, straight down to
the heads-ward line. Then matches, from the crossing to the foot of the
card.

\textbf{Achilles:} \emph{{[}He stands the square-cornered card on its
true corner at the arrow's tip and lets its edge fall straight down onto
the heads-ward line; where the corner meets the line is the foot. From
the crossing out to that foot he lays the matches end to end.{]}} Three.
And the arrow is five. Three parts in five.

\textbf{Tortoise:} \emph{{[}She unties the closing tag, and ties on in
its place, the same way, a square with an arrow five matches long up its
left-hand side.{]}} Tails-ward, and nothing else.

\textbf{Achilles:} \emph{{[}The card drops from the tip to the
tails-ward line this time; the foot comes down further out along it.{]}}
Four. Four parts in five. But that's no coin. A string tagged at both
ends hands over a \emph{number}?

\textbf{Tortoise:} One number. It has nothing left to travel to.

\emph{{[}Somewhere in the shutter's frame the wind has found a gap and
settled into it, and out of that one closed slot it draws a single
steady note, the same however hard or soft it blows. Achilles tips his
head at the sound, then lets it go. Against the shutter the first snow
ticks --- a few grains, no more.{]}}

\textbf{Achilles:} Then those are the shares! Heads three parts in five,
tails four. Which makes seven parts in --- five. They overshoot the
whole, exactly as they did last evening. \emph{{[}He looks along the
table.{]}} And the dark one?

\textbf{Tortoise:} Now.

\begin{center}\rule{0.5\linewidth}{0.5pt}\end{center}

\hypertarget{second-challenge-why-squares}{%
\section{\texorpdfstring{Second Challenge --- \emph{Why
Squares}}{Second Challenge --- Why Squares}}\label{second-challenge-why-squares}}

\textbf{Tortoise:} The same tags, each at the other end. Trace them; a
square can't hang in two places.

\emph{{[}Achilles lays a fresh square over the heads-ward arrow and, the
tracing paper letting the shape through, inks its copy; he traces last
evening's three-and-four arrow the same way. Now the Tortoise gives the
dark string the two squares the pale string had --- but each at the end
the pale string did not use. The heads-ward tracing she ties to the dark
string's left end, where a start tag sends its reach off; the
three-and-four tracing she carries to the right-hand end as it lies and
ties with the knot facing back, to wait. The arrows are the very ones
from the pale string; only their errands have swapped ends --- the arrow
that asked now waits, and the arrow that waited now asks.{]}}

\textbf{Achilles:} The question sets out, and my arrow waits for it.
That's the wrong way about.

\textbf{Tortoise:} Find its number. No ruled sheet this time: lay its
start over its closing tag, tail on tail, as they lie. Drop the card
from the tip of the heads-ward arrow onto the line of the other.

\textbf{Achilles:} \emph{{[}He slides the heads-ward start square over
the three-and-four square until tail sits on tail, the two arrows
springing from one point; the card drops from the heads-ward tip onto
the three-and-four arrow's line, and from that arrow's tail out to the
foot he lays the matches.{]}} Three. And the heads-ward arrow was drawn
five long. Three parts in five --- the pale string's number, arriving
from the other end.

\textbf{Tortoise:} And tails-ward?

\textbf{Achilles:} \emph{{[}He traces the tails-ward arrow onto a fresh
square and swaps it in for the heads-ward start; he lays it over the
three-and-four closing, tail on tail as before, and drops the card from
its tip --- and now the foot lands further along the line, and the
matches reach four.{]}} Four in five on the pale string. Four in five on
the dark. It only repeats what it was told!

\textbf{Tortoise:} Good --- keep hold of that. Those little squares in
the dish, each a single matchstick to a side, are for laying a shape out
flat where you can count it --- nothing cleverer than that. Reach for a
handful when you want them.

\textbf{Tortoise:} Then you have two equal sides. \emph{{[}She tips a
heap of little paper squares onto the table, well clear of the pale
string and the dark, each one match on a side.{]}} One side from the
pale string, one from the dark. Fill it in.

\textbf{Achilles:} \emph{{[}He sets the little squares out where the
reach ran and stacks equal rows up from it, as many high as the reach is
wide, until the block is full; then a second block the same way on the
tails-ward reach.{]}} Three by three on the heads-ward reach: nine. Four
by four on the tails-ward: sixteen. Nine and sixteen. Twenty-five in
all.

\textbf{Tortoise:} And to hold the answer, reach again for the strip
from last evening --- the paper that stood in for the covered coin's
wager.

\textbf{Tortoise:} Now a strip. Fold it in twenty-five. From the knot
end out to the change of shade is heads' share, as it was for the
covered coin.

\textbf{Achilles:} \emph{{[}He folds the strip across into twenty-five
equal parts, then shades it from the knot end inward, part by part, and
stops after the ninth; from there to the far end the strip stays
pale.{]}} Heads, nine in twenty-five; tails, sixteen. And those
\emph{do} make one whole. But this is a wager --- a coin under a foot!
There's no foot here, only an arrow and a question.

\emph{{[}On the sill beyond the shutter the snow has begun to lie, and
it only ever adds --- a grain, another grain, never one taken back ---
so that the few grains of a moment ago have become a covering, and the
covering deepens and evens until the whole ledge carries the same soft
weight from end to end. Achilles fills both cups again from the pot
without looking up; the tea has gone from scalding to merely warm.{]}}

\textbf{Tortoise:} That is the recipe. An arrow, a question, the reach
on the pale string, the same on the dark, the square on the two. If ever
that question were put to that arrow, the strip says how often each
answer would come.

\textbf{Achilles:} But why should the world go by the squares?

\textbf{Tortoise:} That is how often. Nobody has caught it doing
otherwise. \emph{{[}She takes a fresh square of tracing paper and draws
an arrow from its far left-hand corner: one match heads-ward, and one
match towards herself.{]}} One heads-ward. One tails-ward, backwards.

\textbf{Achilles:} A reach of less than nothing?

\begin{center}\rule{0.5\linewidth}{0.5pt}\end{center}

\hypertarget{third-challenge-a-reach-backwards}{%
\section{\texorpdfstring{Third Challenge --- \emph{A Reach
Backwards}}{Third Challenge --- A Reach Backwards}}\label{third-challenge-a-reach-backwards}}

\textbf{Tortoise:} First the honest one. \emph{{[}She draws another from
the near left-hand corner: one match heads-ward, one tails-ward.{]}} The
slant. Tie it on and ask it.

\emph{{[}The slant is tied to the pale string's left end. The Tortoise
marks its length along the card's edge and draws the closing arrows,
heads-ward and tails-ward, exactly so long. Achilles traces; the dark
string is tagged the other way about.{]}}

\textbf{Achilles:} \emph{{[}He lays the ruled sheet square over the
slant, crossing on the tail, and drops the card from the arrow's tip to
each line in turn --- to the heads-ward line, then to the tails-ward ---
laying matches from the crossing to each foot.{]}} Heads-ward, one
match. Tails-ward, one match. And along the slant itself --- \emph{{[}He
turns the card to the slant's own arrow and lays matches up it from the
crossing.{]}} --- one and a bit. Not two. One part in\ldots{} it doesn't
come out in matches.

\textbf{Tortoise:} Then don't say it in matches. And the dark one?

\textbf{Achilles:} One match, and one match. The same.

\textbf{Tortoise:} Squares.

\textbf{Achilles:} One little square on each reach. But the square on
the slant?

\textbf{Tortoise:} \emph{{[}She takes the two little squares --- the one
from the heads-ward reach, the one from the tails-ward --- and cuts each
along its diagonal, corner to corner; the four half-squares she fits
point to point around the slant until they close into one square
standing squarely on it. The paper is all there: none left over, none
wanting.{]}} Made of those two, and nothing over.

\textbf{Achilles:} Then each reach has one square of the two. Heads, one
in two. \emph{{[}He folds a strip in two and shades one part.{]}} Evens:
my half-and-half strip again. Now your back-slant.

\emph{{[}The back-slant is tied on as the pale string's start; its
tracing becomes the dark string's closing tag.{]}}

\emph{{[}He lays the sheet square over the back-slant and drops the card
from its tip to each line. The heads-ward reach lands out along its
line, honest enough; but the tails-ward stroke of this arrow leans back
past its own tail, and the card comes down where no reach has come down
before.{]}}

\textbf{Achilles:} Heads-ward: one match. Tails-ward --- the foot of the
card falls \emph{behind} the crossing. One match, backwards.

\textbf{Tortoise:} And the dark one?

\textbf{Achilles:} \emph{{[}He lays the tails-ward start over the
back-slant. The card's foot misses the arrow; he produces the arrow's
line backwards through its tail with the card's edge, and lays a match
there.{]}} One match, backwards. Backwards, and backwards again. But can
one lay paper on a reach of less than nothing?

\textbf{Tortoise:} \emph{{[}She lays a little square on the two
backwards sides, on the far side of the crossing.{]}} Pick it up. Is
there less of it than of the others?

\textbf{Achilles:} It's a square like any other. One, and one: the four
triangles, the same strip. \emph{{[}He holds the strip against the
back-slant, and then against the slant.{]}} So the backwards reach is
bookkeeping! A mark for the paper, and nothing to the world. Same strip,
same start. \emph{{[}He ties the one strip on at the pale string's left
end, beside the square.{]}} One strip will do for both arrows.

\emph{{[}Outside, the wind has quietly gone round to blow from the other
side of the house --- the chimney, which had been steady, takes up a low
mutter, and the draught that had been coming in at the door slackens and
dies. The fire does not notice; it burns exactly as it burned, and the
one note in the shutter holds its pitch without a waver. Achilles sits
back, satisfied, and warms his hands.{]}}

\textbf{Tortoise:} \emph{{[}letting him{]}} You have asked it one
question.

\emph{{[}She sets a fingertip on the crossing of the ruled sheet,
holding it fast to the arrow's tail, and swings the sheet about that one
point. The crossing does not stir; the two ruled lines wheel away from
heads-ward and tails-ward.{]}}

\begin{center}\rule{0.5\linewidth}{0.5pt}\end{center}

\hypertarget{fourth-challenge-the-turned-question}{%
\section{\texorpdfstring{Fourth Challenge --- \emph{The Turned
Question}}{Fourth Challenge --- The Turned Question}}\label{fourth-challenge-the-turned-question}}

\emph{{[}The ruled sheet has come to rest a half-right-angle round, the
crossing still pinned to the arrow's tail. One of its lines now runs
along the slant; the other, along the back-slant.{]}}

\emph{{[}The wind veers again --- not round to the far side now but a
half-turn, crosswise --- and the note that the shutter had held falls to
a low nothing, while a chink somewhere that had been silent all evening
begins to whistle. The house sounds like another house.{]}}

\textbf{Achilles:} That's no question. Neither line is heads-ward or
tails-ward.

\textbf{Tortoise:} Two lines at right angles, two answers. It needs
closing tags, that is all: an arrow along each line, as long as the
start. You have drawn them already.

\textbf{Achilles:} The slant and the back-slant themselves! \emph{{[}He
traces them. The back-slant is the pale string's start; a tracing of the
slant is carried across as it lies and tied on, knot facing back.{]}}
Back-slant, asked along the slant's line. \emph{{[}He drops the card.
Its foot falls on the crossing.{]}} No matches at all. It has no reach
along that line.

\textbf{Tortoise:} And along the other?

\textbf{Achilles:} It \emph{lies} along the other. The whole arrow, and
nothing left over.

\textbf{Tortoise:} And the dark one?

\textbf{Achilles:} \emph{{[}laying one tracing over the other{]}}
Nothing; and then the whole arrow. The same. And with the slant for a
start it is the other way about: the whole arrow along its own line,
none along the back-slant's. So the strips --- \emph{{[}He shades one
strip from end to end, and leaves another white.{]}} All dark. All
white.

\emph{{[}He unties the half-and-half strip from beside the start tag.
Then he swings the ruled sheet back until it sits square, heads-ward and
tails-ward; then round again until it lies along the slant and the
back-slant; then back to square --- rocking it between the two, the same
two lines telling him first that the arrows are alike and then that they
are utterly apart. He looks up at the print on the wall --- Escher's
\textup{Sky and Water I}.{]}}

\emph{{[}The fire chooses that moment to settle, and for a breath the
flames stand still; the light lies flat and even across the whole print,
birds and fish alike, neither nearer than the other. Then a coal drops,
and the shadows come back.{]}}

\textbf{Achilles:} The fish are only what is left between the birds. And
the birds are what is left between the fish.

\textbf{Tortoise:} And the sheet?

\textbf{Achilles:} \ldots Oh. They were never the same start. One
question hands me the same strip for both, and the next tells them apart
every time --- they stand at right angles, as far apart as heads-ward
from tails-ward. The backwards reach was the whole difference.

\textbf{Tortoise:} Last evening's arrow, under the turned sheet.

\emph{{[}The three-and-four arrow is tied on again. The Tortoise lays
five matches along each turned line and draws the closing arrows so;
Achilles traces, and the dark string is tagged the other way about.{]}}

\textbf{Achilles:} Along the slant's line: just short of five matches.
Along the back-slant's: short of one match, and backwards.

\textbf{Tortoise:} And the dark one?

\textbf{Achilles:} Just short of five; short of one, backwards. So the
shares: all but a sliver --- and the sliver comes from a backwards
reach. Then every backwards reach can be caught, if one turns the
question far enough!

\begin{center}\rule{0.5\linewidth}{0.5pt}\end{center}

\hypertarget{fifth-challenge-the-arrow-reversed}{%
\section{\texorpdfstring{Fifth Challenge --- \emph{The Arrow
Reversed}}{Fifth Challenge --- The Arrow Reversed}}\label{fifth-challenge-the-arrow-reversed}}

\textbf{Tortoise:} \emph{{[}She unties the pale string's start tag and
turns the square end for end on the table --- the tip swung down to
where the tail was, the arrow now pointing the opposite way though it is
the very same arrow of the very same length --- and ties it on again by
what had been its far corner.{]}} Every reach backwards at once.

\textbf{Achilles:} Then it will be caught at once. \emph{{[}He lays the
sheet square and drops the card from the reversed arrow's tip to each
line; both feet come down behind the crossing now, where before they
fell in front of it.{]}} Heads-ward: three, backwards. Tails-ward: four,
backwards.

\textbf{Tortoise:} And the dark one? Its closing tag is a tracing of
this. \emph{{[}She turns that square end for end likewise, and ties it
on, knot facing back; the heads-ward tracing is the dark string's start
again.{]}}

\textbf{Achilles:} \emph{{[}producing the reversed arrow's line back
through its tail{]}} Three, backwards. Four, backwards. \emph{{[}He lays
the little squares on the far side of the crossing.{]}} Nine and
sixteen. My first strip. Well --- the square sheet missed the back-slant
too. \emph{{[}He turns the sheet.{]}} Just short of five, backwards.
Short of one match, forwards. And the dark one?

\textbf{Tortoise:} That was my line.

\textbf{Achilles:} The same, reach for reach. All but a sliver, as
before.

\textbf{Tortoise:} A different arrow. Find me the question that knows
it.

\emph{{[}Achilles turns the ruled sheet to a position of his own
choosing and holds it there. He lays matches along its lines for the
arrow reversed. Then, the sheet held still, he turns the square end for
end beneath it, lays a second row of matches beside the first, and turns
the square back.{]}}

\textbf{Achilles:} The same lengths. Every one on the other side of the
crossing. \emph{{[}The Tortoise lays five matches along one of his lines
and draws the closing arrow; he traces, and lays the dark string's
squares one over the other.{]}} And whatever turns backwards on the pale
string turns backwards on the dark. Two sides, a paper square, the same
square. I can't catch it!

\emph{{[}At the height of it the storm comes round the whole way --- the
wind that had been crosswise swings until it stands head-on against the
shutters, and in one breath every draught in the room trades places, the
cold from the door and the cold from the chimney changing ends. The fire
burns on as it burned; the tea sits cooling in the cups. Every draught
is reversed, and nothing in the room can be seen to have moved.{]}}

\textbf{Tortoise:} One reach turned backwards?

\textbf{Achilles:} A different start, and some question knows it every
time. Every reach turned backwards: nothing to catch.

\textbf{Tortoise:} A different arrow on the paper, even so. And no
question we can tie on tells them apart.

\textbf{Achilles:} Then what \emph{is} fixed? Five of what? Twenty-five
of what?

\begin{center}\rule{0.5\linewidth}{0.5pt}\end{center}

\hypertarget{sixth-challenge-the-whole-square}{%
\section{\texorpdfstring{Sixth Challenge --- \emph{The Whole
Square}}{Sixth Challenge --- The Whole Square}}\label{sixth-challenge-the-whole-square}}

\textbf{Tortoise:} \emph{{[}She turns the reversed squares end for end
again: last evening's arrow, as it was drawn. One she ties to the pale
string's left end.{]}} Trace it.

\emph{{[}Achilles traces it. The Tortoise carries the tracing to the
pale string's right-hand end as it lies, and ties it on with the knot
facing back.{]}}

\textbf{Achilles:} The arrow, asked about --- itself?

\textbf{Tortoise:} \emph{{[}She turns the ruled sheet until one line
lies along the arrow.{]}} Ask.

\textbf{Achilles:} \emph{{[}The card's foot falls on the arrow's own
tip.{]}} Five. Five parts in five: all of it. And along the other line,
no reach at all.

\textbf{Tortoise:} And the dark one?

\textbf{Achilles:} \emph{{[}He traces the arrow twice more, ties one
copy to the dark string's left end and carries the other to its right,
and lays the start square over the closing, tail on tail; the card drops
from the tip clean onto the arrow's own tip.{]}} Five in five. What else
could it say?

\textbf{Tortoise:} \emph{{[}She draws on a sheet of paper a square five
matches on a side.{]}} The square on it. Fill it in.

\textbf{Achilles:} I have no more little squares. Only the nine and the
sixteen from --- \emph{{[}He breaks the three-by-three block into its
nine loose squares and the four-by-four into its sixteen, and sets them
one by one into her five-by-five outline, row filling into row, until
the outline is a full field of squares.{]}} It fills.

\textbf{Tortoise:} Exactly?

\textbf{Achilles:} None over, none short. \ldots Oh. A string tagged at
both ends is a number, and the number may be backwards. The number with
its twin from the dark string is a square of paper, and that never is.
And the squares, whatever the question, make up the square on the arrow
asked about itself: one whole.

\emph{{[}The snow has stopped. Where the shutter stands a little open,
the world outside has gone to one unbroken white, sill and yard and low
roofs filled level and made all of a piece. The wind has dropped away to
nothing, and the house has stopped sounding altogether.{]}}

\textbf{Tortoise:} Twenty-five of that.

\emph{{[}She unties every tag and puts the squares away. The pale string
and the dark string lie bare, end to end, a hand's breadth of table
between them, as they lay at the start. Her foot has not moved from the
coin. Achilles plucks the pale string: one note. He plucks the dark
string: the same note, after it.{]}}

\textbf{Achilles:} \emph{{[}He warms his hands at what is left of the
fire.{]}} The embers are nearly out, and the tea went cold an age ago. A
good hour to stop, this.

\textbf{Tortoise:} \emph{{[}not moving{]}} A very good hour.

\textbf{Achilles:} But the reaches themselves --- three in five, one
match backwards. Not shares. Not facts. What \emph{are} they?

\textbf{Tortoise:} They have a name.

\textbf{Achilles:} Tell me.

\textbf{Tortoise:} It will keep till morning.

\emph{{[}The fire is a low red breathing under grey. Beyond the shutter
the white field keeps its stillness, and will keep it till morning.
Above the hearth the print has gone back to being a print --- birds
along the top, fish along the bottom, dim now and part of the wall. The
two bare strings lie where she left them, and under the Tortoise's foot
the coin lies covered, as it has lain all night, and nobody has
looked.{]}}

\hypertarget{chapter-2-amplitudes}{%
\chapter{Chapter 2: Amplitudes}\label{chapter-2-amplitudes}}

\emph{--- Nocturne in Negative Numbers ---}

\emph{Part I: The Single Wire}

\emph{Escher: ``Sky and Water I'' --- birds above and fish below, and in
the band between them each is only the space the other leaves}

\begin{center}\rule{0.5\linewidth}{0.5pt}\end{center}

How can a number that is allowed to be less than nothing tell us how
often something will happen?

For the covered coin of the last chapter there was no such puzzle. The
two numbers in its column \emph{were} the chances: heads five times in
six, tails once in six. They were never negative, they made up one
whole, and lifting the hand settled the wager with a fact. Whatever the
column said, it said about what would be seen.

Then we wrote a third kind of thing in the same two places. Its entries
reach three fifths one way and four fifths the other, and three fifths
and four fifths make seven fifths, which is no whole at all. Its entries
may be negative. They may fail to be real numbers. Nobody's chances look
like that. And yet when a qubit is finally asked a question with two
answers, an answer comes; and when the same question is put to a great
many qubits prepared in the same way, the answers come in steady
proportions, as reliable as a die. Somewhere between the column and the
counting there must be a rule that turns the one into the other.

This chapter finds the place where that rule lives. We left every wire
of Chapter 1 running off the right-hand side of the page with nobody
asking its passenger anything. Here we close that end. What comes out of
a wire closed at both ends is a single number, and it is not yet a
chance: it has a sign, or worse. The chance appears only when the number
is heard together with its mirror twin. So this is a chapter about
numbers that are never themselves seen --- it is night, the coin stays
covered, nothing is looked at --- and about the one way they have of
showing in the morning.

Those numbers have a name. It is the title of the chapter.

\begin{center}\rule{0.5\linewidth}{0.5pt}\end{center}

\hypertarget{a-name-for-a-column}{%
\section{§2.1 A Name for a Column}\label{a-name-for-a-column}}

Before we close the wire we need a better way of writing what travels on
it. Chapter 1 wrote every state out as a column, in full, every time.
That is honest and it is slow, and it has a deeper fault: it gives us no
way of \emph{naming} a column, so that we could speak of the same thing
lying on its side.

\textbf{On the blackboard.} A \textbf{ket} is a name for a column,
written between a vertical bar and an angle bracket: \(\ket{\psi}\),
read ``ket psi''. What stands inside is a label and nothing more. We fix
two labels for good:

\[\ket{0} \;=\; \begin{pmatrix}1\\0\end{pmatrix}, \qquad\qquad \ket{1} \;=\; \begin{pmatrix}0\\1\end{pmatrix}.\]

These are the states \(0\) and \(1\) of §1.1, with the first (top) entry
belonging to \(0\) as always. A word of warning that every beginner
needs once: \textbf{\(\ket{0}\) is not the zero column.} The \(0\)
inside the ket is a name, as ``heads'' would be; the column it names has
a \(1\) in it. The column with two zero entries is not a qubit state at
all, since its squared moduli sum to \(0\), and we shall never give it a
ket.

Columns can be multiplied by numbers, entry by entry, and added, entry
by entry. Doing so with our two named columns recovers every column
there is:

\[\alpha\ket{0}+\beta\ket{1} \;=\; \alpha\begin{pmatrix}1\\0\end{pmatrix}+\beta\begin{pmatrix}0\\1\end{pmatrix} \;=\; \begin{pmatrix}\alpha\\0\end{pmatrix}+\begin{pmatrix}0\\\beta\end{pmatrix} \;=\; \begin{pmatrix}\alpha\\\beta\end{pmatrix}.\]

So the general qubit state of §1.5 has the ket spelling

\[\ket{\psi} \;=\; \alpha\ket{0}+\beta\ket{1}, \qquad |\alpha|^2+|\beta|^2=1,\]

and Chapter 1's arrow, which we shall use all through this chapter and
call \(T\) for its triangle, is

\[\ket{T} \;=\; \tfrac35\ket{0}+\tfrac45\ket{1} \;=\; \begin{pmatrix}3/5\\4/5\end{pmatrix}.\]

Nothing has been added to Chapter 1 except a way of writing. The
definition of a qubit state is unchanged, and so is every matrix.

\textbf{In the sketchbook.} There is no new element, and nothing is
redrawn. The state boxes of Chapter 1 were already labelled \(0\), \(1\)
and \(\psi\); the ket is the blackboard catching up with a habit the
sketchbook had from the start, which is to put a name in a box and let
the box's \emph{shape} say what kind of thing is named. On the
blackboard the shape is the bar and the bracket; in the sketchbook it is
the rounded box with a wire leaving on its right. The state box labelled
\(\psi\) stands for \(\ket{\psi}\), and the state boxes labelled \(0\)
and \(1\) stand for \(\ket{0}\) and \(\ket{1}\).

\textbf{\emph{Why this definition.}} Why wrap a label in a bar and a
bracket, when a plain letter \(\psi\) would name a column just as well?
Because the wrapping says \emph{which way up} the thing is. A ket is a
column. In the next section we shall need the same label on a row, and
the notation will do it by turning the wrapping round. A plain letter
cannot be turned round.

\begin{quote}
\textbf{Feynman's Notebook.}

\emph{Where the ket comes from.} The notation is Dirac's. It was
introduced in a paper whose title says exactly what it contains, ``A new
notation for quantum mechanics'', published in 1939 in the journal of
the Cambridge Philosophical Society, and it is the reason the algebraic
strand of this book is called \emph{Dirac's Blackboard}. It has long
since stopped being new. Nielsen and Chuang's \emph{Quantum Computation
and Quantum Information}, the textbook our blackboard leans on, develops
in its second chapter the linear algebra a quantum computer scientist
needs, and writes it in kets throughout; so does nearly everything else
the reader will meet.

It is also sturdier than a mere abbreviation. When Abramsky and Coecke
rebuilt the formalism of quantum mechanics from abstract parts in 2004,
one of the things they checked (in their §7.2, ``Bras and kets'') was
that Dirac's notation could be interpreted in the new setting. A
notation that survives having its foundations replaced was well chosen.
\end{quote}

\begin{center}\rule{0.5\linewidth}{0.5pt}\end{center}

\hypertarget{the-other-end-of-the-wire-1}{%
\section{§2.2 The Other End of the
Wire}\label{the-other-end-of-the-wire-1}}

Now to the debt of §1.9: something that closes the \emph{right-hand} end
of a wire.

Begin with the plainest case. In §1.1 we read a column as two answers,
to the questions ``is it \(0\)?'' and ``is it \(1\)?''. Picking out the
first answer is itself something a matrix can do --- not a square matrix
this time, but a \textbf{row}, a matrix with one row and two columns. A
row acts on a column by the same rule as any matrix product, along the
row and down the column, multiplying in pairs and adding; only there is
one row to run along, so the result has a single entry:

\[\begin{pmatrix}c&d\end{pmatrix}\begin{pmatrix}x\\y\end{pmatrix} \;=\; cx+dy .\]

A matrix with one row and one column is a number, and we shall treat it
as one. Now ask the general state the first question:

\[\begin{pmatrix}1&0\end{pmatrix}\begin{pmatrix}\alpha\\\beta\end{pmatrix}=1\cdot\alpha+0\cdot\beta=\alpha, \qquad\qquad \begin{pmatrix}0&1\end{pmatrix}\begin{pmatrix}\alpha\\\beta\end{pmatrix}=\beta .\]

For the arrow \(T\) the two results are \(\tfrac35\) and \(\tfrac45\).
The state went in; a number came out; and there is nothing left for any
further matrix to act on. That is what closing the right-hand end means.

The rows \(\begin{pmatrix}1&0\end{pmatrix}\) and
\(\begin{pmatrix}0&1\end{pmatrix}\) are the columns \(\ket{0}\) and
\(\ket{1}\) laid on their sides. Could \emph{any} state be laid on its
side and used as a question --- ``how much of you lies along \(\phi\)?''
Let us try it, and test the attempt in the one case where we know what
the answer ought to be. A state asked about \emph{itself} should answer
as emphatically as anything can. For \(T\), laid on its side and applied
to itself:

\[\begin{pmatrix}\tfrac35&\tfrac45\end{pmatrix}\begin{pmatrix}3/5\\4/5\end{pmatrix}=\tfrac9{25}+\tfrac{16}{25}=1 .\]

Good. But Chapter 1 planted a second example for exactly this moment,
the state \(\begin{pmatrix}3/5\\4i/5\end{pmatrix}\) with an entry that
is not real. Laid on its side and applied to itself it gives
\(\tfrac9{25}+\bigl(\tfrac{4i}5\bigr)^2=\tfrac9{25}-\tfrac{16}{25}=-\tfrac7{25}\).
A legal state, asked about itself, answers with a negative number
smaller than one. Something is wrong with the recipe, and Chapter 1's
rule says what: the rule was never about the squares of the entries. It
was about their squared \emph{moduli}.

The repair needs one piece of school algebra. The \textbf{complex
conjugate} of \(z=a+bi\) (with \(a\), \(b\) real) is \(\bar z=a-bi\). A
real number is its own conjugate. Conjugating twice does nothing;
conjugation passes through sums and products,
\(\overline{z+w}=\bar z+\bar w\) and \(\overline{zw}=\bar z\,\bar w\);
and, the fact that matters most in this chapter,

\[\bar z\,z \;=\; (a-bi)(a+bi) \;=\; a^2+b^2 \;=\; |z|^2 .\]

A number times its conjugate is its squared modulus: real, and never
negative.

\textbf{On the blackboard.} For a state
\(\ket{\phi}=\begin{pmatrix}\gamma\\\delta\end{pmatrix}\), the
\textbf{bra} \(\bra{\phi}\), read ``bra phi'', is the column laid on its
side \emph{with every entry conjugated}:

\[\bra{\phi} \;=\; \begin{pmatrix}\bar\gamma&\bar\delta\end{pmatrix}.\]

The label is the same; the wrapping has been turned round, angle bracket
on the left and bar on the right, and that is how the notation says
``row''. For real columns the conjugation is invisible:
\(\bra{0}=\begin{pmatrix}1&0\end{pmatrix}\) and
\(\bra{1}=\begin{pmatrix}0&1\end{pmatrix}\), the two rows we began with.

A bra applied to a ket is written with the two bars merged into one,
\(\braket{\phi|\psi}\), and called the \textbf{inner product} of
\(\phi\) and \(\psi\), or simply the \textbf{bracket}:

\[\braket{\phi|\psi} \;=\; \begin{pmatrix}\bar\gamma&\bar\delta\end{pmatrix}\begin{pmatrix}\alpha\\\beta\end{pmatrix} \;=\; \bar\gamma\,\alpha+\bar\delta\,\beta .\]

It is a number, in general a complex one. We call it the
\textbf{amplitude} of \(\psi\) for \(\phi\): what the state \(\psi\)
answers when asked \(\phi\). In particular

\[\braket{0|\psi}=\alpha, \qquad\qquad \braket{1|\psi}=\beta,\]

so the entries of a column, which Chapter 1 could not name, are its
amplitudes for \(0\) and for \(1\).

Exchange the two labels and the bracket is conjugated:

\[\braket{\psi|\phi} \;=\; \bar\alpha\gamma+\bar\beta\delta \;=\; \overline{\alpha\bar\gamma+\beta\bar\delta} \;=\; \overline{\braket{\phi|\psi}} .\]

Asking \(\psi\) about \(\phi\) and asking \(\phi\) about \(\psi\) give
mirror-twin numbers. We shall use this on almost every page.

\begin{quote}
\textbf{Dirac's Blackboard.}

\emph{Why the conjugate.} Take
\(\ket{\chi}=\begin{pmatrix}3/5\\4i/5\end{pmatrix}\). It is a state:
\(|3/5|^2+|4i/5|^2=\tfrac9{25}+\tfrac{16}{25}=1\). Its bra is
\(\bra{\chi}=\begin{pmatrix}\tfrac35&-\tfrac{4i}5\end{pmatrix}\), and

\[\braket{\chi|\chi}=\tfrac35\cdot\tfrac35+\Bigl(-\tfrac{4i}5\Bigr)\Bigl(\tfrac{4i}5\Bigr)=\tfrac9{25}-\tfrac{16\,i^2}{25}=\tfrac9{25}+\tfrac{16}{25}=1 .\]

Without the conjugate the second term is
\(\bigl(\tfrac{4i}5\bigr)^2=-\tfrac{16}{25}\) and the total is
\(-\tfrac7{25}\). The plain row is not wrong as a matrix; it is the
wrong question.

The twin rule, on numbers. With
\(\ket{T}=\begin{pmatrix}3/5\\4/5\end{pmatrix}\):

\[\braket{\chi|T}=\tfrac35\cdot\tfrac35+\Bigl(-\tfrac{4i}5\Bigr)\tfrac45=\tfrac{9-16i}{25}, \qquad \braket{T|\chi}=\tfrac35\cdot\tfrac35+\tfrac45\cdot\tfrac{4i}5=\tfrac{9+16i}{25} .\]

Conjugates, as promised: the real parts agree and the imaginary parts
are opposite, which is all that conjugation ever does. Their product is
\(\tfrac{81+256}{625}=\tfrac{337}{625}\): real, and not negative. Keep
an eye on products of that kind.
\end{quote}

\textbf{In the sketchbook.} One new element, described before it is
drawn.

An \textbf{effect} is a box with one wire coming \emph{in}, on its left,
and \textbf{no wire going out}. It is the mirror image of a state. Hold
a mirror upright at the right of a state box: the box that had nothing
on its left and a wire leaving on its right becomes a box with a wire
arriving on its left and nothing on its right. A state is a way of
starting; an effect is a way of ending --- a question put to whatever
arrives.

We draw an effect with the same rounded ends as a state, and the reader
should be warned how to tell them apart, because the outline of the box
itself is symmetric and looks the same in the mirror. \textbf{Look for
the wire.} A rounded box with its wire on the \emph{right} is a state; a
rounded box with its wire on the \emph{left} is an effect. It is the
same box, met from the other side.

\[\text{a state:}\;\;
\begin{tikzpicture}[sd, baseline=-0.55ex]
  \draw[wire] (0,0) -- (3,0);
  \node[sdstate, text height=1.5ex, text depth=0.5ex] at (0,0) {$\phi$};
\end{tikzpicture}
\qquad\qquad
\text{its mirror image, an effect:}\;\;
\begin{tikzpicture}[sd, baseline=-0.55ex]
  \draw[wire] (0,0) -- (3,0);
  \node[sdeffect, text height=1.5ex, text depth=0.5ex] at (3,0) {$\phi$};
\end{tikzpicture}\]

The effect labelled \(\phi\) stands for the bra \(\bra{\phi}\). So the
mirror does one thing that cannot be seen in the drawing: \textbf{the
mirror conjugates.} The state \(\phi\) carries the entries
\(\gamma,\delta\) in a column; its mirror image carries
\(\bar\gamma,\bar\delta\) in a row.

Now put a state at the left end of a wire and an effect at the right
end. Nothing comes in and nothing goes out: the picture has no open end
anywhere. We call it a \textbf{closed diagram}. By the rule of §1.3 a
row of things on a wire is the matrix product of those things, written
in the opposite order; here the things are a column and a row, and their
product is a number. \textbf{A closed diagram is a number:}

\[\begin{tikzpicture}[sd, baseline=-0.55ex]
  \draw[wire] (0,0) -- (4,0);
  \node[sdstate, text height=1.5ex, text depth=0.5ex] at (0,0) {$\psi$};
  \node[sdeffect, text height=1.5ex, text depth=0.5ex] at (4,0) {$\phi$};
\end{tikzpicture}
\;\;=\;\; \braket{\phi|\psi} .\]

Here is Chapter 1's loud convention again, and it deserves to be loud
again. \textbf{In the picture the effect is on the right. In the bracket
the bra is on the left.} The picture is read left to right --- first the
state \(\psi\), then the question \(\phi\) --- and the bracket
\(\braket{\phi|\psi}\) is a matrix product, read right to left. They say
the same thing. The two calculations that opened this section are now
two picture equations:

\[\begin{tikzpicture}[sd, baseline=-0.55ex]
  \draw[wire] (0,0) -- (4,0);
  \node[sdstate, text height=1.5ex, text depth=0.5ex] at (0,0) {$\psi$};
  \node[sdeffect, text height=1.5ex, text depth=0.5ex] at (4,0) {$0$};
\end{tikzpicture}
\;=\; \alpha
\qquad\qquad
\begin{tikzpicture}[sd, baseline=-0.55ex]
  \draw[wire] (0,0) -- (4,0);
  \node[sdstate, text height=1.5ex, text depth=0.5ex] at (0,0) {$\psi$};
  \node[sdeffect, text height=1.5ex, text depth=0.5ex] at (4,0) {$1$};
\end{tikzpicture}
\;=\; \beta\]

And the twin rule \(\braket{\psi|\phi}=\overline{\braket{\phi|\psi}}\)
is what happens when the mirror is turned on a whole closed diagram: the
state \(\psi\) on the left becomes the effect \(\psi\) on the right, the
effect \(\phi\) on the right becomes the state \(\phi\) on the left, and
the number is conjugated.

\[\begin{tikzpicture}[sd, baseline=-0.55ex]
  \draw[wire] (0,0) -- (4,0);
  \node[sdstate, text height=1.5ex, text depth=0.5ex] at (0,0) {$\phi$};
  \node[sdeffect, text height=1.5ex, text depth=0.5ex] at (4,0) {$\psi$};
\end{tikzpicture}
\;\;=\;\; \braket{\psi|\phi} \;\;=\;\; \overline{\braket{\phi|\psi}}\]

Every picture equation in this chapter is exact: the same number on both
sides, with nothing swept aside.

\begin{quote}
\textbf{Penrose's Sketchbook.}

\emph{The mirror, stated carefully.} Three things the mirror does, and
one it does not yet do. It turns a state into an effect with the same
label, and an effect back into the state: two reflections are none, as
two conjugations are none. It conjugates the entries, which no drawing
shows; the label \(\phi\) on an effect always names the \emph{state}
whose mirror image it is, so its row is \(\bar\gamma,\bar\delta\) and
not \(\gamma,\delta\). And turned on a closed diagram made of one state,
one wire and one effect, it reverses the order, mirrors each end, and
conjugates the number.

What it does not do, in this chapter, is reflect a \emph{box}. We have
said nothing about the mirror image of NOT, or of any process with a
wire on both sides, and nothing here should be read as a rule for
reflecting pictures in general. That is Chapter 3's subject.
\end{quote}

The test that broke the plain row now passes, and not only for our two
examples. It is the headline theorem of the chapter:

\begin{quote}
\textbf{Theorem 2.1 (a state meets its own mirror image in the number
\(1\)).} For every qubit state
\(\ket{\psi}=\alpha\ket{0}+\beta\ket{1}\), \(\;\braket{\psi|\psi}=1\),
exactly. In pictures: the state \(\psi\), a wire, and the effect
\(\psi\) make a closed diagram, and that closed diagram is the number
\(1\).
\end{quote}

\[\begin{tikzpicture}[sd, baseline=-0.55ex]
  \draw[wire] (0,0) -- (4,0);
  \node[sdstate, text height=1.5ex, text depth=0.5ex] at (0,0) {$\psi$};
  \node[sdeffect, text height=1.5ex, text depth=0.5ex] at (4,0) {$\psi$};
\end{tikzpicture}
\;\;=\;\; 1\]

\emph{Proof idea.} The bra was built so that each entry meets its own
conjugate, and an entry times its conjugate is its squared modulus; the
squared moduli of a qubit state sum to \(1\) by definition. §2.8 proves
it twice, in full, and draws from it a consequence that is not a
restated definition. See \texttt{Code/Ch2/Closing.py} for the bra, the
effect and the closed diagram, including the check that the
\emph{un}conjugated mirror fails.

\textbf{\emph{Why this definition.}} Why should the mirror conjugate,
when simply laying the column on its side is simpler? Because only the
conjugated row makes a state meet itself in \(|\alpha|^2+|\beta|^2\),
the one quantity Chapter 1's rule speaks about. With the plain row, the
state \(\ket{\chi}\) of the sidebar meets itself in \(-\tfrac7{25}\),
and the closed diagram of Theorem 2.1 would be \(1\) for some states and
not for others.

\begin{center}\rule{0.5\linewidth}{0.5pt}\end{center}

\hypertarget{why-squares}{%
\section{§2.3 Why Squares}\label{why-squares}}

We have a number, and we must now say what it has to do with anything
that happens.

The first guess is the natural one: the amplitudes are the chances. The
arrow \(T\) answers \(\tfrac35\) when asked \(0\) and \(\tfrac45\) when
asked \(1\); perhaps those are the shares of heads and tails. They
cannot be. A pair of shares must make one whole, and
\(\tfrac35+\tfrac45=\tfrac75\). Shares are never negative, and
amplitudes may be (§2.4 has one). Shares are real numbers, and the
amplitudes of \(\ket{\chi}\) are not. An amplitude is not a chance of
anything.

But we met, a page ago, a way of making a real, non-negative number out
of any complex one: multiply it by its conjugate. And Chapter 1's rule
says that for the two amplitudes of a state, the numbers made this way
sum to exactly \(1\). They are real, they are not negative, they make
one whole. They have every property a pair of shares must have.

\textbf{On the blackboard.} That the world actually uses them so is not
something we can prove. It is a postulate --- in Nielsen and Chuang's
presentation (\emph{Quantum Computation and Quantum Information}, ch.~2)
it stands among the postulates of quantum mechanics --- and it is called
the \textbf{Born rule}:

\begin{quote}
\textbf{The Born rule (postulate).} When a qubit in the state
\(\ket{\psi}\) is asked a question one of whose possible answers is
\(b\), the answer \(b\) comes with probability

\[p(b) \;=\; |\braket{b|\psi}|^2 .\]
\end{quote}

For the moment ``a question'' means the one question we have, with the
two answers \(0\) and \(1\); what other questions there are is the
business of §2.5. Since \(|z|^2=\bar z z\), and since the conjugate of a
bracket is the bracket with its labels exchanged (§2.2), the rule has a
second spelling that we shall find more useful than the first:

\[p(b) \;=\; \overline{\braket{b|\psi}}\;\braket{b|\psi} \;=\; \braket{\psi|b}\,\braket{b|\psi} .\]

\textbf{A probability is an amplitude times its mirror twin.} For the
general state the two values are \(|\alpha|^2\) and \(|\beta|^2\). For
the arrow:

\[p(0)=\bigl(\tfrac35\bigr)^2=\tfrac9{25}, \qquad p(1)=\bigl(\tfrac45\bigr)^2=\tfrac{16}{25}, \qquad \tfrac9{25}+\tfrac{16}{25}=1 .\]

Look at what has been produced. The column
\(\begin{pmatrix}9/25\\16/25\end{pmatrix}\) has real entries, not
negative, summing to \(1\): it is a \emph{stochastic vector}, a covered
coin in the exact sense of §1.4. So the Born rule is a bridge between
the two rules of Chapter 1's table. It takes a qubit state \textbf{and a
question}, and hands back a wager. The normalization
\(|\alpha|^2+|\beta|^2=1\), which Chapter 1 could only state, now reads:
\emph{the chances make one whole.}

\textbf{In the sketchbook.} The second spelling of the rule is a product
of two brackets, and each bracket is a closed diagram. We need to say
what two closed diagrams \emph{in a row} mean, and the answer is already
in our hands. A row is a matrix product whatever the shapes of the
matrices in it (§1.3); in the row ``state \(\psi\), effect \(b\), state
\(b\), effect \(\psi\)'' the product, read right to left, is
\(\bra{\psi}\,\ket{b}\,\bra{b}\,\ket{\psi}\), and since matrix
multiplication may be grouped as we please, this is
\(\bigl(\bra{\psi}\ket{b}\bigr)\bigl(\bra{b}\ket{\psi}\bigr)\): the
product of two numbers. So \textbf{closed diagrams in a row multiply.}
This is a consequence of the row rule, worked out on the blackboard, and
not a new rule of drawing; in particular, where an effect is followed by
a state, \emph{no wire joins them}. The first wire has ended. The second
has not yet begun.

Here, then, is the Born rule as a picture. We shall meet it on most of
the remaining pages, so it gets a name: \textbf{\emph{The Echo}}.

\[\begin{tikzpicture}[sd, baseline=-0.55ex]
  \draw[wire] (0,0) -- (4,0);
  \draw[wire] (8,0) -- (12,0);
  \node[sdstate, text height=1.5ex, text depth=0.5ex] at (0,0) {$\psi$};
  \node[sdeffect, text height=1.5ex, text depth=0.5ex] at (4,0) {$b$};
  \node[sdstate, text height=1.5ex, text depth=0.5ex] at (8,0) {$b$};
  \node[sdeffect, text height=1.5ex, text depth=0.5ex] at (12,0) {$\psi$};
\end{tikzpicture}
\;\;=\;\; \braket{\psi|b}\,\braket{b|\psi} \;\;=\;\; |\braket{b|\psi}|^2\]

Read it from the left, looking for the wires as §2.2 advised. The state
\(\psi\), with its wire on the right; that wire ends in the effect
\(b\). Then a gap: bare paper, no wire. Then the state \(b\), whose wire
ends in the effect \(\psi\). The left half says ``\(\psi\), asked
\(b\)'', and its number is the amplitude \(\braket{b|\psi}\). The right
half is the left half seen in the mirror, every piece reflected and the
order reversed, and its number is the conjugate, \(\braket{\psi|b}\).
The call and its echo; and the chance is the two together.

This is why a Born value is real and never negative, in the sketchbook's
terms: it is a closed diagram followed by its own mirror image, a number
next to its twin. The amplitude alone may have a sign, or an \(i\) in
it. The amplitude \emph{with its echo} cannot.

Let us be plain about what has not been claimed. The picture does not
explain why nature computes chances this way, and it does not derive the
rule; the rule is a postulate, on the blackboard and in the sketchbook
alike. What the picture does is \emph{locate} the square. The question
``why squares?'' with which Chapter 1 ended has become a question one
can point at: \emph{why is there a second wire?} We shall say in §2.8
how much of an answer this chapter can honestly give. See
\texttt{Code/Ch2/Squares.py}.

\textbf{\emph{Why this definition.}} Why \(\bar z z\), and not the
amplitude itself, or its plain square \(z^2\)? The amplitude itself
fails every test for a share, as we saw. The plain square fails for the
same reason the plain row failed in §2.2: for \(\ket{\chi}\) it would
give \(p(1)=\bigl(\tfrac{4i}5\bigr)^2=-\tfrac{16}{25}\), a negative
chance. The product of a number with its conjugate is the simplest
expression that is real and non-negative for \emph{every} amplitude.

\begin{quote}
\textbf{Feynman's Notebook.}

\emph{How a probability is met in a laboratory.} Nobody has ever read an
amplitude off an instrument. What a laboratory can do is prepare a
system in the same way a great many times, put the same two-answer
question to it each time --- Nielsen and Chuang call this step a
measurement, and give it a postulate of its own in their second chapter
--- and count: so many runs answered one way, so many the other. A
single run yields a single answer, never a number like \(\tfrac35\), and
certainly never one with a minus sign or an \(i\) in it.

The Born rule is a claim about the counts: over a long series of runs,
the fraction answering \(b\) settles near \(|\braket{b|\psi}|^2\). It is
a postulate, not a deduction from anything more primitive, and it is
kept because the counting agrees with it. That is the whole of its
authority, and it is a great deal. The amplitudes belong to the night,
when nothing is looked at; the counts are what is there in the morning.
\end{quote}

\begin{center}\rule{0.5\linewidth}{0.5pt}\end{center}

\hypertarget{a-sign-you-cannot-see}{%
\section{§2.4 A Sign You Cannot See}\label{a-sign-you-cannot-see}}

Now, at last, the negative entry that Chapter 1 mentioned and set aside.

\textbf{On the blackboard.} Two states get names of their own, a plus
sign and a minus sign inside a ket:

\[\ket{+} \;=\; \tfrac1{\sqrt2}\begin{pmatrix}1\\1\end{pmatrix}, \qquad\qquad \ket{-} \;=\; \tfrac1{\sqrt2}\begin{pmatrix}1\\-1\end{pmatrix}.\]

Both are legal:
\(\bigl(\tfrac1{\sqrt2}\bigr)^2+\bigl(\pm\tfrac1{\sqrt2}\bigr)^2=\tfrac12+\tfrac12=1\).
As arrows on paper, \(\ket{+}\) reaches equally toward \(0\) and toward
\(1\); \(\ket{-}\) reaches toward \(0\) and \emph{backwards}, away from
\(1\). The amplitude \(\braket{1|-}=-\tfrac1{\sqrt2}\) is a reach of
less than nothing.

Ask each of them the standard question:

\[\ket{+}:\qquad p(0)=\bigl|\tfrac1{\sqrt2}\bigr|^2=\tfrac12, \qquad p(1)=\bigl|\tfrac1{\sqrt2}\bigr|^2=\tfrac12;\]

\[\ket{-}:\qquad p(0)=\bigl|\tfrac1{\sqrt2}\bigr|^2=\tfrac12, \qquad p(1)=\bigl|-\tfrac1{\sqrt2}\bigr|^2=\tfrac12 .\]

The same wager from both: the fair covered coin of §1.4. The square has
eaten the sign. And this is general. The states
\(\alpha\ket{0}+\beta\ket{1}\) and \(\alpha\ket{0}-\beta\ket{1}\) give
the standard Born values \(|\alpha|^2\) and \(|{\pm}\beta|^2=|\beta|^2\)
alike.

Yet they are different columns --- and not slightly different. Ask one
about the other:

\[\braket{+|-} \;=\; \tfrac1{\sqrt2}\cdot\tfrac1{\sqrt2}\bigl(1\cdot1+1\cdot(-1)\bigr) \;=\; \tfrac12\,(1-1) \;=\; 0,\]

whereas \(\braket{+|+}=\tfrac12\,(1+1)=1\).

For real arrows a bracket equal to \(0\) has a familiar meaning: the
arrows are \textbf{perpendicular}; neither has any reach along the
other. Draw \((1,1)\) and \((1,-1)\) on squared paper and they stand at
right angles. We take the word over for all states: \(\ket{\phi}\) and
\(\ket{\psi}\) are perpendicular when \(\braket{\phi|\psi}=0\). By that
test \(\ket{+}\) and \(\ket{-}\) are as different as two states can be,
exactly as different as \(\ket{0}\) and \(\ket{1}\), for which
\(\braket{0|1}=1\cdot0+0\cdot1=0\) too. (Chapter 5 will read a zero like
this one in quite another way; here it is plain geometry.)

\textbf{In the sketchbook.} No new element. Two \emph{Echo} pictures
with different labels in their outer boxes, and the same number:

\[\begin{tikzpicture}[sd, baseline=-0.55ex]
  \draw[wire] (0,0) -- (4,0);
  \draw[wire] (8,0) -- (12,0);
  \node[sdstate, text height=1.5ex, text depth=0.5ex] at (0,0) {$+$};
  \node[sdeffect, text height=1.5ex, text depth=0.5ex] at (4,0) {$1$};
  \node[sdstate, text height=1.5ex, text depth=0.5ex] at (8,0) {$1$};
  \node[sdeffect, text height=1.5ex, text depth=0.5ex] at (12,0) {$+$};
\end{tikzpicture}
\;\;=\;\; \tfrac1{\sqrt2}\cdot\tfrac1{\sqrt2} \;\;=\;\; \tfrac12\]

\[\begin{tikzpicture}[sd, baseline=-0.55ex]
  \draw[wire] (0,0) -- (4,0);
  \draw[wire] (8,0) -- (12,0);
  \node[sdstate, text height=1.5ex, text depth=0.5ex] at (0,0) {$-$};
  \node[sdeffect, text height=1.5ex, text depth=0.5ex] at (4,0) {$1$};
  \node[sdstate, text height=1.5ex, text depth=0.5ex] at (8,0) {$1$};
  \node[sdeffect, text height=1.5ex, text depth=0.5ex] at (12,0) {$-$};
\end{tikzpicture}
\;\;=\;\; \bigl(-\tfrac1{\sqrt2}\bigr)\bigl(-\tfrac1{\sqrt2}\bigr) \;\;=\;\; \tfrac12\]

In the second picture each half is negative, and the echo repeats the
sign faithfully; the product cannot show it. And one closed diagram
whose number is zero:

\[\begin{tikzpicture}[sd, baseline=-0.55ex]
  \draw[wire] (0,0) -- (4,0);
  \node[sdstate, text height=1.5ex, text depth=0.5ex] at (0,0) {$-$};
  \node[sdeffect, text height=1.5ex, text depth=0.5ex] at (4,0) {$+$};
\end{tikzpicture}
\;\;=\;\; \braket{+|-} \;\;=\;\; 0\]

\hypertarget{a-trap-1}{%
\subsection{A trap}\label{a-trap-1}}

Here is the tempting conclusion. Everything that anyone will ever
\emph{count} comes through the Born rule; the Born rule gives
\(\ket{+}\) and \(\ket{-}\) the same wager; so the minus sign is
bookkeeping, a mark on the blackboard with no echo in the world, and the
two states might as well be called one.

Every step of that is sound except the word ``everything''. We have put
\emph{one question} to these two states. §1.4 taught that one state is
not enough to tell two boxes apart. It would be rash to assume that one
question is enough to tell two states apart --- and the bracket
\(\braket{+|-}=0\) is a warning that it is not. See
\texttt{Code/Ch2/HiddenSign.py}.

\begin{center}\rule{0.5\linewidth}{0.5pt}\end{center}

\hypertarget{the-turned-question}{%
\section{§2.5 The Turned Question}\label{the-turned-question}}

What would another question be? The standard question has two answers,
\(0\) and \(1\), and each answer is an effect: the mirror image of a
state. Nothing in the Born rule singles those two states out. What we
need is to say which \emph{pairs} of states can serve as the two answers
of a question.

\textbf{On the blackboard.} An \textbf{orthonormal basis} is a pair of
states \(\ket{b_0},\ket{b_1}\) that are perpendicular to one another:

\[\braket{b_0|b_0}=1, \qquad \braket{b_1|b_1}=1, \qquad \braket{b_0|b_1}=0\]

(and then \(\braket{b_1|b_0}=\bar 0=0\) as well). The first two
conditions say only that both are states; the third is the new one. The
pair \(\ket{0},\ket{1}\) is an orthonormal basis, the \textbf{standard
basis}. By the calculation of §2.4, so is the pair \(\ket{+},\ket{-}\).
On paper it is the standard pair of directions turned through half a
right angle.

From now on \textbf{a question is an orthonormal basis}, its two answers
are the two members, and the Born rule of §2.3 is complete as it stands:
asked the question \(\{b_0,b_1\}\), the state \(\ket{\psi}\) answers
\(b_0\) with probability \(|\braket{b_0|\psi}|^2\) and \(b_1\) with
probability \(|\braket{b_1|\psi}|^2\).

Put the turned question to the two states of the trap:

\[\ket{+}:\qquad p(+)=|\braket{+|+}|^2=1, \qquad p(-)=|\braket{-|+}|^2=0;\]

\[\ket{-}:\qquad p(+)=|\braket{+|-}|^2=0, \qquad p(-)=|\braket{-|-}|^2=1 .\]

Certainty, and opposite certainties. The two states that one question
could not tell apart at all, another tells apart every single time. The
sign was never bookkeeping. It was the whole difference, waiting for the
right question.

The arrow \(T\) under the turned question gives two less tidy numbers,
worth having:

\[\braket{+|T}=\tfrac1{\sqrt2}\bigl(\tfrac35+\tfrac45\bigr)=\tfrac{7}{5\sqrt2}, \qquad \braket{-|T}=\tfrac1{\sqrt2}\bigl(\tfrac35-\tfrac45\bigr)=-\tfrac{1}{5\sqrt2},\]

\[p(+)=\tfrac{49}{50}, \qquad p(-)=\tfrac1{50} .\]

All but a sliver --- and the sliver comes from a negative amplitude. The
conjugate has its turned question too: \(T\) and the non-real state
\(\ket{\chi}\) of §2.2 give the same standard values
\(\tfrac9{25},\tfrac{16}{25}\), but
\(\braket{+|\chi}=\tfrac{3+4i}{5\sqrt2}\), whose squared modulus is
\(\tfrac{9+16}{50}=\tfrac12\), against \(T\)'s \(\tfrac{49}{50}\).

The word \emph{basis} makes a promise, which we now keep: every state
can be written in terms of the two members, and the coefficients are its
amplitudes.

\begin{quote}
\textbf{Lemma 2.3 (expansion in an orthonormal basis).} If
\(\ket{b_0},\ket{b_1}\) is an orthonormal basis, then for every state
\(\ket{\psi}\)

\[\ket{\psi} \;=\; \braket{b_0|\psi}\,\ket{b_0}+\braket{b_1|\psi}\,\ket{b_1} .\]
\end{quote}

For the standard basis this is the first display of §2.1, since
\(\braket{0|\psi}=\alpha\) and \(\braket{1|\psi}=\beta\). For the turned
basis and the arrow \(T\) it is checked by hand in the sidebar below.
\emph{Proof idea:} subtract from \(\ket{\psi}\) its reach along each
member; what is left over has no reach along either, and in a space with
only two places that leaves nothing. \emph{Proof sketch.} Write
\(c_0=\braket{b_0|\psi}\), \(c_1=\braket{b_1|\psi}\), and let \(r\) be
the leftover column \(\ket{\psi}-c_0\ket{b_0}-c_1\ket{b_1}\). A row
times a sum of columns is the sum of the products, so \(\bra{b_0}\)
applied to \(r\) gives \(c_0-c_0\cdot1-c_1\cdot0=0\), and likewise
\(\bra{b_1}\) gives \(0\). Stack the two bras into a \(2\times2\) matrix
\(B\), and set the two kets side by side as the columns of a
\(2\times2\) matrix \(K\). The four entries of \(BK\) are the four
brackets \(\braket{b_i|b_j}\), so orthonormality says precisely
\(BK=I\). For square matrices a one-sided inverse is two-sided, so
\(KB=I\) too, and \(r=K(Br)=K\begin{pmatrix}0\\0\end{pmatrix}\) is the
zero column. \(\;\blacksquare\) See \texttt{Code/Ch2/TurnedQuestion.py}.

\begin{quote}
\textbf{Dirac's Blackboard.}

\emph{The turned basis, by brute force.} Orthonormal:
\(\braket{+|+}=\tfrac12(1+1)=1\), \(\braket{-|-}=\tfrac12(1+1)=1\),
\(\braket{+|-}=\tfrac12(1-1)=0\).

The expansion of \(T\). With \(c_+=\tfrac7{5\sqrt2}\) and
\(c_-=-\tfrac1{5\sqrt2}\),

\[c_+\ket{+}+c_-\ket{-}=\tfrac{7}{10}\begin{pmatrix}1\\1\end{pmatrix}-\tfrac{1}{10}\begin{pmatrix}1\\-1\end{pmatrix}=\tfrac1{10}\begin{pmatrix}6\\8\end{pmatrix}=\begin{pmatrix}3/5\\4/5\end{pmatrix}.\]

The Born values. \(c_+^2=\tfrac{49}{50}\), \(c_-^2=\tfrac1{50}\), and

\[\tfrac{49}{50}+\tfrac1{50}=1 .\]

Both amplitudes are real, so each squared modulus is a plain square, and
the minus sign on \(c_-\) disappears in it.

Compare the standard question: \(\tfrac9{25}+\tfrac{16}{25}=1\). Two
different wagers from one state, and each makes one whole. Nothing in
the definition of a state mentioned the turned basis. Why should the
total come out right there as well? That is Corollary 2.2, in §2.8.
\end{quote}

\textbf{In the sketchbook.} No new element, and that is the point worth
drawing. The effects labelled \(+\) and \(-\) are mirrored states like
any others, and \emph{The Echo} does not care which labels it carries:

\[\begin{tikzpicture}[sd, baseline=-0.55ex]
  \draw[wire] (0,0) -- (4,0);
  \draw[wire] (8,0) -- (12,0);
  \node[sdstate, text height=1.5ex, text depth=0.5ex] at (0,0) {$+$};
  \node[sdeffect, text height=1.5ex, text depth=0.5ex] at (4,0) {$+$};
  \node[sdstate, text height=1.5ex, text depth=0.5ex] at (8,0) {$+$};
  \node[sdeffect, text height=1.5ex, text depth=0.5ex] at (12,0) {$+$};
\end{tikzpicture}
\;\;=\;\; 1\cdot1 \;\;=\;\; 1\]

\[\begin{tikzpicture}[sd, baseline=-0.55ex]
  \draw[wire] (0,0) -- (4,0);
  \draw[wire] (8,0) -- (12,0);
  \node[sdstate, text height=1.5ex, text depth=0.5ex] at (0,0) {$-$};
  \node[sdeffect, text height=1.5ex, text depth=0.5ex] at (4,0) {$+$};
  \node[sdstate, text height=1.5ex, text depth=0.5ex] at (8,0) {$+$};
  \node[sdeffect, text height=1.5ex, text depth=0.5ex] at (12,0) {$-$};
\end{tikzpicture}
\;\;=\;\; 0\cdot0 \;\;=\;\; 0\]

Set these beside the two pictures of §2.4. There the outer labels \(+\)
and \(-\) made no difference to the number; here they make all the
difference there can be.

The expansion of Lemma 2.3 stays on the blackboard. It is a \emph{sum}
of two columns, and the sketchbook does not yet have a way to add two
pictures that have open wires; we shall meet that tool in Chapter 5.
What the sketchbook may do is add \emph{numbers}, and a closed diagram
is a number. We shall use that liberty once, in §2.8, and say so when we
do.

\textbf{\emph{Why this definition.}} Why must the answers of a question
be perpendicular states, rather than any two different states? Because
answers must not be confusable. If the qubit \emph{is} in the state
\(\ket{b_0}\), the question ought to return \(b_0\) with certainty and
\(b_1\) never, and the Born rule makes those two demands into
\(|\braket{b_0|b_0}|^2=1\) and \(|\braket{b_1|b_0}|^2=0\). That is
orthonormality, no more and no less.

\begin{center}\rule{0.5\linewidth}{0.5pt}\end{center}

\hypertarget{the-arrow-reversed}{%
\section{§2.6 The Arrow Reversed}\label{the-arrow-reversed}}

One might now over-learn the lesson: every sign is caught by \emph{some}
question, if the question is turned far enough. Let us test that on the
boldest sign of all, the one in front of the whole column.

Reverse the arrow end for end:
\(-\ket{T}=\begin{pmatrix}-3/5\\-4/5\end{pmatrix}\). It is a state,
since squaring removes the signs, and it is a different column from
\(\ket{T}\). Every amplitude changes sign: \(-\tfrac35\) and
\(-\tfrac45\) for the standard question, \(-\tfrac7{5\sqrt2}\) and
\(+\tfrac1{5\sqrt2}\) for the turned one. And every Born value stays
where it was: \(\tfrac9{25},\tfrac{16}{25}\);
\(\tfrac{49}{50},\tfrac1{50}\). Turn the question to any position you
like, and the same will happen, because the bracket of \emph{any} bra
with \(-\ket{T}\) is minus its bracket with \(\ket{T}\), and the square
does not see the minus.

\textbf{On the blackboard.} The general form of this needs the complex
numbers of modulus \(1\). For a real number \(\theta\), write
\(e^{i\theta}=\cos\theta+i\sin\theta\). Its squared modulus is
\(\cos^2\theta+\sin^2\theta=1\); its conjugate is
\(\cos\theta-i\sin\theta=e^{-i\theta}\); and
\(e^{i\theta}e^{-i\theta}=1\). With \(\theta=\pi\) it is \(-1\). Such a
number is called a \textbf{phase}, and when it multiplies a whole
column, a \textbf{global phase}.

\begin{quote}
\textbf{Proposition 2.4 (a global phase is seen by no question).} Let
\(\ket{\psi}\) be a state, \(\theta\) real, and
\(\ket{\psi'}=e^{i\theta}\ket{\psi}\). Then \(\ket{\psi'}\) is a state,
it is a different column from \(\ket{\psi}\) whenever
\(e^{i\theta}\neq1\), and for \textbf{every} state \(\ket{b}\)

\[|\braket{b|\psi'}|^2 \;=\; |\braket{b|\psi}|^2 .\]
\end{quote}

\emph{Proof idea.} The factor comes out of the amplitude as
\(e^{i\theta}\) and out of its mirror twin as \(e^{-i\theta}\), because
the mirror conjugates; and the two meet in the product. \emph{Proof
sketch.} The entries of \(\ket{\psi'}\) are \(e^{i\theta}\alpha\) and
\(e^{i\theta}\beta\), with squared moduli \(|\alpha|^2\) and
\(|\beta|^2\), so \(\ket{\psi'}\) is a state. A row applied to a
multiple of a column is the multiple of the result:
\(\braket{b|\psi'}=e^{i\theta}\braket{b|\psi}\). The bra of
\(\ket{\psi'}\) has the conjugated entries
\(e^{-i\theta}\bar\alpha,\,e^{-i\theta}\bar\beta\), that is,
\(\bra{\psi'}=e^{-i\theta}\bra{\psi}\), so
\(\braket{\psi'|b}=e^{-i\theta}\braket{\psi|b}\). Therefore

\[\braket{\psi'|b}\,\braket{b|\psi'} \;=\; e^{-i\theta}\,e^{i\theta}\;\braket{\psi|b}\,\braket{b|\psi} \;=\; \braket{\psi|b}\,\braket{b|\psi} . \qquad\blacksquare\]

See \texttt{Code/Ch2/Reversed.py}, which tries effects drawn at random
as well as our own.

\textbf{In the sketchbook.} A number may multiply a state picture, as it
may multiply a column: \(e^{i\theta}\) times the state \(\psi\) is the
state box labelled \(e^{i\theta}\psi\), and its mirror image is the
effect with the same label, whose row is \(e^{-i\theta}\bra{\psi}\). The
two state pictures are different,

\[\begin{tikzpicture}[sd, baseline=-0.55ex]
  \draw[wire] (0,0) -- (4,0);
  \node[sdstate, inner xsep=0.55em, text height=1.5ex, text depth=0.5ex] at (0,0) {$e^{i\theta}\psi$};
\end{tikzpicture}
\;\;\neq\;\;
\begin{tikzpicture}[sd, baseline=-0.55ex]
  \draw[wire] (0,0) -- (3,0);
  \node[sdstate, text height=1.5ex, text depth=0.5ex] at (0,0) {$\psi$};
\end{tikzpicture}
\qquad (e^{i\theta}\neq1),\]

and their two \emph{Echo} pictures are the same number, for every label
\(b\):

\[\begin{tikzpicture}[sd, baseline=-0.55ex]
  \draw[wire] (0,0) -- (5,0);
  \draw[wire] (9,0) -- (14,0);
  \node[sdstate, inner xsep=0.55em, text height=1.5ex, text depth=0.5ex] at (0,0) {$e^{i\theta}\psi$};
  \node[sdeffect, text height=1.5ex, text depth=0.5ex] at (5,0) {$b$};
  \node[sdstate, text height=1.5ex, text depth=0.5ex] at (9,0) {$b$};
  \node[sdeffect, inner xsep=0.55em, text height=1.5ex, text depth=0.5ex] at (14,0) {$e^{i\theta}\psi$};
\end{tikzpicture}
\;\;=\;\; e^{-i\theta}\braket{\psi|b}\;\cdot\;e^{i\theta}\braket{b|\psi}\]

\[=\;\; \braket{\psi|b}\,\braket{b|\psi} \;\;=\;\;
\begin{tikzpicture}[sd, baseline=-0.55ex]
  \draw[wire] (0,0) -- (4,0);
  \draw[wire] (8,0) -- (12,0);
  \node[sdstate, text height=1.5ex, text depth=0.5ex] at (0,0) {$\psi$};
  \node[sdeffect, text height=1.5ex, text depth=0.5ex] at (4,0) {$b$};
  \node[sdstate, text height=1.5ex, text depth=0.5ex] at (8,0) {$b$};
  \node[sdeffect, text height=1.5ex, text depth=0.5ex] at (12,0) {$\psi$};
\end{tikzpicture}\]

The factor leaves the left half as \(e^{i\theta}\) and the right half as
\(e^{-i\theta}\). Whatever the call gains, the echo gives back.

We must be careful about what this says. It does \textbf{not} say that
\(e^{i\theta}\ket{\psi}\) and \(\ket{\psi}\) are the same state, or the
same picture, in some loosened sense of ``same''. They are different
columns and different pictures; the closed diagram of one with an effect
\(b\) is a different number from that of the other. What the proposition
says is that their \emph{Born values} agree for every effect, so that no
question of the kind we have defined can tell them apart.

Set this beside §2.4. A sign --- or any phase other than \(1\) --- on
\emph{one} entry of a column whose entries are both non-zero makes a
different state, and some question gives the two different chances (for
\(\ket{+}\) and \(\ket{-}\), opposite certainties). The same phase on
\emph{both} entries makes a different column that no question detects at
all. The difference between those two cases is the most important thing
to carry out of this chapter. (We have shown it for every question asked
directly of the state. That it survives whatever boxes are placed on the
wire first is true as well, and belongs to Chapter 3.)

\begin{center}\rule{0.5\linewidth}{0.5pt}\end{center}

\hypertarget{the-rosetta-table-1}{%
\section{§2.7 The Rosetta Table}\label{the-rosetta-table-1}}

This chapter adds eight rows to the running dictionary, numbered on from
Chapter 1's seven. Rows 8 and 9 are new spellings, on the algebra side
only, of Chapter 1's third and fourth rows. Rows 10 and 11 introduce the
chapter's two new picture elements, the effect and the closed diagram.
The remaining four are things those elements say.

\begingroup\renewcommand{\arraystretch}{1.5}

\begin{longtable}[]{@{}
  >{\raggedright\arraybackslash}p{(\columnwidth - 4\tabcolsep) * \real{0.0909}}
  >{\raggedright\arraybackslash}p{(\columnwidth - 4\tabcolsep) * \real{0.4545}}
  >{\raggedright\arraybackslash}p{(\columnwidth - 4\tabcolsep) * \real{0.4545}}@{}}
\toprule\noalign{}
\begin{minipage}[b]{\linewidth}\raggedright
\#
\end{minipage} & \begin{minipage}[b]{\linewidth}\raggedright
Algebra (Dirac)
\end{minipage} & \begin{minipage}[b]{\linewidth}\raggedright
Picture (Penrose)
\end{minipage} \\
\midrule\noalign{}
\endhead
\bottomrule\noalign{}
\endlastfoot
8 & \(\ket{\psi}=\alpha\ket{0}+\beta\ket{1}\), the ket spelling of a
column & the state box labelled \(\psi\) (no new element) \\
9 & \(\ket{0}\), \(\ket{1}\) & the state boxes labelled \(0\) and \(1\)
(no new element) \\
10 & the bra
\(\bra{\phi}=\begin{pmatrix}\bar\gamma&\bar\delta\end{pmatrix}\) & an
\textbf{effect}: a rounded box with its wire on the left and no output,
the mirror image of the state \(\phi\) (the mirror conjugates);
\emph{new element} \\
11 & the bracket \(\braket{\phi\vert\psi}\), bra on the left & state
\(\psi\), wire, effect \(\phi\), effect on the right: a \textbf{closed
diagram}, which is a number --- \emph{new element} \\
12 & a product of numbers
\(\braket{\chi\vert\omega}\,\braket{\phi\vert\psi}\) & closed diagrams
in a row (the row rule, with no open end) \\
13 & the Born rule
\(\lvert\braket{b\vert\psi}\rvert^2=\braket{\psi\vert b}\,\braket{b\vert\psi}\)
& \textbf{\emph{The Echo}}: a closed diagram followed by its mirror
image \\
14 & \(\braket{\psi\vert\psi}=1\) & a state meeting its own mirror image
\(=1\) \\
15 & \(e^{i\theta}\ket{\psi}\neq\ket{\psi}\), yet
\(\lvert\braket{b\vert e^{i\theta}\psi}\rvert^2=\lvert\braket{b\vert\psi}\rvert^2\)
for every \(b\) & two different state pictures; the same \emph{Echo}
number for every effect \\
\end{longtable}

\endgroup

All eight rows are exact. In row 12 the labels \(\chi,\omega\) stand for
any two further states; in row 15 the label \(e^{i\theta}\psi\) inside a
bracket stands for the state \(e^{i\theta}\ket{\psi}\).

\begin{center}\rule{0.5\linewidth}{0.5pt}\end{center}

\hypertarget{the-theorem-twice-1}{%
\section{§2.8 The Theorem, Twice}\label{the-theorem-twice-1}}

We have met the theorem in examples three times: \(\ket{T}\) met itself
in \(\tfrac9{25}+\tfrac{16}{25}=1\) (§2.2); so did the non-real
\(\ket{\chi}\), once its row was conjugated; so did \(\ket{+}\) in §2.4.
Here is the statement for all states at once.

\begin{quote}
\textbf{Theorem 2.1 (a state meets its own mirror image in the number
\(1\)).} Let \(\ket{\psi}=\alpha\ket{0}+\beta\ket{1}\) be a qubit state:
\(\alpha,\beta\) complex with \(|\alpha|^2+|\beta|^2=1\). Let
\(\bra{\psi}=\begin{pmatrix}\bar\alpha&\bar\beta\end{pmatrix}\) be its
bra, the conjugated row. Then

\[\braket{\psi|\psi} \;=\; 1,\]

exactly. In pictures: the state \(\psi\), a wire, and the effect
\(\psi\) form a closed diagram, and that closed diagram is the number
\(1\).
\end{quote}

\hypertarget{diracs-blackboard-1}{%
\subsection{Dirac's Blackboard}\label{diracs-blackboard-1}}

Write out the row and the column and multiply, along the row and down
the column:

\[\braket{\psi|\psi} \;=\; \begin{pmatrix}\bar\alpha&\bar\beta\end{pmatrix}\begin{pmatrix}\alpha\\\beta\end{pmatrix} \;=\; \bar\alpha\,\alpha+\bar\beta\,\beta .\]

For any complex number \(z=a+bi\),
\(\;\bar z z=(a-bi)(a+bi)=a^2+b^2=|z|^2\) (§2.2). Apply this to each
term:

\[\bar\alpha\,\alpha+\bar\beta\,\beta \;=\; |\alpha|^2+|\beta|^2 .\]

And \(\ket{\psi}\) is a qubit state, which by the definition of §1.5
means

\[|\alpha|^2+|\beta|^2 \;=\; 1 .\]

One row times one column, two products, one appeal to the definition.
\(\;\blacksquare\)

\hypertarget{penroses-sketchbook-1}{%
\subsection{Penrose's Sketchbook}\label{penroses-sketchbook-1}}

The sketchbook owns three things that bear on this picture, and the
proof uses each once. We name them. \emph{Read the closed diagram} (row
11 of the table): a state, a wire and an effect are the bracket of the
effect's row with the state's column. \emph{The mirror conjugates} (row
10): the effect labelled \(\psi\) carries the row
\(\bar\alpha,\bar\beta\). \emph{Chapter 1's rule}: the label \(\psi\) in
a state box stands for a column whose squared moduli sum to \(1\).

\textbf{Step 1 --- \emph{read the closed diagram}.} The picture has no
open end, so it is a number: the row of the effect \(\psi\) times the
column of the state \(\psi\).

\[\begin{tikzpicture}[sd, baseline=-0.55ex]
  \draw[wire] (0,0) -- (4,0);
  \node[sdstate, text height=1.5ex, text depth=0.5ex] at (0,0) {$\psi$};
  \node[sdeffect, text height=1.5ex, text depth=0.5ex] at (4,0) {$\psi$};
\end{tikzpicture}
\;\;\overset{\raisebox{0pt}[\height][1.2ex]{\scriptsize read the closed diagram}}{=}\;\;
\bigl(\text{row of the effect }\psi\bigr)\begin{pmatrix}\alpha\\\beta\end{pmatrix}\]

\textbf{Step 2 --- \emph{the mirror conjugates}.} The effect \(\psi\) is
the mirror image of the state \(\psi\), so its row is the conjugated
one.

\[\bigl(\text{row of the effect }\psi\bigr)\begin{pmatrix}\alpha\\\beta\end{pmatrix}
\;\;\overset{\raisebox{0pt}[\height][1.2ex]{\scriptsize the mirror conjugates}}{=}\;\;
\begin{pmatrix}\bar\alpha&\bar\beta\end{pmatrix}\begin{pmatrix}\alpha\\\beta\end{pmatrix}
\;=\; |\alpha|^2+|\beta|^2\]

\textbf{Step 3 --- \emph{Chapter 1's rule}.} The box on the left is a
state box, so the sum of squared moduli is \(1\).

\[|\alpha|^2+|\beta|^2
\;\;\overset{\raisebox{0pt}[\height][1.2ex]{\scriptsize Chapter 1's rule}}{=}\;\; 1,
\qquad\text{and so}\qquad
\begin{tikzpicture}[sd, baseline=-0.55ex]
  \draw[wire] (0,0) -- (4,0);
  \node[sdstate, text height=1.5ex, text depth=0.5ex] at (0,0) {$\psi$};
  \node[sdeffect, text height=1.5ex, text depth=0.5ex] at (4,0) {$\psi$};
\end{tikzpicture}
\;=\; 1 . \qquad\blacksquare\]

Notice what Step 2 guards against. Had the mirror \emph{not} conjugated,
the middle expression would have been \(\alpha^2+\beta^2\), and for
\(\ket{\chi}\) the picture would have come to \(-\tfrac7{25}\).

\hypertarget{what-the-theorem-is-for}{%
\subsection{What the theorem is for}\label{what-the-theorem-is-for}}

Taken alone, Theorem 2.1 is Chapter 1's definition in a new dress, and
we should not pretend otherwise. Its worth is in what follows from it.
The sidebar of §2.5 found that the arrow's Born values make one whole
under the standard question, \(\tfrac9{25}+\tfrac{16}{25}\), and again
under the turned one, \(\tfrac{49}{50}+\tfrac1{50}\) --- though nothing
in the definition of a state mentions the turned basis. That cannot be
left as a coincidence. If the Born rule is to hand back an honest wager
for \emph{every} question, the total must be \(1\) for every orthonormal
basis, and it is.

\begin{quote}
\textbf{Corollary 2.2 (the chances make one whole, for every question).}
For every state \(\ket{\psi}\) and every orthonormal basis
\(\ket{b_0},\ket{b_1}\),

\[|\braket{b_0|\psi}|^2+|\braket{b_1|\psi}|^2 \;=\; \braket{\psi|\psi} \;=\; 1 .\]
\end{quote}

\emph{Proof idea.} Expand the ket inside \(\braket{\psi|\psi}\) in the
basis, by Lemma 2.3; the bracket breaks into two terms, and each term is
an amplitude times its mirror twin. \emph{Proof.} By Lemma 2.3,
\(\ket{\psi}=\braket{b_0|\psi}\ket{b_0}+\braket{b_1|\psi}\ket{b_1}\).
Apply the row \(\bra{\psi}\) to both sides; a row times a sum of columns
is the sum of the products, and numbers may be moved to the front:

\[\braket{\psi|\psi} \;=\; \braket{b_0|\psi}\,\braket{\psi|b_0}+\braket{b_1|\psi}\,\braket{\psi|b_1} \;=\; |\braket{b_0|\psi}|^2+|\braket{b_1|\psi}|^2,\]

using \(\braket{\psi|b}=\overline{\braket{b|\psi}}\) in the last step.
The left side is \(1\) by Theorem 2.1. \(\;\blacksquare\) See
\texttt{Code/Ch2/WholeSquare.py}, which checks the theorem, the failure
of the unconjugated mirror, and the corollary for the standard basis,
the turned basis and a basis drawn at random.

This proof is the blackboard's alone, and we say so plainly. Its one
real step is the expansion, a sum of two kets, which the sketchbook
cannot yet draw. What the sketchbook can do is state the result, as
arithmetic on pictures already evaluated: the number of \emph{The Echo}
with \(b=b_0\), plus the number of \emph{The Echo} with \(b=b_1\),
equals the number of the closed diagram of Theorem 2.1. Three pictures,
three numbers, and one sum of numbers borrowed from the blackboard. It
is not a rewrite.

Here, too, is the honest answer to the question Chapter 1 left us with,
\emph{why the rule must use squares}. We have not shown that nature had
no other choice; the Born rule is a postulate. What we have shown is
narrower and still worth having. The amplitudes themselves cannot be
shares (§2.3). Their plain squares cannot either. The products
\(\bar z z\) can: they are real and never negative, and --- this is the
corollary --- they make exactly one whole not for one favoured question
but for every question there is, and for no deeper reason than the rule
\(|\alpha|^2+|\beta|^2=1\) by which a qubit state was defined. The
definition of Chapter 1 and the postulate of this chapter were made for
each other. That is the sense, and the only sense, in which the rule
``must'' use squares.

\hypertarget{the-verdict-1}{%
\subsection{The verdict}\label{the-verdict-1}}

A draw, and a modest one; if anything, the algebra has it. Both proofs
unfold the same definition. The blackboard does so in one line. The
sketchbook's three steps are that same line, reached by first reading
the picture; it owns no rewrite rule that \emph{produces} the number
\(1\), and its Step 3 is Chapter 1's blackboard rule under another name.

What the picture gives is not a proof but a place. On the blackboard the
normalization is a rule about the two entries of a column, and so about
the standard basis. In the sketchbook it is one closed picture with a
\(\psi\) at each end and \textbf{no basis in it anywhere} --- which is
the right way to hold it if one wants to see why, by Corollary 2.2, the
total is one whole for every question. The picture locates the square in
the second wire of \emph{The Echo}. It does not explain the square, and
it does not derive the Born rule.

\begin{center}\rule{0.5\linewidth}{0.5pt}\end{center}

\hypertarget{the-echo}{%
\section{§2.9 The Echo}\label{the-echo}}

Here is the chapter's master diagram, first drawn in §2.3. We call it
\textbf{The Echo}. It is a row of four rounded boxes on one line: the
state \(\psi\), whose wire runs to the effect \(b\); then bare paper;
then the state \(b\), whose wire runs to the effect \(\psi\). Each of
its wires is closed at both ends, and the right half is the left half in
the mirror. Its evaluation is the Born rule:

\[\begin{tikzpicture}[sd, baseline=-0.55ex]
  \draw[wire] (0,0) -- (4,0);
  \draw[wire] (8,0) -- (12,0);
  \node[sdstate, text height=1.5ex, text depth=0.5ex] at (0,0) {$\psi$};
  \node[sdeffect, text height=1.5ex, text depth=0.5ex] at (4,0) {$b$};
  \node[sdstate, text height=1.5ex, text depth=0.5ex] at (8,0) {$b$};
  \node[sdeffect, text height=1.5ex, text depth=0.5ex] at (12,0) {$\psi$};
\end{tikzpicture}
\;\;=\;\; \braket{\psi|b}\,\braket{b|\psi} \;\;=\;\; |\braket{b|\psi}|^2\]

The equation is exact. The wire on the left contributes the amplitude
\(\braket{b|\psi}\); the wire on the right contributes its conjugate
\(\braket{\psi|b}\); nothing joins the effect \(b\) to the state \(b\),
and the two numbers multiply. We prove nothing again. Instead, as with
Chapter 1's master diagram, we \emph{read} the one picture three times,
changing only the labels.

\textbf{First reading: the standard question.} Take \(\psi=T\) and
\(b=0\). Each half is \(\tfrac35\).

\[\begin{tikzpicture}[sd, baseline=-0.55ex]
  \draw[wire] (0,0) -- (4,0);
  \draw[wire] (8,0) -- (12,0);
  \node[sdstate, text height=1.5ex, text depth=0.5ex] at (0,0) {$T$};
  \node[sdeffect, text height=1.5ex, text depth=0.5ex] at (4,0) {$0$};
  \node[sdstate, text height=1.5ex, text depth=0.5ex] at (8,0) {$0$};
  \node[sdeffect, text height=1.5ex, text depth=0.5ex] at (12,0) {$T$};
\end{tikzpicture}
\;\;=\;\; \tfrac35\cdot\tfrac35 \;\;=\;\; \tfrac9{25}\]

With \(b=1\) it is \(\tfrac45\cdot\tfrac45=\tfrac{16}{25}\), and the two
numbers are the covered coin \(\begin{pmatrix}9/25\\16/25\end{pmatrix}\)
of §2.3.

\textbf{Second reading: the turned question.} The same state, with
\(b=+\). Each half is \(\tfrac7{5\sqrt2}\).

\[\begin{tikzpicture}[sd, baseline=-0.55ex]
  \draw[wire] (0,0) -- (4,0);
  \draw[wire] (8,0) -- (12,0);
  \node[sdstate, text height=1.5ex, text depth=0.5ex] at (0,0) {$T$};
  \node[sdeffect, text height=1.5ex, text depth=0.5ex] at (4,0) {$+$};
  \node[sdstate, text height=1.5ex, text depth=0.5ex] at (8,0) {$+$};
  \node[sdeffect, text height=1.5ex, text depth=0.5ex] at (12,0) {$T$};
\end{tikzpicture}
\;\;=\;\; \tfrac7{5\sqrt2}\cdot\tfrac7{5\sqrt2} \;\;=\;\; \tfrac{49}{50}\]

With \(b=-\) each half is \(-\tfrac1{5\sqrt2}\), negative on both wires,
and the product is \(\tfrac1{50}\). One state, a different question, a
different coin: \(\begin{pmatrix}49/50\\1/50\end{pmatrix}\). For real
states such as \(T\) the two halves are equal, sign and all. For a state
with non-real entries they are conjugates and in general not equal: with
\(\psi=\chi\) and \(b=+\) the left half is \(\tfrac{3+4i}{5\sqrt2}\),
the right half \(\tfrac{3-4i}{5\sqrt2}\), and the product
\(\tfrac{25}{50}=\tfrac12\).

\textbf{Third reading: the arrow reversed.} Replace \(\psi\) by
\(e^{i\theta}\psi\) and leave \(b\) alone.

\[\begin{tikzpicture}[sd, baseline=-0.55ex]
  \draw[wire] (0,0) -- (5,0);
  \draw[wire] (9,0) -- (14,0);
  \node[sdstate, inner xsep=0.55em, text height=1.5ex, text depth=0.5ex] at (0,0) {$e^{i\theta}\psi$};
  \node[sdeffect, text height=1.5ex, text depth=0.5ex] at (5,0) {$b$};
  \node[sdstate, text height=1.5ex, text depth=0.5ex] at (9,0) {$b$};
  \node[sdeffect, inner xsep=0.55em, text height=1.5ex, text depth=0.5ex] at (14,0) {$e^{i\theta}\psi$};
\end{tikzpicture}
\;\;=\;\; |\braket{b|\psi}|^2\]

A different picture from the one at the head of the section --- two of
its four labels have changed --- and the same number, for every \(b\)
(Proposition 2.4). With \(\theta=\pi\) and \(\psi=T\) this is the
reversed arrow of §2.6: each half of the first reading becomes
\(-\tfrac35\), and the product is still \(\tfrac9{25}\).

There is a fourth reading, which we have already had: put \(b=\psi\),
and each half is the closed diagram of Theorem 2.1, so that \emph{The
Echo} is \(1\cdot1=1\). A state asked whether it is itself says yes,
every time.

One picture; the question changed, the state changed, and each time the
square was found in the same place. See \texttt{Code/Ch2/Squares.py} for
the picture and its first reading, \texttt{Code/Ch2/TurnedQuestion.py}
for the second, and \texttt{Code/Ch2/Reversed.py} for the third.

\begin{quote}
\textbf{Penrose's Sketchbook.}

\emph{What the picture does not say.} \emph{The Echo} says where the
square is. It does not say why the world counts that way, and no amount
of looking at it will.

The idea that a picture with no open wire is a number has a precise
home. Abramsky and Coecke (2004, §6) call the maps from the trivial
system to itself \emph{scalars}, and note that for finite-dimensional
vector spaces these are just the elements of the field; ``no open wire''
is our gloss on that, not their wording. They go on to interpret Dirac's
notation in their setting (§7.2), to observe that on the scalars of
finite-dimensional Hilbert spaces the operation that turns a map round
is complex conjugation (§7.3), and to write the probability of an
outcome as a scalar composed with its turned-round twin, a scalar they
say can be thought of as the ``probability amplitude'' (§8). In their
framework the Born rule emerges from the axioms. In this chapter it is a
postulate, and the picture is a place to keep it.
\end{quote}

\begin{center}\rule{0.5\linewidth}{0.5pt}\end{center}

\hypertarget{what-may-sit-on-the-wire}{%
\section{§2.10 What May Sit on the
Wire}\label{what-may-sit-on-the-wire}}

Let us gather what we have. A ket is a name for a column, and a bra is
the same name on the conjugated row (§2.1, §2.2). In the sketchbook the
bra is an effect, the mirror image of a state, and the mirror
conjugates; a state and an effect on one wire make a closed diagram,
which is a number, an amplitude (§2.2). An amplitude is not a chance.
The chance is the amplitude times its mirror twin, by a postulate that
the picture locates and does not explain (§2.3). A sign on one entry of
a column is invisible to one question and visible to another ---
decisive, for \(\ket{+}\) against \(\ket{-}\) (§2.4, §2.5); a phase on
the whole column is a different column that no question can detect
(§2.6). And a state meets its own mirror image in the number \(1\), from
which it follows that the chances make one whole whatever is asked
(§2.8). Through all of it, one picture with two labels to change (§2.9).

Now look at the pictures once more and notice what none of them has.
Every wire in this chapter runs \emph{bare} from its state to its
effect. Chapter 1 hung boxes on wires and never closed them; this
chapter closed wires and never hung a box on them. Put the two together
and a question presents itself at once. Place a box between the state
\(\psi\) and the effect \(\psi\), and then the same between the pieces
of the echo: is the closed picture still \(1\)? Do the chances still
make one whole? For NOT, which only exchanges the two entries, one can
check that they do. For SET-0 we already know they do not: §1.5 watched
it turn the arrow into \(\begin{pmatrix}7/5\\0\end{pmatrix}\), which
meets its own mirror image in \(\tfrac{49}{25}\).

So some boxes keep the promise of Theorem 2.1 and some break it, and a
wire that carries a qubit can bear only the first kind. To say which
they are we shall need something this chapter was careful not to
provide: the mirror image of a \emph{box}. We have reflected states into
effects and left every process alone. What is the mirror image of a box,
and which boxes keep a closed picture equal to \(1\)? We leave both
questions open here, with the wire bare and the coin still covered. The
next chapter turns the mirror on the boxes themselves.

\hypertarget{luxe4ndler-on-a-sphere}{%
\chapter{Ländler on a Sphere}\label{luxe4ndler-on-a-sphere}}

\begin{center}\rule{0.5\linewidth}{0.5pt}\end{center}

\emph{足跡を}\\
\emph{踏んで戻るや}\\
\emph{深雪道}

\emph{ashiato o}\\
\emph{funde modoru ya}\\
\emph{miyukimichi}

\emph{Deep snow on the path ---}\\
\emph{going back, a foot set down}\\
\emph{in each waiting print.}

\begin{center}\rule{0.5\linewidth}{0.5pt}\end{center}

\emph{{[}The valley has been white for a day and a night, and the snow
lies now, folded in drifts over the one narrow path to the Crab's door.
It is late; the cold has teeth. A single track of prints goes up that
path --- the Tortoise ahead, Achilles set down in her prints behind ---
and nothing else has passed all evening. At the end of it, one curtained
window holds a low, unhurried glow, the only warm thing in all the
dark.{]}}

\emph{{[}Deep winter; a lamp-lit evening; the snow lies. Achilles and
the Tortoise stand on the Crab's doorstep. The door-knob turns.{]}}

\textbf{Achilles:} When one watches a turn from the other side, which
way round does it go?

\textbf{Tortoise:} You are looking at the door-knob, Achilles.

\textbf{Achilles:} It turned by itself. Someone has hold of it on the
other side.

\emph{{[}The door stands open, and the Crab in it.{]}}

\textbf{Crab:} Good evening, good evening! Mind the step; it is the same
step either way.

\textbf{Tortoise:} Good evening, Mr Crab.

\textbf{Crab:} One set of footprints between you, I hope. My path is
narrow.

\textbf{Achilles:} The Tortoise goes first and I go in her prints. It is
slow, but it is tidy.

\textbf{Tortoise:} Shall we, Achilles?

\emph{{[}Over the threshold.{]}}

\textbf{Crab:} In, in, both of you, and let the warm keep what it has.
Coats on the hook that isn't the clock --- everyone tries the clock, and
then the hour is wrong until midnight.

\emph{{[}They come in stamping, and the cold comes off them in a grey
breath the room swallows. The Tortoise lays her wet mittens on the
fender to steam, the one pulled off last dropped on top of the other. In
the grate a new fire has only just caught, spitting and ducking, not yet
sure of its draught; and over it all the clock knocks out its patient
three.{]}}

\textbf{Achilles:} \emph{{[}warming his hands{]}} I had forgotten there
was such a thing as warm. The path did its best to keep us out of doors.

\textbf{Tortoise:} A short path, well kept. Is that your grandmother's
clock still, Mr Crab, keeping its three-time?

\textbf{Crab:} It ticks in threes and always has; I gave up correcting
it years ago. There is tea in the brown pot, and something warmer in the
other if the evening runs long.

\textbf{Achilles:} Tea, before anything. My hands have stopped believing
in themselves.

\emph{{[}The room is low and warm and full of the fire's moving light.
On the wall beside the dresser hangs a print the Crab has had for years
--- Escher's \textup{Hand with Reflecting Sphere}: a hand holds up a
mirrored ball, and in the ball the man himself, and the whole curved
room bending round him to the corners. The lamp lays a soft sheen along
it.{]}}

\hypertarget{first-challenge-one-turn-and-its-mirror}{%
\section{\texorpdfstring{First Challenge --- \emph{One Turn and Its
Mirror}}{First Challenge --- One Turn and Its Mirror}}\label{first-challenge-one-turn-and-its-mirror}}

\emph{{[}Indoors. On the Crab's table, on a felt-lined ring on a low
pedestal, rests a ball the size of an orange, of dull lamp-blacked
brass.{]}}

\textbf{Crab:} There. From a binnacle-maker's widow in Bergen, who had
no notion what she was selling. Go on, Achilles: ask me why it doesn't
shine.

\textbf{Achilles:} Why doesn't it shine? A brass ball ought to.

\textbf{Crab:} Lamp-black, and meant. It is not there to be looked at;
it is there to be turned. Now --- \emph{{[}he touches a small bead on
top of the ball{]}} --- the glass needle. It runs right through and out
underneath, down the hollow of the stand, with a bead on this end only.
At rest it stands upright, bead on top.

\textbf{Achilles:} And that one, pointing at my waistcoat?

\textbf{Crab:} The steel needle. Level, through the middle, bead towards
the company --- by which I mean the two of you; I am furniture. It
passes the glass one by a hair's breadth and touches nothing. I have
lain awake over that hair's breadth.

\textbf{Tortoise:} And the one between, Mr Crab?

\textbf{Crab:} The brass needle, exactly midway between the others: bead
up and towards the company. We shan't want it yet. Glass bead up, steel
bead towards you --- that is \emph{home}, and the brass one follows of
itself.

\textbf{Achilles:} And inside?

\textbf{Crab:} Something inside keeps its bearing, however the case is
turned --- as a compass card does in a turned binnacle. The beads tell
only how the \emph{case} has been turned.

\textbf{Achilles:} Then open it, and see how the inside stands!

\textbf{Crab:} \emph{{[}a claw over a small shutter in the ball's
side{]}} There is a shutter. Open it once and look, and the instrument
is spoiled for good. I have never opened it. Eleven winters, Achilles.

\textbf{Tortoise:} How does one turn it, Mr Crab?

\textbf{Crab:} One way only. Lift the ball clear of the stand; take a
needle by its two ends and hold it still; and with your thumbs turn the
ball about the needle, as a bead turns on a wire. The other needles go
round with the ball. My grandmother danced the ländler so: the couple
turn, and nobody lets go.

\textbf{Achilles:} How far? Which way?

\textbf{Crab:} The collar where each needle goes in has four notches:
one click is a quarter of the way round, two clicks half. With the clock
or against it, as seen from that needle's bead. For two clicks it is all
one --- halfway round is halfway round. Gently!

\textbf{Achilles:} \emph{{[}taking the glass needle by its two ends{]}}
Glass, one click, with the clock. \emph{{[}He lifts the ball just clear
of its stand and holds the glass needle fast at both ends. The glass
bead is the pin the turn is made about: it does not stir. A quarter-turn
with the clock carries the steel and brass needles round with the ball,
and the steel bead, that had pointed at the company, swings a quarter of
the way round. A click; the Crab winces.{]}} The glass bead is on top,
as it was. The steel bead lies level and points at the stove.

\textbf{Crab:} Stove on the one hand, dresser on the other. In this room
that is all the sideways there is.

\textbf{Tortoise:} \emph{{[}She has brought tracing paper, as usual, and
no string. On a strip of it she draws a curved arrow running with the
clock, and beside it a dot of pale blue.{]}} Pale blue for glass; grey
will be steel, and yellow brass. One arrowhead for one click, two for
two. If the arrow runs with the clock on the strip as it lies before the
player, the turn is with the clock as seen from the bead. One plays from
left to right.

\textbf{Achilles:} A score! For one note.

\textbf{Tortoise:} \emph{{[}She holds the strip by its short upright
edge and turns it over towards her, as one turns a page. It is the same
strip --- the blue dot has not moved, the length has not changed --- but
the printed side now faces the paper, and read through it the one arrow
that ran with the clock runs against.{]}} I turn it like a page. Seen
through the paper, the last arrow would come first, and every arrow runs
the other way round. Now I play what I see. \emph{{[}She lifts the ball
and holds the glass needle fast: a quarter-turn against the clock. The
glass bead keeps the top; the steel bead swings back the quarter it had
gone.{]}}

\textbf{Achilles:} Glass bead up. Steel bead towards the company. Home.

\textbf{Crab:} Thank heaven.

\textbf{Achilles:} One note is nothing. Let me write a real one.
\emph{{[}He draws two arrows on a fresh strip, each with two heads: a
blue dot, then a grey.{]}} Glass, two clicks; then steel, two clicks.
\emph{{[}He lifts the ball. Glass, held fast: the glass bead keeps its
place on top while the steel bead is carried the whole half-way round,
off the company and out the far side. Then the steel needle, held where
it now lies: the steel bead keeps that far place while the glass bead is
carried the half-way round, from the top down into the hollow of the
stand. The Crab winces four times.{]}} Glass bead straight down, in the
hollow of the stand. Steel bead away, towards Mr Crab.

\textbf{Tortoise:} And to undo it?

\textbf{Achilles:} I turn the strip round on the table --- flat,
half-way round, not turned over. The last arrow comes first, and each
runs as it was drawn. Steel, two clicks; glass, two clicks. \emph{{[}He
plays what now stands first: the steel needle held, the glass bead
carried back up out of the hollow; then the glass needle held, the steel
bead carried back to the company.{]}} Glass bead up, steel bead towards
the company. Home!

\textbf{Tortoise:} Is that the only way to read a strip backwards?

\textbf{Achilles:} There's another. Flip it top over bottom, about its
long edge: the order stays, and every arrow runs the other way round.

\emph{{[}Twice more the score is played from home --- glass bead down,
steel bead away, each time --- and twice more undone: once from the
strip flipped top over bottom, once from the strip turned like a page.
Home; and home. Then Achilles plays it steel first, glass second ---
glass bead down, steel bead away --- and plays the same again: home.{]}}

\emph{{[}The Crab has crossed to the dresser for the other pot and comes
back the way he went, sidelong, the same step carrying him out and in
--- a crab is at no more trouble to go one way than the other. He tops
up the cups without being asked and settles again into his chair.{]}}

\textbf{Achilles:} The same place, whichever goes first. So what you're
saying is: to undo a score, play it again from the other end --- any
other end --- and the order never mattered anyway.

\textbf{Tortoise:} I said nothing. But it is a handsome rule, and it has
not failed yet.

\textbf{Achilles:} Nor will it.

\textbf{Tortoise:} Then one click each will be no harder.

\begin{center}\rule{0.5\linewidth}{0.5pt}\end{center}

\hypertarget{second-challenge-the-score-read-backwards}{%
\section{\texorpdfstring{Second Challenge --- \emph{The Score Read
Backwards}}{Second Challenge --- The Score Read Backwards}}\label{second-challenge-the-score-read-backwards}}

\textbf{Tortoise:} \emph{{[}On Achilles's strip she rubs out one head of
each arrow.{]}} Glass, one click; then steel, one click. Both with the
clock, as you drew them. Mr Crab?

\textbf{Crab:} \emph{{[}He plays it, wincing: glass a quarter, the glass
bead keeping the top while the steel bead leaves the company; then the
steel needle where it now lies, held, the steel bead keeping its new
place while the glass bead is carried a quarter down from the top.{]}}
If I must.

\textbf{Achilles:} Glass bead level, pointing at the company. Steel bead
level, pointing at the stove. Now, my rule. The strip round on the
table: steel one click, then glass one click, each with the clock as
drawn. \emph{{[}From where the beads now stand he plays it: the steel
needle held at the stove while the glass bead swings a quarter, from the
company down into the hollow; then the glass needle held while the steel
bead swings a quarter round to the company.{]}} Glass bead\ldots{}
straight down. Steel bead towards the company. The steel is home and the
glass is standing on its head!

\textbf{Crab:} Allow me. \emph{{[}He backs out of Achilles's undoing,
last click first and each the other way: the glass needle held while the
steel bead swings a quarter back, then the steel needle held while the
glass bead swings a quarter back --- and the beads stand as they did
after the score was played.{]}} Glass bead at the company, steel bead at
the stove. As it was before you began.

\textbf{Achilles:} Then my other rule. The strip top over bottom: glass
first, then steel, each against the clock. \emph{{[}He plays it: the
glass needle held, the steel bead swinging a quarter the other way; then
the steel needle held, the glass bead swinging a quarter the other
way.{]}} Glass bead level, pointing at the stove. Steel bead straight
down. That's not home either. It isn't even the same mistake!

\textbf{Crab:} Allow me. \emph{{[}He backs out again, last click first
and each the other way: the steel needle held while the glass bead
swings a quarter back, then the glass needle held while the steel bead
swings a quarter back.{]}} Glass at the company, steel at the stove. As
it was.

\textbf{Achilles:} But both rules worked. Every time!

\textbf{Tortoise:} On arrows with two heads. \emph{{[}She turns the
strip over about its short edge. Now the grey arrow stands first and the
blue second --- the order reversed --- and, read through the paper, each
single head points the other way; the two dots have not moved, nor the
colours.{]}} The page. Read me what you see.

\textbf{Achilles:} Steel first --- that's my first rule. Each against
the clock --- that's my second. \emph{{[}He plays it from where the
beads stand: the steel needle held while the glass bead swings a quarter
back up out of the hollow; then the glass needle held while the steel
bead swings a quarter back to the company.{]}} Glass bead up. Steel bead
towards the company. Home. \ldots Oh. Both at once. Last arrow first
\emph{and} every arrow the other way round, and neither half will do
alone.

\emph{{[}On the fender the mittens have gone quiet, done steaming. The
one pulled off last lies on top of the other and inside-out, as it left
the hand.{]}}

\textbf{Tortoise:} And the two-headed arrows?

\textbf{Achilles:} They hid both halves. Two clicks have no with or
against, and those two turns didn't mind which went first. And Mr Crab
has been doing it under my nose. Twice!

\textbf{Crab:} Doing what? I went out the way you came in. I always
leave a room by the way I came in, only backwards. I see nothing
remarkable in it.

\textbf{Tortoise:} Are there scores that the page does not change?

\begin{center}\rule{0.5\linewidth}{0.5pt}\end{center}

\hypertarget{third-challenge-the-sandwich}{%
\section{\texorpdfstring{Third Challenge --- \emph{The
Sandwich}}{Third Challenge --- The Sandwich}}\label{third-challenge-the-sandwich}}

\emph{{[}The fire has its draught now and burns tall and even, and the
room has stopped feeling like a place one has just come in from the
cold. The clock ticks on over everything, unhurried, as it has all
evening.{]}}

\textbf{Tortoise:} \emph{{[}On a fresh strip she draws a two-headed
arrow with a yellow dot, a two-headed arrow with a blue dot, a
two-headed arrow with a yellow dot.{]}} Brass, two clicks. Glass, two
clicks. Brass, two clicks. \emph{{[}She turns the strip over about its
short edge, and back again. Through the paper the order is reversed ---
yellow, blue, yellow, which is yellow, blue, yellow still --- and each
arrow would run the other way, but a two-headed arrow has no other way;
so the strip in the glass is the strip on the table, nothing
altered.{]}} The same score, either side of the paper.

\textbf{Crab:} The brass needle at last.

\textbf{Achilles:} Before anyone touches it --- Mr Crab, the maker's
card that came in the case. May I?

\textbf{Crab:} For each needle's turn a little square of four numbers. I
never could read it.

\textbf{Achilles:} I can, and I needn't even multiply. Brass twice over
is no turn at all: two half-turns. So the brass arrows cancel, near
enough, and what is left is the glass half-turn in the middle. Glass
bead up, steel bead away. Nearly home.

\textbf{Tortoise:} There is a glass turn standing between your two brass
ones.

\textbf{Achilles:} With two-headed arrows the order never mattered. We
tried it! Nearly home, I say.

\textbf{Crab:} \emph{{[}He takes the brass needle by its ends: two
clicks. The brass bead is the pin now --- it holds fast while the glass
and steel needles swing a half-turn right past one another.{]}} Brass.

\textbf{Tortoise:} Look where the glass needle lies now.

\textbf{Achilles:} Level, bead at the company --- where the steel one
lay. And the steel one stands upright, bead on top, in the glass one's
place. They have changed places, bead for bead.

\textbf{Crab:} \emph{{[}The glass needle, as it now lies level: two
clicks --- the glass bead keeps its place while the steel swings the
half-way round; then the brass needle: two clicks, the brass bead
holding while the glass and steel trade ends once more.{]}} Glass.
Brass.

\textbf{Achilles:} Glass bead straight down. Steel bead towards the
company, where it began. Brass bead down and towards. That isn't nearly
home. I said glass up and steel away, and it is exactly the other thing!

\textbf{Tortoise:} \emph{{[}She turns the strip like a page once
more.{]}} The page, Mr Crab.

\emph{{[}The Crab plays the same score again, the ball lifted clear in
his claw: brass, two clicks --- glass and steel trade ends; glass, two
clicks --- the glass bead holds while the steel swings; brass, two
clicks --- and they trade back. Glass bead up, steel bead towards the
company. With the dull ball still held up before him he glances at the
print on the wall --- Escher's \textup{Hand with Reflecting Sphere}.{]}}

\emph{{[}The clock knocks out its three, and the room goes still with
it. Held up so, between the lamp and the wall, the ball throws a soft
round shadow across the print, and for a moment the little curved room
inside the mirrored sphere seems to lean out toward the lamplight.{]}}

\textbf{Crab:} That fellow holds up his ball and it shows him himself,
and the whole room behind him. Mine shows me nothing at all --- which is
the whole point of it.

\textbf{Tortoise:} Home, Achilles. Remember where the beads were.

\textbf{Crab:} \emph{{[}He sets the ball down, takes it up again by the
two ends of the steel needle: two clicks. The steel needle holds fast
while everything else swings a half-turn about it --- the steel bead
does not stir, and the glass and brass beads come round to the far
side.{]}} One turn of the steel --- and there they all are again.

\textbf{Achilles:} Glass bead straight down. Steel bead towards the
company. Brass bead down and towards. Bead for bead the same!

\textbf{Tortoise:} And your numbers?

\textbf{Achilles:} Honestly, this time. \emph{{[}He takes the maker's
card and a scrap of paper.{]}} The brass square, then the glass square,
then the brass square: row against column, twice over. The halves double
up into wholes, the signs fall where they should --- and what is left is
the steel needle's square, number for number. And it was quicker than Mr
Crab's six clicks.

\emph{{[}He sets the scrap down and sits back, pleased with it, and for
a moment only the fire and the clock are working.{]}}

\textbf{Tortoise:} It was. Did it tell you why?

\textbf{Achilles:} No.~Only that it is so. But after the first brass
turn the glass needle lay where the steel one had lain\ldots{}

\textbf{Crab:} \emph{{[}The steel needle again: two clicks.{]}} And two
clicks more of the steel, and we are home. Glass up, steel towards you.

\textbf{Achilles:} But the brass needle is a luxury, surely --- could
one do without it?

\begin{center}\rule{0.5\linewidth}{0.5pt}\end{center}

\hypertarget{fourth-challenge-three-turns-make-any-turn}{%
\section{\texorpdfstring{Fourth Challenge --- \emph{Three Turns Make Any
Turn}}{Fourth Challenge --- Three Turns Make Any Turn}}\label{fourth-challenge-three-turns-make-any-turn}}

\textbf{Tortoise:} Glass, steel, glass. One click each, with the clock.

\textbf{Achilles:} \emph{{[}He plays it: glass, one click --- the glass
bead holds while the steel bead swings a quarter; steel, one click ---
the steel bead holds while the glass bead swings a quarter; glass, one
click again.{]}} Glass bead level, pointing at the company. Steel bead
on top. But that is where the first brass turn left them!

\textbf{Tortoise:} Now say the score from the other end.

\textbf{Achilles:} Glass, steel, glass. It is the same from both ends.
But on the page each arrow would run against the clock, so that is its
undoing. \emph{{[}He plays it, each click running against the clock this
time: glass, steel, glass --- every quarter-turn walking back the way it
came.{]}} Glass bead up, steel bead towards the company. Home, as the
page promised.

\textbf{Tortoise:} Again. The page again, from home.

\textbf{Achilles:} From \emph{home}? \emph{{[}He plays it: glass, steel,
glass, each against the clock.{]}} Glass bead at the company. Steel bead
on top. The same place! The score and its page go to the same place.

\textbf{Tortoise:} As a brass turn should.

\textbf{Crab:} \emph{{[}He takes the brass needle by its ends: two
clicks --- brass bead holding, glass and steel trading ends.{]}} Then
this will take it back. Glass up. Steel towards the company. Home.

\textbf{Achilles:} So the brass needle \emph{is} a luxury. And the
others? Could one get anywhere at all without it?

\textbf{Tortoise:} Mr Crab, would you shut your eyes and tumble it?

\textbf{Crab:} Tumble it! \emph{{[}Eyes shut, he turns the ball over and
over in his claws and sets it down anyhow.{]}} It can be coaxed to rest
between the notches of a collar, if one is gentle. Which I beg you to
be.

\textbf{Achilles:} Glass bead low, somewhere between the dresser and
away. Steel bead high, towards the stove. Nowhere with a name. I'll take
the glass needle. \emph{{[}He holds it by its ends and turns the ball
slowly about it.{]}} The steel bead goes round and round. The glass bead
doesn't stir. Of course it doesn't: I am holding it.

\textbf{Tortoise:} Glass, until the steel needle lies level. Stop. Now
steel, until the glass bead is on top. Stop. Now glass again, until the
steel bead points at the company.

\textbf{Achilles:} \emph{{[}He does each as she says it: the glass
needle held fast while the steel bead swings round until it lies level;
then the steel needle held while the glass bead climbs to the top; then
the glass needle held again while the steel bead comes round to the
company --- each time the held bead standing still and the far one
travelling only as far as she says, then stopping.{]}} Level. On top. At
the company. Glass bead up, steel bead towards the company: home, from
nowhere.

\textbf{Tortoise:} Glass, steel, glass. It is always enough.

\emph{{[}In the grate a log that has burned through at the middle gives
way, drops, and rolls a little until it finds a flat seat among the
coals and holds there --- settled, from wherever it happened to
fall.{]}}

\textbf{Crab:} And now nothing sits in its notch. Hold the ball still on
its ring, Achilles --- still! --- while I roll each needle alone until
it clicks into its collar. \emph{{[}He rolls the glass needle between
his claws: a click. The steel: a click. The brass: a click.{]}} There.
The ball has not moved, and a click is a quarter of the way round again.

\textbf{Achilles:} So everything can be undone, if one is patient.

\textbf{Tortoise:} Everything on the strip.

\begin{center}\rule{0.5\linewidth}{0.5pt}\end{center}

\hypertarget{fifth-challenge-what-has-no-mirror}{%
\section{\texorpdfstring{Fifth Challenge --- \emph{What Has No
Mirror}}{Fifth Challenge --- What Has No Mirror}}\label{fifth-challenge-what-has-no-mirror}}

\textbf{Achilles:} What else is there? \emph{{[}He looks over the
stand.{]}} This little brass catch on the pedestal?

\textbf{Crab:} Ah. Put the ball somewhere first, and tell us where.

\textbf{Achilles:} Steel, one click with the clock. Glass bead level,
pointing at the dresser. Steel bead towards the company.

\textbf{Crab:} That little catch is the ball's valet --- when a guest
leaves it in a muddle, one press sets it to rights: home again, however
it was left.

\textbf{Crab:} \emph{{[}He presses the catch. Wherever the beads stood,
they are gone from there in the one instant: the glass bead up, the
steel bead at the company --- home, as though no turn had been made at
all.{]}} Home.

\textbf{Achilles:} Glass up, steel towards the company --- so it is!
Again. Glass, one click against the clock; then steel, two clicks. Glass
bead straight down. Steel bead pointing at the dresser.

\textbf{Crab:} \emph{{[}He presses the catch. From this second place too
the beads leap to exactly the same home --- glass up, steel at the
company --- though they set out from nowhere near the first place.{]}}
Home.

\textbf{Achilles:} From there as well! How does it do it?

\textbf{Crab:} A weight inside, I dare say. Will anyone take more tea?

\emph{{[}He pours, and tips a little milk after it; the pale cloud turns
over once in the cup and thins away to one even brown. Achilles takes
the cup and warms both hands round it.{]}}

\textbf{Tortoise:} \emph{{[}She hands Achilles a blank strip.{]}} Draw
the catch.

\textbf{Achilles:} One arrow\ldots{} From the first place it did what
steel, one click against the clock, would have done. From the second it
did something else altogether. It isn't one turn. It is whichever turn
happens to be wanted.

\textbf{Tortoise:} Draw it anyway. Then turn the page.

\textbf{Achilles:} The page would take the ball from home back to where
it was. But it was in two different places, and the catch took both to
the same one. The page would have to lead to two places at once.
\ldots Oh. It can't be drawn.

\textbf{Tortoise:} And what can?

\textbf{Achilles:} What the page can undo --- a turn, and nothing but a
turn. The catch is not on the strip.

\textbf{Tortoise:} \emph{{[}She takes the evening's strips --- her one
blue arrow; Achilles's strip with the heads rubbed out; brass, glass,
brass --- lays them end to end in that order, and gums them into
one.{]}} Then let us see how much the page can undo. From home, Mr Crab,
left to right.

\textbf{Crab:} \emph{{[}He plays the long score, wincing at every click:
glass, one click; glass, one click, and steel, one click; brass, two
clicks, glass, two clicks, brass, two clicks.{]}} All of it. My poor
collars.

\textbf{Achilles:} Glass bead level, pointing at the dresser --- and the
little score sent everything sideways to the stove. Steel bead level,
away. Brass bead level too, away and towards the dresser. Nowhere near
home.

\textbf{Tortoise:} \emph{{[}She turns the whole long strip over like a
page.{]}} Read it to him.

\textbf{Achilles:} Brass, two clicks. Glass, two clicks. Brass, two
clicks. Steel, one click against the clock. Glass, one click against.
Glass, one click against.

\textbf{Crab:} \emph{{[}He plays it as it is read.{]}} And the last.

\textbf{Achilles:} Glass bead up. Steel bead towards the company. Home
--- by turns alone, to the last click.

\textbf{Crab:} Home, and nothing owing. Every click of the evening has
been paid back. \emph{{[}His claw rests on the shutter.{]}} There is one
thing in this house for which nobody can write a page.

\textbf{Tortoise:} No, Mr Crab. Nobody can.

\textbf{Crab:} Eleven winters. I find I should rather know.

\textbf{Achilles:} Whatever it shows you, Mr Crab, we are glad we came.

\emph{{[}The Crab opens the shutter, and looks.{]}}

\emph{{[}On the wall the print is only a print again: the mirrored
sphere holds its same curved room and nothing that has happened tonight.
The fire has burned down to a bed of coals with one red eye still open
in the grey, and the clock ticks on, softer now.{]}}

\emph{{[}For a moment no one speaks.{]}}

\textbf{Tortoise:} Shall we, Achilles?

\emph{{[}The doorstep. The snow lies; one set of footprints runs down
the path, and every print points at the house.{]}}

\textbf{Achilles:} The Tortoise goes first and I go in her prints. It is
slow, but it is tidy.

\textbf{Crab:} One set of footprints between you, I hope. My path is
narrow.

\textbf{Tortoise:} Good evening, Mr Crab.

\textbf{Crab:} Good evening, good evening! Mind the step; it is the same
step either way.

\emph{{[}The door closes. The door-knob turns.{]}}

\textbf{Achilles:} It turned by itself. Someone has hold of it on the
other side.

\textbf{Tortoise:} You are looking at the door-knob, Achilles.

\textbf{Achilles:} When one watches a turn from the other side, which
way round does it go?

\emph{{[}Outside, the one track stays as it was --- the Tortoise's
prints and Achilles's set in them --- and the frost is setting the snow
hard about their edges, to keep them till morning. The window holds its
low glow behind the drawn curtain; behind that, a banked fire the Crab
will build up again once the cold has the house to itself.{]}}

\hypertarget{chapter-3-turning-the-sphere}{%
\chapter{Chapter 3: Turning the
Sphere}\label{chapter-3-turning-the-sphere}}

\emph{--- Ländler on a Sphere ---}

\emph{Part I: The Single Wire}

\emph{Escher: ``Hand with Reflecting Sphere'' --- a hand holds up a
mirrored ball, and in the ball are the man who holds it and the whole
room behind him}

\begin{center}\rule{0.5\linewidth}{0.5pt}\end{center}

When you watch a turn from the other side, which way round does it go
--- and why should that be the way to undo it?

Stand in front of a shop window while someone inside turns a globe on
its stand. From where they stand it goes with the clock. From where you
stand, looking through the glass, it goes against. Nothing about the
globe is in dispute between you; you are watching one turn from its two
sides. Now suppose you wanted to put the globe back as it was. You would
do what \emph{you} saw, from where \emph{they} stand. The undoing of a
turn is the same turn, seen from behind.

That is easy enough for one turn. It becomes a real question when there
are several, one after another, about different axes. Do you undo them
in the order they were made, or in the other order? Do you reverse each
of them, or only the whole? People who have packed a suitcase know the
answer in their hands --- last in, first out, and every movement
backwards --- and still get it wrong on paper.

For us the question has an edge to it. Two chapters have now said that a
qubit is a thing nobody looks at: a covered coin, an arrow in the dark.
Chapter 1 hung boxes on its wire and never closed the wire; Chapter 2
closed the wire and never hung a box on it. If a box is to act on
something that may not be inspected before, during or after, then there
had better be a guarantee, written into the box itself, that it has not
spoiled what it carries. We shall find that the guarantee and the view
from the other side of the glass are the same thing. The boxes that may
sit on a qubit's wire are exactly the ones undone by their own mirror
image. And when we ask what those boxes \emph{do}, the answer will be
the plainest picture in this book's subject: they turn a sphere.

\begin{center}\rule{0.5\linewidth}{0.5pt}\end{center}

\hypertarget{a-box-on-a-closed-wire}{%
\section{§3.1 A Box on a Closed Wire}\label{a-box-on-a-closed-wire}}

Chapter 2 ended (§2.10) with an instruction and two questions:
\emph{place a box between the state \(\psi\) and the effect \(\psi\),
and then the same between the pieces of the echo: is the closed picture
still \(1\)? Do the chances still make one whole?} Before we can answer
we must learn to read a closed picture that has a box on it, and then we
must be careful about which closed picture the question means.

One change of name first. Chapter 2 called Chapter 1's arrow \(T\), for
its triangle. This chapter needs the letter \(T\) for a box that the
whole literature calls by that name (§3.4), so here the arrow is
\(\ket{a}\):

\[\ket{a} \;=\; \tfrac35\ket{0}+\tfrac45\ket{1} \;=\; \begin{pmatrix}3/5\\4/5\end{pmatrix}.\]

And the box NOT will from now on carry the name of its matrix, \(X\).

\textbf{On the blackboard.} A bra is a row, a \(2\times2\) matrix is a
matrix, a ket is a column, and the product of the three, in that order,
is a row times a column: a number. We write it with the matrix between
the bars,

\[\bra{\phi}A\ket{\psi},\]

and compute it in whichever grouping is convenient, since matrix
products may be grouped as we please: either the row \(\bra{\phi}\)
applied to the column \(A\ket{\psi}\), or the row \(\bra{\phi}A\)
applied to the column \(\ket{\psi}\).

The plainest case tells us what these numbers are. Number the entries of
a matrix from zero, row first,

\[A \;=\; \begin{pmatrix}A_{00}&A_{01}\\A_{10}&A_{11}\end{pmatrix}.\]

The column \(A\ket{j}\) is the \(j\)-th column of \(A\) (§1.1), and the
row \(\bra{i}\) picks out the \(i\)-th entry of whatever column it meets
(§2.2). So, for \(i\) and \(j\) each \(0\) or \(1\),

\[\bra{i}A\ket{j} \;=\; A_{ij} .\]

\textbf{The four numbers in a matrix are its four closed pictures on the
standard states.} For NOT, \(\bra{1}X\ket{0}=1\) and
\(\bra{0}X\ket{0}=0\): the state \(0\), after NOT, answers ``yes'' to
the question \(1\) and ``no'' to the question \(0\).

\textbf{In the sketchbook.} No new element yet. A state, then a box,
then an effect, all on one wire, is a row of three things with no open
end, so by the row rule (§1.3) and the reading of closed diagrams (§2.2)
it is a number, and the number is the product just written. As always
the picture runs left to right and the product right to left:

\[\begin{tikzpicture}[sd, baseline=-0.55ex]
  \draw[wire] (0,0) -- (6,0);
  \node[sdstate, text height=1.5ex, text depth=0.5ex] at (0,0) {$\psi$};
  \node[sdbox, text height=1.5ex, text depth=0.5ex] at (3,0) {$A$};
  \node[sdeffect, text height=1.5ex, text depth=0.5ex] at (6,0) {$\phi$};
\end{tikzpicture}
\;\;=\;\; \bra{\phi}A\ket{\psi}
\qquad\qquad
\begin{tikzpicture}[sd, baseline=-0.55ex]
  \draw[wire] (0,0) -- (6,0);
  \node[sdstate, text height=1.5ex, text depth=0.5ex] at (0,0) {$j$};
  \node[sdbox, text height=1.5ex, text depth=0.5ex] at (3,0) {$A$};
  \node[sdeffect, text height=1.5ex, text depth=0.5ex] at (6,0) {$i$};
\end{tikzpicture}
\;\;=\;\; A_{ij}\]

Every picture equation in this chapter is exact: the same number, or the
same matrix, on both sides, with no factor swept aside.

Now the question of §2.10, and a warning about how to read it. Taken
literally, ``a box between the state \(\psi\) and the effect \(\psi\)''
is the picture on the left above with \(\phi=\psi\). It is a perfectly
good number. For the arrow and NOT,

\[\begin{tikzpicture}[sd, baseline=-0.55ex]
  \draw[wire] (0,0) -- (6,0);
  \node[sdstate, text height=1.5ex, text depth=0.5ex] at (0,0) {$a$};
  \node[sdbox, text height=1.5ex, text depth=0.5ex] at (3,0) {$X$};
  \node[sdeffect, text height=1.5ex, text depth=0.5ex] at (6,0) {$a$};
\end{tikzpicture}
\;\;=\;\; \begin{pmatrix}\tfrac35&\tfrac45\end{pmatrix}\begin{pmatrix}4/5\\3/5\end{pmatrix} \;=\; \tfrac{12}{25}+\tfrac{12}{25} \;=\; \tfrac{24}{25},\]

which is not \(1\), and nothing has gone wrong: the arrow was turned
over and then asked whether it was still the \emph{old} arrow. (Keep the
number \(\tfrac{24}{25}\); it returns in §3.5 with a new meaning.) That
is not the promise of Theorem 2.1. The promise was that a state meets
\emph{its own} mirror image in \(1\), and the state on the wire after
the box is no longer \(\ket{a}\); it is \(X\ket{a}\). What §2.10
computed for SET-0, and what we must ask of every box, is whether the
state \emph{after the box} still meets its own mirror image in \(1\):

\[X\ket{a}=\begin{pmatrix}4/5\\3/5\end{pmatrix}:\;\; \tfrac{16}{25}+\tfrac{9}{25}=1;\]

\[S_0\ket{a}=\begin{pmatrix}1&1\\0&0\end{pmatrix}\begin{pmatrix}3/5\\4/5\end{pmatrix}=\begin{pmatrix}7/5\\0\end{pmatrix}:\;\; \tfrac{49}{25} .\]

NOT keeps the promise for the arrow; SET-0 breaks it, as §1.5 foresaw.
See \texttt{Code/Ch3/NoMirror.py}.

Here the sketchbook is stuck, and the place where it sticks is the
subject of the chapter. We can draw the state ``\(a\), then \(X\)''. To
draw \emph{its mirror image} we could cheat, work out the column on the
blackboard and put the answer in a rounded box. But a mirror ought to
work on the picture as it stands, piece by piece. Chapter 2 taught the
mirror to reflect a state into an effect. It was careful to say (§2.2)
that it had not taught it to reflect a box.

\textbf{\emph{Why this definition.}} Why read §2.10's question as ``the
state after the box meets its own mirror image'', and not as the literal
picture with \(\psi\) at both ends? Because only the first is the
promise of Theorem 2.1, on which the Born rule's ``one whole'' rests
(Corollary 2.2). The literal picture compares the new state with the old
one, and a box that does anything at all is entitled to fail that
comparison.

\begin{center}\rule{0.5\linewidth}{0.5pt}\end{center}

\hypertarget{the-mirror-of-a-box}{%
\section{§3.2 The Mirror of a Box}\label{the-mirror-of-a-box}}

We do not get to choose what the mirror image of a box is. Chapter 2 has
already decided it.

Hold the mirror up to the closed picture ``state \(j\), box \(A\),
effect \(i\)'', whose number is \(A_{ij}\). §2.2 told us three things
about what we must see. The state \(j\) on the left becomes the effect
\(j\) on the right. The effect \(i\) on the right becomes the state
\(i\) on the left. And the number is conjugated. In between there is
some box, the reflection of \(A\); call it \(A'\) for a moment. The
mirrored picture reads ``state \(i\), box \(A'\), effect \(j\)'', which
is the entry \(A'_{ji}\). So

\[A'_{ji} \;=\; \overline{A_{ij}} \qquad\text{for all four choices of } i \text{ and } j,\]

and four entries are a whole matrix. There is nothing left to decide.

\textbf{On the blackboard.} The \textbf{conjugate transpose} of a matrix
\(A\), written \(A^{\dagger}\) and read ``\(A\) dagger'', is the matrix
whose entry in row \(j\) and column \(i\) is the conjugate of the entry
of \(A\) in row \(i\) and column \(j\):

\[\bigl(A^{\dagger}\bigr)_{ji} \;=\; \overline{A_{ij}}, \qquad\qquad \begin{pmatrix}p&q\\r&s\end{pmatrix}^{\dagger} \;=\; \begin{pmatrix}\bar p&\bar r\\\bar q&\bar s\end{pmatrix}.\]

Two things happen at once. The matrix is flipped about its diagonal,
rows becoming columns: that alone is the \textbf{transpose}, \(A^{T}\).
And every entry is conjugated: that alone is the \textbf{conjugate},
\(\overline{A}\). The dagger is both. Done twice it does nothing,
\(\bigl(A^{\dagger}\bigr)^{\dagger}=A\), since two flips are none and
two conjugations are none; and \(I^{\dagger}=I\).

The same recipe applies to matrices of any shape, and we have been using
it since §2.2 without the name: the conjugate transpose of a column is
the conjugated row. \textbf{The bra is the dagger of the ket},
\(\bra{\psi}=\bigl(\ket{\psi}\bigr)^{\dagger}\), and the ket is the
dagger of the bra.

To see that the three operations really differ we need a matrix that is
neither real nor symmetric about its diagonal; for real matrices the
conjugate is invisible, for symmetric ones the transpose is. Here are
two matrices that §3.4 will introduce properly. For now they are only
arrays of numbers:

\[S=\begin{pmatrix}1&0\\0&i\end{pmatrix}, \qquad V=\tfrac12\begin{pmatrix}1+i&1-i\\1-i&1+i\end{pmatrix},\]

\[U \;=\; VS \;=\; \tfrac12\begin{pmatrix}1+i&1+i\\1-i&-1+i\end{pmatrix}.\]

(Multiplying on the right by \(S\) multiplies the second column by
\(i\), and \(i(1-i)=1+i\), \(\;i(1+i)=-1+i\).) The matrix \(U\) is the
row ``\(S\), then \(V\)''. It is not real, and it is not symmetric: its
off-diagonal entries are \(\tfrac{1+i}2\) and \(\tfrac{1-i}2\). It will
be our test piece for the whole chapter. Side by side:

\[U^{T}=\tfrac12\begin{pmatrix}1+i&1-i\\1+i&-1+i\end{pmatrix}, \qquad \overline{U}=\tfrac12\begin{pmatrix}1-i&1-i\\1+i&-1-i\end{pmatrix},\]

\[U^{\dagger}=\tfrac12\begin{pmatrix}1-i&1+i\\1-i&-1-i\end{pmatrix}.\]

Three different matrices, and none of them is \(U\).

\textbf{In the sketchbook.} One new element, described before it is
drawn; it is the only new element of the chapter.

The \textbf{mirror image of a box} \(A\) is the box seen in a mirror
held upright across the wire. The wire that came in on the left now
leaves on the right, and the wire that left on the right now comes in on
the left: \emph{input and output change places}. And, as for states in
§2.2, \textbf{the mirror conjugates.} The reflected box stands for the
matrix \(A^{\dagger}\).

We have to be candid about the drawing. A box on one wire has one wire
on each side, and its outline is a plain rectangle, so the reflected box
has exactly the outline of the original. With a state, §2.2 could say
``look for the wire''; here the wire is on both sides, before the mirror
and after. \textbf{On a single wire the mirror shows in only two places:
in the label, which we write \(A^{\dagger}\), and in the order of a row}
(below). So we draw it like this:

\[\text{a box:}\;\;
\begin{tikzpicture}[sd, baseline=-0.55ex]
  \draw[wire] (0,0) -- (5,0);
  \node[sdbox, text height=1.5ex, text depth=0.5ex] at (2.5,0) {$A$};
\end{tikzpicture}
\qquad\qquad
\text{its mirror image:}\;\;
\begin{tikzpicture}[sd, baseline=-0.55ex]
  \draw[wire] (0,0) -- (5,0);
  \node[sdbox, text height=1.5ex, text depth=0.5ex] at (2.5,0) {$A^{\dagger}$};
\end{tikzpicture}\]

Those who draw such pictures for a living meet the same difficulty, and
solve it by putting a mark in one corner of every box, so that the
reflection can be told from the original by where the mark has gone (the
\emph{Feynman's Notebook} of §3.8 has the reference). Our mark is the
dagger on the label.

With the new element the opening argument of this section becomes a
picture equation, true for all four choices of \(i\) and \(j\):

\[\begin{tikzpicture}[sd, baseline=-0.55ex]
  \draw[wire] (0,0) -- (6,0);
  \node[sdstate, text height=1.5ex, text depth=0.5ex] at (0,0) {$i$};
  \node[sdbox, text height=1.5ex, text depth=0.5ex] at (3,0) {$A^{\dagger}$};
  \node[sdeffect, text height=1.5ex, text depth=0.5ex] at (6,0) {$j$};
\end{tikzpicture}
\;\;=\;\; \bra{j}A^{\dagger}\ket{i} \;\;=\;\; \overline{A_{ij}} .\]

\textbf{\emph{Why this definition.}} Why both the flip and the
conjugation, when either alone is simpler? Because Chapter 2 left no
room. The mirror exchanges the two ends of a closed picture, and that is
the flip, \(ij\) becoming \(ji\); and the mirror conjugates the number,
and that is the bar. A reflection that did only one of the two would
disagree with §2.2 about the four closed pictures that make up the box.

\hypertarget{the-mirror-of-a-row}{%
\subsection{The mirror of a row}\label{the-mirror-of-a-row}}

Our test piece \(U\) is a row of two boxes, \(S\) then \(V\). So there
are two ways to find its mirror image: reflect the matrix \(U\) as a
whole, which we have done; or hold the mirror up to the row. A mirror
held up to a row of things reverses their order: what was leftmost is
rightmost in the glass. So the mirror image of ``\(S\), then \(V\)''
should be ``\(V^{\dagger}\), then \(S^{\dagger}\)'', whose matrix
(rightmost box first in the product, as always) is
\(S^{\dagger}V^{\dagger}\). Let us see. Both \(S\) and \(V\) happen to
be symmetric, so their daggers are their conjugates:

\[S^{\dagger}V^{\dagger} \;=\; \begin{pmatrix}1&0\\0&-i\end{pmatrix}\cdot\tfrac12\begin{pmatrix}1-i&1+i\\1+i&1-i\end{pmatrix} \;=\; \tfrac12\begin{pmatrix}1-i&1+i\\1-i&-1-i\end{pmatrix} \;=\; U^{\dagger} .\]

(Multiplying on the left by \(S^{\dagger}\) multiplies the second row by
\(-i\).) It is. And the two cheaper recipes fail. Reflect each piece but
keep the order, ``\(S^{\dagger}\), then \(V^{\dagger}\)'': the matrix is
\(V^{\dagger}S^{\dagger}=\overline{V}\,\overline{S}=\overline{U}\),
which is not \(U^{\dagger}\). Reverse the order but do not reflect the
pieces, ``\(V\), then \(S\)'': the matrix is
\(SV=S^{T}V^{T}=(VS)^{T}=U^{T}\), which is not \(U^{\dagger}\) either.
(That these two wrong answers are exactly the conjugate and the
transpose of \(U\) is an accident of our having chosen symmetric pieces;
that they are \emph{wrong} is not.)

\begin{quote}
\textbf{Theorem 3.1 (The Mirror of a Row).} For matrices \(A\) and \(B\)
whose product \(BA\) is defined, of any shapes,

\[(BA)^{\dagger} \;=\; A^{\dagger}B^{\dagger} .\]

In particular, for every \(2\times2\) matrix \(A\) and all states
\(\ket{\psi}\), \(\ket{\phi}\):

\[\bigl(A\ket{\psi}\bigr)^{\dagger}=\bra{\psi}A^{\dagger}, \qquad\qquad \overline{\bra{\phi}A\ket{\psi}} \;=\; \bra{\psi}A^{\dagger}\ket{\phi} .\]

In pictures: the mirror image of a row is the row of the mirror images,
\textbf{in the reverse order}; and the mirror image of a closed picture
is the conjugate number.
\end{quote}

\emph{Proof idea.} An entry of a product is a sum of products of
entries, one from each factor. The flip exchanges the roles of rows and
columns in both factors, which is what reverses their order; the
conjugation passes through sums and products (§2.2) and so can be done
entry by entry. \emph{Proof sketch.} The index calculation is two lines
and is carried out in the sidebar. The first special case is the theorem
with the column \(\ket{\psi}\) in the place of \(A\), and \(A\) in the
place of \(B\). For the second, a number is a matrix with one row and
one column, and its dagger is its conjugate; apply the theorem twice to
the product of three,
\(\bigl(\bra{\phi}\,A\,\ket{\psi}\bigr)^{\dagger}=\bigl(\ket{\psi}\bigr)^{\dagger}A^{\dagger}\bigl(\bra{\phi}\bigr)^{\dagger}=\bra{\psi}A^{\dagger}\ket{\phi}\).
\(\;\blacksquare\) See \texttt{Code/Ch3/Backwards.py}, which checks the
theorem on matrices drawn at random and checks that the two cheaper
recipes fail; and \texttt{Code/Ch3/Mirror.py} for the mirror of a single
box.

\begin{quote}
\textbf{Dirac's Blackboard.}

\emph{The mirror of a product, entry by entry.} Let \(A\) and \(B\) be
\(2\times2\); other shapes change only the ranges of the indices. The
entry of \(BA\) in row \(i\), column \(j\) is, by the rule ``along the
row and down the column'',

\[(BA)_{ij} \;=\; B_{i0}A_{0j}+B_{i1}A_{1j} .\]

By definition the entry of \((BA)^{\dagger}\) in row \(j\), column \(i\)
is the conjugate of this, and conjugation passes through sums and
products:

\[\bigl((BA)^{\dagger}\bigr)_{ji} \;=\; \overline{B_{i0}}\;\overline{A_{0j}}+\overline{B_{i1}}\;\overline{A_{1j}} .\]

Now read the right side again.
\(\overline{A_{kj}}=\bigl(A^{\dagger}\bigr)_{jk}\) and
\(\overline{B_{ik}}=\bigl(B^{\dagger}\bigr)_{ki}\), so it is

\[\bigl(A^{\dagger}\bigr)_{j0}\bigl(B^{\dagger}\bigr)_{0i}+\bigl(A^{\dagger}\bigr)_{j1}\bigl(B^{\dagger}\bigr)_{1i} \;=\; \bigl(A^{\dagger}B^{\dagger}\bigr)_{ji} .\]

Numbers may be multiplied in either order; matrices may not. The index
\(k\) that joined the \emph{back} of \(B\) to the \emph{front} of \(A\)
now joins the back of \(A^{\dagger}\) to the front of \(B^{\dagger}\),
and that is the whole reason for the reversal. \(\;\blacksquare\)
\end{quote}

In the sketchbook the theorem is the rule for reflecting any picture
that this book has so far drawn. \textbf{Mirror every piece, and reverse
the order.} A state becomes an effect (§2.2), an effect a state, a box
its mirror image, and what was first is last. Here it is for a state
followed by two boxes,

\[\text{the mirror image of}\;\;
\begin{tikzpicture}[sd, baseline=-0.55ex]
  \draw[wire] (0,0) -- (7.5,0);
  \node[sdstate, text height=1.5ex, text depth=0.5ex] at (0,0) {$\psi$};
  \node[sdbox, text height=1.5ex, text depth=0.5ex] at (2.8,0) {$A$};
  \node[sdbox, text height=1.5ex, text depth=0.5ex] at (5.4,0) {$B$};
\end{tikzpicture}
\;\;\text{is}\;\;
\begin{tikzpicture}[sd, baseline=-0.55ex]
  \draw[wire] (0,0) -- (7.5,0);
  \node[sdbox, text height=1.5ex, text depth=0.5ex] at (2.1,0) {$B^{\dagger}$};
  \node[sdbox, text height=1.5ex, text depth=0.5ex] at (4.7,0) {$A^{\dagger}$};
  \node[sdeffect, text height=1.5ex, text depth=0.5ex] at (7.5,0) {$\psi$};
\end{tikzpicture}\]

(on the blackboard:
\(\bigl(BA\ket{\psi}\bigr)^{\dagger}=\bra{\psi}A^{\dagger}B^{\dagger}\)),
and for a closed picture:

\[\begin{tikzpicture}[sd, baseline=-0.55ex]
  \draw[wire] (0,0) -- (6,0);
  \node[sdstate, text height=1.5ex, text depth=0.5ex] at (0,0) {$\phi$};
  \node[sdbox, text height=1.5ex, text depth=0.5ex] at (3,0) {$A^{\dagger}$};
  \node[sdeffect, text height=1.5ex, text depth=0.5ex] at (6,0) {$\psi$};
\end{tikzpicture}
\;\;=\;\; \bra{\psi}A^{\dagger}\ket{\phi} \;\;=\;\; \overline{\bra{\phi}A\ket{\psi}} .\]

A worked instance, with a number that is not real so that the bar can be
seen. Take the arrow, the test piece and the effect \(+\). First
\(U\ket{a}=\tfrac1{10}\begin{pmatrix}7+7i\\-1+i\end{pmatrix}\), and then

\[\bra{+}U\ket{a} \;=\; \tfrac1{\sqrt2}\cdot\tfrac1{10}\bigl((7+7i)+(-1+i)\bigr) \;=\; \tfrac{6+8i}{10\sqrt2} \;=\; \tfrac{3+4i}{5\sqrt2} .\]

The mirrored picture is ``state \(+\), box \(U^{\dagger}\), effect
\(a\)''. From the matrix \(U^{\dagger}\) above,
\(U^{\dagger}\ket{+}=\tfrac1{2\sqrt2}\begin{pmatrix}2\\-2i\end{pmatrix}=\tfrac1{\sqrt2}\begin{pmatrix}1\\-i\end{pmatrix}\),
and

\[\bra{a}U^{\dagger}\ket{+} \;=\; \tfrac1{\sqrt2}\bigl(\tfrac35-\tfrac45 i\bigr) \;=\; \tfrac{3-4i}{5\sqrt2} .\]

Conjugates, as the theorem requires. With \(U^{T}\) or \(\overline{U}\)
in the mirrored picture they would not be.

\begin{center}\rule{0.5\linewidth}{0.5pt}\end{center}

\hypertarget{boxes-that-keep-the-promise}{%
\section{§3.3 Boxes That Keep the
Promise}\label{boxes-that-keep-the-promise}}

Now the sketchbook of §3.1 is no longer stuck. The state ``\(\psi\),
then \(A\)'' has the mirror image ``\(A^{\dagger}\), then the effect
\(\psi\)'', and a state meeting its own mirror image is one closed
picture with two boxes on it:

\[\begin{tikzpicture}[sd, baseline=-0.55ex]
  \draw[wire] (0,0) -- (9,0);
  \node[sdstate, text height=1.5ex, text depth=0.5ex] at (0,0) {$\psi$};
  \node[sdbox, text height=1.5ex, text depth=0.5ex] at (3,0) {$A$};
  \node[sdbox, text height=1.5ex, text depth=0.5ex] at (6,0) {$A^{\dagger}$};
  \node[sdeffect, text height=1.5ex, text depth=0.5ex] at (9,0) {$\psi$};
\end{tikzpicture}
\;\;=\;\; \bra{\psi}A^{\dagger}A\ket{\psi} .\]

This is the picture §2.10 was asking about. If it is \(1\) for every
state \(\psi\), the box \(A\) \textbf{keeps the promise} of Theorem 2.1:
whatever state goes in, a state comes out.

Try the two boxes we know. For NOT: \(X\) is real and symmetric, so
\(X^{\dagger}=X\), and \(X^{\dagger}X=XX=I\) by Theorem 1.1. The two
boxes in the middle of the picture are a bare wire, and the picture is
\(\braket{\psi|\psi}=1\) for every \(\psi\) at once, not only for the
arrow. For SET-0:

\[S_0^{\dagger}S_0 \;=\; \begin{pmatrix}1&0\\1&0\end{pmatrix}\begin{pmatrix}1&1\\0&0\end{pmatrix} \;=\; \begin{pmatrix}1&1\\1&1\end{pmatrix}, \qquad \bra{a}S_0^{\dagger}S_0\ket{a} \;=\; \bigl(\tfrac35+\tfrac45\bigr)^2 \;=\; \tfrac{49}{25} .\]

And for the test piece, which is neither real nor symmetric, multiply
out \(U^{\dagger}U\) with the matrices of §3.2. The four entries are

\[\tfrac14\bigl((1-i)(1+i)+(1+i)(1-i)\bigr)=1,\]

\[\tfrac14\bigl((1-i)(1+i)+(1+i)(-1+i)\bigr)=\tfrac14(2-2)=0,\]

and likewise \(0\) and \(1\) in the second row: \(U^{\dagger}U=I\). The
cheaper reflections do \emph{not} give the bare wire:

\[U^{T}U \;=\; \begin{pmatrix}0&i\\i&0\end{pmatrix}, \qquad\qquad \overline{U}\,U \;=\; \tfrac12\begin{pmatrix}1-i&1+i\\-1+i&1+i\end{pmatrix}.\]

In the examples that work, the promise is kept for the strongest
possible reason: the box followed by its mirror image \emph{is} the bare
wire. Could a box keep the promise for some weaker reason? It could not.

\textbf{On the blackboard.} A \(2\times2\) matrix \(U\) is
\textbf{unitary} when \(U^{\dagger}U=I\).

\begin{quote}
\textbf{Theorem 3.2 (Boxes That Keep the Promise).} For a \(2\times2\)
complex matrix \(U\) the following are equivalent:

\begin{enumerate}
\def\labelenumi{\arabic{enumi}.}
\tightlist
\item
  \(U^{\dagger}U=I\) (the box, then its mirror image, is the bare wire);
\item
  for every state \(\ket{\psi}\), the column \(U\ket{\psi}\) is a state:
  \(\bra{\psi}U^{\dagger}U\ket{\psi}=1\).
\end{enumerate}

When they hold, also \(UU^{\dagger}=I\) (the mirror image, then the box,
is the bare wire); the two columns of \(U\) are an orthonormal basis;
and \(U\) preserves every bracket: for all states \(\ket{\phi}\),
\(\ket{\psi}\), the bracket of \(U\ket{\phi}\) with \(U\ket{\psi}\) is
\(\braket{\phi|\psi}\).
\end{quote}

\emph{Proof idea.} One direction is a single application of \emph{The
Mirror of a Row}. For the other, a matrix is pinned down by a handful of
its closed pictures: two standard states fix the diagonal of
\(U^{\dagger}U\), and two further states, one real and one not, fix the
real and imaginary parts of what is off the diagonal.

\emph{Proof sketch.} (1 \(\Rightarrow\) 2, and the bracket.) By Theorem
3.1 the bra of \(U\ket{\phi}\) is \(\bra{\phi}U^{\dagger}\), so the
bracket of \(U\ket{\phi}\) with \(U\ket{\psi}\) is
\(\bra{\phi}U^{\dagger}U\ket{\psi}=\bra{\phi}I\ket{\psi}=\braket{\phi|\psi}\).
With \(\phi=\psi\) this is statement 2.

(2 \(\Rightarrow\) 1.) Write \(M=U^{\dagger}U\) with entries \(m_{ij}\),
and put four states into \(\bra{\psi}M\ket{\psi}=1\). The states
\(\ket{0}\) and \(\ket{1}\) give \(m_{00}=1\) and \(m_{11}=1\). The
state \(\ket{+}\) gives \(\tfrac12(m_{00}+m_{01}+m_{10}+m_{11})=1\),
hence \(m_{01}+m_{10}=0\). The state
\(\ket{\omega}=\tfrac1{\sqrt2}\bigl(\ket{0}+i\ket{1}\bigr)\), whose bra
is \(\tfrac1{\sqrt2}\begin{pmatrix}1&-i\end{pmatrix}\), gives
\(\tfrac12(m_{00}+i\,m_{01}-i\,m_{10}+m_{11})=1\), hence
\(m_{01}-m_{10}=0\). So \(m_{01}=m_{10}=0\) and \(M=I\).

(The rest.) For square matrices a one-sided inverse is two-sided (we
used this in Lemma 2.3), so \(U^{\dagger}U=I\) gives \(UU^{\dagger}=I\).
And the entry \(m_{ij}\) of \(U^{\dagger}U\) is
\(\bra{i}U^{\dagger}U\ket{j}\), the bracket of the \(i\)-th column of
\(U\) with the \(j\)-th; \(M=I\) says these brackets are \(1,0,0,1\),
which is orthonormality (§2.5). \(\;\blacksquare\)

See \texttt{Code/Ch3/Mirror.py}, which checks all of this on a unitary
drawn at random and checks that the transpose and the conjugate are
\emph{not} its undoing, and \texttt{Code/Ch3/NoMirror.py} for SET-0.

For the test piece, take the factor \(\tfrac{1+i}{\sqrt2}=e^{i\pi/4}\)
out of each column of \(U\), using \(\tfrac{1-i}{1+i}=-i\) and
\(\tfrac{-1+i}{1+i}=i\): the columns are
\(e^{i\pi/4}\cdot\tfrac1{\sqrt2}\bigl(\ket{0}-i\ket{1}\bigr)\) and
\(e^{i\pi/4}\cdot\tfrac1{\sqrt2}\bigl(\ket{0}+i\ket{1}\bigr)=e^{i\pi/4}\ket{\omega}\).
Their bracket is \(\tfrac12\bigl(1+(i)(i)\bigr)=0\): an orthonormal
basis, and not one we have met before, each member carrying a global
phase.

Two remarks keep the theorem honest. First, ``has an undoing'' is weaker
than ``is undone by its mirror image''. The matrix
\(D=\begin{pmatrix}1&0\\0&2\end{pmatrix}\) is undone by
\(\begin{pmatrix}1&0\\0&1/2\end{pmatrix}\), but
\(D^{\dagger}D=\begin{pmatrix}1&0\\0&4\end{pmatrix}\), and
\(\bra{a}D^{\dagger}D\ket{a}=\tfrac9{25}+\tfrac{64}{25}=\tfrac{73}{25}\).
It breaks the promise, and may not sit on a qubit's wire. SET-0 is
worse: it sends both \(\ket{0}\) and \(\ket{1}\) to \(\ket{0}\), so
nothing placed after it can recover which one went in (§1.3), and it has
no undoing of any kind.

Second, §2.6 left a debt. It showed that a global phase is seen by no
question asked directly of the state, and said that this survives
whatever boxes are placed on the wire first. It does. Matrices act on a
multiple of a column by giving the same multiple of the result, so
\(U\bigl(e^{i\theta}\ket{\psi}\bigr)=e^{i\theta}\,U\ket{\psi}\); if
\(U\) is unitary then \(U\ket{\psi}\) is a state, and Proposition 2.4,
applied to \emph{that} state, says no question separates it from its
multiple. The same argument covers a phase on a box: \(e^{i\theta}U\) is
a different matrix from \(U\) when \(e^{i\theta}\neq1\), it is again
unitary, and it sends each state to \(e^{i\theta}\) times where \(U\)
sends it. Different boxes; and no question put afterwards can tell which
of them was on the wire.

\textbf{In the sketchbook.} Unitarity is a rewrite rule, the first this
book has had since Theorem 1.1, and it contains that theorem as a
special case:

\[\begin{tikzpicture}[sd, baseline=-0.55ex]
  \draw[wire] (0,0) -- (7,0);
  \node[sdbox, text height=1.5ex, text depth=0.5ex] at (2.1,0) {$U$};
  \node[sdbox, text height=1.5ex, text depth=0.5ex] at (4.9,0) {$U^{\dagger}$};
\end{tikzpicture}
\;\;=\;\;
\begin{tikzpicture}[sd, baseline=-0.55ex]
  \draw[wire] (0,0) -- (4,0);
\end{tikzpicture}
\;\;=\;\;
\begin{tikzpicture}[sd, baseline=-0.55ex]
  \draw[wire] (0,0) -- (7,0);
  \node[sdbox, text height=1.5ex, text depth=0.5ex] at (2.1,0) {$U^{\dagger}$};
  \node[sdbox, text height=1.5ex, text depth=0.5ex] at (4.9,0) {$U$};
\end{tikzpicture}\]

\textbf{A box meeting its reflection, in either order, leaves a bare
wire.} We call the rule \textbf{\emph{a box meets its mirror}}. It is
licensed by Theorem 3.2 and by nothing else, and it may be used only for
a box known to be unitary.

The second half of §2.10's question was about \emph{The Echo}. Put the
box between the pieces of each half, remembering that the right half is
the mirror image of the left, so that by \emph{The Mirror of a Row} it
carries \(U^{\dagger}\), on the other side of the gap:

\[\begin{tikzpicture}[sd, baseline=-0.55ex]
  \draw[wire] (0,0) -- (6,0);
  \draw[wire] (9,0) -- (15,0);
  \node[sdstate, text height=1.5ex, text depth=0.5ex] at (0,0) {$\psi$};
  \node[sdbox, text height=1.5ex, text depth=0.5ex] at (3,0) {$U$};
  \node[sdeffect, text height=1.5ex, text depth=0.5ex] at (6,0) {$b$};
  \node[sdstate, text height=1.5ex, text depth=0.5ex] at (9,0) {$b$};
  \node[sdbox, text height=1.5ex, text depth=0.5ex] at (12,0) {$U^{\dagger}$};
  \node[sdeffect, text height=1.5ex, text depth=0.5ex] at (15,0) {$\psi$};
\end{tikzpicture}
\;=\; \bra{\psi}U^{\dagger}\ket{b}\,\bra{b}U\ket{\psi}\]

The number is \(\lvert\bra{b}U\ket{\psi}\rvert^2\): by the Born rule,
the chance of the answer \(b\) when the box has acted first. Do the
chances still make one whole? Add the numbers for \(b=0\) and \(b=1\) (a
sum of \emph{numbers}, the liberty §2.5 allowed). They are the squared
moduli of the two entries of the column \(U\ket{\psi}\), and their sum
is that column meeting its own mirror image:
\(\bra{\psi}U^{\dagger}U\ket{\psi}\), the closed picture at the head of
this section. For a unitary box that is \(1\), and then \(U\ket{\psi}\)
is a state and Corollary 2.2 makes the total \(1\) for every other
question too. For SET-0 and the arrow the two numbers are
\(\tfrac{49}{25}\) and \(0\). ``Chances'' that come to nearly two are
not chances.

\begin{quote}
\textbf{Penrose's Sketchbook.}

\emph{A reflection that is not an undoing.} The rule \emph{a box meets
its mirror} is not a fact about mirrors. Every box has a mirror image;
only some are undone by it. Here is SET-0 meeting its reflection:

\[\begin{tikzpicture}[sd, baseline=-0.55ex]
  \draw[wire] (0,0) -- (7,0);
  \node[sdbox, text height=1.5ex, text depth=0.5ex] at (2.1,0) {$S_0$};
  \node[sdbox, text height=1.5ex, text depth=0.5ex] at (4.9,0) {$S_0^{\dagger}$};
\end{tikzpicture}
\;\;=\;\;
\begin{tikzpicture}[sd, baseline=-0.55ex]
  \draw[wire] (0,0) -- (5,0);
  \node[sdbox, inner xsep=0.5em] at (2.5,0) {$\begin{smallmatrix}1&1\\1&1\end{smallmatrix}$};
\end{tikzpicture}
\;\;\neq\;\;
\begin{tikzpicture}[sd, baseline=-0.55ex]
  \draw[wire] (0,0) -- (4,0);
\end{tikzpicture}\]

The middle picture is a box whose matrix has all four entries \(1\). In
the other order, \(S_0^{\dagger}\) then \(S_0\), the matrix is
\(\begin{pmatrix}2&0\\0&0\end{pmatrix}\): not a bare wire either, and
not even the same failure. So when a row contains a box and its
reflection side by side, do not cancel them by eye. Cancel them because
the box is unitary, and say so. The sketchbook cannot tell a unitary box
from any other by its outline; that knowledge comes from the blackboard,
through Theorem 3.2.
\end{quote}

\textbf{\emph{Why this definition.}} Why demand \(U^{\dagger}U=I\), when
the requirement we actually care about is only that states go to states?
Because they are the same requirement (Theorem 3.2), and the first can
be checked by one matrix product, while the second speaks of infinitely
many states. And why the mirror image, not merely \emph{some} undoing?
Because \(D\) above has an undoing and still turns the arrow into a
column of squared length \(\tfrac{73}{25}\).

\begin{center}\rule{0.5\linewidth}{0.5pt}\end{center}

\hypertarget{six-boxes}{%
\section{§3.4 Six Boxes}\label{six-boxes}}

It is time to meet the boxes that the rest of the book is built from.
There are six, and we already know the first.

\textbf{On the blackboard.}

\[X=\begin{pmatrix}0&1\\1&0\end{pmatrix}, \qquad Y=\begin{pmatrix}0&-i\\i&0\end{pmatrix}, \qquad Z=\begin{pmatrix}1&0\\0&-1\end{pmatrix},\]

\[H=\tfrac1{\sqrt2}\begin{pmatrix}1&1\\1&-1\end{pmatrix}, \qquad S=\begin{pmatrix}1&0\\0&i\end{pmatrix}, \qquad T=\begin{pmatrix}1&0\\0&e^{i\pi/4}\end{pmatrix},\]

where \(e^{i\pi/4}=\tfrac{1+i}{\sqrt2}\).

The first three are called the \textbf{Pauli matrices}; \(H\) is the
\textbf{Hadamard} box; \(S\) and \(T\) are the \textbf{phase} boxes.
(\(S\) is the \(S\) of §3.2, and the matrix \(V\) there was \(HSH\), as
the reader can now check.)

\emph{Four are their own mirror images.} \(X\), \(Z\) and \(H\) are real
and symmetric, so neither the bar nor the flip changes them. \(Y\) is
more interesting: it is \emph{not} real and \emph{not} symmetric. Its
transpose is \(-Y\) and its conjugate is \(-Y\); the two signs cancel,
and \(Y^{\dagger}=Y\). Each of the four is unitary, and a unitary box
that is its own mirror image is its own undoing:

\[XX=I, \qquad YY=\begin{pmatrix}(-i)(i)&0\\0&(i)(-i)\end{pmatrix}=I, \qquad ZZ=I,\]

\[HH=\tfrac12\begin{pmatrix}1+1&1-1\\1-1&1+1\end{pmatrix}=I .\]

\emph{Two are not.}
\(S^{\dagger}=\begin{pmatrix}1&0\\0&-i\end{pmatrix}\neq S\) and
\(T^{\dagger}=\begin{pmatrix}1&0\\0&e^{-i\pi/4}\end{pmatrix}\neq T\).
Both are unitary, since \(\bar\imath\, i=1\) and
\(e^{-i\pi/4}e^{i\pi/4}=1\). But done twice they are not the bare wire.
They climb a ladder instead:

\[T^2=S, \qquad S^2=Z, \qquad Z^2=I .\]

\emph{They are related.} Multiply out:
\(XZ=\begin{pmatrix}0&-1\\1&0\end{pmatrix}\) and
\(ZX=\begin{pmatrix}0&1\\-1&0\end{pmatrix}\), so

\[XZ \;=\; -ZX, \qquad\qquad Y \;=\; iXZ .\]

The row ``\(Z\), then \(X\)'' and the row ``\(X\), then \(Z\)'' are
different boxes; they differ by a sign on the whole box. By the second
remark of §3.3, no question asked afterwards can tell them apart.

\emph{What \(H\) and \(Z\) do to the four states we have names for.}
These eight small facts are used constantly, and the sketchbook proof of
§3.8 will borrow them:

\[H\ket{0}=\tfrac1{\sqrt2}\begin{pmatrix}1\\1\end{pmatrix}=\ket{+}, \qquad H\ket{1}=\tfrac1{\sqrt2}\begin{pmatrix}1\\-1\end{pmatrix}=\ket{-},\]

\[Z\ket{+}=\ket{-}, \qquad Z\ket{-}=\ket{+},\]

the first two because \(H\ket{j}\) is the \(j\)-th column of \(H\), the
last two because \(Z\) changes the sign of the second entry. Apply \(H\)
to the first two and use \(HH=I\): \(H\ket{+}=\ket{0}\) and
\(H\ket{-}=\ket{1}\). And \(Z\ket{0}=\ket{0}\), \(Z\ket{1}=-\ket{1}\).
So \(H\) carries the standard basis to the turned basis of §2.5 and
back; \textbf{\(Z\) does to the turned basis what NOT does to the
standard one}; and to the standard basis \(Z\) does nothing that a
standard question can see, which is §2.4's hidden sign, now in a box.

Take the dagger of \(H\ket{0}=\ket{+}\). By \emph{The Mirror of a Row}
and \(H^{\dagger}=H\):

\[\bra{0}H=\bra{+}, \qquad\qquad \bra{1}H=\bra{-} .\]

\textbf{In the sketchbook.} No new element: six labelled boxes. Four of
them are their own mirror images (\(X^{\dagger}\) \emph{is} the box
\(X\), and we shall never write the dagger on it) and two are not. The
facts just computed are picture equations between states, in the manner
of Chapter 1's \emph{NOT on 0}, and we name them the same way.
\textbf{\emph{H on 0}} and \textbf{\emph{H on 1}}:

\[\begin{tikzpicture}[sd, baseline=-0.55ex]
  \draw[wire] (0,0) -- (5,0);
  \node[sdstate, text height=1.5ex, text depth=0.5ex] at (0,0) {$0$};
  \node[sdbox, text height=1.5ex, text depth=0.5ex] at (2.8,0) {$H$};
\end{tikzpicture}
\;=\;
\begin{tikzpicture}[sd, baseline=-0.55ex]
  \draw[wire] (0,0) -- (2.6,0);
  \node[sdstate, text height=1.5ex, text depth=0.5ex] at (0,0) {$+$};
\end{tikzpicture}
\qquad\qquad
\begin{tikzpicture}[sd, baseline=-0.55ex]
  \draw[wire] (0,0) -- (5,0);
  \node[sdstate, text height=1.5ex, text depth=0.5ex] at (0,0) {$1$};
  \node[sdbox, text height=1.5ex, text depth=0.5ex] at (2.8,0) {$H$};
\end{tikzpicture}
\;=\;
\begin{tikzpicture}[sd, baseline=-0.55ex]
  \draw[wire] (0,0) -- (2.6,0);
  \node[sdstate, text height=1.5ex, text depth=0.5ex] at (0,0) {$-$};
\end{tikzpicture}\]

\textbf{\emph{Z on +}} and \textbf{\emph{Z on \(-\)}}:

\[\begin{tikzpicture}[sd, baseline=-0.55ex]
  \draw[wire] (0,0) -- (5,0);
  \node[sdstate, text height=1.5ex, text depth=0.5ex] at (0,0) {$+$};
  \node[sdbox, text height=1.5ex, text depth=0.5ex] at (2.8,0) {$Z$};
\end{tikzpicture}
\;=\;
\begin{tikzpicture}[sd, baseline=-0.55ex]
  \draw[wire] (0,0) -- (2.6,0);
  \node[sdstate, text height=1.5ex, text depth=0.5ex] at (0,0) {$-$};
\end{tikzpicture}
\qquad\qquad
\begin{tikzpicture}[sd, baseline=-0.55ex]
  \draw[wire] (0,0) -- (5,0);
  \node[sdstate, text height=1.5ex, text depth=0.5ex] at (0,0) {$-$};
  \node[sdbox, text height=1.5ex, text depth=0.5ex] at (2.8,0) {$Z$};
\end{tikzpicture}
\;=\;
\begin{tikzpicture}[sd, baseline=-0.55ex]
  \draw[wire] (0,0) -- (2.6,0);
  \node[sdstate, text height=1.5ex, text depth=0.5ex] at (0,0) {$+$};
\end{tikzpicture}\]

And the mirror image of \emph{H on 0}, every piece reflected and the
order reversed, is an equation between effects, with \(H\) its own
reflection:

\[\begin{tikzpicture}[sd, baseline=-0.55ex]
  \draw[wire] (0,0) -- (5,0);
  \node[sdbox, text height=1.5ex, text depth=0.5ex] at (2.2,0) {$H$};
  \node[sdeffect, text height=1.5ex, text depth=0.5ex] at (5,0) {$0$};
\end{tikzpicture}
\;=\;
\begin{tikzpicture}[sd, baseline=-0.55ex]
  \draw[wire] (0,0) -- (2.6,0);
  \node[sdeffect, text height=1.5ex, text depth=0.5ex] at (2.6,0) {$+$};
\end{tikzpicture}\]

This last equation says what \(H\) is \emph{for}. To ask the turned
question of §2.5, put an \(H\) on the wire and then ask the standard
one.

\textbf{\emph{Why this definition.}} Why these six, out of all the
unitary matrices there are? For this chapter the answer is modest:
\(X\), \(Y\) and \(Z\) will be the three axes of the sphere (§3.5);
\(H\) passes between two of them; \(S\) and \(T\) are the simplest boxes
that are \emph{not} their own mirror images, without which nothing in
this chapter could be tested. That a great deal can be built from them
is a claim for later chapters, and we do not make it here.

\begin{quote}
\textbf{Feynman's Notebook.}

\emph{Names that everybody uses.} The letters in this section are not
ours. Van de Wetering's survey, written for working quantum computer
scientists, opens its account of circuits (§2.1) by recording that the
community has mostly agreed on its naming conventions: \(X\), \(Y\) and
\(Z\) are the Pauli matrices; the \(S\) gate and the \(T\) gate are the
quarter- and eighth-turn members of the family of \(Z\)-rotations (§3.6
below says what the turning is); \(H\) is the Hadamard gate; and a
dagger on a gate's name means the inverse gate. That last convention is
Theorem 3.2 in everyday use: nobody who writes \(S^{\dagger}\) in a
circuit stops to say whether they mean the mirror image or the undoing,
because for a unitary box there is no difference. ``Gate'' is the
engineers' word for what we call a box. The textbook treatment of the
Pauli matrices, of \(H\) and of the sphere they turn is Nielsen \&
Chuang, ch.~4.
\end{quote}

\hypertarget{a-trap-2}{%
\subsection{A trap}\label{a-trap-2}}

Here are two tempting rules. \emph{To undo a row, play it again from the
other end. And the order of two boxes never really matters.}

Both are true of everything a cautious person might try first. Take any
row made of \(X\), \(Z\) and \(H\) --- say ``\(H\), then \(Z\)''. Its
mirror image is ``\(Z^{\dagger}\), then \(H^{\dagger}\)'', which is
``\(Z\), then \(H\)'': the same boxes from the other end, because each
is its own reflection. And \(X\) with \(Z\) in the two orders differ
only by a sign on the box, which no question sees.

Now try \(S\). ``Play it again'' gives \(S\), then \(S\), which is \(Z\)
and not the bare wire. Play our test row ``\(S\), then \(V\)'' again
from the other end, ``\(V\), then \(S\)'': the four boxes multiply to
\(S\,V\,V\,S=U^{T}U=iX\) (§3.3), nowhere near a bare wire. There is a
subtler version. Wrap \(X\) between \(S\) and its mirror image, and
between \(S\) and itself:

\[S\,X\,S^{\dagger} \;=\; \begin{pmatrix}0&-i\\i&0\end{pmatrix} \;=\; Y, \qquad\qquad S\,X\,S \;=\; \begin{pmatrix}0&i\\i&0\end{pmatrix} \;=\; iX .\]

The first is a new box; the second is \(X\) again with a phase on it.
Someone who had only ever wrapped things in \(H\), where
\(H^{\dagger}=H\) hides the distinction, would not know that there was a
choice to make.

As for order:
\(HZ=\tfrac1{\sqrt2}\begin{pmatrix}1&-1\\1&1\end{pmatrix}\) and
\(ZH=\tfrac1{\sqrt2}\begin{pmatrix}1&1\\-1&1\end{pmatrix}\), and no sign
or phase turns one into the other, since their top-left entries agree
and their top-right entries do not. The moral is the moral of §1.4 over
again. \(X\), \(Z\) and \(H\) are real and symmetric, and \textbf{a real
symmetric box can never test a claim about mirrors.} See
\texttt{Code/Ch3/Mirror.py} and \texttt{Code/Ch3/Backwards.py}.

\begin{center}\rule{0.5\linewidth}{0.5pt}\end{center}

\hypertarget{the-sphere}{%
\section{§3.5 The Sphere}\label{the-sphere}}

A qubit state is two complex numbers, which is four real ones; the rule
\(|\alpha|^2+|\beta|^2=1\) removes one, and the global phase, which no
question sees, removes another. Two real numbers are left. That is as
many as it takes to say where you are on the surface of the Earth. This
section shows that the resemblance is exact. It is a blackboard section:
the sphere is a way of reading numbers, and we add no sphere to the
sketchbook's elements.

\textbf{On the blackboard.} For a state
\(\ket{\psi}=\alpha\ket{0}+\beta\ket{1}\) define three numbers, one for
each Pauli matrix:

\[x=\bra{\psi}X\ket{\psi}, \qquad y=\bra{\psi}Y\ket{\psi}, \qquad z=\bra{\psi}Z\ket{\psi} .\]

Multiplying out,

\[x=\bar\alpha\beta+\bar\beta\alpha, \qquad y=-i\,\bar\alpha\beta+i\,\bar\beta\alpha, \qquad z=|\alpha|^2-|\beta|^2 .\]

Write \(w=\bar\alpha\beta\). Then \(\bar\beta\alpha=\bar w\), and a
number plus its conjugate is twice its real part, while \(-i(w-\bar w)\)
is twice its imaginary part. So

\[x=2\,\mathrm{Re}(\bar\alpha\beta), \qquad y=2\,\mathrm{Im}(\bar\alpha\beta), \qquad z=|\alpha|^2-|\beta|^2 .\]

\begin{quote}
\textbf{Proposition 3.3 (the point of a state).} For every state
\(\ket{\psi}\) the numbers \(x,y,z\) are real, and

\[x^2+y^2+z^2 \;=\; 1 .\]

The states \(\ket{\psi}\) and \(e^{i\theta}\ket{\psi}\) have the same
three numbers. And every point \((x,y,z)\) with \(x^2+y^2+z^2=1\)
belongs to some state.
\end{quote}

So each state marks a point on the sphere of radius \(1\) about the
origin --- it is called the \textbf{Bloch sphere} (Nielsen \& Chuang,
ch.~4) --- and the point does not notice a global phase.

\emph{Proof idea.} Reality is a mirror argument; the sum of squares is
the identity \((p-q)^2+4pq=(p+q)^2\). \emph{Proof sketch.} The formulas
above show \(x,y,z\) real. Next,
\(x^2+y^2=4\bigl(\mathrm{Re}(w)^2+\mathrm{Im}(w)^2\bigr)=4|w|^2=4|\alpha|^2|\beta|^2\),
so

\[x^2+y^2+z^2 \;=\; 4|\alpha|^2|\beta|^2+\bigl(|\alpha|^2-|\beta|^2\bigr)^2 \;=\; \bigl(|\alpha|^2+|\beta|^2\bigr)^2 \;=\; 1 .\]

A global phase multiplies \(\alpha\) and \(\beta\) by \(e^{i\theta}\)
and their conjugates by \(e^{-i\theta}\), and in each of
\(\bar\alpha\beta\), \(|\alpha|^2\), \(|\beta|^2\) the two factors
cancel. Last, given a point of the sphere, write it as
\((\sin\vartheta\cos\varphi,\;\sin\vartheta\sin\varphi,\;\cos\vartheta)\),
as one does with latitude and longitude, and take

\[\ket{\psi} \;=\; \cos\tfrac\vartheta2\,\ket{0}+e^{i\varphi}\sin\tfrac\vartheta2\,\ket{1} .\]

Then
\(\bar\alpha\beta=\cos\tfrac\vartheta2\sin\tfrac\vartheta2\,e^{i\varphi}=\tfrac12\sin\vartheta\,e^{i\varphi}\)
and \(z=\cos^2\tfrac\vartheta2-\sin^2\tfrac\vartheta2=\cos\vartheta\),
which is the given point. \(\;\blacksquare\) See
\texttt{Code/Ch3/Sandwich.py}, which checks the proposition on states
drawn at random.

Where do the states we know sit?

\[\ket{0}:\;(0,0,1), \qquad \ket{1}:\;(0,0,-1),\]

\[\ket{+}:\;(1,0,0), \qquad \ket{-}:\;(-1,0,0), \qquad \ket{\omega}:\;(0,1,0),\]

where \(\ket{\omega}=\tfrac1{\sqrt2}(\ket{0}+i\ket{1})\) is the fourth
test state of §3.3. The arrow has \(\bar\alpha\beta=\tfrac{12}{25}\) and
\(|\alpha|^2-|\beta|^2=-\tfrac7{25}\):

\[\ket{a}:\;\bigl(\tfrac{24}{25},\;0,\;-\tfrac7{25}\bigr) .\]

There is the \(\tfrac{24}{25}\) of §3.1. The ``literal'' picture with
the arrow at both ends and \(X\) between was never a broken promise; it
was the arrow's \(x\)-coordinate. Chapter 2's non-real state
\(\ket{\chi}=\begin{pmatrix}3/5\\4i/5\end{pmatrix}\) has
\(\bar\alpha\beta=\tfrac{12i}{25}\) and sits at
\(\bigl(0,\tfrac{24}{25},-\tfrac7{25}\bigr)\): the same latitude as the
arrow, a quarter of the way round. The real states all lie on the one
great circle \(y=0\); the \(i\) in an amplitude is what takes a state
off it.

\textbf{An honest word about Chapter 2's paper arrows.} On squared
paper, \(\ket{0}\) and \(\ket{1}\) were perpendicular, and so were
\(\ket{+}\) and \(\ket{-}\). On the sphere each pair sits at
\emph{opposite poles}. The arrows \(\ket{0}\) and \(\ket{+}\) were half
a right angle apart on paper; their points \((0,0,1)\) and \((1,0,0)\)
are a full right angle apart. The sphere doubles every angle of the
paper picture. That is the half-angle \(\tfrac\vartheta2\) in the
formula above, and it is why perpendicular states (the two answers of a
question, §2.5) are opposite points. Neither picture is wrong. The paper
arrow kept the sign of the whole column, which no question sees, and
could show only real states. The sphere shows every state and nothing
that cannot be seen.

\textbf{In the sketchbook.} Each of the three numbers is a closed
picture of the kind §3.1 taught us to read:

\[\begin{tikzpicture}[sd, baseline=-0.55ex]
  \draw[wire] (0,0) -- (6,0);
  \node[sdstate, text height=1.5ex, text depth=0.5ex] at (0,0) {$\psi$};
  \node[sdbox, text height=1.5ex, text depth=0.5ex] at (3,0) {$X$};
  \node[sdeffect, text height=1.5ex, text depth=0.5ex] at (6,0) {$\psi$};
\end{tikzpicture}
\;\;=\;\; x\]

and the same with \(Y\) for \(y\), and with \(Z\) for \(z\). The
sketchbook can say, in its own terms, why they are real. Hold the mirror
up to the picture above. The state \(\psi\) and the effect \(\psi\)
change places, and \(X\) is its own mirror image: the reflected picture
\emph{is the picture we started with}. But the number of a reflected
picture is the conjugate (Theorem 3.1). A number equal to its own
conjugate is real. The same argument serves for any box that is its own
reflection, and it is a better reason than the formulas, which only
happen to come out real.

One caution. The three numbers are three separate pictures. We do not
set them side by side in a row, because a row of closed pictures is a
\emph{product} (§2.3), and the product \(xyz\) means nothing here. The
sum of their squares is \(1\); that is arithmetic on three numbers, done
on the blackboard.

\textbf{\emph{Why this definition.}} Why these three numbers and not,
say, the two amplitudes themselves? Because the amplitudes change with
the global phase and these do not: each is a closed picture with the
state at one end and its mirror image at the other, and whatever phase
the state gains its reflection gives back (§2.6). The point on the
sphere is the state with exactly the unseeable part removed.

\begin{center}\rule{0.5\linewidth}{0.5pt}\end{center}

\hypertarget{turning}{%
\section{§3.6 Turning}\label{turning}}

What does a box do to the point?

Start with an example. The arrow sits at
\(\bigl(\tfrac{24}{25},0,-\tfrac7{25}\bigr)\). Send it through \(H\):
§2.5 already computed the amplitudes,
\(H\ket{a}=\tfrac1{5\sqrt2}\begin{pmatrix}7\\-1\end{pmatrix}\), so
\(\bar\alpha\beta=-\tfrac7{50}\) and
\(|\alpha|^2-|\beta|^2=\tfrac{49}{50}-\tfrac1{50}\), and the new point
is

\[H\ket{a}:\;\bigl(-\tfrac7{25},\;0,\;\tfrac{24}{25}\bigr) .\]

The first and third coordinates have changed places. That is too tidy to
be luck, and the reason can be found without choosing a state at all.
The new \(x\)-coordinate is, by \emph{The Mirror of a Row} and
\(H^{\dagger}=H\),

\[\bigl(\bra{\psi}H\bigr)\,X\,\bigl(H\ket{\psi}\bigr) \;=\; \bra{\psi}\,HXH\,\ket{\psi},\]

the closed picture of the \emph{old} state with the row \(H\), \(X\),
\(H\) in the middle. If the new \(x\) is always the old \(z\), the row
\(H\), \(X\), \(H\) must be the box \(Z\); and if the new \(z\) is the
old \(x\), the row \(H\), \(Z\), \(H\) must be the box \(X\).

\begin{quote}
\textbf{Theorem 3.4 (The Sandwich).} \(HZH=X\), exactly. In pictures:
the row of boxes \(H\), \(Z\), \(H\) on one wire is the box \(X\).
\end{quote}

\[\begin{tikzpicture}[sd, baseline=-0.55ex]
  \draw[wire] (0,0) -- (9,0);
  \node[sdbox, text height=1.5ex, text depth=0.5ex] at (2,0) {$H$};
  \node[sdbox, text height=1.5ex, text depth=0.5ex] at (4.5,0) {$Z$};
  \node[sdbox, text height=1.5ex, text depth=0.5ex] at (7,0) {$H$};
\end{tikzpicture}
\;\;=\;\;
\begin{tikzpicture}[sd, baseline=-0.55ex]
  \draw[wire] (0,0) -- (5,0);
  \node[sdbox, text height=1.5ex, text depth=0.5ex] at (2.5,0) {$X$};
\end{tikzpicture}\]

\emph{Proof idea.} \(H\) turns the standard question into the turned one
and back again; in between, \(Z\) does to the turned basis what NOT does
to the standard one (§3.4). This is the headline theorem of the chapter,
and §3.8 proves it twice in full.

\begin{quote}
\textbf{Corollary 3.5 (the other two sandwiches).} \(HXH=Z\) and
\(HYH=-Y\), exactly.
\end{quote}

\emph{Proof.} Wrap both sides of Theorem 3.4 in \(H\) and use \(HH=I\):
\(Z=H(HZH)H=HXH\). For \(Y=iXZ\), slip \(HH=I\) between the factors:
\(HYH=i\,(HXH)(HZH)=i\,ZX=-i\,XZ=-Y\), using \(ZX=-XZ\).
\(\;\blacksquare\)

\textbf{On the blackboard.} Now we can say what several boxes do to the
point of \emph{every} state. Two families of boxes are needed, with a
real number \(\theta\) for a dial:

\[R_z(\theta)=\begin{pmatrix}e^{-i\theta/2}&0\\0&e^{i\theta/2}\end{pmatrix},\]

\[R_x(\theta)=H\,R_z(\theta)\,H=\begin{pmatrix}\cos\tfrac\theta2&-i\sin\tfrac\theta2\\-i\sin\tfrac\theta2&\cos\tfrac\theta2\end{pmatrix}.\]

Both are unitary: \(R_z(\theta)^{\dagger}=R_z(-\theta)\) undoes
\(R_z(\theta)\), and the mirror image of the row \(H\), \(R_z(\theta)\),
\(H\) is the row \(H\), \(R_z(-\theta)\), \(H\), which undoes it box by
box.

\begin{quote}
\textbf{Proposition 3.6 (turns).} Let \(\ket{\psi}\) be any state, with
point \((x,y,z)\).

\begin{enumerate}
\def\labelenumi{\arabic{enumi}.}
\tightlist
\item
  The point of \(H\ket{\psi}\) is \((z,\,-y,\,x)\).
\item
  The point of \(R_z(\theta)\ket{\psi}\) is
  \((x\cos\theta-y\sin\theta,\;\; x\sin\theta+y\cos\theta,\;\; z)\): the
  sphere is turned through the angle \(\theta\) about the \(z\)-axis.
\item
  The point of \(R_x(\theta)\ket{\psi}\) is
  \((x,\;\; y\cos\theta-z\sin\theta,\;\; y\sin\theta+z\cos\theta)\): a
  turn through \(\theta\) about the \(x\)-axis.
\item
  If \(U\) is unitary, \(e^{i\delta}U\) moves every point exactly as
  \(U\) does.
\end{enumerate}
\end{quote}

\emph{Proof idea.} The new coordinate for the Pauli matrix \(P\) is
\(\bra{\psi}U^{\dagger}PU\ket{\psi}\): wrap each Pauli matrix in the box
and its mirror image, and read off the result in terms of the old Pauli
matrices. \emph{Proof sketch.} (1) is Theorem 3.4 and Corollary 3.5. (2)
\(R_z(\theta)^{\dagger}\,X\,R_z(\theta)\) has zero diagonal and
off-diagonal entries \(e^{i\theta}\) (top right) and \(e^{-i\theta}\)
(bottom left); so has \(\cos\theta\,X-\sin\theta\,Y\), since
\(\cos\theta+i\sin\theta=e^{i\theta}\). Hence the new \(x\) is
\(x\cos\theta-y\sin\theta\). In the same way
\(R_z(\theta)^{\dagger}\,Y\,R_z(\theta)=\sin\theta\,X+\cos\theta\,Y\);
and \(Z\), being diagonal like \(R_z(\theta)\), is left alone. (3)
\(R_x(\theta)\) is the row \(H\), \(R_z(\theta)\), \(H\), so apply (1),
then (2), then (1): \((x,y,z)\) goes to \((z,-y,x)\), then to
\((z\cos\theta+y\sin\theta,\;z\sin\theta-y\cos\theta,\;x)\), then to
\((x,\;y\cos\theta-z\sin\theta,\;y\sin\theta+z\cos\theta)\). (4) In
\(\bigl(e^{i\delta}U\bigr)^{\dagger}P\bigl(e^{i\delta}U\bigr)\) the
factors \(e^{-i\delta}\) and \(e^{i\delta}\) cancel. \(\;\blacksquare\)
See \texttt{Code/Ch3/Sandwich.py}.

So \(H\) is a \textbf{half-turn about the direction midway between the
\(x\)- and \(z\)-axes}: that half-turn exchanges the two axes and
reverses the one perpendicular to both, which is what \((z,-y,x)\) says.
And our diagonal boxes are \(z\)-turns with a phase in front:

\[Z=e^{i\pi/2}R_z(\pi), \qquad S=e^{i\pi/4}R_z\bigl(\tfrac\pi2\bigr), \qquad T=e^{i\pi/8}R_z\bigl(\tfrac\pi4\bigr),\]

\[X=e^{i\pi/2}R_x(\pi), \qquad V=HSH=e^{i\pi/4}R_x\bigl(\tfrac\pi2\bigr).\]

\(Z\) is a half-turn about the \(z\)-axis, \(S\) a quarter-turn, \(T\)
an eighth; \(X\) is a half-turn about the \(x\)-axis and \(V\) a
quarter-turn about it. The ladder \(T^2=S\), \(S^2=Z\) of §3.4 is an
eighth and an eighth making a quarter, a quarter and a quarter making a
half. The mirror image \(S^{\dagger}\) is the quarter-turn the other way
--- the turn seen through the shop window. For a half-turn the two
senses agree, and that is how \(X\), \(Z\) and \(H\) come to be their
own mirror images, and why they could never test a claim about mirrors.

\emph{The Sandwich} can now be read on the sphere. A half-turn about
\(z\), performed between two exchanges of the \(x\)- and \(z\)-axes, is
a half-turn about \(x\): follow the point,
\((x,y,z)\mapsto(z,-y,x)\mapsto(-z,\,y,\,x)\mapsto(x,-y,-z)\), and the
last is what \(X\) does. And the sign in \(XZ=-ZX\) is explained: a
half-turn about \(x\) and a half-turn about \(z\), in either order, are
the same turn of the sphere, \((x,y,z)\mapsto(-x,y,-z)\), the half-turn
about \(y\); the two orders give boxes that differ only by the phase
\(-1\). For \(H\) and \(Z\) the two orders are different turns: \(ZH\)
sends the point of \(\ket{0}\) to \((-1,0,0)\), and \(HZ\) sends it to
\((1,0,0)\).

One honest sentence about the half-angle. \(R_z(2\pi)=-I\): a full turn
of the sphere brings every point home, and the box that performs it is
not the bare wire but the bare wire with a sign; it takes two full turns
to get the box itself back. The sphere cannot see the difference, by
part 4. Later chapters, with more than one wire, will.

\hypertarget{three-turns}{%
\subsection{Three turns}\label{three-turns}}

How much can be done with \(z\)-turns alone? Not much.
\(R_z(\theta)\ket{0}=e^{-i\theta/2}\ket{0}\): the point \((0,0,1)\)
never moves. With \(x\)-turns alone, \((1,0,0)\) never moves. But take
one of each kind, a quarter-turn apiece, in the row \(S\), \(V\), \(S\).
The first two boxes are our test piece, \(VS=U\), so the product is
\(SU\); multiplying on the left by \(S\) multiplies the second row by
\(i\):

\[S\,V\,S \;=\; S\,U \;=\; \tfrac12\begin{pmatrix}1+i&1+i\\i(1-i)&i(-1+i)\end{pmatrix} \;=\; \tfrac{1+i}2\begin{pmatrix}1&1\\1&-1\end{pmatrix} \;=\; e^{i\pi/4}\,H .\]

\textbf{A quarter-turn about \(z\), a quarter-turn about \(x\), a
quarter-turn about \(z\), is the half-turn \(H\)} --- exactly, with the
scalar \(e^{i\pi/4}\) written out. Removing the phases carried by \(S\)
and \(V\), the same equation reads
\(H=e^{i\pi/2}R_z(\tfrac\pi2)R_x(\tfrac\pi2)R_z(\tfrac\pi2)\). Reflect
the whole equation: \(S^{\dagger}V^{\dagger}S^{\dagger}=e^{-i\pi/4}H\),
three quarter-turns the other way, and the same half-turn. So a box
about a slanting axis can be had from boxes about the two upright ones.
It needs three, and the row reads the same from either end. The general
fact is this.

\begin{quote}
\textbf{Theorem 3.7 (Three Turns; cited, not proved here).} For every
\(2\times2\) unitary matrix \(U\) there are real numbers
\(\alpha,\beta,\gamma,\delta\) with

\[U \;=\; e^{i\delta}\,R_z(\alpha)\,R_x(\beta)\,R_z(\gamma) .\]
\end{quote}

We prove only the instance above. The general statement is the
decomposition of a rotation into \emph{Euler angles}; it is standard
textbook material (Nielsen \& Chuang, ch.~4), and the
\(z\)--\(x\)--\(z\) form we use, together with the instance \(U=H\) with
all three angles equal to \(\tfrac\pi2\), is stated in van de Wetering's
survey (§3.6). The companion code does not prove it either, but it does
test it: \texttt{Code/Ch3/ThreeTurns.py} computes the four angles from
the four entries of a hundred unitaries drawn at random and checks the
equation each time.

Granting the theorem, Proposition 3.6 says what every unitary box does
to the sphere: three turns, one after another. Turning is all that a box
on a single qubit's wire can do.

\textbf{In the sketchbook.} The instance is a picture equation with a
scalar in it. A number beside a picture multiplies it, as in §2.6:

\[\begin{tikzpicture}[sd, baseline=-0.55ex]
  \draw[wire] (0,0) -- (9,0);
  \node[sdbox, text height=1.5ex, text depth=0.5ex] at (2,0) {$S$};
  \node[sdbox, text height=1.5ex, text depth=0.5ex] at (4.5,0) {$V$};
  \node[sdbox, text height=1.5ex, text depth=0.5ex] at (7,0) {$S$};
\end{tikzpicture}
\;\;=\;\; e^{i\pi/4}\;\;
\begin{tikzpicture}[sd, baseline=-0.55ex]
  \draw[wire] (0,0) -- (5,0);
  \node[sdbox, text height=1.5ex, text depth=0.5ex] at (2.5,0) {$H$};
\end{tikzpicture}\]

The scalar is part of the equation. Without it the two sides are
different boxes that turn the sphere alike, and ``turn the sphere
alike'' is a sentence about the sphere, not an equals sign between
pictures. The sketchbook has, at this stage, no rewrite that would
\emph{find} such an equation; it records what the blackboard has
multiplied out.

\textbf{\emph{Why this definition.}} Why write \(e^{-i\theta/2}\) and
\(e^{i\theta/2}\) in \(R_z(\theta)\), when
\(\begin{pmatrix}1&0\\0&e^{i\theta}\end{pmatrix}\) is simpler and turns
the sphere identically? For the symmetry: with the half-angles shared
out evenly, \(R_x(\theta)=HR_z(\theta)H\) comes out with plain cosines
and sines, and the mirror image is simply \(\theta\mapsto-\theta\). The
simpler matrix is the one the phase boxes \(S\) and \(T\) use, and the
price is the scalars \(e^{i\pi/4}\) and \(e^{i\pi/8}\) above. Both
conventions are in use; one must say which, every time.

\begin{center}\rule{0.5\linewidth}{0.5pt}\end{center}

\hypertarget{the-rosetta-table-2}{%
\section{§3.7 The Rosetta Table}\label{the-rosetta-table-2}}

This chapter adds seven rows to the running dictionary, numbered on from
Chapter 2's fifteen. Row 16 is the chapter's one new picture element.
Row 17 is the rule for reflecting a whole picture. Rows 18 to 21 are
things that can now be said without any new element, and row 22 is the
headline theorem.

\begingroup\renewcommand{\arraystretch}{1.5}

\begin{longtable}[]{@{}
  >{\raggedright\arraybackslash}p{(\columnwidth - 4\tabcolsep) * \real{0.0909}}
  >{\raggedright\arraybackslash}p{(\columnwidth - 4\tabcolsep) * \real{0.4545}}
  >{\raggedright\arraybackslash}p{(\columnwidth - 4\tabcolsep) * \real{0.4545}}@{}}
\toprule\noalign{}
\begin{minipage}[b]{\linewidth}\raggedright
\#
\end{minipage} & \begin{minipage}[b]{\linewidth}\raggedright
Algebra (Dirac)
\end{minipage} & \begin{minipage}[b]{\linewidth}\raggedright
Picture (Penrose)
\end{minipage} \\
\midrule\noalign{}
\endhead
\bottomrule\noalign{}
\endlastfoot
16 & \(A^{\dagger}\), the conjugate transpose & the \textbf{mirror image
of the box} \(A\), labelled \(A^{\dagger}\): input and output change
places, and the mirror conjugates --- \emph{new element} \\
17 & \((BA)^{\dagger}=A^{\dagger}B^{\dagger}\);
\(\;\bigl(A\ket{\psi}\bigr)^{\dagger}=\bra{\psi}A^{\dagger}\) & the
mirror image of a row: every piece mirrored, the order reversed \\
18 & \(\bra{i}A\ket{j}=A_{ij}\) & state \(j\), box \(A\), effect \(i\):
a closed picture, a number \\
19 & \(U^{\dagger}U=I=UU^{\dagger}\) (unitary) & \textbf{\emph{a box
meets its mirror}}: in either order, a bare wire \\
20 & \(A^{\dagger}=A\) (here \(X\), \(Y\), \(Z\), \(H\)) & a box that is
its own mirror image \\
21 & \(H\ket{0}=\ket{+}\), \(H\ket{1}=\ket{-}\); \(\;\bra{0}H=\bra{+}\),
\(\bra{1}H=\bra{-}\) & state \(0\), then \(H\) \(=\) state \(+\);
\(\;H\), then effect \(0\) \(=\) effect \(+\) (likewise \(1\) and
\(-\)) \\
22 & \(HZH=X\) & the row \(H\), \(Z\), \(H\) \(=\) the box \(X\) \\
\end{longtable}

\endgroup

All seven rows are exact. There is no row for the sphere: it is a way of
reading three closed pictures of the kind in row 18, and not a picture
element.

\begin{center}\rule{0.5\linewidth}{0.5pt}\end{center}

\hypertarget{the-theorem-twice-2}{%
\section{§3.8 The Theorem, Twice}\label{the-theorem-twice-2}}

We have met the theorem through its consequence: the arrow's first and
third coordinates changed places under \(H\) (§3.6). Here is the
statement.

\begin{quote}
\textbf{Theorem 3.4 (The Sandwich).} Let
\(H=\tfrac1{\sqrt2}\begin{pmatrix}1&1\\1&-1\end{pmatrix}\),
\(Z=\begin{pmatrix}1&0\\0&-1\end{pmatrix}\) and
\(X=\begin{pmatrix}0&1\\1&0\end{pmatrix}\). Then

\[HZH \;=\; X,\]

exactly, scalar included. In pictures: the row of boxes \(H\), \(Z\),
\(H\) on one wire is the box \(X\).
\end{quote}

\hypertarget{diracs-blackboard-2}{%
\subsection{Dirac's Blackboard}\label{diracs-blackboard-2}}

Two matrix products. First \(ZH\); multiplying on the left by \(Z\)
changes the sign of the second row:

\[ZH \;=\; \begin{pmatrix}1&0\\0&-1\end{pmatrix}\cdot\tfrac1{\sqrt2}\begin{pmatrix}1&1\\1&-1\end{pmatrix} \;=\; \tfrac1{\sqrt2}\begin{pmatrix}1&1\\-1&1\end{pmatrix}.\]

Then \(H\) times that, along the rows and down the columns, with
\(\tfrac1{\sqrt2}\cdot\tfrac1{\sqrt2}=\tfrac12\) in front:

\[H(ZH) \;=\; \tfrac12\begin{pmatrix}1&1\\1&-1\end{pmatrix}\begin{pmatrix}1&1\\-1&1\end{pmatrix}\]

\[=\; \tfrac12\begin{pmatrix}1\cdot1+1\cdot(-1)&1\cdot1+1\cdot1\\1\cdot1+(-1)(-1)&1\cdot1+(-1)\cdot1\end{pmatrix} \;=\; \tfrac12\begin{pmatrix}0&2\\2&0\end{pmatrix} \;=\; X .\]

Two products, eight multiplications in the second, and the two factors
\(\tfrac1{\sqrt2}\) meet to cancel the \(2\). \(\;\blacksquare\)

\hypertarget{penroses-sketchbook-2}{%
\subsection{Penrose's Sketchbook}\label{penroses-sketchbook-2}}

The sketchbook owns no rewrite that turns three boxes into one. What it
owns, and where each thing comes from, is this. \emph{H on 0}, \emph{H
on 1}, \emph{Z on +} and \emph{Z on \(-\)} (§3.4) are \textbf{lent by
the blackboard}: each is a column computed there and recorded as a
picture. \emph{A box meets its mirror} (row 19) is this chapter's own
rule, and with row 20 (\(H\) is its own mirror image) it says that two
\(H\) boxes in a row are a bare wire. From Chapter 1 we have \emph{NOT
on 0} and \emph{NOT on 1} (§1.1), which with NOT labelled \(X\) read
``state \(0\), then \(X\), is the state \(1\)'' and ``state \(1\), then
\(X\), is the state \(0\)''; and the principle \textbf{\emph{two states
suffice}} (§1.7), also lent by the blackboard: two rows of boxes that
agree after the state \(0\) and after the state \(1\) are equal.

\textbf{Step 0 --- \emph{H on +}.} One more small equation is needed,
and the chapter's own rule supplies it. Read \emph{H on 0} backwards to
replace the state \(+\); then two \(H\) boxes are a bare wire.

\[\begin{tikzpicture}[sd, baseline=-0.55ex]
  \draw[wire] (0,0) -- (5,0);
  \node[sdstate, text height=1.5ex, text depth=0.5ex] at (0,0) {$+$};
  \node[sdbox, text height=1.5ex, text depth=0.5ex] at (2.8,0) {$H$};
\end{tikzpicture}
\;\overset{\raisebox{0pt}[\height][1.2ex]{\scriptsize H on 0}}{=}\;
\begin{tikzpicture}[sd, baseline=-0.55ex]
  \draw[wire] (0,0) -- (7.4,0);
  \node[sdstate, text height=1.5ex, text depth=0.5ex] at (0,0) {$0$};
  \node[sdbox, text height=1.5ex, text depth=0.5ex] at (2.8,0) {$H$};
  \node[sdbox, text height=1.5ex, text depth=0.5ex] at (5.2,0) {$H$};
\end{tikzpicture}
\;\overset{\raisebox{0pt}[\height][1.2ex]{\scriptsize a box meets its mirror}}{=}\;
\begin{tikzpicture}[sd, baseline=-0.55ex]
  \draw[wire] (0,0) -- (2.6,0);
  \node[sdstate, text height=1.5ex, text depth=0.5ex] at (0,0) {$0$};
\end{tikzpicture}\]

In the same way, from \emph{H on 1}: the state \(-\), then \(H\), is the
state \(1\). Call these \textbf{\emph{H on +}} and \textbf{\emph{H on
\(-\)}}.

\textbf{Step 1 --- follow the state \(0\).} Place the state \(0\) before
the row and replace, three times, a state and the box after it by a
single state.

\[\begin{tikzpicture}[sd, baseline=-0.55ex]
  \draw[wire] (0,0) -- (9.8,0);
  \node[sdstate, text height=1.5ex, text depth=0.5ex] at (0,0) {$0$};
  \node[sdbox, text height=1.5ex, text depth=0.5ex] at (2.8,0) {$H$};
  \node[sdbox, text height=1.5ex, text depth=0.5ex] at (5.2,0) {$Z$};
  \node[sdbox, text height=1.5ex, text depth=0.5ex] at (7.6,0) {$H$};
\end{tikzpicture}
\;\overset{\raisebox{0pt}[\height][1.2ex]{\scriptsize H on 0}}{=}\;
\begin{tikzpicture}[sd, baseline=-0.55ex]
  \draw[wire] (0,0) -- (7.4,0);
  \node[sdstate, text height=1.5ex, text depth=0.5ex] at (0,0) {$+$};
  \node[sdbox, text height=1.5ex, text depth=0.5ex] at (2.8,0) {$Z$};
  \node[sdbox, text height=1.5ex, text depth=0.5ex] at (5.2,0) {$H$};
\end{tikzpicture}\]

\[\overset{\raisebox{0pt}[\height][1.2ex]{\scriptsize Z on +}}{=}\;
\begin{tikzpicture}[sd, baseline=-0.55ex]
  \draw[wire] (0,0) -- (5,0);
  \node[sdstate, text height=1.5ex, text depth=0.5ex] at (0,0) {$-$};
  \node[sdbox, text height=1.5ex, text depth=0.5ex] at (2.8,0) {$H$};
\end{tikzpicture}
\;\overset{\raisebox{0pt}[\height][1.2ex]{\scriptsize H on $-$}}{=}\;
\begin{tikzpicture}[sd, baseline=-0.55ex]
  \draw[wire] (0,0) -- (2.6,0);
  \node[sdstate, text height=1.5ex, text depth=0.5ex] at (0,0) {$1$};
\end{tikzpicture}
\;\overset{\raisebox{0pt}[\height][1.2ex]{\scriptsize NOT on 0}}{=}\;
\begin{tikzpicture}[sd, baseline=-0.55ex]
  \draw[wire] (0,0) -- (5,0);
  \node[sdstate, text height=1.5ex, text depth=0.5ex] at (0,0) {$0$};
  \node[sdbox, text height=1.5ex, text depth=0.5ex] at (2.8,0) {$X$};
\end{tikzpicture}\]

\textbf{Step 2 --- follow the state \(1\).} The same, with the signs the
other way round.

\[\begin{tikzpicture}[sd, baseline=-0.55ex]
  \draw[wire] (0,0) -- (9.8,0);
  \node[sdstate, text height=1.5ex, text depth=0.5ex] at (0,0) {$1$};
  \node[sdbox, text height=1.5ex, text depth=0.5ex] at (2.8,0) {$H$};
  \node[sdbox, text height=1.5ex, text depth=0.5ex] at (5.2,0) {$Z$};
  \node[sdbox, text height=1.5ex, text depth=0.5ex] at (7.6,0) {$H$};
\end{tikzpicture}
\;\overset{\raisebox{0pt}[\height][1.2ex]{\scriptsize H on 1}}{=}\;
\begin{tikzpicture}[sd, baseline=-0.55ex]
  \draw[wire] (0,0) -- (7.4,0);
  \node[sdstate, text height=1.5ex, text depth=0.5ex] at (0,0) {$-$};
  \node[sdbox, text height=1.5ex, text depth=0.5ex] at (2.8,0) {$Z$};
  \node[sdbox, text height=1.5ex, text depth=0.5ex] at (5.2,0) {$H$};
\end{tikzpicture}\]

\[\overset{\raisebox{0pt}[\height][1.2ex]{\scriptsize Z on $-$}}{=}\;
\begin{tikzpicture}[sd, baseline=-0.55ex]
  \draw[wire] (0,0) -- (5,0);
  \node[sdstate, text height=1.5ex, text depth=0.5ex] at (0,0) {$+$};
  \node[sdbox, text height=1.5ex, text depth=0.5ex] at (2.8,0) {$H$};
\end{tikzpicture}
\;\overset{\raisebox{0pt}[\height][1.2ex]{\scriptsize H on +}}{=}\;
\begin{tikzpicture}[sd, baseline=-0.55ex]
  \draw[wire] (0,0) -- (2.6,0);
  \node[sdstate, text height=1.5ex, text depth=0.5ex] at (0,0) {$0$};
\end{tikzpicture}
\;\overset{\raisebox{0pt}[\height][1.2ex]{\scriptsize NOT on 1}}{=}\;
\begin{tikzpicture}[sd, baseline=-0.55ex]
  \draw[wire] (0,0) -- (5,0);
  \node[sdstate, text height=1.5ex, text depth=0.5ex] at (0,0) {$1$};
  \node[sdbox, text height=1.5ex, text depth=0.5ex] at (2.8,0) {$X$};
\end{tikzpicture}\]

\textbf{Step 3 --- \emph{two states suffice}.} The row \(H\), \(Z\),
\(H\) and the box \(X\) give equal pictures after the state \(0\) (Step
1) and after the state \(1\) (Step 2). So they are equal:

\[\begin{tikzpicture}[sd, baseline=-0.55ex]
  \draw[wire] (0,0) -- (9,0);
  \node[sdbox, text height=1.5ex, text depth=0.5ex] at (2,0) {$H$};
  \node[sdbox, text height=1.5ex, text depth=0.5ex] at (4.5,0) {$Z$};
  \node[sdbox, text height=1.5ex, text depth=0.5ex] at (7,0) {$H$};
\end{tikzpicture}
\;\;\overset{\raisebox{0pt}[\height][1.2ex]{\scriptsize two states suffice}}{=}\;\;
\begin{tikzpicture}[sd, baseline=-0.55ex]
  \draw[wire] (0,0) -- (5,0);
  \node[sdbox, text height=1.5ex, text depth=0.5ex] at (2.5,0) {$X$};
\end{tikzpicture}
\qquad\blacksquare\]

Every equals sign above is exact. The states passed through were
\(0,+,-,1\) in Step 1 and \(1,-,+,0\) in Step 2, the same four in the
opposite order. §3.9 draws that.

\hypertarget{the-verdict-2}{%
\subsection{The verdict}\label{the-verdict-2}}

The algebra wins, and it is not close. The blackboard needs two products
of \(2\times2\) matrices. The sketchbook needs eight replacements and a
closing principle, and of the facts it uses only one, \emph{a box meets
its mirror}, is its own; the columns of \(H\), the action of \(Z\) on
\(\pm\), and \emph{two states suffice} are all on loan from the
blackboard. At this stage the sketchbook has no rule that acts on \(H\),
\(Z\), \(H\) as boxes. A later chapter brings a notation in which this
equation is a single move; we do not have it yet, and we shall not
pretend to.

What the longer proof gives is a reason, which the matrix product does
not. The first \(H\) turns the standard question into the turned one. In
the turned basis, \(Z\) is a NOT. The second \(H\) turns back. So
``NOT'' is not the name of a matrix but of a job, exchanging the two
answers of a question, and \(Z\) and \(X\) are the same job done for two
different questions. The blackboard proves the theorem and shows none of
this.

\begin{quote}
\textbf{Feynman's Notebook.}

\emph{The adjoint as a mirror image.} That the dagger of a picture is
its reflection is not this book's conceit. Selinger's survey of
graphical languages (§7) sets out the general notion, a \emph{dagger
category}, in which every map has an adjoint going the other way, with
exactly the three laws of our §3.2: the adjoint of doing nothing is
doing nothing, the adjoint of a composite is the composite of the
adjoints in the reverse order, and the adjoint of an adjoint is the
original. He states the graphical rule in one line: the adjoint of any
diagram is its mirror image. He also notes the difficulty we met, that a
box's reflection can be told from the box only because the boxes carry a
mark in one corner. The same section defines a unitary map as an
invertible one whose inverse is its adjoint. Abramsky and Coecke (2004,
§7.1) have the same definition and prove from it, in their abstract
setting, that unitary maps preserve the inner product: our Theorem 3.2,
without a matrix in sight.
\end{quote}

\begin{center}\rule{0.5\linewidth}{0.5pt}\end{center}

\hypertarget{the-turn-in-three-steps}{%
\section{§3.9 The Turn in Three Steps}\label{the-turn-in-three-steps}}

Here is the chapter's master diagram. We call it \textbf{The Turn in
Three Steps}. It is Step 1 of the sketchbook proof with every
intermediate state \emph{asked whether it is what we said it was}. On
one row, from the left: the state \(0\), the box \(H\), the effect
\(+\); bare paper; the state \(+\), the box \(Z\), the effect \(-\);
bare paper; the state \(-\), the box \(H\), the effect \(1\). Each wire
is closed at both ends and carries one box, and nothing joins an effect
to the state that follows it.

\[\begin{tikzpicture}[sd, baseline=-0.55ex]
  \draw[wire] (0,0) -- (4,0);
  \draw[wire] (6.5,0) -- (10.5,0);
  \draw[wire] (13,0) -- (17,0);
  \node[sdstate,  text height=1.5ex, text depth=0.5ex] at (0,0)    {$0$};
  \node[sdbox] at (2,0) {$H$};
  \node[sdeffect, text height=1.5ex, text depth=0.5ex] at (4,0)    {$+$};
  \node[sdstate,  text height=1.5ex, text depth=0.5ex] at (6.5,0)  {$+$};
  \node[sdbox] at (8.5,0) {$Z$};
  \node[sdeffect, text height=1.5ex, text depth=0.5ex] at (10.5,0) {$-$};
  \node[sdstate,  text height=1.5ex, text depth=0.5ex] at (13,0)   {$-$};
  \node[sdbox] at (15,0) {$H$};
  \node[sdeffect, text height=1.5ex, text depth=0.5ex] at (17,0)   {$1$};
\end{tikzpicture} \;\;=\;\; 1\cdot 1\cdot 1 \;=\; 1\]

The equation is exact. Closed pictures in a row multiply (§2.3), and the
three factors are

\[\bra{+}H\ket{0}=\braket{+|+}=1, \qquad \bra{-}Z\ket{+}=\braket{-|-}=1, \qquad \bra{1}H\ket{-}=\braket{1|1}=1,\]

by \emph{H on 0}, \emph{Z on +} and \emph{H on \(-\)} followed each time
by Theorem 2.1. Each \(1\) is a Born certainty: the state \(0\), after
\(H\), is \(+\) for sure; \(+\), after \(Z\), is \(-\) for sure; \(-\),
after \(H\), is \(1\) for sure. The number \(1\) is also the entry
\(\bra{1}X\ket{0}\) of the box \(X\), as it must be by \emph{The
Sandwich}. But we draw no equals sign between this row and the closed
picture with all three boxes on a single wire between the state \(0\)
and the effect \(1\). That is a different picture, and at this stage
only the blackboard can say why it has the same value: between its boxes
it carries whatever state there is, and here we happen to have named
each one. As with \emph{The Echo}, we prove nothing again, and read the
picture three times.

\textbf{First reading: as drawn.} \(0\) to \(+\) to \(-\) to \(1\): the
first column of \(X\), in three certainties. Value \(1\).

\textbf{Second reading: in the mirror.} Reflect the whole row by the
rule of §3.2: every piece mirrored, the order reversed. States become
effects and effects states; \(H\) and \(Z\) are their own mirror images;
the last wire becomes the first.

\[\begin{tikzpicture}[sd, baseline=-0.55ex]
  \draw[wire] (0,0) -- (4,0);
  \draw[wire] (6.5,0) -- (10.5,0);
  \draw[wire] (13,0) -- (17,0);
  \node[sdstate,  text height=1.5ex, text depth=0.5ex] at (0,0)    {$1$};
  \node[sdbox] at (2,0) {$H$};
  \node[sdeffect, text height=1.5ex, text depth=0.5ex] at (4,0)    {$-$};
  \node[sdstate,  text height=1.5ex, text depth=0.5ex] at (6.5,0)  {$-$};
  \node[sdbox] at (8.5,0) {$Z$};
  \node[sdeffect, text height=1.5ex, text depth=0.5ex] at (10.5,0) {$+$};
  \node[sdstate,  text height=1.5ex, text depth=0.5ex] at (13,0)   {$+$};
  \node[sdbox] at (15,0) {$H$};
  \node[sdeffect, text height=1.5ex, text depth=0.5ex] at (17,0)   {$0$};
\end{tikzpicture} \;\;=\;\; \bar1\cdot\bar1\cdot\bar1 \;=\; 1\]

Its value is the conjugate, which for the number \(1\) is \(1\) again.
And look at what it says: \(1\) to \(-\) to \(+\) to \(0\). That is Step
2 of the sketchbook proof, the \emph{other} column of \(X\). The picture
and its reflection together are the whole proof of \emph{The Sandwich}.
The theorem is a row that reads the same from either end, and so is its
proof.

\textbf{Third reading: any three boxes that keep the promise.} Replace
\(0\) by any state \(\psi\), and \(H\), \(Z\), \(H\) by any three
unitary boxes \(A\), \(B\), \(C\); label each effect with the state that
ought to arrive there.

\[\begin{tikzpicture}[sd, baseline=-0.55ex]
  \draw[wire] (0,0) -- (4.8,0);
  \draw[wire] (8.4,0) -- (13.9,0);
  \draw[wire] (18.1,0) -- (24.4,0);
  \node[sdstate,  text height=1.5ex, text depth=0.5ex] at (0,0)    {$\psi$};
  \node[sdbox, text height=1.5ex, text depth=0.5ex] at (2.2,0) {$A$};
  \node[sdeffect, inner xsep=0.5em, text height=1.5ex, text depth=0.5ex] at (4.8,0)  {$A\psi$};
  \node[sdstate,  inner xsep=0.5em, text height=1.5ex, text depth=0.5ex] at (8.4,0)  {$A\psi$};
  \node[sdbox, text height=1.5ex, text depth=0.5ex] at (11,0) {$B$};
  \node[sdeffect, inner xsep=0.5em, text height=1.5ex, text depth=0.5ex] at (13.9,0) {$BA\psi$};
  \node[sdstate,  inner xsep=0.5em, text height=1.5ex, text depth=0.5ex] at (18.1,0) {$BA\psi$};
  \node[sdbox, text height=1.5ex, text depth=0.5ex] at (21,0) {$C$};
  \node[sdeffect, inner xsep=0.5em, text height=1.5ex, text depth=0.5ex] at (24.4,0) {$CBA\psi$};
\end{tikzpicture}\]

\[=\;\; 1\cdot1\cdot1 \;\;=\;\; 1\]

Here a rounded box labelled \(A\psi\) stands for the state
\(A\ket{\psi}\), or for its mirror image when the wire is on its left.
Why is the first factor \(1\)? The effect labelled \(A\psi\) is the
mirror image of ``state \(\psi\), then \(A\)'', which by \emph{The
Mirror of a Row} is ``\(A^{\dagger}\), then the effect \(\psi\)''. So
the first wire is the closed picture at the head of §3.3,

\[\bra{\psi}A^{\dagger}A\ket{\psi} \;=\; 1,\]

by \emph{a box meets its mirror} and Theorem 2.1. The other two wires
are the same argument with \(A\ket{\psi}\) and \(BA\ket{\psi}\) in the
place of \(\ket{\psi}\); they are states because \(A\) and \(B\) keep
the promise. To make the reading an honest test of the dagger, take
boxes that are not their own reflections: \(A=S\), \(B=U\) (our test
piece, neither real nor symmetric), \(C=T\), and \(\psi\) any state at
all. Every factor is still \(1\). A cheaper reflection fails at the
first wire: with the transpose in the place of the mirror, and the arrow
for \(\psi\), the first factor would be
\(\bra{a}S^{T}S\ket{a}=\bra{a}Z\ket{a}=-\tfrac7{25}\).

Now take a box that does not keep the promise. With \(A=S_0\) and the
arrow, the first wire is ``state \(a\), box \(S_0\), effect \(S_0a\)'',
and its value is \(\tfrac{49}{25}\): not a certainty, and not a chance
of any kind. This third reading is the answer to §2.10. \emph{The boxes
that keep every wire of this picture equal to \(1\) are the unitary
ones.} See \texttt{Code/Ch3/Sandwich.py}, which evaluates all three
readings.

\begin{quote}
\textbf{Penrose's Sketchbook.}

\emph{What the gaps are.} The bare paper between an effect and the next
state is part of the picture. Left of a gap, a wire has ended in a
question; right of it, a new wire begins in a state. That the new state
has the same label as the old question is our choice, made because the
answer was certain. It is not something the picture enforces. Nothing
flows across a gap.

This is why the row is not the sandwich itself. Join the wires into one
by erasing the inner effects and states, and you have the closed picture
``state \(0\), then \(H\), \(Z\), \(H\), then effect \(1\)'', in which
the boxes pass along whatever is on the wire, certain or not. For these
labels, and for the third reading, the two pictures have the same
number. For boxes in general, with intermediate questions whose answers
are \emph{not} certain, cutting a wire and asking changes the number.
What cutting and asking does to a wire is a later chapter's subject. We
have been careful to cut only where the answer was already sure.
\end{quote}

\begin{center}\rule{0.5\linewidth}{0.5pt}\end{center}

\hypertarget{a-second-wire}{%
\section{§3.10 A Second Wire}\label{a-second-wire}}

Let us gather what we have. A box on a closed wire is a number, and four
such numbers are the box (§3.1). The mirror image of a box is forced by
what Chapter 2 said about mirrors: the conjugate transpose, and neither
of the two cheaper things (§3.2). The mirror image of a row reverses the
order \emph{and} reflects each piece (§3.2). The boxes that send every
state to a state are exactly those that, meeting their reflection, leave
a bare wire (§3.3). Six such boxes have names; four are their own
reflections, and for that very reason can test nothing about reflections
(§3.4). A state, stripped of the phase no question sees, is a point on a
sphere (§3.5). Each of our boxes turns that sphere; \(H\) exchanges two
of its axes; and, granting one cited theorem, every box on a qubit's
wire is three turns (§3.6). The row \(H\), \(Z\), \(H\) is the box
\(X\); the blackboard proves it in two products, and the sketchbook,
with borrowed facts, in three steps and their reflection (§3.8, §3.9).

All of this comes to one sentence, and it answers the question with
which Chapter 2 closed. \textbf{A box may sit on a qubit's wire exactly
when its mirror image is its undoing.} Every such box can be run
backwards by whoever stands on the other side of the glass. SET-0 is not
such a box. It sends many states to one, and its reflection would have
to send that one state back to many. We said in Chapter 1 that a qubit
is a thing not looked at, and we can now add that everything done to it
on the wire is a thing that could be undone. Whatever looking is --- and
we have still not said --- it is not a box of this kind, for no one has
ever unlooked.

Now notice what none of our pictures has. So far every diagram in this
book has lain along a single line. When we needed more than one wire, in
§3.9, we set them end to end, with bare paper between, and their numbers
multiplied; nothing passed from one to the next. But a coin is seldom
alone. Put a second coin on the table beside the first. Is its wire
drawn beside the first wire --- above it, below it? What is the column
for two coins: two entries and two more, or something else? If a box
sits on one wire and nothing on the other, what is the matrix of the
pair? And if two boxes sit on two wires, does it matter which we think
of as acting first?

One wire, we have learned, is a sphere, and all that can be done to it
is to turn it. We leave the sphere turning. The next chapter lays a
second wire beside the first.

\hypertarget{yoked-inventions}{%
\chapter{Yoked Inventions}\label{yoked-inventions}}

\begin{center}\rule{0.5\linewidth}{0.5pt}\end{center}

\emph{糸と糸}\\
\emph{並べて張るや}\\
\emph{春隣}

\emph{ito to ito}\\
\emph{narabete haru ya}\\
\emph{haru tonari}

\emph{A string, and a string,}\\
\emph{stretched taut and laid side by side ---}\\
\emph{spring, one door away.}

\begin{center}\rule{0.5\linewidth}{0.5pt}\end{center}

\emph{{[}The snow is going. It has drawn back from the foot of the
Tortoise's wall and from the sunny side of every post, and what is left
of it in the yard is grey and full of holes; the path to the gate is wet
earth. Along the porch eave hangs a row of icicles, and at the end the
sun reached first one of them has begun to drip --- a drop, a long wait,
another drop --- as though the thaw had not yet made up its mind. The
rest hang dry. Out in the yard the old pear tree is bare; across the
yard the woodshed roof, and the roof of the lean-to under it, still
carry their snow.{]}}

\emph{{[}Achilles comes up the porch steps; the Tortoise is unfolding a
wooden frame onto the table.{]}}

\textbf{Achilles:} The whole hillside is running water this morning ---
your snow has gone to soup, and I came the last stretch more sliding
than walking.

\textbf{Tortoise:} It does that, about now; by noon the path will be
firmer than it was at dawn. Sit and warm your hands, the kettle is hot.
Though I did hope you might come early enough to be useful.

\textbf{Achilles:} Useful is never a restful word from you. What is this
you are unfolding?

\textbf{Tortoise:} A frame for stretching strings across --- two
uprights, a row of hooks down each, and a string with a loop at either
end to go between a facing pair. It lets us set one thing beside another
and follow them together, in place of one at a time. And here beside it
a stiff card, to stand a coin behind, so that we speak of its honest
shares and not of the face it happens to show.

\textbf{Achilles:} One string and one coin I can manage before the tea
has cooled. Hook it up, and let me begin on ground I know.

\begin{center}\rule{0.5\linewidth}{0.5pt}\end{center}

\hypertarget{first-challenge-a-string-below}{%
\section{\texorpdfstring{First Challenge --- \emph{A String
Below}}{First Challenge --- A String Below}}\label{first-challenge-a-string-below}}

\emph{{[}The porch is open to the yard along its front and has the house
wall at its back. The door has a window in it, a bar across and a bar
down, all clear glass but the upper pane on the hinge side, which is old
green bottle-glass; where Achilles stamped his boots before it there is
a small puddle on the boards. On the wall beside the door hangs Escher's
\textup{Three Spheres II} in a narrow frame: three balls in a row on a
table top, the left-hand one letting the light straight through it, the
right-hand one giving none of it back. The shadow of the corner post
lies across it just now. Against the end rail leans a glazed light off
the cold frame, waiting for warmer nights, and on the rail itself stands
a seed tray parted into cells by slats. The table is at the front, in
the sun.{]}}

\emph{{[}Late winter, the first thaw; a clear morning; the Tortoise's
porch. On the table stands a small folding wooden frame out of her
shell: two uprights with a column of hooks down each, and a silk string,
a loop at either end, stretched left to right between the top hooks.
Achilles walks a coin along the table beside the string, from the left
upright to the right.{]}}

\emph{{[}The frame does not stand bolt upright. A strut behind lets it
lean far back, as a reading-slope leans, so that a coin on the table
lies under its string. Its left upright is the one towards the porch
steps; Achilles sits on the side of the top hooks and the Tortoise on
the side of the bottom ones, and there are more hooks down each upright
than the morning looks likely to want. The silk is the finest thing on
the porch: drawn tight, it takes the sun along its whole length and
shows as hardly more than a line of light.{]}}

\textbf{Achilles:} Heads at the left, heads at the right. A bare string
says nothing; I remember that much. One string, one coin, and a list of
two lines for it --- a line for heads, a line for tails. I could do this
in my sleep.

\emph{{[}Achilles settles back, both hands round the hot cup. The kettle
ticks as it cools on its trivet.{]}}

\textbf{Tortoise:} \emph{{[}She loops another string over the next pair
of hooks down, sets up a card on edge as a screen, takes Achilles's coin
behind it, and there tosses it, and two pennies after it.{]}} A string
below: hemp. A coin on it, behind the card: tails, if both my pennies
fell tails, and heads otherwise. And what may honestly be said of that
coin: heads, three times in four.

\textbf{Achilles:} \emph{{[}She has laid the tossed coin behind the card
at the silk string's left end.{]}} A string above: silk. A coin on it,
behind the card: as it fell. And what may honestly be said of that coin:
heads, one time in two.

\emph{{[}The hemp is the colour of a doormat, and hairy where the silk
is smooth; the sun makes nothing of it. It runs a hand's breadth below
the silk from hook to hook and touches it nowhere.{]}}

\textbf{Tortoise:} Write down what you know of the pair.

\textbf{Achilles:} \emph{{[}writing{]}} Two lines for the upper coin: a
half, a half. Two lines for the lower: three-quarters, a quarter. I keep
two little lists, one under the other, as the strings are.

\textbf{Tortoise:} Then read me off the share of heads above and tails
below, together.

\textbf{Achilles:} Together. Well --- a half from the upper list, a
quarter from the lower: three-quarters?

\textbf{Tortoise:} And heads above with heads below?

\textbf{Achilles:} A half and three-quarters: one and a quarter. More
than the whole morning. No.~My two lists have no line for it, and adding
won't make one.

\textbf{Tortoise:} Then we want something with a place for every such
pairing at once. This plain sheet will serve: folded across and across,
it falls into panels, and we can give one panel to each way the two
coins might come up.

\textbf{Tortoise:} \emph{{[}She takes the blank sheet --- one whole face
--- and folds it once across the middle; the sheet is no smaller folded,
but the crease has scored it into an upper half and a lower half.{]}}
For the silk: above the crease it shows heads, below it tails.
\emph{{[}She folds it again, this time end to end, and opens it flat:
the outline is unchanged, yet the two creases have quartered the single
face into four panels.{]}} For the hemp: left of the crease heads, right
of it tails. Four panels. Fill them.

\textbf{Achilles:} Upper right is the one you asked for. In half of all
mornings the silk coin shows heads; in a quarter of \emph{those} the
hemp shows tails. A half of a quarter: one-eighth.

\textbf{Tortoise:} And its neighbours?

\textbf{Achilles:} \emph{{[}writing in each panel{]}} Upper left, a half
of three-quarters: three-eighths. Below them the same again:
three-eighths, one-eighth. Eight eighths. They make one whole, which my
two lists never did!

\textbf{Tortoise:} What would you call it?

\textbf{Achilles:} A patchwork. Every panel is its row's share times its
column's share, and the two little lists sit along the edges like a
border.

\emph{{[}On the rail the seed tray stands in the sun. The Tortoise has
stuck a label in the earth at the head of each column, for what was sown
there, and one at the end of each row, for the week it was sown in. The
labels go round two edges and are soon read; the cells fill the tray.
Nothing is up in any of them yet.{]}}

\textbf{Tortoise:} And had I given you reaches, not shares?

\textbf{Achilles:} Reach times reach in every panel, and the squaring
after. It comes to the same patchwork.

\emph{{[}The Tortoise takes down the card. The silk coin shows tails;
the hemp coin, heads. Achilles puts a finger on the lower left
panel.{]}}

\textbf{Achilles:} So what you're saying is, the pair wants a line for
every panel, and the panels go by strings. Two and two --- four. So a
third string makes six.

\textbf{Tortoise:} Fold it.

\begin{center}\rule{0.5\linewidth}{0.5pt}\end{center}

\hypertarget{second-challenge-sixteen-panels}{%
\section{\texorpdfstring{Second Challenge --- \emph{Sixteen
Panels}}{Second Challenge --- Sixteen Panels}}\label{second-challenge-sixteen-panels}}

\textbf{Achilles:} \emph{{[}The Tortoise hands him a string out of her
shell; he loops it over the next hooks down, as she did.{]}} A string
below: gut. And a fold for it. \emph{{[}He takes a fresh sheet and folds
it, the last crease turning the whole folded square back over on
itself.{]}} Once across for the silk. Once the other way for the hemp.
Once more for the gut. Six panels. \emph{{[}He opens it flat and runs a
finger down the panels, counting.{]}} Two, four, six\ldots{} eight.

\textbf{Tortoise:} Count them again, by all means.

\textbf{Achilles:} Eight. \ldots Oh. The new fold didn't add two panels
at the edge. It went through every panel that was already there.

\emph{{[}Out in the yard the pear tree is bare, and with the snow off it
one can see how the Tortoise has pruned it all these years. The trunk
parts in two at the height of the wall; each limb parts in two above
that, and each branch again, and again. Nobody has ever thought of
counting the twigs.{]}}

\textbf{Tortoise:} \emph{{[}She loops another string over the next free
hooks; it settles straight beneath the others, and not one of the
strings already hooked so much as stirs.{]}} A string below: horsehair.
And a fold for it.

\textbf{Achilles:} \emph{{[}folding the sheet over on itself again, and
once more; it is growing thick, and each fold splits every panel there
is anew{]}} A fourth fold. \emph{{[}He opens it out and counts along the
rows.{]}} Sixteen. And a fifth would be thirty-two, and I can no longer
crease the paper. Not two more for each string, then: twice as many for
each string.

\emph{{[}The gut is yellow and hard, a little clear at the edges like a
sliver of horn; it gave a small dry note when Achilles's thumbnail
caught it in the hooking. The horsehair under it is dark and dull and
gives no note at all. Below the horsehair the hooks go on down the
uprights, empty.{]}}

\textbf{Tortoise:} And if the coins are separate, as ours were --- each
tossed by itself behind the card?

\textbf{Achilles:} Then a little list for each would do, along the
borders. Two lines for the silk coin, two for the hemp, two for the gut,
two for the horsehair: eight lines, and I multiply across to fill the
panels.

\textbf{Tortoise:} Eight lines in all. The sheet has sixteen panels. And
not every sheet is a patchwork.

\textbf{Achilles:} What else could a sheet be? You folded it; I
multiplied.

\emph{{[}The Tortoise does not answer. She puts the blank sheet by,
keeps the sheet of the first folding at hand, and takes from her shell a
small box of wooden clips. Each has a rule written on it.{]}}

\textbf{Tortoise:} You have met these --- the little clips, each with
its rule. This one says turn the coin over; that one says set it heads,
whatever it showed. Pinch a clip onto a string, and it does its rule to
that coin as it passes, and to no other coin at all.

\emph{{[}The sun has come along the eave. Where one icicle dripped there
are now a dozen at it, each in its own time --- quick ones, slow ones,
and one that only gathers and gathers --- and together they make a
patter that nobody could beat time to. Under them the snow along the
foot of the porch is beginning to be pitted.{]}}

\textbf{Achilles:} And a clip? A clip knows nothing of panels. It sits
on one string.

\begin{center}\rule{0.5\linewidth}{0.5pt}\end{center}

\hypertarget{third-challenge-one-clip-nothing-on-the-other}{%
\section{\texorpdfstring{Third Challenge --- \emph{One Clip, Nothing on
the
Other}}{Third Challenge --- One Clip, Nothing on the Other}}\label{third-challenge-one-clip-nothing-on-the-other}}

\textbf{Tortoise:} \emph{{[}She takes one clip and pinches it onto the
silk string, halfway along; the hemp string below it is left bare.{]}} A
clip on the silk: ``turn it over''. Nothing on the hemp. What becomes of
the sheet?

\textbf{Achilles:} The hemp coin is untouched, so half the sheet is
untouched. And better: the silk coin was half-and-half, and turned over
it is half-and-half still. Three-eighths, one-eighth above;
three-eighths, one-eighth below. The whole sheet is untouched!

\textbf{Tortoise:} The numbers are alike. Are the panels? Lay the coins
where you can see them, and walk them. \emph{{[}She gives him a blue
pencil.{]}} Mark where each panel goes.

\textbf{Achilles:} Heads above, heads below. \emph{{[}He sets a coin on
each string and walks the pair along together, level, from the left
upright; at the clip the silk coin is flipped and the hemp coin below
passes it untouched.{]}} Tails above, heads below. Upper left goes to
lower left. \emph{{[}He draws a blue arrow on the sheet.{]}} Heads
above, tails below: arrives tails above, tails below. Upper right goes
to lower right.

\textbf{Tortoise:} And from the lower panels?

\textbf{Achilles:} Up. Each to the panel above it. \emph{{[}Two more
walks, two more arrows; and every arrow he has drawn runs straight up or
straight down, not one of them sideways.{]}} Every panel has moved! The
halves made by the first fold have changed places, upper for lower. The
shares looked the same only because the two halves were twins.

\emph{{[}By the door, where Achilles stamped his boots, the puddle has
the door's window in it, standing on its head. The green pane, which is
at the top on the hinge side, lies at the bottom on the hinge side, and
every clear pane has gone over the bar across with it. A drop blown in
from the eave shakes the whole window to pieces, and it comes back the
same.{]}}

\textbf{Tortoise:} Your turn to say it.

\textbf{Achilles:} A clip on the hemp: ``turn it over''. Nothing on the
silk.

\emph{{[}The Tortoise lifts the clip off the silk and pinches it onto
the hemp instead; now it is the lower coin the rule will turn. She gives
him a red pencil, and he walks the coins from each panel in turn.{]}}

\textbf{Achilles:} Heads above, heads below: arrives heads above, tails
below. Upper left goes to upper right, and upper right comes back to
upper left. Below the crease, the same. \emph{{[}Red arrows.{]}} Every
panel moves again --- but this time nothing crosses the first fold. The
panels change places \emph{inside} each half.

\emph{{[}The glazed light leaning on the end rail has the same window in
it, right way up this time: the green pane at the top still, and on the
latch side. In the puddle every pane went over the bar across; in the
glass every pane has gone over the bar down.{]}}

\textbf{Tortoise:} One clip; one bare string; and the sheet?

\textbf{Achilles:} Moved all over, both times, and differently. There
isn't a panel anywhere that belongs to one coin. \emph{{[}The Tortoise
moves the clip back to the silk string.{]}} Then put one on each. ---
Which goes first?

\begin{center}\rule{0.5\linewidth}{0.5pt}\end{center}

\hypertarget{fourth-challenge-which-first}{%
\section{\texorpdfstring{Fourth Challenge --- \emph{Which
First?}}{Fourth Challenge --- Which First?}}\label{fourth-challenge-which-first}}

\textbf{Tortoise:} \emph{{[}She clips a second clip to the hemp string:
``heads up, whatever it was''.{]}} One on each.

\textbf{Achilles:} But which do the coins meet first? It makes a
difference; I have it in my notes. On one string, ``turn it over'' and
then ``heads up, whatever it was'' arrives heads; the same two in the
other order arrive tails. I shall write out both orders. \emph{{[}He
begins a column on the back of the blank sheet.{]}}

\textbf{Tortoise:} \emph{{[}She slides her clip along the hemp string
until it stands well to the left of his; it runs freely the whole way,
for his clip rides the silk above and there is nothing on the hemp to
check it.{]}} Walk them. Heads above, tails below.

\textbf{Achilles:} \emph{{[}walking the coins, level{]}} The hemp coin
meets its clip first: heads up. Then the silk coin meets its own: over.
They arrive tails above, heads below.

\textbf{Tortoise:} \emph{{[}She slides her clip rightward and it passes
clean under his --- the two never so much as graze, being on strings one
above the other --- and comes to rest at the far end.{]}} I slide it
past yours. Again.

\textbf{Achilles:} The silk coin first, this time: over. Then the hemp:
heads up. Tails above, heads below. The same!

\textbf{Tortoise:} \emph{{[}She slides her clip back until it stands
exactly under his.{]}} And level.

\textbf{Achilles:} Both at once. Tails above, heads below.

\emph{{[}He walks the coins from every other pair of starts, with her
clip left, right and level. The same faces arrive each time. His column
of orders is left unfinished.{]}}

\emph{{[}At the corner of the eave two icicles hang side by side, and
each has worn its own hollow in the snow beneath. Now the left one lets
go a drop first, now the right, now the two together, and there is no
telling which it will be; the hollows deepen alike. All along the foot
of the porch the pitting has become a row of such hollows, one to an
icicle.{]}}

\textbf{Achilles:} Those were single clips. With more of them there must
be an order somewhere.

\textbf{Tortoise:} \emph{{[}She adds a second clip to each string, to
the right of the first: on the silk ``heads up, whatever it was'', on
the hemp ``turn it over''. The first clips stand level, and the second
clips stand level.{]}} Take the silk; I shall take the hemp. We walk
together, and each of us says what is done.

\textbf{Achilles:} Silk. The coin starts tails up.

\textbf{Tortoise:} Hemp. The coin starts heads up.

\textbf{Achilles:} Silk, the first clip: turn it over. Heads.

\textbf{Tortoise:} Hemp, the first clip: heads up, whatever it was.
Heads.

\textbf{Achilles:} Silk, the second clip: heads up, whatever it was.
Heads still.

\textbf{Tortoise:} Hemp, the second clip: turn it over. Tails.

\textbf{Achilles:} Silk arrives heads.

\textbf{Tortoise:} Hemp arrives tails.

\textbf{Achilles:} Both first clips, and then both second clips. So
there \emph{was} an order, after all.

\textbf{Tortoise:} Set the coins back. Now walk all of the silk before I
stir.

\textbf{Achilles:} Tails; over, heads; heads up, heads. Heads, as
before.

\textbf{Tortoise:} Heads; heads up, heads; over, tails. Tails, as
before.

\textbf{Achilles:} The same again. First both firsts and then both
seconds; or all of the silk and then all of the hemp. But between the
two walks nobody touched a clip. \ldots Oh. Nothing was moved. It was
one frame, and I read it two ways.

\emph{{[}He looks up from the frame. On the wall beside the door hangs a
print --- Escher's \textup{Three Spheres II}.{]}}

\emph{{[}The shadow of the corner post, which lay over the print when
they sat down, has moved off along the wall, and the sun is full on it.
For the length of his look no drop falls anywhere along the eave.{]}}

\textbf{Achilles:} Your three balls stand side by side on their table
just so, and nobody asks which of them comes first. Though the middle
one has both its neighbours in its shine.

\textbf{Tortoise:} \emph{{[}On the silk string she pushes the ``turn it
over'' clip to the right, towards its neighbour.{]}} Now I slide this
one past that one.

\textbf{Achilles:} It won't go. They knock together. They are on the
same string!

\emph{{[}Down any one icicle the water goes in file. The drop behind
cannot get by the drop in front, and they come to the tip in the order
they set out.{]}}

\textbf{Tortoise:} \emph{{[}Because both clips share the one silk
string, neither can pass the other; she lifts the ``turn it over'' clip
clean off the silk and pinches it on again beyond its neighbour ---
where a moment ago, on two strings, the sliding had needed no lifting at
all.{]}} It must come off, and go back on the far side. Walk the silk.
Tails up.

\textbf{Achilles:} Heads up: heads. Over: tails. Silk arrives
\emph{tails}, and a moment ago it arrived heads. So the old fact stands:
along one string there is an order, and it is part of what the string
says.

\textbf{Tortoise:} \emph{{[}She takes the clip off once more and
restores it to the left of its neighbour.{]}} And between one string and
another?

\textbf{Achilles:} None to find. So strings keep to themselves. They
never so much as touch.

\emph{{[}Achilles sits back, pleased with himself, and lets the sun have
his hands a moment. The boards beneath him, where it has lain all
morning, have gone dry and warm.{]}}

\emph{{[}The Tortoise reaches down to the right-hand upright and unhooks
a loop.{]}}

\begin{center}\rule{0.5\linewidth}{0.5pt}\end{center}

\hypertarget{fifth-challenge-the-crossing}{%
\section{\texorpdfstring{Fifth Challenge --- \emph{The
Crossing}}{Fifth Challenge --- The Crossing}}\label{fifth-challenge-the-crossing}}

\textbf{Tortoise:} Take the clips off, Achilles. We want bare strings.

\emph{{[}Achilles takes the clips off the silk and the hemp and puts
them back in their box. On the lower pair the gut runs straight along
its own row to its own hook, the horsehair straight along the row
beneath to its. The Tortoise takes hold of the two right-hand loops, the
gut's in one hand and the horsehair's in the other.{]}}

\textbf{Tortoise:} I unhook the right-hand loops of the gut and the
horsehair. The upper comes down in front, the lower goes up behind, and
each takes the other's hook. \emph{{[}She does so: the left-hand ends
never move, and neither string is cut or joined to the other, yet the
two right ends trade hooks --- the gut's onto the lower hook, the
horsehair's onto the upper --- so the pair exchange their two rows and
cross once, flat, the one lying against the other.{]}} Walk a coin on
each. A coin keeps to its string.

\textbf{Achilles:} \emph{{[}He sets a coin heads up on the gut and a
coin tails up on the horsehair, and walks them along, each following its
string through the crossing.{]}} The heads coin arrives on the
horsehair's hook. The tails coin on the gut's. They have changed hooks,
and neither has changed its face.

\emph{{[}Two sparrows have the porch rail to themselves, a cock and a
hen. The cock hops past in front of the hen, the hen shuffles along
behind him, and when they have settled each is at the other's end of the
rail, and the same bird as before.{]}}

\textbf{Tortoise:} \emph{{[}She clips ``turn it over'' to the gut
string, to the left of the crossing.{]}} Both heads up, this time.

\textbf{Achilles:} The gut coin is turned, then goes down through the
crossing: tails, on the horsehair's hook. The other arrives heads, on
the gut's hook.

\textbf{Tortoise:} \emph{{[}She slides the clip rightward along the gut;
at the crossing she lifts the horsehair a finger's breadth, so the clip
rides under it and stays with the gut, coming out on the lower row ---
its row changed, its string the same.{]}} And now?

\textbf{Achilles:} Now the clip is on the lower row, so it will turn the
horsehair coin.

\textbf{Tortoise:} Which string is it clipped to?

\textbf{Achilles:} \emph{{[}following the string back with a finger{]}}
The gut. Still the gut. \emph{{[}He walks the coins, both heads up.{]}}
The gut coin comes down through the crossing and \emph{then} is turned:
tails, on the horsehair's hook; the other heads, on the gut's. As before
the slide, face for face. The clip kept its string; the string changed
its place.

\textbf{Tortoise:} \emph{{[}She takes the clip off and puts it in the
box.{]}} Your turn to say it.

\emph{{[}The lower pair still lies crossed once: at the right-hand end
the horsehair is uppermost now, the gut beneath it.{]}}

\textbf{Achilles:} I unhook the right-hand loops of the gut and the
horsehair. The upper comes down in front, the lower goes up behind, and
each takes the other's hook. \emph{{[}He does so, the same hand-movement
as before --- but the loop now uppermost is the horsehair, so this time
the horsehair passes in front and the gut behind; each string has now
gone in front just once. The right ends come home to their own hooks,
and still the pair will not lie flat.{]}} There: back on their own
hooks, and --- they are wound round each other.

\textbf{Tortoise:} Pull them straight.

\textbf{Achilles:} \emph{{[}He draws the strings taut. He plucks at
them.{]}} It won't pull out. And I said what you said and did what you
did!

\textbf{Tortoise:} You did. Walk a coin on each.

\textbf{Achilles:} \emph{{[}One coin heads up, one tails up; each
follows its string through the one crossing, and through the next.{]}}
Heads arrives on the gut's hook, tails on the horsehair's. Their own
hooks, their own faces. \emph{{[}He walks two coins along the bare silk
and hemp.{]}} And along a straight pair: own hooks, own faces. The same
report, word for word.

\emph{{[}On the rail the sparrows have changed ends again, the hen going
by in front this time. Each is back where it began.{]}}

\textbf{Tortoise:} \emph{{[}On a scrap of paper she draws two lines side
by side that change places, flat, and further along change places again.
She rubs the middle out, and draws the two lines straight.{]}} I have
only the one kind of crossing on paper, and I keep a rule for it: two of
them, one after the other, I may draw straight.

\textbf{Achilles:} But the strings are still wound! The paper doesn't
make them straight, and no pulling does. Your rule is false of the
string.

\textbf{Tortoise:} Of the string, yes. The string remembers which went
in front. The coin cannot tell. I draw for the coin.

\textbf{Achilles:} Then it is a rule one has to earn, by walking the
coins; the paper won't vouch for it. \emph{{[}He reaches for the
loops.{]}} Shall I unwind them?

\textbf{Tortoise:} Leave them wound.

\textbf{Achilles:} Very well. Crossed or not, each coin still minds only
its own string.

\emph{{[}Achilles flexes his fingers, stiff from the cool of the boards.
The sun by now has come round far enough to warm his back.{]}}

\emph{{[}The Tortoise takes a short wooden bar out of her shell. It has
a clip at either end.{]}}

\textbf{Tortoise:} A bar of wood, a clip at either end --- long enough
to reach across two neighbouring strings, and no further. Where a clip
lays a rule on one string, this lays a single rule across the pair at
once.

\begin{center}\rule{0.5\linewidth}{0.5pt}\end{center}

\hypertarget{sixth-challenge-the-yoke}{%
\section{\texorpdfstring{Sixth Challenge --- \emph{The
Yoke}}{Sixth Challenge --- The Yoke}}\label{sixth-challenge-the-yoke}}

\textbf{Tortoise:} \emph{{[}She lays the bar across the silk and the
hemp near the left upright, a clip biting each string --- just long
enough to span the two and no more; both strings run on beneath it
unbroken, a single line through each clip. She reads what is burnt into
the wood.{]}} ``If the upper coin shows tails, turn the lower one
over.''

\textbf{Achilles:} \emph{{[}He walks the pairs through it.{]}} Heads
above: the lower coin goes through as it came, heads or tails. Tails
above: the lower coin is turned. The upper coin is never touched. That
is two clips nailed to a stick.

\textbf{Tortoise:} \emph{{[}She unclips the bar and pushes the box of
clips towards him.{]}} Then do it with two clips.

\textbf{Achilles:} On the silk, nothing; the upper coin is never
touched. On the hemp: ``turn it over''. \emph{{[}He clips it on, and
walks.{]}} Tails above, heads below: arrives tails, tails. Right! Heads
above, heads below: arrives heads, tails. Wrong. \emph{{[}He takes the
clip off.{]}} And with no clip, right for heads above and wrong for
tails. I want a cleverer clip than you keep in the box.

\textbf{Tortoise:} \emph{{[}She clips the bar on again and stands the
card on edge between the silk and the hemp, running it from the left
upright only as far as the bar --- for the bar itself blocks the way the
card would go on. He stands on the silk side of it, she on the hemp
side; from her seat the silk coin cannot be seen at all, and from his
the hemp coin is hidden the same.{]}} Write me its rule. It sits on the
hemp string, on my side of the card.

\textbf{Achilles:} ``If the upper coin---''

\textbf{Tortoise:} Which coin is that? I see one string from here.

\textbf{Achilles:} ``If\ldots{}'' No.~Whatever I write on it, it must
turn your coin on some mornings and leave it on others, and it cannot
see which morning it is. And the card cannot go any further along,
because the bar goes through where the card would stand. \ldots Oh. The
stick isn't holding two clips together. The stick is how the lower one
looks up.

\emph{{[}Across the yard the woodshed roof lets go its snow at one end,
all in a piece. It comes down on the lean-to roof below; and the
lean-to, which had held its own snow all morning, lets that go too, and
the whole lot lands in the yard with a thump that makes Achilles look
round. At the other end the upper roof has kept its snow, and the
lean-to under it keeps its.{]}}

\textbf{Tortoise:} Then walk with me. Toss yours; mine starts heads up.

\emph{{[}Achilles tosses the silk coin and lays it as it fell. They walk
their coins along, the card between them. At the bar the Tortoise
stops.{]}}

\textbf{Tortoise:} What does yours show, Achilles?

\textbf{Achilles:} Tails.

\textbf{Tortoise:} Then mine goes over. Tails.

\textbf{Achilles:} You waited for me. All morning you never once waited
for me.

\emph{{[}Achilles sits back, and for a moment neither of them says
anything. He warms his hands on the cup, tepid now, and lets her have
the last word.{]}}

\textbf{Tortoise:} All morning I had a string to myself. \emph{{[}She
folds a fresh sheet twice, as the first was folded, and gives it to
him.{]}} A mark in the panel where each pair arrives.

\emph{{[}A dozen tosses. Each time she stops at the bar and asks, and he
answers, and only then does she move.{]}}

\textbf{Achilles:} Seven marks in the upper left, heads with heads. Five
in the lower right, tails with tails. In the other two panels, not one.
As shares: a half, nothing, nothing, a half. Now the border lists.

\textbf{Tortoise:} Begin with an empty panel.

\textbf{Achilles:} \emph{{[}He writes, and stops.{]}} The upper right
panel is nothing, so one of its border lines must be nothing: the silk's
heads line or the hemp's tails line. But the first would empty the upper
left panel, and the second the lower right, and those are full. It will
not multiply out. It isn't a patchwork!

\textbf{Tortoise:} A pair of gloves, posted in two parcels. Open either
one: left or right, evens. Open one, and you know the other.

\textbf{Achilles:} And no pair of clips could do it. Two clips make a
patchwork of a patchwork.

\textbf{Tortoise:} \emph{{[}She takes down the card and, unhooking
nothing, slides the winding of the lower pair apart along the strings
until its two crossings stand well clear of one another. She takes a
second bar from her shell and clips it across the silk and the hemp to
the right of the first; below, the wound pair's first crossing now lies
between the two bars, its second beyond them --- the two orders
straddling each other along the one frame, and no coin yet touched.{]}}
The same rule, burnt into the same wood. Walk the pairs.

\textbf{Achilles:} Heads above: nothing happens, twice. Tails above,
heads below: the lower coin goes over at the first bar, and over again
at the second. Heads. Tails above, tails below: heads, then tails again.

\textbf{Tortoise:} And a tossed one?

\textbf{Achilles:} \emph{{[}He tosses the silk coin, and walks; and
tosses again.{]}} The tossed one too: mine arrives as it fell, and the
hemp coin arrives heads every time. Every pair comes out as it went in.
Twice yoked, and nothing shows!

\textbf{Tortoise:} Below, more room on the sheet than the strings have
lines. On different strings, no first and no second. Crossed, a change
of place and nothing else. Yoked, one string minding another. And all of
it at once, on one frame.

\emph{{[}The thaw has the whole eave now. Every icicle keeps its own
time over its own hollow, and the water from all of them has found the
path and goes down it to the gate in one bright runnel, past the pear
tree, past the woodshed's bare end. In the seed tray the labels stand
round their two edges in the sun.{]}}

\emph{{[}Achilles sets a coin on each string at the left upright: on the
silk, the hemp, the gut, the horsehair.{]}}

\textbf{Achilles:} Shall I walk the upper pair first, or the lower? ---
No.~I know better than to ask.

\emph{{[}He walks them, the upper pair through the two bars, the lower
pair through their two crossings. Every coin arrives on its own hook,
showing the face it set out with.{]}}

\textbf{Tortoise:} Every one home. And your tea gone stone cold, never
touched --- warm your hands a moment; the eaves have all but stopped
dripping, and the path will bear you well when you go.

\textbf{Achilles:} As they left, every one. A whole frame of bars and
windings, and it comes to bare strings. Tortoise --- this morning I had
one string and a list of two lines, and I thought I knew the rest. If
one had begun with only the one string, how would one ever have guessed
all this?

\textbf{Tortoise:} Draw a second string below the first, and we shall
see.

\emph{{[}The sun has climbed: the porch roof's shadow has come down the
wall and over the print beside the door, and the three balls stand in
their row in the shade, a picture on a porch wall. On the table the
frame stands full in the light, the silk a line of light as it was when
the morning began, the hemp taking none, the gut yellow, the horsehair
dark; below them the empty hooks go on down the uprights. Along the eave
the icicles have dripped themselves thin, and the drops come slowly now,
and evenly, each into its own hollow. Two cups stand cold on the table,
and beside the frame lies a fresh sheet of paper, not yet folded, warm
from the sun.{]}}

\hypertarget{chapter-4-uniting-wires}{%
\chapter{Chapter 4: Uniting Wires}\label{chapter-4-uniting-wires}}

\emph{--- Yoked Inventions ---}

\emph{Part II: Wires Together}

\emph{Escher: ``Three Spheres II'' --- three spheres side by side on a
table top, the middle one holding its neighbours and the room in its
reflection}

\begin{center}\rule{0.5\linewidth}{0.5pt}\end{center}

Put a second coin on the table beside the first: is there now twice as
much to say, or more --- and if something is done to each of them, which
was done first?

Two travellers set out from two different towns. To describe either
journey you need a timetable, and a timetable is a list. For the pair of
them, then, two lists? That is what anyone would say, and for a great
many purposes it is right. But ask an innocent question about the two of
them \emph{together} --- what is the chance that she is delayed and he
is not? --- and the two lists fall silent. Neither of them has a line
for it. A description of two things must have room for every combination
of what the two might do, and combinations do not add. They multiply.
With two coins the difference hides, since two and two are four and two
times two are four. It stops hiding at the third coin, and by the tenth
it is the difference between twenty and a thousand.

The second half of the question sounds like a riddle. Two musicians play
in two different rooms, each from her own part. We may ask which bar of
the first part comes before which other bar of the first part: the page
says. But is the third bar of one part before or after the fifth bar of
the other? We could decide to say so; we could record the two rooms and
lay the tapes side by side any way we pleased. Nothing in the music
would change. Between two things that do not touch, ``first'' and
``second'' are not facts; they are habits of the person writing it down.

This chapter is about both halves, and about how each looks in our two
notations. On the blackboard the second coin brings a new kind of
multiplication, and an honest half page of indices to show that it
behaves. In the sketchbook it brings almost nothing: a second line,
drawn under the first. We shall find that the drawing had already
answered the riddle, before anyone asked, simply by being a drawing.
Whether it was \emph{entitled} to answer is another matter, and we shall
have to be honest about who pays for that. Then we make the two wires
notice each other, in the only two ways there are: they can change
places, and a box can be laid across them both.

\begin{center}\rule{0.5\linewidth}{0.5pt}\end{center}

\hypertarget{a-second-wire-below}{%
\section{§4.1 A Second Wire, Below}\label{a-second-wire-below}}

Chapter 3 ended with four questions, and this chapter answers them in
the order in which they were asked. The first: \emph{is the second
coin's wire drawn beside the first wire --- above it, below it?}

Below. And since we shall have to refer to the two wires constantly, we
fix once and for all which of them is mentioned first: \textbf{the top
wire is the first.} Nothing deep is being decided. But it must be
decided once and then never varied, because a good deal of what follows
looks different, though it is not, if the two wires are silently
exchanged.

\textbf{On the blackboard.} A single coin, when looked at, gives one of
two answers, \(0\) or \(1\). Two coins give one of four \emph{pairs} of
answers, and we write a pair as two digits, the first coin's answer
first: \(00\), \(01\), \(10\), \(11\). Each pair gets a ket of its own,

\[\ket{00}, \qquad \ket{01}, \qquad \ket{10}, \qquad \ket{11},\]

and each ket is a column of four entries, with a single \(1\) in the
place that belongs to its pair. The places are taken in the order just
written, which is the order of counting in binary:

\[\ket{00}=\begin{pmatrix}1\\0\\0\\0\end{pmatrix}, \quad \ket{01}=\begin{pmatrix}0\\1\\0\\0\end{pmatrix}, \quad \ket{10}=\begin{pmatrix}0\\0\\1\\0\end{pmatrix}, \quad \ket{11}=\begin{pmatrix}0\\0\\0\\1\end{pmatrix}.\]

We shall also write \(\ket{a,b}\) for \(\ket{ab}\) when the comma helps
the eye. In \(\ket{ab}\) the digit \(a\) is the answer of the first
coin, the top wire; \(b\) is the answer of the second, the bottom wire.
So \(\ket{01}\) says: top \(0\), bottom \(1\).

\textbf{In the sketchbook.} One new element, described before it is
drawn. To \textbf{stack} two pictures is to draw one below the other,
with clear paper between them, each on its own wire. The simplest stack
is two bare wires, one under the other. The next simplest is two states,
each a rounded box with its wire leaving to the right, exactly as in
Chapter 1, the one on the upper wire and the other on the lower. The
stack of the state \(a\) over the state \(b\) is the picture of the
column \(\ket{ab}\):

\[\begin{tikzpicture}[sd]
  \draw[wire] (0,2.2) -- (3,2.2);
  \draw[wire] (0,0) -- (3,0);
  \node[sdstate, text height=1.5ex, text depth=0.5ex] at (0,2.2) {$a$};
  \node[sdstate, text height=1.5ex, text depth=0.5ex] at (0,0) {$b$};
\end{tikzpicture}
\;\;=\;\; \ket{ab}
\qquad\qquad
\begin{tikzpicture}[sd]
  \draw[wire] (0,2.2) -- (3,2.2);
  \draw[wire] (0,0) -- (3,0);
  \node[sdstate, text height=1.5ex, text depth=0.5ex] at (0,2.2) {$0$};
  \node[sdstate, text height=1.5ex, text depth=0.5ex] at (0,0) {$1$};
\end{tikzpicture}
\;\;=\;\; \ket{01} \;=\; \begin{pmatrix}0\\1\\0\\0\end{pmatrix}\]

Pictures are still read from left to right. What is new is that the page
now has a second direction, and it means something different from the
first. Left to right is \emph{then}. Top to bottom is \emph{and}. Every
picture equation in this chapter is exact: the same number, or the same
matrix, on both sides, with no factor swept aside.

\textbf{\emph{Why this definition.}} Why below, and why ``top first''?
For no reason except that a choice is needed and conventions differ:
this book reads the top wire first, and the reference for the general
theory that we shall cite in §4.9 happens to stack the other way. What
is \emph{not} a matter of taste is that the choice is made once. The
column \(\ket{01}\) and the column \(\ket{10}\) are different columns,
and a reader who exchanged the wires without saying so would be
exchanging them. See \texttt{Code/Ch4/SecondString.py}, which asserts
the order.

\begin{center}\rule{0.5\linewidth}{0.5pt}\end{center}

\hypertarget{the-column-for-two-coins}{%
\section{§4.2 The Column for Two Coins}\label{the-column-for-two-coins}}

The second question of §3.10: \emph{what is the column for two coins ---
two entries and two more, or something else?}

Begin with coins we understand, the covered coins of §1.4, whose columns
are lists of shares. Let the upper coin be heads and tails in equal
shares, \(\begin{pmatrix}1/2\\1/2\end{pmatrix}\), and the lower one
heads three times in four, \(\begin{pmatrix}3/4\\1/4\end{pmatrix}\), and
suppose the two coins have nothing to do with each other. ``Two entries
and two more'' would be a list of four numbers,
\(\tfrac12,\tfrac12,\tfrac34,\tfrac14\), and it has a flaw that shows at
once: its entries sum to \(2\). It has a worse flaw. Ask for the share
of \emph{heads above and tails below, together}, and no line of it
answers. A description of the pair needs one line for each pair of
answers, and when the coins are separate the line for a pair is the
product of the two shares:

\[\begin{pmatrix}\tfrac12\cdot\tfrac34\\[2pt]\tfrac12\cdot\tfrac14\\[2pt]\tfrac12\cdot\tfrac34\\[2pt]\tfrac12\cdot\tfrac14\end{pmatrix} \;=\; \begin{pmatrix}3/8\\1/8\\3/8\\1/8\end{pmatrix}, \qquad \tfrac38+\tfrac18+\tfrac38+\tfrac18=1 .\]

So it is four entries in either case, but they are not two and two more.
They are two \emph{times} two.

\textbf{On the blackboard.} The \textbf{Kronecker product} of two
columns of two entries is the column of four entries

\[\begin{pmatrix}a\\b\end{pmatrix}\otimes\begin{pmatrix}c\\d\end{pmatrix} \;=\; \begin{pmatrix}ac\\ad\\bc\\bd\end{pmatrix}.\]

The sign \(\otimes\) is read ``tensor'', and physicists call the
operation the \textbf{tensor product}. The recipe: take each entry of
the \emph{left} column in turn, and multiply the whole right column by
it. The entry in the place \(ik\) is the \(i\)-th entry of the left
column times the \(k\)-th entry of the right. The order of the places is
thus the order of §4.1, and the basis kets are themselves products,
\(\ket{ab}=\ket{a}\otimes\ket{b}\). \textbf{The left factor is the top
wire.} This is fixed here and does not vary again.

The product passes through sums and multiples in each factor separately,
because every entry of it is one entry from the left times one from the
right:

\[\bigl(\ket{\psi}+\ket{\psi'}\bigr)\otimes\ket{\phi}=\ket{\psi}\otimes\ket{\phi}+\ket{\psi'}\otimes\ket{\phi},\]

\[\bigl(\lambda\ket{\psi}\bigr)\otimes\ket{\phi}=\lambda\bigl(\ket{\psi}\otimes\ket{\phi}\bigr)=\ket{\psi}\otimes\bigl(\lambda\ket{\phi}\bigr),\]

and likewise on the right. Note where the number \(\lambda\) ends up in
the last equation: it belongs to neither factor. It belongs to the pair.

For two \emph{qubits} the entries are amplitudes and not shares, and we
take the same product. Here is one, with entries that are not real so
that nothing can hide. The upper qubit is the arrow of Chapter 3,
\(\ket{a}=\tfrac35\ket{0}+\tfrac45\ket{1}\), and the lower is
\(\ket{\omega}=\tfrac1{\sqrt2}\bigl(\ket{0}+i\ket{1}\bigr)\) from §3.3:

\[\ket{a}\otimes\ket{\omega} \;=\; \begin{pmatrix}3/5\\4/5\end{pmatrix}\otimes\tfrac1{\sqrt2}\begin{pmatrix}1\\i\end{pmatrix} \;=\; \tfrac1{5\sqrt2}\begin{pmatrix}3\\3i\\4\\4i\end{pmatrix}.\]

The squared moduli are
\(\tfrac9{50},\tfrac9{50},\tfrac{16}{50},\tfrac{16}{50}\), and they sum
to \(1\). Each is the product of the two separate chances,
\(\tfrac9{25}\cdot\tfrac12\) and \(\tfrac{16}{25}\cdot\tfrac12\), as it
should be for two coins that have nothing to do with each other: the
squared modulus of a product is the product of the squared moduli. And
the order of the factors matters:
\(\ket{\omega}\otimes\ket{a}=\tfrac1{5\sqrt2}(3,\;4,\;3i,\;4i)^{T}\), a
different column.

In general, \textbf{a state of two qubits is a column of four complex
numbers whose squared moduli sum to \(1\)}; the set of all columns of
four complex numbers, thought of as belonging to two wires, is written
\(\mathbb{C}^2\otimes\mathbb{C}^2\). A state of the form
\(\ket{\psi}\otimes\ket{\phi}\) is called a \textbf{product state}. That
a product of two states \emph{is} a state --- that its squared moduli
sum to \(1\) --- is a special case of the next fact, which is the first
thing one wants to know about the new product: how it meets the old
bracket.

The bra of a column is the row of its conjugates (§2.2), and the
conjugate of a product of numbers is the product of the conjugates. So
the bra of \(\ket{\psi}\otimes\ket{\phi}\) is the row
\(\begin{pmatrix}\bar a\bar c&\bar a\bar d&\bar b\bar c&\bar b\bar d\end{pmatrix}\),
which is the same recipe applied to the two rows: it is
\(\bra{\psi}\otimes\bra{\phi}\).

An example first. Ask the turned question of §2.5 on both wires: what is
the bracket of \(\ket{+}\otimes\ket{+}=\tfrac12(1,1,1,1)^{T}\) with the
state above? The row \(\tfrac12\begin{pmatrix}1&1&1&1\end{pmatrix}\)
adds up the entries and halves them:

\[\tfrac12\cdot\tfrac1{5\sqrt2}\,(3+3i+4+4i) \;=\; \tfrac{7+7i}{10\sqrt2} .\]

On the separate wires,
\(\braket{+|a}=\tfrac1{\sqrt2}\bigl(\tfrac35+\tfrac45\bigr)=\tfrac7{5\sqrt2}\)
and
\(\braket{+|\omega}=\tfrac1{\sqrt2}\cdot\tfrac1{\sqrt2}(1+i)=\tfrac{1+i}2\),
and their product is \(\tfrac{7(1+i)}{10\sqrt2}\): the same number. Is
that luck?

\begin{quote}
\textbf{Proposition 4.1 (brackets multiply).} For all columns
\(\ket{\psi},\ket{\psi'},\ket{\phi},\ket{\phi'}\) of two complex
entries,

\[\bigl(\bra{\psi'}\otimes\bra{\phi'}\bigr)\bigl(\ket{\psi}\otimes\ket{\phi}\bigr) \;=\; \braket{\psi'|\psi}\;\braket{\phi'|\phi} .\]

In particular the Kronecker product of two states is a state. In
pictures: two closed diagrams stacked are the product of their numbers.
\end{quote}

\emph{Proof idea.} A sum of four products that factorizes into two sums
of two. \emph{Proof sketch.} Write \(\ket{\psi}=(a,b)^{T}\),
\(\ket{\phi}=(c,d)^{T}\), and the primed columns with primed letters.
The left side is a row times a column,

\[\bar a'\bar c'\,ac+\bar a'\bar d'\,ad+\bar b'\bar c'\,bc+\bar b'\bar d'\,bd \;=\; \bigl(\bar a'a+\bar b'b\bigr)\bigl(\bar c'c+\bar d'd\bigr),\]

as one sees by multiplying out the right side; and the two brackets on
the right are \(\braket{\psi'|\psi}\) and \(\braket{\phi'|\phi}\). For
the particular case take \(\psi'=\psi\) and \(\phi'=\phi\) and use
Theorem 2.1 twice: \(1\cdot1=1\). \(\;\blacksquare\) See
\texttt{Code/Ch4/SecondString.py}. We shall see in §4.9 that this
proposition is the smallest case of the chapter's main theorem.

\textbf{In the sketchbook.} No element beyond stacking. Two states
stacked stand for the Kronecker product of their columns, the upper
state on the left of the \(\otimes\):

\[\begin{tikzpicture}[sd]
  \draw[wire] (0,2.2) -- (3,2.2);
  \draw[wire] (0,0) -- (3,0);
  \node[sdstate, text height=1.5ex, text depth=0.5ex] at (0,2.2) {$\psi$};
  \node[sdstate, text height=1.5ex, text depth=0.5ex] at (0,0) {$\phi$};
\end{tikzpicture}
\;\;=\;\; \ket{\psi}\otimes\ket{\phi}\]

Effects stack in the same way, and so do closed diagrams. A closed
diagram is a number (§2.2), and Proposition 4.1 says what number a stack
of two of them is:

\[\begin{tikzpicture}[sd]
  \draw[wire] (0,2.2) -- (4.5,2.2);
  \draw[wire] (0,0) -- (4.5,0);
  \node[sdstate, text height=1.5ex, text depth=0.5ex] at (0,2.2) {$\psi$};
  \node[sdeffect, text height=1.5ex, text depth=0.5ex] at (4.5,2.2) {$\psi'$};
  \node[sdstate, text height=1.5ex, text depth=0.5ex] at (0,0) {$\phi$};
  \node[sdeffect, text height=1.5ex, text depth=0.5ex] at (4.5,0) {$\phi'$};
\end{tikzpicture}
\;\;=\;\; \braket{\psi'|\psi}\;\braket{\phi'|\phi}
\;\;=\;\;
\begin{tikzpicture}[sd, baseline=-0.55ex]
  \draw[wire] (0,0) -- (4.5,0);
  \draw[wire] (6.5,0) -- (11,0);
  \node[sdstate, text height=1.5ex, text depth=0.5ex] at (0,0) {$\psi$};
  \node[sdeffect, text height=1.5ex, text depth=0.5ex] at (4.5,0) {$\psi'$};
  \node[sdstate, text height=1.5ex, text depth=0.5ex] at (6.5,0) {$\phi$};
  \node[sdeffect, text height=1.5ex, text depth=0.5ex] at (11,0) {$\phi'$};
\end{tikzpicture}\]

The picture on the right is the old one: closed diagrams \emph{in a row}
multiply (row 12 of the Rosetta Table). So for closed diagrams it makes
no difference whether they are set in a row or stacked. That is as it
should be. A number has no wire by which to be before or after anything,
or above or below it.

\textbf{\emph{Why this definition.}} Why products of entries, and not
something simpler? For shares there was never a choice: the chance of
two independent things both happening is the product of their chances.
For amplitudes the requirement is that the chances \emph{computed from
them} by the Born rule should multiply in that way, for every question
asked separately on the two wires, and Proposition 4.1 says they do: the
chance of the answers \(\psi'\) above and \(\phi'\) below is
\(\lvert\braket{\psi'|\psi}\braket{\phi'|\phi}\rvert^2\), the product of
the two separate chances. Why not the simpler ``two entries and two
more''? Because a column is a list of answers to questions, and the
questions about a pair are about pairs.

\begin{center}\rule{0.5\linewidth}{0.5pt}\end{center}

\hypertarget{doubling}{%
\section{§4.3 Doubling}\label{doubling}}

There was a trap in the last section, and the numbers set it. Two wires,
four entries; and \(2+2=4\). A hasty reader could still believe that
columns lengthen by two for each wire, and would expect six entries for
three wires.

\textbf{On the blackboard.} Three coins give one of eight triples of
answers, \(000,001,\dots,111\), and the column has a place for each. The
product of three columns is formed by taking the product twice, and it
does not matter how the three are bracketed:
\(\bigl(\ket{\psi}\otimes\ket{\phi}\bigr)\otimes\ket{\chi}\) and
\(\ket{\psi}\otimes\bigl(\ket{\phi}\otimes\ket{\chi}\bigr)\) both have,
in the place \(ikm\), the number \(\psi_i\phi_k\chi_m\) --- the \(i\)-th
entry of the first column times the \(k\)-th of the second times the
\(m\)-th of the third --- and both list the places in the order of
counting. So we write \(\ket{\psi}\otimes\ket{\phi}\otimes\ket{\chi}\)
without brackets. Each further wire doubles the column:

\begin{longtable}[]{@{}
  >{\centering\arraybackslash}p{(\columnwidth - 4\tabcolsep) * \real{0.3333}}
  >{\centering\arraybackslash}p{(\columnwidth - 4\tabcolsep) * \real{0.3333}}
  >{\centering\arraybackslash}p{(\columnwidth - 4\tabcolsep) * \real{0.3333}}@{}}
\toprule\noalign{}
\begin{minipage}[b]{\linewidth}\centering
wires \(n\)
\end{minipage} & \begin{minipage}[b]{\linewidth}\centering
entries of the column, \(2^n\)
\end{minipage} & \begin{minipage}[b]{\linewidth}\centering
numbers needed for \(n\) separate states, \(2n\)
\end{minipage} \\
\midrule\noalign{}
\endhead
\bottomrule\noalign{}
\endlastfoot
\(1\) & \(2\) & \(2\) \\
\(2\) & \(4\) & \(4\) \\
\(3\) & \(8\) & \(6\) \\
\(4\) & \(16\) & \(8\) \\
\end{longtable}

\textbf{A column for \(n\) wires has \(2^n\) entries}, one for each
string of \(n\) answers. The third column of the table counts something
else: a product of \(n\) one-wire states is fixed by naming the \(n\)
states, two complex numbers each, \(2n\) in all. At one wire and at two
the columns of the table agree, which is what made the trap possible;
from three wires on they part, and they never meet again. At ten wires
it is \(1024\) against \(20\); at forty the column has more than a
million million entries.

The gap has a plain meaning, which we state and leave. The columns that
are products of one-wire states are a thin family inside the space of
all columns. \textbf{Most columns are not products, and what that means
is the subject of Chapter 6.} Here we need only the count.

\textbf{In the sketchbook.} Stacking goes on downwards: a third wire
under the second, a fourth under the third, the top wire always first.

\[\begin{tikzpicture}[sd]
  \draw[wire] (0,4.4) -- (3,4.4);
  \draw[wire] (0,2.2) -- (3,2.2);
  \draw[wire] (0,0) -- (3,0);
  \node[sdstate, text height=1.5ex, text depth=0.5ex] at (0,4.4) {$a$};
  \node[sdstate, text height=1.5ex, text depth=0.5ex] at (0,2.2) {$b$};
  \node[sdstate, text height=1.5ex, text depth=0.5ex] at (0,0) {$c$};
\end{tikzpicture}
\;\;=\;\; \ket{abc} \;=\; \ket{a}\otimes\ket{b}\otimes\ket{c}, \qquad \text{a column of eight entries.}\]

Here the two notations part company in a way worth noticing. On the
blackboard each new wire doubles the writing. In the sketchbook each new
wire costs one more line of ink: \(n\) wires, \(n\) lines. The picture
is compact. But let us not claim too much for it. The picture is short
because it does not write the column out; it does not make the column
any shorter. Whoever wants the numbers must still find \(2^n\) of them,
and no drawing in this chapter will spare anyone that labour. A notation
that is brief is not thereby a calculation that is cheap.

\textbf{\emph{Why this definition.}} Why one entry for each string of
answers, when \(2n\) numbers describe \(n\) coins so much more cheaply?
Because \(2n\) numbers describe only coins that have nothing to do with
one another. Already for covered coins this is too little: a list of
shares for a pair must be able to say ``both heads or both tails, never
one of each'' (we shall meet exactly this list in §4.7), and no two
separate lists can say it. The column has a place for every combination
because every combination is a question that can be asked. See
\texttt{Code/Ch4/Doubling.py}.

\begin{center}\rule{0.5\linewidth}{0.5pt}\end{center}

\hypertarget{a-box-on-one-wire}{%
\section{§4.4 A Box on One Wire}\label{a-box-on-one-wire}}

The third question of §3.10: \emph{if a box sits on one wire and nothing
on the other, what is the matrix of the pair?}

One is tempted to answer that a box on the upper wire leaves the lower
coin alone, and so leaves half of the column alone. It does leave the
lower coin alone. It leaves \emph{no} entry of the column alone, because
every entry of the column is about both coins. Let us see what must
happen. Put NOT on the top wire and nothing on the bottom. The pair
\(\ket{ab}\) must go to the pair with \(a\) turned over and \(b\) as it
was: \(\ket{00}\mapsto\ket{10}\), \(\ket{01}\mapsto\ket{11}\),
\(\ket{10}\mapsto\ket{00}\), \(\ket{11}\mapsto\ket{01}\). The column of
a matrix is where it sends the corresponding basis ket (§1.1), so the
matrix is known, and it is \(4\times4\):

\[\begin{pmatrix}0&0&1&0\\0&0&0&1\\1&0&0&0\\0&1&0&0\end{pmatrix}.\]

With NOT on the \emph{bottom} wire instead, \(\ket{00}\mapsto\ket{01}\),
\(\ket{01}\mapsto\ket{00}\), \(\ket{10}\mapsto\ket{11}\),
\(\ket{11}\mapsto\ket{10}\), and the matrix is a different one:

\[\begin{pmatrix}0&1&0&0\\1&0&0&0\\0&0&0&1\\0&0&1&0\end{pmatrix}.\]

Cut each of them into four blocks of size \(2\times2\), and a pattern
appears. The first is \(\begin{pmatrix}0&I\\I&0\end{pmatrix}\): it is
the matrix \(X\) with each of its entries multiplied into a copy of
\(I\). The second is \(\begin{pmatrix}X&0\\0&X\end{pmatrix}\): it is the
matrix \(I\) with each of its entries multiplied into a copy of \(X\).

\textbf{On the blackboard.} Number the entries of a \(2\times2\) matrix
from zero, row first, as in §3.1. The \textbf{Kronecker product} of two
\(2\times2\) matrices is the \(4\times4\) matrix

\[A\otimes B \;=\; \begin{pmatrix}A_{00}\,B&A_{01}\,B\\A_{10}\,B&A_{11}\,B\end{pmatrix},\]

each block being the whole matrix \(B\) multiplied by one entry of
\(A\). The rows and the columns of \(A\otimes B\) are labelled by pairs,
in the order \(00,01,10,11\), and the entry in the row \(ik\) and the
column \(jl\) is

\[(A\otimes B)_{ik,\,jl} \;=\; A_{ij}\,B_{kl} .\]

The first digit of each label goes to \(A\) and the second to \(B\):
\textbf{the left factor is the top wire}, as for columns. The same
recipe serves for matrices of any shapes, the labels running over
whatever ranges the shapes allow; the product of two columns in §4.2 is
the case in which each factor has a single column, and the product of
two rows is the case of a single row.

The definition was chosen to make one thing true, and it does. The
column \(jl\) of \(A\otimes B\) has the entries \(A_{ij}B_{kl}\), and
that is the Kronecker product of the \(j\)-th column of \(A\) with the
\(l\)-th column of \(B\):

\[(A\otimes B)\bigl(\ket{j}\otimes\ket{l}\bigr) \;=\; A\ket{j}\otimes B\ket{l} .\]

\textbf{Each factor acts on its own wire.} (The same holds with any
columns \(\ket{\psi}\) and \(\ket{\phi}\) in the places of \(\ket{j}\)
and \(\ket{l}\); that will be a special case of Theorem 4.2, whose proof
does not depend on it.)

Now the question has its answer. A box \(A\) on the top wire, with
nothing on the bottom wire, is \(A\otimes I\); on the bottom wire, with
nothing on the top, it is \(I\otimes A\):

\[A\otimes I=\begin{pmatrix}A_{00}\,I&A_{01}\,I\\A_{10}\,I&A_{11}\,I\end{pmatrix}, \qquad\qquad I\otimes A=\begin{pmatrix}A&0\\0&A\end{pmatrix}.\]

Think of the column of four as two halves: the upper half holds the
entries whose first digit is \(0\), the lower half those whose first
digit is \(1\). Then \(A\otimes I\) \textbf{mixes the two halves with
each other}, as whole halves, in the proportions that \(A\) prescribes;
and \(I\otimes A\) \textbf{acts inside each half} and never mixes the
two. For NOT these are the two matrices we found: \(X\otimes I\) makes
the halves change places (entries \(00\leftrightarrow10\) and
\(01\leftrightarrow11\)), and \(I\otimes X\) exchanges the two entries
inside each half (\(00\leftrightarrow01\) and \(10\leftrightarrow11\)).
In both cases every entry moves. In no case is ``half the column left
alone''.

The two are different matrices unless \(A\) is a multiple of \(I\). For
if \(A\otimes I=I\otimes A\), compare blocks: the off-diagonal blocks
give \(A_{01}I=0\) and \(A_{10}I=0\), and the upper left block gives
\(A_{00}I=A\); so \(A=A_{00}I\). With a test matrix that is neither real
nor symmetric, \(A=\begin{pmatrix}1&i\\0&1\end{pmatrix}\):

\[A\otimes I=\begin{pmatrix}1&0&i&0\\0&1&0&i\\0&0&1&0\\0&0&0&1\end{pmatrix}, \qquad\qquad I\otimes A=\begin{pmatrix}1&i&0&0\\0&1&0&0\\0&0&1&i\\0&0&0&1\end{pmatrix}.\]

\textbf{In the sketchbook.} No new element; stacking does it all. A box
on one wire with a bare wire above or below it, and two boxes stacked:

\[\begin{tikzpicture}[sd]
  \draw[wire] (0,2.2) -- (5,2.2);
  \draw[wire] (0,0) -- (5,0);
  \node[sdbox, text height=1.5ex, text depth=0.5ex] at (2.5,2.2) {$A$};
  \path (0,-1) -- (0,3.2);
\end{tikzpicture}
\;=\; A\otimes I
\qquad
\begin{tikzpicture}[sd]
  \draw[wire] (0,2.2) -- (5,2.2);
  \draw[wire] (0,0) -- (5,0);
  \node[sdbox, text height=1.5ex, text depth=0.5ex] at (2.5,0) {$A$};
  \path (0,-1) -- (0,3.2);
\end{tikzpicture}
\;=\; I\otimes A\]

\[\begin{tikzpicture}[sd]
  \draw[wire] (0,2.2) -- (5,2.2);
  \draw[wire] (0,0) -- (5,0);
  \node[sdbox, text height=1.5ex, text depth=0.5ex] at (2.5,2.2) {$A$};
  \node[sdbox, text height=1.5ex, text depth=0.5ex] at (2.5,0) {$B$};
\end{tikzpicture}
\;=\; A\otimes B\]

The bare wire is the matrix \(I\) (the sixth row of Chapter 1's table,
row 6), so the first two pictures are the third with a bare wire in the
place of one box: nothing on a wire \emph{is} the box \(I\) on it. The
sketchbook's answer to §3.10's question was therefore available before
the blackboard's: draw the box, and draw the other wire. What the
sketchbook cannot do is tell us that the first picture is a matrix that
moves whole halves and the second a matrix that works inside them. The
two pictures differ by where one box sits; the two matrices differ in
nearly every entry.

\begin{quote}
\textbf{Dirac's Blackboard.}

\emph{A Kronecker product written out in full.} Two test matrices,
neither of them real, symmetric, or its own mirror image:

\[A=\begin{pmatrix}1&i\\0&1\end{pmatrix}, \qquad B=\begin{pmatrix}0&1\\i&0\end{pmatrix}.\]

The four blocks of \(A\otimes B\) are \(1\cdot B\), \(\;i\cdot B\),
\(\;0\cdot B\), \(\;1\cdot B\), and
\(iB=\begin{pmatrix}0&i\\-1&0\end{pmatrix}\):

\[A\otimes B=\begin{pmatrix}0&1&0&i\\i&0&-1&0\\0&0&0&1\\0&0&i&0\end{pmatrix}.\]

Check one entry against the formula: row \(01\), column \(10\) is
\(A_{01}B_{10}=i\cdot i=-1\). In the other order the blocks are
\(0\cdot A\), \(\;1\cdot A\), \(\;i\cdot A\), \(\;0\cdot A\):

\[B\otimes A=\begin{pmatrix}0&0&1&i\\0&0&0&1\\i&-1&0&0\\0&i&0&0\end{pmatrix}.\]

The same sixteen products \(A_{ij}B_{kl}\), in different places. The
upper left block of the first is \(B\); of the second, zero. The
Kronecker product is \emph{not} commutative, and §4.6 will say exactly
how the two orders are related.
\end{quote}

\textbf{\emph{Why this definition.}} Why blocks of \(B\) scaled by the
entries of \(A\), and not the other way about, or something more
even-handed? Because there is nothing left to choose once §4.2 is fixed.
A matrix is known by its columns; the pair of boxes must send
\(\ket{j}\otimes\ket{l}\) to \(A\ket{j}\otimes B\ket{l}\), each box
minding its own wire; and those four columns are the matrix above. The
other arrangement of blocks is \(B\otimes A\) --- the right matrix for a
reader who takes the bottom wire first. See
\texttt{Code/Ch4/OneClip.py}.

\begin{center}\rule{0.5\linewidth}{0.5pt}\end{center}

\hypertarget{which-first}{%
\section{§4.5 Which First?}\label{which-first}}

The fourth question of §3.10: \emph{if two boxes sit on two wires, does
it matter which we think of as acting first?}

On one wire, order is everything. Chapter 1 made the point with NOT and
SET-0, the box \(S_0=\begin{pmatrix}1&1\\0&0\end{pmatrix}\) that sends
both states to the state \(0\): NOT and then SET-0 is \(S_0X=S_0\),
while SET-0 and then NOT is \(XS_0=S_1\) (§1.3). Different boxes. Now
put the same two on \emph{different} wires, NOT on the top and SET-0 on
the bottom, and compute both orders. Matrices cut into \(2\times2\)
blocks multiply by the usual rule, ``along the row and down the
column'', with blocks in the places of entries, so long as the order of
the factors inside each product of blocks is kept. First SET-0, then
NOT:

\[(X\otimes I)(I\otimes S_0) \;=\; \begin{pmatrix}0&I\\I&0\end{pmatrix}\begin{pmatrix}S_0&0\\0&S_0\end{pmatrix} \;=\; \begin{pmatrix}0&S_0\\S_0&0\end{pmatrix}.\]

First NOT, then SET-0:

\[(I\otimes S_0)(X\otimes I) \;=\; \begin{pmatrix}S_0&0\\0&S_0\end{pmatrix}\begin{pmatrix}0&I\\I&0\end{pmatrix} \;=\; \begin{pmatrix}0&S_0\\S_0&0\end{pmatrix}.\]

The same matrix, and it is \(X\otimes S_0\): the two boxes acting ``at
once''. The pair that quarrelled on one wire has nothing to quarrel
about on two.

That was two boxes, each with a bare wire beside it. The real test has
four: on the top wire \(C\) and then \(A\), on the bottom wire \(D\) and
then \(B\). We may think of this as two steps taken by the pair, the
step \(C\otimes D\) and then the step \(A\otimes B\). Or we may think of
it as two separate wires, the top one carrying the box \(AC\) and the
bottom one the box \(BD\). Are they the same \(4\times4\) matrix? And if
they are, is the pairing the natural one, or must something be reversed
or reflected, as so often in Chapter 3?

\begin{quote}
\textbf{Dirac's Blackboard.}

\emph{The example in full.} With \(A\) and \(B\) as in §4.4, take also

\[C=\begin{pmatrix}1&0\\1&i\end{pmatrix}, \qquad D=\begin{pmatrix}i&2\\0&1\end{pmatrix}; \qquad C\otimes D=\begin{pmatrix}i&2&0&0\\0&1&0&0\\i&2&-1&2i\\0&1&0&i\end{pmatrix}.\]

Multiply the two \(4\times4\) matrices, sixty-four products of entries:

\[(A\otimes B)(C\otimes D)=\begin{pmatrix}0&1+i&0&-1\\-1-i&-2+2i&1&-2i\\0&1&0&i\\-1&2i&-i&-2\end{pmatrix}.\]

Now the small products:

\[AC=\begin{pmatrix}1+i&-1\\1&i\end{pmatrix},\qquad BD=\begin{pmatrix}0&1\\-1&2i\end{pmatrix}.\]

The blocks of \(AC\otimes BD\) are \((1+i)BD\), \(\;-BD\), \(\;BD\),
\(\;i\,BD\): exactly the matrix above, block for block. The cheaper
guesses fail. With the orders reversed, \(CA\otimes DB\) has the upper
left entry \((CA)_{00}(DB)_{00}=1\cdot2i\), where the true answer has
\(0\). Transposes, daggers and conjugates on the factors fail likewise.
The pairing that works is the plain one: top wire with top wire, bottom
with bottom, each in the order its own wire carries.
\end{quote}

\begin{quote}
\textbf{Theorem 4.2 (The Interchange Law).} For any four \(2\times2\)
complex matrices \(A,B,C,D\),

\[(A\otimes B)(C\otimes D) \;=\; AC\otimes BD,\]

exactly, scalar included, the Kronecker product being taken with the top
wire as the left factor. The same holds for matrices of any shapes for
which the products \(AC\) and \(BD\) are defined. In pictures: with
\(C\) and then \(A\) on the top wire, and \(D\) and then \(B\) on the
bottom wire, the picture read column by column and the picture read wire
by wire are one and the same drawing.
\end{quote}

\begin{quote}
\textbf{Corollary 4.3 (Sliding).} For all \(2\times2\) matrices \(A\)
and \(B\),

\[(A\otimes I)(I\otimes B) \;=\; A\otimes B \;=\; (I\otimes B)(A\otimes I) .\]

In pictures: boxes on different wires slide past each other.
\end{quote}

\emph{Proof idea.} An entry of \(A\otimes B\) is a product of one entry
of \(A\) and one of \(B\); an entry of a matrix product is a sum over a
middle label. For the pair, the middle label is a pair of labels, and a
sum over pairs of a product of two independent things is the product of
two sums. Nothing crosses from one factor to the other at any point, and
so nothing is reversed. The proof is carried out twice in §4.9.
\emph{Proof of the corollary.} Put \(I\) in the places of \(B\) and
\(C\) in the theorem:
\((A\otimes I)(I\otimes B)=AI\otimes IB=A\otimes B\). Put \(I\) in the
places of \(A\) and \(D\), and rename:
\((I\otimes B)(A\otimes I)=IA\otimes BI=A\otimes B\). \(\;\blacksquare\)
See \texttt{Code/Ch4/Interchange.py}, which checks the law on the
matrices above and on matrices drawn at random, checks that each of the
wrong pairings fails, and checks one case of larger shapes.

Two consequences at once. The theorem with columns in the places of
\(C\) and \(D\) is the general form of §4.4's rule,
\((A\otimes B)\bigl(\ket{\psi}\otimes\ket{\phi}\bigr)=A\ket{\psi}\otimes B\ket{\phi}\):
each factor acts on its own wire, for all columns and not only the basis
kets. And boxes that keep the promise of Chapter 3, stacked, keep it
still. The mirror image of a stack is the stack of the mirror images,
\((A\otimes B)^{\dagger}=A^{\dagger}\otimes B^{\dagger}\), since the
entry of either side in the row \(jl\) and the column \(ik\) is
\(\overline{A_{ij}}\;\overline{B_{kl}}\); so for unitary \(U\) and
\(V\),

\[(U\otimes V)^{\dagger}(U\otimes V) \;=\; U^{\dagger}U\otimes V^{\dagger}V \;=\; I\otimes I,\]

and \(I\otimes I\) is the \(4\times4\) identity.

\textbf{In the sketchbook.} The corollary is a rewrite rule, the first
that needs two wires, and we name it \textbf{\emph{sliding}}: \textbf{a
box may be moved along its own wire past any box that sits on a
different wire.}

\[\begin{tikzpicture}[sd]
  \draw[wire] (0,2.2) -- (7,2.2);
  \draw[wire] (0,0) -- (7,0);
  \node[sdbox, text height=1.5ex, text depth=0.5ex] at (4.8,2.2) {$A$};
  \node[sdbox, text height=1.5ex, text depth=0.5ex] at (2.2,0) {$B$};
\end{tikzpicture}
\;=\;
\begin{tikzpicture}[sd]
  \draw[wire] (0,2.2) -- (5,2.2);
  \draw[wire] (0,0) -- (5,0);
  \node[sdbox, text height=1.5ex, text depth=0.5ex] at (2.5,2.2) {$A$};
  \node[sdbox, text height=1.5ex, text depth=0.5ex] at (2.5,0) {$B$};
\end{tikzpicture}
\;=\;
\begin{tikzpicture}[sd]
  \draw[wire] (0,2.2) -- (7,2.2);
  \draw[wire] (0,0) -- (7,0);
  \node[sdbox, text height=1.5ex, text depth=0.5ex] at (2.2,2.2) {$A$};
  \node[sdbox, text height=1.5ex, text depth=0.5ex] at (4.8,0) {$B$};
\end{tikzpicture}\]

From left to right these are \((A\otimes I)(I\otimes B)\), then
\(A\otimes B\), then \((I\otimes B)(A\otimes I)\). And the theorem
itself is the picture with four boxes:

\[\begin{tikzpicture}[sd]
  \draw[wire] (0,2.2) -- (7,2.2);
  \draw[wire] (0,0) -- (7,0);
  \node[sdbox, text height=1.5ex, text depth=0.5ex] at (2.2,2.2) {$C$};
  \node[sdbox, text height=1.5ex, text depth=0.5ex] at (4.8,2.2) {$A$};
  \node[sdbox, text height=1.5ex, text depth=0.5ex] at (2.2,0) {$D$};
  \node[sdbox, text height=1.5ex, text depth=0.5ex] at (4.8,0) {$B$};
\end{tikzpicture}
\;\;=\;\; (A\otimes B)(C\otimes D) \;\;=\;\; AC\otimes BD .\]

Look at what the sketchbook has done here. The blackboard has two
expressions and an equation between them to prove. The sketchbook has
one drawing. We did not decide that the two readings should agree; we
never noticed, while drawing, that there were two. This is the strength
of the notation, and also the place where it must be watched, and §4.9
weighs the matter carefully.

\begin{quote}
\textbf{Penrose's Sketchbook.}

\emph{What the two directions of the page mean.} A picture on two wires
has two kinds of position, and they are not alike. \emph{Up and down} is
meaningful in full: which wire a box sits on is part of what the picture
says, and \(A\) over \(B\) is not \(B\) over \(A\). \emph{Left and
right} is meaningful only along a single wire. There it is the order of
a row, and it matters as much as it ever did:

\[\begin{tikzpicture}[sd, baseline=-0.55ex]
  \draw[wire] (0,0) -- (7,0);
  \node[sdbox, text height=1.5ex, text depth=0.5ex] at (2.1,0) {$X$};
  \node[sdbox, text height=1.5ex, text depth=0.5ex] at (4.9,0) {$S_0$};
\end{tikzpicture}
\;\;\neq\;\;
\begin{tikzpicture}[sd, baseline=-0.55ex]
  \draw[wire] (0,0) -- (7,0);
  \node[sdbox, text height=1.5ex, text depth=0.5ex] at (2.1,0) {$S_0$};
  \node[sdbox, text height=1.5ex, text depth=0.5ex] at (4.9,0) {$X$};
\end{tikzpicture}\]

Two boxes on one wire cannot pass each other; each is in the other's
way. Between \emph{different} wires, left and right mean nothing at all.
Whether \(A\) is drawn a little before \(B\), level with it, or a little
after, is an accident of the pencil. So \emph{sliding} is a rule about
what the picture never said. Use it freely between wires, and never
along one.
\end{quote}

\textbf{\emph{Why this definition.}} Why is the law stated for four
matrices, when \emph{sliding}, with two, is what one uses every day?
Because the law with four is what makes pictures on two wires well
defined. A drawing with boxes on two wires can be cut into steps in more
than one way, and the law says that all the ways give one matrix.
\emph{Sliding} is the case in which half the boxes are bare wire.

\begin{center}\rule{0.5\linewidth}{0.5pt}\end{center}

\hypertarget{the-crossing}{%
\section{§4.6 The Crossing}\label{the-crossing}}

So far the two wires keep entirely to themselves. Each carries its own
boxes in its own order, and neither knows that the other is there. There
are two ways of making them take notice of each other, and this section
has the milder: they can change places.

An example shows why we should want it. §4.4's sidebar computed
\(A\otimes B\) and \(B\otimes A\) and found two different matrices. Yet
they describe the same two boxes on the same two coins; all that differs
is which coin is drawn on top. There ought to be a matrix that does the
redrawing, that carries one description into the other. What is it?

\textbf{On the blackboard.} \textbf{SWAP} is the matrix that exchanges
the two digits of every label, \(\ket{ab}\mapsto\ket{ba}\). It fixes
\(\ket{00}\) and \(\ket{11}\) and exchanges \(\ket{01}\) with
\(\ket{10}\); its columns are therefore

\[\mathrm{SWAP} \;=\; \begin{pmatrix}1&0&0&0\\0&0&1&0\\0&1&0&0\\0&0&0&1\end{pmatrix}.\]

On a column of four it exchanges the two middle entries, and that is
exactly what exchanges the factors of a product:

\[\mathrm{SWAP}\begin{pmatrix}ac\\ad\\bc\\bd\end{pmatrix}=\begin{pmatrix}ac\\bc\\ad\\bd\end{pmatrix}=\begin{pmatrix}c\\d\end{pmatrix}\otimes\begin{pmatrix}a\\b\end{pmatrix},\]

that is,
\(\mathrm{SWAP}\bigl(\ket{\psi}\otimes\ket{\phi}\bigr)=\ket{\phi}\otimes\ket{\psi}\).

For the state of §4.2: \(\mathrm{SWAP}\) carries
\(\tfrac1{5\sqrt2}(3,\;3i,\;4,\;4i)^{T}\) to
\(\tfrac1{5\sqrt2}(3,\;4,\;3i,\;4i)^{T}\), which is
\(\ket{\omega}\otimes\ket{a}\), as computed there.

\begin{quote}
\textbf{Proposition 4.4 (the two rules of the crossing).} For all
\(2\times2\) matrices \(A\) and \(B\),

\[\mathrm{SWAP}\,(A\otimes B) \;=\; (B\otimes A)\,\mathrm{SWAP}, \qquad\qquad \mathrm{SWAP}\cdot\mathrm{SWAP} \;=\; I .\]

And \(A\otimes B=B\otimes A\) only if one of \(A\), \(B\) is a multiple
of the other.
\end{quote}

\emph{Proof idea.} A matrix is known by what it does to the four basis
kets, and each basis ket is a product, on which we know what everything
does. \emph{Proof sketch.} Apply both sides of the first equation to
\(\ket{j}\otimes\ket{l}\). On the left, \(A\otimes B\) gives
\(A\ket{j}\otimes B\ket{l}\) (§4.4) and SWAP turns this into
\(B\ket{l}\otimes A\ket{j}\). On the right, SWAP gives
\(\ket{l}\otimes\ket{j}\) and then \(B\otimes A\) gives
\(B\ket{l}\otimes A\ket{j}\). The four columns agree, so the matrices
are equal. For the second equation, exchanging the digits of a label
twice restores it. For the last statement, suppose
\(A_{ij}B_{kl}=B_{ij}A_{kl}\) for all labels. If \(B\) is zero it is a
multiple of \(A\). If not, fix \(k,l\) with \(B_{kl}\neq0\); then
\(A_{ij}=\bigl(A_{kl}/B_{kl}\bigr)B_{ij}\) for all \(i,j\), so \(A\) is
a multiple of \(B\). \(\;\blacksquare\) See
\texttt{Code/Ch4/Crossing.py}.

The first equation is the answer to our example. Multiply it on the
right by SWAP and use the second:
\(\mathrm{SWAP}\,(A\otimes B)\,\mathrm{SWAP}=B\otimes A\). The two
matrices of §4.4's sidebar are the same sixteen numbers, rearranged by
SWAP on both sides.

Could the crossing be had more cheaply, as some box \(P\) on the top
wire and some box \(Q\) on the bottom? It could not. The blocks of
\(P\otimes Q\) are the multiples \(P_{00}Q\), \(P_{01}Q\), \(P_{10}Q\),
\(P_{11}Q\) of one and the same matrix \(Q\). The upper left block of
SWAP is \(\begin{pmatrix}1&0\\0&0\end{pmatrix}\) and its lower right
block is \(\begin{pmatrix}0&0\\0&1\end{pmatrix}\). From the first,
\(P_{00}\neq0\) and \(Q\) is a multiple of
\(\begin{pmatrix}1&0\\0&0\end{pmatrix}\); but then so is \(P_{11}Q\),
and it cannot be \(\begin{pmatrix}0&0\\0&1\end{pmatrix}\). \textbf{SWAP
is not a Kronecker product of two boxes.} It is the first such matrix we
have met.

\textbf{In the sketchbook.} A new element, described before it is drawn.
A \textbf{crossing} is a place where two neighbouring wires exchange
rows: the wire that came in on the upper row leaves on the lower, and
the wire that came in on the lower row leaves on the upper. The two
lines pass through each other on the page. Neither is broken and neither
is joined to the other; each is one unbroken line from the left edge to
the right. We draw the crossing \emph{flat}: the two lines simply cross,
and the drawing does not say that either passes in front. This book has
one kind of crossing, and this is it.

\[\begin{tikzpicture}[sd]
  \draw[wire] (0,2.2) -- (1,2.2) .. controls (2.5,2.2) and (2.5,0) .. (4,0) -- (5,0);
  \draw[wire] (0,0) -- (1,0) .. controls (2.5,0) and (2.5,2.2) .. (4,2.2) -- (5,2.2);
\end{tikzpicture}
\;\;=\;\; \mathrm{SWAP}
\qquad\qquad
\begin{tikzpicture}[sd]
  \draw[wire] (0,2.2) -- (1.5,2.2) .. controls (3,2.2) and (3,0) .. (4.5,0) -- (5.5,0);
  \draw[wire] (0,0) -- (1.5,0) .. controls (3,0) and (3,2.2) .. (4.5,2.2) -- (5.5,2.2);
  \node[sdstate, text height=1.5ex, text depth=0.5ex] at (0,2.2) {$\psi$};
  \node[sdstate, text height=1.5ex, text depth=0.5ex] at (0,0) {$\phi$};
\end{tikzpicture}
\;\;=\;\;
\begin{tikzpicture}[sd]
  \draw[wire] (0,2.2) -- (3,2.2);
  \draw[wire] (0,0) -- (3,0);
  \node[sdstate, text height=1.5ex, text depth=0.5ex] at (0,2.2) {$\phi$};
  \node[sdstate, text height=1.5ex, text depth=0.5ex] at (0,0) {$\psi$};
\end{tikzpicture}\]

The crossing comes with two rules, and they are the two equations of
Proposition 4.4. The first we call \textbf{\emph{a box slides through a
crossing}}: a box pushed along its wire through a crossing stays on its
own wire, and arrives on the other \emph{row}.

\[\begin{tikzpicture}[sd]
  \draw[wire] (0,2.2) -- (4,2.2) .. controls (5.5,2.2) and (5.5,0) .. (7,0) -- (8,0);
  \draw[wire] (0,0) -- (4,0) .. controls (5.5,0) and (5.5,2.2) .. (7,2.2) -- (8,2.2);
  \node[sdbox, text height=1.5ex, text depth=0.5ex] at (2,2.2) {$A$};
  \node[sdbox, text height=1.5ex, text depth=0.5ex] at (2,0) {$B$};
\end{tikzpicture}
\;\;=\;\;
\begin{tikzpicture}[sd]
  \draw[wire] (0,2.2) -- (1,2.2) .. controls (2.5,2.2) and (2.5,0) .. (4,0) -- (8,0);
  \draw[wire] (0,0) -- (1,0) .. controls (2.5,0) and (2.5,2.2) .. (4,2.2) -- (8,2.2);
  \node[sdbox, text height=1.5ex, text depth=0.5ex] at (6,2.2) {$B$};
  \node[sdbox, text height=1.5ex, text depth=0.5ex] at (6,0) {$A$};
\end{tikzpicture}\]

On the left \(A\) is on the upper row; follow its wire through the
crossing and the same wire is the lower one on the right, where \(A\)
now sits. The second rule we call \textbf{\emph{two crossings are two
bare wires}}:

\[\begin{tikzpicture}[sd]
  \draw[wire] (0,2.2) -- (1,2.2) .. controls (2.5,2.2) and (2.5,0) .. (4,0) -- (5,0) .. controls (6.5,0) and (6.5,2.2) .. (8,2.2) -- (9,2.2);
  \draw[wire] (0,0) -- (1,0) .. controls (2.5,0) and (2.5,2.2) .. (4,2.2) -- (5,2.2) .. controls (6.5,2.2) and (6.5,0) .. (8,0) -- (9,0);
\end{tikzpicture}
\;\;=\;\;
\begin{tikzpicture}[sd]
  \draw[wire] (0,2.2) -- (5,2.2);
  \draw[wire] (0,0) -- (5,0);
\end{tikzpicture}\]

Both rules are exact equations between \(4\times4\) matrices, and both
need a word of warning.

\begin{quote}
\textbf{Penrose's Sketchbook.}

\emph{Two rules, and what the paper does not do.} It would be pleasant
to say that these two equations are facts about lines on paper, and that
anyone can see them. They are not, and one cannot. Take two real
strings, stretched side by side; exchange their right-hand ends, and
then exchange them again the same way round. The strings are now wound
about each other, and no pulling will straighten them. Paper and string
do not enforce ``two crossings are two bare wires''. We \emph{adopt} it,
as a rule about our flat crossing, and the licence for it is the
blackboard line \(\mathrm{SWAP}\cdot\mathrm{SWAP}=I\) and nothing else.
Likewise the box sliding through: the licence is
\(\mathrm{SWAP}\,(A\otimes B)=(B\otimes A)\,\mathrm{SWAP}\). A string
remembers which strand passed in front. Our drawing has no way of
recording it, and deliberately so, because the matrix SWAP has nothing
to record: it exchanges two labels, and exchanging them twice restores
them. The drawing is made for the thing that travels along the wire, and
that thing cannot tell a wound pair from a straight one.
\end{quote}

\textbf{\emph{Why this definition.}} Why a flat crossing, with no strand
in front? Because the picture should record what the algebra records and
no more, and the algebra of §4.2 has exactly one way of exchanging two
factors. A drawing that distinguished two kinds of crossing would be
promising a distinction that none of our matrices could honour.

\begin{quote}
\textbf{Feynman's Notebook.}

\emph{A crossing need not undo itself.} That our crossing is its own
undoing is a special property. Selinger's survey of graphical languages
(§3.3) first treats the general case, a \emph{braiding}: a natural
family of isomorphisms from \(A\otimes B\) to \(B\otimes A\), drawn as
wires crossing. (``Natural'' is the technical word for a condition
which, written out for our matrices, is the first equation of
Proposition 4.4: a box slides through.) The survey notes there that a
braiding followed by its inverse is the identity, but that a braiding
followed by another braiding in general is \emph{not}. Only in §3.5 does
it single out the case in which the braiding is self-inverse. It is then
called a \emph{symmetry}, and it is this that the survey draws as a
plain crossing. Its example on vector spaces with the tensor product is
the map \(x\otimes y\mapsto y\otimes x\), which is our SWAP. For that
case the survey states a coherence theorem, its Theorem 3.12, citing
Joyal and Street: an equation follows from the axioms exactly when the
two diagrams are isomorphic, with no reference, as the survey puts it,
to two- or three-dimensional structure.
\end{quote}

\begin{center}\rule{0.5\linewidth}{0.5pt}\end{center}

\hypertarget{a-box-on-two-wires}{%
\section{§4.7 A Box on Two Wires}\label{a-box-on-two-wires}}

A crossing moves the wires and does nothing to what they carry. The
stronger way for two wires to take notice of each other is for what
happens on one to \emph{depend} on the other. Here is the plainest rule
of that kind: \emph{if the upper coin shows \(1\), turn the lower coin
over; if it shows \(0\), leave the lower coin alone.} The upper coin is
not changed in either case.

\textbf{On the blackboard.} Write \(a\oplus b\) for the sum of two
digits modulo \(2\): \(0\oplus0=0\), \(\;0\oplus1=1\oplus0=1\),
\(\;1\oplus1=0\). Adding \(1\) in this sense turns a digit over and
adding \(0\) leaves it. The rule is then
\(\ket{a,b}\mapsto\ket{a,\,a\oplus b}\), or

\[\ket{00}\mapsto\ket{00}, \qquad \ket{01}\mapsto\ket{01}, \qquad \ket{10}\mapsto\ket{11}, \qquad \ket{11}\mapsto\ket{10},\]

and the matrix with these columns is called \textbf{CNOT}, the
\emph{controlled} NOT. The top wire is the \textbf{control} and the
bottom wire the \textbf{target}:

\[\mathrm{CNOT} \;=\; \begin{pmatrix}1&0&0&0\\0&1&0&0\\0&0&0&1\\0&0&1&0\end{pmatrix} \;=\; \begin{pmatrix}I&0\\0&X\end{pmatrix}.\]

The blocks say what the words said. In the upper half of the column,
where the first digit is \(0\), do nothing; in the lower half, where it
is \(1\), apply NOT. Set it beside the two matrices of §4.4 that it
resembles: \(I\otimes I=\begin{pmatrix}I&0\\0&I\end{pmatrix}\) does
nothing in both halves, and
\(I\otimes X=\begin{pmatrix}X&0\\0&X\end{pmatrix}\) applies NOT in both.
CNOT does the one in one half and the other in the other, and this is
just what a Kronecker product cannot do.

\begin{quote}
\textbf{Proposition 4.5 (CNOT is not a pair of boxes).} There are no
\(2\times2\) matrices \(P\) and \(Q\) with \(P\otimes Q=\mathrm{CNOT}\).
Moreover \(\mathrm{CNOT}\cdot\mathrm{CNOT}=I\).
\end{quote}

\emph{Proof idea.} All four blocks of a Kronecker product are multiples
of one matrix; two blocks of CNOT are not multiples of any one matrix.
\emph{Proof sketch.} Suppose \(P\otimes Q=\mathrm{CNOT}\). The upper
left blocks give \(P_{00}Q=I\), so \(P_{00}\neq0\) and \(Q=I/P_{00}\).
The lower right blocks give \(P_{11}Q=X\), that is,
\(X=\bigl(P_{11}/P_{00}\bigr)I\). But \(X\) has zeros on its diagonal
and is not zero, and no multiple of \(I\) is like that. For the second
statement, multiply by blocks:
\(\begin{pmatrix}I&0\\0&X\end{pmatrix}\begin{pmatrix}I&0\\0&X\end{pmatrix}=\begin{pmatrix}I&0\\0&XX\end{pmatrix}\),
and \(XX=I\) by Theorem 1.1. \(\;\blacksquare\) See
\texttt{Code/Ch4/Yoke.py}.

So CNOT is its own undoing, for the same reason that NOT is; it is,
after all, NOT done or not done. It is real and symmetric, hence its own
mirror image; being also its own undoing, it is undone by its mirror
image, which is the test of §3.3 for a unitary box, here with
\(4\times4\) matrices. So it may sit on a pair of qubit wires. The same
is true of SWAP.

What can such a box do that a pair of boxes cannot? Here is one thing,
and we take care to show it on ground we fully understand: covered coins
and their shares. CNOT has entries that are not negative and columns
that each sum to \(1\), so it is a stochastic matrix (§1.4), and it acts
on lists of shares. Let the upper coin be tossed --- shares
\(\begin{pmatrix}1/2\\1/2\end{pmatrix}\) --- and the lower coin be laid
heads up, \(\begin{pmatrix}1\\0\end{pmatrix}\). The two have nothing to
do with each other, and their joint list is the product. Then CNOT
exchanges its last two entries:

\[\begin{pmatrix}1/2\\1/2\end{pmatrix}\otimes\begin{pmatrix}1\\0\end{pmatrix}=\begin{pmatrix}1/2\\0\\1/2\\0\end{pmatrix} \;\;\longmapsto\;\; \begin{pmatrix}1/2\\0\\0\\1/2\end{pmatrix}.\]

Both heads, or both tails, in equal shares, and never one of each. This
list is \emph{not} the product of two lists. For suppose
\((pr,\;ps,\;qr,\;qs)=\bigl(\tfrac12,0,0,\tfrac12\bigr)\). From \(ps=0\)
either \(p=0\) or \(s=0\); but \(p=0\) would empty the first entry,
\(pr=\tfrac12\), and \(s=0\) would empty the last, \(qs=\tfrac12\).
There is nothing mysterious in it. It is the list of a matched pair of
gloves, posted in two parcels: either parcel alone is left or right in
equal shares, and to open one is to know the other. It is ordinary
correlation, and a box across two wires is how correlation is
\emph{made}. A pair of separate boxes cannot make it: \((P\otimes Q)\)
applied to a product list gives the product list
\(P\binom{p}{q}\otimes Q\binom{r}{s}\), by §4.5.

\textbf{In the sketchbook.} The chapter's last new element, described
before it is drawn. A \textbf{box that spans two wires} is a box tall
enough to cover two neighbouring rows. Two wires enter it on the left
and two leave it on the right, and it carries a label, which names a
\(4\times4\) matrix. It does not join the two wires into one, and it
does not break them: \textbf{two lines run through it}, the upper one in
at the upper left and out at the upper right, the lower one likewise.
CNOT is drawn in this way and in no other in this book until further
notice, and always with the control on top:

\[\begin{tikzpicture}[sd]
  \draw[wire] (0,2.2) -- (6,2.2);
  \draw[wire] (0,0) -- (6,0);
  \node[sdbox, minimum height=4.1em] at (3,1.1) {$\mathrm{CNOT}$};
\end{tikzpicture}
\;\;=\;\; \begin{pmatrix}I&0\\0&X\end{pmatrix}\]

The rule \(\ket{a,b}\mapsto\ket{a,a\oplus b}\) is four picture
equations, of which here are the two in which something happens:

\[\begin{tikzpicture}[sd]
  \draw[wire] (0,2.2) -- (7,2.2);
  \draw[wire] (0,0) -- (7,0);
  \node[sdstate, text height=1.5ex, text depth=0.5ex] at (0,2.2) {$1$};
  \node[sdstate, text height=1.5ex, text depth=0.5ex] at (0,0) {$0$};
  \node[sdbox, minimum height=4.1em] at (4,1.1) {$\mathrm{CNOT}$};
\end{tikzpicture}
\;=\;
\begin{tikzpicture}[sd]
  \draw[wire] (0,2.2) -- (2.6,2.2);
  \draw[wire] (0,0) -- (2.6,0);
  \node[sdstate, text height=1.5ex, text depth=0.5ex] at (0,2.2) {$1$};
  \node[sdstate, text height=1.5ex, text depth=0.5ex] at (0,0) {$1$};
\end{tikzpicture}
\qquad\quad
\begin{tikzpicture}[sd]
  \draw[wire] (0,2.2) -- (7,2.2);
  \draw[wire] (0,0) -- (7,0);
  \node[sdstate, text height=1.5ex, text depth=0.5ex] at (0,2.2) {$1$};
  \node[sdstate, text height=1.5ex, text depth=0.5ex] at (0,0) {$1$};
  \node[sdbox, minimum height=4.1em] at (4,1.1) {$\mathrm{CNOT}$};
\end{tikzpicture}
\;=\;
\begin{tikzpicture}[sd]
  \draw[wire] (0,2.2) -- (2.6,2.2);
  \draw[wire] (0,0) -- (2.6,0);
  \node[sdstate, text height=1.5ex, text depth=0.5ex] at (0,2.2) {$1$};
  \node[sdstate, text height=1.5ex, text depth=0.5ex] at (0,0) {$0$};
\end{tikzpicture}\]

With the state \(0\) on the upper wire the lower state comes out as it
went in. And \(\mathrm{CNOT}\cdot\mathrm{CNOT}=I\) is a rewrite rule,
which we call \textbf{\emph{CNOT twice is nothing}}:

\[\begin{tikzpicture}[sd]
  \draw[wire] (0,2.2) -- (11,2.2);
  \draw[wire] (0,0) -- (11,0);
  \node[sdbox, minimum height=4.1em] at (3,1.1) {$\mathrm{CNOT}$};
  \node[sdbox, minimum height=4.1em] at (8,1.1) {$\mathrm{CNOT}$};
\end{tikzpicture}
\;\;=\;\;
\begin{tikzpicture}[sd]
  \draw[wire] (0,2.2) -- (5,2.2);
  \draw[wire] (0,0) -- (5,0);
\end{tikzpicture}\]

We must be plain about the standing of these equations. They are
\textbf{lent by the blackboard}, every one. At this stage the sketchbook
has a labelled rectangle and no rule that looks inside it; it cannot
derive \emph{CNOT twice is nothing}, it can only record it. Nor does the
outline of a two-wire box tell us whether it is, or is not, a pair of
boxes in disguise. A two-wire box labelled \(A\otimes B\) is a pair of
boxes; one labelled CNOT is not, by Proposition 4.5; the rectangles look
the same. \emph{Sliding} still applies to the new box, with the natural
reading: a two-wire box cannot pass a box on either of its own two
wires, and it slides freely past anything on wires that it does not
touch. The second half of that is Theorem 4.2 for larger shapes, and
§4.10 uses it.

\textbf{\emph{Why this definition.}} Why draw CNOT as a plain labelled
rectangle, when the literature has a special symbol for it? Because the
symbol is a small picture of how the box works inside, and we do not yet
have the elements from which to build that inside honestly. Chapter 8
will have them. Until then a rectangle with a name is a truthful
statement of what we know: a \(4\times4\) matrix, two wires in, two
wires out, and a proof that it does not come apart.

\begin{center}\rule{0.5\linewidth}{0.5pt}\end{center}

\hypertarget{the-rosetta-table-3}{%
\section{§4.8 The Rosetta Table}\label{the-rosetta-table-3}}

This chapter adds ten rows to the running dictionary, numbered on from
Chapter 3's twenty-two. It is a long list because the page has acquired
a second direction, and each row is a separate element or rule. Rows 23,
28 and 29 are the three new picture elements: stacking, the box that
spans two wires, and the crossing. Rows 27, 30 and 31 are rewrite rules.

\begingroup\renewcommand{\arraystretch}{1.5}

\begin{longtable}[]{@{}
  >{\raggedright\arraybackslash}p{(\columnwidth - 4\tabcolsep) * \real{0.0909}}
  >{\raggedright\arraybackslash}p{(\columnwidth - 4\tabcolsep) * \real{0.4545}}
  >{\raggedright\arraybackslash}p{(\columnwidth - 4\tabcolsep) * \real{0.4545}}@{}}
\toprule\noalign{}
\begin{minipage}[b]{\linewidth}\raggedright
\#
\end{minipage} & \begin{minipage}[b]{\linewidth}\raggedright
Algebra (Dirac)
\end{minipage} & \begin{minipage}[b]{\linewidth}\raggedright
Picture (Penrose)
\end{minipage} \\
\midrule\noalign{}
\endhead
\bottomrule\noalign{}
\endlastfoot
23 & \(\mathbb{C}^2\otimes\mathbb{C}^2\); the kets \(\ket{ab}\), with
\(a\) first & two wires \textbf{stacked}, one below the other; the top
wire is the first --- \emph{new element} \\
24 & \(\ket{\psi}\otimes\ket{\phi}\), the Kronecker column & two states
stacked, \(\psi\) above \\
25 & \(\braket{\psi'\vert\psi}\cdot\braket{\phi'\vert\phi}\) & two
closed diagrams stacked multiply (as they do in a row, row 12) \\
26 & \(A\otimes I\), \(\;I\otimes A\), \(\;A\otimes B\) & a box on one
wire with a bare wire below it, or above it; two boxes stacked \\
27 & \((A\otimes B)(C\otimes D)=AC\otimes BD\);
\(\;(A\otimes I)(I\otimes B)=(I\otimes B)(A\otimes I)\) &
\textbf{\emph{sliding}}: boxes on different wires slide past each
other \\
28 & a \(4\times4\) matrix; CNOT, \(\ket{a,b}\mapsto\ket{a,a\oplus b}\)
& a \textbf{box that spans two wires}, labelled; two lines run through
it --- \emph{new element} \\
29 & SWAP: \(\ket{ab}\mapsto\ket{ba}\) & a \textbf{crossing}, drawn flat
--- \emph{new element} \\
30 & \(\mathrm{SWAP}\,(A\otimes B)=(B\otimes A)\,\mathrm{SWAP}\) &
\textbf{\emph{a box slides through a crossing}} and changes row (a
rule) \\
31 & \(\mathrm{SWAP}\cdot\mathrm{SWAP}=I\) & \textbf{\emph{two crossings
are two bare wires}} (a rule) \\
32 & \(n\) wires: a column of \(2^n\) entries & \(n\) lines of ink \\
\end{longtable}

\endgroup

All ten rows are exact. \emph{CNOT twice is nothing} has no row of its
own: it is a fact about one labelled box, lent by the blackboard, as
\emph{NOT on 0} was in Chapter 1.

\begin{center}\rule{0.5\linewidth}{0.5pt}\end{center}

\hypertarget{the-theorem-twice-3}{%
\section{§4.9 The Theorem, Twice}\label{the-theorem-twice-3}}

We have met the theorem through an example in which sixty-four products
of entries came out, block for block, as two small products side by side
(§4.5). Here is the statement again.

\begin{quote}
\textbf{Theorem 4.2 (The Interchange Law).} For any four \(2\times2\)
complex matrices \(A,B,C,D\),

\[(A\otimes B)(C\otimes D) \;=\; AC\otimes BD,\]

exactly, scalar included, where \(A\otimes B\) is the matrix with
entries \((A\otimes B)_{ik,\,jl}=A_{ij}B_{kl}\), the top wire being the
left factor. The same holds for matrices of any shapes for which \(AC\)
and \(BD\) are defined. In pictures: \(C\) and then \(A\) on the top
wire, \(D\) and then \(B\) on the bottom wire, read column by column or
wire by wire, is one drawing.
\end{quote}

\hypertarget{diracs-blackboard-3}{%
\subsection{Dirac's Blackboard}\label{diracs-blackboard-3}}

Everything is in the entry formula. The rows and columns of the
\(4\times4\) matrices are labelled by pairs of digits. We compute the
entry of the left side in the row \(ik\) and the column \(mn\).

\emph{Step 1: the rule for a matrix product.} The entry of a product is
the sum, over a middle label, of an entry of the left factor times an
entry of the right. Here the middle label is a pair \(jl\), and it runs
over \(00,01,10,11\):

\[\bigl((A\otimes B)(C\otimes D)\bigr)_{ik,\,mn} \;=\; \sum_{jl}\,(A\otimes B)_{ik,\,jl}\,(C\otimes D)_{jl,\,mn} .\]

\emph{Step 2: the entry formula, twice.}
\((A\otimes B)_{ik,\,jl}=A_{ij}B_{kl}\) and
\((C\otimes D)_{jl,\,mn}=C_{jm}D_{ln}\):

\[=\; \sum_{j}\sum_{l}\,A_{ij}\,B_{kl}\,C_{jm}\,D_{ln} .\]

\emph{Step 3: numbers may be reordered.} Each term is a product of four
complex numbers. Bring the two that carry \(j\) together, and the two
that carry \(l\):

\[=\; \sum_{j}\sum_{l}\,\bigl(A_{ij}C_{jm}\bigr)\bigl(B_{kl}D_{ln}\bigr) .\]

\emph{Step 4: the double sum factorizes.} The first bracket does not
contain \(l\) and the second does not contain \(j\), so the sum over all
pairs is the product of the two separate sums:

\[=\; \Bigl(\sum_{j}A_{ij}C_{jm}\Bigr)\Bigl(\sum_{l}B_{kl}D_{ln}\Bigr) .\]

\emph{Step 5: the rule for a matrix product, backwards, twice; then the
entry formula, backwards.}

\[=\; (AC)_{im}\,(BD)_{kn} \;=\; \bigl(AC\otimes BD\bigr)_{ik,\,mn} .\]

The two sides agree in every entry. Nothing in the five steps used that
the labels run over two values; with other ranges for \(i,j,k,l,m,n\)
the same lines prove the law for any shapes for which \(AC\) and \(BD\)
exist. \(\;\blacksquare\)

Observe what did \emph{not} happen. In Step 3 we reordered numbers; we
never reordered matrices. \(A\) stays to the left of \(C\), and \(B\) to
the left of \(D\), from the first line to the last, because the label
\(j\) that joins the back of \(A\) to the front of \(C\) is never
anything else. This is why there is no reversal and no dagger in the
law, where Chapter 3's \emph{Mirror of a Row} had both.

Trace one entry of the example of §4.5 through the five steps: the row
\(01\) and the column \(01\). The row \(01\) of \(A\otimes B\) is
\((i,\;0,\;-1,\;0)\) and the column \(01\) of \(C\otimes D\) is
\((2,\;1,\;2,\;1)^{T}\), so Step 1 gives \(2i+0-2+0=-2+2i\). Step 5
gives \((AC)_{00}(BD)_{11}=(1+i)(2i)=-2+2i\). They agree.

The smallest case is worth a look. Let the right-hand factors be columns
and the left-hand factors rows: \(A=\bra{\psi'}\), \(B=\bra{\phi'}\),
\(C=\ket{\psi}\), \(D=\ket{\phi}\). Then \(AC\) and \(BD\) are numbers,
the Kronecker product of two numbers is their product, and the law reads
\(\bigl(\bra{\psi'}\otimes\bra{\phi'}\bigr)\bigl(\ket{\psi}\otimes\ket{\phi}\bigr)=\braket{\psi'|\psi}\braket{\phi'|\phi}\).
That is Proposition 4.1; its proof in §4.2 was Steps 2 to 5 with no free
labels.

\hypertarget{penroses-sketchbook-3}{%
\subsection{Penrose's Sketchbook}\label{penroses-sketchbook-3}}

The sketchbook's tools are three dictionary entries. \textbf{\emph{In a
row}} (the fifth row of Chapter 1's table, row 5): two boxes in a row on
a wire, \(C\) on the left and then \(A\), are the box \(AC\).
\textbf{\emph{Stacked boxes}} (row 26): a box above a box is the
Kronecker product, the upper box on the left. And a two-wire box with a
label (row 28) stands for the matrix that its label names. Each side of
the theorem is a two-wire box; we unfold both.

\textbf{Step 1 --- the left side, unfolded by columns.} The left side is
a row of two two-wire boxes. Replace each by the stack that its label
names (\emph{stacked boxes}, used twice):

\[\begin{tikzpicture}[sd]
  \draw[wire] (0,2.2) -- (11,2.2);
  \draw[wire] (0,0) -- (11,0);
  \node[sdbox, minimum height=4.1em] at (3,1.1) {$C\otimes D$};
  \node[sdbox, minimum height=4.1em] at (8,1.1) {$A\otimes B$};
\end{tikzpicture}
\;\overset{\raisebox{0pt}[\height][1.2ex]{\scriptsize stacked boxes}}{=}\;
\begin{tikzpicture}[sd]
  \draw[wire] (0,2.2) -- (7,2.2);
  \draw[wire] (0,0) -- (7,0);
  \node[sdbox, text height=1.5ex, text depth=0.5ex] at (2.2,2.2) {$C$};
  \node[sdbox, text height=1.5ex, text depth=0.5ex] at (4.8,2.2) {$A$};
  \node[sdbox, text height=1.5ex, text depth=0.5ex] at (2.2,0) {$D$};
  \node[sdbox, text height=1.5ex, text depth=0.5ex] at (4.8,0) {$B$};
\end{tikzpicture}\]

\textbf{Step 2 --- the right side, unfolded by wires.} The right side is
one two-wire box. Replace it by the stack that its label names
(\emph{stacked boxes}), and then, on each wire separately, replace the
box by the row that its label names (\emph{in a row}, used twice):

\[\begin{tikzpicture}[sd]
  \draw[wire] (0,2.2) -- (7,2.2);
  \draw[wire] (0,0) -- (7,0);
  \node[sdbox, minimum height=4.1em] at (3.5,1.1) {$AC\otimes BD$};
\end{tikzpicture}
\;\overset{\raisebox{0pt}[\height][1.2ex]{\scriptsize stacked boxes}}{=}\;
\begin{tikzpicture}[sd]
  \draw[wire] (0,2.2) -- (6,2.2);
  \draw[wire] (0,0) -- (6,0);
  \node[sdbox, text height=1.5ex, text depth=0.5ex] at (3,2.2) {$AC$};
  \node[sdbox, text height=1.5ex, text depth=0.5ex] at (3,0) {$BD$};
\end{tikzpicture}\]

\[\overset{\raisebox{0pt}[\height][1.2ex]{\scriptsize in a row}}{=}\;
\begin{tikzpicture}[sd]
  \draw[wire] (0,2.2) -- (7,2.2);
  \draw[wire] (0,0) -- (7,0);
  \node[sdbox, text height=1.5ex, text depth=0.5ex] at (2.2,2.2) {$C$};
  \node[sdbox, text height=1.5ex, text depth=0.5ex] at (4.8,2.2) {$A$};
  \node[sdbox, text height=1.5ex, text depth=0.5ex] at (2.2,0) {$D$};
  \node[sdbox, text height=1.5ex, text depth=0.5ex] at (4.8,0) {$B$};
\end{tikzpicture}\]

\textbf{Step 3 --- compare.} The picture at the end of Step 1 and the
picture at the end of Step 2 are the same drawing: the same four boxes,
on the same wires, in the same order along each. There is nothing to
rewrite. \(\;\blacksquare\)

And \emph{sliding} follows inside the sketchbook, with no further help.
In the four-box drawing let \(B\) and \(C\) be bare wire: that is \(A\)
to the right on the top wire and \(D\) to the left on the bottom. Let
\(A\) and \(D\) be bare wire instead: the boxes have changed sides. Both
are the stack of the two boxes.

\hypertarget{the-verdict-3}{%
\subsection{The verdict}\label{the-verdict-3}}

The picture wins, and by a margin that is almost embarrassing: five
steps of indices against the remark that a drawing is equal to itself.
But the picture wins on credit, and the creditor should be named. Step 3
quietly assumed that a drawing of four boxes on two wires \emph{means
one matrix}, whichever way it is cut up to be read: into columns, as in
Step 1, or into wires, as in Step 2. Nothing about ink guarantees that.
For our matrices it is guaranteed by the blackboard proof and by nothing
else; the five steps of indices are what the freedom of the drawing
costs, paid once, here. The same is true of the crossing's two rules in
§4.6, each of which stands on its own blackboard line.

What each strand shows that the other hides: the blackboard shows
\emph{why} nothing is reversed --- the labels of the two wires never
meet --- and it shows the price, sixty-four products at a time. The
sketchbook shows that the question ``which first?'' was never a question
about the boxes. It was a question about the notation
\((\;\cdot\;)(\;\cdot\;)\), which forces a writer to put one step before
another even when the steps do not touch. The drawing is not so forced,
and so it had the law built in before we knew there was a law.

\begin{quote}
\textbf{Feynman's Notebook.}

\emph{Who says a picture may be read both ways?} The general theory
treats our theorem not as something proved but as something demanded.
Selinger's survey of graphical languages (§3.1) describes the setting in
which wires are drawn in parallel, a \emph{monoidal category}: besides
composing two maps one may form their tensor, drawn by stacking the two
diagrams. (The survey stacks from the bottom upwards; this book reads
the top wire first.) One of its axioms is that the tensor is a
\emph{bifunctor},
\((k\otimes h)\circ(g\otimes f)=(k\circ g)\otimes(h\circ f)\), which is
our interchange law word for word; vector spaces with the tensor product
are among the examples it lists. The survey derives the sliding law from
this axiom and remarks that, drawn, it holds up to deformation of the
diagram. Then comes the converse, its Theorem 3.1, which the survey
attributes to Joyal and Street: an equation follows from the axioms if
and only if its two diagrams can be deformed into one another in the
plane. It records one caveat (Caveat 3.2): the proof covers deformations
in which every wire keeps running from left to right throughout.
\end{quote}

Readers who want the pictures first and the matrices afterwards will
find that road taken, at book length, in Coecke and Kissinger's
\emph{Picturing Quantum Processes}.

\begin{center}\rule{0.5\linewidth}{0.5pt}\end{center}

\hypertarget{twice-yoked-twice-crossed}{%
\section{§4.10 Twice Yoked, Twice
Crossed}\label{twice-yoked-twice-crossed}}

Everything in the chapter can be put into one picture. Take two pairs of
wires, an upper pair and a lower pair. On the upper pair put a CNOT, and
further to the right a second CNOT. On the lower pair put a crossing,
and further to the right a second crossing, each placed between the
columns of the boxes above, so that, reading from left to right, one
meets a CNOT, a crossing, a CNOT, a crossing. We call the picture
\textbf{\emph{Twice Yoked, Twice Crossed}}:

\medskip

\begin{center}
\begin{tikzpicture}[sd]
  \draw[wire] (0,6.6) -- (18,6.6);
  \draw[wire] (0,4.4) -- (18,4.4);
  \draw[wire] (0,2.2) -- (5,2.2) .. controls (6.5,2.2) and (6.5,0) .. (8,0)
              -- (13,0) .. controls (14.5,0) and (14.5,2.2) .. (16,2.2) -- (18,2.2);
  \draw[wire] (0,0) -- (5,0) .. controls (6.5,0) and (6.5,2.2) .. (8,2.2)
              -- (13,2.2) .. controls (14.5,2.2) and (14.5,0) .. (16,0) -- (18,0);
  \node[sdbox, minimum height=4.1em] at (2.5,5.5)  {$\mathrm{CNOT}$};
  \node[sdbox, minimum height=4.1em] at (10.5,5.5) {$\mathrm{CNOT}$};
\end{tikzpicture}
\par\medskip
\end{center}

Follow the lines. The uppermost runs straight through the upper half of
both boxes: it is the control of each. The next runs through their lower
halves: the target. The third starts on its own row, crosses down, runs
along the lowest row, and crosses back. The last does the opposite. No
line is joined to another and none is broken: two lines run through each
box, and a crossing changes the row a line is on and nothing else about
it.

The picture is open, with a wire for every row entering on the left and
leaving on the right, so it stands for a matrix, and by §4.3 the matrix
is \(16\times16\): two hundred and fifty-six entries, which nobody would
wish to write out. We shall not need to.

\textbf{On the blackboard.} Write \(I_4\) for the \(4\times4\) identity,
which is two bare wires, and \(I_{16}\) for the \(16\times16\) identity.
The Kronecker product of two \(4\times4\) matrices is formed by the
recipe of §4.4 with labels that are pairs, and the upper pair of wires
is the left factor. The four columns of the picture, left to right, are
\(\mathrm{CNOT}\otimes I_4\), then \(I_4\otimes\mathrm{SWAP}\), then the
same two again; so, reading the product from right to left,

\[(I_4\otimes\mathrm{SWAP})(\mathrm{CNOT}\otimes I_4)(I_4\otimes\mathrm{SWAP})(\mathrm{CNOT}\otimes I_4) .\]

By the interchange law for \(4\times4\) factors, each neighbouring pair
of brackets is
\((I_4\cdot\mathrm{CNOT})\otimes(\mathrm{SWAP}\cdot I_4)=\mathrm{CNOT}\otimes\mathrm{SWAP}\),
and so the whole is

\[\begin{aligned}&(\mathrm{CNOT}\otimes\mathrm{SWAP})(\mathrm{CNOT}\otimes\mathrm{SWAP})\\ &\qquad=\; (\mathrm{CNOT}\cdot\mathrm{CNOT})\otimes(\mathrm{SWAP}\cdot\mathrm{SWAP}) \;=\; I_4\otimes I_4 \;=\; I_{16},\end{aligned}\]

using the law once more, then Propositions 4.5 and 4.4. Three uses of
one theorem and two small facts, in the place of a product of four
\(16\times16\) matrices.

\textbf{In the sketchbook.} First \emph{sliding}, in its form for
two-wire boxes (§4.7): the second CNOT touches neither wire of the lower
pair, so it slides to the left past the first crossing; and the first
crossing, which touches neither wire of the upper pair, may be drawn
further to the right.

\[\begin{tikzpicture}[sd]
  \draw[wire] (0,6.6) -- (18,6.6);
  \draw[wire] (0,4.4) -- (18,4.4);
  \draw[wire] (0,2.2) -- (5,2.2) .. controls (6.5,2.2) and (6.5,0) .. (8,0)
              -- (13,0) .. controls (14.5,0) and (14.5,2.2) .. (16,2.2) -- (18,2.2);
  \draw[wire] (0,0) -- (5,0) .. controls (6.5,0) and (6.5,2.2) .. (8,2.2)
              -- (13,2.2) .. controls (14.5,2.2) and (14.5,0) .. (16,0) -- (18,0);
  \node[sdbox, minimum height=4.1em] at (2.5,5.5)  {$\mathrm{CNOT}$};
  \node[sdbox, minimum height=4.1em] at (10.5,5.5) {$\mathrm{CNOT}$};
\end{tikzpicture}\]

\[\overset{\raisebox{0pt}[\height][1.2ex]{\scriptsize sliding}}{=}\;\;
\begin{tikzpicture}[sd]
  \draw[wire] (0,6.6) -- (18,6.6);
  \draw[wire] (0,4.4) -- (18,4.4);
  \draw[wire] (0,2.2) -- (10,2.2) .. controls (11.5,2.2) and (11.5,0) .. (13,0)
              .. controls (14.5,0) and (14.5,2.2) .. (16,2.2) -- (18,2.2);
  \draw[wire] (0,0) -- (10,0) .. controls (11.5,0) and (11.5,2.2) .. (13,2.2)
              .. controls (14.5,2.2) and (14.5,0) .. (16,0) -- (18,0);
  \node[sdbox, minimum height=4.1em] at (2.5,5.5)  {$\mathrm{CNOT}$};
  \node[sdbox, minimum height=4.1em] at (7,5.5) {$\mathrm{CNOT}$};
\end{tikzpicture}\]

Now the upper pair carries two CNOT boxes in a row and the lower pair
two crossings in a row. Apply \emph{CNOT twice is nothing} above:

\[\overset{\raisebox{0pt}[\height][1.2ex]{\scriptsize CNOT twice is nothing}}{=}\;\;
\begin{tikzpicture}[sd]
  \draw[wire] (0,6.6) -- (9,6.6);
  \draw[wire] (0,4.4) -- (9,4.4);
  \draw[wire] (0,2.2) -- (1,2.2) .. controls (2.5,2.2) and (2.5,0) .. (4,0)
              -- (5,0) .. controls (6.5,0) and (6.5,2.2) .. (8,2.2) -- (9,2.2);
  \draw[wire] (0,0) -- (1,0) .. controls (2.5,0) and (2.5,2.2) .. (4,2.2)
              -- (5,2.2) .. controls (6.5,2.2) and (6.5,0) .. (8,0) -- (9,0);
\end{tikzpicture}\]

and then \emph{two crossings are two bare wires} below:

\[\overset{\raisebox{0pt}[\height][1.2ex]{\scriptsize two crossings}}{=}\;\;
\begin{tikzpicture}[sd]
  \draw[wire] (0,6.6) -- (7,6.6);
  \draw[wire] (0,4.4) -- (7,4.4);
  \draw[wire] (0,2.2) -- (7,2.2);
  \draw[wire] (0,0) -- (7,0);
\end{tikzpicture}
\;\;=\;\; I_{16} .\]

\textbf{\emph{Twice Yoked, Twice Crossed} is bare wire on every row.}
Three moves. And it is right to say where each came from. \emph{Sliding}
is this chapter's theorem, in the form for larger shapes. \emph{CNOT
twice is nothing} is a fact about one matrix, lent by the blackboard.
\emph{Two crossings are two bare wires} is a rule we adopted for the
flat crossing, licensed by \(\mathrm{SWAP}\cdot\mathrm{SWAP}=I\). The
sketchbook supplied the ease; the blackboard supplied every licence.

The picture says a little more than its evaluation. Put any two-wire
boxes \(P\) and then \(Q\) in the places of the two CNOTs. The same
first move applies, the lower pair becomes two bare wires as before, and
the whole is \(QP\otimes I_4\): whatever is done on the upper pair, the
twice-crossed pair below contributes nothing to it and takes nothing
from it. See \texttt{Code/Ch4/Yoke.py}, which checks the evaluation in
both orders of the columns and checks the general statement on
\(4\times4\) matrices drawn at random.

\begin{center}\rule{0.5\linewidth}{0.5pt}\end{center}

\hypertarget{two-roads}{%
\section{§4.11 Two Roads}\label{two-roads}}

Let us gather what we have. The second wire is drawn below the first,
and the top wire is mentioned first (§4.1). The column for two coins has
an entry for each pair of answers, and for separate coins each entry is
a product; brackets of products multiply (§4.2). Each further wire
doubles the column while adding one line to the drawing, and the
drawing's brevity is not the calculation's (§4.3). A box on one wire of
two is \(A\otimes I\) or \(I\otimes A\), and it moves every entry of the
column (§4.4). Boxes on different wires have no order between them; on
one wire they still have (§4.5). Wires can cross, and the crossing has
two rules, each standing on a line of algebra (§4.6). A box can span two
wires, and CNOT, which does, is not a pair of boxes; from a tossed coin
and a plain one it makes a matched pair (§4.7). The interchange law is
five steps of indices on the blackboard and a single drawing in the
sketchbook, and the drawing's freedom is paid for by the indices (§4.9).
Two CNOTs above two crossings are bare wire on every row (§4.10).

It comes to one sentence, and it answers the last of Chapter 3's
questions. \textbf{Side by side there is no first and no second: what is
drawn on different wires happens in no order at all, and the only order
a picture has is the order along each wire.}

We have spent the chapter looking \emph{across} the wires. Now look once
more \emph{along} one. The box \(H\) sends the state \(0\) to the state
\(+\) (§3.4), which answers the standard question either way with equal
chances: one start, headed towards both answers. Put a second \(H\)
after the first. From the start \(0\) to the end \(0\) there are now two
roads, one through the answer \(0\) in the middle and one through the
answer \(1\); the rule for a matrix product says as much, for an entry
of \(HH\) is a sum of two terms, one for each middle label. We know the
total already, since \(HH=I\): the state \(0\) comes home with
certainty, and never arrives at \(1\). But a chance cannot be taken away
by adding another chance to it. What, then, is being added along the two
roads? Do the amplitudes add --- and can they cancel?

That is the subject of Chapter 5. It sets a box between the two \(H\)'s,
to mark one of the roads, and watches the arrivals change: the
arrangement is called an interferometer. With it come a small problem of
Deutsch's about a function one is allowed to consult only once, and a
curious effect called phase kickback, in which a box meant to act on one
wire leaves its mark on the other. We shall need everything in this
chapter to say what that means. The wires are laid side by side; the
next chapter sends something down two roads at once.

\hypertarget{two-path-invention}{%
\chapter{Two-Path Invention}\label{two-path-invention}}

\begin{center}\rule{0.5\linewidth}{0.5pt}\end{center}

\emph{春の夜や}\\
\emph{輪と輪重なり}\\
\emph{月の凪}

\emph{haru no yo ya}\\
\emph{wa to wa kasanari}\\
\emph{tsuki no nagi}

\emph{spring night ---}\\
\emph{two rings cross on the water,}\\
\emph{the moon lies still}

\begin{center}\rule{0.5\linewidth}{0.5pt}\end{center}

\begin{center}\rule{0.5\linewidth}{0.5pt}\end{center}

\emph{{[}Spring dusk on the Tortoise's porch, at the edge of the pond. A
cat is already there when Achilles comes up the steps --- stretched
along the railing, or most of it is; the end that ought to be tail is
more a suggestion than a tail. It is neither quite asleep nor quite
awake.{]}}

\emph{{[}The porch gives straight onto the water, and the water gives
back the last of the sky. It is the first evening warm enough to sit
out; the pond is glass, not a breath on it yet, and along the near bank
the frogs are only clearing their throats. On the table a lamp is lit
against the coming dark, a pot of tea beside it, two chairs drawn up.
Through the open door behind, the room keeps its own small light: on the
far wall, level with the lamp, hangs a little framed print of a still
pond under a low moon --- Escher's \textup{Rippled Surface} --- black
water crossed by broken rings of white, catching a dull shine off the
glass. The cat has settled into the one warm patch the day has left
along the railing, or the patch is only where the sun was; from the
steps you could not have said which.{]}}

\textbf{Achilles:} Tortoise! I walked the long way round, by the water.
I'd forgotten how loud the frogs get the moment the light goes.

\textbf{Tortoise:} They rehearse all week for one evening and then
overplay it. Sit --- the tea's still hot. You've timed it well; the moon
isn't up yet, but it will be, and from this chair it comes up straight
out of the pond.

\textbf{Achilles:} I came out hungry for a puzzle. And --- good evening.
I didn't see the cat there. Have you been sitting long?

\textbf{Cat:} \emph{{[}without troubling to open both eyes{]}} Coming
and going. Mostly the one, a little of the other.

\textbf{Achilles:} That is no sort of answer.

\textbf{Tortoise:} It's the only sort you'll get to a question shaped so
loosely. Ask it something baggy and it hands you something baggy back.
Drink your tea. You'll grow used to it.

\emph{{[}Achilles takes the offered cup and folds both hands round it;
the warmth is welcome, the evening cooler than it looked from the road.
Out over the water the frogs have found their nerve at last, two of them
now trading a phrase back and forth across the dark.{]}}

\begin{center}\rule{0.5\linewidth}{0.5pt}\end{center}

\hypertarget{first-challenge-two-roads}{%
\section{\texorpdfstring{First Challenge --- \emph{Two
Roads}}{First Challenge --- Two Roads}}\label{first-challenge-two-roads}}

\textbf{Achilles:} Here is what's been nagging me since we last talked.
You showed me a single start sent toward \emph{both} answers at once,
and a second gate set after it, and all of a sudden there were two roads
leading from the one start back to the very same end. Two roads to a
town are two chances of reaching it; you don't arrive \emph{less} often
for being given a choice. So the numbers on the two roads simply add up,
surely. Do the amplitudes add --- and can they cancel?

\textbf{Tortoise:} Let's not decide it by weighing the word
``amplitude.'' Let's walk it.

\emph{{[}She unrolls a folding board the length of the table.
Brass-lined channels run down it, side by side: the willow lane nearest
the pond, then alder, then hazel, then rowan --- and, set a little apart
from the rest and cut deeper, the reed lane at the far edge.{]}}

\emph{{[}The lamp finds the board and runs the length of it, a warm
gleam down each brass channel: the willow nearest the water, the reed
farthest off and set a little apart, with alder and hazel and rowan
ranked between. There is no moon yet to reach them --- only the lamp,
and the dark settling over the pond beyond.{]}}

\textbf{Tortoise:} This is a conduit-board. Each brass channel is a
lane, and a walker sent down a lane runs it straight from one end to the
other --- a lane never forks, never splits into two. Tonight we want
only this nearest one, the willow; the others can lie idle a while.

\textbf{Tortoise:} When a walker reaches the end of a lane it drops a
little disc --- a token --- into the lane's arrival pocket, to record
that it came. A token has two faces, and shows one at a time: silver up,
or slate up. And each lane ends not in one pocket but two --- a
\emph{home} pocket, for a walker that has come back to the value it
started from, and an \emph{away} pocket, for one turned aside to the
other value.

\textbf{Tortoise:} At the head of the willow lane, and again at its
foot, stands a little turnstile --- the workshop calls it a spreader. It
won't let a walker home until the lane has been walked once for
\emph{each} value the lane can carry. That is how one walker comes to
owe two roads, without the channel ever splitting.

\textbf{Achilles:} Two roads from the one walker. Send it, then --- down
the low road only, and I'll watch that.

\textbf{Cat:} \emph{{[}uncurling onto the head of the board{]}} I don't
take the low road only. I take both, or I take none.

\textbf{Tortoise:} You may ask it, afterwards, which road it took, if
you care to. You won't enjoy the answer. Watch it walk.

\emph{{[}The Cat walks the willow lane. At the head spreader it goes two
ways at once, as though it were two cats of one substance; down the lane
the two roads run together, and each lets fall its token at the foot.
Both come up silver. The foot spreader deals them out into the two
pockets: into the home pocket it lays both just as they fell --- silver
upon silver, and they pile; into the away pocket it lays them too, but
turns the high road's token over as it sets it down --- silver beside
slate, unlike faces, and the pair lifts out together and leaves that
pocket bare. Home heaped, away empty. The walker has come home.{]}}

\emph{{[}Out on the water the two frogs come down on the very same beat,
and the doubled note carries clear across the pond and back off the far
bank, fuller than either could manage alone. Achilles grins at it
despite himself.{]}}

\textbf{Achilles:} There --- the home pocket full, the away pocket bare.
The two roads met at home, both silver, and piled up. The walker comes
home every time. I \emph{said} the chances only add.

\textbf{Tortoise:} Watch once more, with something in the lane.

\emph{{[}She takes a small pin from her shell.{]}}

\textbf{Tortoise:} This is a phase-peg. I press it down into the willow
lane itself --- into the brass, midway between the two spreaders --- and
it does one single thing: it turns the face of any token that passes it,
but only along the road that carries the \emph{high} value. The low road
runs past it untouched.

\emph{{[}She presses the peg into the willow lane. The lane is otherwise
just as it was; only now a single pin stands upright in the brass,
halfway down.{]}}

\emph{{[}The Cat walks again. Two roads again from the one head
spreader; but this time, as the high road passes the peg, its token
turns over --- silver going down, slate coming up --- while the low road
slips past the pin and stays silver. At the foot, the spreader deals as
before: into home, both as they fell --- silver and slate now, unlike
faces, and the pair lifts out and leaves the home pocket bare; into
away, the high road's token turned again as it is set down --- slate
back to silver, silver beside silver, and these pile. Home bare, away
heaped. The same walk, the same lane; one pin between the spreaders, and
the walker is turned away.{]}}

\emph{{[}The two frogs are still going, but they have drifted out of
step --- one filling exactly the gap the other leaves, on and on, so
evenly that the porch seems to fall quiet though neither has stopped for
a moment. Achilles tilts his head, listening for a silence that is not
quite there.{]}}

\textbf{Achilles:} But nothing touched the roads except a pin that they
--- no. The low road went by it untouched.

\textbf{Tortoise:} Read me the low road.

\textbf{Achilles:} The low road came down silver, and lay silver in the
home pocket, exactly as it always has.

\textbf{Tortoise:} And the high road came down turned by the peg, and
lay in the home pocket slate against your silver --- and where a moment
ago the two piled up, now they lift out together and leave nothing
behind.

\textbf{Achilles:} Two roads into one pocket, one silver, one slate.
\ldots Oh. They didn't add. They \emph{cancelled}.

\emph{{[}The Tortoise looks up from the board to the pond. The moon has
begun to clear the water, and two rings of ripple spreading from out in
the dark cross one another --- standing high and bright where they ride
up together, flattening away to nothing where crest meets trough ---
while the moon's reflection breaks along the seams and closes and breaks
again.{]}}

\textbf{Tortoise:} \emph{Rippled Surface}. \emph{{[}She turns back to
the board.{]}}

\emph{{[}Through the open door, the first of the risen moon has reached
the room and lies now on the glass over the little frame on the far
wall.{]}}

\textbf{Achilles:} But the walker went both ways at once. Where did the
other road go --- can you draw the both-at-once?

\begin{center}\rule{0.5\linewidth}{0.5pt}\end{center}

\hypertarget{second-challenge-a-sum-of-pictures}{%
\section{\texorpdfstring{Second Challenge --- \emph{A Sum of
Pictures}}{Second Challenge --- A Sum of Pictures}}\label{second-challenge-a-sum-of-pictures}}

\emph{{[}The Tortoise draws the peg up out of the willow lane; the pin
comes free, and the brass lies unbroken again, the lane bare from
spreader to spreader.{]}}

\textbf{Tortoise:} We're done turning tokens for the moment --- the
willow runs clean again. Now: draw me your both-at-once.

\emph{{[}The moon is properly up now, clear of the water and climbing,
and its light has just reached the near edge of the board --- laid first
along the willow lane, the one nearest the pond, while the rest of the
board keeps to the lamp. The Tortoise sets a fresh pot between them and
fills both cups; Achilles wraps his hands round his and leans in over
the paper.{]}}

\textbf{Achilles:} Gladly. The walker comes in here, and then---
\emph{{[}he begins to sketch, on a scrap of paper, a lane dividing into
two{]}} ---the lane divides, one branch high, one branch low, and---

\textbf{Tortoise:} Stop your pencil there. A lane that splits into two
is not a lane; it's a different creature altogether, one road in and two
out, and we own no such thing on this board --- nor shall for a good
while yet. A walker never comes to a fork.

\emph{{[}His pencil stops where she stopped it, over the half-drawn
fork. His hand throws two shadows across the paper at once --- one from
the lamp at his elbow, one from the moon now clear over the water ---
two whole hands, each cast entire, crossing and darkening where they
fall together and nowhere splitting in two.{]}}

\emph{{[}She lays two flat printed cards on the table.{]}}

\textbf{Tortoise:} These are one-road tiles. Each is a still picture of
a single, whole road with a definite value --- this one the low road,
this one the high --- with a face printed on it, silver or slate, for
the sign it carries. They're only cards: they lay no brass, they add no
lane, and nobody walks them. But laid together in a pocket, they say
what the walk says.

\emph{{[}She sets the two tiles into the willow's home pocket, one over
the other. For the plain walk both are printed silver: silver over
silver, and the pocket reads one doubled silver --- a pile. Then she
lifts the high-road tile, turns it to its slate face, and lays it back:
silver beneath slate, unlike faces, and the pair lifts out as one and
the pocket reads bare. Nothing was walked; two still cards, summed, gave
the same full pocket and the same empty one as the Cat's two roads
did.{]}}

\textbf{Achilles:} So ``both at once'' was never a fork in the road.
It's two whole roads, each with its own face, laid into the one pocket
and added. And if I set a second spreader after the first, then every
road through the pair is one thing: in at the start, through \emph{one}
middle value, out at the end --- and I add up those roads over every
middle value they might pass through.

\textbf{Tortoise:} That is the whole of it. A walk is a sum of roads;
each road a run of definite values, each carrying its face.

\emph{{[}The Tortoise sits back a little, satisfied, and lets him turn
it over. Along the railing behind them the warm patch is still warm, or
warm again; whether the cat is stretched there or up on the board's
head, the lamp and the moon between them make it hard to say. Beyond the
rail a frog plops off into the water and sets a ring going that widens
and thins and is gone.{]}}

\textbf{Achilles:} So every road carries a sign. What if I could make a
road's sign depend on a \emph{question} --- on some rule I couldn't see
the inside of --- so that walking the road \emph{asked} the question,
and the road came home marked by the answer?

\begin{center}\rule{0.5\linewidth}{0.5pt}\end{center}

\hypertarget{third-challenge-the-mark-on-the-other-wire}{%
\section{\texorpdfstring{Third Challenge --- \emph{The Mark on the Other
Wire}}{Third Challenge --- The Mark on the Other Wire}}\label{third-challenge-the-mark-on-the-other-wire}}

\textbf{Tortoise:} For a question you need something to ask, and
something to ask it \emph{of}. Bring the reed lane into it --- the deep
one at the far edge, past the alder and the hazel and the rowan that
have lain idle all evening. The willow is the road you ask \emph{from};
the reed is the road you ask \emph{about}.

\emph{{[}The moon has climbed enough to lay its light further across the
board now --- past the willow and onto the alder and the hazel, so that
only the rowan and the deep reed still lie in the lamp's yellow. The
reed, set apart and cut deeper than the rest, holds the shadow
longest.{]}}

\emph{{[}She lifts a long brass beam and lays it square across all the
lanes at once, from willow to reed.{]}}

\textbf{Tortoise:} This is the bar. It is no lane --- it lies
\emph{across} every lane, a solid beam, not a channel of its own, so
you'll never mistake it for a road. It carries a rule hidden inside it:
where a road runs under it at the high value the bar acts, and where a
road runs the low value the bar does nothing. To ask it a plain question
I want a plain, single road --- so up comes the willow's head spreader
for a moment.

\emph{{[}She lifts the head spreader off the willow lane. The turnstile
comes away and the lane lies open at its head; now the willow owes not
two roads but one --- the single road she chooses to walk by hand. She
drops a token into the reed's pocket, silver up and lying flat, and
walks a token down the willow's high road herself. As the hand-walked
token passes beneath the bar, the bar reaches across and turns the reed
token --- silver going down, slate coming up --- while the willow token
runs on and comes home unmarked, silver as it set out. The mark has
landed on the reed, the lane asked \emph{about}, exactly where Achilles
would have put it.{]}}

\textbf{Achilles:} Just as I'd expect. The bar wrote its mark on the
reed --- the lane I asked about --- and left my willow token quite
alone.

\textbf{Tortoise:} Now set the reed's token differently. Not lying flat
with one face up, but wedged upright on its edge, so it shows silver to
one side and slate to the other, both at once. Two-faced. That is the
difference setting.

\emph{{[}She stands the reed's token up on its edge in the deep lane,
wedged there, silver to one hand and slate to the other, and walks a
token down the willow's high road once more. The bar reaches across to
turn the reed token --- but a two-faced token turned end for end shows
the very same two faces, and stands entirely unchanged. The turn has
nothing there to spend itself on. So it jumps: back along the bar to the
willow road that called it, and turns \emph{that} token instead ---
silver to slate --- while the reed token stands exactly as it was
wedged.{]}}

\emph{{[}Out on the pond a frog launches itself at the moon's broken
reflection and comes down short in a smack of water. The moon on the
water does not so much as flinch; but the ring the frog threw runs out
into the dark, meets the bank, and comes rocking back to lift the very
pad it leapt from.{]}}

\textbf{Achilles:} \emph{{[}staring{]}} You laid the bar across the
\emph{reed}. The reed didn't change. And my willow token turned over.
The mark lands on the road I asked from, not the road I asked about.

\textbf{Tortoise:} And a walker carries that turned token home into its
pocket, where it can meet another and pile or cancel.

\textbf{Achilles:} So one walk under the bar comes home with the rule's
mark on my own road. Could one question's mark tell me something about
the rule \emph{as a whole} --- not this answer or that one, but the
whole of it at once?

\begin{center}\rule{0.5\linewidth}{0.5pt}\end{center}

\hypertarget{fourth-challenge-the-one-question}{%
\section{\texorpdfstring{Fourth Challenge --- \emph{The One
Question}}{Fourth Challenge --- The One Question}}\label{fourth-challenge-the-one-question}}

\textbf{Tortoise:} \emph{{[}setting the willow's head spreader back in
place, so the lane owes both its roads again{]}} Here is my wager. You
may consult the bar exactly \emph{once} --- one walk of the board
beneath it --- and from that one walk you must tell me whether the
hidden rule gives the willow the \emph{same} answer on both its roads,
or \emph{opposite} answers. Same, or different. Not the rule you've just
watched, mind --- a fresh one, hidden.

\emph{{[}The moon is high enough now to have found the rowan as well;
the greater part of the board lies in its light, and only the reed at
the far edge still keeps to the lamp, deep and apart. The frogs have
thinned to a pair again, trading their phrase across the water.{]}}

\emph{{[}She sets the bar to a new rule and lays her hand flat over its
face, so the rule cannot be read.{]}}

\textbf{Achilles:} The plain way first. Up comes the head spreader, and
I walk one road by hand --- the low one --- beneath your hidden bar.

\emph{{[}He lifts the head spreader, walks a single token down the
willow's low road under the covered bar, and reads the face it brings
home.{]}}

\textbf{Achilles:} One answer, there in the pocket. Now to
\emph{compare}, I must walk the high road as well and read that too---

\textbf{Tortoise:} ---a second walk. Two, the plodding way; and you have
spent your one already, on a single answer.

\textbf{Achilles:} That's no fair test, then.

\textbf{Tortoise:} It is a game, not an execution. \emph{{[}She resets
the bar beneath her hand.{]}} A fresh rule. Try again --- better.

\emph{{[}She says it lightly, and there is tea still warm in the pot;
Achilles laughs under his breath, half at himself, and squares up to the
board again.{]}}

\textbf{Achilles:} Then I won't plod. \emph{{[}He sets the head spreader
back, so both roads open again.{]}} Let the Cat take both roads in the
one walk, and I'll read just the road I want out of it. \emph{{[}to the
Cat{]}} Which road brought the mark home?

\textbf{Cat:} \emph{{[}eyes drifting half-shut{]}} The one, and the
other, and a little of the way between.

\emph{{[}The two frogs are at it still, and it is the oddest thing ---
you can hear plainly whether they are keeping step or drifting apart,
the whole shape of it clear across the water, and yet you could not say
for your life which of the two it was you last heard.{]}}

\textbf{Achilles:} That tells me nothing whatever. Then I'll lift off
the \emph{foot} spreader --- let the two roads simply \emph{arrive},
unsorted, and read whichever token lands.

\emph{{[}He lifts the foot spreader clear of the willow. The Cat walks
under the covered bar; the two roads come down and drop their tokens
into the pocket with nothing to sort them --- silver, then slate, then
silver, no telling which walk gave which.{]}}

\emph{{[}A moth has found the lamp and beats at the hot glass, struck
away and falling and climbing back, and there is no calling in advance
which side of it will catch the light as it turns.{]}}

\textbf{Achilles:} \ldots a coin. Perfectly fair and perfectly empty.
Reading one raw road tells me nothing about the rule either.

\textbf{Tortoise:} \emph{{[}resetting the bar once more{]}} Another
fresh rule, then. Now stop trying to read a single road. Put the foot
spreader back, leave the reed wedged two-faced, and simply walk the
whole thing once --- head spreader, bar, foot spreader --- and read the
pockets.

\emph{{[}He sets the foot spreader back. The Cat walks. At the head
spreader both roads open; each runs under the hidden bar and comes back
turned by the reed's two-faced kick --- the low road carrying the rule's
answer for the low value, the high road its answer for the high value.
At the foot the spreader brings the two together. This rule has given
them opposite faces: silver and slate meet in the home pocket and lift
out as a pair, leaving it bare, while the away pocket fills. The walker
is turned away.{]}}

\textbf{Achilles:} Turned away --- the home pocket empty. So the two
answers were \emph{opposite}, the roads cancelled, and I never once had
to read either on its own. \emph{{[}a pause{]}} And yet --- which was
the low answer, silver or slate? I couldn't tell you. I only know the
two disagreed.

\textbf{Achilles:} \ldots Oh. I never learned either value, only whether
they agree. One walk, one consultation of the bar, and the pocket told
me same-or-different --- the very thing that no single answer holds by
itself.

\textbf{Tortoise:} Count what a plodding walker needs. To learn
same-or-different by \emph{reading answers}, he must read them both ---
two consultations, as you found when you spent one on a single road. You
did it in one.

\textbf{Tortoise:} One walk has settled a two-answer question that costs
two answers to settle the plodding way. Now stretch the question over
more roads --- and let us see how far the cancelling reaches.

\begin{center}\rule{0.5\linewidth}{0.5pt}\end{center}

\hypertarget{fifth-challenge-the-sixteen-roads}{%
\section{\texorpdfstring{Fifth Challenge --- \emph{The Sixteen
Roads}}{Fifth Challenge --- The Sixteen Roads}}\label{fifth-challenge-the-sixteen-roads}}

\textbf{Tortoise:} Bring the others in beside the willow --- the alder,
the hazel, the rowan --- each one a query road, each with its own head
and foot spreader; and the reed stays wedged two-faced at the edge, as
before. Now one walk owes a road for every choice of values across the
whole board.

\emph{{[}She opens a small tally-book beside the board.{]}}

\textbf{Tortoise:} The road-book --- a tally of the roads a single walk
owes. We've had no need of it till now. Count them with me as they fall.

\textbf{Achilles:} Two roads on the willow alone --- high and low. Bring
in the alder, and each of those two doubles: four. The hazel doubles
them again: eight. The rowan doubles them once more: sixteen. Sixteen
roads owed in the one walk.

\textbf{Achilles:} And for once I can tell you what will happen before
the Cat sets out. If the bar gives the \emph{same} answer down every one
of the sixteen roads, all sixteen tokens come home the same face and
pile up --- the walker comes home. But if the bar \emph{splits} them ---
half one answer, half the other --- then eight tokens meet eight of the
opposite face in the home pocket, and lift out pair by pair, and leave
nothing at all.

\textbf{Tortoise:} Then here is a rule, and I shan't tell you which it
is. \emph{{[}She sets the bar to a hidden rule and lays her hand over
its face.{]}} Walk it, and read the home pocket alone.

\emph{{[}The moon stands high over the pond at last, and its light has
crept the whole width of the board and down into the reed lane too, the
deep one at the edge --- so that every channel lies lit together now,
lamp and moon at once along the brass. Out on the water a dozen frogs
are going, and the rings they throw run out and cross and cross again,
standing up bright where they climb together and going flat to nothing
where they meet against the grain.{]}}

\emph{{[}The Cat walks --- sixteen roads at once, the whole board alight
from head to foot. At the foot, read the home pocket alone: eight tokens
come home silver and eight come home slate; face meets opposite face,
and pair after pair lifts out, until the home pocket stands bare. The
road-book's tally reads sixteen roads walked; the home pocket holds none
of them.{]}}

\textbf{Achilles:} Empty --- not a token left at home. Eight cancelled
eight, so the rule \emph{split} them, and I read that straight off the
bare pocket without ever seeing the bar.

\textbf{Tortoise:} There it is. \emph{{[}She sits back. The moon is well
up now, laid out whole across the pond, and the frogs have quieted to
the odd remark in the dark.{]}} Well walked, Achilles. You saw that one
coming.

\textbf{Achilles:} I did, didn't I. \emph{{[}He watches the empty board
a moment, then turns to the Cat.{]}}

\textbf{Achilles:} \emph{{[}quietly{]}} All evening I've asked you loose
things and had loose answers back. Let me ask you one straight one, now
the board has settled it. The roads you just walked --- are the roads
all of one kind, or split? Same, or different?

\textbf{Cat:} Different --- the roads cancel.

\emph{{[}The Cat is awake now, or was; the railing is empty, the lanes
are empty, and the last of it, when everything else of it has faded from
the porch, is a grin above the emptied board.{]}}

\emph{{[}The pond lies whole and quiet under the risen moon, the frogs
gone down to the odd remark in the dark, and the brass channels of the
board keep the light along them --- willow, alder, hazel, rowan, and the
deep reed at the edge --- lamp and moon together, going cool now that no
one walks them. In the room through the open door the little frame on
the far wall has gone dark, a plain rectangle back in the wall where it
always hung. The tea has gone cold in the pot, the lamp burns steady,
and over the water the two ring-systems have long since run out and the
surface has closed over where they crossed.{]}}

\hypertarget{chapter-5-roads-that-cancel}{%
\chapter{Chapter 5: Roads That
Cancel}\label{chapter-5-roads-that-cancel}}

\emph{--- Two-Path Invention ---}

\emph{Part II: Wires Together}

\emph{Escher: ``Rippled Surface'' --- a still pond holding the moon and
the bare trees, and two sets of rings spreading across the reflection,
crossing}

\begin{center}\rule{0.5\linewidth}{0.5pt}\end{center}

If a walker can take two roads to the same place at once, can the two
roads erase each other so that the walker never arrives?

Nothing in ordinary life prepares us for a yes. Two roads to a town are
two chances of reaching it; block one and the other still serves; open
both and you are more likely to arrive, not less. Chances add, and
adding makes things larger. A traveller who could take both roads at
once would, if anything, be doubly sure of the town.

Chapter 4 left a wire in a strange condition. A single start, sent
through the gate \(H\), was pointed at \emph{both} answers at once; and
if a second \(H\) came after it, there were now two ways to get from the
start \(\ket0\) back to the end \(\ket0\) --- one through the middle
value \(\ket0\), one through the middle value \(\ket1\). Two roads, one
start, one finish. The question we shirked there is the question of this
chapter. When the walker travels both roads, what reaches the far end is
not two chances but two \emph{amplitudes}, and amplitudes are not
chances. They are numbers that may be positive or negative, and two
numbers of opposite sign do not pile up. They cancel.

This is the whole of the chapter, and it is more than a curiosity. If a
road's sign could be made to depend on some unknown function \(f\) ---
if travelling the road \emph{asked a question} and came back marked
\(+\) or \(-\) by the answer --- then two roads carrying two answers
would, on meeting, report not either answer but their \emph{relation}:
whether they agree or differ. That is a fact about \(f\) as a whole,
wrung from a single walk. By the end of the chapter one walk, past one
oracle, will settle with certainty a question that any road-by-road
interrogation needs two answers to decide; and a longer board, sixteen
roads wide, will let eight roads cancel eight and leave nothing at all.
The algebra will be three short lines. We shall draw it too --- every
road, every sign --- and be honest at the close about which of the two
strands did the proving and which only made the cancellation something
you could see.

\begin{center}\rule{0.5\linewidth}{0.5pt}\end{center}

\hypertarget{two-roads-1}{%
\section{§5.1 Two Roads}\label{two-roads-1}}

Return to the picture Chapter 4 left open. The gate \(H\) is the
labelled box of Chapter 3, the one that sends a definite value to an
even mixture:

\[H\ket0=\tfrac1{\sqrt2}\bigl(\ket0+\ket1\bigr),\qquad H\ket1=\tfrac1{\sqrt2}\bigl(\ket0-\ket1\bigr),\]

and it is its own undoing, \(HH=I\) --- a fact we proved in Chapter 3
and shall now \emph{see}. Put two \(H\)'s in a row and feed in
\(\ket0\).

\textbf{On the blackboard.} Take it in two steps, naming the middle
value each time.

\[\ket0 \;\xrightarrow{\;H\;}\; \tfrac1{\sqrt2}\bigl(\ket0+\ket1\bigr) \;\xrightarrow{\;H\;}\; \tfrac1{\sqrt2}\bigl(H\ket0+H\ket1\bigr) \;=\; \tfrac12\bigl(\ket0+\ket1\bigr)+\tfrac12\bigl(\ket0-\ket1\bigr).\]

Now read the amplitude of each ending. The end \(\ket0\) is reached by
two roads --- through the middle value \(\ket0\) (amplitude
\(\tfrac12\)) and through the middle value \(\ket1\) (amplitude
\(\tfrac12\)) --- and the two have the \emph{same} sign, so they pile up
to \(1\). The end \(\ket1\) is also reached by two roads, of amplitudes
\(+\tfrac12\) and \(-\tfrac12\), and these have \emph{opposite} sign, so
they cancel to \(0\). The walker comes home to \(\ket0\) with certainty.
Two \(H\)'s undo each other, and the reason is a cancellation.

Now insert, between the two \(H\)'s, a half-turn of phase. This is the
gate \(Z\) of Chapter 3: it leaves \(\ket0\) alone and sends
\(\ket1\to-\ket1\), so it flips the sign of the \(\ket1\)-value road and
nothing else. The middle line above becomes

\[\tfrac1{\sqrt2}\bigl(H\ket0-H\ket1\bigr)=\tfrac12\bigl(\ket0+\ket1\bigr)-\tfrac12\bigl(\ket0-\ket1\bigr)=\ket1.\]

The signs have swapped fortunes: now the two roads to \(\ket0\) cancel
and the two roads to \(\ket1\) pile up. The walker is turned away. In
one matrix,

\[HZH=X,\]

the bit-flip of Chapter 3 --- a wire that carries every value straight
across to the \emph{other} value. A half-turn of phase between two
mirrors is exactly a flip. This is the Mach--Zehnder interferometer,
written in gates: two half-silvered mirrors (\(H\) and \(H\)) with a
phase in one arm, and whether the particle leaves by the near door or
the far door is decided entirely by whether the two internal roads agree
or disagree in sign.

\textbf{In the sketchbook.} Two boxes on one wire; a state at the start,
the wire running out. With no phase between the mirrors, the picture is
a wire with nothing on it:

\[\begin{tikzpicture}[sd]
  \draw[wire] (0,0) -- (6,0);
  \node[sdstate] at (0,0) {$0$};
  \node[sdbox] at (2,0) {$H$};
  \node[sdbox] at (4,0) {$H$};
\end{tikzpicture}
\;=\;
\begin{tikzpicture}[sd]
  \draw[wire] (0,0) -- (3,0);
  \node[sdstate] at (0,0) {$0$};
\end{tikzpicture}\]

With the half-turn \(Z\) between them, the two mirrors and the phase are
a single bit-flip, and the start \(\ket0\) leaves as \(\ket1\):

\[\begin{tikzpicture}[sd]
  \draw[wire] (0,0) -- (8,0);
  \node[sdstate] at (0,0) {$0$};
  \node[sdbox] at (2,0) {$H$};
  \node[sdbox] at (4,0) {$Z$};
  \node[sdbox] at (6,0) {$H$};
\end{tikzpicture}
\;=\;
\begin{tikzpicture}[sd]
  \draw[wire] (0,0) -- (5,0);
  \node[sdstate] at (0,0) {$0$};
  \node[sdbox] at (2,0) {$X$};
\end{tikzpicture}
\;=\;
\begin{tikzpicture}[sd]
  \draw[wire] (0,0) -- (3,0);
  \node[sdstate] at (0,0) {$1$};
\end{tikzpicture}\]

Both equations are exact, on the nose --- the same matrix on each side,
number for number. But as pictures they only \emph{report} what the
blackboard found. They do not yet show the roads, and it is the roads
that carry the story: two of them, meeting, one pocket filling and one
emptying. See \texttt{Code/Ch5/TwoRoads.py}.

\textbf{\emph{Why this picture.}} We drew the phase as a box \(Z\)
sitting \emph{on the wire}, not as a fork in it, because a wire that
splits into two roads and rejoins would be a new kind of drawing, and we
have no rule for it yet. The next section supplies the honest Stage-A
way to draw ``both roads at once'' --- a way that never cuts the wire.

\begin{quote}
\textbf{Feynman's Notebook.}

The Mach--Zehnder interferometer is a real optical instrument, and
Cleve, Ekert, Macchiavello and Mosca open their 1997 review --- the
paper this chapter leans on throughout --- by describing it. A particle,
a photon say, meets a half-silvered mirror and takes two paths at once;
a phase shifter sits in each arm; a second half-silvered mirror brings
the two paths back together and directs the particle to one of two
detectors. That second mirror, they stress, erases all record of which
path was taken --- which is exactly what lets the two paths interfere.
They then identify the half-silvered mirror with the Hadamard gate
\(H\), and argue that a quantum computer is at bottom a many-particle
interferometer whose phase shift is computed by a logic gate. The whole
chapter is that one idea, turned over and over.
\end{quote}

\begin{center}\rule{0.5\linewidth}{0.5pt}\end{center}

\hypertarget{a-sum-of-pictures}{%
\section{§5.2 A Sum of Pictures}\label{a-sum-of-pictures}}

Achilles wants to draw the walker going both ways. His first instinct is
to split the wire --- to draw it forking, one branch for each middle
value. The Tortoise will not let him. A wire that forks would be a box
with one input and \emph{two} outputs, a different creature altogether
from the boxes we have (we shall meet it, and its name, much later). In
Stage A a wire never splits. So how is ``both at once'' to be drawn?

\textbf{On the blackboard.} The answer is already in the algebra of
§5.1, and it needs a name. A box applied to a definite value is a
\emph{weighted list} of definite values:

\[A\ket{x}=\sum_{y}A_{yx}\ket{y},\]

where \(A_{yx}\) is the entry of \(A\) in row \(y\) and column \(x\) ---
the amplitude for the road from \(x\) to \(y\). Nothing here is a fork.
It is a \emph{sum}: the output is so many parts of \(\ket0\) plus so
many parts of \(\ket1\), each part a definite value with a number in
front of it. And boxes are \textbf{linear} --- a box applied to a sum is
the sum of the box applied to each part:

\[B\bigl(A\ket{x}\bigr)=\sum_{y}A_{yx}\,B\ket{y}=\sum_{y}A_{yx}\sum_{z}B_{zy}\ket{z}=\sum_{z}\Bigl(\sum_{y}B_{zy}A_{yx}\Bigr)\ket{z}.\]

Read the double sum before it is collected. Each term \(B_{zy}A_{yx}\)
is one \textbf{road}: in the front door at \(x\), through the middle
value \(y\), out at \(z\), carrying the product of the two amplitudes it
met. The output amplitude for \(z\) is the sum of these roads over every
middle value \(y\) they could pass through. This is the picture Chapter
4 gestured at and this chapter is named for: \emph{a computation is a
sum over roads}, each road a definite choice at every step, each
carrying a signed weight, and what you observe at the end is their sum.
It is the book's own way of saying what physicists call a sum over
paths.

\textbf{In the sketchbook.} So we draw ``both at once'' as a \textbf{sum
of pictures}. A picture equals a sum of \emph{whole} pictures, one for
each definite output value, each with its wire unbroken and a number ---
a sign, in the cases we care about --- written in front. Here is the
first and simplest, the gate \(H\) acting on \(\ket0\):

\[\begin{tikzpicture}[sd]
  \draw[wire] (0,0) -- (3,0);
  \node[sdstate] at (0,0) {$0$};
  \node[sdbox] at (2,0) {$H$};
\end{tikzpicture}
\;=\;\tfrac1{\sqrt2}\;\Biggl(\;
\begin{tikzpicture}[sd]
  \draw[wire] (0,0) -- (2,0);
  \node[sdstate] at (0,0) {$0$};
\end{tikzpicture}
\;+\;
\begin{tikzpicture}[sd]
  \draw[wire] (0,0) -- (2,0);
  \node[sdstate] at (0,0) {$1$};
\end{tikzpicture}
\;\Biggr)\]

Each summand on the right is a complete, legal Stage-A picture --- a
definite value on an unbroken wire. No wire was cut. The equality is
exact, on the nose: the same vector on each side, amplitude for
amplitude. From \(\ket1\) the same box gives a sum with a \emph{minus}
on the second road, \(H\ket1=\tfrac1{\sqrt2}(\ket0-\ket1)\): the
\(\ket1\)-road is signed. And because boxes are linear, a sum of
pictures may always be carried forward through whatever box comes next
--- you push each summand through and add. That one permission is what
turns a stack of still pictures into a moving calculation, and every
step of this chapter's picture-proofs will use it.

\textbf{\emph{Why a sum, not a fork.}} Two reasons, and they agree. A
fork would change the \emph{type} of the diagram --- one wire in, two
out --- and would need a rule we do not have. A sum changes nothing: the
wire stays a wire, the drawing stays inside the grammar we built in
Chapters 1--4, and the new content is carried entirely by the little
\(+\) signs and the numbers in front. The picture pays for ``both
roads'' in ink, not in new machinery. See
\texttt{Code/Ch5/SumOfPictures.py}.

\begin{quote}
\textbf{Penrose's Sketchbook.}

A caution about the sum of pictures, because the drawing invites a wrong
reading. The two roads on the right are not two walkers you could stop
and question one at a time. They are summands of a \emph{single}
picture, and the whole of interference lives in that word ``single.''
Put an effect on one road alone --- read the middle value --- and the
sum collapses to that one road; the cancellation of §5.4 needs both
roads, untouched, and so it never survives a peek. Nor is the number
written in front --- the \(\tfrac1{\sqrt2}\), or a kicked
\((-1)^{f(x)}\) --- a box. It sits on no wire and passes through no
gate; row 37 takes the oracle \emph{away} and leaves a bare sign, so the
picture ends with fewer boxes, not a new one. And when two roads of
opposite sign meet, what is left is not some special empty box but the
plain number \(0\).
\end{quote}

\begin{center}\rule{0.5\linewidth}{0.5pt}\end{center}

\hypertarget{the-mark-on-the-other-wire}{%
\section{§5.3 The Mark on the Other
Wire}\label{the-mark-on-the-other-wire}}

A road now carries a sign. Achilles asks the question that unlocks the
chapter: could a road's sign be made to depend on a \emph{question} ---
on some function \(f\)? Then walking the road would interrogate \(f\),
and the sign coming back would be the answer.

To ask a question we need something to ask it of. Bring in a second wire
--- call the first the \textbf{query} wire and the second the
\textbf{answer} wire --- and the oracle that couples them. An oracle for
a function \(f:\{0,1\}\to\{0,1\}\) is the reversible box

\[U_f:\ket{x}\ket{y}\longmapsto\ket{x}\ket{y\oplus f(x)},\]

which leaves the query alone and adds \(f(x)\), modulo \(2\), into the
answer. We draw it as a \emph{tall box} laid across both wires --- not a
thin line, so that it is never mistaken for a wire of its own. Feed it a
definite query \(\ket{x}\) and a definite answer \(\ket{y}\) and it does
the obvious thing: it writes \(f(x)\) into the answer. If the answer
starts at \(\ket0\), it comes out reading \(f(x)\), and the mark lands
where we expect --- on the answer wire.

\textbf{On the blackboard.} Now do the one thing that is not obvious:
set the answer wire not to a definite value but to the \emph{difference}
state

\[\ket-=\tfrac1{\sqrt2}\bigl(\ket0-\ket1\bigr).\]

Watch what \(U_f\) does to \(\ket{x}\ket-\), taking the two cases of
\(f(x)\) together. Since \(U_f\) sends
\(\ket{x}\ket0\mapsto\ket{x}\ket{f(x)}\) and
\(\ket{x}\ket1\mapsto\ket{x}\ket{1\oplus f(x)}\),

\[U_f\,\ket{x}\ket- \;=\; \tfrac1{\sqrt2}\ket{x}\bigl(\ket{f(x)}-\ket{1\oplus f(x)}\bigr) \;=\; (-1)^{f(x)}\,\ket{x}\ket-.\]

If \(f(x)=0\) the bracket is \(\ket0-\ket1\), unchanged; if \(f(x)=1\)
it is \(\ket1-\ket0\), which is \(-(\ket0-\ket1)\). Either way the
answer wire comes back \emph{exactly} as it went in, \(\ket-\), and the
entire action of the oracle has become a single sign, \((-1)^{f(x)}\),
standing in front. The answer wire is an eigenstate of the oracle; the
mark the oracle would have made on it has nowhere to land there, so it
lands on the \emph{query} wire instead. This is \textbf{phase kickback}
(Cleve and colleagues, 1997, Eq. 2.2). The bar touched the answer wire;
the mark appeared on the query.

\textbf{In the sketchbook.} With the answer wire held in \(\ket-\), the
tall box \(U_f\) across a definite query and that held answer equals a
\emph{number}, \((-1)^{f(x)}\), written before the same picture with the
oracle taken away:

\[\begin{tikzpicture}[sd]
  \draw[wire] (0,2) -- (4,2);
  \draw[wire] (0,0) -- (4,0);
  \node[sdstate] at (0,2) {$x$};
  \node[sdstate] at (0,0) {$-$};
  \node[sdbox, minimum height=3.9em] at (2,1) {$U_f$};
\end{tikzpicture}
\;=\;(-1)^{f(x)}\;
\begin{tikzpicture}[sd]
  \draw[wire] (0,2) -- (3,2);
  \draw[wire] (0,0) -- (3,0);
  \node[sdstate] at (0,2) {$x$};
  \node[sdstate] at (0,0) {$-$};
\end{tikzpicture}\]

exact, sign included. The oracle has spent itself: the two wires are as
they were, and all that is left of the query is a signed number. See
\texttt{Code/Ch5/Kickback.py}.

\textbf{\emph{Why the difference state.}} Put a definite \(\ket0\) on
the answer wire and the oracle writes \(f(x)\) there and leaves the
query untouched --- the mark stays on the wire you asked \emph{about}.
Only when the answer wire is an eigenstate of the flip, \(\ket-\), does
the oracle have nothing to change there, and the only bookkeeping left
is a sign on the query --- the wire you asked \emph{from}. The
difference state is precisely the state that turns an action into a
mark.

\begin{quote}
\textbf{Dirac's Blackboard.}

\emph{Why \(\ket-\) and not \(\ket+\).} The flip on the answer wire is
the box \(X\), and its eigenvectors are
\(\ket+=\tfrac1{\sqrt2}(\ket0+\ket1)\) and
\(\ket-=\tfrac1{\sqrt2}(\ket0-\ket1)\):
\[X\ket+=+\ket+,\qquad X\ket-=-\ket-.\] When \(f(x)=1\) the oracle
applies \(X\) to the answer wire; when \(f(x)=0\) it does nothing. So on
the eigenvector \(\ket-\) the oracle multiplies by \((-1)\) exactly when
\(f(x)=1\), i.e.~by \((-1)^{f(x)}\), and \(\ket+\) would give \((+1)\)
always --- no mark at all. The kicked sign is the eigenvalue \(-1\) of
\(X\), carried out of the answer wire and set down on the query. Nothing
is created: the phase was \(X\)'s all along, and \(\ket-\) merely
refuses to hide it. See \texttt{Code/Ch5/Kickback.py}.
\end{quote}

\begin{center}\rule{0.5\linewidth}{0.5pt}\end{center}

\hypertarget{the-one-question}{%
\section{§5.4 The One Question}\label{the-one-question}}

We can now ask \(f\) one question and get back, not \(f(0)\), not
\(f(1)\), but the one thing a single walk can decide with certainty:
whether they \emph{agree}.

First feel the need. Suppose you are allowed to consult the oracle
exactly once, and you want to know whether \(f\) is \textbf{constant}
(\(f(0)=f(1)\)) or \textbf{balanced} (\(f(0)\neq f(1)\)). Classically
you are stuck: one evaluation tells you one value, \(f(0)\) \emph{or}
\(f(1)\), and either value alone is consistent with both a constant
\(f\) and a balanced one. You would need two evaluations. Nor does the
quantum machine rescue you by letting you read a single road: send in
\(\ket0\), kick a sign, and \emph{read} the query wire, and you get a
fair coin that tells you nothing about \(f\) --- each single road
answers in superposition. The trap looks airtight. It is not, and the
escape is interference.

Here is a concrete run to make the mechanism visible. Take the balanced
function \(f(0)=0,\,f(1)=1\) (the identity). Prepare the query in
\(\ket0\), the answer in \(\ket-\), and put an \(H\) on the query before
and after the oracle.

\begin{itemize}
\tightlist
\item
  After the first \(H\): the query is \(\tfrac1{\sqrt2}(\ket0+\ket1)\)
  --- both roads open.
\item
  Through the oracle: kickback stamps each road with the oracle's sign
  --- the \(\ket0\)-road by \((-1)^{f(0)}\), the \(\ket1\)-road by
  \((-1)^{f(1)}\) --- which for this \(f\) leaves
  \(\tfrac1{\sqrt2}(\ket0-\ket1)\). The two roads now carry
  \emph{opposite} signs, because \(f\) sent them opposite answers.
\item
  Through the second \(H\): this is \(H\ket-=\ket1\). The roads to
  \(\ket0\) cancel; the walker is turned away to \(\ket1\).
\end{itemize}

Read the query and you get \(\ket1\), with certainty --- the balanced
verdict --- from a single consultation. Try a constant \(f\) and the two
roads come back with the \emph{same} sign, pile up, and the walker comes
home to \(\ket0\). The one walk did what two evaluations do.

\textbf{On the blackboard.} In general, for any of the four functions
\(f\),

\[\begin{aligned}
(H\otimes I)\,U_f\,(H\otimes I)\,\ket0\ket-
&= (H\otimes I)\,U_f\,\tfrac1{\sqrt2}\bigl(\ket0+\ket1\bigr)\ket- \\
&= (H\otimes I)\,\tfrac1{\sqrt2}\bigl((-1)^{f(0)}\ket0+(-1)^{f(1)}\ket1\bigr)\ket- \\
&= (-1)^{f(0)}\,(H\otimes I)\,\tfrac1{\sqrt2}\bigl(\ket0+(-1)^{f(0)\oplus f(1)}\ket1\bigr)\ket- \\
&= (-1)^{f(0)}\,\ket{f(0)\oplus f(1)}\,\ket-.
\end{aligned}\]

The last step is the two cases of §5.1 read backwards: if
\(f(0)\oplus f(1)=0\) the query is \(H\ket+=\ket0\); if it is \(1\), the
query is \(H\ket-=\ket1\). So the query wire ends in
\(\ket{f(0)\oplus f(1)}\) --- \(\ket0\) for constant, \(\ket1\) for
balanced --- read off with a single certain measurement, one oracle
call. \textbf{This is Deutsch's algorithm} (Theorem 5.1), and §5.7
proves it twice.

The factor \((-1)^{f(0)}\) out front is a \textbf{global phase}. It
multiplies the whole state, so it cannot change any measured outcome ---
the query still reads \(f(0)\oplus f(1)\) with certainty. But it is
genuinely there, and the equation is only true with it: drop it and you
have written down a falsehood. We keep it. See
\texttt{Code/Ch5/Deutsch.py}, which checks the equation for all four
\(f\) with the sign intact.

\textbf{\emph{Why one query can do it.}} Classically, one call yields
one bit and that bit is one value of \(f\); the relation
\(f(0)\oplus f(1)\) is a second bit that no single value determines.
Quantumly, the single call is spent on \emph{both} roads at once,
kickback writes the two answers into the two roads' signs, and the
second \(H\) reads out not either sign but whether they match. You never
learn \(f(0)\) or \(f(1)\) --- the global phase is exactly the
information you threw away --- and in exchange you learn the one thing
their \emph{difference} holds.

\begin{quote}
\textbf{Feynman's Notebook.}

A careful word about who did what, because it is easy to get wrong.
David Deutsch posed this problem in 1985 and gave the first quantum
algorithm for it --- but that original algorithm was
\emph{probabilistic}. As Cleve, Ekert, Macchiavello and Mosca describe
it, Deutsch's 1985 procedure had three possible outcomes, ``constant,''
``balanced,'' and ``inconclusive'': with probability one-half it
returned the right answer, and with probability one-half it returned
``inconclusive'' and told you nothing. Still remarkable --- no classical
single evaluation can even do that --- but not a certainty.

The \emph{deterministic} one-query algorithm, the one proved above that
always answers correctly, came later: it is due to Cleve, Ekert,
Macchiavello and Mosca (1997), who record that Alain Tapp found a
similar algorithm independently the same year (private communication,
1997). So the version this chapter calls ``Deutsch's algorithm'' is
really the 1997 improvement of Deutsch's 1985 idea. The credit for the
certain single-query answer belongs to them, not to the 1985 paper.
\end{quote}

\begin{center}\rule{0.5\linewidth}{0.5pt}\end{center}

\hypertarget{the-sixteen-roads}{%
\section{§5.5 The Sixteen Roads}\label{the-sixteen-roads}}

One walk decided a two-valued question. Stretch the question over more
roads and ask how far the cancelling reaches.

Let \(f\) now take four bits in, \(f:\{0,1\}^4\to\{0,1\}\), with a
promise: \(f\) is either \textbf{constant} (the same value on all
sixteen inputs) or \textbf{balanced} (value \(0\) on exactly eight
inputs and \(1\) on the other eight). Which is it? Classically the worst
case is cruel: you may see eight equal values in a row and still not
know, so you may need \(2^{4-1}+1=9\) evaluations before you are
certain. This is the problem Deutsch and Jozsa posed and solved in 1992
for inputs of any width; the one-oracle-call circuit we are about to
draw is the streamlined form given by Cleve and colleagues in 1997.

\textbf{On the blackboard.} Put four query wires, each starting at
\(\ket0\) with an \(H\) before and after the oracle, and one answer wire
held in \(\ket-\). After the first four \(H\)'s the query register is an
equal sum over all sixteen inputs,
\(H^{\otimes4}\ket{0000}=\tfrac14\sum_{x}\ket{x}\) (four factors of
\(\tfrac1{\sqrt2}\) make \(\tfrac14\)). Kickback stamps each input's
road with its sign, and the second layer of \(H\)'s brings the roads
together. We ask for one number: the amplitude with which the whole
thing returns to \(\ket{0000}\), which is what a closed diagram
computes. It is

\[\bra{0000}\,H^{\otimes4}\,U_f\,H^{\otimes4}\,\ket{0000}\;=\;\frac{1}{16}\sum_{x\in\{0,1\}^4}(-1)^{f(x)},\]

with the answer wire riding in \(\ket-\) at both ends, its kickback
already absorbed into the signs and \(\braket{-|-}=1\) leaving no stray
factor behind. (We keep \(\ket-\) \emph{normalized} at both the start
and the finish precisely so that this number comes out clean; an
unnormalized answer wire would multiply it by \(2\).) Read the two
cases:

\begin{itemize}
\tightlist
\item
  \(f\) \textbf{constant}: all sixteen signs agree, the sum is
  \(\pm16\), the number is \(\pm1\). Every road piles into the home
  pocket.
\item
  \(f\) \textbf{balanced}: eight signs are \(+\) and eight are \(-\),
  the sum is \(0\), the number is \(0\). Eight roads cancel eight, and
  the home pocket empties.
\end{itemize}

The chapter's title, made literal. See
\texttt{Code/Ch5/SixteenRoads.py}, which builds a genuine balanced \(f\)
(checking the tally is eight-and-eight before trusting the zero) and a
constant one, and confirms both values.

\textbf{In the sketchbook.} Draw the whole thing closed --- states at
the head of every wire, effects at the foot --- so that the picture is a
single number. This is the master diagram of the chapter, \textbf{The
Sixteen Roads}: four query wires, each running
\(\ket0\to H\to U_f\to H\to\bra0\), and one answer wire running
\(\ket-\to U_f\to\bra-\) with no \(H\), since it is the kickback
eigenstate. The tall box \(U_f\) is laid across them all.

\begin{center}
\begin{tikzpicture}[sd]
  \draw[wire] (0,8) -- (11,8);
  \draw[wire] (0,6) -- (11,6);
  \draw[wire] (0,4) -- (11,4);
  \draw[wire] (0,2) -- (11,2);
  \draw[wire] (0,0) -- (11,0);
  \node[sdstate] at (0,8) {$0$};
  \node[sdstate] at (0,6) {$0$};
  \node[sdstate] at (0,4) {$0$};
  \node[sdstate] at (0,2) {$0$};
  \node[sdstate] at (0,0) {$-$};
  \node[sdbox] at (2.5,8) {$H$};
  \node[sdbox] at (2.5,6) {$H$};
  \node[sdbox] at (2.5,4) {$H$};
  \node[sdbox] at (2.5,2) {$H$};
  \node[sdbox, minimum height=9.2em] at (5.5,4) {$U_f$};
  \node[sdbox] at (8.5,8) {$H$};
  \node[sdbox] at (8.5,6) {$H$};
  \node[sdbox] at (8.5,4) {$H$};
  \node[sdbox] at (8.5,2) {$H$};
  \node[sdeffect] at (11,8) {$0$};
  \node[sdeffect] at (11,6) {$0$};
  \node[sdeffect] at (11,4) {$0$};
  \node[sdeffect] at (11,2) {$0$};
  \node[sdeffect] at (11,0) {$-$};
\end{tikzpicture}
\end{center}

Each query wire runs unbroken from its state to its effect, straight
through the tall oracle; the answer wire runs likewise beneath. The
oracle crosses them all and joins none --- a wide box is not a wire.
This closed picture \emph{is} the number
\(\tfrac1{16}\sum_x(-1)^{f(x)}\): sixteen roads through it, each a
definite choice of the four query values, each carrying a sign, summed.
When they agree it is \(\pm1\); when they balance it is \(0\). The
equality is exact, on the nose.

\textbf{\emph{Why draw it closed.}} An open diagram is a state or a
process --- a thing with loose ends. Closing both ends turns it into a
single number (the rule of Chapter 2), and it is the number, not the
process, that answers the constant-or-balanced question. The picture
makes the answer a \emph{pile height}: sixteen tokens in a pocket, or
none.

\begin{center}\rule{0.5\linewidth}{0.5pt}\end{center}

\hypertarget{the-rosetta-table-4}{%
\section{§5.6 The Rosetta Table}\label{the-rosetta-table-4}}

The rows this chapter adds to the running dictionary. Rows 33 and 34
introduce the new picture element --- the \emph{sum of pictures} --- in
its first and general form; row 35 reads a product of boxes as a sum
over roads; row 36 is the rule that lets two roads pile up or cancel;
rows 37 and 38 record kickback and the closed diagram whose value the
signs decide. Every equality is exact, scalar included.

\begingroup\renewcommand{\arraystretch}{1.6}

\begin{longtable}[]{@{}
  >{\raggedright\arraybackslash}p{(\columnwidth - 4\tabcolsep) * \real{0.2727}}
  >{\raggedright\arraybackslash}p{(\columnwidth - 4\tabcolsep) * \real{0.3636}}
  >{\raggedright\arraybackslash}p{(\columnwidth - 4\tabcolsep) * \real{0.3636}}@{}}
\toprule\noalign{}
\begin{minipage}[b]{\linewidth}\raggedright
\#
\end{minipage} & \begin{minipage}[b]{\linewidth}\raggedright
Algebra (Dirac)
\end{minipage} & \begin{minipage}[b]{\linewidth}\raggedright
Picture (Penrose)
\end{minipage} \\
\midrule\noalign{}
\endhead
\bottomrule\noalign{}
\endlastfoot
33 & \(H\ket0=\tfrac1{\sqrt2}(\ket0+\ket1)\);
\(H\ket1=\tfrac1{\sqrt2}(\ket0-\ket1)\) & a state through the box \(H\)
\(=\) a \(\ket0\)-road plus a \(\ket1\)-road; from \(\ket0\) both roads
carry \(+\), from \(\ket1\) the \(\ket1\)-road is signed \(-\) --- the
first \emph{sum of pictures} (\emph{new element}) \\
34 & \(A\ket{x}=\sum_y A_{yx}\ket{y}\) & a state through \textbf{any}
box \(=\) a sum of one-road pictures, one per output value, each
weighted by its entry, every wire unbroken --- the general \emph{sum of
pictures} \\
35 &
\(BA\ket{x}=\sum_{y}(BA)_{yx}\ket{y}=\sum_{y,z}B_{zy}A_{yx}\ket{z}\) &
two boxes in a row on a state \(=\) a sum over the middle value \(=\)
the \textbf{roads} (the sum over paths) \\
36 & \(\ket\psi+\ket\phi\); opposite amplitudes give \(0\) & one-road
pictures of the \textbf{same} sign make one picture of doubled weight;
of \textbf{opposite} sign they make the number \(0\) (the roads
cancel) \\
37 & \(U_f\ket{x}\ket{-}=(-1)^{f(x)}\ket{x}\ket{-}\), \(\ket-\) an
eigenstate & the tall box \(U_f\) across a definite query and a
\(\ket-\) answer \(=\) the number \((-1)^{f(x)}\) before the same
picture with \(U_f\) removed --- \textbf{phase kickback} \\
38 &
\(\bra{0}^{\otimes n}\!\bra{-}\,(H^{\otimes n}\!\otimes I)U_f(H^{\otimes n}\!\otimes I)\,\ket{0}^{\otimes n}\!\ket{-}=\tfrac1{2^n}\sum_x(-1)^{f(x)}\)
& a fully closed diagram \(=\) a \textbf{number}: the signed sum of its
one-road pictures, \(0\) when the \(+\) and \(-\) roads balance \\
\end{longtable}

\endgroup

All six rows are exact equalities. The sum of pictures (rows 33--36) is
the last new element of Stage A; from the next chapter the wire itself
will be allowed to bend.

\begin{center}\rule{0.5\linewidth}{0.5pt}\end{center}

\hypertarget{the-theorem-twice-4}{%
\section{§5.7 The Theorem, Twice}\label{the-theorem-twice-4}}

\begin{quote}
\textbf{Theorem 5.1 (Deutsch's algorithm).} For \(f:\{0,1\}\to\{0,1\}\)
given as the oracle
\(U_f:\ket{x}\ket{y}\mapsto\ket{x}\ket{y\oplus f(x)}\),
\[(H\otimes I)\,U_f\,(H\otimes I)\,\ket0\ket- \;=\; (-1)^{f(0)}\,\ket{f(0)\oplus f(1)}\,\ket-,\]
exactly, global phase included --- so one oracle call decides constant
(\(f(0)=f(1)\)) versus balanced (\(f(0)\neq f(1)\)) with certainty,
where any classical strategy needs two evaluations.
\end{quote}

\hypertarget{diracs-blackboard-4}{%
\subsection{Dirac's Blackboard}\label{diracs-blackboard-4}}

Line by line, with the answer wire in \(\ket-\) throughout. Write
\(a=f(0)\oplus f(1)\).

\[\begin{aligned}
\ket0\ket- \;&\xrightarrow{H\otimes I}\; \tfrac1{\sqrt2}\bigl(\ket0+\ket1\bigr)\ket-
 && \text{(both roads open)}\\
 &\xrightarrow{U_f}\; \tfrac1{\sqrt2}\bigl((-1)^{f(0)}\ket0+(-1)^{f(1)}\ket1\bigr)\ket-
 && \text{(kickback)}\\
 &=\; (-1)^{f(0)}\,\tfrac1{\sqrt2}\bigl(\ket0+(-1)^{a}\ket1\bigr)\ket-
 && \text{(common sign out)}\\
 &\xrightarrow{H\otimes I}\; (-1)^{f(0)}\,\ket{a}\,\ket-.
 && \text{(}H\ket-=\ket1\text{)}
\end{aligned}\]

Three lines of work. The middle two are the whole content: kickback
turns the oracle into a pair of signs on the two roads, and the last
\(H\) reads whether those signs agree (\(a=0\), query \(\ket0\)) or
differ (\(a=1\), query \(\ket1\)). The global phase \((-1)^{f(0)}\)
rides through untouched and is still there at the end.
\(\;\blacksquare\)

\hypertarget{penroses-sketchbook-4}{%
\subsection{Penrose's Sketchbook}\label{penroses-sketchbook-4}}

The same calculation, drawn --- tracking the query wire's road-sum, with
the answer wire held in \(\ket-\) beside it and supplying every sign
through kickback (row 37).

\textbf{Step 1 --- split the first \(H\) (row 33).} The start \(\ket0\)
through \(H\) becomes two roads of equal, positive weight:

\[\begin{tikzpicture}[sd]
  \draw[wire] (0,0) -- (3,0);
  \node[sdstate] at (0,0) {$0$};
  \node[sdbox] at (2,0) {$H$};
\end{tikzpicture}
\;=\;\tfrac1{\sqrt2}\Biggl(\;
\begin{tikzpicture}[sd]
  \draw[wire] (0,0) -- (1.3,0);
  \node[sdstate] at (0,0) {$0$};
\end{tikzpicture}
\;+\;
\begin{tikzpicture}[sd]
  \draw[wire] (0,0) -- (1.3,0);
  \node[sdstate] at (0,0) {$1$};
\end{tikzpicture}
\;\Biggr)\]

\textbf{Step 2 --- carry the sum through the oracle (linearity + row
37).} Push each road through \(U_f\). The answer wire is \(\ket-\), so
the oracle spends itself as a sign on each query road: the
\(\ket0\)-road picks up \((-1)^{f(0)}\), the \(\ket1\)-road picks up
\((-1)^{f(1)}\).

\[\tfrac1{\sqrt2}\Biggl(\,(-1)^{f(0)}\!\!
\begin{tikzpicture}[sd]
  \draw[wire] (0,0) -- (1.3,0);
  \node[sdstate] at (0,0) {$0$};
\end{tikzpicture}
\;+\;(-1)^{f(1)}\!\!
\begin{tikzpicture}[sd]
  \draw[wire] (0,0) -- (1.3,0);
  \node[sdstate] at (0,0) {$1$};
\end{tikzpicture}
\,\Biggr)
\;=\;(-1)^{f(0)}\,\tfrac1{\sqrt2}\Biggl(
\begin{tikzpicture}[sd]
  \draw[wire] (0,0) -- (1.3,0);
  \node[sdstate] at (0,0) {$0$};
\end{tikzpicture}
\;+\;(-1)^{a}\!\!
\begin{tikzpicture}[sd]
  \draw[wire] (0,0) -- (1.3,0);
  \node[sdstate] at (0,0) {$1$};
\end{tikzpicture}
\,\Biggr)\]

with the common sign \((-1)^{f(0)}\) pulled to the front and
\(a=f(0)\oplus f(1)\).

\textbf{Step 3 --- split the second \(H\) and collect (row 33, then row
36).} Push each road through the last \(H\): the \(\ket0\)-road becomes
\(\tfrac1{\sqrt2}(\ket0+\ket1)\), the \(\ket1\)-road becomes
\(\tfrac1{\sqrt2}(\ket0-\ket1)\). Four one-road pictures now land in two
pockets. In the \(\ket0\) pocket the weight is
\(\tfrac12\bigl(1+(-1)^a\bigr)\); in the \(\ket1\) pocket,
\(\tfrac12\bigl(1-(-1)^a\bigr)\). By row 36 the two roads in a pocket
pile up or cancel:

\[(-1)^{f(0)}\Biggl(\tfrac12\bigl(1+(-1)^{a}\bigr)\!\!
\begin{tikzpicture}[sd]
  \draw[wire] (0,0) -- (1.3,0);
  \node[sdstate] at (0,0) {$0$};
\end{tikzpicture}
\;+\;\tfrac12\bigl(1-(-1)^{a}\bigr)\!\!
\begin{tikzpicture}[sd]
  \draw[wire] (0,0) -- (1.3,0);
  \node[sdstate] at (0,0) {$1$};
\end{tikzpicture}
\Biggr)
\;=\;(-1)^{f(0)}\!\!
\begin{tikzpicture}[sd]
  \draw[wire] (0,0) -- (1.3,0);
  \node[sdstate] at (0,0) {$a$};
\end{tikzpicture}\]

For \(a=0\) (constant) the \(\ket0\) pocket fills to \(1\) and the
\(\ket1\) pocket empties; for \(a=1\) (balanced) the reverse. Either way
one pocket carries a single road of weight \(\pm1\) and the other holds
nothing: the query is \(\ket{a}=\ket{f(0)\oplus f(1)}\), with the global
phase \((-1)^{f(0)}\) still in front, and the answer wire is \(\ket-\)
as it always was. \(\;\blacksquare\)

\hypertarget{the-verdict-4}{%
\subsection{The verdict}\label{the-verdict-4}}

The two proofs are the same length and the same calculation --- line for
line, the picture translates the algebra: split, carry, kick, split,
collect. This is honest to report and worth stating plainly:
\textbf{there is no Stage-A rewrite that derives Deutsch's algorithm
more cheaply than the kets do.} The blackboard is, if anything, the more
economical of the two, because it writes a superposition as one short
sum where the picture must draw each road as a separate tile.

What the picture earns is not brevity but \emph{sight}. On the
blackboard the cancellation is a line of arithmetic,
\(\tfrac12(1-(-1)^a)\), that happens to vanish; in the sketchbook it is
two tokens of opposite face dropped into one pocket, which you watch
empty. The strand that proves the theorem here is the algebra. The
strand that shows you \emph{why the answer is certain} --- why one
pocket must be full and the other bare --- is the picture. Neither
claims the other's virtue.

\begin{center}\rule{0.5\linewidth}{0.5pt}\end{center}

\hypertarget{the-sign-that-was-not-there}{%
\section{§5.8 The Sign That Was Not
There}\label{the-sign-that-was-not-there}}

Gather the chapter. Two roads to one place can add or cancel, and a
half-turn of phase between two mirrors decides which (§5.1). ``Both
roads at once'' is drawn not by cutting a wire but by summing whole
pictures, each a signed road (§5.2). With the answer wire held in
\(\ket-\), an oracle's action becomes a sign kicked onto the query
(§5.3); two such signs, brought together by a second \(H\), report
whether \(f\) agrees with itself --- one query, a certain verdict
(§5.4); and sixteen roads under a balanced oracle cancel to nothing at
all (§5.5). Throughout, the deciding quantity was a \emph{sign} --- a
thing with no size, invisible to any single road, felt only where two
roads meet.

That is the last work Stage A will do with a straight wire. Every road
in this chapter ran unbroken from one end of the board to the other; the
only new thing we allowed was to \emph{sum} such roads. Notice what we
never once needed: a wire that leaves and returns. Notice, too, what the
wires never did. The whole trick of kickback was that a mark made on one
wire showed up on another; yet the answer wire went in as \(\ket-\) and
came out as \(\ket-\), and at every step it could be described on its
own, apart from the query --- the hopping sign was only a number
standing in front. The wires talked, and still they came apart. In the
next chapter we meet a pair that does not come apart. We let a single
wire bend back on itself so that both its ends point the same way, and
we find that this bent wire --- going out and coming home --- is by
itself a shared fact between two places, held in neither of them. The
roads that cancel were a rehearsal. What is kept in the bend is the
thing itself.

\hypertarget{obbligato-for-six-music-boxes}{%
\chapter{Obbligato for Six Music
Boxes}\label{obbligato-for-six-music-boxes}}

\begin{center}\rule{0.5\linewidth}{0.5pt}\end{center}

\emph{長き夜や}\\
\emph{離れて一つ}\\
\emph{オルゴール}

\emph{nagaki yo ya}\\
\emph{hanarete hitotsu}\\
\emph{orugōru}

\emph{The long autumn night ---}\\
\emph{music boxes, far apart,}\\
\emph{one note between them.}

\begin{center}\rule{0.5\linewidth}{0.5pt}\end{center}

\begin{center}\rule{0.5\linewidth}{0.5pt}\end{center}

\emph{{[}Achilles and the Tortoise have come to the Crab's for the
afternoon; the Crab meets them at the door, squinting into the low
sun.{]}}

\textbf{Crab:} Come in, come in, both of you --- mind the step; at this
hour you come up to it in your own shadow, and it hides there. You found
the lane?

\textbf{Achilles:} Tortoise found it. Left to me, I'd have walked us
straight into the canal.

\textbf{Tortoise:} He was giving directions and telling a story at the
same time. One of them had to go.

\textbf{Achilles:} A concert, it was --- a quartet, and the cellist a
whole bar behind the other three the entire evening. I have never heard
anything so wrong stay so cheerful about it.

\textbf{Crab:} Then you have come to exactly the right house. Come
through; there's a pot just drawn. I have something I've been aching to
show a pair of people who'd know what they were looking at --- and I
have waited all week for tonight.

\textbf{Tortoise:} We are all attention, Mr Crab. Achilles most of all;
he has promised to keep step with me for once, which is more than his
cellist managed.

\emph{{[}They go through to the warm room and take the chairs he waves
them to, the lane shut out behind them. The Crab pours, glad of company,
and hands the cups round; Achilles, still beating time for his cellist
with one hand, has to be offered his a second time.{]}}

\begin{center}\rule{0.5\linewidth}{0.5pt}\end{center}

\hypertarget{first-challenge-two-recipe-cards}{%
\section{\texorpdfstring{First Challenge --- \emph{Two Recipe
Cards}}{First Challenge --- Two Recipe Cards}}\label{first-challenge-two-recipe-cards}}

\emph{{[}The Crab's music room, late on an autumn afternoon. Nearly
everything in it was bought as a pair and has stayed one: two upright
pianos side by side against the long wall, each tuned to the other, so
that it is no use asking which of them is flat; two music stands sharing
the two halves of one duet; a lamp at either end of the table, not yet
lit; and on the mantelpiece two clocks by one maker, which do not keep
step --- their ticks fall anyhow, like two people walking on gravel.
Over the mantel, between the clocks, hangs Escher's lithograph
\textup{Bond of Union} in a narrow black frame: two heads, face to face,
hollow as peeled rind, with small globes adrift in the dark around them
and inside them. The sun, low in the one tall window, comes through a
pair of net curtains and has not reached it yet.{]}}

\emph{{[}The Crab's table. On a tray stand his music boxes, none bigger
than a matchbox, set out two by two: a pair in pearwood, a pair in
walnut, a pair in ebony. Each has a lid on top and a small door in its
side. Beside the tray lies a ledger.{]}}

\textbf{Crab:} Now, before anyone touches anything. These came to me
from a little workshop in Delft, behind the Oude Kerk, where a father
and his daughter make nothing else. They are made in pairs, and only in
pairs --- the pearwood together, the walnut together, the ebony together
--- each pair voiced at one bench on one afternoon. And the first thing
the daughter tells every customer is---

\emph{{[}The left pearwood box sits shut on the tray, one note folded up
inside it, whole. Achilles's thumb catches the lid and lifts it --- the
box itself does not move, only the little hinged lid swings up --- and
the note springs out into the room: a single high tone, clear as a
struck glass. Then it is gone, and the lid settles back on a box with
nothing left inside to sound. What was held is spent; what stays is a
small pearwood box that will never sing again.{]}}

\textbf{Achilles:} Oh, that's lovely. How do you wind it again?

\textbf{Crab:} ---is that one doesn't.

\textbf{Achilles:} Doesn't what?

\textbf{Crab:} Wind it again. Each box plays one note, once. High or
low; you have had high. It is now a very small pearwood box with nothing
further to say. \emph{{[}He writes a single word in the left-hand column
of the ledger --- the first of a pair --- and leaves the space beside it
in the right-hand column blank; then he lifts the right pearwood box off
the tray and folds it against his shell, shut and full, while its spent
partner sits shut and silent on the tray{]}} And nobody touches this
one.

\textbf{Achilles:} I'm terribly sorry. One note? It seems a great deal
of cabinet-making for one note.

\textbf{Tortoise:} It was a very good note.

\emph{{[}The left pearwood box is the Crab's marvel no longer: the same
close grain, the same hinge no bigger than a grain of rice, and about as
much left to say for itself as a box one keeps stamps in.{]}}

\textbf{Crab:} It was half of something. There are two ways in, by the
way, since you didn't wait to be told. The lid, which you have found.
And the side door, which strikes the reed from the other side: by the
door one hears not high or low but one of two other notes --- the
workshop calls them bright and dark. Either way: one note, once, and the
box is spent.

\textbf{Achilles:} And which does a box play? By the lid, say. High or
low?

\textbf{Crab:} Nobody can tell you. I have weighed them. I have held
them up to the lamp. The daughter in Delft says she doesn't know either,
and I believe her; she has an honest face.

\textbf{Achilles:} Let me put it on paper the way one writes out a
recipe --- a card for each box, and a card meant to hold everything
there is to say about the thing it describes. If a box really is nothing
but a coin, one card will catch it whole.

\textbf{Achilles:} Then I can tell you, in the only sense that's left.
\emph{{[}He takes two blank cards from the Crab's recipe tin and
writes.{]}} A recipe card for each box. Left box: ``plays high half the
time, low half the time.'' Right box: the same. There --- the pair,
fully described. Two boxes, two cards, and a coin toss on each.

\textbf{Tortoise:} A recipe ought to be cookable. Would you play your
cards for us? A coin for each card; heads for high.

\textbf{Achilles:} Gladly. \emph{{[}He tosses two coins together.{]}}
High, high. \emph{{[}Again.{]}} High, low. Low, low. Low, high. High,
low\ldots{}

\emph{{[}The Anteater is shown in while the coins are in the air.{]}}

\textbf{Anteater:} Good evening, good evening --- don't get up, you are
in the middle of a tally. I bring Aunt Hillary's regards. She has moved
her nurseries to the south slope of the hill this week, and is very
pleased with herself.

\textbf{Crab:} Delighted. Sit by the ledger, Doctor; it is about to be
wanted.

\emph{{[}Fresh tea comes in behind him: two pots of one pattern, one
holding tea and the other hot water, which the household has never
learned to tell apart except by pouring. The Anteater sits where he is
told. There is a little reddish earth dried on his cuff.{]}}

\textbf{Achilles:} \ldots low, high. High, high. That's twenty throws.
The coins agreed eleven times and disagreed nine. About half and half,
which is as it should be: a coin doesn't care what its neighbour does.

\textbf{Tortoise:} And the boxes, Mr Crab? How often have the two boxes
of a pair disagreed?

\textbf{Crab:} \emph{{[}opening the ledger{]}} I have been spending
pairs on my friends for eleven years. Left box in the left column, right
box in the right column, a line for each pair. Read across, and keep
going.

\textbf{Achilles:} High, high. Low, low. Low, low. High, high. Low,
low\ldots{} \emph{{[}He turns a page. Another.{]}} Here's a ``bright,
bright''.

\textbf{Crab:} A side-door evening. Both boxes by the door. One opens a
pair by the same way in; I was told so very firmly.

\textbf{Achilles:} \ldots dark, dark. Bright, bright. High, high. They
never disagree. Not once, in --- it must be hundreds of lines. My coins
couldn't keep it up for twenty.

\textbf{Tortoise:} Then your cards are the recipe for some other dish.

\textbf{Achilles:} But I'd still swear to each of them! Whatever is true
of that left box, it's on my left card, and the same for the right. It's
only when I lay the two cards side by side that they start telling lies.
How can two honest cards make a dishonest pair?

\textbf{Tortoise:} What did you leave off them?

\textbf{Achilles:} Each other. \emph{{[}He tears both cards across,
takes one fresh card, and writes along its whole width.{]}} ``The pair:
both high half the time, both low half the time.'' Two cards won't do
it. It needs one card for the pair.

\textbf{Crab:} And that card will not cut in two. I have thought about
it on many evenings. Whichever way you take the scissors to it, you lose
the word ``both''.

\emph{{[}At the window the two net curtains hang one behind the other,
and where they overlap the low sun shows a watered pattern, broad slow
bands like the figure in walnut, which is woven into neither of them. A
draught lifts the inner curtain by a finger's breadth; the bands race
the whole height of the window, and settle.{]}}

\textbf{Achilles:} But now I've a card for the pair and no card for a
box. And there it sits in your claw, Mr Crab --- one box, a solid
object, with a lid. Then what does \emph{one} box know?

\begin{center}\rule{0.5\linewidth}{0.5pt}\end{center}

\hypertarget{second-challenge-the-ledger}{%
\section{\texorpdfstring{Second Challenge --- \emph{The
Ledger}}{Second Challenge --- The Ledger}}\label{second-challenge-the-ledger}}

\textbf{Anteater:} At which level are you asking?

\textbf{Achilles:} I beg your pardon?

\textbf{Anteater:} You asked what one box knows. It is a fair question,
but questions have addresses, and I am not sure yours has been delivered
to the right one. Mr Crab, might the blotter be laid over one of your
columns?

\textbf{Crab:} The blotter is here by the inkwell; I keep it to dry a
wet page. It will hide a column of the ledger as neatly as it dries one.

\textbf{Crab:} I'll do it. Nobody handles the ledger but me. \emph{{[}He
lays the sheet of blotting paper down the right-hand column, edge to the
centre rule: the right column vanishes under it, the left column stays
open to the lamp, and eleven years of left-hand boxes stand there with
nothing beside them to agree or disagree with{]}} There. The left boxes
by themselves, eleven years of them.

\textbf{Achilles:} \emph{{[}running a finger down the page and keeping a
tally on the back of a torn card{]}} High, low, low, high, low, high,
high\ldots{} Forty-one highs and thirty-nine lows on this page. Over the
leaf, a door evening: bright, dark, dark, bright --- no better. It's a
coin. It's as good a coin as mine.

\textbf{Crab:} \emph{{[}sliding the blotter across to bury the left
column instead, and uncovering the right; the same eleven years, read
down the other side now{]}} And the right boxes.

\textbf{Achilles:} Low, high, high, high, low\ldots{} Another coin. Just
as fair and just as empty. If you had handed me either column alone, I'd
have said the workshop in Delft had sold you a very expensive way of
tossing a penny.

\textbf{Crab:} \emph{{[}lifting the blotter clean off, so that both
columns lie open to the lamp at once, side by side --- neither column
changed by a single mark from what it was under cover{]}} And now.

\textbf{Achilles:} \emph{{[}reading across{]}} High, high. Low, low.
Low, low. High, high. It's the same column twice. Line for line, lid
evenings and door evenings alike.

\emph{{[}Beyond the curtains the starlings have come in over the roofs
to roost, as they do at this hour. Any one bird, followed by eye, goes
every way by turns and nowhere in particular; the flock draws itself
out, folds, and pours down the sky in a single movement, like a cloth
shaken once.{]}}

\textbf{Anteater:} That is what I meant by an address. Ask any ant of
Aunt Hillary's where the colony is moving this week, and as far as I
have ever been able to tell you get a coin toss: this one is off east
with a crumb, that one west with a grub. Ask Aunt Hillary, and she says
``the south slope'' at once, and she is right. Mind, I don't say the
ants are not there --- I am extremely fond of ants --- only that they do
not hold the answer.

\textbf{Tortoise:} A box knows what an ant knows.

\emph{{[}The flock crosses what is left of the sun, and the room dims
for the length of a breath. When the light comes back it comes level,
along the wall over the mantel, and lies for a moment on the glass of
the print.{]}}

\textbf{Achilles:} \emph{{[}looking up from the twin columns to the
print on the Crab's wall --- Escher's \textup{Bond of Union}{]}} Two
heads --- and when you follow it round, only the one ribbon. I'd always
taken it for a portrait of two people.

\textbf{Crab:} More tea, anyone?

\textbf{Achilles:} No --- wait, this is much simpler than you're all
making it. The father and daughter in Delft write the note inside
beforehand: the same note in both boxes of a pair, high in this pair,
low in that, as the fancy takes them. And a second note for the side
door, the same in both. Each box alone looks like a coin, because nobody
told \emph{us} what they wrote; the pair always agrees, because it was
written to. There's your whole ledger --- and every box knows perfectly
well what it is going to play.

\textbf{Crab:} She said she didn't know.

\textbf{Achilles:} She needn't. The box knows.

\textbf{Anteater:} It is a respectable theory. I cannot find a line in
the ledger that contradicts it.

\textbf{Tortoise:} Nor can I, Achilles, and I shan't try. A theory like
that deserves better than an argument; it deserves a game. Bring it with
you another evening --- notes written inside, as many as you please,
composed any way you please --- and we shall play it against the boxes
and see which of them wins.

\textbf{Achilles:} You're being very gracious about losing.

\textbf{Tortoise:} I am always gracious.

\emph{{[}On the mantel the two clocks strike the hour, the one beginning
before the other has done, so that the strokes fall in between one
another and nobody in the room could have counted them. Then they go
back to their ticking, which has become a limp --- long, short, long,
short --- that changes if one listens to it.{]}}

\textbf{Crab:} There is a quicker way than a game. \emph{{[}He has the
right pearwood box in both claws, one claw-tip under the edge of the
lid.{]}} We all heard high. One lift, and we should know whether---

\textbf{Tortoise:} Whether what, Mr Crab? You know what it will play.

\textbf{Crab:} I know what the ledger says it will play.

\textbf{Tortoise:} And when you have lifted it you will own one more
line of the ledger, of which you have hundreds, and one fewer box that
can still sing, of which you have very few.

\textbf{Crab:} \emph{{[}setting the box down at the far corner of the
tray, shut, and keeping a claw beside it{]}} It is a hard thing, to own
something whose whole value lies in its not having been looked at.

\textbf{Achilles:} I'll play your game, Tortoise, and I'll win it. But
let me take the other side for a moment, purely for the sport. Suppose
it \emph{isn't} written in the boxes. Then where is it kept? The
agreement is there on the page; it has to be somewhere on the table.

\emph{{[}Tortoise reaches into her shell and brings out a piece of
string.{]}}

\begin{center}\rule{0.5\linewidth}{0.5pt}\end{center}

\hypertarget{third-challenge-the-string}{%
\section{\texorpdfstring{Third Challenge --- \emph{The
String}}{Third Challenge --- The String}}\label{third-challenge-the-string}}

\textbf{Tortoise:} May I have the walnut pair, Mr Crab? I shan't open
them. I only want them to stand somewhere.

\textbf{Crab:} String. Near my boxes. \emph{{[}He slides the walnut pair
across the cloth, shut, with both claws.{]}} If a lid so much as stirs,
Madam Tortoise---

\emph{{[}The lamps have been lit, one at either end of the table, and
everything on the cloth has two shadows. The walnut pair stand out in
the middle of it, off the tray for the first time since Delft. Behind
them the ebony pair have the tray almost to themselves: the lamplight
lies in the left lid and in the right, and there is still a wisp of
wood-wool caught in the hinge of the left one.{]}}

\textbf{Tortoise:} This is nothing but a piece of string. I keep it
because it is the plainest thing I know that can be in two places and
still be all one thing --- and your agreement, Achilles, has to be kept
somewhere that is in two places at once. Watch what the two ends of it
have to do with each other.

\textbf{Tortoise:} \emph{{[}she takes the loose straight string and lays
it on the cloth in a single hairpin --- one end set down at the left
walnut box, the string carried away from the pair, round in one soft
bend, and back again to rest its other end at the right walnut box: two
ends now at two boxes, joined behind by one unbroken length{]}} There is
a pair. Achilles, put your hand over one end and tell me what the other
is the end of.

\textbf{Achilles:} Of a string. From here I couldn't say more; it might
run anywhere.

\textbf{Tortoise:} And with your hand away?

\textbf{Achilles:} They're the two ends of the \emph{same} string. So
that's where it's kept? In the bend --- in the one part that isn't at
either box? Whereabouts in the bend? This inch, or that one?

\textbf{Tortoise:} \emph{{[}she pushes the bend further out across the
cloth, then draws it back close to the pair, then shoves it crooked off
to one side --- the bend wanders all over the table, but through every
move the two ends never leave their boxes and the string never once
comes apart{]}} Does it matter where I put it?

\textbf{Achilles:} No.~It only matters that it goes out and comes back.

\textbf{Tortoise:} Then follow the string, and let us see what a bend
will bear. \emph{{[}She takes the string up and lays it out longer this
time: out past the left walnut box, round in a first bend, back between
the two boxes, round in a second bend that faces the other way about,
and out again past the right one --- three legs and two bends now, out,
back, out, and still one string{]}}

\textbf{Achilles:} Two bends. Let me work it out. The first bend is a
pair: a coin toss on each leg, and the legs agreeing. The second bend
faces the other way about, so it's a pair being\ldots{} listened for? By
somebody who wants two legs to agree, which by my cards they'd do one
time in --- no, wait, the middle leg belongs to both bends, so I must
tally it twice, or else halve it. A half of a half of\ldots{} Tortoise,
I have three legs, two bends and four coins, and I've lost one of the
coins.

\textbf{Tortoise:} Take that end. I have this one. Pull.

\emph{{[}Each takes an end and draws it away from the other. The slack
runs out of the loop; the two bends slide toward one another, meet, and
pull through into nothing --- out and back and out drawing straight.
What was three legs and two bends is one taut line from Achilles's hand
to the Tortoise's, and the two walnut boxes it bent around a moment ago
sit untouched to either side. The kinks are gone; the string is the same
string, and just as long as it ever was.{]}}

\textbf{Achilles:} \ldots Oh. It's a string. It was only ever a string
with a kink in it. Out, back and out again is just --- \emph{out}. There
was nothing to tally.

\emph{{[}The draught under the door has dropped. The steam from the
spout of the hot-water pot, which has been leaning, doubling back on
itself and leaning again, stands up in a single thread, as if it had
never done anything else.{]}}

\textbf{Tortoise:} \emph{{[}laying it out again the other way about, the
first bend now facing down where it faced up --- the same S turned into
a Z{]}} And bent like this?

\textbf{Achilles:} \emph{{[}taking his end and pulling without being
asked; the two bends run out and cancel exactly as before, the string
coming straight{]}} Straight again. So a bend out and a bend back undo
each other, whichever comes first, and whatever goes in at one end of
all that travelling comes out at the other exactly as it went in.

\textbf{Crab:} You moved the cloth. The left walnut box has shifted by a
hair.

\textbf{Tortoise:} It has not opened.

\textbf{Crab:} It has been \emph{disturbed}.

\emph{{[}In the silence that follows this remark the clocks can be
heard. The limp has evened out: a tick from the left of the mantel, a
tock from the right, turn and turn about --- the sound of one clock,
made by two.{]}}

\textbf{Achilles:} But of course a bare string pulls straight; there's
nothing on it to catch. What if something were threaded on it?

\begin{center}\rule{0.5\linewidth}{0.5pt}\end{center}

\hypertarget{fourth-challenge-the-bead}{%
\section{\texorpdfstring{Fourth Challenge --- \emph{The
Bead}}{Fourth Challenge --- The Bead}}\label{fourth-challenge-the-bead}}

\textbf{Tortoise:} \emph{{[}bringing a bead out of her shell{]}}
Something like this. Look at it well before I thread it, because I shall
ask you about it afterwards.

\textbf{Tortoise:} It is only a passenger. It rides wherever the string
carries it and can change nothing about the string itself; so it will
show us honestly what a bend does to whatever is threaded on it ---
which is just what you wanted to know.

\textbf{Achilles:} A flat bead; it lies on the cloth like a little tile.
One end comes to a point. And there's a notch cut into one flank, like a
bite out of a biscuit. Lopsided twice over.

\textbf{Tortoise:} Twice over is what we shall need. \emph{{[}She bends
the string once round the left walnut box, so that two legs run back
from the bend side by side: a far leg, over nearer the window, and a
near leg, nearer the company. She slides the bead onto the far leg and
settles it face-up on the cloth, its point aimed along the string toward
the bend, its notched flank turned to face the window.{]}} Point towards
the bend. Notch towards the window. \emph{{[}On a napkin she draws the
bead exactly as it lies now: a long tile, a point at its left-hand end,
a bite out of the flank on the window side --- a fixed record of how it
went in.{]}} That is how it goes in. I am going to slide it along the
string, round the bend, and onto the near leg. How does it arrive?

\textbf{Achilles:} As it left. It's the same bead on the same string,
and you're doing nothing to it but moving it. Point to the left, notch
to the window.

\textbf{Tortoise:} \emph{{[}sliding the bead along the far leg, round
the bend behind the box, and out along the near leg --- flat on the
cloth the whole way, never once lifted from the table or turned over by
hand, only carried round the corner by the string it rides on{]}} It
went in point first. Look how it comes out.

\textbf{Achilles:} Point first still --- but pointing to the
\emph{right}. Well, naturally; it's travelling the other way now. Very
well, I was hasty. It's reversed: it arrives mirrored. That's the
looking-glass bead.

\textbf{Tortoise:} Mr Crab, might we have your mirror?

\textbf{Crab:} My mirror is a watchmaker's inspection glass from Le
Locle, silvered on its front face, and it has never been in the same
room as string. \emph{{[}He hands it over.{]}} By the edges.

\textbf{Tortoise:} \emph{{[}standing the silvered glass upright on the
napkin, its foot set across the nose of the drawn bead, so the inked
bead and its reflection meet point to point behind the glass{]}} Here is
the bead as it went in, and there behind it is your looking-glass bead.
Read it to me.

\textbf{Achilles:} In the glass the point is turned about: pointing
right. Good. And the notch is\ldots{} towards the window. \emph{{[}He
looks at the bead on the near leg.{]}} But the real one has its notch
towards \emph{me}. Away from the window. The point has changed sides and
the notch has changed sides too. In the glass only the point has.

\textbf{Tortoise:} So the glass shows you a different bead. There is no
such bead on my string.

\textbf{Achilles:} Then what did happen to it?

\textbf{Tortoise:} It never left the table. It went round a corner, flat
on the cloth, and a thing carried round a corner arrives facing the
other way: turned end for end. A half-turn, Achilles --- point and notch
go round together. To make the looking-glass bead I should have had to
pick it up and lay it down on its face, and I did not.

\textbf{Crab:} May I have the glass back now.

\emph{{[}The glass goes back into its wash-leather, by the edges. On the
rim of Achilles's saucer a ladybird, woken by the lamps from wherever
ladybirds spend the autumn, is walking: along the far side towards the
door, round the curve, and back along the near side towards the
fireplace. It has not, so far as it knows, turned round.{]}}

\textbf{Tortoise:} Now call the bead anything one might do to a box
before it is opened. Done on the far leg, to one box of the pair. I slid
it round, and I changed neither the string nor the bead by a hair: it is
the same table I started with. So whatever you do to one box of a pair,
you could as well have done --- turned half round --- to the other.

\textbf{Achilles:} And nobody holding the pair could tell which, because
there's only the one string. \emph{{[}He pushes the bead back round the
bend with a fingertip --- along the near leg, round, onto the far leg
--- and watches it turn end for end once more, point and notch swinging
together back to where they began{]}} Tortoise. The bead went round the
bend. Could the bend go round the bead?

\begin{center}\rule{0.5\linewidth}{0.5pt}\end{center}

\hypertarget{fifth-challenge-a-key-in-the-post}{%
\section{\texorpdfstring{Fifth Challenge --- \emph{A Key in the
Post}}{Fifth Challenge --- A Key in the Post}}\label{fifth-challenge-a-key-in-the-post}}

\textbf{Tortoise:} Mr Crab, you have a tuning key for these boxes.

\textbf{Crab:} I have \emph{the} tuning key. \emph{{[}He opens a case
lined in blue velvet.{]}} Cut by the father in Delft from a single piece
of drill-rod, with a bit as fine as a fish bone. A quarter-turn in the
socket at the back, before a box is opened, and a box that would have
sung high will sing low, and the other way about. My cousin in Ostend
has a pair from the same workshop, and has been begging the loan of it
for a year.

\textbf{Achilles:} Then post it to him.

\textbf{Crab:} \emph{Post} it. Achilles, the boxes came from Delft by
post, packed in wood-wool, and you could drop them down a well. The key
would not survive a sorting office. The key does not survive being
looked at sternly.

\textbf{Tortoise:} May I thread it? \emph{{[}She passes the string
through the bow of the key and draws it out straight on the cloth, the
key riding at the middle of one taut line: a note comes in along the
string at the left end, passes through the key, and leaves at the right
end --- a wire in, a wire out, and the key at work between{]}} Here is
the key as a thing one \emph{does}. A box's note comes along the string
from this end; the key turns it; it goes out at that end, changed. Too
delicate to travel. Now. \emph{{[}She takes the near end of the string
and swings it round until both ends point the same way: the one straight
line closes into a hairpin, and where a moment ago there was one wire in
and one wire out there are now two legs lying side by side --- one
running to the left walnut box with the key still upon it, the other
running bare to the right. The key has not moved a hair; but it sits on
a pair now, not on a through-line.{]}}

\textbf{Achilles:} That's the bead's picture again. A bend, two legs, a
box at the end of each --- and something sitting on one leg.

\textbf{Tortoise:} So what have I made?

\textbf{Achilles:} A pair of boxes, with the key already turned in one
of them. Let me see --- in the ledger that pair wouldn't read high-high
and low-low like the others. On a lid evening it would read high-low,
low-high, all the way down: never agreeing, and just as reliably. It
isn't a key any more. It's the key's whole effect, written out as the
way two boxes answer each other.

\textbf{Crab:} And a pair can go in the post.

\emph{{[}The lamplight, coming low across the floor, shows what daylight
hid: a track in the carpet from the door to the table and from the table
to the cabinet, the pile laid all one way, about the width of a crab
walking sideways. Years of errands made it. Anyone crossing the room now
goes that way without choosing to.{]}}

\textbf{Anteater:} A thing done and a thing kept: the same thing at two
levels. Aunt Hillary makes no distinction between a habit and a
monument, I may say. Her trails are both.

\textbf{Achilles:} \ldots Oh. Every key is a pair. Any key at all ---
anything one could do to a box --- bend the string behind it and it's a
pair of boxes that behave just so. And every pair is a key, if you can
bend it back?

\textbf{Tortoise:} \emph{{[}bending the near end back the way it came,
so the two legs open out into one line again, and pulling it taut: the
hairpin gone, the string straight, the key riding at the middle exactly
as before{]}} The key again. Nothing lost on the way.

\textbf{Crab:} And my cousin performs that last manoeuvre how,
precisely? He has no string. He has a parcel from me and a breakfast
table.

\textbf{Tortoise:} With a pair of his own, a good ear, and some luck.
How much luck, I think Dr Anteater is about to tell us.

\textbf{Anteater:} Since you speak of parcels. That reminds me of two
friends of Aunt Hillary's who never met.

\begin{center}\rule{0.5\linewidth}{0.5pt}\end{center}

\hypertarget{sixth-challenge-the-relay}{%
\section{\texorpdfstring{Sixth Challenge --- \emph{The
Relay}}{Sixth Challenge --- The Relay}}\label{sixth-challenge-the-relay}}

\textbf{Anteater:} Could I have the whole tray, Mr Crab, set out in a
row. It is a story with a good deal of furniture.

\textbf{Crab:} \emph{{[}lifting the boxes off the tray one at a time and
standing them in a single line along the far side of the cloth, left to
right, the three pairs opened out and set end to end into one long
row{]}} Of course. The pearwood: the spent one, and beside it its
partner --- \emph{shut}, and I am watching everybody. The walnut pair.
The ebony pair, which nobody has so much as breathed on.

\textbf{Anteater:} Excellent. At the left-hand end of the row lives
Alice, a friend of Aunt Hillary's from the east side of the hill. At the
right-hand end, a long way off, lives Bob, a friend from the west. They
have never met, never corresponded, and never owned boxes from the same
bench. Alice has the left pearwood box. Bob has the right ebony one.

\textbf{Achilles:} Careful --- Alice's is spent.

\textbf{Anteater:} Kindly recall that it is a story; there, hers is not.
Stories are kinder than evenings.

\textbf{Achilles:} Even so --- the boxes in between?

\textbf{Anteater:} Between them stand two relay stations. Alice's pair
was made at a bench near her, and its other box sent along the road to
the first station. The walnut pair was made between the stations: one
box to the first, one to the second. The ebony pair was made near Bob;
one box went back up the road to the second station, and Bob kept the
other.

\emph{{[}The Tortoise slips the tuning key off the string and hands it
back to the Crab, who lays it in its open case. The whole row now stands
ready along the far edge of the cloth, every box shut and silent; the
freed string lies loose and bare in front of them, nothing threaded on
it and nothing yet bent.{]}}

\textbf{Tortoise:} \emph{{[}laying the string on the cloth in front of
the row as he speaks, a leg of string for each box, like a garden path
to each door: up to the left pearwood box, round in a bend, and back
from the right pearwood box{]}} Let Bob draw a path to the first pair;
one bend. \emph{{[}The string turns about at the near side and runs up
to the left walnut box; round in a bend; back from the right walnut
box.{]}} Let Bob draw on to the first station; the next pair.
\emph{{[}It turns about again at the near side and runs up to the left
ebony box; round; and back from the right ebony box, where the string
ends.{]}} Let Bob draw on to the second station; the last pair.

\emph{{[}The string now runs the length of the row: up to the first pair
and round, back down and up to the second pair and round, back down and
up to the third pair and round --- three bends reaching up to the boxes,
two turns doubling back between them, and two free ends left lying at
the near edge, one below the far-left box and one below the
far-right.{]}}

\textbf{Achilles:} So the first station is holding the right pearwood
box and the left walnut one --- boxes from different benches, with
nothing to say to each other. Two pennies. And Alice and Bob have less
than nothing: between her box and his there are all those separate
bends, and a gap at each station.

\textbf{Tortoise:} Are they gaps? Follow the string.

\textbf{Achilles:} They're\ldots{} turns. The string turns about in
front of each station and goes out again. Out, back, out, back, out. But
a turn of string on a tablecloth isn't something a relay station can
\emph{do}.

\textbf{Anteater:} No.~This is what it can do. Each station has an
instrument into which its two boxes are set side by side, and which puts
to them a single question: not ``what does each of you play?'' but ``do
you two play as one?'' Most evenings it gives one of several other
answers. Now and then it gives a certain chord; and on the evening a
station hears that chord, its two boxes have been joined exactly as
Madam Tortoise's string is turned.

\textbf{Achilles:} How often is now and then?

\textbf{Anteater:} One evening in four.

\textbf{Achilles:} Then for both stations together: one in four, and one
in four of that. One evening in sixteen. Your friends will need
patience.

\textbf{Anteater:} Aunt Hillary's friends have little else. On the other
evenings something else happens, which is not a failure; but that is
another evening's story. Tonight, both stations hear the chord.

\emph{{[}Along the cornice runs the wire from the street-door pull, on
its way to the kitchen by way of two small brass cranks, one at each
corner of the room. Somebody at the street door pulls. The wire
twitches, each crank nods once and is still, and a long way off, beyond
the baize door, a bell rings. Nobody gets up; it is only the evening
post.{]}}

\textbf{Tortoise:} Let Bob draw both ends in. Achilles, one finger on
the bend in front of the walnut pair --- lightly. Mr Crab, I am not
going to touch a box. \emph{{[}She takes the two free ends of the
string, the one below Alice's box and the one below Bob's, and pulls
them towards her. The turns at the stations run out; the bends draw into
one another. She spreads what is left with a claw: one long bend, up to
the left pearwood box, across the front of the whole row, and back from
the right ebony box.{]}}

\emph{{[}Nothing on the table has been opened or moved: the same boxes
stand in the same row. Only the string has changed --- every turn at the
stations pulled out, and the two end legs, still resting at the far-left
and far-right boxes, joined now by one unbroken bend across the whole
front of the row.{]}}

\textbf{Achilles:} One bend. From Alice's box to Bob's. But that's a
\emph{pair} --- it's the hairpin you laid for the walnut boxes, only
wider. Her box and his would fill a line in the ledger like any other: a
coin on the left, a coin on the right, and the same coin.

\textbf{Anteater:} They never shared a bench.

\textbf{Achilles:} Then the agreement was never in the boxes. That's
what I asked you, Tortoise --- where is it kept? --- and it isn't kept
in any box, or at any bench. It's in how they are joined. And joining
can be handed along: the stations passed it down the row, and neither of
them kept any.

\textbf{Tortoise:} What is left at the stations?

\textbf{Achilles:} Nothing. Spent boxes, and the memory of a chord.
\emph{{[}He looks along the row.{]}} And Alice and Bob at the two ends,
each holding half of something --- and neither half is anything at all
without the other.

\textbf{Anteater:} \emph{{[}gathering himself, and folding his
napkin{]}} Well. It is a long way round the hill, and Aunt Hillary will
want her regards carried home. You have both been very patient with an
old creature and all his furniture.

\textbf{Tortoise:} Stay for the last of the tea, Doctor. There is one
thing still on the table, and I think Mr Crab would rather we saw it out
together than alone afterwards.

\emph{{[}The last of the tea goes round, the pot lighter now, and the
Anteater, half risen, settles again. Nobody hurries. From the mantel
come the tick and the tock, so even by now that a visitor would look for
one clock and find two.{]}}

\textbf{Crab:} \emph{{[}after a silence{]}} That is the most
uncomfortable account of my evening I have yet heard. I have been
sitting beside half of something since before the tea. \emph{{[}He takes
his claw away from the right pearwood box.{]}} Madam Tortoise. You are
the only person in this room who has not wanted to. You had better do
it.

\textbf{Tortoise:} By the lid, I think; as Achilles did.

\emph{{[}The right pearwood box has sat shut and full since the first
lines of the evening, its note folded up inside it. The Tortoise catches
the lid and lifts it --- the box does not move, only the little hinged
lid swings up, as its partner's did --- and the note springs out: a
single high tone, the same as the first, clear as a struck glass, and
gone. The Crab writes one word in the right-hand column of the ledger,
beside the word that has waited there alone since the evening began; and
now the two words stand side by side, the same word twice, one line of a
pair completed a whole evening late.{]}}

\textbf{Achilles:} High. Of course it's high. But was it in there all
along?

\textbf{Tortoise:} Bring that to the game.

\begin{center}\rule{0.5\linewidth}{0.5pt}\end{center}

\emph{{[}The note has gone where the first one went. On the cloth the
row stands as the Crab set it --- pearwood and pearwood, walnut and
walnut, ebony and ebony --- and before it lies the string in its one
long bend, an end at either end of the row. In the ledger the ink of the
second word is drying beside the first. The lamps burn level, the steam
has stopped, and on the mantel the two clocks go on handing the time
back and forth between them; over them the print has gone back into the
wall, a dark frame with a little lamplight on its glass; it is the long
part of the night. Nobody at the table reaches for anything, and for the
moment nothing on it wants looking into.{]}}

\hypertarget{chapter-6-entanglement}{%
\chapter{Chapter 6: Entanglement}\label{chapter-6-entanglement}}

\emph{--- Obbligato for Six Music Boxes ---}

\emph{Part II: Wires Together}

\emph{Escher: ``Bond of Union'' --- two heads formed from one continuous
ribbon}

\begin{center}\rule{0.5\linewidth}{0.5pt}\end{center}

Can two things, far apart, share a fact that neither of them holds?

Toss two coins in two different towns and compare notes afterwards. Each
list of results is noise, and the two lists have nothing to say to each
other. Now imagine a stranger pair of coins. Each one, watched on its
own, is still perfectly fair: heads and tails, half and half, no pattern
that any statistician could find. But when the two lists are laid side
by side they are the same list, line for line.

The everyday explanation comes at once: somebody fixed the coins in
advance. Perhaps that is right. This chapter will not decide it, and we
ask the reader to keep the suspicion alive, because the next chapter has
plans for it. What this chapter does is look at how quantum theory
itself describes such a pair --- and the description is odd in a precise
way. The theory has a complete account of the pair. It has no account at
all of either member. It is as if a duet had been written down, in full,
with no part for the first voice and no part for the second, and yet
neither voice could be left out.

So where is the agreement kept, if not in the things that agree? On the
blackboard the answer will be a sum that refuses to factorize, and we
shall prove that it refuses. In the sketchbook the answer is more
surprising, and it is the reason this chapter is a turning point in the
book: the agreement is kept in a \emph{bend}. A wire that goes out and
comes back is, all by itself, the pair we are looking for. Once a wire
may bend, three things can be done with it. It can be pulled straight. A
box can be slid along it, round the corner, and we shall watch carefully
how the box arrives. And a box can be bent into a state, so that
something one \emph{does} becomes something one \emph{has}.

Those three moves are the whole chapter. At the end we thread one wire
back and forth a few times, pull, and find that two parties who never
met have been joined.

\begin{center}\rule{0.5\linewidth}{0.5pt}\end{center}

\hypertarget{two-boxes-one-description}{%
\section{§6.1 Two Boxes, One
Description}\label{two-boxes-one-description}}

In Chapter 4 we learned to put two wires side by side. If the top wire
carries \(\ket{\alpha}=\alpha_0\ket0+\alpha_1\ket1\) and the bottom wire
carries \(\ket{\beta}=\beta_0\ket0+\beta_1\ket1\), the pair carries the
tensor product

\[\ket{\alpha}\otimes\ket{\beta} \;=\; \alpha_0\beta_0\ket{00}+\alpha_0\beta_1\ket{01}+\alpha_1\beta_0\ket{10}+\alpha_1\beta_1\ket{11},\]

where, as always in this book, the first symbol inside a ket belongs to
the top wire. The two wires of Deutsch's algorithm in Chapter 5 were
always of this kind: from start to finish the answer wire stayed in
\(\ket-\), a factor by itself, and the kicked sign was only a number in
front. But the pair of wires is a four-dimensional space, and \emph{any}
vector in it is a legitimate state:

\[\ket{\psi} \;=\; c_{00}\ket{00}+c_{01}\ket{01}+c_{10}\ket{10}+c_{11}\ket{11}.\]

The question of this section is whether every such \(\ket\psi\) is
secretly an \(\ket\alpha\otimes\ket\beta\). It would be a comfort if it
were: a description of the pair would then be nothing more than a
description of each wire.

\textbf{On the blackboard.} A state of two wires is a \textbf{product
state} if it can be written as \(\ket\alpha\otimes\ket\beta\) for some
one-wire states \(\ket\alpha\) and \(\ket\beta\). A state that cannot be
so written is \textbf{entangled}. Here is the state that will occupy us
for the rest of the chapter, the \textbf{Bell state}

\[\ket{\Phi^+} \;=\; \frac{1}{\sqrt2}\bigl(\ket{00}+\ket{11}\bigr).\]

Let us try to factorize it. We need \(\alpha_0\beta_0=\tfrac1{\sqrt2}\)
and \(\alpha_1\beta_1=\tfrac1{\sqrt2}\), so none of the four numbers
\(\alpha_0,\alpha_1,\beta_0,\beta_1\) is zero. But we also need
\(\alpha_0\beta_1=0\). That is a contradiction, and no cleverness with
complex numbers gets round it. The Bell state is entangled.

It is not an exotic state. The two gates of Chapters 3 and 4 make it
from the plainest possible start: apply the Hadamard gate \(H\) to the
top wire of \(\ket{00}\), giving \(\tfrac1{\sqrt2}(\ket{00}+\ket{10})\),
then apply CNOT, which flips the bottom wire exactly when the top wire
reads \(1\):

\[\mathrm{CNOT}\,(H\otimes I)\,\ket{00} \;=\; \frac{1}{\sqrt2}\bigl(\ket{00}+\ket{11}\bigr) \;=\; \ket{\Phi^+}.\]

\textbf{In the sketchbook.} A product state is two separate state boxes,
one on each wire. Nothing joins the top half of the page to the bottom
half: the picture is \emph{disconnected}. The Bell state, by contrast,
has to be drawn as a single state box with two wires leaving it. We have
no new notation for this yet; it is an ordinary box of the kind we have
drawn since Chapter 4, labelled \(\Phi^+\). Here is the preparation
above, as a picture equation:

\[\begin{tikzpicture}[sd]
  \draw[wire] (0,2.2) -- (10,2.2);
  \draw[wire] (0,0) -- (10,0);
  \node[sdstate] at (0,2.2) {$0$};
  \node[sdstate] at (0,0) {$0$};
  \node[sdbox] at (3,2.2) {$H$};
  \node[sdbox, minimum height=4.1em] at (7,1.1) {$\mathrm{CNOT}$};
\end{tikzpicture}
\;=\;
\begin{tikzpicture}[sd]
  \draw[wire] (0,2.2) -- (2.5,2.2);
  \draw[wire] (0,0) -- (2.5,0);
  \node[sdstate, minimum height=4.1em] at (0,1.1) {$\Phi^+$};
\end{tikzpicture}\]

and here is the claim that it is entangled, for every choice of
\(\alpha\) and \(\beta\):

\[\begin{tikzpicture}[sd]
  \draw[wire] (0,2.2) -- (2.5,2.2);
  \draw[wire] (0,0) -- (2.5,0);
  \node[sdstate, minimum height=4.1em] at (0,1.1) {$\Phi^+$};
\end{tikzpicture}
\;\neq\;
\begin{tikzpicture}[sd]
  \draw[wire] (0,2.2) -- (2.5,2.2);
  \draw[wire] (0,0) -- (2.5,0);
  \node[sdstate] at (0,2.2) {$\alpha$};
  \node[sdstate] at (0,0) {$\beta$};
\end{tikzpicture}\]

\textbf{\emph{Why this definition.}} We define entanglement negatively
--- \emph{not} a product --- because the product states are the ones we
understand: they are exactly the states for which ``the state of the top
wire'' means something. Entanglement is what is left when that phrase
stops making sense.

\hypertarget{a-test-that-always-works}{%
\subsection{A test that always works}\label{a-test-that-always-works}}

Our argument for \(\ket{\Phi^+}\) used its two zeros. What about a state
with no zeros at all, such as

\[\ket{\chi}=\tfrac12\bigl(\ket{00}+\ket{01}+\ket{10}-\ket{11}\bigr)\,?\]

One can stare at this for some time. It would be good to have a test
that needs no staring. Arrange the four amplitudes of \(\ket\psi\) in a
square, top wire indexing the rows:

\[C \;=\; \begin{pmatrix} c_{00} & c_{01}\\ c_{10} & c_{11}\end{pmatrix}.\]

\begin{quote}
\textbf{Theorem 6.1 (the product test).} A two-wire state
\(\ket\psi=\sum_{ij}c_{ij}\ket{ij}\) is a product state if and only if
\(\det C = c_{00}c_{11}-c_{01}c_{10}=0\).
\end{quote}

\emph{Proof idea.} For a product state the square is a multiplication
table --- row labels \(\alpha_i\), column labels \(\beta_j\) --- and the
two rows of a multiplication table are proportional. The determinant of
a \(2\times2\) matrix is precisely the measure of whether its rows are
proportional.

\emph{Proof sketch.} If \(\ket\psi=\ket\alpha\otimes\ket\beta\) then
\(c_{ij}=\alpha_i\beta_j\) and
\(\det C=\alpha_0\beta_0\alpha_1\beta_1-\alpha_0\beta_1\alpha_1\beta_0=0\).
Conversely, suppose \(\det C=0\). The state is not the zero vector, so
some row of \(C\) is not zero; say it is the top row \((c_{00},c_{01})\)
(otherwise exchange the roles of the rows). A vanishing determinant says
the bottom row is a multiple of it:
\((c_{10},c_{11})=\lambda\,(c_{00},c_{01})\). Then
\(\ket\alpha=\ket0+\lambda\ket1\) and
\(\ket\beta=c_{00}\ket0+c_{01}\ket1\) give
\(\ket\alpha\otimes\ket\beta=\ket\psi\). \(\;\blacksquare\)

For \(\ket{\Phi^+}\) the determinant is
\(\tfrac1{\sqrt2}\cdot\tfrac1{\sqrt2}-0=\tfrac12\): entangled, as we
found by hand. For \(\ket\chi\) it is
\(\tfrac14\cdot(-1)-\tfrac14=-\tfrac12\): entangled too, which was far
from evident. And for \(\tfrac12(\ket{00}+\ket{01}+\ket{10}+\ket{11})\),
which differs from \(\ket\chi\) by one sign, it is \(0\): that state is
just \(\ket{+}\otimes\ket{+}\). See \texttt{Code/Ch6/NotAProduct.py}.

Notice that the test is pure blackboard. The sketchbook can \emph{say}
``this box does not come apart'', but at this stage it has no way to see
why. We shall return to this determinant in §6.5 and find that it was
the determinant of something.

\begin{quote}
\textbf{Feynman's Notebook.}

In May 1935 Albert Einstein, Boris Podolsky and Nathan Rosen published a
four-page paper in the \emph{Physical Review} with a question for a
title: ``Can Quantum-Mechanical Description of Physical Reality Be
Considered Complete?'' Their example was a pair of particles that had
interacted and flown apart, described by a state of just the kind above
--- one description for the pair, none for the members. They took this
as a sign that the theory was leaving something out.

Erwin Schrödinger read the paper and answered within months, in German
and in English. He gave the phenomenon its name ---
\emph{Verschränkung}, entanglement --- and in a paper communicated to
the Cambridge Philosophical Society that year he wrote of it: ``I would
not call that \emph{one} but rather \emph{the} characteristic trait of
quantum mechanics, the one that enforces its entire departure from
classical lines of thought.''

Both reactions were reasonable. It took nearly thirty years before
anyone found a way to put the disagreement to an experiment; that is
Chapter 7.
\end{quote}

\begin{center}\rule{0.5\linewidth}{0.5pt}\end{center}

\hypertarget{random-parts-definite-whole}{%
\section{§6.2 Random Parts, Definite
Whole}\label{random-parts-definite-whole}}

If the pair has a description and the parts have none, what happens when
somebody looks at a part?

\textbf{On the blackboard.} Recall the rule of Chapter 4 for looking at
one wire of several. To ask the top wire ``are you \(\ket b\)?'', apply
the effect \(\bra b\) to the top wire and leave the bottom wire alone;
that is the matrix \(\bra b\otimes I\). What comes out is an
unnormalized state of the bottom wire. Its squared length is the
probability of the answer ``yes'', and its direction is the state the
bottom wire is left in.

Ask the top half of a Bell pair whether it is \(\ket0\):

\[(\bra0\otimes I)\,\ket{\Phi^+} \;=\; \frac1{\sqrt2}\bigl(\braket{0|0}\ket0+\braket{0|1}\ket1\bigr) \;=\; \frac1{\sqrt2}\,\ket0 .\]

The squared length is \(\tfrac12\). The same calculation with \(\bra1\)
gives \(\tfrac1{\sqrt2}\ket1\), again with squared length \(\tfrac12\).
So the top wire, asked for its bit, answers like a fair coin --- and
whichever answer it gives, the bottom wire is left in that same state,
so the same question put to the bottom wire is now certain to get the
same answer. By the symmetry of \(\ket{00}+\ket{11}\) the story is
identical with the wires exchanged. In a table, the joint outcomes
\(00,01,10,11\) have probabilities \(\tfrac12,0,0,\tfrac12\).

Perhaps that is special to the basis in which we wrote the state down.
Try \(\ket{\pm}=\tfrac1{\sqrt2}(\ket0\pm\ket1)\):

\[(\bra{+}\otimes I)\,\ket{\Phi^+} \;=\; \frac1{\sqrt2}\cdot\frac1{\sqrt2}\bigl(\ket0+\ket1\bigr) \;=\; \frac1{\sqrt2}\,\ket{+},\]

\[(\bra{-}\otimes I)\,\ket{\Phi^+} \;=\; \frac1{\sqrt2}\cdot\frac1{\sqrt2}\bigl(\ket0-\ket1\bigr) \;=\; \frac1{\sqrt2}\,\ket{-}.\]

A fair coin again, and agreement again. In fact
\(\ket{\Phi^+}=\tfrac1{\sqrt2}(\ket{++}+\ket{--})\), as one checks by
multiplying out: the state looks the same in the new basis as in the
old.

\textbf{In the sketchbook.} Plug an effect, labelled \(b\), into the top
wire of the \(\Phi^+\) box. One wire is still open, so what remains is a
state of one wire. The calculations above say which:

\[\begin{tikzpicture}[sd]
  \draw[wire] (0,1.5) -- (2.8,1.5);
  \draw[wire] (0,0) -- (4,0);
  \node[sdstate, minimum height=3.4em] at (0,0.75) {$\Phi^+$};
  \node[sdeffect] at (2.8,1.5) {$b$};
\end{tikzpicture}
\;=\;\frac{1}{\sqrt2}\;\;
\begin{tikzpicture}[sd]
  \draw[wire] (0,0) -- (2.8,0);
  \node[sdstate] at (0,0) {$b$};
\end{tikzpicture}
\qquad \text{for } b\in\{0,1,+,-\}.\]

As a picture this is a bare report of the algebra. Why should an effect
on the top wire turn up as a state on the bottom wire? The sketchbook
cannot yet say. It will in §6.4, where this equation becomes a corollary
of the chapter's main theorem: the effect has slid round a corner.

\begin{quote}
\textbf{Dirac's Blackboard.}

\emph{Every basis, and a conjugate.} Let \(\ket b=b_0\ket0+b_1\ket1\) be
any unit vector, so \(\bra b=\bar b_0\bra0+\bar b_1\bra1\), the bar
denoting the complex conjugate. Then

\[(\bra b\otimes I)\ket{\Phi^+}=\tfrac1{\sqrt2}\bigl(\bar b_0\ket0+\bar b_1\ket1\bigr)=\tfrac1{\sqrt2}\ket{b^*},\]

where \(\ket{b^*}\) is \(\ket b\) with every amplitude conjugated. Its
squared length is \(\tfrac12\bigl(|b_0|^2+|b_1|^2\bigr)=\tfrac12\).
\textbf{Every question one can put to one half of a Bell pair is
answered by a fair coin.}

The other half is left in \(\ket{b^*}\), not \(\ket b\). For a basis
with real amplitudes --- \(\ket0,\ket1\) or \(\ket\pm\) --- there is no
difference, and the halves agree. For
\(\ket{b}=\tfrac1{\sqrt2}(\ket0+i\ket1)\) we get
\(\ket{b^*}=\tfrac1{\sqrt2}(\ket0-i\ket1)\), which is orthogonal to
\(\ket b\): asked \emph{that} question, the two halves always disagree.
Perfectly correlated still; only the labels are swapped. The conjugate
is not a nuisance. It is the first footprint of the transpose of §6.4.
See \texttt{Code/Ch6/Marginals.py}.
\end{quote}

So each part is as random as anything can be, in every direction one
looks, and the whole is as definite as anything can be: the two answers
are locked together. That is what the title of this section means, and
it is what the opening page described.

A caution, in plain words. Nothing here shows that the answers were
\emph{not} written into the two halves beforehand. Two identical sealed
notes, prepared at the source, reproduce both calculations above: each
note read alone looks like a coin toss, and the two always agree. We
have tested very little so far. The matter is taken up in Chapter 7,
with sharper tools; until then the hidden-note theory stands.

\textbf{\emph{Why this way of computing.}} We compute the statistics of
a part directly from amplitudes --- apply an effect, take a squared
length --- because at this point in the book there is no object that
could be called ``the state of the top wire''. Chapter 9 constructs one.
It will not be a ket.

\begin{center}\rule{0.5\linewidth}{0.5pt}\end{center}

\hypertarget{the-cup-the-cap-and-the-yank}{%
\section{§6.3 The Cup, the Cap, and the
Yank}\label{the-cup-the-cap-and-the-yank}}

We have a state that is one thing on two wires. Let us draw it as what
it looks like.

\textbf{In the sketchbook.} Take a wire and bend it back on itself, so
that both of its ends point to the right, like the letter C. Read from
left to right, nothing comes in and two wires go out: by the grammar of
Chapter 4 it is a \emph{state on two wires}. We call it the
\textbf{cup}. Now bend a wire the other way, so that both ends point to
the left, like a reversed C. Two wires come in and nothing goes out: an
\emph{effect on two wires}. We call it the \textbf{cap}.

\[\text{cup}\;=\;
\begin{tikzpicture}[sd]
  \sdcup{0}{1.5}{0}
  \draw[wire] (0,1.5) -- (2,1.5);
  \draw[wire] (0,0) -- (2,0);
\end{tikzpicture}
\qquad\qquad
\text{cap}\;=\;
\begin{tikzpicture}[sd]
  \draw[wire] (0,1.5) -- (2,1.5);
  \draw[wire] (0,0) -- (2,0);
  \sdcap{2}{1.5}{0}
\end{tikzpicture}\]

\textbf{On the blackboard.} We \emph{define}

\[\text{cup} \;=\; \ket{00}+\ket{11}, \qquad\qquad \text{cap} \;=\; \bra{00}+\bra{11}.\]

The cup is a column of four numbers, \((1,0,0,1)\); the cap is the same
four numbers as a row. Look closely: there is no \(\tfrac1{\sqrt2}\).
\textbf{The cup is not normalized.} Its length is \(\sqrt2\), and its
relation to the Bell state is

\[\ket{\Phi^+} \;=\; \frac{1}{\sqrt2}\,\text{cup}, \qquad\qquad \bra{\Phi^+} \;=\; \frac1{\sqrt2}\,\text{cap}.\]

This convention holds for the whole book, and we shall repeat it
wherever a cup and a Bell state stand close enough to be confused. In a
picture, the Bell state is a cup with the number \(\tfrac1{\sqrt2}\)
written beside it.

The cap deserves an honest word at once. As an effect it is
\(\sqrt2\,\bra{\Phi^+}\): the number \(\sqrt2\) times the answer ``yes''
to the question ``are these two wires in the state \(\ket{\Phi^+}\)?''
That is one of four possible answers to a two-wire question (the other
three are for Chapter 10), and nobody can \emph{make} a question come
out ``yes''. A diagram with a cap in it describes what happens
\textbf{when} that outcome occurs --- it is \emph{postselected}. With
that understood, caps are perfectly good mathematics, and we shall use
them freely and count their cost at the end.

\hypertarget{an-example-that-wants-to-be-a-rule}{%
\subsection{An example that wants to be a
rule}\label{an-example-that-wants-to-be-a-rule}}

Put a cup and a cap together so that they share one wire. Draw three
rows. A wire comes in on the top row; a cup sits on the middle and
bottom rows; a cap then joins the top and middle rows; the bottom row
goes out. The wire travels right, doubles back, and goes right again, in
the shape of a Z. As a matrix this is
\((\text{cap}\otimes I)(I\otimes\text{cup})\). Feed it \(\ket0\):

\[\begin{aligned}
\ket0 \;&\xrightarrow{\;I\otimes\text{cup}\;}\; \ket{000}+\ket{011}\\
 &\xrightarrow{\;\text{cap}\otimes I\;}\; \bigl(\braket{00|00}+\braket{11|00}\bigr)\ket0+\bigl(\braket{00|01}+\braket{11|01}\bigr)\ket1 \;=\; \ket0,
\end{aligned}\]

where the cap has eaten the first two symbols of each ket. Of four
brackets one survives. Feed it \(\ket1\) and, in the same way, \(\ket1\)
comes out. A bent wire that returns every basis state unchanged is no
different from a wire that was never bent.

\begin{quote}
\textbf{Theorem 6.2 (yanking).} A wire bent into a Z, or into an S, is a
straight wire:

\[(\text{cap}\otimes I)(I\otimes\text{cup}) \;=\; I \;=\; (I\otimes\text{cap})(\text{cup}\otimes I),\]

exactly.
\end{quote}

\[\begin{tikzpicture}[sd]
  \draw[wire] (0,3) -- (3.5,3);
  \sdcap{3.5}{3}{1.5}
  \draw[wire] (1.5,1.5) -- (3.5,1.5);
  \sdcup{1.5}{1.5}{0}
  \draw[wire] (1.5,0) -- (5,0);
\end{tikzpicture}
\;=\;
\begin{tikzpicture}[sd]
  \draw[wire] (0,0) -- (3,0);
\end{tikzpicture}
\;=\;
\begin{tikzpicture}[sd]
  \draw[wire] (0,0) -- (3.5,0);
  \sdcap{3.5}{0}{1.5}
  \draw[wire] (1.5,1.5) -- (3.5,1.5);
  \sdcup{1.5}{1.5}{3}
  \draw[wire] (1.5,3) -- (5,3);
\end{tikzpicture}\]

\emph{Proof idea.} The cup says ``my two ends carry the same symbol'';
the cap says ``I pass only if my two ends carry the same symbol''. Chain
them and the symbol at the far end must equal the symbol at the near
end. That is the identity.

\emph{Proof sketch.} For the Z-bend on the left, take a general input
\(a_0\ket0+a_1\ket1\). Then \(I\otimes\text{cup}\) produces
\(\sum_{i,k}a_i\ket{i\,k\,k}\), four terms. The cap sends \(\ket{i\,k}\)
to \(\braket{00|ik}+\braket{11|ik}\), which is \(1\) if \(i=k\) and
\(0\) otherwise; two terms survive, \(a_0\ket0+a_1\ket1\). For the
S-bend on the right the same computation runs with the rows in the other
order. \(\;\blacksquare\) See \texttt{Code/Ch6/Yanking.py}.

This is the chapter's first \textbf{rewrite rule}: wherever a wire makes
an S or a Z, we may pull it straight, and wherever a wire is straight,
we may put a kink in it. The equation is exact --- the same matrix on
both sides, number for number --- and so is every other picture equation
in this chapter.

\textbf{\emph{Why this definition.}} Why no \(\tfrac1{\sqrt2}\) in the
cup? Try it with one. Replace cup by \(\ket{\Phi^+}\) and cap by
\(\bra{\Phi^+}\) and the same computation gives

\[(\bra{\Phi^+}\otimes I)(I\otimes\ket{\Phi^+}) \;=\; \tfrac1{\sqrt2}\cdot\tfrac1{\sqrt2}\;I \;=\; \tfrac12\,I .\]

A wire with a kink in it would be \emph{half} a wire, and every picture
would have to carry a tally of its kinks. We prefer that a bent wire be
a wire. The price is that the cup is not a unit vector, and we pay it by
remembering one number, \(\tfrac1{\sqrt2}\), at the moments when
probabilities are wanted.

\begin{quote}
\textbf{Coecke's Sketchbook.}

\emph{Where the bent wire comes from.} In 1971 Roger Penrose published
``Applications of negative dimensional tensors'', setting out a way of
calculating with tensors by drawing them: a blob for each tensor, a line
for each index, and joined lines for summed indices. He had, by his own
account, used it privately for years. In that notation a bent line is an
old friend --- it is how an index is raised or lowered, how an ``in''
becomes an ``out''.

Samson Abramsky and Bob Coecke, computer scientists at Oxford, read the
same bent line in quantum terms: the bend that turns an input into an
output \emph{is} the Bell state, and the yanking equation is the
skeleton of protocols such as teleportation. The bent wire is Penrose's;
the quantum-process reading that unfolds from here --- the calculus
whose name this Sketchbook now carries --- is Coecke's, with Abramsky.

One caution, to carry through the book: a picture rule is only as good
as the algebra that licenses it. We may pull a wire straight because of
Theorem 6.2, and for no other reason.
\end{quote}

\begin{center}\rule{0.5\linewidth}{0.5pt}\end{center}

\hypertarget{sliding-a-box-around-the-bend}{%
\section{§6.4 Sliding a Box Around the
Bend}\label{sliding-a-box-around-the-bend}}

A bare wire can be pulled straight. What if there is a box on it?

Let us experiment, with a matrix chosen to have no helpful symmetries:

\[A \;=\; \begin{pmatrix} 1 & 2\\ i & 3\end{pmatrix}.\]

It is not a gate --- it is not unitary --- and it need not be: a box may
hold any linear map. Apply \(A\) to the top leg of a cup. Since
\(A\ket0=\ket0+i\ket1\) and \(A\ket1=2\ket0+3\ket1\) (the columns of
\(A\)),

\[(A\otimes I)\,\text{cup} \;=\; A\ket0\otimes\ket0+A\ket1\otimes\ket1 \;=\; \ket{00}+2\ket{01}+i\ket{10}+3\ket{11}.\]

Now the question. The two legs of a cup are one wire. Could the same
state have been made by a box on the \emph{bottom} leg? We want a matrix
\(B\) with
\((I\otimes B)\,\text{cup}=\ket0\otimes B\ket0+\ket1\otimes B\ket1\)
equal to the state above. Collect the terms that begin with \(\ket0\)
and those that begin with \(\ket1\):

\[\ket0\otimes\bigl(\ket0+2\ket1\bigr) \;+\; \ket1\otimes\bigl(i\ket0+3\ket1\bigr),\]

so \(B\ket0=\ket0+2\ket1\) and \(B\ket1=i\ket0+3\ket1\). These are the
columns of \(B\):

\[B \;=\; \begin{pmatrix} 1 & i\\ 2 & 3\end{pmatrix}.\]

Compare with the three usual suspects:

\[A^{T}=\begin{pmatrix} 1 & i\\ 2 & 3\end{pmatrix},\qquad
A^{\dagger}=\begin{pmatrix} 1 & -i\\ 2 & 3\end{pmatrix},\qquad
\bar A=\begin{pmatrix} 1 & 2\\ -i & 3\end{pmatrix}.\]

It is the \textbf{transpose} --- rows and columns exchanged, and
\emph{no} complex conjugation. The example suggests a rule: a box pushed
round the bend of a cup arrives transposed. Before we can state that as
a picture equation, the sketchbook needs to know what a transposed box
looks like.

\textbf{In the sketchbook.} Take the box \(A\) and bend \emph{both} of
its wires round. Draw three rows. The input comes in on the top row and
runs to the right, where a cap joins the top row to the middle row. The
box \(A\) sits on the middle row. On the left, a cup joins the middle
row to the bottom row, and the output leaves on the bottom row. The wire
makes a Z, as in the yanking equation, but now with \(A\) on its middle
stroke. Follow the wire from the input to the output and you pass
through \(A\) \emph{backwards}, from its output side to its input side.
The box has been turned through a \textbf{half-turn}. We take this
picture as the sketchbook's definition of the box labelled \(A^{T}\):

\[\begin{tikzpicture}[sd]
  \draw[wire] (0,0) -- (3.4,0);
  \node[sdbox] at (1.7,0) {$A^{T}$};
\end{tikzpicture}
\;:=\;
\begin{tikzpicture}[sd]
  \draw[wire] (0,3) -- (5,3);
  \sdcap{5}{3}{1.5}
  \draw[wire] (1.5,1.5) -- (5,1.5);
  \node[sdbox] at (3.25,1.5) {$A$};
  \sdcup{1.5}{1.5}{0}
  \draw[wire] (1.5,0) -- (6.5,0);
\end{tikzpicture}\]

\textbf{On the blackboard.} A definition by picture had better agree
with the matrix whose name it borrows.

\begin{quote}
\textbf{Lemma 6.3 (the half-turn is the transpose).} For every
\(2\times2\) complex matrix \(A\),

\[(\text{cap}\otimes I)(I\otimes A\otimes I)(I\otimes\text{cup}) \;=\; A^{T},\]

where \((A^{T})_{ki}=A_{ik}\).
\end{quote}

\emph{Proof.} Write \(A_{ik}=\bra iA\ket k\) for the entry in row \(i\)
and column \(k\). Feed in a basis state \(\ket i\):

\[\ket i \;\longmapsto\; \sum_k \ket{i}\ket{k}\ket{k} \;\longmapsto\; \sum_k \ket{i}\,\bigl(A\ket{k}\bigr)\,\ket{k} \;\longmapsto\; \sum_k \bra{i}A\ket{k}\;\ket k \;=\; \sum_k A_{ik}\ket k .\]

In the last step the cap, acting on the first two wires, keeps exactly
the component of \(A\ket k\) along \(\ket i\). So the entry of our
three-row matrix in row \(k\), column \(i\) is \(A_{ik}\). There are
four entries and we have checked them all: it is \(A^{T}\). No conjugate
appeared at any step, because no bra of \(A\)'s entries was ever taken.
\(\;\blacksquare\)

\textbf{\emph{Why this definition.}} We could have \emph{declared} a new
box called \(A^{T}\) and given it a rule. Defining it as \(A\) with its
wires bent costs nothing new --- only cups and caps --- and it makes the
theorem below almost draw itself. But note which way the debt runs: the
picture is entitled to the name ``transpose'' only because of the index
check just made.

\begin{quote}
\textbf{Coecke's Sketchbook.}

\emph{A half-turn is not a mirror image.} In Chapter 3 we met the dagger
\(A^{\dagger}\), and its picture was the \emph{mirror image} of the box:
the diagram reflected left to right, with every number conjugated. The
transpose is a different motion. Nothing is reflected; the box stays in
the plane of the paper and is \emph{rotated} through a half-turn,
carried round by its own wires. A rotation cannot conjugate anything.

The two are easily confused because for many favourite boxes they agree.
If every entry of \(A\) is real, then \(A^{T}=A^{\dagger}\). If \(A\) is
symmetric, \(A^{T}=A\) and the box does not seem to have turned at all
--- \(H\), \(X\) and \(Z\) are all like this. The Pauli matrix \(Y\) is
a good test: \(Y^{\dagger}=Y\) but \(Y^{T}=-Y\).

This is why our example has an \(i\) in one corner and a \(2\) in the
other. With anything tamer, a wrong rule would have passed the test.
\end{quote}

With the half-turned box defined, the rule suggested by the example can
be drawn.

\begin{quote}
\textbf{Theorem 6.4 (sliding; the transpose trick).} For every
\(2\times2\) complex matrix \(A\),
\(\;(A\otimes I)\,\text{cup}=(I\otimes A^{T})\,\text{cup}\), exactly:
\end{quote}

\[\begin{tikzpicture}[sd]
  \sdcup{0}{1.5}{0}
  \draw[wire] (0,1.5) -- (4,1.5);
  \draw[wire] (0,0) -- (4,0);
  \node[sdbox] at (2,1.5) {$A$};
\end{tikzpicture}
\;=\;
\begin{tikzpicture}[sd]
  \sdcup{0}{1.5}{0}
  \draw[wire] (0,1.5) -- (4,1.5);
  \draw[wire] (0,0) -- (4,0);
  \node[sdbox] at (2,0) {$A^{T}$};
\end{tikzpicture}\]

This is the headline theorem of the chapter, and §6.7 proves it twice.
See \texttt{Code/Ch6/Sliding.py}, which also checks that the equation
\emph{fails} with \(A^{\dagger}\), with \(\bar A\), and with \(A\)
itself on the bottom leg.

Two consequences are worth having now.

\emph{States and effects slide too.} Neither proof needs the box to have
one wire in and one wire out. An effect \(\bra b\) is a \(1\times2\)
matrix, \((\bar b_0\;\;\bar b_1)\); its transpose is the \(2\times1\)
column with the same entries, which is the ket \(\ket{b^*}\). So an
effect on the top leg of a cup slides round to become a state:

\[\begin{tikzpicture}[sd]
  \sdcup{0}{1.5}{0}
  \draw[wire] (0,1.5) -- (2.8,1.5);
  \draw[wire] (0,0) -- (4,0);
  \node[sdeffect] at (2.8,1.5) {$b$};
\end{tikzpicture}
\;=\;
\begin{tikzpicture}[sd]
  \draw[wire] (0,0) -- (2.8,0);
  \node[sdstate] at (0,0) {$b^{*}$};
\end{tikzpicture}\]

Multiply both sides by \(\tfrac1{\sqrt2}\), remembering that
\(\ket{\Phi^+}=\tfrac1{\sqrt2}\,\)cup, and this is exactly the Dirac's
Blackboard of §6.2. The picture that was a bare report there now has a
reason: asking one half of the pair a question \emph{is} handing the
other half the answer, turned half round.

\emph{A rotation of the whole pair.} For a unitary \(U\), slide \(U\)
round to the bottom leg:
\((U\otimes\bar U)\,\text{cup}=(I\otimes\bar U U^{T})\,\text{cup}=\text{cup}\),
because \(\bar UU^{T}\) is the conjugate of \(UU^{\dagger}=I\). Turning
one half by \(U\) and the other by \(\bar U\) leaves the pair as it was.
For real \(U\) this is why the Bell state looked the same in the \(\pm\)
basis as in the \(0,1\) basis.

\begin{center}\rule{0.5\linewidth}{0.5pt}\end{center}

\hypertarget{bending-every-process-is-a-state}{%
\section{§6.5 Bending: Every Process Is a
State}\label{bending-every-process-is-a-state}}

In §6.4 the state \((A\otimes I)\,\text{cup}\) was a stepping stone.
Look at it again for its own sake:

\[(A\otimes I)\,\text{cup} \;=\; \ket{00}+2\ket{01}+i\ket{10}+3\ket{11}
\qquad\text{for}\qquad
A=\begin{pmatrix} 1 & 2\\ i & 3\end{pmatrix}.\]

The four amplitudes are the four entries of \(A\), in reading order.

\textbf{On the blackboard.} To every \(2\times2\) matrix \(A\) assign
the two-wire state

\[\ket{\psi_A} \;:=\; (A\otimes I)\,\text{cup} \;=\; \sum_{k} A\ket k\otimes\ket k \;=\; \sum_{i,k} A_{ik}\,\ket{ik}.\]

The amplitude of \(\ket{ik}\) is the entry \(A_{ik}\): the top wire is
the row index and the bottom wire the column index. Nothing is lost and
nothing is added; \(\ket{\psi_A}\) is \(A\) with its four numbers stood
in a column.

\textbf{In the sketchbook.} The box \(A\) has an input wire on its left.
Take hold of that wire and bend it round with a cup, so that it points
to the right. A process with one wire in and one wire out has become a
\emph{state} on two wires:

\[\begin{tikzpicture}[sd]
  \draw[wire] (0,1.5) -- (2.5,1.5);
  \draw[wire] (0,0) -- (2.5,0);
  \node[sdstate, minimum height=3.4em] at (0,0.75) {$\psi_A$};
\end{tikzpicture}
\;:=\;
\begin{tikzpicture}[sd]
  \sdcup{0}{1.5}{0}
  \draw[wire] (0,1.5) -- (4,1.5);
  \draw[wire] (0,0) -- (4,0);
  \node[sdbox] at (2,1.5) {$A$};
\end{tikzpicture}\]

Can the box be recovered from the state? Bend the wire back: put a cap
on the bottom wire of \(\ket{\psi_A}\) and on a fresh input wire below
it.

\begin{quote}
\textbf{Theorem 6.5 (bending).} The assignment \(A\mapsto\ket{\psi_A}\)
is a one-to-one linear correspondence between \(2\times2\) matrices and
two-wire vectors, and the box is recovered by
\(\;(I\otimes\text{cap})(\ket{\psi_A}\otimes I)=A\):
\end{quote}

\[\begin{tikzpicture}[sd]
  \sdcup{1.5}{3}{1.5}
  \draw[wire] (1.5,3) -- (7.5,3);
  \node[sdbox] at (3.5,3) {$A$};
  \draw[wire] (1.5,1.5) -- (5.5,1.5);
  \sdcap{5.5}{1.5}{0}
  \draw[wire] (0,0) -- (5.5,0);
\end{tikzpicture}
\;=\;
\begin{tikzpicture}[sd]
  \draw[wire] (0,0) -- (3.4,0);
  \node[sdbox] at (1.7,0) {$A$};
\end{tikzpicture}\]

\emph{Proof idea.} On the left the wire comes in at the bottom, goes
round a cap, round a cup, and out through \(A\): an S-bend with a box on
its last stroke. Pull it straight.

\emph{Proof sketch.} In the sketchbook: the box \(A\) sits on the top
row to the right of the cup, so it can be set aside while the S formed
by the cup and the cap is straightened by yanking (Theorem 6.2); what is
left is a straight wire into \(A\). On the blackboard:
\((I\otimes\text{cap})\bigl(\sum_{i,k}A_{ik}\ket{ik}\otimes\ket j\bigr)=\sum_i A_{ij}\ket i=A\ket j\),
since the cap keeps only \(k=j\). That the correspondence is linear and
one-to-one is visible from \(\ket{\psi_A}=\sum A_{ik}\ket{ik}\): four
free entries on one side, four free amplitudes on the other.
\(\;\blacksquare\) See \texttt{Code/Ch6/Bending.py}.

Three entries in this new dictionary are worth reading off.

\begin{itemize}
\tightlist
\item
  The identity matrix corresponds to \(\ket{\psi_I}=\ket{00}+\ket{11}\):
  the cup itself. \textbf{A Bell pair is a plain wire, bent} --- with
  the usual reminder that the normalized Bell state is
  \(\tfrac1{\sqrt2}\) of it.
\item
  The Hadamard matrix corresponds to
  \(\tfrac1{\sqrt2}(\ket{00}+\ket{01}+\ket{10}-\ket{11})\), which is
  \(\sqrt2\) times the puzzling state \(\ket\chi\) of §6.1.
\item
  The square of amplitudes \(C\) in the product test of §6.1 was the
  matrix of the state all along: \(\ket\psi=\ket{\psi_C}\). Theorem 6.1
  now reads: \emph{a state is a product exactly when its matrix is not
  invertible.} A product state \(\ket\alpha\otimes\ket\beta\) is the
  bent form of the matrix \(\ket\alpha\bra{\beta^*}\), which squashes
  everything onto the single direction \(\ket\alpha\). An entangled
  state is the bent form of an invertible box, one that forgets nothing.
  The determinant of §6.1 was the determinant of a process.
\end{itemize}

\textbf{\emph{Why this definition.}} Why the top leg? Because then
amplitudes equal entries, with nothing to remember. Had we put the box
on the bottom leg, sliding tells us what we would get:
\((I\otimes A)\,\text{cup}=(A^{T}\otimes I)\,\text{cup}=\ket{\psi_{A^{T}}}\),
the state of the transposed matrix. Both conventions are in use in the
literature. This book keeps the box on the top leg throughout.

And the honest scope of the two directions. Bending a gate \(U\) into a
state is something one can do in a laboratory: make a Bell pair, apply
\(U\) to the top half, and one holds \(\tfrac1{\sqrt2}\ket{\psi_U}\).
Bending back needs a cap, and a cap is a postselected outcome --- the
answer ``\(\Phi^+\)'' to a question with four possible answers. The
equation of Theorem 6.5 is exact; what it describes happens only when
that answer comes. Chapter 10 will show what to do about the other
three, and the result will be teleportation.

\begin{quote}
\textbf{Dirac's Blackboard.}

\emph{A state with the numbers of a matrix.} The correspondence between
processes and states is named after Andrzej Jamiołkowski (1972) and
Man-Duen Choi (1975). Their theorems concern a grander kind of process
than ours, which we are not equipped to meet before Chapter 9; what we
have here is the seed of it.

Stated exactly, in our conventions: with the box on the \emph{top} leg
of the \emph{unnormalized} cup,

\[\ket{\psi_A}=\begin{pmatrix}A_{00}\\A_{01}\\A_{10}\\A_{11}\end{pmatrix},\qquad A=\begin{pmatrix}A_{00}&A_{01}\\A_{10}&A_{11}\end{pmatrix}.\]

Composition of boxes becomes something visible too. By Theorem 6.4,
\((A\otimes B)\,\text{cup}=(AB^{T}\otimes I)\,\text{cup}=\ket{\psi_{AB^{T}}}\):
a box on each leg of a cup is a single box on one leg. One last
bookkeeping point: the squared length of \(\ket{\psi_A}\) is
\(\sum_{i,k}|A_{ik}|^2\), which for a unitary \(U\) is \(2\), not \(1\).
The missing \(\tfrac1{\sqrt2}\) is the cup's normalization again, and
nothing else.
\end{quote}

\begin{center}\rule{0.5\linewidth}{0.5pt}\end{center}

\hypertarget{the-rosetta-table-5}{%
\section{§6.6 The Rosetta Table}\label{the-rosetta-table-5}}

The rows this chapter adds to the running dictionary, numbered on from
Chapter 5's thirty-eight. Three of them introduce new picture elements
--- the cup, the cap, and the half-turned box; three are rewrite rules;
one records a normalization.

\begingroup\renewcommand{\arraystretch}{1.5}

\begin{longtable}[]{@{}
  >{\raggedright\arraybackslash}p{(\columnwidth - 4\tabcolsep) * \real{0.2727}}
  >{\raggedright\arraybackslash}p{(\columnwidth - 4\tabcolsep) * \real{0.3636}}
  >{\raggedright\arraybackslash}p{(\columnwidth - 4\tabcolsep) * \real{0.3636}}@{}}
\toprule\noalign{}
\begin{minipage}[b]{\linewidth}\raggedright
\#
\end{minipage} & \begin{minipage}[b]{\linewidth}\raggedright
Algebra (Dirac)
\end{minipage} & \begin{minipage}[b]{\linewidth}\raggedright
Picture (Coecke)
\end{minipage} \\
\midrule\noalign{}
\endhead
\bottomrule\noalign{}
\endlastfoot
39 & cup \(=\ket{00}+\ket{11}\) (unnormalized; length \(\sqrt2\)) & a
wire bent so that both ends point right: a state on two wires ---
\emph{new element} \\
40 & cap \(=\bra{00}+\bra{11}\) (unnormalized; as an effect, the
postselected outcome \(\sqrt2\,\bra{\Phi^+}\)) & a wire bent so that
both ends point left: an effect on two wires --- \emph{new element} \\
41 & \(\ket{\Phi^+}=\tfrac{1}{\sqrt2}\,\)cup & the cup with the number
\(\tfrac1{\sqrt2}\) written beside it \\
42 &
\((\text{cap}\otimes I)(I\otimes\text{cup})=I=(I\otimes\text{cap})(\text{cup}\otimes I)\)
& an S-bend or a Z-bend pulled straight --- \textbf{yanking} \\
43 & \(A^{T}\), with entries \((A^{T})_{ki}=A_{ik}\) & the box \(A\)
with both wires bent round: a half-turn (\emph{not} the mirror image,
which is \(A^{\dagger}\)) --- \emph{new element} \\
44 & \((A\otimes I)\,\text{cup}=(I\otimes A^{T})\,\text{cup}\) & a box
slid round the bend arrives turned half round --- \textbf{sliding} \\
45 & \(\ket{\psi_A}=\sum_{i,k}A_{ik}\ket{ik}\) and
\(A=(I\otimes\text{cap})(\ket{\psi_A}\otimes I)\) & bend the input wire
round: a box becomes a state, and back --- \textbf{bending} \\
\end{longtable}

\endgroup

All seven rows are exact equalities of matrices.

\begin{center}\rule{0.5\linewidth}{0.5pt}\end{center}

\hypertarget{the-theorem-twice-5}{%
\section{§6.7 The Theorem, Twice}\label{the-theorem-twice-5}}

\begin{quote}
\textbf{Theorem 6.4 (sliding; the transpose trick).} For every
\(2\times2\) complex matrix \(A\) --- unitary or not ---

\[(A\otimes I)\,\text{cup} \;=\; (I\otimes A^{T})\,\text{cup},\qquad \text{cup}=\ket{00}+\ket{11},\]

exactly, where \(A^{T}\) is the transpose: not the dagger, and not the
conjugate.
\end{quote}

The example of §6.4 made the theorem plausible for one matrix. Here are
two proofs for all of them.

\hypertarget{diracs-blackboard-5}{%
\subsection{Dirac's Blackboard}\label{diracs-blackboard-5}}

Write \(A=\sum_{i,k}A_{ik}\ket i\bra k\), so that
\(A\ket k=\sum_i A_{ik}\ket i\), and recall \((A^{T})_{ki}=A_{ik}\), so
that \(A^{T}\ket i=\sum_k (A^{T})_{ki}\ket k=\sum_k A_{ik}\ket k\).
Remember also that cup \(=\sum_k\ket{kk}\), whatever letter the sum runs
over. Then

\[\begin{aligned}
(A\otimes I)\,\text{cup}
  &= \sum_{k} A\ket k\otimes\ket k
   && \text{($A$ on the top wire)}\\
  &= \sum_{k}\sum_{i} A_{ik}\,\ket i\otimes\ket k
   && \text{(expand $A\ket k$)}\\
  &= \sum_{i}\ket i\otimes\Bigl(\sum_{k} A_{ik}\ket k\Bigr)
   && \text{(collect by $\ket i$)}\\
  &= \sum_{i}\ket i\otimes A^{T}\ket i
   && \text{(recognize $A^{T}\ket i$)}\\
  &= (I\otimes A^{T})\,\text{cup}.
   && \text{($A^{T}$ on the bottom wire)}
\end{aligned}\]

Every step is an equality of four-entry columns; for the matrix of §6.4
each line is the column \((1,\,2,\,i,\,3)\). The entire content is in
the middle line. The double sum \(\sum_{i,k}A_{ik}\ket{ik}\) does not
care whether we group it by \(k\), which reads \(A\) column by column,
or by \(i\), which reads it row by row. Reading a matrix by rows instead
of columns \emph{is} transposing it. No conjugate is ever taken, because
nothing is ever turned from a ket into a bra. \(\;\blacksquare\)

\hypertarget{coeckes-sketchbook}{%
\subsection{Coecke's Sketchbook}\label{coeckes-sketchbook}}

Start from the right-hand side, \((I\otimes A^{T})\,\text{cup}\): a cup
with the box \(A^{T}\) on its bottom leg.

\textbf{Step 1 --- substitution (the definition of the half-turned box,
§6.4).} Replace the box \(A^{T}\) by its picture. We need four rows of
wire, counted from the top of the picture. The outer cup joins rows 1
and 2, and row 1 runs straight out to the right. Row 2 is the input of
the half-turn picture: it runs right to a cap joining rows 2 and 3;
\(A\) sits on row 3; a cup on the left joins rows 3 and 4; row 4 runs
out to the right.

\[\begin{tikzpicture}[sd]
  \sdcup{0}{1.5}{0}
  \draw[wire] (0,1.5) -- (4,1.5);
  \draw[wire] (0,0) -- (4,0);
  \node[sdbox] at (2,0) {$A^{T}$};
\end{tikzpicture}
\;\;\overset{\text{def.}}{=}\;\;
\begin{tikzpicture}[sd]
  \sdcup{0}{4.5}{3}
  \draw[wire] (0,4.5) -- (7,4.5);
  \draw[wire] (0,3) -- (5,3);
  \sdcap{5}{3}{1.5}
  \draw[wire] (1.5,1.5) -- (5,1.5);
  \node[sdbox] at (3.25,1.5) {$A$};
  \sdcup{1.5}{1.5}{0}
  \draw[wire] (1.5,0) -- (7,0);
\end{tikzpicture}\]

\textbf{Step 2 --- yanking (Theorem 6.2, the S-bend).} Follow the wire
from the output side of \(A\). It leaves \(A\) on row 3, goes round the
cap to row 2, runs left, goes round the outer cup to row 1, and runs
out. That is an S-bend with nothing on it,
\((I\otimes\text{cap})(\text{cup}\otimes I)\), and it equals a straight
wire. Pull it straight: the output of \(A\) now runs directly out on the
top.

\[\begin{tikzpicture}[sd]
  \sdcup{0}{4.5}{3}
  \draw[wire] (0,4.5) -- (7,4.5);
  \draw[wire] (0,3) -- (5,3);
  \sdcap{5}{3}{1.5}
  \draw[wire] (1.5,1.5) -- (5,1.5);
  \node[sdbox] at (3.25,1.5) {$A$};
  \sdcup{1.5}{1.5}{0}
  \draw[wire] (1.5,0) -- (7,0);
\end{tikzpicture}
\;\;\overset{\text{yank}}{=}\;\;
\begin{tikzpicture}[sd]
  \sdcup{0}{1.5}{0}
  \draw[wire] (0,1.5) -- (4,1.5);
  \draw[wire] (0,0) -- (4,0);
  \node[sdbox] at (2,1.5) {$A$};
\end{tikzpicture}\]

What remains is the inner cup with \(A\) on its top leg and a bare
bottom leg: \((A\otimes I)\,\text{cup}\). The two outputs are in the
right order --- what was row 1 is on top, what was row 4 is below ---
and no wire has crossed another. One substitution, one yank.
\(\;\blacksquare\)

\hypertarget{the-verdict-5}{%
\subsection{The verdict}\label{the-verdict-5}}

The sketchbook proof is shorter, and it shows something the blackboard
hides: \emph{why} the theorem is true. The box never changes; it is
carried round a corner by its own wire, and a thing carried round a
corner arrives facing the other way. One sees at a glance that the same
must hold for states, for effects, and for boxes of any size.

But the picture proof has a debt, and it should be named. It proves that
\(A\) on the top leg equals \emph{the half-turned box} on the bottom leg
--- and what it calls \(A^{T}\) is a drawing. That this drawing is the
matrix transpose, and not the dagger or the conjugate, is Lemma 6.3: a
four-entry index check, which is algebra. The blackboard proof is five
lines, asks nothing of anyone, and tells us exactly which numbers go
where. So: the picture wins on length and on insight, and it wins
honestly only because the blackboard has signed for the name. In this
chapter neither strand could have done without the other --- which,
considering the subject, seems fair.

\begin{center}\rule{0.5\linewidth}{0.5pt}\end{center}

\hypertarget{the-chain}{%
\section{§6.8 The Chain}\label{the-chain}}

We have been bending one wire. Let us bend it more than once.

\hypertarget{one-junction}{%
\subsection{One junction}\label{one-junction}}

Take two cups, on wires 1--2 and 3--4: two separate Bell pairs (each
\(\sqrt2\) times a normalized one; cups are unnormalized, as ever).
Wires 1 and 4 have nothing to do with each other: the state is cup
\(\otimes\) cup, a product across the middle of the page. Now put a cap
on the two \emph{inner} wires, 2 and 3. On the blackboard, writing the
two cups as \(\sum_{i,j}\ket{i\,i\,j\,j}\),

\[(I\otimes\text{cap}\otimes I)\sum_{i,j}\ket{i\,i\,j\,j} \;=\; \sum_{i,j}\bigl(\braket{00|ij}+\braket{11|ij}\bigr)\,\ket{i\,j} \;=\; \sum_{i}\ket{i\,i} \;=\; \text{cup},\]

because the cap passes only \(i=j\). The outer wires, 1 and 4, which
began in separate pairs, are now a pair. In the sketchbook there is
nothing to calculate: cup, cap, cup in a column is an S-bend sitting on
the top leg of the lower cup, and by Theorem 6.2 it pulls straight:

\[\begin{tikzpicture}[sd]
  \sdcup{0}{4.5}{3}
  \sdcup{0}{1.5}{0}
  \draw[wire] (0,4.5) -- (5,4.5);
  \draw[wire] (0,3) -- (3,3);
  \draw[wire] (0,1.5) -- (3,1.5);
  \draw[wire] (0,0) -- (5,0);
  \sdcap{3}{3}{1.5}
  \node[sdlabel] at (1.5,4.95) {1};
  \node[sdlabel] at (1.5,3.45) {2};
  \node[sdlabel] at (1.5,1.95) {3};
  \node[sdlabel] at (1.5,0.45) {4};
\end{tikzpicture}
\;\overset{\text{yank}}{=}\;
\begin{tikzpicture}[sd]
  \sdcup{0}{4.5}{0}
  \draw[wire] (0,4.5) -- (3,4.5);
  \draw[wire] (0,0) -- (3,0);
  \node[sdlabel] at (1.5,4.95) {1};
  \node[sdlabel] at (1.5,0.45) {4};
\end{tikzpicture}\]

As always the cap is a postselected outcome; we count its cost below.

This is called \textbf{entanglement swapping}. It deserves to be strung
out further.

\hypertarget{the-master-diagram}{%
\subsection{The master diagram}\label{the-master-diagram}}

Here is the Anteater's story from the dialogue, as a diagram. Alice
holds wire 1. A first relay station holds wires 2 and 3. A second relay
station holds wires 4 and 5. Bob holds wire 6. Three sources have each
made a pair and sent the halves to neighbours: a cup on wires 1--2, a
cup on wires 3--4, a cup on wires 5--6. Each station asks its two wires
the Bell question and --- in the case we draw --- hears ``\(\Phi^+\)'':
a cap on wires 2--3 and a cap on wires 4--5.

\begin{center}
\begin{tikzpicture}[sd]
  \sdcup{0}{7.5}{6}
  \sdcup{0}{4.5}{3}
  \sdcup{0}{1.5}{0}
  \draw[wire] (0,7.5) -- (6,7.5);
  \draw[wire] (0,6) -- (3.5,6);
  \draw[wire] (0,4.5) -- (3.5,4.5);
  \draw[wire] (0,3) -- (3.5,3);
  \draw[wire] (0,1.5) -- (3.5,1.5);
  \draw[wire] (0,0) -- (6,0);
  \sdcap{3.5}{6}{4.5}
  \sdcap{3.5}{3}{1.5}
  \node[sdlabel] at (1.75,7.95) {1};
  \node[sdlabel] at (1.75,6.45) {2};
  \node[sdlabel] at (1.75,4.95) {3};
  \node[sdlabel] at (1.75,3.45) {4};
  \node[sdlabel] at (1.75,1.95) {5};
  \node[sdlabel] at (1.75,0.45) {6};
  \node[sdlabel, anchor=west] at (6.4,7.5) {Alice};
  \node[sdlabel, anchor=west] at (6.4,0) {Bob};
\end{tikzpicture}
\end{center}

This is the master diagram of the chapter, and we call it \textbf{The
Chain}.\footnote{The reader may wish to count the wires in this figure.
  We offer no explanation at this time.} It is one wire, threaded out
and back: out round a cup, back round a cap, out, back, out.

\begin{quote}
\textbf{Theorem 6.6 (the chain).} Three cups and two caps, threaded as
in The Chain, are one cup on wires 1 and 6:

\[(I\otimes\text{cap}\otimes\text{cap}\otimes I)\,(\text{cup}\otimes\text{cup}\otimes\text{cup}) \;=\; \text{cup},\]

exactly.
\end{quote}

\emph{Proof idea.} One junction turns two pairs into one. Two junctions
turn three pairs into one. Each junction is a single yank.

\emph{Proof, in the sketchbook.} \textbf{First yank (Theorem 6.2,
S-bend):} the cup on 1--2, the cap on 2--3 and the top leg of the cup on
3--4 form an S; pull it straight, and a single cup joins wires 1 and 4.
\textbf{Second yank (the same rule):} that cup, the cap on 4--5 and the
top leg of the cup on 5--6 form another S; pull.

\[\begin{tikzpicture}[sd]
  \sdcup{0}{7.5}{6}
  \sdcup{0}{4.5}{3}
  \sdcup{0}{1.5}{0}
  \draw[wire] (0,7.5) -- (5,7.5);
  \draw[wire] (0,6) -- (3,6);
  \draw[wire] (0,4.5) -- (3,4.5);
  \draw[wire] (0,3) -- (3,3);
  \draw[wire] (0,1.5) -- (3,1.5);
  \draw[wire] (0,0) -- (5,0);
  \sdcap{3}{6}{4.5}
  \sdcap{3}{3}{1.5}
  \node[sdlabel] at (1.5,7.95) {1};
  \node[sdlabel] at (1.5,6.45) {2};
  \node[sdlabel] at (1.5,4.95) {3};
  \node[sdlabel] at (1.5,3.45) {4};
  \node[sdlabel] at (1.5,1.95) {5};
  \node[sdlabel] at (1.5,0.45) {6};
\end{tikzpicture}
\;\overset{\text{yank}}{=}\;
\begin{tikzpicture}[sd]
  \sdcup{0}{7.5}{3}
  \sdcup{0}{1.5}{0}
  \draw[wire] (0,7.5) -- (5,7.5);
  \draw[wire] (0,3) -- (3,3);
  \draw[wire] (0,1.5) -- (3,1.5);
  \draw[wire] (0,0) -- (5,0);
  \sdcap{3}{3}{1.5}
  \node[sdlabel] at (1.5,7.95) {1};
  \node[sdlabel] at (1.5,3.45) {4};
  \node[sdlabel] at (1.5,1.95) {5};
  \node[sdlabel] at (1.5,0.45) {6};
\end{tikzpicture}
\;\overset{\text{yank}}{=}\;
\begin{tikzpicture}[sd]
  \sdcup{0}{7.5}{0}
  \draw[wire] (0,7.5) -- (3,7.5);
  \draw[wire] (0,0) -- (3,0);
  \node[sdlabel] at (1.5,7.95) {1};
  \node[sdlabel] at (1.5,0.45) {6};
\end{tikzpicture}\]

\emph{Proof, on the blackboard.} Three cups are
\(\sum_{i,j,k}\ket{i\,i\,j\,j\,k\,k}\). The cap on wires 2--3 passes
only \(i=j\); the cap on wires 4--5 passes only \(j=k\). What is left is
\(\sum_i\ket{i\,i}\) on wires 1 and 6. \(\;\blacksquare\) See
\texttt{Code/Ch6/Chain.py}.

Here the two strands finish level --- one line each --- but the picture
has the better of it in one respect: it makes plain that the length of
the chain is of no consequence. Ten stations or a thousand, the wire
pulls straight --- as an equality of processes; what each pull costs
comes next.

\hypertarget{what-the-pull-costs}{%
\subsection{What the pull costs}\label{what-the-pull-costs}}

Now the \(\tfrac1{\sqrt2}\) that we have been carrying must be paid in.
The sources make normalized Bell states, and a station's answer
``\(\Phi^+\)'' is the normalized effect \(\bra{\Phi^+}\). Three states
and two effects contribute \(\bigl(\tfrac1{\sqrt2}\bigr)^{5}\), and the
cup that comes out is \(\sqrt2\,\ket{\Phi^+}\):

\[(I\otimes\bra{\Phi^+}\otimes\bra{\Phi^+}\otimes I)\;\ket{\Phi^+}\otimes\ket{\Phi^+}\otimes\ket{\Phi^+} \;=\; \frac14\,\ket{\Phi^+}.\]

The amplitude \(\tfrac14\) is \(\tfrac12\cdot\tfrac12\), one factor for
each junction --- the same \(\tfrac12\) that a normalized kink cost in
§6.3. Its square, \(\tfrac1{16}\), is the probability that \emph{both}
stations hear ``\(\Phi^+\)'': each does so with probability
\(\tfrac14\), fairly enough for one answer out of four. When they do,
Alice and Bob hold a perfect, normalized Bell pair.

So the diagram is an exact equation about one case in sixteen. That is
what a cap means, and we do not wish to sell it as more. It is an
equality of processes; it is not a claim about what a network of such
stations achieves, how fast, or how reliably. What happens in the other
fifteen cases is not failure --- each of the other answers leaves Alice
and Bob with a pair that is a known one-box correction away from
\(\ket{\Phi^+}\) --- but seeing that requires classical messages and a
little more notation. That is Chapter 10.

\begin{quote}
\textbf{Feynman's Notebook.}

\emph{Particles that never met.} Entanglement swapping was proposed in
1993 by Marek Żukowski, Anton Zeilinger, Michael Horne and Artur Ekert,
in a paper whose title spoke of ``event-ready detectors'': the click of
the middle station would announce that two far-apart particles, from
independent sources, were now entangled and ready to be tested. It was
first carried out in 1998 by Jian-Wei Pan, Dik Bouwmeester, Harald
Weinfurter and Zeilinger, with photons; their title says it plainly ---
``entangling photons that never interacted''.

In the laboratory the middle station is a half-silvered mirror where two
photons, one from each pair, arrive at the same moment, followed by
detectors. Certain patterns of clicks identify a particular Bell
outcome; the experimenters keep those runs and set the others aside.
That is postselection with the lid off, and it is the kind of thing our
cap draws. With plain optics at most two of the four outcomes can be
told apart, and \(\Phi^+\) is not among them; the 1998 experiment caught
another, which serves as well.

Alice's and Bob's particles share no past. What they share is a wire
that was threaded through other people's hands.
\end{quote}

\begin{center}\rule{0.5\linewidth}{0.5pt}\end{center}

\hypertarget{was-it-in-there-all-along}{%
\section{§6.9 Was It in There All
Along?}\label{was-it-in-there-all-along}}

Let us gather what we have. A Bell pair has a description although its
halves have none, and a determinant says so (§6.1). Each half, asked
anything, answers like a fair coin, and the two answers are locked
together (§6.2). Drawn as it deserves, the pair is a bent wire, and a
bent wire can be pulled straight (§6.3). A box pushed round the bend
arrives transposed --- turned half round, never mirrored --- which we
proved twice (§6.4, §6.7). Bending turns every box into a state and
back, so that the pair is nothing but the identity, bent (§6.5). And a
wire bent several times still pulls straight, which joins parties who
never met (§6.8). In every picture the agreement between the two ends
was kept in the same place: not at either end, but in the fact that they
are ends of one wire.

That is how the theory \emph{describes} the pair. It does not yet answer
the question a sensible person asks first, the one we set aside at the
very beginning. Two halves that always agree: were the answers simply
written inside them at the source, one sealed note in each, identical?
Nothing in this chapter says no. The sealed notes reproduce the fair
coins; they reproduce the agreement; with a little care they reproduce
every number we have computed. If there is a difference between a bent
wire and a pair of sealed notes, it must show itself somewhere we have
not yet looked --- in questions put to the two halves in
\emph{different} directions, by people too far apart to confer.

In 1964 John Bell found where to look. He wrote down a quantity that any
pair of sealed notes must keep below a certain bound, and asked whether
a bent wire must keep below it too. The clearest way to meet his
discovery is as a game, played by partners who may agree on anything
beforehand and may not speak once it starts. Bring your best strategy.
Bring the sealed notes. The next chapter deals the cards.

\hypertarget{partita-for-three-players-far-apart}{%
\chapter{Partita for Three Players, Far
Apart}\label{partita-for-three-players-far-apart}}

\begin{center}\rule{0.5\linewidth}{0.5pt}\end{center}

\emph{初霜や}\\
\emph{三つの窓に}\\
\emph{紋ひとつ}

\emph{hatsushimo ya}\\
\emph{mittsu no mado ni}\\
\emph{mon hitotsu}

\emph{The first frost ---}\\
\emph{on three separate windows,}\\
\emph{one pattern.}

\begin{center}\rule{0.5\linewidth}{0.5pt}\end{center}

\emph{{[}The stair hall of the Crab's house, on the first hard morning
of the winter; the low sun comes level through the tall windows and lies
white on the floor. The Crab is away at his cousin's in Ostend. On the
hall table stand a travelling-case with its lid still down, three sealed
envelopes laid in a row, and a folded sheet of paper held flat under a
glass paperweight; a longer letter, already unfolded, lies beside them.
Achilles and the Tortoise come in from the lane, stamping the cold off
their feet; the Tortoise has the house-key in her glove.{]}}

\emph{{[}The hall is built around its stair. One broad flight climbs to
the first landing, and there the ways divide --- on up under the roof,
down to the cellar, off by a passage to the garden room --- so that
three rooms hang from the one hall and no two of them share a floor; the
hall table stands close under the foot of it. The Crab has left a fire
banked in the grate, enough to take the worst edge off; the cold still
comes up out of the flagstones as a surprise. The three tall windows are
blind with the night's frost, bloomed white from corner to corner, and
the frozen town beyond has gone to the same milk --- the low sun stands
behind the panes and does not yet burn through. On the wall at the first
landing, where the stair turns, hangs Escher's \textup{Relativity} in a
plain frame: stairways that cross it at every angle, a walled garden
showing through one archway, small figures forever climbing and
descending who never seem to arrive; the frost-light has not reached it.
The three envelopes lie in their row with a name to each; the folded
slip keeps still under its glass; the case waits with its lid down.{]}}

\textbf{Achilles:} I have never seen the canal stand still before. There
were boys walking on it as if it owed them money.

\textbf{Tortoise:} It will bear a boy this morning. It would not have
borne one yesterday.

\textbf{Achilles:} And our host bears none of it --- off to the seaside
the one week the whole town turns to glass. \emph{{[}He warms his
hands.{]}} You have his key, his hall and his breakfast fire, and the
man himself is watching the grey waves at Ostend.

\emph{{[}He stands close to the grate a moment; even indoors, this near
the frozen glass, his breath still shows. He rubs a thumb-print clear in
the frost and looks out: the canal lies flat and white under the boys,
and the low sun strikes no glitter off it at all.{]}}

\textbf{Tortoise:} He left the key with me and a letter with the key. He
is sorry to miss us; he says so at length. \emph{{[}She sets down her
glove and takes up the unfolded letter.{]}} And he begs --- twice,
underlined --- that we do not wait for the evening.

\textbf{Achilles:} Why ever not? Last autumn it was all ``another
evening, bring it another evening.''

\textbf{Tortoise:} \emph{{[}reading{]}} ``A trio does not keep as a pair
keeps. Two boxes will sit on a shelf and wait for you; three made
together grow restless apart, and I will not answer for them past a day
or two. Open them by daylight, and soon, or the whole afternoon is
wasted.'' \emph{{[}She lays the letter down.{]}} You were never going to
wait for the dark in any case. You have been carrying that card since
October.

\textbf{Achilles:} I have. \emph{{[}He produces a stiff ruled card from
his coat.{]}} You told me, when we parted, to bring my theory to the
game. Here it is, written out. Notes sealed inside, composed any way I
please --- your own words --- and I mean to play them against whatever
is in that case and win.

\emph{{[}The Anteater is shown in behind them, red winter earth still on
one cuff.{]}}

\textbf{Anteater:} Good morning, good morning --- don't let the cold in
on my account, I bring enough of it. Aunt Hillary sends her regards and
a loan. \emph{{[}He sets a worn ruled book on the table.{]}} Her friends
over the hill kept it; I thought it might earn its keep today.

\textbf{Tortoise:} It will. Sit where you like, Doctor; the Crab has
left us the whole hall and strict instructions to spoil nothing.

\emph{{[}With the Doctor in and the door shut on the cold, the hall
settles. The Tortoise fills three cups from the pot the Crab has left
warming on the hob by the grate and hands them round. The sun has
climbed enough to lay one shaft across the table through a thinning
place in the frost, and the dust the three of them have stirred up hangs
turning in it --- each mote going its own way, the shaft keeping its
shape whole; the Anteater watches it a moment, the way a man watches a
column of gnats, before he sits.{]}}

\begin{center}\rule{0.5\linewidth}{0.5pt}\end{center}

\hypertarget{first-challenge-the-sealed-notes-praeludium}{%
\section{\texorpdfstring{First Challenge --- \emph{The Sealed Notes}
(Praeludium)}{First Challenge --- The Sealed Notes (Praeludium)}}\label{first-challenge-the-sealed-notes-praeludium}}

\textbf{Tortoise:} Before the case is opened at all, Achilles, make your
theory as strong as you can on paper. There is a game the Crab
describes; play your card against it dry, and see how it does.
\emph{{[}She turns to the letter and finds the place.{]}} The rules
first, in his own hand.

\textbf{Tortoise:} \emph{{[}reading{]}} ``The pair game. Set the two
boxes of one pair in two rooms, far apart, and let a referee ask each
one a single question --- the lid, or the door --- chosen only after the
players are parted, so that neither can know what the other was asked.
Each box answers with its one note. They win the round if the two notes
come out alike --- save in the one case where the left box is asked the
door and the right box the lid, when, for reasons the daughter would not
tell me, they win only if the two are unlike. Three of the four askings
want alike; the odd one wants unlike.''

\textbf{Achilles:} Easy. If the answer is written in the box before it
leaves the bench, I only have to write it well. \emph{{[}He fetches a
handful of reversi counters from the Crab's games table.{]}} These will
do for the notes --- white face up for a bright note, black for a dark
one; a counter in a square is what the box has decided to say.
\emph{{[}He rules his card into four squares and names them under his
breath.{]}} Left box, lid. Left box, door. Right box, lid. Right box,
door. Four squares, four sealed notes, the whole of what a pair can be
told to do.

\textbf{Achilles:} \emph{{[}laying a white counter in each{]}} And I
say: let both boxes always answer bright. All white. Then whenever they
are asked to be alike, they are --- three askings in four, won on the
spot.

\textbf{Tortoise:} The Crab keeps a packet of masks with the boxes ---
strips of card with windows cut in them; laid over your card, a mask
shows only the two squares one asking cares about, and a word to say
what that asking wants. \emph{{[}She takes them from the packet.{]}} One
asking to a mask. Let me lay them over what you have written, one at a
time.

\emph{{[}The card as it stands: four white counters. She lays the first
mask; its two windows sit over the left-lid and the right-door squares,
both white, and the word on it is ALIKE.{]}}

\textbf{Tortoise:} Left lid, right door: alike, and alike it is.
\emph{{[}She lays the second, windows over the two lids, word ALIKE ---
white and white.{]}} The two lids: alike. \emph{{[}The third, over the
two doors, ALIKE --- white and white.{]}} The two doors: alike.
\emph{{[}She lays the last, its windows over the left-door and right-lid
squares, and its word is UNLIKE; both showing are white.{]}} And here
--- left door, right lid --- this one wants unlike, and you have given
it alike.

\emph{{[}From her pocket she takes three coat buttons and stands one on
the fourth mask.{]}}

\textbf{Tortoise:} A button, where a rule is broken. One broken, three
kept.

\textbf{Achilles:} A slip, only. \emph{{[}He turns the right-lid counter
over.{]}} Make the right box answer dark when it is asked the lid. Now
the unlike mask has a white and a black --- unlike, mended.

\emph{{[}The card as it was: four white counters, one button on the
unlike mask. He turns one counter, the right-lid square, white to black;
the other three do not move. The unlike mask now reads white-and-black
-- unlike, kept -- and its button has nothing to sit on; but the mask
over the two lids, which read white-and-white and was kept, now reads
white-and-black, broken. The button crosses from the one mask to the
other. One counter turned, two masks changed their verdict, and there is
still exactly one button on the card.{]}}

\textbf{Achilles:} \ldots The button has only moved. It has gone to the
two-lids mask.

\textbf{Tortoise:} Mend that one.

\textbf{Achilles:} \emph{{[}turning the left-lid counter to black{]}}
Both lids dark, then --- alike again. \emph{{[}He watches.{]}} And now
the left-lid-right-door mask has gone wrong instead, and the button is
sitting on \emph{that}. It walks. Whatever I flip, it walks to the next
square along and sits down again.

\emph{{[}Achilles has been fidgeting over the card, one heel up on the
bottom step of the stair behind him, and the treads answer him --- one
gives with a small complaint as he leans his weight on it, and when he
shifts off it that one hushes and the tread above takes up the creak.
There is always one of them talking, and it moves when he moves.{]}}

\textbf{Achilles:} Two cards, then --- one for a lid morning and one for
a door morning --- and I toss a coin to choose which card the pair
carries out with them.

\textbf{Tortoise:} You may. But the boxes part before the questions are
asked; the coin is spent before anyone knows what will be wanted. All
the coin does is choose which of your two cards is on the table when the
mask comes down --- and each card, alone, still has its one broken mask.
The coin only picks which asking loses.

\textbf{Achilles:} Then where is the fault? Not in the left box --- its
two notes are as good as I please. Not in the right. Not in any one
square on this card.

\textbf{Anteater:} At which level are you asking? The fault is real, but
you are right that it is in no single square. Look where each square
lies. \emph{{[}He sets a fingertip on the left-door counter.{]}} This
one shows through two masks at once --- the two-doors mask and the
unlike mask.

\textbf{Achilles:} Every square, the same way?

\textbf{Anteater:} Every square. Each note you write is read by two of
the four askings, so turn one counter and you never change one verdict
--- you change two. An even change, to a thing that must come out odd.

\textbf{Achilles:} Odd?

\textbf{Anteater:} Count the broken masks, ever. One, or three. Never
none, never two. To win all four you would have to break none, and none
is even, and you can only ever move by twos across an odd total. The
button has nowhere to vanish to.

\textbf{Achilles:} Then three in four is simply the most anyone can do.

\textbf{Anteater:} I know two people who do better.

\emph{{[}The sun has come far enough round now to rake the lowest window
from the side, low and glancing, and where its light runs along the
frost the flat white breaks open into fern and feather and whole
branching forests. Achilles, waiting on the Doctor, follows one frond of
it with his eye until it loses itself in the next.{]}}

\begin{center}\rule{0.5\linewidth}{0.5pt}\end{center}

\hypertarget{second-challenge-the-wedge-allemande}{%
\section{\texorpdfstring{Second Challenge --- \emph{The Wedge}
(Allemande)}{Second Challenge --- The Wedge (Allemande)}}\label{second-challenge-the-wedge-allemande}}

\textbf{Achilles:} Better? Nobody does better than the best. Three in
four is the ceiling; I have just watched it hold.

\textbf{Anteater:} It is the ceiling for written notes. \emph{{[}He
draws Aunt Hillary's book to himself.{]}} This is a score-book --- a
hundred ruled rows, a hundred evenings, kept by two of Aunt Hillary's
friends who have never once met. Alice lives east of the hill; Bob lives
west; the workshop sent them a hundred pairs between them, one box of
each pair to each, and every evening for a hundred evenings they played
the Crab's pair game across the whole width of the hill and marked down
whether they won. Let me tell you how Bob's end of it goes, because Bob
does a small thing that changes everything.

\textbf{Anteater:} Consider Bob, west of the hill, and the little brass
wedge that came in his parcel along with the boxes.

\textbf{Achilles:} Only a wedge? A doorstop with ideas?

\textbf{Anteater:} Each evening, before he opens his box, he stands it
up on that wedge, so that the box is put its question not square-on but
at a slant.

\textbf{Achilles:} Cutting the difference between the two ways of asking
--- I begin to see it. But between which two?

\textbf{Anteater:} Keep the picture and hold your question one moment;
the wedge is the whole of his trick, and I want the Tortoise to name it.

\textbf{Achilles:} Even a slant is a definite thing, though. He still
hears one note and writes one mark.

\textbf{Tortoise:} The wedge asks the box a question tilted halfway
between lid and door --- square between the two, halfway from the one to
the other, and that is the whole of it.

\emph{{[}Held at that slant, the glancing light picks every fern out of
the frost bright, edge and vein and all; the same pane, taken square-on
from where the Tortoise sits, shows nothing but a blank white. She tilts
her head to catch it the sidelong way, and the whole forest comes up out
of the glass at once.{]}}

\textbf{Tortoise:} Let Bob draw the two columns, then, east and west,
with a space for the mark between. \emph{{[}In the tale, Bob rules the
page; on the table the Tortoise sets a finger to the head of each
column.{]}} What Alice was asked, what Bob was asked.

\textbf{Tortoise:} Let Bob draw the hundred rows down the page, one to
an evening. \emph{{[}Bob, in the tale, rules the lines across; the
Tortoise runs her finger down the ready-ruled page.{]}} East asked, west
asked, and a stroke where they won.

\textbf{Anteater:} And here is the book he and Alice filled. Count it
yourself; I would rather you did.

\textbf{Achilles:} \emph{{[}running down the strokes column{]}}
Strokes\ldots{} a great many strokes. \emph{{[}He tallies on the back of
a torn card.{]}} Eighty-five. Eighty-five wins in the hundred.
\emph{{[}He stops.{]}} That is not possible. My card could not pass
seventy-five however I wrote it --- you all watched the button refuse to
leave.

\textbf{Anteater:} Let the Tortoise show you the drawing before you
decide it is cheating.

\textbf{Tortoise:} Let Bob draw a stroke for every evening they won,
sorted by which pair of questions was asked. \emph{{[}Bob, in the tale,
gathers the strokes into four heaps; the Tortoise blocks the page into
four with her hand.{]}} Lid-and-door here, lid-and-lid here,
door-and-door, door-and-lid. Count each heap.

\textbf{Achilles:} \emph{{[}counting{]}} About the same in each. A
little over four in five, every heap. \emph{{[}Slowly.{]}} That is the
thing my card could never do. My card always had one asking it lost
\emph{every single time} --- the button's square. Theirs loses no asking
always. There is no heap they gave up.

\textbf{Achilles:} But eighty-five in a hundred --- they must be
signalling. A lamp over the hill, a flag; Bob learns what Alice was
asked and leans his wedge to suit.

\textbf{Anteater:} Ask the book. Take Bob's column alone, the way I once
laid a blotter over half the Crab's ledger. \emph{{[}He covers Alice's
column with his hand.{]}} Bob's own notes, whatever Alice was asked that
evening: bright and dark, near enough half and half, all the way down
--- and Alice's alone, just the same. Neither column, read by itself, so
much as twitches when the other is asked a different question. There is
nothing in Bob's own results that could tell him what Alice heard;
nothing crosses the hill.

\emph{{[}The shaft of sun still stands over the table, its dust turning.
Follow any single mote and it wanders --- up, aside, back, no telling
where it means to go; take the whole shaft at once and it holds its slow
steady shape against the dark, and not a grain of it carries a word to
the grain beside it.{]}}

\textbf{Achilles:} \emph{{[}to the Tortoise{]}} But this breaks the
Crab's own rule, and you were the one who told it to him. He was made to
promise, last autumn, that a pair is opened by \emph{the same way in} on
both sides. Bob is standing his box on a wedge and Alice is not. They
are opening differently.

\textbf{Tortoise:} They are --- but that rule was for a different
purpose. ``Same way in'' was the recipe for making a pair \emph{agree}:
open both by the lid and they match, both by the door and they match.
This game does not ask them to agree; it asks them to be alike sometimes
and unlike once, on questions chosen after they part. A pair that only
ever agreed would lose the unlike asking every time --- it would be your
all-white card. The wedge is for winning a harder game, not for
agreeing.

\textbf{Achilles:} Eighty-five is still not a hundred. They lose fifteen
evenings in the hundred; they \emph{do} lose. Perhaps the notes are
written after all, and only smudged --- good enough for eighty-five,
spoilt on the rest.

\emph{{[}The Tortoise lifts the lid of the Crab's travelling-case.
Inside, in three shaped nests, lie three small brass boxes.{]}}

\textbf{Tortoise:} These never lose.

\emph{{[}The three little boxes lie in their three nests, shut, giving
off nothing --- no tick, no shine, only cold brass and the low sun
caught in it. Nobody reaches to wake them. Achilles draws his chair
nearer the grate and warms his hands round the cup, watching the open
case the way one watches a thing one has been told not to touch.{]}}

\begin{center}\rule{0.5\linewidth}{0.5pt}\end{center}

\hypertarget{third-challenge-three-from-one-bench-courante}{%
\section{\texorpdfstring{Third Challenge --- \emph{Three from One Bench}
(Courante)}{Third Challenge --- Three from One Bench (Courante)}}\label{third-challenge-three-from-one-bench-courante}}

\textbf{Achilles:} Never? Everything loses sometimes. The wedge stopped
at eighty-five.

\textbf{Tortoise:} The wedge is a pair. This is a trio. \emph{{[}She
lifts the three boxes from their nests, one, and one, and one, and
stands them on the cloth, shut.{]}} Made together at one bench on one
afternoon, sent in one case. No lids on these --- two ways in only: a
little door in the side, and a drawer underneath. I shall not open one;
I only want them standing where you can see them.

\textbf{Achilles:} And the guarantee?

\emph{{[}The three stand out on the bare cloth now, out of their case
for the first time since Delft, each in its own small pool of shadow
thrown by the low sun. Three of one make, and nothing to choose between
them --- alike as three notes of one chord not yet sounded.{]}}

\textbf{Tortoise:} \emph{{[}from the letter{]}} ``The trio, and my
workshop's word upon it. Open all three by the door, and the number of
dark notes will be even --- none dark, or two. Open one by the door and
the other two by the drawer, and the dark notes will be odd --- one, or
three. My workshop's ledger shows no trio that ever broke it, not
once.'' \emph{{[}She sets the letter down.{]}} His word, and the
bench's, that it holds every time.

\textbf{Achilles:} Then this is my case made for me, and made better
than I could make it. \emph{{[}He sweeps the four counters off the front
of his card.{]}} Watch. If I open two of them by the door and hear their
notes, I can tell you the third before it is touched. Two brights so
far, darks even --- the third must be bright, or the count goes odd and
the guarantee is broken. I would know it without laying a finger on the
box. And what can be known for certain, without touching the thing ---
Doctor, you taught me the words yourself last autumn --- was a real
property of the thing all along. It was written inside.

\textbf{Tortoise:} Say it on the card, then, where we can check it.
\emph{{[}She turns his card over and rules the back into six squares ---
three columns, one per box, and two rows, the door and the drawer.{]}}
Six squares: what each of the three would say, asked each of the two
ways. And these same masks read a trio, not a pair --- three windows
now, one to a box, and a word for what the three together must come to.
Now --- turn your back a moment.

\textbf{Anteater:} Lend me the Crab's egg-cup from the dresser, first
--- this one. Upturned, it covers a single counter whole; whatever I
stand it on is fixed but hidden, which is exactly what I shall want.

\emph{{[}Achilles turns to the window. The Tortoise lays a single
trio-mask over the empty card: three windows, one to a column, and the
word EVEN across it --- the all-doors asking, whose darks must be even.
The Anteater sets a counter in each window, obeying the word: two
bright, and a third he keeps to himself; then he stands the upturned
egg-cup on the flat of the mask, over the third window. What can be seen
now: two white faces and a small white dome; what is fixed and hidden:
the one note under the cup.{]}}

\textbf{Anteater:} Turn round. The mask says EVEN --- all three by the
door, darks even. Two of the three are shown you. The third is under the
cup. Name it.

\textbf{Achilles:} Two brights showing, and darks must be even --- none
or two. Two brights is no darks yet; a dark under the cup would make it
one, which is odd and forbidden. So: bright.

\emph{{[}The Anteater lifts the egg-cup. The counter beneath shows
white.{]}}

\textbf{Anteater:} Bright. Again --- turn round.

\emph{{[}Achilles turns away. The same EVEN mask stays where it is; the
Anteater lifts two of the counters, sets one white and one black in
their windows, and covers a different window with the egg-cup. Two faces
show -- one white, one black -- and the dome stands over the third.{]}}

\textbf{Anteater:} Same mask, same word. One bright, one dark shown.
Name the hidden one.

\textbf{Achilles:} One dark already, and even is wanted; a bright under
the cup leaves it at one, odd, forbidden --- so it must be dark, to make
two. Dark.

\emph{{[}The egg-cup comes up; the counter is black.{]}}

\textbf{Achilles:} Foretold without touching it: then it was there.
Twice, dead certain, and my hand never near the third box. The note was
inside it, waiting; the guarantee only reads out what was already
written.

\emph{{[}The sun stands where it stands, and its height settles exactly
how far the glass paperweight must throw its shadow across the table.
There the shadow lies at just that reach --- told over before anyone
stooped to measure it, and no hand near the glass.{]}}

\textbf{Anteater:} It is the strongest form your theory can take, and
you have made it well. But one small thing, at another level, before you
are too pleased. \emph{{[}He touches the boxes.{]}} Open just one of the
three by its door: if it sings bright, the other two --- this very
instant --- are a pair, exactly like last autumn's.

\textbf{Achilles:} A pair? Out of a trio?

\textbf{Anteater:} A whole become a pair, by our closing one leg of it
--- but only if it sings bright. If it sings dark, the two that are left
are a stranger pair, with a turn put on one of them. The trio is not
three notes lying side by side; it is one thing on three boxes.

\textbf{Achilles:} \emph{{[}reaching to lift the mask{]}} Then let me
write all six and be done --- a full sealed theory, both ways in for all
three.

\textbf{Tortoise:} Lift the mask, and set your six. \emph{{[}He raises
the trio-mask off the card and lays six counters into the six squares
--- all white, every box bright by the door and bright by the drawer, as
plain and clean a theory as he can write.{]}} There. You have checked
one mask. There are four.

\emph{{[}The sun is as high as it will climb today and lies steady on
the middle window; the ferns there have burned down to a plainer thing,
a clean ruled lattice of lines that cross and cross again. The Tortoise
traces one of them a finger's breadth off the cold glass, corner to
corner, not touching it.{]}}

\begin{center}\rule{0.5\linewidth}{0.5pt}\end{center}

\hypertarget{fourth-challenge-six-notes-four-masks-sarabande}{%
\section{\texorpdfstring{Fourth Challenge --- \emph{Six Notes, Four
Masks}
(Sarabande)}{Fourth Challenge --- Six Notes, Four Masks (Sarabande)}}\label{fourth-challenge-six-notes-four-masks-sarabande}}

\textbf{Achilles:} Four? There is the all-doors asking I just did. And
the others are --- one door, two drawers, in its three arrangements.

\textbf{Tortoise:} Four askings, four masks. Lay them over your six as
they come.

\emph{{[}The card as it stands: six white counters --- every box bright,
by the door and by the drawer. The Tortoise lays the EVEN mask, its
three windows over the three door-squares: three whites, no dark, and
even it wants --- kept. Then the first ODD mask, one window on a
door-square and two on drawer-squares: three whites again, no dark, but
this asking wants an odd count of darks --- broken. The second ODD mask,
its windows over a different door and the two remaining drawers: broken
the same way. The third: broken. Three of the four masks stand wanting a
dark they have not been given.{]}}

\textbf{Tortoise:} \emph{{[}standing a button on each of the three{]}}
Three broken, one kept.

\textbf{Achilles:} Then a dark is wanted. \emph{{[}The card as it
stands: three buttons, on the three ODD masks. He turns the middle box's
drawer counter, white to black; the five others hold still. That one
drawer-square lies under two of the ODD masks --- and both of them,
wanting a dark and now given one, read odd and are kept: two buttons
lift off together. The third ODD mask cannot see this square and still
wants its dark; the EVEN mask, untouched, is still kept. Three buttons
were on the card; two come off at once; one is left standing. Three
become one.{]}} Two gone at a stroke!

\textbf{Achilles:} Now mend the last. \emph{{[}He turns another drawer
counter --- the one under the single surviving button --- white to
black. That mask has its dark now and is kept; but the same square lies
under a second ODD mask that was already kept, and that one, given a
second dark, reads even where it wanted odd. The button does not lift;
it crosses to the mask next door. One counter turned, two verdicts
changed, and still exactly one button on the card.{]}} \ldots And it
only walks. The very same walk as the front of the card, six squares
wide instead of four.

\textbf{Anteater:} Because it is the same machine, only larger. One mask
wants even and three want odd, so the four together must come out odd.
But every counter you have set shows through exactly two of the masks;
so add the four together, whatever you have written, and they come out
even. Odd is not even. One mask is always broken --- or three; never all
four kept.

\emph{{[}Achilles turns one more counter, watches the single button
cross from mask to mask and settle where it always settles, and takes
his hand off the card.{]}}

\textbf{Achilles:} \ldots Oh. It is not that I have written the six
notes badly. It is that there are no six notes to write. No arrangement
of them can keep all four guarantees at once; the arithmetic will not
have it.

\emph{{[}The thumb-print Achilles rubbed clear this morning has long
since bloomed shut again --- from the cold edge inward --- so that by
now, had one set out to wipe the whole pane, it would have frosted over
behind the hand before the hand was done. The pane is never all clear at
once; it cannot be made to be.{]}}

\textbf{Achilles:} Then the workshop's guarantee is simply false. No
three boxes can do this. The Crab has been sold a promise that the
counters prove impossible.

\textbf{Tortoise:} The counters prove that no \emph{written notes} can
do it. \emph{{[}She glances at the case, and then at the hall table.{]}}
There are three envelopes on the table with our names on them.

\emph{{[}The three envelopes lie where they have lain all morning, each
with its name turned up to the light; the folded slip keeps still beside
them under its glass, and the letter lies folded back within reach. The
fire has settled to a low red bed, and the hall is very quiet.{]}}

\begin{center}\rule{0.5\linewidth}{0.5pt}\end{center}

\hypertarget{fifth-challenge-the-round-gigue}{%
\section{\texorpdfstring{Fifth Challenge --- \emph{The Round}
(Gigue)}{Fifth Challenge --- The Round (Gigue)}}\label{fifth-challenge-the-round-gigue}}

\textbf{Achilles:} Our names? \emph{{[}He takes up the three envelopes
and reads the faces.{]}} One for each of us. And the last rule in the
letter, I suppose.

\textbf{Tortoise:} \emph{{[}reading{]}} ``The round, to be played apart.
Give each of three players one box of the trio and one of these sealed
envelopes; the envelope names the way that player must open --- the
door, or the drawer, and no choice about it. The tally slip, under the
glass, tells the referee alone what the dark notes must come to: even,
if all three envelopes say door; odd, if one says door and two say
drawer. Each player goes to a room by himself, opens his box the named
way, hears his one note, and comes back. Lay a counter for each note;
read the tally last. Match it, and the round is won.'' \emph{{[}She
looks up.{]}} We cannot confer once we have parted. And there is nobody
left in the house to carry a word between the rooms; the Crab saw to
that by going to Ostend.

\textbf{Achilles:} Three rooms, then. I will take the attic.

\textbf{Anteater:} The garden room for me.

\textbf{Tortoise:} The cellar. The slip stays here under its glass;
nobody's hand on it, all our hands full of box. \emph{{[}Each takes one
box and one envelope; the slip lies untouched under the paperweight.{]}}
Reader --- you are the referee. You alone will stand where all three
rooms can be seen.

\emph{{[}They climb the stairs together as far as the first landing,
where the three ways part. Achilles stops a moment at Escher's
\textup{Relativity} on the wall.{]}}

\textbf{Achilles:} Look at that --- one house of stairs, and every soul
on it walking to a different down, passing each other on the same step
and never once sharing a floor. \emph{{[}Nobody answers him; they have
already turned away, each to a separate stair.{]}}

\emph{{[}The climbing sun reaches the first landing at last and runs up
the print's crossing stairways; the hall goes still behind them.{]}}

\emph{{[}Achilles, alone in the attic.{]}}

\textbf{Achilles:} You can see all three of us; I can see only this
room, and only what is in my two hands. Here is my envelope. \emph{{[}He
breaks the seal.{]}} ``Door.'' So I swing the door. \emph{{[}The box
sits shut and full on his palm, one note folded up inside it. One hand
steadies it; the other catches the little hinged door in its side and
swings it open --- the box itself does not stir, only the panel turns
back --- and the note springs out into the attic, a single struck tone,
and is gone. What was held is spent; what stays is a small brass box
with nothing further to say, and no lid to try instead.{]}} Dark. There
--- you have mine, whoever is watching from where the whole house shows:
I was asked the door, and the door gave me dark. And do you know, the
word in the envelope told me nothing I could use. Door or drawer, I
should have heard whatever this box was going to say. It might as well
have been a coin that sent me up here.

\emph{{[}The Tortoise, alone in the cellar.{]}}

\textbf{Tortoise:} Mine says drawer. \emph{{[}Her box stands shut and
full on the cellar bench. She holds it down with one claw and, with the
other, draws out the little drawer from underneath --- the box holds its
place, only the drawer slides free --- and as it comes clear one note
lifts from it, and is spent. The box is silent now, the drawer empty,
and the door in its side will never be worth trying.{]}} Bright. That is
all of mine. You have two of the three now; hold on to them.

\emph{{[}The Anteater, alone in the garden room.{]}}

\textbf{Anteater:} My envelope says drawer as well. \emph{{[}He draws it
out from below; a note sounds and is gone.{]}} Bright. You alone have
heard all three now --- a dark and two brights --- and you alone can say
whether we have won, for not one of us in these three rooms knows what
another was asked, or what another heard. I know only my own room. The
round, if it is anything at all, is in none of the three rooms by
itself; it is in the three of them at once --- and only you are standing
where that can be seen.

\emph{{[}The same low sun is standing in all three rooms at once ---
square in the attic's little gable, thin through the cellar grating,
broad on the garden room's frosted glass --- though no one of the rooms
can see into another, and nothing goes along the stairs between them.
Down in the hall a whole pane has come clear at last, and through it the
three frozen roads out of the town lie plain, each striking off its own
way from the one square.{]}}

\emph{{[}The three come back down to the hall, each with a spent
box.{]}}

\textbf{Achilles:} Now the counters. One for each note, laid by each
spent box. \emph{{[}He lays a black counter by his own.{]}} Dark, for
mine. \emph{{[}The Tortoise lays white by hers, the Anteater white by
his.{]}} Bright, and bright. One dark among the three. \emph{{[}He
slides the glass paperweight off the slip and unfolds it.{]}} Odd. One
black. Won.

\textbf{Achilles:} We won it. Every one of us opened blind, and it came
out right to the letter.

\textbf{Tortoise:} One round proves very little by itself, Achilles; be
honest with the reader. Sealed notes win the pair game three deals in
four --- you could win a single round with your card and crow. It is not
this one round. It is the workshop's ledger, that never once shows a
trio breaking the guarantee, and it is your own counters through there,
that say no card of written notes could keep the guarantee even in
principle. Those two together are the thing.

\textbf{Achilles:} Put like that\ldots{} The button would not vanish ---
written notes stop at three in four, and cannot keep all four
trio-guarantees at all. And the boxes keep them, always, and win apart
with nobody carrying a word between the rooms. So the most that written
notes can do, the boxes go past; and if the boxes go past what any
writing could do, then nothing was written. The notes were never in the
boxes.

\emph{{[}The winter sun has come round the hall and lies long and low
across the floor; there are still hours of light in it. The three spent
boxes go back into their three nests in the case, the counters are
gathered into a little heap, and the last of the pot goes round while
nobody hurries.{]}}

\textbf{Achilles:} But that is the part I cannot hold. If it wasn't in
there, and nothing crossed the house --- where was it?

\textbf{Tortoise:} In neither box. Draw the three of them, and look at
the place where the lines meet.

\emph{{[}On the table the case stands with its lid down, the three
envelopes open where their names had faced the light, the slip unfolded
beside them, and the paperweight set aside with nothing left under it to
hold. On the tall windows the frost has given up the middles and drawn
back into the corners, so that the clear glass shows the afternoon
plainly through. On the wall at the first landing the print has quieted
--- a plain frame with a little sun caught in its glass, its crossing
stairs gone dim. The fire ticks as it cools. Nobody at the table reaches
for anything, and for the moment there is nothing on it that wants
opening.{]}}

\hypertarget{chapter-7-inequalities-nature-breaks}{%
\chapter{Chapter 7: Inequalities Nature
Breaks}\label{chapter-7-inequalities-nature-breaks}}

\emph{--- Partita for Three Players, Far Apart ---}

\emph{Part II: Wires Together}

\emph{Escher: ``Relativity'' --- one house of stairs, three directions
of down; the walkers pass and never meet}

\begin{center}\rule{0.5\linewidth}{0.5pt}\end{center}

At the end of the last chapter we let a suspicion stand: perhaps two
things that always agree agree because the answers were written into
them beforehand. Can that suspicion survive an honest test?

The suspicion is the most reasonable one in physics. Two halves of a
Bell pair, questioned far apart, gave matching answers every time; and
the plainest account of matching answers is a pair of sealed notes,
written at the source and carried away, one in each half. Chapter 6
could not tell the two accounts apart, and said so. Every number we
computed --- the fair coin in each half, the perfect agreement between
them --- the sealed notes reproduce exactly. If a bent wire and a pair
of sealed notes differ at all, the difference is hiding somewhere we
have not yet looked.

This chapter looks there. The place to look, John Bell saw in 1964, is
not in one question repeated but in \emph{several} questions, put in
different combinations to parties who are too far apart to confer once
the questioning starts. A sealed note must have an answer ready for
every question it might be asked, decided in advance, the same whoever
does the asking. That is a strong constraint, and Bell found a way to
measure it: a quantity that any collection of sealed notes must keep
below a fixed bound, and that a shared quantum state can push past.

We shall meet the bound first as a game. Two partners, then three, are
asked questions in separate rooms and must answer without conferring;
they win or lose by a rule fixed in advance; and we count how often the
best possible preparation lets them win. With sealed notes the count has
a ceiling. With a shared quantum state --- a bent wire, and then a new
thing, a single point where three wires meet --- the ceiling is broken,
and in the three-party game it is broken \emph{completely}: they win not
more often, but every single time. No collection of sealed notes can do
that. So the notes were not written.

On the blackboard this will be a short, stubborn piece of arithmetic
about the numbers \(+1\) and \(-1\). In the sketchbook it will be a
single new drawing --- a dot with three legs --- and one new rule for
pushing boxes through it. Neither strand, we should say plainly at the
outset, proves that nature is nonlocal; each proves an equality of
processes, and the arithmetic does the rest. But between them they close
the door that Chapter 6 left open.

\begin{center}\rule{0.5\linewidth}{0.5pt}\end{center}

\hypertarget{the-hidden-note-theory}{%
\section{§7.1 The Hidden-Note Theory}\label{the-hidden-note-theory}}

Achilles ended the last chapter as a realist, and a realist has a
theory. It is this: when a source prepares two halves and sends them
apart, it writes into each half, in advance, an answer for every
question that half might be asked. The halves then part. Each carries
its answers like a sealed note; whatever is later asked, the note is
simply read. Nothing needs to pass between the halves when they are
questioned, because everything was settled at the source.

This is worth stating precisely, because it is exactly what we are going
to test.

\textbf{On the blackboard.} A \textbf{deterministic local strategy} for
two parties is a pair of functions: one, \(a(x)\), that gives the first
party's answer to each question \(x\) it might receive, and another,
\(b(y)\), for the second party. Both functions are fixed before the
parties separate. For three parties there are three: \(a(x)\), \(b(y)\),
\(c(z)\). The parties may also share, before parting, a table of such
strategies and a rule for choosing among them at random --- but a random
choice among fixed strategies is a \emph{mixture}, and a mixture of
strategies can do no better, on average, than the best single strategy
in it. So to find the ceiling of the hidden-note theory it is enough to
examine the deterministic strategies one by one.

Behind this is a criterion that Einstein, Podolsky and Rosen made famous
in 1935, and that we shall use in the reader's own words: \emph{if the
answer to a question can be predicted with certainty, without touching
the thing that will answer it, then that answer was a real property of
the thing all along.} A sealed note is exactly a thing whose every
answer can be predicted without touching it --- by reading the copy at
the source, or, as we shall see in §7.4, by asking the other halves. The
hidden-note theory is the claim that a Bell pair, and a trio, are
collections of sealed notes in this precise sense.

\textbf{In the sketchbook.} We need no new element for the local theory.
Two parties who share nothing but a table agreed in advance are, as a
picture, two separate state boxes --- one on each party's wire --- with
nothing joining the top of the page to the bottom. It is the
\emph{product state} of §6.1, drawn as a picture that falls into
disconnected pieces. Whatever effects are later applied to the two
halves, the closed picture factors into a top piece times a bottom
piece, and the number it yields is a product of two independent numbers.
That is the whole content of ``the parties cannot confer'': the picture
does not connect them.

\textbf{\emph{Why this definition.}} We phrase the local theory as fixed
answer-functions, and not as ``the particles have definite values'',
because answer-functions are the thing a game can catch. They commit the
theory to a number for every question in advance, and a commitment made
in advance can be checked against a question chosen at the last moment.

\begin{center}\rule{0.5\linewidth}{0.5pt}\end{center}

\hypertarget{the-game-and-the-sealed-notes}{%
\section{§7.2 The Game and the Sealed
Notes}\label{the-game-and-the-sealed-notes}}

Here is the game. Two partners, far apart, each receive one bit of
question --- call the first partner's question \(x\) and the second's
\(y\), each either \(0\) or \(1\). Each answers with one bit, \(a\) and
\(b\). They win the round exactly when

\[a\oplus b \;=\; x\cdot y,\]

where \(\oplus\) is addition modulo \(2\) and \(x\cdot y\) is the
ordinary product, which is \(1\) only when both questions are \(1\). In
words: for three of the four question pairs the partners must give
\emph{equal} answers, and for the one remaining pair, \(x=y=1\), they
must give \emph{unequal} answers. They may plan anything beforehand;
they may not communicate once the questions are handed out. (This is the
game of Clauser, Horne, Shimony and Holt, in the form used by Brunner
and colleagues, \(a\oplus b=x\cdot y\), whose one ``unequal'' pair is
\(x=y=1\). We have only chosen which physical question --- lid or door
--- each label names, to suit the story to come; the game is the same.)

\hypertarget{the-example-that-fixes-the-ceiling}{%
\subsection{The example that fixes the
ceiling}\label{the-example-that-fixes-the-ceiling}}

Suppose the partners agree: \emph{always answer \(0\)}, whatever the
question. Then \(a\oplus b=0\) always. They win the three pairs that ask
for equal answers and lose the one, \(x=y=1\), that asks for unequal
answers. Three wins out of four. Can they patch the fourth without
spoiling the other three? Change the second partner to answer \(1\) when
\(y=1\). Now the pair \(x=0,y=1\) has answers \(0\) and \(1\) ---
unequal --- and it was supposed to be equal. They have mended one pair
and broken another. The count stays at three. Every patch we try does
this: it moves the single loss around the table without removing it.

That is not an accident, and here is why.

\begin{quote}
\textbf{Theorem 7.1 (the sealed-note ceiling).} No deterministic local
strategy wins more than three of the four question pairs. Equivalently,
the CHSH quantity \[S \;=\; E_{00}+E_{01}+E_{10}-E_{11}\] --- where
\(E_{xy}=+1\) if the strategy answers pair \((x,y)\) with equal bits and
\(-1\) if unequal --- satisfies \(S\le 2\) for every strategy, and a
mixture of strategies can only lower \(|S|\).
\end{quote}

\emph{Proof idea.} The four winning conditions are four equations in the
same four answer-bits. Add them modulo \(2\): the left sides cancel in
pairs and vanish, but the right sides do not. The four conditions
therefore cannot all hold at once; one must fail.

\emph{Proof (on the blackboard).} A deterministic strategy is four bits:
\(a_0=a(0)\), \(a_1=a(1)\), \(b_0=b(0)\), \(b_1=b(1)\). The four winning
conditions are

\[a_0\oplus b_0=0,\quad a_0\oplus b_1=0,\quad a_1\oplus b_0=0,\quad a_1\oplus b_1=1.\]

Add all four left-hand sides modulo \(2\). Each of \(a_0,a_1,b_0,b_1\)
appears exactly twice, so the sum is \(0\). Add the four right-hand
sides: \(0\oplus0\oplus0\oplus1=1\). If all four conditions held, we
would have \(0=1\). So at most three hold, and the best strategy wins
exactly three of the four pairs. There are \(16\) deterministic
strategies; each wins three pairs at most, so each has \(S\le 2\) (three
agreements at \(+1\) and one disagreement, or fewer). A random mixture
assigns probabilities to strategies and averages their \(S\) values,
which cannot exceed the largest, \(2\). \(\;\blacksquare\) See
\texttt{Code/Ch7/SealedNotes.py}, which enumerates all \(16\) strategies
and confirms the best wins three of four, that every strategy has
\(|S|=2\), and that changing one answer-bit always flips exactly two of
the four verdicts --- the reason a patch never helps.

\textbf{In the sketchbook.} The sketchbook has nothing new to add here,
and it is honest to say so. A sealed-note strategy is a product picture;
the number for each question pair is a product of a top number (the
first partner's) and a bottom number (the second's); and no product
picture can beat the ceiling, because the ceiling is a fact about
products of \(\pm1\) numbers, which the disconnected picture faithfully
computes and cannot escape. The picture will earn its keep in this
chapter, but not here: the classical bound is a blackboard result
through and through.

\begin{quote}
\textbf{Dirac's Blackboard.}

\emph{Why one loss cannot be removed.} Write the four verdicts as the
parities \(v_{xy}=a_x\oplus b_y\oplus(x\cdot y)\); a pair is \emph{won}
when \(v_{xy}=0\). Sum all four modulo \(2\):
\[\begin{aligned}&v_{00}\oplus v_{01}\oplus v_{10}\oplus v_{11}\\&\quad=(a_0\oplus a_0\oplus a_1\oplus a_1)\oplus(b_0\oplus b_1\oplus b_0\oplus b_1)\oplus 1=1.\end{aligned}\]
An odd number of the \(v_{xy}\) must therefore equal \(1\): one loss, or
three, never zero and never two. Flip a single answer-bit, say \(a_0\):
it sits in \(v_{00}\) and \(v_{01}\) and in no others, so exactly two
verdicts flip. A loss can be shuffled to a neighbour or split into
three; it cannot be made to vanish, because vanishing would need an even
change to an odd total. Everything the classical player can do lives
inside this one equation.
\end{quote}

\textbf{\emph{Why this game.}} The game is the same statement as Bell's
inequality \(S\le2\), in newer dress. We lead with the game because it
makes the constraint vivid: a sealed note must answer \emph{before} it
knows which of its questions was chosen, and the game is built to punish
exactly that. The number \(S\) is a bookkeeper's summary; the odd number
of losses is the thing itself.

\begin{quote}
\textbf{Feynman's Notebook.}

\emph{From a paradox to an inequality.} The suspicion this chapter tests
was sharpened in 1935 by Einstein, Podolsky and Rosen, who argued that
if a distant measurement can predict a result with certainty, that
result must already be a real property --- so quantum theory, they held,
was incomplete. For three decades the argument stayed philosophical.
Then in 1964 Bell turned it into arithmetic: a quantity that any local
theory must keep below a bound. Clauser, Horne, Shimony and Holt cast
that bound in the experiment-ready form of 1969 our game rehearses. The
first laboratory tests followed --- Freedman and Clauser in 1972, Aspect
and colleagues in 1982 --- each finding the bound broken, though each
still leaning on an extra assumption. We take the history through
Brunner and colleagues' review.
\end{quote}

\begin{center}\rule{0.5\linewidth}{0.5pt}\end{center}

\hypertarget{the-wedge}{%
\section{§7.3 The Wedge}\label{the-wedge}}

Sealed notes stop at three in four. Now give the partners a bent wire
instead of a table of answers, and one small tool: a wedge, under the
second partner's box, that tilts the question it is asked to a direction
halfway between the two.

\textbf{On the blackboard.} The partners share the Bell state
\(\ket{\Phi^+}=\tfrac1{\sqrt2}\,\text{cup}\), the cup being
unnormalized, \(\text{cup}=\ket{00}+\ket{11}\), as always in this book.
Each partner has two questions. We take the first partner's question
\(x=0\) to be the \emph{lid} question, measured by the observable \(Z\),
and \(x=1\) the \emph{door} question, measured by \(X\); each is a
\(\pm1\)-valued observable, \(+1\) for a bright answer (bit \(0\)),
\(-1\) for a dark one (bit \(1\)). The second partner's questions are
the door (\(y=0\)) and the lid (\(y=1\)) --- named in this order so that
a single wedge serves both.

The wedge is a real rotation of the sphere of Chapter 3. We introduce
the box

\[R_{\theta}=\begin{pmatrix}\cos\theta & -\sin\theta\\[2pt] \sin\theta & \cos\theta\end{pmatrix},\]

a real matrix, \emph{not symmetric}, with
\(R_{\theta}^{T}=R_{-\theta}\). The second partner applies \(R_{\pi/8}\)
to their half before asking. The matrix turns through \(\pi/8\); the
sphere turns through \emph{twice} that, \(\pi/4\), which is halfway from
the lid (\(Z\)) to the door (\(X\)) --- a right angle apart on the
sphere. The effect, as one checks directly, is to replace the second
partner's two plain observables by tilted ones:

\[R_{\pi/8}^{T}\,Z\,R_{\pi/8}=\frac{Z-X}{\sqrt2},\qquad R_{\pi/8}^{T}\,X\,R_{\pi/8}=\frac{Z+X}{\sqrt2}.\]

The relevant number for the game is the \textbf{correlation}
\(E_{xy}=\braket{\Phi^+|A_x\otimes B_y|\Phi^+}\), where \(A_x\) is the
first partner's observable (\(Z\) or \(X\)) and \(B_y\) is the second
partner's tilted observable. For the Bell state there is a quick way to
compute it, and it is Chapter 6's sliding trick:
\(\braket{\Phi^+|A\otimes B|\Phi^+}=\tfrac12\braket{\text{cup}|A\otimes B|\text{cup}}\),
and by Theorem 6.4,
\((A\otimes B)\,\text{cup}=(I\otimes BA^{T})\,\text{cup}\), so the whole
thing collapses to
\(\tfrac12\operatorname{Tr}(BA^{T})=\tfrac12\operatorname{Tr}(A^{T}B)\).
Since \(Z\) and \(X\) are symmetric,
\(E_{xy}=\tfrac12\operatorname{Tr}(A_xB_y)\). Now read off the four
numbers, using \(\operatorname{Tr}(ZZ)=\operatorname{Tr}(XX)=2\) and
\(\operatorname{Tr}(ZX)=0\):

\[E_{00}=\tfrac12\operatorname{Tr}\!\Big(Z\cdot\tfrac{Z+X}{\sqrt2}\Big)=\tfrac1{\sqrt2},\quad
E_{01}=\tfrac12\operatorname{Tr}\!\Big(Z\cdot\tfrac{Z-X}{\sqrt2}\Big)=\tfrac1{\sqrt2},\]
\[E_{10}=\tfrac12\operatorname{Tr}\!\Big(X\cdot\tfrac{Z+X}{\sqrt2}\Big)=\tfrac1{\sqrt2},\quad
E_{11}=\tfrac12\operatorname{Tr}\!\Big(X\cdot\tfrac{Z-X}{\sqrt2}\Big)=-\tfrac1{\sqrt2}.\]

So \(E=(+,+,+,-)\tfrac1{\sqrt2}\), and

\[S=E_{00}+E_{01}+E_{10}-E_{11}=\frac{4}{\sqrt2}=2\sqrt2\;>\;2.\]

The winning probability for each pair is
\(p_{\text{win}}=\tfrac12(1+(-1)^{x\cdot y}E_{xy})=\tfrac12\big(1+\tfrac1{\sqrt2}\big)=\cos^2(\pi/8)\approx0.8536\)
--- the same for all four pairs, and above the sealed-note ceiling of
\(\tfrac34\). See \texttt{Code/Ch7/Wedge.py}, which builds all sixteen
amplitudes and confirms \(p_{\text{win}}=\cos^2(\pi/8)\) for each of the
four pairs and \(S=2\sqrt2\), and that each partner's own answer is a
fair coin for every question --- the marginal is \(\tfrac12\), so
neither partner can read the other's question from their own results.

\begin{quote}
\textbf{Theorem 7.2 (the wedge strategy).} With the shared Bell state
and the wedge \(R_{\pi/8}\) on the second partner's half, every one of
the four question pairs is won with probability
\(\cos^2(\pi/8)=\tfrac12\big(1+\tfrac1{\sqrt2}\big)\), and
\(S=2\sqrt2\). No quantum strategy does better: \(S\le 2\sqrt2\) for
every state and every choice of \(\pm1\)-observables.
\end{quote}

The last sentence is \textbf{Tsirelson's bound} (Cirel'son, 1980; we
take it, and its history, through Brunner's review). We \emph{state} it
and do not prove it here; the proof is a fact about the norm of an
operator built from the four observables, and it is not needed for
anything this chapter does. What we need is only that \(2\sqrt2\) is
more than \(2\).

\textbf{In the sketchbook.} The whole game is one closed picture. The
pair is the Bell state \(\ket{\Phi^+}=\tfrac1{\sqrt2}\,\)cup --- the cup
unnormalized, \(\text{cup}=\ket{00}+\ket{11}\), so the
\(\tfrac1{\sqrt2}\) rides in front of the picture; the second partner's
lower leg carries the wedge \(R_{\pi/8}\); and each leg ends in an
effect for that partner's answer, \(a\) on top and \(b\) below:

\begin{center}
$\tfrac1{\sqrt2}\;$\begin{tikzpicture}[sd]
  \sdcup{1}{2.5}{0}
  \draw[wire] (1,2.5) -- (5.4,2.5);
  \draw[wire] (1,0) -- (5.4,0);
  \node[sdeffect] at (5.4,2.5) {$a$};
  \node[sdbox] at (3.2,0) {$R_{\pi/8}$};
  \node[sdeffect] at (5.4,0) {$b$};
\end{tikzpicture}
\end{center}

Read left to right and closed at both ends, it is a number --- the
amplitude of the answer pair \((a,b)\) --- and its squared modulus is
their probability, by Chapter 5's rule. The one move worth savouring is
where the wedge may stand. Sliding it round the cup (Theorem 6.4, with
\(R_{\pi/8}^{T}=R_{-\pi/8}\)) carries it from the lower leg to the
upper, turned the other way:

\[\begin{tikzpicture}[sd]
  \sdcup{0.8}{2.5}{0}
  \draw[wire] (0.8,2.5) -- (4.2,2.5);
  \draw[wire] (0.8,0) -- (4.2,0);
  \node[sdbox] at (2.6,0) {$R_{\pi/8}$};
\end{tikzpicture}
\;=\;
\begin{tikzpicture}[sd]
  \sdcup{0.8}{2.5}{0}
  \draw[wire] (0.8,2.5) -- (4.4,2.5);
  \draw[wire] (0.8,0) -- (4.4,0);
  \node[sdbox] at (2.6,2.5) {$R_{-\pi/8}$};
\end{tikzpicture}\]

so the wedge may sit under either box, tilted oppositely, and the game
is unchanged. \texttt{Wedge.py} checks the two placements give the same
closed pictures exactly, and --- because \(R_{\theta}\) is real, so that
a transpose could hide behind a dagger --- first checks the sliding rule
with a random complex, non-symmetric box, where transpose and dagger
genuinely differ.

\begin{quote}
\textbf{Dirac's Blackboard.}

\emph{One pair, in full, with the \(\tfrac1{\sqrt2}\) shown.} Take
\((x,y)=(0,0)\): the first partner asks the lid (\(Z\)-basis, answers
\(\ket0,\ket1\)), the second the door, after the wedge. The pair, before
either asks, is
\((I\otimes R_{\pi/8})\ket{\Phi^+}=\tfrac1{\sqrt2}\big(\ket0\otimes R_{\pi/8}\ket0+\ket1\otimes R_{\pi/8}\ket1\big)\),
with \(R_{\pi/8}\ket0=\cos\tfrac\pi8\ket0+\sin\tfrac\pi8\ket1\). The
amplitude for the bright--bright outcome \((\ket0,\ket+)\) is
\[\begin{aligned}\big(\bra0\otimes\bra+\big)(I\otimes R_{\pi/8})\ket{\Phi^+}&=\tfrac1{\sqrt2}\,\bra+R_{\pi/8}\ket0\\&=\tfrac1{\sqrt2}\cdot\tfrac{\cos\frac\pi8+\sin\frac\pi8}{\sqrt2}=\tfrac{\cos\frac\pi8+\sin\frac\pi8}{2}.\end{aligned}\]
Its square is
\(\tfrac14\big(1+\sin\tfrac\pi4\big)=\tfrac14\big(1+\tfrac1{\sqrt2}\big)\).
The dark--dark amplitude is
\(-\tfrac12(\cos\tfrac\pi8+\sin\tfrac\pi8)\), the same square. This pair
wins on agreement, so
\(p_{\text{win}}=\tfrac14\big(1+\tfrac1{\sqrt2}\big)+\tfrac14\big(1+\tfrac1{\sqrt2}\big)=\tfrac12\big(1+\tfrac1{\sqrt2}\big)=\cos^2\tfrac\pi8\).
Every \(\tfrac1{\sqrt2}\) in sight is the Bell state's normalization.
\end{quote}

\textbf{\emph{Why this definition of \(R_{\theta}\).}} We use a plain
real rotation, not a fraction of a bit-flip, because we want a genuinely
\emph{intermediate} question --- one tilted between lid and door --- and
only a rotation of the sphere tilts one axis toward another. A fraction
of \(X\) would leave the question pointing the wrong way.

\begin{center}\rule{0.5\linewidth}{0.5pt}\end{center}

\hypertarget{a-node-with-three-legs}{%
\section{§7.4 A Node with Three Legs}\label{a-node-with-three-legs}}

Two boxes from one bench were a pair. Three boxes from one bench ---
made together, sent together --- are something the wire notation cannot
yet draw. We need one new element, and it is the only one this chapter
adds.

\textbf{In the sketchbook.} Imagine a single point on the page from
which \emph{three} wires leave, all to the right: not three separate
wires laid side by side, and not a wire that branches into a fork, but
one small solid dot where three legs meet as equals. We draw it as a
filled dot with three legs:

\begin{center}
\begin{tikzpicture}[sd]
  \node[sdnode] at (0,1.5) {};
  \draw[wire] (0,1.5) to[out=90,in=180] (1.2,3) -- (4,3);
  \draw[wire] (0,1.5) -- (4,1.5);
  \draw[wire] (0,1.5) to[out=-90,in=180] (1.2,0) -- (4,0);
\end{tikzpicture}
\end{center}

Read left to right, nothing comes in and three wires go out: it is a
state on three wires. We call it, plainly and only, \textbf{the node},
or the three-legged node. It has no label, no colour, and no number
written on it; it is always the same thing.

\textbf{On the blackboard.} We \emph{define} the node to be the
unnormalized state

\[N \;=\; \ket{000}+\ket{111},\]

whose squared length is \(\braket{N|N}=2\). Exactly as the cup is
\(\sqrt2\) times the Bell state, the node is \(\sqrt2\) times the state
a physicist would normalize:

\[\tfrac1{\sqrt2}\,N \;=\; \tfrac1{\sqrt2}\big(\ket{000}+\ket{111}\big),\]

the three-party state of Greenberger, Horne and Zeilinger. We keep the
node unnormalized for the same reason we kept the cup unnormalized: so
that the rules for pushing things through it hold with no stray scalars.
Wherever the node is drawn in this book it means
\(\ket{000}+\ket{111}\), and we write the \(\tfrac1{\sqrt2}\) beside it
whenever the normalized state is wanted.

The node is not three cups, and it is not a product. Its three legs
agree --- every basis term is \(000\) or \(111\) --- but they agree
three at a time, in a way no pair of them can hold alone.
(\texttt{Code/Ch7/Trio.py} confirms \(N\) is not a product across any
way of splitting the three legs.)

\textbf{\emph{Why this definition.}} We could have tried to build a trio
from cups, joining pairs. No joining of cups gives
\(\ket{000}+\ket{111}\); the all-agree-together state is genuinely new,
and it deserves its own mark on the page. The dot is the least we can
draw: a place where wires meet, carrying no data of its own.

\hypertarget{what-the-node-lets-a-party-predict}{%
\subsection{What the node lets a party
predict}\label{what-the-node-lets-a-party-predict}}

Here is the guarantee the trio comes with, and the realist's best
argument. The three boxes are made together; when all three are opened
the same way (all by the door), the number of dark notes is always
\emph{even}; when one is opened by the door and the other two by the
drawer, the number of dark notes is always \emph{odd}. Never broken,
says the workshop's ledger.

Grant the realist the guarantee. Then any one party can \emph{predict}
their own answer with certainty from the other two --- for the all-door
case, two answers fix the third by evenness --- without touching their
own box. By EPR's criterion, then, the answer was written inside. This
is the theory at its strongest, and §7.6 will need every bit of that
strength to overturn it. For now, two facts about the node make the
guarantee precise, and both are exact equalities.

\begin{quote}
\textbf{Theorem 7.3 (sliding and merging through the node).} A
\emph{diagonal} box slides freely from one leg of the node to any other,
unchanged, and two diagonal boxes on one leg merge into their product.
For diagonal \(D\) and \(D'\),
\[\begin{aligned}(D\otimes I\otimes I)\,N&=(I\otimes D\otimes I)\,N=(I\otimes I\otimes D)\,N,\\ (D\otimes D'\otimes I)\,N&=(I\otimes I\otimes DD')\,N.\end{aligned}\]
A box that is not diagonal does \textbf{not} slide.
\end{quote}

\emph{Proof idea.} The node lives on just two terms, \(\ket{000}\) and
\(\ket{111}\). A diagonal box multiplies each basis state by a number
depending only on the bit on its leg; but on those two terms all three
legs carry the same bit, so it does not matter which leg the box sits
on. A non-diagonal box moves weight between \(\ket0\) and \(\ket1\),
which breaks the all-same pattern, and then the leg matters.

\emph{Proof (on the blackboard).} Let
\(D=\operatorname{diag}(d_0,d_1)\). On \(N\)'s two terms,
\((D\otimes I\otimes I)\ket{000}=d_0\ket{000}\) and
\((D\otimes I\otimes I)\ket{111}=d_1\ket{111}\); putting \(D\) on the
second or third leg gives the very same \(d_0\ket{000}+d_1\ket{111}\),
because each term carries one bit across all three legs. For two boxes,
\((D\otimes D'\otimes I)N=d_0d'_0\ket{000}+d_1d'_1\ket{111}=(I\otimes I\otimes DD')N\).
\(\;\blacksquare\) \texttt{Trio.py} checks both with random
\emph{complex} diagonal boxes, and checks that a random non-diagonal box
fails to slide.

\textbf{In the sketchbook.} The rule is one picture. A diagonal box
\(D\) on any leg is the same drawing wherever it sits, and two such
boxes collect onto a single leg as their product:

\[\begin{tikzpicture}[sd]
  \node[sdnode] at (0,1.5) {};
  \draw[wire] (0,1.5) to[out=90,in=180] (0.9,3) -- (3.4,3);
  \draw[wire] (0,1.5) -- (3.4,1.5);
  \draw[wire] (0,1.5) to[out=-90,in=180] (0.9,0) -- (3.4,0);
  \node[sdbox] at (2.3,3) {$D$};
\end{tikzpicture}
\;=\;
\begin{tikzpicture}[sd]
  \node[sdnode] at (0,1.5) {};
  \draw[wire] (0,1.5) to[out=90,in=180] (0.9,3) -- (3.4,3);
  \draw[wire] (0,1.5) -- (3.4,1.5);
  \draw[wire] (0,1.5) to[out=-90,in=180] (0.9,0) -- (3.4,0);
  \node[sdbox] at (2.3,1.5) {$D$};
\end{tikzpicture}
\;=\;
\begin{tikzpicture}[sd]
  \node[sdnode] at (0,1.5) {};
  \draw[wire] (0,1.5) to[out=90,in=180] (0.9,3) -- (3.4,3);
  \draw[wire] (0,1.5) -- (3.4,1.5);
  \draw[wire] (0,1.5) to[out=-90,in=180] (0.9,0) -- (3.4,0);
  \node[sdbox] at (2.3,0) {$D$};
\end{tikzpicture}\]

\[\begin{tikzpicture}[sd]
  \node[sdnode] at (0,1.5) {};
  \draw[wire] (0,1.5) to[out=90,in=180] (0.9,3) -- (3.4,3);
  \draw[wire] (0,1.5) -- (3.4,1.5);
  \draw[wire] (0,1.5) to[out=-90,in=180] (0.9,0) -- (3.4,0);
  \node[sdbox] at (2.3,3) {$D$};
  \node[sdbox] at (2.3,1.5) {$D'$};
\end{tikzpicture}
\;=\;
\begin{tikzpicture}[sd]
  \node[sdnode] at (0,1.5) {};
  \draw[wire] (0,1.5) to[out=90,in=180] (0.9,3) -- (3.6,3);
  \draw[wire] (0,1.5) -- (3.6,1.5);
  \draw[wire] (0,1.5) to[out=-90,in=180] (0.9,0) -- (3.6,0);
  \node[sdbox] at (2.4,0) {$DD'$};
\end{tikzpicture}\]

A box that is not diagonal cannot be moved this way; the drawing would
change its value.

\begin{quote}
\textbf{Theorem 7.4 (closing one leg).} Closing one leg of the node with
a bright \emph{door}-effect leaves the other two as a cup, and a dark
door-effect a cup with a \(Z\) box on one leg; a bright
\emph{lid}-effect leaves a plain product:
\[\begin{aligned}(\bra+\otimes I\otimes I)\,N&=\tfrac1{\sqrt2}\,\text{cup}=\ket{\Phi^+},\\ (\bra-\otimes I\otimes I)\,N&=\tfrac1{\sqrt2}(Z\otimes I)\,\text{cup},\\ (\bra0\otimes I\otimes I)\,N&=\ket{00}.\end{aligned}\]
\end{quote}

\emph{Proof.} With \(\bra+=\tfrac1{\sqrt2}(\bra0+\bra1)\), applying it
to the first leg of \(N=\ket{000}+\ket{111}\) keeps
\(\tfrac1{\sqrt2}(\ket{00}+\ket{11})=\tfrac1{\sqrt2}\text{cup}\); with
\(\bra-=\tfrac1{\sqrt2}(\bra0-\bra1)\) the second term picks up a minus,
giving
\(\tfrac1{\sqrt2}(\ket{00}-\ket{11})=\tfrac1{\sqrt2}(Z\otimes I)\text{cup}\);
with the lid-effect \(\bra0\), only \(\ket{00}\) survives.
\(\;\blacksquare\) (See \texttt{Trio.py}.)

\textbf{In the sketchbook.} A bright door-effect on one leg bends the
other two into a cup --- the first cup of this chapter, and unnormalized
as ever, \(\text{cup}=\ket{00}+\ket{11}\):

\[\begin{tikzpicture}[sd]
  \node[sdnode] at (0,1.5) {};
  \draw[wire] (0,1.5) to[out=90,in=180] (0.9,3) -- (3.2,3);
  \draw[wire] (0,1.5) -- (3.2,1.5);
  \draw[wire] (0,1.5) to[out=-90,in=180] (0.9,0) -- (3.2,0);
  \node[sdeffect] at (3.2,3) {$+$};
\end{tikzpicture}
\;=\;\tfrac1{\sqrt2}\;
\begin{tikzpicture}[sd]
  \sdcup{0.6}{1.5}{0}
  \draw[wire] (0.6,1.5) -- (2.8,1.5);
  \draw[wire] (0.6,0) -- (2.8,0);
\end{tikzpicture}\]

so the trio, with one door opened bright, is a Bell pair on the wires
that remain.

This is the Anteater's remark, at which level to ask: open one of the
three by its door, and if it sings bright, the other two are a Bell pair
--- a \emph{pair}, though a moment ago there was no pair, only a trio.
If it sings dark, the other two are a pair with a \(Z\) turned on one
leg. The whole was one thing on three wires; closing a leg leaves one
thing on two.

\hypertarget{the-drawer-is-the-y-question}{%
\subsection{\texorpdfstring{The drawer is the \(Y\)
question}{The drawer is the Y question}}\label{the-drawer-is-the-y-question}}

The trio's second way in --- the drawer --- is the \(Y\) question, and
it is where the conjugate of Chapter 6 comes back to matter. The
drawer-effect asks along the states
\(\ket{\pm i}=\tfrac1{\sqrt2}(\ket0\pm i\ket1)\), and its bra is a
door-effect behind a box:

\[\bra{\pm i}=\bra{\pm}\,S^{\dagger},\qquad S^{\dagger}=\begin{pmatrix}1&0\\0&-i\end{pmatrix}.\]

This is an equality that a real test would miss: \(\ket{\pm i}\) has an
imaginary amplitude, and only against a non-real state does
\(\bra{\pm}S^{\dagger}\) differ from \(\bra{\pm}\).
(\texttt{Code/Ch7/SixNotes.py} checks it with the genuinely non-real
\(\ket{\pm i}\), and checks that dropping the conjugation fails.) In a
picture, a drawer-effect on a leg is an \(S^{\dagger}\) box followed by
a door-effect --- a fact we shall spend in §7.6.

It also settles a small debt from Chapter 6, where the Crab was told to
build his pairs with only one way in besides the lid, and warned that
not every third way would agree. Here is why: a pair opened by the
\emph{drawer} on both sides would \emph{disagree}, because
\((Y\otimes Y)\,\text{cup}=-\text{cup}\) (checked in
\texttt{SixNotes.py}). The minus sign that will bring down the realist
in §7.6 is already visible here, in a single pair.

\begin{quote}
\textbf{Coecke's Sketchbook.}

\emph{A dot is not a cup, and its normalization is its own.} Draw the
node and it is tempting to read it as a wire that forks --- one line
splitting into three. It is not. A forked wire would copy a state, and
no box copies an unknown quantum state; the node does no such thing. It
is a \emph{state} on three legs, born equal, with no ``first'' leg the
others come from. Its three legs meet as three, and that is the whole of
it. One caution to carry: the node is unnormalized,
\(N=\ket{000}+\ket{111}\) with \(\braket{N|N}=2\), exactly as the cup is
\(\ket{00}+\ket{11}\) with length \(\sqrt2\). When a probability is
wanted, the \(\tfrac1{\sqrt2}\) must be paid in, once, just as for the
cup. Draw the dot freely; remember the number when you cash out.
\end{quote}

\begin{center}\rule{0.5\linewidth}{0.5pt}\end{center}

\hypertarget{the-rosetta-table-6}{%
\section{§7.5 The Rosetta Table}\label{the-rosetta-table-6}}

The rows this chapter adds to the running dictionary, numbered on from
Chapter 6's forty-five. One introduces a new picture element --- the
node; one is a new rewrite rule --- sliding and merging through it; the
rest record boxes and equalities we shall lean on.

\begingroup\renewcommand{\arraystretch}{1.5}

\begin{longtable}[]{@{}
  >{\raggedright\arraybackslash}p{(\columnwidth - 4\tabcolsep) * \real{0.2727}}
  >{\raggedright\arraybackslash}p{(\columnwidth - 4\tabcolsep) * \real{0.3636}}
  >{\raggedright\arraybackslash}p{(\columnwidth - 4\tabcolsep) * \real{0.3636}}@{}}
\toprule\noalign{}
\begin{minipage}[b]{\linewidth}\raggedright
\#
\end{minipage} & \begin{minipage}[b]{\linewidth}\raggedright
Algebra (Dirac)
\end{minipage} & \begin{minipage}[b]{\linewidth}\raggedright
Picture (Coecke)
\end{minipage} \\
\midrule\noalign{}
\endhead
\bottomrule\noalign{}
\endlastfoot
46 & \(N=\ket{000}+\ket{111}\) (unnormalized; length \(\sqrt2\)); GHZ
state \(=\tfrac1{\sqrt2}N\) & a solid dot with three legs --- a state on
three wires --- \emph{new element} \\
47 & \((D\otimes I\otimes I)N=(I\otimes D\otimes I)N\), \(D\) diagonal;
\((D\otimes D'\otimes I)N=(I\otimes I\otimes DD')N\) & a diagonal box
slides through the dot from leg to leg; two on a leg merge into one ---
\textbf{sliding through the node} \\
48 & \((\bra+\otimes I\otimes I)N=\tfrac1{\sqrt2}(\ket{00}+\ket{11})\);
a dark leg leaves it with a \(Z\) box & a bright door-effect on one leg
leaves \(\tfrac1{\sqrt2}\) cup \\
49 & \(\bra{\pm i}=\bra{\pm}\,S^{\dagger}\) & a drawer-effect is the box
\(S^{\dagger}\) followed by a door-effect \\
50 & \(R_{\theta}\) real, \(R_{\theta}^{T}=R_{-\theta}\);
\((I\otimes R_{\theta})\text{cup}=(R_{-\theta}\otimes I)\text{cup}\) &
the box \(R_{\theta}\); slid round the cup it arrives turned the other
way \\
51 &
\(P=\bigl(\begin{smallmatrix}H&0\\0&HS^{\dagger}\end{smallmatrix}\bigr)\);
\((\bra q\otimes I)\,P\,(\ket q\otimes I)\) is its \(q\)-th block & a
two-wire box whose upper wire, pinned between state \(q\) and effect
\(q\), leaves the chosen box on the lower wire \\
52 &
\(\sum_{\text{win}}\big|\bra{ab}(I\otimes R_{\pi/8})\ket{\Phi^+}\big|^2=\cos^2(\pi/8)\)
& the game as a closed picture: the \(\tfrac1{\sqrt2}\) cup,
\(R_{\pi/8}\), two effects \\
\end{longtable}

\endgroup

Every row is an exact equality of matrices; the box \(P\) of row 51
arrives in §7.7.

\begin{center}\rule{0.5\linewidth}{0.5pt}\end{center}

\hypertarget{the-theorem-twice-6}{%
\section{§7.6 The Theorem, Twice}\label{the-theorem-twice-6}}

Now the trio's guarantee is turned against the realist. If the answers
are sealed notes --- six of them, a door-answer and a drawer-answer for
each of the three boxes --- then the four cases of the guarantee are
four equations those six numbers must satisfy. We shall see they cannot.

\hypertarget{the-example-that-fixes-the-ceiling-1}{%
\subsection{The example that fixes the
ceiling}\label{the-example-that-fixes-the-ceiling-1}}

Write each box's two sealed answers as \(\pm1\): for box \(A\), a
door-answer \(x_A\) and a drawer-answer \(y_A\), and likewise
\(x_B,y_B,x_C,y_C\). A bright note is \(+1\), a dark note \(-1\), so
``the number of dark notes is even'' means ``their product is \(+1\)'',
and ``odd'' means ``their product is \(-1\)''. The four cases of the
guarantee become

\[x_A x_B x_C=+1,\qquad x_A y_B y_C=-1,\qquad y_A x_B y_C=-1,\qquad y_A y_B x_C=-1.\]

Can six numbers satisfy all four? Try to satisfy the last three first:
set \(x_A=x_B=x_C=+1\) and \(y_A=+1\), and read off from
\(x_A y_B y_C=-1\) that \(y_By_C=-1\); take \(y_B=-1,y_C=+1\). Check
\(y_Ax_By_C=(+1)(+1)(+1)=+1\) --- but this was required to be \(-1\). It
fails. Adjust \(y_A=-1\): now \(y_Ax_By_C=(-1)(+1)(+1)=-1\), good, and
\(y_Ay_Bx_C=(-1)(-1)(+1)=+1\) --- required \(-1\). It fails again. Every
repair breaks another case. Three of the four we can always meet; the
fourth refuses. This is the same stubbornness as in §7.2, and it has the
same one-line reason.

\hypertarget{the-theorem}{%
\subsection{The theorem}\label{the-theorem}}

The eigen-relations behind the guarantee are exact. Recall the bridge: a
bright answer is \(+1\), a dark answer \(-1\); the door question is the
observable \(X\), the drawer question the observable \(Y\). Then
``\(X\otimes X\otimes X\) acting on \(N\) gives \(+N\)'' \emph{is} the
statement ``all doors: darks even'', and ``\(X\otimes Y\otimes Y\) gives
\(-N\)'' \emph{is} ``one door, two drawers: darks odd''.

\begin{quote}
\textbf{Theorem 7.5 (the GHZ paradox).} Let \(N=\ket{000}+\ket{111}\).
\textbf{(i)} The node satisfies, exactly,
\[\begin{aligned}X\otimes X\otimes X\,N&=+N,\\ X\otimes Y\otimes Y\,N=Y\otimes X\otimes Y\,N=Y\otimes Y\otimes X\,N&=-N.\end{aligned}\]
\textbf{(ii)} There are no six numbers
\(x_A,y_A,x_B,y_B,x_C,y_C\in\{\pm1\}\) with
\[x_Ax_Bx_C=+1,\quad x_Ay_By_C=y_Ax_By_C=y_Ay_Bx_C=-1.\] Part (i) says
the trio meets all four guarantees at once; part (ii) says no collection
of sealed notes can. So the trio is not a collection of sealed notes.
\end{quote}

Part (ii) is a piece of arithmetic and belongs to neither strand:
multiply the four left-hand sides. Every letter appears in exactly two
of the four products --- \(x_A\) in the first and second, \(y_A\) in the
third and fourth, and so on --- so the total is a product of squares,
\((+1)\). Multiply the four right-hand sides: \((+1)(-1)(-1)(-1)=-1\).
If all four held, \(+1=-1\). So at most three can hold.
(\texttt{SixNotes.py} runs all \(64\) assignments: none satisfies four,
the maximum is three, and flipping one note flips exactly two of the
four verdicts.)

Part (i) is what we prove twice.

\hypertarget{diracs-blackboard-6}{%
\subsection{Dirac's Blackboard}\label{diracs-blackboard-6}}

The Pauli boxes act on basis states by \(X\ket0=\ket1,\ X\ket1=\ket0\)
and \(Y\ket0=i\ket1,\ Y\ket1=-i\ket0\). Apply each operator to \(N\)'s
two terms in turn.

For \(X\otimes X\otimes X\):
\[X\!\otimes\!X\!\otimes\!X\,\ket{000}=\ket{111},\qquad X\!\otimes\!X\!\otimes\!X\,\ket{111}=\ket{000},\]
so \(X\otimes X\otimes X\,N=\ket{111}+\ket{000}=+N\).

For \(X\otimes Y\otimes Y\):
\[X\!\otimes\!Y\!\otimes\!Y\,\ket{000}=\ket1\otimes i\ket1\otimes i\ket1=i^2\ket{111}=-\ket{111},\]
\[X\!\otimes\!Y\!\otimes\!Y\,\ket{111}=\ket0\otimes(-i\ket0)\otimes(-i\ket0)=(-i)^2\ket{000}=-\ket{000},\]
so \(X\otimes Y\otimes Y\,N=-(\ket{111}+\ket{000})=-N\). The two
remaining relations are identical with the \(X\) in a different slot;
the two factors of \(\pm i\) multiply to \(-1\) each time, wherever the
lone \(X\) sits. Each minus sign is a single line of algebra:
\(Y\ket0=i\ket1\), \(Y\ket1=-i\ket0\), and \(i\cdot i=(-i)(-i)=-1\).
\(\;\blacksquare\)

\hypertarget{coeckes-sketchbook-1}{%
\subsection{Coecke's Sketchbook}\label{coeckes-sketchbook-1}}

We prove the same four relations as one drawing. The door-effect for
answer bit \(a\) is \(\bra aH\) --- so \(\bra0H=\bra+\) and
\(\bra1H=\bra-\) --- and we mark it simply \(a\) at a leg's end. Start
from the all-door case, the three door-effects \(a,b,c\) on the node's
legs; its value is the base case

\[\bra{abc}\,(H\otimes H\otimes H)\,N=\tfrac1{\sqrt2}\quad\text{when }a+b+c\text{ is even},\qquad 0\quad\text{when odd},\]

a closed picture we compute once (below). Now build the mixed case, door
on \(A\) and drawer on \(B\) and \(C\), by the rule of §7.4: each
drawer-effect is an \(S^{\dagger}\) box followed by a door-effect. So
the mixed case is the all-door picture with an \(S^{\dagger}\) on leg
\(B\) and an \(S^{\dagger}\) on leg \(C\):

\[\begin{tikzpicture}[sd]
  \node[sdnode] at (0,1.5) {};
  \draw[wire] (0,1.5) to[out=90,in=180] (1.2,3) -- (5.2,3);
  \draw[wire] (0,1.5) -- (5.2,1.5);
  \draw[wire] (0,1.5) to[out=-90,in=180] (1.2,0) -- (5.2,0);
  \node[sdeffect] at (5.2,3) {$a$};
  \node[sdbox] at (3,1.5) {$S^{\dagger}$};
  \node[sdeffect] at (5.2,1.5) {$b$};
  \node[sdbox] at (3,0) {$S^{\dagger}$};
  \node[sdeffect] at (5.2,0) {$c$};
\end{tikzpicture}\]

\textbf{Step 1 --- sliding and merging (Theorem 7.3).} Both
\(S^{\dagger}\) boxes are diagonal, so they slide through the node onto
one leg and merge; since \(S^{\dagger}S^{\dagger}=Z\), the two of them
become a single \(Z\) box, which we may keep on leg \(A\):

\[\begin{tikzpicture}[sd]
  \node[sdnode] at (0,1.5) {};
  \draw[wire] (0,1.5) to[out=90,in=180] (1.2,3) -- (5.2,3);
  \draw[wire] (0,1.5) -- (5.2,1.5);
  \draw[wire] (0,1.5) to[out=-90,in=180] (1.2,0) -- (5.2,0);
  \node[sdbox] at (3,3) {$Z$};
  \node[sdeffect] at (5.2,3) {$a$};
  \node[sdeffect] at (5.2,1.5) {$b$};
  \node[sdeffect] at (5.2,0) {$c$};
\end{tikzpicture}\]

\textbf{Step 2 --- the \(Z\) flips one door-answer.} A \(Z\) before a
door-effect turns it into the other: since \(HZ=XH\) and
\(\bra aX=\bra{a\oplus1}\), we have
\(\bra aHZ=\bra aXH=\bra{a\oplus1}H\) --- the door-answer bit \(a\) is
flipped. So the mixed picture equals the all-door picture with answer
\(a\) replaced by its opposite. Its value is therefore
\(\tfrac1{\sqrt2}\) exactly when \(a+b+c\) is \emph{odd} --- the
parities have exchanged. That is precisely ``darks odd'', the relation
\(-N\).

All three minus-sign relations are this one move: whichever two legs
carry drawers, their two \(S^{\dagger}\) boxes merge into one \(Z\), and
the \(Z\) flips one door-answer, reversing the parity. The picture says
not only \emph{that} the sign appears but \emph{where} it lives --- on
the leg the merged box sits on. \(\;\blacksquare\) (\texttt{SixNotes.py}
runs the rewrite step by step as exact picture equalities and confirms
the base-case values.)

\emph{The base case, computed.}
\(H^{\otimes3}N=\ket{+{+}{+}}+\ket{-{-}{-}}\), and
\(\braket{abc|+{+}{+}}=\tfrac1{2\sqrt2}\) while
\(\braket{abc|-{-}{-}}=\tfrac1{2\sqrt2}(-1)^{a+b+c}\), so the sum is
\(\tfrac1{2\sqrt2}\big(1+(-1)^{a+b+c}\big)\), which is
\(\tfrac1{\sqrt2}\) for even \(a+b+c\) and \(0\) for odd.

\hypertarget{the-verdict-6}{%
\subsection{The verdict}\label{the-verdict-6}}

The algebra wins on length, and narrowly. Each minus sign is one line on
the blackboard --- \(Y\ket0=i\ket1\), \(Y\ket1=-i\ket0\), and
\(i\cdot i=-1\) --- and there is nothing to set up first. The sketchbook
must first own a new rule, sliding-and-merging through the node, whose
own justification is a two-line index check (Theorem 7.3); once that is
paid for, all three minus signs are a \emph{single} move --- two
\(S^{\dagger}\) merge to a \(Z\), and the \(Z\) flips one answer --- and
the picture shows \emph{where} the sign lives, which the blackboard does
not. There is a second debt the picture cannot dodge: its base case, the
all-door value \(\tfrac1{\sqrt2}\)-on-even, is itself a closed
calculation, and that calculation is algebra. So the picture explains
and the blackboard proves, and the blackboard is shorter.

And a warning against overselling either. Part (ii) --- that no six
numbers work --- is arithmetic, and belongs to neither strand: it is the
multiplication of four \(\pm1\) products. Part (i) is what the two
strands prove, and all it establishes is an \emph{equality of
processes}: that the node meets all four guarantees at once. That the
sealed notes cannot, and hence that the trio is not sealed notes, is the
arithmetic. Neither the picture nor the algebra proves that nature is
nonlocal; they prepare the contradiction, and the counting closes it. No
experiment is a rewrite, and no rewrite is an experiment.

\begin{center}\rule{0.5\linewidth}{0.5pt}\end{center}

\hypertarget{the-round}{%
\section{§7.7 The Round}\label{the-round}}

We can now play the whole thing as one game, apart, and draw it as a
single picture.

Three players take one box each and go to three rooms; each also carries
a sealed envelope naming which way to open --- door or drawer. A referee
--- the reader --- holds a fourth sealed slip that says what the answers
must total: \emph{even} if all three envelopes say door, \emph{odd} if
one says door and two say drawer. Each player opens their box the named
way, hears bright or dark, and returns; the answers are compared with
the slip. The round is won when the answers match the slip.

One master diagram lays this out as a single picture. It needs one new
box.

\textbf{The player's box \(P\).} Each player's question and their leg of
the node meet in a two-wire box

\[P=\begin{pmatrix}H&0\\0&HS^{\dagger}\end{pmatrix},\]

built in Chapter 4's block form, exactly as
\(\mathrm{CNOT}=\bigl(\begin{smallmatrix}I&0\\0&X\end{smallmatrix}\bigr)\)
is. The upper wire carries the question, \(\ket0\) for door or \(\ket1\)
for drawer; the lower wire carries the player's leg. When the upper wire
is pinned between the state \(\ket q\) and the effect \(\bra q\) of the
\emph{same} value, the box collapses \textbf{exactly} to its \(q\)-th
block --- \(H\) for the door, \(HS^{\dagger}\) for the drawer --- an
equality, not a postselection with a probability attached. Following the
chosen block by a door-effect \(\bra a\) gives \(\bra aH=\bra{\pm}\)
(the door question) or \(\bra aHS^{\dagger}=\bra{\pm i}\) (the drawer
question, by §7.4). One box, both questions, chosen by the wire.

\hypertarget{the-diagram}{%
\subsection{The diagram}\label{the-diagram}}

Here is \emph{The Round}. Rows, top to bottom: the question \(q_A\), leg
\(A\), then \(q_B\), leg \(B\), then \(q_C\), leg \(C\), and at the foot
the referee's tally wire.

\begin{center}
\begin{tikzpicture}[sd]
  % the three legs (lines 1-3), from the node to the open right-hand ends
  \draw[wire] (0,6.6) to[out=90,in=180] (2.2,11) -- (24.6,11);
  \draw[wire] (0,6.6) -- (24.6,6.6);
  \draw[wire] (0,6.6) to[out=-90,in=180] (2.2,2.2) -- (24.6,2.2);
  % the three question wires (lines 4-6): state, through the player's box, effect
  \draw[wire] (2.6,13.2) -- (7.6,13.2);
  \draw[wire] (2.6,8.8) -- (7.6,8.8);
  \draw[wire] (2.6,4.4) -- (7.6,4.4);
  % the tally wire (line 7): one unbroken line, two crossings
  \draw[wire] (6.4,0) -- (11.4,0) to[out=0,in=180] (13.6,4.4) -- (16.6,4.4)
               to[out=0,in=180] (18.8,8.8) -- (22.9,8.8);
  \node[sdnode] at (0,6.6) {};
  \node[sdstate] at (2.6,13.2) {$q_A$};
  \node[sdstate] at (2.6,8.8) {$q_B$};
  \node[sdstate] at (2.6,4.4) {$q_C$};
  \node[sdbox, minimum height=4.1em] at (5.1,12.1) {$P$};
  \node[sdbox, minimum height=4.1em] at (5.1,7.7) {$P$};
  \node[sdbox, minimum height=4.1em] at (5.1,3.3) {$P$};
  \node[sdeffect] at (7.6,13.2) {$q_A$};
  \node[sdeffect] at (7.6,8.8) {$q_B$};
  \node[sdeffect] at (7.6,4.4) {$q_C$};
  \node[sdstate] at (6.4,0) {$r$};
  \node[sdbox, minimum height=4.1em] at (10,1.1) {$\scriptstyle\mathrm{CNOT}$};
  \node[sdbox, minimum height=4.1em] at (15.1,5.5) {$\scriptstyle\mathrm{CNOT}$};
  \node[sdbox, minimum height=4.1em] at (20.2,9.9) {$\scriptstyle\mathrm{CNOT}$};
  \node[sdeffect] at (22.9,8.8) {$1$};
\end{tikzpicture}
\end{center}

The three legs of the node run out to the right, one to each player;
each passes through its player's box \(P\) (lower wire of the box) and
then acts as the control, on top, of one \(\mathrm{CNOT}\), before
ending at an open right-hand end --- the player's answer. Each question
wire begins in the state \(\ket{q}\), runs through the upper wire of its
\(P\), and ends in the effect \(\bra q\) of the same label: this is the
pinning that collapses \(P\) to its block. The tally wire begins in the
referee's state \(\ket r\), is flipped by each leg's \(\mathrm{CNOT}\)
in turn, and ends in the effect \(\bra1\) --- the \emph{losing} slip.
The two places where the tally wire steps up past a leg are crossings,
and a crossing is Chapter 4's swap drawn with one strand held straight;
bending a wire is free (Chapter 6). A crossing neither joins nor breaks
a line, so the tally is one unbroken line from the state \(\ket r\) to
the effect \(\bra1\).

\begin{quote}
\textbf{Theorem 7.6 (the round never loses).} For each of the four dealt
question triples \((q_A,q_B,q_C)\in\{000,011,101,110\}\), with
\(r=q_A\lor q_B\lor q_C\), \emph{The Round} with the losing effect
\(\bra1\) on the tally wire evaluates to the \textbf{zero vector} on the
three open answer wires --- exactly. The amplitude of every losing
outcome is zero. With the winning effect \(\bra0\) instead, the picture
equals the same diagram with the tally wire and the three
\(\mathrm{CNOT}\)s erased: four answers of the correct parity, each of
amplitude \(\tfrac1{\sqrt2}\).
\end{quote}

\emph{Proof idea.} Pinning each \(P\) turns the diagram into the node
with the three chosen questions on its legs and a running parity check
on the tally. For a dealt triple the guarantee of Theorem 7.5(i) says
the parity of the answers is fixed; the losing slip demands the opposite
parity, which never occurs, so its amplitude is zero everywhere.

\emph{Proof (on the blackboard).} Pin the three question wires: each
\(P\) becomes its block, so the three legs of \(N\) carry \(H\), or
\(HS^{\dagger}\), according to the deal, and then their door-effects ---
that is, the node measured in the door or drawer question per the deal.
By Theorem 7.5(i) the surviving answers all have the parity the
guarantee names, and the three \(\mathrm{CNOT}\)s write that parity onto
the tally wire, which started at \(\ket r\) with \(r\) the guaranteed
value. The tally therefore reaches \(\ket0\) on every surviving branch,
and the losing effect \(\bra1\) annihilates all of them: the result is
the zero vector. Replace \(\bra1\) by \(\bra0\) and every surviving
branch passes; erasing the (now trivial) parity machinery leaves four
equal-amplitude answers, squared norm \(2=\braket{N|N}\).
\(\;\blacksquare\) See \texttt{Code/Ch7/Round.py}, which builds the full
diagram, confirms the zero vector for the four deals with the losing
slip and a non-zero result with the winning slip, and --- the point of
playing apart --- that each player's own answer is a fair coin for every
one of the eight triples, so nothing about the other players' questions
travels to any room.

\textbf{What it evaluates to, and its cost.} The winning picture's open
end has squared norm \(2\), the node's \(\braket{N|N}\); the honest
probabilities are these divided by \(2\), so each of the four correct
answers occurs with probability \(\tfrac14\). The contrast with §7.3 is
the whole point of the chapter. There, the wedge won \emph{more often}
--- with probability \(\cos^2(\pi/8)\), past the sealed-note ceiling but
short of certainty. Here the players win \emph{every time}, and no
collection of sealed notes can, by Theorem 7.5(ii). A game won with
certainty by a shared state and never by any local strategy is what
Brunner's review calls pseudo-telepathy (Brassard and colleagues, 2005);
we shall not need the name again.

\begin{quote}
\textbf{Feynman's Notebook.}

\emph{Three particles, and a campus.} The three-party argument of this
chapter is due to Greenberger, Horne and Zeilinger (1989), with Shimony
added in a fuller 1990 account and a much-read exposition by Mermin the
same year; we take the attribution, as we take the game, through Brunner
and colleagues' review. Their original construction used \emph{four}
particles, the three-particle version noted as a possibility; the crisp
three-party form in \(X\) and \(Y\) is the one the review presents, and
the one drawn above.

Closing every loophole at once took until 2015, when Hensen and
colleagues, on the campus of Delft University of Technology, entangled
two electron spins in diamond \(1.3\) km apart and ran \(245\) trials of
the CHSH game. They found \(S=2.42\pm0.20\) --- above the sealed-note
ceiling of \(2\) --- winning \(k=196\) trials, with probability
\(0.039\) that a local-realist model with memory could fake it
(\(0.019\) under their more conventional analysis). The bent wire had
broken the inequality in the open.
\end{quote}

\begin{center}\rule{0.5\linewidth}{0.5pt}\end{center}

\hypertarget{where-was-it}{%
\section{§7.8 Where Was It?}\label{where-was-it}}

Let us gather what we have. Sealed notes, questioned in the several
combinations of a game, meet a ceiling: they win at most three of the
four pairs, and \(S\le2\) (§7.2). A shared bent wire, tilted by a wedge,
breaks the ceiling --- \(S=2\sqrt2\), a win rate of \(\cos^2(\pi/8)\)
(§7.3). A single new drawing, a dot where three wires meet, carries a
guarantee that any one answer can be foretold from the other two, which
is the realist's whole case (§7.4). And that same guarantee, written out
as six sealed numbers, contradicts itself --- no six numbers meet all
four cases, though the node meets them exactly, which we proved twice
(§7.6). Played apart, the trio wins with certainty what no notes could
(§7.7). In every picture the agreement was kept in the same place it was
kept in Chapter 6: not in either box, but in the fact that the boxes are
legs of one thing.

So the answer to the chapter's opening question is no. The notes were
not written. Two halves that always agree, questioned only one way,
could have carried sealed notes; three boxes questioned in four
combinations cannot, because the four combinations demand of the notes
something arithmetic forbids. Where was the answer, then, if not in the
boxes and not crossing the house? In neither box: in the place where the
three legs meet.

And that dot has more to tell. We have used it as one fixed thing,
unlabelled, always \(\ket{000}+\ket{111}\). But a dot where wires meet
is the beginning of a language, not the end of one: allow it a family
--- dots that count in a second way, dots that carry a turn of the
sphere --- and the pictures we have drawn box by box will start to
rewrite themselves, without ever passing back through the algebra. The
next chapter gives the dot its family, and its name.

\hypertarget{octet-for-spiders}{%
\chapter{Octet for Spiders}\label{octet-for-spiders}}

\begin{center}\rule{0.5\linewidth}{0.5pt}\end{center}

\emph{短日や}\\
\emph{糸八本の}\\
\emph{一結び}

\emph{tanjitsu ya}\\
\emph{ito happon no}\\
\emph{hitomusubi}

\emph{The short day ---}\\
\emph{eight threads,}\\
\emph{one knot.}

\begin{center}\rule{0.5\linewidth}{0.5pt}\end{center}

\emph{{[}Achilles comes in out of the winter dark, stamping the snow off
his boots; the Crab's scullery is warm and low-lit, a fire in the grate.
A stout beam crosses the room overhead, and from hooks along it hang
lengths of kitchen cord, some strung through small curtain-rings. The
short day is already going grey at the one window. The Tortoise is there
before him, in a chair drawn near the fire.{]}}

\emph{{[}Along the long wall, running the whole length of the room under
the beam, hangs Escher's \textup{Metamorphosis II} --- a narrow strip of
a woodcut, no taller than a hand and about twice a man in length, so
that the eye has to walk it end to end. The lamp on the table throws its
light along the strip: part way along it a little town stands by the
water with a square tower, and elsewhere the print goes to a chequer of
light and dark squares. Above it, from the beam, the strung cords hang
still in the warm air. At the one window the afternoon has gone the flat
grey of coming snow, and the first flakes have begun --- few and fitful,
in no hurry, settling and starting again.{]}}

\textbf{Achilles:} I had forgotten a person could be this cold and still
be alive to complain of it. The canal's frozen hard as the road --- I
could not tell you where the one stopped and the other began.

\textbf{Crab:} Come to the fire, come to the fire. There's a pot on the
hob and it's been waiting on you. \emph{{[}He pours.{]}} I keep this
room warm for the work; the cords go slack if it's damp and sulk if it's
cold.

\textbf{Achilles:} The work being all this washing-line you've hung the
ceiling with. \emph{{[}He takes the cup gratefully.{]}} I walked past
your window twice thinking you'd gone to bed and left the string out.

\textbf{Tortoise:} He's been at it since noon. I came early and have
watched him tie knots and untie them and stand back and admire them like
a man hanging pictures.

\textbf{Crab:} And now I have my audience, I shall show the pair of you
the finest thing to come out of a workshop in Delft in twenty years, and
neither of you will leave until you have understood it. Warm your hands
first. It keeps till you're thawed.

\emph{{[}Achilles crouches to the grate a moment, turning his hands to
it until the ache goes out of them; the cup steams under his chin. The
scullery closes warm around the three of them, the snow settling against
the grey pane, and the cords overhead barely stirring in the heat off
the fire.{]}}

\begin{center}\rule{0.5\linewidth}{0.5pt}\end{center}

\hypertarget{first-challenge-a-name-a-colour-a-turn}{%
\section{\texorpdfstring{First Challenge --- \emph{A Name, a Colour, a
Turn}}{First Challenge --- A Name, a Colour, a Turn}}\label{first-challenge-a-name-a-colour-a-turn}}

\textbf{Achilles:} \emph{{[}looking up at the beam, where one ring hangs
with three cords running out of it{]}} Wait --- I know that. Or I knew
its plainer cousin. Last autumn you showed us a bare dot where three
cords met, that already knew things about itself: any one of its cords
foretold by the other two. But that one had no colour, and no more than
the three legs. This has a colour, and --- how many legs is it wearing?

\textbf{Crab:} As many as I care to give it. \emph{{[}He runs a fourth
cord through the ring, then a fifth.{]}} That is the whole point of the
family. The dot has grown up.

\textbf{Achilles:} It has a family now. Has it a name?

\textbf{Crab:} It has. This --- a ring with cords drawn through it and
knotted into one --- I call a \textbf{green ring}, for the plainest of
reasons: the ring is green. It asks its cords a single question, the
same question your autumn dot asked: \emph{do you all say the same?} All
bright together, or all dark together --- it passes those; anything else
it stops. \emph{{[}He touches a small bead threaded on the ring
itself.{]}} And this bead is its \textbf{turn} --- it slides round the
ring to mark how far the thing is turned, the very sphere-turn you
learned in the spring. Bead at the top, no turn at all.

\textbf{Achilles:} Then let me guess what the simplest of them does. A
green ring with one cord coming in and one going out, and its bead at
the top --- no turn. It asks the one cord whether it agrees with\ldots{}
itself, which it can hardly fail to do. So it does nothing. It's a bare
cord with a bead on a ring, and might as well not be there.

\textbf{Crab:} \emph{{[}pleased{]}} Right, and right the first time ---
which I gather is not always your way. A green ring, one in, one out, no
turn, is nothing but wire.

\textbf{Tortoise:} So the interest is in the ones with more legs, and in
what happens when two of them meet. Watch. \emph{{[}She takes a loose
cord and threads two green rings onto it, a little apart, each with its
bead set some way round.{]}} Here are two, on one cord, each turned a
little.

\emph{{[}The two green rings sit on the loose cord, a hand's breadth
between them, one bead a third of the way round, the other a quarter.
She sets a finger to the near ring and slides it along the cord until it
meets the far one; the cord does not move and its far end does not stir,
only the two rings travel together and knot into a single green ring
where they meet. The two beads, sliding round with their rings, come to
rest as one bead, at a turn that is the two turns added -- a third and a
quarter together. Two rings have become one; the far end of the cord is
exactly where it was.{]}}

\textbf{Tortoise:} Two green rings sharing a cord are one green ring,
and the turns add. Slide them together; that is all.

\textbf{Crab:} And it makes no difference to anything downstream.
\emph{{[}He threads a test cord through the fused ring and gives its far
end a tug; the end answers exactly as it did when there were two
rings.{]}} The far end can't tell whether it was one ring or two.
\emph{{[}He lifts the demo cord and its ring off and sets them
aside.{]}}

\emph{{[}In the grate two flames stand up side by side off the one log,
wavering apart and together, until a draught leans them into each other
and they run up as a single taller flame, its two heats gathered in the
one tongue. The chimney draws it exactly as it drew the two; nothing
above the pot warming on the hob is any the wiser.{]}}

\textbf{Achilles:} So a whole chain of green rings is really just one
green ring wearing the sum of all their turns. Tidy. \emph{{[}He looks
up at the beam.{]}} And you want to build something. Where do we start?

\textbf{Crab:} Here. \emph{{[}He fixes two green rings to the beam, one
above the other on the left, and runs a cord into each from the left
side.{]}} Two green rings, and a cord fed to each --- one for you to
hold responsible, one for the Tortoise. The beginning of the thing.

\emph{{[}Achilles steadies the cord into the upper green ring; the
Tortoise sets the cord into the lower one. Two cords now run in from the
left, one to each green ring.{]}}

\textbf{Achilles:} \emph{{[}stepping back, taking it in{]}} Two green
rings, two cords in, and all of it green. Is green all there is?

\begin{center}\rule{0.5\linewidth}{0.5pt}\end{center}

\hypertarget{second-challenge-the-other-colour}{%
\section{\texorpdfstring{Second Challenge --- \emph{The Other
Colour}}{Second Challenge --- The Other Colour}}\label{second-challenge-the-other-colour}}

\textbf{Crab:} \emph{{[}with relish{]}} It is not. \emph{{[}He lifts a
heavy flat gadget onto the table --- a griddle-like press with a slotted
rim and, standing beside it, a second row of empty hooks.{]}} This is my
\textbf{koekpan}. It will do the heavy work later, and I'll thank you
not to stare at it while it does --- it's a shy thing. Two parts to know
now. Along its rim, this \textbf{flip-slot}: a cord passed through it
comes out the other side changed. And this second set of hooks ---
whatever the pan makes, it hangs \emph{there}, and leaves the first
cords be.

\emph{{[}It is a homely thing to have such a fuss made of it: a broad
flat press of old Delft iron, its face worn to a soft shine by years of
hands, the handle darkened where a claw has gripped it a thousand times.
Set down on the table it gives off a warmth of its own now, and --- if
nobody is turned its way --- a low sound at the very edge of hearing,
like a kettle a long way from the boil. Achilles leans in to catch it;
the sound is not there; he sits back and it steals up again.{]}}

\textbf{Achilles:} Changed into what?

\textbf{Crab:} Watch. \emph{{[}He takes a loose cord feeding a green
ring and passes the cord through the flip-slot.{]}} A green ring counts
one way --- all-agree, bright-together or dark-together. Feed its cord
through the slot, and out comes a ring that counts the \emph{other} way.
A \textbf{red ring}.

\emph{{[}The green ring sits on its loose cord. He draws the cord
through the flip-slot; the cord runs through and the ring at its end, as
it comes clear of the slot, is red where it was green. Its legs and its
bead have not moved -- same cords, same turn -- only the colour, and
with the colour the question it asks.{]}}

\textbf{Achilles:} One slot, and it changes sides. But that's only the
one cord flipped. What of a ring with a whole fistful of legs?

\textbf{Crab:} Then you flip them all. \emph{{[}He takes a green ring
with four cords and threads every one of its cords through a flip-slot
in turn.{]}} Every leg through a slot, and the ring itself turns.

\emph{{[}The green ring hangs with its four cords. One by one he draws
each cord through the flip-slot; as the last of the four comes through,
the ring itself is red -- not touched, not re-tied, only every one of
its legs sent through the slot, and the colour has followed. The ring's
four cords still run where they ran; only green has become red.{]}}

\textbf{Crab:} \emph{{[}handing the flipped-round ring to Achilles{]}}
There. Handle it. It answers to every question exactly as a red ring
does --- it \emph{is} one. \emph{{[}Achilles tries a cord or two against
it, nods; the Crab takes the demo ring and its cords down and sets them
by the first.{]}}

\textbf{Achilles:} So it isn't that I \emph{called} it red. Dressing a
green ring's every leg in flip-slots truly makes it count the other way.
The colour earned it.

\textbf{Tortoise:} A red ring and a green ring see different things.
Green asks whether the cords all say the same the plain way; red asks it
the sideways way --- the way the spring's other pair of alternatives
said it. Two honest counters, one cord, two answers.

\emph{{[}The snow has thickened at the window while they talked. Where
the flakes fall against the bare black boughs of the tree by the bridge
they show white, each one plain; where they cross the pale grey of the
sky behind, the very same flakes go dark, small specks against the
light. Achilles's eye drifts from the black branches up to the sky, and
the whole falling field turns over at once --- white to dark, the one
snow answering to whichever ground he sets it against.{]}}

\textbf{Crab:} \emph{{[}fixing two red rings to the beam on the right,
and running a cord out of each to the right{]}} And here are two reds,
on the right, where the answers will come out --- a cord drawn from
each. \emph{{[}Two cords now hang out to the right, one from each red
ring.{]}} Greens on the left, reds on the right. That is my frame.

\textbf{Achilles:} Greens counting one way on the left, reds the other
way on the right. \emph{{[}He looks from side to side.{]}} A red ring
and a green ring see different things --- you've just made me agree to
it. So what happens where they meet?

\begin{center}\rule{0.5\linewidth}{0.5pt}\end{center}

\hypertarget{third-challenge-each-copies-the-others-points}{%
\section{\texorpdfstring{Third Challenge --- \emph{Each Copies the
Other's
Points}}{Third Challenge --- Each Copies the Other's Points}}\label{third-challenge-each-copies-the-others-points}}

\textbf{Tortoise:} That is exactly the question. Let us make them meet
on the smallest scale first: one point of one colour, fed into a single
ring of the other. \emph{{[}She takes a short loose cord ending in a red
ring with a single leg.{]}} This is a red point --- a red ring with just
the one cord, the plainest red thing there is. And here is a green ring
made to fork: one cord in, two cords out, its bead at the top, no turn.

\textbf{Achilles:} Then I can tell you what happens without your pan.
The green ring's business is to ask \emph{do my cords agree?} --- but it
asks in \emph{green}, and a red point is not a green kind of point; it's
a stranger to that question. So the ring will have nothing to say to it.
It passes straight through and comes out on one leg, unchanged. One cord
in, one meaningful cord out.

\textbf{Tortoise:} Say it to the pan, then. \emph{{[}She feeds the red
point into the green fork, and presses it on the koekpan.{]}}

\emph{{[}The red point-cord runs into the green fork-ring: one cord in
on the left, two green legs out on the right. She sets the little tangle
on the koekpan and presses. The Crab turns his face to the fire while it
works.{]}}

\emph{{[}Off the pan, on the second set of hooks, hangs not what
Achilles foretold. There is no green ring left at all -- and instead of
one cord out, there are two, and each is a red point, the very twin of
the one that went in. One red point in; two red points out; the green
ring gone.{]}}

\textbf{Achilles:} \emph{{[}staring{]}} Two. There are two of them, and
both red, and the green ring's --- gone. But I sent in \emph{one}.
\ldots Oh. It didn't pass the point along. It \emph{copied} it. The
green ring can't read a red point, so it does the only thing it knows
--- it hands out a matching one on every leg. Each colour copies the
other's points.

\emph{{[}Achilles's cup stands where he set it on the table, between the
lamp at his elbow and the fire at his back. It throws two shadows across
the cloth, one for each light, and they are the very same cup twice over
--- one bowl, one handle, the one thing laid down in two identical
prints, and no reaching for it makes a third.{]}}

\textbf{Crab:} \emph{{[}without looking round{]}} It does. Mind you
don't lean over it; the last man who watched it work jammed it for a
week. \emph{{[}He presses once more, head still turned; the two red
points come off clean again, and he lifts the copies down.{]}}

\textbf{Tortoise:} Now a companion trick, and this one wants no pan ---
just my two hands. \emph{{[}She sets a green ring on a taut cord, its
bead turned well off the top, and threads a red half-turn onto the cord
before it --- a red ring with its bead set square at the halfway
mark.{]}} A green ring with a turn on it, and coming up the cord towards
it, a red half-turn. Watch what each does to the other as I push the red
one through.

\emph{{[}The green ring sits with its bead a third of the way round; the
red half-turn waits on the cord before it. She pushes the red half-turn
along the cord and through the green ring to the far side. What passes
through is unchanged -- still a red half-turn, its bead square at the
mark. But the green ring it passed through has its bead swung clean over
to the mirror side, the same turn the other way about. The traveller is
the same; the thing it travelled through is reversed.{]}}

\textbf{Achilles:} The red one came through untouched --- still a
half-turn. But the green ring's turn has flipped to the far side, the
mirror of what it was. A half-turn of one colour, pushed through the
other, turns the other's turn around.

\emph{{[}Beyond the glass the fall has come on in earnest, whirling and
crossing against the last of the grey daylight; no single flake can be
followed for more than a moment before it loses itself in the rest of
them. The room has drawn in close and warm about the small work on the
table, and Achilles, without looking away from it, holds his free hand
out to the fire.{]}}

\textbf{Tortoise:} Just so. Now --- green copies red's points, and red
copies green's; a half-turn of one flips the turn of the other.
\emph{{[}She gathers the loose cords.{]}} Put two green rings and two
red together, and let all of that happen at once.

\begin{center}\rule{0.5\linewidth}{0.5pt}\end{center}

\hypertarget{fourth-challenge-the-cradle-comes-undone}{%
\section{\texorpdfstring{Fourth Challenge --- \emph{The Cradle Comes
Undone}}{Fourth Challenge --- The Cradle Comes Undone}}\label{fourth-challenge-the-cradle-comes-undone}}

\textbf{Achilles:} All at once. That sounds like a knot.

\textbf{Crab:} It is a beauty of a knot. \emph{{[}At the beam: the two
green rings on the left, the two red on the right, feed-cords in,
draw-cords out. Now he strings the middle.{]}} Every green tied to every
red --- the top green to the top red, and across to the bottom red; the
bottom green to the top red, and across to the bottom. A full cross
between the colours. I call it the \textbf{cradle}.

\emph{{[}He runs the four cords of the cross, one after another: top
green to top red; top green down across to bottom red; bottom green up
across to top red; bottom green to bottom red. The two feed-cords still
run in on the left, the two draw-cords still run out on the right;
between them the four crossed cords hang taut. The cradle stands strung
on the beam.{]}}

\textbf{Achilles:} Four cords in the cross, every green speaking to
every red. \emph{{[}He counts the crossings under his breath, tracing
them with a finger.{]}} It's a proper thicket. What does a thicket like
that \emph{do}?

\textbf{Crab:} Ask the pan. \emph{{[}He sets the cradle's shape onto the
koekpan --- a copy of it, the strung cradle staying put on the beam ---
and presses, his eyes on the fire.{]}}

\emph{{[}The four-cord cross is copied onto the pan; the pan cooks it
down. On the beam the cradle itself does not stir -- every cord where it
was. But on the second set of hooks comes off something far barer: a
single red ring gathering the two cords that came in, one cord to a
single green ring, and the two cords going out from the green. One
knot-pair, doing what the whole cross did. Four crossed cords have
become one; the cradle on the beam is untouched.{]}}

\textbf{Achilles:} \emph{{[}working it over aloud{]}} The whole cross
--- all four crossings of it --- comes off as one red ring, one cord,
one green ring. It reads whether the two cords coming in agree, and
writes that same answer onto both going out. \emph{{[}slowly{]}} And the
cradle does \emph{exactly} that, I'll wager, if I trace it
through\ldots{} yes. Fewer cords for the same job.

\textbf{Crab:} \emph{{[}lifting the pressed pair off the second
hooks{]}} Fewer cords, same job. That is the whole trade.

\emph{{[}Out at the window the snow has been at work. Where this morning
the frozen canal ran its one way and the road crossed it another --- a
tangle of edges and banks and a little bridge --- the fall has filled it
all level and laid one unbroken white sheet across the lot, so that the
crossing is a single smooth place now, and yet a body could still walk
it corner to corner and fetch up exactly where the old crossing would
have taken him.{]}}

\textbf{Achilles:} \emph{{[}frowning{]}} Though your pan is a bit of a
liar about the bookkeeping. When I trace the cross through by hand and
match it against that little pair, they do the same \emph{work} --- but
the pair comes out fainter than it ought, by some steady factor the
picture just\ldots{} loses count of. The shapes agree; a number has gone
missing.

\textbf{Tortoise:} The number is real and the picture drops it; that is
honest to say. The pictures will tell you \emph{that two things are
equal} and decline to keep the change. Now --- the cross's twin.
\emph{{[}She ties a green ring to a red ring with two separate cords
between them.{]}} A green and a red, joined not by one cord but by two.
Guess.

\textbf{Achilles:} Two cords must bind them twice as tightly.

\textbf{Tortoise:} Press it. \emph{{[}She sets the double-linked pair on
the koekpan; the Crab presses, looking away.{]}}

\emph{{[}The green ring and the red ring hang joined by two distinct
cords. The pan takes a copy and cooks it. Off the second hooks come the
two rings -- but with nothing between them at all: the two cords gone,
the green loose on its side, the red loose on its own. Two cords between
the colours have come to the same as none.{]}}

\textbf{Achilles:} They fell apart. Two cords between a green and a red
is the same as \emph{no} cord --- the same as never having tied them.
\emph{{[}He shakes his head.{]}} Twice tied is untied.

\emph{{[}The steam off the pot on the hob has been wavering all
afternoon --- leaning toward the thread of cold that creeps under the
door, then the other way toward the colder air sliding down off the
frozen pane. Now, for a moment, the two draughts come on together and
meet under the spout, and the steam that each alone bent aside stands up
dead straight between them, as still as if the room held no draught at
all.{]}}

\textbf{Tortoise:} \emph{{[}taking the two freed rings down{]}} The two
facts are one fact seen twice. The cross coming undone, and two cords
being as good as none --- the same law, that the two colours reorganise
each other. It is the engine of everything the pan does.

\textbf{Achilles:} \emph{{[}looking back up at the cradle on the
beam{]}} Fewer cords for the same job. Could that untangle a real gate?

\begin{center}\rule{0.5\linewidth}{0.5pt}\end{center}

\hypertarget{fifth-challenge-three-crossings-one-swap}{%
\section{\texorpdfstring{Fifth Challenge --- \emph{Three Crossings, One
Swap}}{Fifth Challenge --- Three Crossings, One Swap}}\label{fifth-challenge-three-crossings-one-swap}}

\textbf{Crab:} A real gate is where I was going. \emph{{[}He ties a
single green ring to a single red ring, sharing one cord --- green up on
the top line, red down on the bottom.{]}} One green, one red, one cord
between them: this is a gate you already know. If the top line runs a
mark, the bottom flips; if not, not. The autumn's gate, the commonest in
all of computing --- the one that copies a bit and carries it.

\textbf{Achilles:} I remember it. The controlled flip.

\textbf{Crab:} Now three of them in a row, and I'll alternate which line
holds the control. \emph{{[}He strings three such gates along the two
lines as a loose tangle --- green-red, then red-green, then green-red
--- the middle one reversed.{]}} Top, then bottom, then top. Three
crossings' worth of gate.

\textbf{Achilles:} \emph{{[}quick, tracing it in the air{]}} Let me run
a mark through before you touch that pan. Mark the top line, bottom
bare: the first gate flips the bottom, so both are marked; the middle
gate, its control on the bottom now, flips the top bare; the last does
nothing. So the mark came in on top and leaves on the bottom --- and the
other way about, it crosses back. The two lines have traded ends: a
swap, plain and exact, and I did it in my head.

\textbf{Crab:} You did. Now watch the pan do it the slow, deep way.
\emph{{[}He copies the three-gate tangle onto the koekpan and presses,
eyes shut.{]}}

\emph{{[}The three gates hang as a loose tangle across the two lines.
The pan takes its copy and cooks it -- gathering the first two gates
into a cross, undoing the cross to a pair, fusing along the line,
letting a doubled link fall away. Off the second hooks comes no tangle
at all: just the two lines, each running from its own side to the other
-- a bare crossing. The top line comes out at the bottom, the bottom at
the top. Three gates have cooked down to one plain cross.{]}}

\textbf{Achilles:} A bare crossing. The same swap I traced --- the two
lines simply changed places.

\emph{{[}The whirl at the window has settled by now to a steady, even
fall, straight down and unhurried, and through it two muffled figures
have come out onto the white crossing from opposite banks. They meet in
the middle, sidestep one another with the little bob strangers make, and
go on --- and each fetches up on the bank the other set out from, the
two of them swapped end for end without either breaking stride.{]}}

\textbf{Tortoise:} So which won? Be fair to both. Your head got there
faster, and got it \emph{exactly} --- no missing number, three gates
traced to a swap in a breath. The pan was slower, and dropped its factor
again.

\textbf{Achilles:} Then the pan lost.

\textbf{Tortoise:} On length, yes. But it showed you \emph{why}: the
swap is no coincidence of three particular gates --- it falls straight
out of the two colours undoing each other, the cradle law stood on its
side. Change the gates for others of that shape, and the same crossing
drops out. The quick way tells you \emph{that}; the slow way tells you
\emph{why}.

\textbf{Achilles:} \emph{{[}quietly, looking up at the cradle again{]}}
Three of them, and they just\ldots{} swap. \emph{{[}He lifts the tangle
and the crossing-copy off, and stands looking at the strung cradle on
the beam.{]}} How many cords went into that cradle, anyway?

\begin{center}\rule{0.5\linewidth}{0.5pt}\end{center}

\hypertarget{sixth-challenge-octet}{%
\section{\texorpdfstring{Sixth Challenge ---
\emph{Octet}}{Sixth Challenge --- Octet}}\label{sixth-challenge-octet}}

\emph{{[}The short day is almost spent; the window is going from grey
toward dark, and the steady fall beyond it has hushed the whole town so
that not a sound comes up from the street. The fire has burned down to a
low, even bed of red, and the lamp is taking the room, its light lying
warm along the beam where the cords hang dead still in the quiet.{]}}

\textbf{Tortoise:} Count them and see. Nobody has touched it --- it has
hung there through everything. Every loose thing we tried is down; only
the frame is left.

\emph{{[}The cradle stands alone on the beam now, every demonstration
cord taken down: the two feed-cords in from the left, the four cords of
the cross, the two draw-cords out to the right, and nothing else strung
anywhere.{]}}

\textbf{Achilles:} \emph{{[}reaching up, touching each cord as he
counts{]}} One\ldots{} two --- the two coming in on the left.
\emph{{[}his hand moves to the cross{]}} Three, four, five, six --- the
whole cross of them, every green to every red. \emph{{[}and to the
right{]}} Seven\ldots{} eight, going out. \emph{{[}He lets his hand
fall.{]}} Eight cords. Eight, and not one of them spare --- every one
carrying its part, and the whole thing meaning nothing if you pull a
single one out. \emph{{[}He is quiet a moment.{]}} Eight. It's a good
number to end on.

\textbf{Crab:} It is the whole ensemble, playing at once. \emph{{[}He
does not press it; he only looks, content, at the thing he has
built.{]}}

\textbf{Tortoise:} And here is the last thing worth seeing. \emph{{[}She
takes hold of the cradle's cords and begins to reshape them --- sliding
the rings along the beam, bending the cross into new curves, drawing the
whole tangle into a shape that looks nothing like it did.{]}}

\emph{{[}The cradle as it stood: two greens left, two reds right, a neat
cross between. Her hands move the rings along the beam and coax the four
crossed cords into long looping bends, until the shape on the beam is a
different picture entirely -- crooked where it was square, curved where
it was straight. But she has untied nothing and joined nothing new:
every cord still runs from the ring it left to the ring it reached. Only
where the rings sit, and the path each cord takes to get there, has
changed.{]}}

\emph{{[}The snow-hush deepens until the room is wholly still, and the
last of the grey light, coming level off the pane, runs the whole length
of the long print on the wall --- the little town, the chequered
squares, the whole strip lit end to end for a moment, so that it seems
to lean in a little out of the dusk.{]}}

\textbf{Tortoise:} \emph{{[}glancing at the wall, where Escher's
\textup{Metamorphosis II} hangs --- birds becoming fish becoming a town,
the strip flowing on and never once broken{]}} Like the print there: the
shape flows on and on, and the thing stays the same. Only what joins to
what ever mattered. \emph{{[}She sets the reshaped tangle on the
koekpan.{]}} Watch.

\emph{{[}The Crab presses, his eyes on the fire. Off the second hooks
comes the very same knot-pair the neat cradle gave: one red gathering
the two in, one cord, one green issuing the two out. A shape drawn
crooked and a shape drawn square cook down to the one pair. What changed
was only the drawing; what held was every connection.{]}}

\emph{{[}Beyond the pane, in the last pale glow the snow throws up of
itself, the town has gone all one white: roof into road into frozen
canal, the sharp corners rounded off and the crossings smoothed away,
until the whole of it is a single unbroken sheet. And yet not a house
has stirred --- each still stands where it always stood, every lane
still runs where it ran, the whole place exactly itself under a coat
that has quite changed the look of it.{]}}

\textbf{Achilles:} The same little pair, out of that mare's nest. It
didn't matter how you drew it.

\textbf{Crab:} \emph{{[}lifting the pressed pair down; the day at the
window has gone from grey to dark{]}} Not a whit. \emph{{[}He looks
along the eight cords of the cradle, then at the two of them.{]}} Well.
We've tied these, and pressed them, and you've counted them at last ---
a whole short afternoon of it, and the fire nearly down. And do you know
what strikes me? Not once, in all of it, did any of us \emph{look}. We
never once watched a cord to see which way it truly runs --- we only
ever pressed a copy and read that.

\textbf{Tortoise:} \emph{{[}gathering the last of the loose cord as the
light fails{]}} No.~Tomorrow we look.

\emph{{[}Outside, the snow has stopped. The town lies under it all one
unbroken white, and the window is dark now but for that faint pallor
thrown back off the ground. Inside, the lamp burns steady and the fire
has gone to a warm red bed that no one troubles to feed; the koekpan
cools on the table where it did its work, only an old iron pan again. On
the beam the cradle hangs exactly as it was tied and counted through the
whole short afternoon, every cord still and dark against the last light.
Along the wall the long print has gone quiet, a little lamplight caught
in it, the town and the squares receding into the dim; and nobody in the
warm room reaches to disturb anything.{]}}

\hypertarget{chapter-8-spiders-red-and-green}{%
\chapter{Chapter 8: Spiders, Red and
Green}\label{chapter-8-spiders-red-and-green}}

\emph{--- Octet for Spiders ---}

\emph{Part III: Pictures Become Proofs}

\emph{Escher: ``Metamorphosis II'' --- one form flowing into another,
birds into fish into a town, while the strip stays one continuous thing}

\begin{center}\rule{0.5\linewidth}{0.5pt}\end{center}

Can a proof be a drawing --- can we settle a fact about real quantum
gates by sliding coloured dots along a few wires, never once going back
to the algebra to check?

For seven chapters the pictures have travelled beside the algebra as a
faithful companion, and never once out in front. Every time a wire was
pulled straight or a box slid round a bend, the move was
\emph{licensed}: some theorem on the blackboard signed for it, and the
picture cashed the cheque. The sketchbook could show you a thing
beautifully, but it could not, on its own, \emph{establish} one. At the
very end of the last chapter that balance began to tip. We had drawn a
single new mark --- a dot where three wires meet --- and found that it
already knew things: that any one of its legs could be predicted from
the other two, a fact we needed six sealed numbers to contradict. A dot
that knows things is the beginning of a language.

This chapter builds that language and then turns it loose. The dot from
Chapter 7 was one fixed thing, always the state \(\ket{000}+\ket{111}\),
unlabelled and alone. We now give it a family. It gets a \textbf{colour}
--- and there will be a second colour, which counts the world in a
different way. It gets a \textbf{turn}, the sphere-rotation of Chapter
3, marked on the dot itself. It gets any number of legs. And once there
are two colours of dot with turns and legs, a handful of rules for
pushing them past one another appears --- rules that never mention a
matrix. With those rules a diagram can be \emph{rewritten} into another
diagram of equal value, the way \(3+4\) is rewritten into \(7\),
entirely inside the pictures.

At the end we will take three of the commonest gates in all of
computing, tangle them together, and cook the tangle down to a single
crossing --- a proof, done with coloured dots, that three of one gate
make another. And then we will count the wires in the picture we drew,
and something will fall quietly into place.

The calculus we are about to draw has a name and an author, given
plainly at the chapter's end. For now, only the colours.

\begin{center}\rule{0.5\linewidth}{0.5pt}\end{center}

\hypertarget{colour-comes-to-the-dot}{%
\section{§8.1 Colour Comes to the Dot}\label{colour-comes-to-the-dot}}

The three-legged dot of Chapter 7 was defined by a single habit: on the
legs it read \(\ket{000}+\ket{111}\) --- all-zeros or all-ones, agreeing
across every leg, nothing in between. That habit does not care how many
legs there are, nor which are inputs and which outputs. Let us free it.

\textbf{On the blackboard.} For any number of input legs \(m\) and
output legs \(k\), and any angle \(\alpha\), define the \textbf{Z
spider}

\[Z(m,k,\alpha) \;=\; \ket{0\cdots0}\bra{0\cdots0} \;+\; e^{i\alpha}\,\ket{1\cdots1}\bra{1\cdots1},\]

with \(k\) zeros (or ones) in the ket and \(m\) in the bra. In words: it
passes the all-agree inputs to the all-agree outputs, and it stamps the
all-ones branch with the phase \(e^{i\alpha}\) --- the very turn of the
sphere we met in Chapter 3, now written on the dot. The angle \(\alpha\)
is the spider's \textbf{phase}. Chapter 7's node is exactly
\(Z(0,3,0)\): no inputs, three outputs, no turn --- the state
\(\ket{000}+\ket{111}\). And the plainest spider of all,
\(Z(1,1,0)=\ket0\bra0+\ket1\bra1\), is the identity matrix: one leg in,
one leg out, no turn, a bare wire.

\textbf{In the sketchbook.} We draw the Z spider as a \textbf{green dot}
with its legs fanning out --- \(m\) wires arriving on the left, \(k\)
leaving on the right --- and its phase written beside it. A dot with no
number written is understood to carry phase \(0\).

\begin{center}
\begin{tikzpicture}[sd,
  zsp/.style={circle,draw,line width=0.9pt,fill=green!55,minimum size=1.4em,inner sep=0pt}]
  \draw[wire] (-1.6,0.7)--(0,0); \draw[wire] (-1.6,-0.7)--(0,0);
  \draw[wire] (0,0)--(1.6,0.7); \draw[wire] (0,0)--(1.6,0); \draw[wire] (0,0)--(1.6,-0.7);
  \node[zsp,label=above:{$\alpha$}] at (0,0) {};
\end{tikzpicture}
\end{center}

Everything that made the Chapter 7 dot useful survives: it is symmetric
in its legs, and it does not distinguish input from output --- a leg can
be bent from one side to the other with the cups and caps of Chapter 6,
and the dot does not mind.

Two green dots joined by a wire are asking the same all-agree question
twice in a row, which is asking it once. That is the first rule.

\begin{quote}
\textbf{Theorem 8.1 (fusion).} Two Z spiders that share at least one leg
are a single Z spider, and their phases add. On the blackboard, for
spiders that meet along one wire,
\[Z(1,1,\beta)\,Z(m,1,\alpha) \;=\; Z(m,1,\alpha+\beta),\] and the same
for any number of shared and free legs. The equality is \textbf{exact}
--- no scalar.
\end{quote}

\emph{Proof idea.} The all-zeros branch of the first spider feeds only
the all-zeros branch of the second, and the all-ones feeds only the
all-ones; the two phases sit on the same all-ones branch and so
multiply, \(e^{i\alpha}e^{i\beta}=e^{i(\alpha+\beta)}\).

\emph{Proof (blackboard).} Compose
\(\bigl(\ket0\bra0+e^{i\beta}\ket1\bra1\bigr)\bigl(\ket{0}\bra{0\cdots0}+e^{i\alpha}\ket{1}\bra{1\cdots1}\bigr)\).
The cross terms vanish because \(\braket{0|1}=\braket{1|0}=0\); what
survives is
\(\ket0\bra{0\cdots0}+e^{i(\alpha+\beta)}\ket1\bra{1\cdots1}=Z(m,1,\alpha+\beta)\).
\(\;\blacksquare\)

In the picture the rule is even plainer: two green dots slide together
into one, and you add the numbers written beside them.

\[\begin{ZX} \zxZ{\alpha} \rar & \zxZ{\beta} \end{ZX}
\;\;=\;\;
\begin{ZX} \zxZ{\alpha+\beta} \end{ZX}\]

See \texttt{Code/Ch8/Fusion.py}, which checks fusion for random phases
\(\alpha,\beta\neq0\) and confirms \(Z(1,1,\pi)=\mathrm{diag}(1,-1)\).

\textbf{\emph{Why this definition.}} We could have kept the dot rigid,
one fixed state. But a dot that carries a phase and any number of legs
turns a \emph{fact about one entangled state} into a \emph{piece of
grammar}: fusion holds for every \(m\), \(k\), \(\alpha\) at once. One
rule now does what a page of tensor bookkeeping did before.

\begin{quote}
\textbf{Coecke's Sketchbook.}

\emph{The calculus, and whose it is.} From here the pictures stop being
a companion to the algebra and become a calculus in their own right ---
one that can be computed with, not merely read. It is the work of Bob
Coecke and his colleagues at Oxford, built on the categorical reading of
quantum processes that he and Samson Abramsky began; the green--red
diagrams themselves are due to Coecke and Ross Duncan. Its governing
slogan, which we will earn in §8.5, is startling the first time:
\emph{only the connectivity matters}. Slide a dot along its wire, bend a
leg from one side to the other, redraw the whole tangle --- so long as
the same dots stay joined to the same dots, the value does not change. A
picture, in this calculus, is not a diagram \emph{of} a computation. It
\emph{is} the computation, and its shape is negotiable.
\end{quote}

\begin{center}\rule{0.5\linewidth}{0.5pt}\end{center}

\hypertarget{one-wire-counted-twice}{%
\section{§8.2 One Wire, Counted Twice}\label{one-wire-counted-twice}}

The green dot reads a wire in one particular way: it asks whether the
legs all say \(0\) or all say \(1\). That is a choice of basis --- the
\(\ket0,\ket1\) basis. But Chapter 6 taught us there is nothing sacred
about it; the state \(\ket\pm=\tfrac1{\sqrt2}(\ket0\pm\ket1)\) makes an
equally good pair of alternatives. A dot could just as well count in
\emph{that} basis. This is the second colour.

\textbf{On the blackboard.} Define the \textbf{X spider} intrinsically
--- not by borrowing the green one, but in its own basis:

\[\begin{aligned}
X(m,k,\beta) &\;=\; \ket{+\cdots+}\bra{+\cdots+} \;+\; e^{i\beta}\,\ket{-\cdots-}\bra{-\cdots-}, \\
\ket\pm &\;=\; \tfrac1{\sqrt2}\bigl(\ket0\pm\ket1\bigr).
\end{aligned}\]

It is the same shape of definition as the green spider, with \(\ket\pm\)
everywhere the green one had \(\ket0,\ket1\). It passes the all-plus
inputs to all-plus outputs and the all-minus to all-minus, stamping the
minus branch with \(e^{i\beta}\). Because \(\ket\pm\) are unit vectors,
so is each term; we build the red spider this way, from normalized
\(\ket\pm\), precisely so that the colour-change rule below is a
\emph{theorem} we prove, and not a definition we have quietly assumed.

Two small facts pin the red dot to earth. Its single-leg, phase-free
state is \(X(0,1,0)=\ket++\ket-=\sqrt2\,\ket0\) --- an unnormalized
\(\ket0\), carrying the same \(\sqrt2\) our cups have carried since
Chapter 6. And \(X(1,1,\pi)=\ket+\bra+-\ket-\bra-\) is the bit-flip
\(\bigl(\begin{smallmatrix}0&1\\1&0\end{smallmatrix}\bigr)\), the NOT
gate.

\textbf{In the sketchbook.} We draw the X spider as a \textbf{red dot},
otherwise exactly like the green one: any number of legs, a phase
written beside it, symmetric in its legs, indifferent to which side a
leg leaves on.

To pass between the two colours we need one more box, the one gate that
turns the \(\ket0,\ket1\) basis into the \(\ket\pm\) basis and back: the
Hadamard
\(H=\tfrac1{\sqrt2}\bigl(\begin{smallmatrix}1&1\\1&-1\end{smallmatrix}\bigr)\)
of Chapter 3, with \(H\ket0=\ket+\), \(H\ket1=\ket-\), and \(H^2=I\). In
this calculus it gets its own mark: a small \textbf{yellow box} on a
wire.

\begin{center}
\begin{tikzpicture}[sd,
  hbox/.style={draw,line width=0.9pt,fill=yellow!85,minimum size=1.0em,inner sep=1.5pt}]
  \draw[wire] (0,0)--(3,0);
  \node[hbox] at (1.5,0) {$H$};
\end{tikzpicture}
\end{center}

A Hadamard on a wire is so common between two dots that the calculus
abbreviates it: a \textbf{dashed wire} --- the \emph{Hadamard-edge} ---
means exactly ``an ordinary wire with a yellow box on it''. Drawn
between two dots, the two are the same picture:

\[\begin{tikzpicture}[sd,
  zsp/.style={circle,draw,line width=0.9pt,fill=green!55,minimum size=1.15em,inner sep=0pt},
  xsp/.style={circle,draw,line width=0.9pt,fill=red!75,minimum size=1.15em,inner sep=0pt},
  hbox/.style={draw,line width=0.9pt,fill=yellow!85,minimum size=0.9em,inner sep=1pt}]
  \draw[wire] (-1.0,0)--(0,0); \node[zsp] at (0,0){};
  \draw[wire] (0,0)--(2.0,0); \node[hbox] at (1.0,0){};
  \node[xsp] at (2.0,0){}; \draw[wire] (2.0,0)--(3.0,0);
\end{tikzpicture}
\;\;=\;\;
\begin{tikzpicture}[sd,
  zsp/.style={circle,draw,line width=0.9pt,fill=green!55,minimum size=1.15em,inner sep=0pt},
  xsp/.style={circle,draw,line width=0.9pt,fill=red!75,minimum size=1.15em,inner sep=0pt}]
  \draw[wire] (-1.0,0)--(0,0); \node[zsp] at (0,0){};
  \draw[wire,dashed] (0,0)--(2.0,0);
  \node[xsp] at (2.0,0){}; \draw[wire] (2.0,0)--(3.0,0);
\end{tikzpicture}\]

a shorthand later chapters will lean on when tangles grow crowded. Now
the bridge between the colours.

\hypertarget{an-example-that-forces-the-rule}{%
\subsection{An example that forces the
rule}\label{an-example-that-forces-the-rule}}

Take the plainest green dot with a turn,
\(Z(1,1,\alpha)=\ket0\bra0+e^{i\alpha}\ket1\bra1\), and put a Hadamard
on each of its two legs. On the blackboard \(H\,Z(1,1,\alpha)\,H\) sends
\(\ket+\mapsto\ket0\mapsto\ket0\mapsto\ket+\) and
\(\ket-\mapsto\ket1\mapsto e^{i\alpha}\ket1\mapsto e^{i\alpha}\ket-\);
that is \(\ket+\bra++e^{i\alpha}\ket-\bra-\), which is the \emph{red}
spider \(X(1,1,\alpha)\) --- exactly, for every \(\alpha\). The colour
did not change by fiat; it changed because dressing the green counter in
Hadamards makes it count in the \(\ket\pm\) basis, which is what red
\emph{means}.

\begin{quote}
\textbf{Theorem 8.2 (colour change).} A Hadamard on every leg of a Z
spider turns it into an X spider of the same phase (and the reverse,
since \(H^2=I\)):
\[H^{\otimes k}\,Z(m,k,\alpha)\,H^{\otimes m} \;=\; X(m,k,\alpha),\]
\textbf{exactly}, for every \(m\), \(k\), \(\alpha\).
\end{quote}

\emph{Proof idea.} \(H\) carries \(\ket0,\ket1\) to \(\ket+,\ket-\);
conjugating the green spider by \(H\) on all legs rewrites its all-agree
action in the \(\ket\pm\) basis, which is the red spider by definition.

\emph{Proof (blackboard).}
\(H^{\otimes k}\ket{0\cdots0}=\ket{+\cdots+}\) and
\(H^{\otimes k}\ket{1\cdots1}=\ket{-\cdots-}\), and likewise for the
bras with \(H^{\otimes m}\). Substituting into \(Z(m,k,\alpha)\) turns
every \(\ket{0\cdots0}\) into \(\ket{+\cdots+}\) and every
\(\ket{1\cdots1}\) into \(\ket{-\cdots-}\), giving \(X(m,k,\alpha)\)
term for term. The phase rides untouched on the second branch.
\(\;\blacksquare\) See \texttt{Code/Ch8/ColourChange.py}, which builds
\(X\) intrinsically from \(\ket\pm\) and checks the equality for
\(\alpha\in\{0,0.7,1.3,2.9\}\).

In the picture: a green dot with a yellow box on each leg is a red dot.

\[\begin{tikzpicture}[sd,
  zsp/.style={circle,draw,line width=0.9pt,fill=green!55,minimum size=1.3em,inner sep=0pt},
  hbox/.style={draw,line width=0.9pt,fill=yellow!85,minimum size=0.9em,inner sep=1pt}]
  \draw[wire] (-2.2,0.6)--(0,0); \draw[wire] (-2.2,-0.6)--(0,0);
  \draw[wire] (0,0)--(2.2,0.6); \draw[wire] (0,0)--(2.2,-0.6);
  \node[hbox] at (-1.3,0.36){}; \node[hbox] at (-1.3,-0.36){};
  \node[hbox] at (1.3,0.36){}; \node[hbox] at (1.3,-0.36){};
  \node[zsp] at (0,0){};
\end{tikzpicture}
\;\;=\;\;
\begin{tikzpicture}[sd,
  xsp/.style={circle,draw,line width=0.9pt,fill=red!75,minimum size=1.3em,inner sep=0pt}]
  \draw[wire] (-1.6,0.6)--(0,0); \draw[wire] (-1.6,-0.6)--(0,0);
  \draw[wire] (0,0)--(1.6,0.6); \draw[wire] (0,0)--(1.6,-0.6);
  \node[xsp] at (0,0){};
\end{tikzpicture}\]

\textbf{\emph{Why this definition.}} Building red from its own
\(\ket\pm\) costs us one honest theorem --- colour change --- instead of
getting it free. That is the point. Had we \emph{defined} red as ``green
wrapped in Hadamards'', the rule would be a tautology and would prove
nothing about the world. Defined intrinsically, colour change is a real
fact: the two counters see the same wire two genuinely different ways.

\begin{quote}
\textbf{Dirac's Blackboard.}

\emph{Spiders are matrices.} Every dot is an honest linear map, and it
is worth writing four of them down. The one-in-one-out green dots are
diagonal:
\(Z(1,1,0)=\bigl(\begin{smallmatrix}1&0\\0&1\end{smallmatrix}\bigr)=I\),
and
\(Z(1,1,\pi)=\bigl(\begin{smallmatrix}1&0\\0&-1\end{smallmatrix}\bigr)\),
the phase-flip. The red ones are their Hadamard conjugates:
\(X(1,1,0)=I\) again, and
\(X(1,1,\pi)=\bigl(\begin{smallmatrix}0&1\\1&0\end{smallmatrix}\bigr)\),
the bit-flip. The legs change the shape:
\(Z(1,2,0)=\ket{00}\bra0+\ket{11}\bra1\) is a \(4\times2\) matrix that
copies a basis bit, and \(Z(0,1,0)=\ket0+\ket1\) is the column
\((1,1)^{T}\). Reading a spider off as a matrix is always this
mechanical: put a \(1\) (or \(e^{i\alpha}\)) where all legs agree, a
\(0\) everywhere else. The pictures never lie about the numbers; they
only decline to write them all out.
\end{quote}

\begin{center}\rule{0.5\linewidth}{0.5pt}\end{center}

\hypertarget{each-colour-copies-the-other}{%
\section{§8.3 Each Colour Copies the
Other}\label{each-colour-copies-the-other}}

We have two counters that read the same wire in incompatible bases. What
happens where a red input meets a green dot --- where the two ways of
counting collide?

\textbf{On the blackboard.} A green copier is
\(Z(1,2,0)=\ket{00}\bra0+\ket{11}\bra1\). On its own green basis points
it does what its name says: \(Z(1,2,0)\ket0=\ket{00}\) and
\(Z(1,2,0)\ket1=\ket{11}\) --- one point in, two identical points out.
Now feed it a \emph{red} point. The red one-leg states are the ones we
met in §8.2: \(\ket{r_0}=X(0,1,0)=\sqrt2\,\ket0\) and
\(\ket{r_1}=X(0,1,\pi)=\sqrt2\,\ket1\), each carrying the cup's
\(\sqrt2\). Watch the green copier act on \(\ket{r_0}\):

\[\begin{aligned}
Z(1,2,0)\,\ket{r_0} &\;=\; \sqrt2\,Z(1,2,0)\ket0 \;=\; \sqrt2\,\ket{00} \\
&\;=\; \tfrac1{\sqrt2}\bigl(\sqrt2\ket0\bigr)\otimes\bigl(\sqrt2\ket0\bigr) \;=\; \tfrac1{\sqrt2}\,\ket{r_0}\otimes\ket{r_0}.
\end{aligned}\]

So green \emph{does} copy a red point --- into two red points --- but
with a factor \(\tfrac1{\sqrt2}\) out front. The honest statement is
that a phase-free green spider copies a red point \textbf{up to the
scalar \(\sqrt2\)}:
\(Z(1,2,0)\,\ket{r}=\tfrac1{\sqrt2}\,\ket{r}\otimes\ket{r}\) for the red
points \(\ket r\), and symmetrically a phase-free red spider copies
green points. Each colour is copied cleanly by the other.

\textbf{In the sketchbook.} A red point (a red dot with one leg) fed
into a phase-free green dot comes out as two red points. This is the
picture Achilles did not expect: one cord in, two cords out.

\[\begin{tikzpicture}[sd,
  zsp/.style={circle,draw,line width=0.9pt,fill=green!55,minimum size=1.3em,inner sep=0pt},
  xsp/.style={circle,draw,line width=0.9pt,fill=red!75,minimum size=1.1em,inner sep=0pt}]
  \node[xsp] at (0,0){};
  \draw[wire] (0,0)--(1.6,0);
  \node[zsp] at (1.6,0){};
  \draw[wire] (1.6,0)--(3.0,0.55); \draw[wire] (1.6,0)--(3.0,-0.55);
\end{tikzpicture}
\;\;=\;\;\tfrac1{\sqrt2}\;\;
\begin{tikzpicture}[sd,
  xsp/.style={circle,draw,line width=0.9pt,fill=red!75,minimum size=1.1em,inner sep=0pt}]
  \node[xsp] at (0,0.55){}; \draw[wire] (0,0.55)--(1.4,0.55);
  \node[xsp] at (0,-0.55){}; \draw[wire] (0,-0.55)--(1.4,-0.55);
\end{tikzpicture}\]

The factor \(\tfrac1{\sqrt2}\) is drawn because from this chapter on we
allow ourselves to state a picture equality \textbf{up to a non-zero
scalar}, writing the scalar in only when it matters. Copy differs from a
clean duplication by a factor of \(\sqrt2\); we say so and move on.

A companion fact concerns a half-turn --- a phase of \(\pi\) --- pushed
through the other colour.

\begin{quote}
\textbf{Theorem 8.3 (copy).} A phase-free spider of one colour copies
the basis points of the other:
\(\;Z(1,2,0)\,\ket r=\tfrac1{\sqrt2}\,\ket{rr}\) for red points
\(\ket r\), and symmetrically. \emph{(Up to \(\sqrt2\).)}

\textbf{Theorem 8.4 (\(\pi\)-commutation).} A half-turn of one colour
passes through the other colour and reverses the other's turn:
\[X(1,1,\pi)\,Z(1,1,\alpha) \;=\; Z(1,1,-\alpha)\,X(1,1,\pi),\] up to a
non-zero scalar (and exactly when \(\alpha=0\)).
\end{quote}

\emph{Proof idea (copy).} The green copier duplicates whatever agrees in
\emph{its} basis; a red point, seen from the green side, is an even
superposition that the copier spreads across the diagonal
\(\ket{00}+\ket{11}\) --- which is two red points, once the \(\sqrt2\)
is accounted.

\emph{Proof idea (\(\pi\)-commutation).} The red half-turn is the
bit-flip \(X(1,1,\pi)\); conjugating a green phase gate by it sends
\(\ket1\mapsto\ket0\) and back, turning \(e^{i\alpha}\) on the \(\ket1\)
branch into \(e^{i\alpha}\) on the \(\ket0\) branch --- a global phase
times \(e^{-i\alpha}\) on \(\ket1\), which is \(Z(1,1,-\alpha)\) up to
that phase.

In the picture, \(\pi\)-commutation is a red \(\pi\)-dot sliding through
a green dot and flipping the green dot's turn:

\[\begin{ZX} \zxX{\pi} \rar & \zxZ{\alpha} \end{ZX}
\;\;=\;\;
\begin{ZX} \zxZ{-\alpha} \rar & \zxX{\pi} \end{ZX}
\qquad\text{(up to a scalar).}\]

See \texttt{Code/Ch8/Copy.py}, which checks copy up to \(\sqrt2\) and
\(\pi\)-commutation for \(\alpha\in\{0.7,1.3\}\) --- not merely
\(\alpha=0\), where it would hold trivially.

\textbf{\emph{Why this definition.}} We insist on \emph{phase-free}
copiers because a spider with a turn does not copy cleanly: the turn
would have to be duplicated too, and \(e^{i\alpha}\) does not split into
two honest points. The rule is exactly as strong as it is, and no
stronger --- which is why the scalar and the phase-free clause are both
written down.

\begin{center}\rule{0.5\linewidth}{0.5pt}\end{center}

\hypertarget{cradle-that-comes-undone}{%
\section{§8.4 Cradle That Comes Undone}\label{cradle-that-comes-undone}}

Copy says green duplicates red points and red duplicates green points.
Put two of each together and the two habits collide in the most
productive way in the whole calculus.

\hypertarget{the-example-a-tangle-that-ought-to-simplify}{%
\subsection{The example: a tangle that ought to
simplify}\label{the-example-a-tangle-that-ought-to-simplify}}

Take two green dots feeding two red dots, with \textbf{every} green
joined to \textbf{every} red --- a complete little cross, four cords,
the shape a child strings between two pairs of nails. Two wires come in
on the greens; two leave from the reds. On the blackboard this is
\(\bigl[X(2,1)\otimes X(2,1)\bigr]\circ(\text{cross})\circ\bigl[Z(1,2)\otimes Z(1,2)\bigr]\),
a \(4\times4\) matrix. Multiply it out (the calculation is short but
joyless) and it equals

\[\tfrac12\Bigl(\ket{00}\bigl(\bra{00}+\bra{11}\bigr)+\ket{11}\bigl(\bra{01}+\bra{10}\bigr)\Bigr),\]

which reads the parity of its two inputs and writes that parity onto
\emph{both} outputs. Now here is the surprise: the very same map, up to
a scalar, is made by a single green dot tied to a single red dot --- no
cross at all.

\begin{quote}
\textbf{Theorem 8.5 (bialgebra).} Two Z spiders joined to two X spiders
by a complete cross of four wires equal one Z spider joined to one X
spider by a single wire --- the inputs gathered on the red, the outputs
issuing from the green:
\[\bigl[X(2,1)\otimes X(2,1)\bigr]\,(\text{cross})\,\bigl[Z(1,2)\otimes Z(1,2)\bigr] \;=\; Z(1,2)\,X(2,1),\]
\textbf{up to the scalar \(\sqrt2\)}.
\end{quote}

\[\begin{tikzpicture}[sd,
  zsp/.style={circle,draw,line width=0.9pt,fill=green!55,minimum size=1.25em,inner sep=0pt},
  xsp/.style={circle,draw,line width=0.9pt,fill=red!75,minimum size=1.25em,inner sep=0pt}]
  \coordinate (gt) at (0,0.9); \coordinate (gb) at (0,-0.9);
  \coordinate (rt) at (2.6,0.9); \coordinate (rb) at (2.6,-0.9);
  \draw[wire] (-1.3,0.9)--(gt); \draw[wire] (-1.3,-0.9)--(gb);
  \draw[wire] (rt)--(3.9,0.9); \draw[wire] (rb)--(3.9,-0.9);
  \draw[wire] (gt)--(rt); \draw[wire] (gb)--(rb);
  \draw[wire] (gt)--(rb); \draw[wire] (gb)--(rt);
  \node[zsp] at (gt){}; \node[zsp] at (gb){};
  \node[xsp] at (rt){}; \node[xsp] at (rb){};
\end{tikzpicture}
\;\;=\;\;\tfrac1{\sqrt2}\;\;
\begin{tikzpicture}[sd,
  zsp/.style={circle,draw,line width=0.9pt,fill=green!55,minimum size=1.25em,inner sep=0pt},
  xsp/.style={circle,draw,line width=0.9pt,fill=red!75,minimum size=1.25em,inner sep=0pt}]
  \coordinate (r) at (0,0); \coordinate (g) at (1.8,0);
  \draw[wire] (-1.3,0.7)--(r); \draw[wire] (-1.3,-0.7)--(r);
  \draw[wire] (r)--(g);
  \draw[wire] (g)--(3.1,0.7); \draw[wire] (g)--(3.1,-0.7);
  \node[xsp] at (r){}; \node[zsp] at (g){};
\end{tikzpicture}\]

\emph{Proof idea.} Copy is the engine. Push each green copier through
each red adder: the green wants to duplicate, the red wants to combine,
and when a whole cross of them does this at once, the duplications and
combinations rearrange into a single duplicate-then-combine. The cross
is the copying written twice over.

\emph{Proof (blackboard).} This is the multiplication above: both sides
send \(\ket{xy}\) to \(\ket{pp}\) with \(p=x\oplus y\) the parity, and
the scalars are \(\tfrac12\) on the left and \(\tfrac1{\sqrt2}\) on the
right, a ratio of \(\sqrt2\). Every ``in general'' here is a finite
check on the four inputs \(\ket{00},\ket{01},\ket{10},\ket{11}\); see
\texttt{Code/Ch8/Complementarity.py}. \(\;\blacksquare\)

The bialgebra has a twin, about \emph{two} cords rather than four. If a
green dot and a red dot are joined by two separate wires, they might as
well be joined by none.

\begin{quote}
\textbf{Theorem 8.6 (Hopf law).} A green spider and an X spider linked
by two distinct wires equal the two spiders disconnected --- a green
point beside a red point:
\[X(2,1)\,Z(1,2) \;=\; \bigl(\ket{r}\bra{g}\bigr),\] the rank-one
disconnected map, \textbf{up to the scalar \(2\)}.
\end{quote}

\emph{Proof idea.} Two wires between the colours let each copy the
other's point and then immediately re-add it; a copy followed by its own
undoing sends everything to a single fixed direction on each side, which
is exactly what ``disconnected'' means. Complementarity of the two bases
\emph{is} this law --- the two facts, bialgebra and Hopf, are one
phenomenon seen twice (Coecke--Duncan, §8).

\[\begin{tikzpicture}[sd,
  zsp/.style={circle,draw,line width=0.9pt,fill=green!55,minimum size=1.25em,inner sep=0pt},
  xsp/.style={circle,draw,line width=0.9pt,fill=red!75,minimum size=1.25em,inner sep=0pt}]
  \draw[wire] (-1.2,0)--(0,0);
  \node[zsp] (g) at (0,0){};
  \draw[wire] (g) to[out=55,in=125] (2,0);
  \draw[wire] (g) to[out=-55,in=-125] (2,0);
  \node[xsp] (r) at (2,0){};
  \draw[wire] (r)--(3.2,0);
\end{tikzpicture}
\;\;=\;\;\tfrac12\;\;
\begin{tikzpicture}[sd,
  zsp/.style={circle,draw,line width=0.9pt,fill=green!55,minimum size=1.25em,inner sep=0pt},
  xsp/.style={circle,draw,line width=0.9pt,fill=red!75,minimum size=1.25em,inner sep=0pt}]
  \draw[wire] (-1.0,0)--(0,0); \node[zsp] at (0,0){};
  \draw[wire] (1.6,0)--(2.6,0); \node[xsp] at (1.6,0){};
\end{tikzpicture}\]

\emph{Proof (blackboard).} With the intrinsic red spider,
\(X(2,1)Z(1,2)=\tfrac1{\sqrt2}\bigl(\begin{smallmatrix}1&1\\0&0\end{smallmatrix}\bigr)=\ket0\bra+\)
--- a rank-one map onto a single direction, hence disconnected. The
disconnected picture on the right --- a red state \(\sqrt2\ket0\) beside
a green effect \(\bra0+\bra1\) --- is
\(\sqrt2\ket0\bigl(\bra0+\bra1\bigr)=2\ket0\bra+\); the two agree up to
the scalar \(2\). See \texttt{Code/Ch8/Complementarity.py}, which builds
the genuine two-wire link (not one wire drawn twice) and derives the
disconnection with the ZX rules. \(\;\blacksquare\)

\textbf{\emph{Why this definition.}} These two are the heart of the
chapter, and the reason the pictures can now outrun the algebra. Fusion
and colour change only tidy dots of one colour at a time; the bialgebra
and Hopf laws are the first rules that let the two colours
\emph{reorganise each other}, turning a dense tangle into a sparse one.
Every reduction to come is these two, applied and reapplied.

\begin{center}\rule{0.5\linewidth}{0.5pt}\end{center}

\hypertarget{keeping-only-the-topology}{%
\section{§8.5 Keeping Only the
Topology}\label{keeping-only-the-topology}}

We have been redrawing diagrams freely --- sliding dots, bending legs,
stretching wires --- and asserting that the value is unchanged. It is
time to say plainly what that freedom is, and what backs it.

\textbf{The claim.} In these diagrams \emph{only the connectivity
matters}: which dot is joined to which, and by how many wires. The
lengths of the wires, the positions of the dots on the page, the order
in which two wires happen to cross with no dot at the crossing --- none
of it changes the map the picture denotes. Two diagrams that are the
same graph, with the same dots joined the same way, are equal.

This is the slogan the Escher woodcut on the scullery wall makes
visible: in \emph{Metamorphosis II} the birds become fish become a town,
the shape flowing on and on, and yet the strip is one continuous thing
throughout --- what changes is never the connection, only the drawing.

We \textbf{state} this and do not prove it here; it is a coherence
theorem, and the honest citation is Peter Selinger's survey. Because our
diagrams have crossings --- swaps, from Chapter 4 --- the backing is
coherence for \emph{symmetric monoidal} categories (Selinger, Theorem
3.12, resting on Joyal and Street): an equation holds in the graphical
language exactly when the two diagrams have the same connectivity. It is
\textbf{not} the planar coherence of Selinger's Theorem 3.1, which he
flags in Caveat 3.2 as proved only for deformations that never turn a
wire back on itself --- too weak once wires cross. And it is worth
keeping one distinction sharp: the extra freedom to \emph{permute a
spider's legs at will} is a separate fact, the spider theorem of Coecke
and Duncan (§7), not part of the coherence result, whose boxes keep
their legs in order. So the slogan is ``only the connectivity matters''
--- never ``bend anything however you like.''

\textbf{How far the calculus reaches.} One naturally asks what these two
colours of dot can express.

\begin{quote}
\textbf{Theorem 8.7 (universality and completeness; stated, not
proved).} Every pure quantum state, every unitary gate, and every
post-selected projective measurement on any number of qubits can be
written as a diagram of green and red spiders and Hadamard boxes
(Coecke--Duncan). And within the \emph{stabilizer} fragment --- the part
of quantum theory built from Hadamards, phase gates and CNOTs --- the
calculus is \textbf{complete}: any two diagrams denoting the same map
can be turned into one another by the rules alone, with no return to
matrices (Backens).
\end{quote}

We will not prove these; the citations discharge them. But they are the
licence for the chapter's ambition. Universality says the pictures miss
nothing; completeness says that, in a large and important fragment, the
rules miss nothing either --- every true equation of matrices is
reachable by pushing dots. What we do next is one small, concrete
instance of that promise.

\begin{quote}
\textbf{Coecke's Sketchbook.}

\emph{What ``only topology'' does not say.} The slogan is intoxicating
and easy to overstate, so a pencil-user's caution. It says that a
rewrite which preserves the connections preserves the \emph{map} --- an
equality of processes, no more. It does not say the shorter diagram is
faster to run, cheaper to build, or in any way optimal; a picture with
three dots and one with thirty can denote the very same gate. Nor does
redrawing prove anything about security, or about how a physical device
behaves when watched. When we cook a tangle down to a crossing in the
next section, we will have shown that two processes are \emph{equal} ---
and we will be careful to claim exactly that, and not a syllable more.
\end{quote}

\begin{center}\rule{0.5\linewidth}{0.5pt}\end{center}

\hypertarget{every-rule-on-one-page-the-rosetta-table}{%
\section{§8.6 Every Rule on One Page --- The Rosetta
Table}\label{every-rule-on-one-page-the-rosetta-table}}

Before the table, one last assembly to name, because we are about to
lean on it. Take a green dot on a top wire and a red dot on a bottom
wire, and join them by a single vertical cord --- the green with legs
going in and out on the top, the red with legs in and out on the bottom.

\begin{center}
\begin{tikzpicture}[sd,
  zsp/.style={circle,draw,line width=0.9pt,fill=green!55,minimum size=1.25em,inner sep=0pt},
  xsp/.style={circle,draw,line width=0.9pt,fill=red!75,minimum size=1.25em,inner sep=0pt}]
  \draw[wire] (-1.3,1)--(1.3,1);
  \draw[wire] (-1.3,-1)--(1.3,-1);
  \draw[wire] (0,1)--(0,-1);
  \node[zsp] at (0,1){}; \node[xsp] at (0,-1){};
\end{tikzpicture}
\end{center}

On the blackboard this green--red pair is the CNOT gate of Chapter 4 ---
control on the green (top) wire, target on the red (bottom) --- divided
by \(\sqrt2\): the bare picture is \(\mathrm{CNOT}/\sqrt2\), the missing
factor being the cup convention we have carried all along. The green
copies the control bit; the red adds it into the target. We verify this
in the next section; here it earns its Rosetta row.

The rows this chapter adds to the running dictionary, numbered on from
Chapter 7's fifty-two. This is the largest single instalment in the
book, and deliberately so: the whole ZX vocabulary --- four new picture
elements, six named rules, the topology principle, and the spider-CNOT
--- arrives together, because it is a language and a language must be
handed over whole.

\begingroup\renewcommand{\arraystretch}{1.5}

\begin{longtable}[]{@{}
  >{\raggedright\arraybackslash}p{(\columnwidth - 4\tabcolsep) * \real{0.2727}}
  >{\raggedright\arraybackslash}p{(\columnwidth - 4\tabcolsep) * \real{0.3636}}
  >{\raggedright\arraybackslash}p{(\columnwidth - 4\tabcolsep) * \real{0.3636}}@{}}
\toprule\noalign{}
\begin{minipage}[b]{\linewidth}\raggedright
\#
\end{minipage} & \begin{minipage}[b]{\linewidth}\raggedright
Algebra (Dirac)
\end{minipage} & \begin{minipage}[b]{\linewidth}\raggedright
Picture (Coecke)
\end{minipage} \\
\midrule\noalign{}
\endhead
\bottomrule\noalign{}
\endlastfoot
53 &
\(Z(m,k,\alpha)=\ket{0\cdots0}\bra{0\cdots0}+e^{i\alpha}\ket{1\cdots1}\bra{1\cdots1}\)
& a green dot, any legs, phase \(\alpha\) --- \emph{new element} \\
54 &
\(X(m,k,\beta)=\ket{+\cdots+}\bra{+\cdots+}+e^{i\beta}\ket{-\cdots-}\bra{-\cdots-}\)
(normalized \(\ket\pm\)) & a red dot, any legs, phase \(\beta\) ---
\emph{new element} \\
55 & the phase \(\alpha\) (Chapter 3's sphere-turn) & a turn written
beside a dot --- \emph{new element} \\
56 &
\(H=\tfrac1{\sqrt2}\bigl(\begin{smallmatrix}1&1\\1&-1\end{smallmatrix}\bigr)\);
\(H\,Z(\alpha)\,H=X(\alpha)\) & a yellow box; a dashed wire
(Hadamard-edge) --- \emph{new element} \\
57 & same-colour spiders on a shared leg merge; phases add (exact) & two
like dots slide together into one --- \textbf{fusion} \\
58 & \(H\) on every leg flips \(Z\!\leftrightarrow\!X\) (exact) & a dot
dressed in Hadamards changes colour --- \textbf{colour change} \\
59 & a phase-free spider copies the other colour's points (up to
\(\sqrt2\)) & a point of one colour into the other comes out doubled ---
\textbf{copy} \\
60 & \(X(\pi)\,Z(\alpha)=Z(-\alpha)\,X(\pi)\) (up to a scalar) & a
half-turn pushed through the other colour flips the turn ---
\textbf{\(\pi\)-commutation} \\
61 & \(K_{2,2}\) of two greens into two reds \(=\) one green--red pair
(up to \(\sqrt2\)) & the cross comes undone to a single knot-pair ---
\textbf{bialgebra} \\
62 & green--(two wires)--red \(=\) green\(\otimes\)red disconnected (up
to \(2\)) & two wires between the colours \(=\) none ---
\textbf{Hopf} \\
63 & value depends only on which wire joins which & slide, bend, stretch
freely --- only the connectivity counts \\
64 & green\((1\!\to\!2)\) control, red\((2\!\to\!1)\) target, one shared
wire; bare \(=\mathrm{CNOT}/\sqrt2\) & a green dot and a red dot joined
by one cord --- the CNOT \\
\end{longtable}

\endgroup

From Chapter 8 onward, picture equalities hold \textbf{up to a non-zero
scalar} unless the scalar is written in. Fusion (57) and colour change
(58) are exact; copy carries \(\sqrt2\), bialgebra \(\sqrt2\), Hopf
\(2\).

\begin{center}\rule{0.5\linewidth}{0.5pt}\end{center}

\hypertarget{the-theorem-twice-7}{%
\section{§8.7 The Theorem, Twice}\label{the-theorem-twice-7}}

Here is the chapter's headline, chosen because it is small enough to
prove in full and large enough to matter: it is the identity every
student of quantum computing meets, that a swap can be built from
nothing but CNOTs.

\begin{quote}
\textbf{Theorem 8.8 (three CNOTs make a swap).} Let \(\mathrm{CNOT}\)
have its control on the top wire and target on the bottom, and
\(\mathrm{CNOT}'\) the reverse --- control bottom, target top. Then
\[\mathrm{CNOT}\cdot\mathrm{CNOT}'\cdot\mathrm{CNOT} \;=\; \mathrm{SWAP},\]
which exchanges the two wires. On the blackboard this is \textbf{exact}.
In the pictures the three spider-CNOTs cook down to a bare crossing
\textbf{up to the scalar \(\tfrac1{2\sqrt2}\)}.
\end{quote}

\hypertarget{diracs-blackboard-7}{%
\subsection{Dirac's Blackboard}\label{diracs-blackboard-7}}

Write the three gates as \(4\times4\) permutation matrices in the basis
\(\ket{00},\ket{01},\ket{10},\ket{11}\). The CNOT with control on top
flips the bottom bit when the top is \(1\); its reverse flips the top
bit when the bottom is \(1\):

\[\begin{aligned}
\mathrm{CNOT}&=\begin{pmatrix}1&0&0&0\\0&1&0&0\\0&0&0&1\\0&0&1&0\end{pmatrix},
&\quad
\mathrm{CNOT}'&=\begin{pmatrix}1&0&0&0\\0&0&0&1\\0&0&1&0\\0&1&0&0\end{pmatrix}, \\[0.6em]
\mathrm{SWAP}&=\begin{pmatrix}1&0&0&0\\0&0&1&0\\0&1&0&0\\0&0&0&1\end{pmatrix}.
\end{aligned}\]

Follow a basis vector through, right to left. Start with \(\ket{10}\).
The first \(\mathrm{CNOT}\) (control top \(=1\)) flips the bottom:
\(\ket{11}\). Then \(\mathrm{CNOT}'\) (control bottom \(=1\)) flips the
top: \(\ket{01}\). Then \(\mathrm{CNOT}\) (control top \(=0\)) does
nothing: \(\ket{01}\). So \(\ket{10}\mapsto\ket{01}\). The same march
sends \(\ket{01}\mapsto\ket{10}\), and fixes \(\ket{00}\) and
\(\ket{11}\). That is the swap, exactly --- three sparse matrices
multiplied in a line, no scalar to explain. See
\texttt{Code/Ch8/SwapFromCNOTs.py}.

\hypertarget{coeckes-sketchbook-2}{%
\subsection{Coecke's Sketchbook}\label{coeckes-sketchbook-2}}

Each CNOT is a green--red pair (row 64). Draw the three in a row: on the
top wire green, red, green; on the bottom wire red, green, red; three
vertical cords joining them.

\[\begin{tikzpicture}[sd,
  zsp/.style={circle,draw,line width=0.9pt,fill=green!55,minimum size=1.1em,inner sep=0pt},
  xsp/.style={circle,draw,line width=0.9pt,fill=red!75,minimum size=1.1em,inner sep=0pt}]
  \draw[wire] (-0.8,0.9)--(5.0,0.9);
  \draw[wire] (-0.8,-0.9)--(5.0,-0.9);
  \draw[wire] (0.6,0.9)--(0.6,-0.9);
  \draw[wire] (2.1,0.9)--(2.1,-0.9);
  \draw[wire] (3.6,0.9)--(3.6,-0.9);
  \node[zsp] at (0.6,0.9){}; \node[xsp] at (0.6,-0.9){};
  \node[xsp] at (2.1,0.9){}; \node[zsp] at (2.1,-0.9){};
  \node[zsp] at (3.6,0.9){}; \node[xsp] at (3.6,-0.9){};
\end{tikzpicture}\]

Now cook it down, each move named. Call the top-wire dots, left to
right, \(g_1\) (green), \(r_2\) (red), \(g_3\) (green), and the
bottom-wire dots \(r_1\) (red), \(g_2\) (green), \(r_3\) (red); the
three vertical cords are \(g_1r_1\), \(r_2g_2\) and \(g_3r_3\).

\textbf{Step 1 --- only connectivity (§8.5).} Look at the four dots of
the first two CNOTs --- the greens \(g_1,g_2\) and the reds \(r_1,r_2\).
Each green is joined to each red: \(g_1\) to \(r_1\) (its cord) and to
\(r_2\) (along the top wire); \(g_2\) to \(r_2\) (its cord) and to
\(r_1\) (along the bottom wire). That is a complete cross between the
colours --- a bialgebra cradle standing on its side. Redraw it as one
without moving an endpoint.

\textbf{Step 2 --- bialgebra (Theorem 8.5).} Collapse that cross to a
single green--red pair. The rule sends the two greens' external legs
onto the new red and the two reds' external legs onto the new green ---
not the other way about. Outside the cradle, the greens \(g_1,g_2\)
carry the top boundary input (on \(g_1\)) and the wire running on toward
the third CNOT's target \(r_3\) (on \(g_2\)); the reds \(r_1,r_2\) carry
the bottom boundary input (on \(r_1\)) and the wire running on toward
the third CNOT's control \(g_3\) (on \(r_2\)). So a \textbf{new red}
gathers the top boundary input together with the wire toward \(r_3\), a
\textbf{new green} gathers the bottom boundary input together with the
wire toward \(g_3\), and a single cord joins the new red to the new
green.

\textbf{Step 3 --- fusion (Theorem 8.1), on the green.} The new green
carries a leg running toward \(g_3\), the third CNOT's green; two greens
on a shared wire, so they fuse into a single green dot (phases add, both
\(0\)). This green is now joined to the third CNOT's red \(r_3\) by one
wire --- the surviving cord \(g_3r_3\) --- and to the new red by
another; but \(r_3\) and the new red are still two \emph{different} red
dots, so there is as yet no double edge.

\textbf{Step 4 --- fusion (Theorem 8.1), on the red.} The new red
carries a leg running toward \(r_3\); two reds on a shared wire, so they
fuse into a single red dot (phases add, both \(0\)). Now a single green
and a single red remain, and they are joined by \emph{two} separate
wires --- the cord left by the bialgebra pair and the third CNOT's own
cord \(g_3r_3\).

\textbf{Step 5 --- Hopf (Theorem 8.6).} Two wires between a green and a
red are the same as none. The link falls away, leaving only phase-free
spiders with one wire in and one out.

\textbf{Step 6 --- a phase-free \(1\!\to\!1\) dot is a bare wire
(\(Z(1,1,0)=I\), §8.1).} Erase those dots and follow the surviving wires
from boundary to boundary: the top input emerges at the bottom and the
bottom input at the top --- a bare crossing.

\[\begin{tikzpicture}[sd]
  \draw[wire] (0,0.8) to[out=0,in=180] (2.4,-0.8);
  \draw[wire] (0,-0.8) to[out=0,in=180] (2.4,0.8);
\end{tikzpicture}
\;\;=\;\;\text{SWAP}.\]

The bare crossing is SWAP exactly; the scalar rides on the reduction
that brought us here. The three spider-CNOTs cook down to this crossing
up to \(\bigl(\tfrac1{\sqrt2}\bigr)^{3}=\tfrac1{2\sqrt2}\), one factor
of \(\tfrac1{\sqrt2}\) for each spider-CNOT (row 64). See
\texttt{Code/Ch8/SwapFromCNOTs.py}, which assembles the three
spider-CNOTs, reduces the diagram by the rules, and confirms it equals
the bare swap up to exactly \(\tfrac1{2\sqrt2}\).

\hypertarget{the-verdict-7}{%
\subsection{The verdict}\label{the-verdict-7}}

Be honest: for this theorem the blackboard wins. Three sparse matrices
multiply to the swap in a single line, exactly, with no scalar to
apologise for --- shorter and cleaner than the six-step rewrite. What
the picture buys is not brevity but \emph{understanding}. It shows that
the swap is not a coincidence of three permutations but a consequence of
complementarity: the swap appears the instant the two colours are made
to reorganise each other by the bialgebra law. And it shows the result
is robust --- the same cradle collapses whatever gates of this shape are
strung, which a specific matrix product cannot suggest. The algebra
tells you \emph{that}; the picture tells you \emph{why}, and carries a
\(\tfrac1{2\sqrt2}\) for the privilege.

\begin{center}\rule{0.5\linewidth}{0.5pt}\end{center}

\hypertarget{the-octet}{%
\section{§8.8 The Octet}\label{the-octet}}

Here is the master diagram of the chapter. Count its cords first: two
coming in on the left, four crossing in the middle, two going out on the
right --- \textbf{eight}.

\begin{center}
\begin{tikzpicture}[sd,
  zsp/.style={circle,draw,line width=0.9pt,fill=green!55,minimum size=1.35em,inner sep=0pt},
  xsp/.style={circle,draw,line width=0.9pt,fill=red!75,minimum size=1.35em,inner sep=0pt}]
  \coordinate (gt) at (0,1.3); \coordinate (gb) at (0,-1.3);
  \coordinate (rt) at (3.4,1.3); \coordinate (rb) at (3.4,-1.3);
  % 1,2: feed cords (left boundary to greens)
  \draw[wire] (-1.8,1.3)--(gt);
  \draw[wire] (-1.8,-1.3)--(gb);
  % 3,4,5,6: the K_{2,2} cradle
  \draw[wire] (gt)--(rt);
  \draw[wire] (gt)--(rb);
  \draw[wire] (gb)--(rt);
  \draw[wire] (gb)--(rb);
  % 7,8: draw cords (reds to right boundary)
  \draw[wire] (rt)--(5.2,1.3);
  \draw[wire] (rb)--(5.2,-1.3);
  \node[zsp] at (gt){}; \node[zsp] at (gb){};
  \node[xsp] at (rt){}; \node[xsp] at (rb){};
  \node[sdlabel,anchor=east] at (-1.9,1.3){in};
  \node[sdlabel,anchor=east] at (-1.9,-1.3){in};
  \node[sdlabel,anchor=west] at (5.3,1.3){out};
  \node[sdlabel,anchor=west] at (5.3,-1.3){out};
\end{tikzpicture}
\end{center}

Two green dots on the left, two red on the right, and between them a
complete cross --- every green tied to every red. It is the bialgebra
cradle of §8.4 drawn at full size, and it is called \textbf{The Octet}.

\begin{quote}
\textbf{Theorem 8.9 (the Octet).} The Octet is a genuine eight-wire
diagram --- two feed cords, four cradle cords, two draw cords --- and it
evaluates to the parity map
\[\tfrac12\Bigl(\ket{00}\bigl(\bra{00}+\bra{11}\bigr)+\ket{11}\bigl(\bra{01}+\bra{10}\bigr)\Bigr),\]
which reads the parity of its two inputs and writes it onto both
outputs. Up to the scalar \(\sqrt2\), this is a single green--red pair
--- the bialgebra rule of §8.4, drawn once in the large.
\end{quote}

\emph{Proof.} The eight cords are enumerated above. The companion file
\texttt{Code/Ch8/Octet.py} builds the diagram and tallies its edges to
exactly eight; it then derives the collapse to the single pair by the
named bialgebra rule and confirms the parity map up to \(\sqrt2\). The
evaluation is the blackboard multiplication of Theorem 8.5.
\(\;\blacksquare\)

And the Tortoise, reshaping the cradle into a different-looking tangle
that presses down to the very same pair, glances at the Escher and says
it: the shape flows on, but the thing stays the same. Only the
connectivity matters.

\begin{quote}
\textbf{Coecke's Sketchbook.}

\emph{A confession about the numbering.} The reader who counts the cords
of the Octet finds eight --- and the chapter is the eighth. That is not
a coincidence, and it has not been one since the beginning. Chapter 1
turned on a single wire; Chapter 2 on two; the six music boxes of
Chapter 6 were six wires bent into cups and caps; the trio of Chapter 7
met at a node whose legs, with the tally, came to seven. Every chapter
of this book has been built on its own number of wires, and its number
\emph{is} that count. We hinted once, in a footnote in Chapter 6, and
said we would not explain. This is the explanation. If you wish to see
where it starts, turn back to the very beginning --- before the first
chapter, to the Prelude, a page with no wires drawn on it at all: the
empty diagram, which counts zero. We will not mention the numbering
again; from here you can carry it yourself.
\end{quote}

\begin{quote}
\textbf{Feynman's Notebook.}

\emph{Whose calculus.} The diagrammatic reading of quantum processes ---
the bent wire as a Bell state, yanking, the whole apparatus of a
\(\dagger\)-compact category --- is due to Samson Abramsky and Bob
Coecke, \emph{A categorical semantics of quantum protocols}, presented
at the 19th IEEE Symposium on Logic in Computer Science (LiCS'04) in
2004 (arXiv:quant-ph/0402130). The two colours of dot --- the
ZX-calculus proper, with fusion, colour change and the bialgebra --- are
Coecke and Ross Duncan, \emph{Interacting quantum observables}, first
given at ICALP 2008 and published in full in 2009 (arXiv:0906.4725).
Completeness for the stabilizer fragment was proved by Miriam Backens in
2013 (arXiv:1307.7025). All of this rests on, but is distinct from,
Roger Penrose's 1971 diagrammatic notation for tensors, which drew the
wires but did not paint them.
\end{quote}

\begin{center}\rule{0.5\linewidth}{0.5pt}\end{center}

\hypertarget{now-the-pictures-can-prove}{%
\section{§8.9 Now the Pictures Can
Prove}\label{now-the-pictures-can-prove}}

Look back at what changed. For seven chapters a diagram was a way of
\emph{displaying} a truth the algebra had established. In this one the
diagram became a way of \emph{establishing} one. We gave the Chapter 7
dot a colour, a second colour, a turn and a family of legs; we found the
rules by which the two colours push through and reorganise one another
--- fusion, colour change, copy, the bialgebra and its twin the Hopf
law; and with those rules alone we cooked three CNOTs down to a
crossing, never once returning to a matrix to check. The calculus is
Coecke's; the promise it makes --- universality, and completeness on the
stabilizer fragment --- is that this was no lucky example.

Everything we did, though, we did with the lights off. The koekpan
pressed its cradles by feel; the cords were tied and crossed and
counted, but never once \emph{watched} to see which way a cord truly
ran. We drew states and effects and post-selected caps, and each time
quietly kept only the outcome we wanted. That cannot last. A real device
does not let you keep only the answer you like; it forces one of several
on you, at random, and destroys the rest --- and the act of looking is
itself a thing that must be drawn. To draw it we will need a new kind of
wire, and a partner in the corner who is neither asleep nor awake until
someone looks. Tomorrow we look.

\hypertarget{lullaby-for-a-cat}{%
\chapter{Lullaby for a Cat}\label{lullaby-for-a-cat}}

\begin{center}\rule{0.5\linewidth}{0.5pt}\end{center}

\emph{猫の目や}\\
\emph{寝るも覚むるも}\\
\emph{冬灯}

\emph{neko no me ya}\\
\emph{neru mo samuru mo}\\
\emph{fuyu tomoshi}

\emph{A cat, one eye open ---}\\
\emph{asleep and awake at once,}\\
\emph{the winter lamp}

\begin{center}\rule{0.5\linewidth}{0.5pt}\end{center}

\emph{{[}The Crab's parlour, the short winter afternoon already on the
turn. Achilles comes in out of the snow, stamping his boots; the
Tortoise is settled by the lamp; and in the corner, on the sill,
something is curled that is not quite asleep and not quite awake.{]}}

\textbf{Achilles:} I ran the last of the lane --- it's the kind of cold
that makes your teeth ring. You said today. Yesterday you said,
\emph{tomorrow we look.} Well. It's tomorrow.

\textbf{Crab:} It is, and I have been up since it was dark, nerving
myself to it. Come to the fire; there's tea. I have kept a thing in a
drawer eleven years that I have never once dared use in company. Today I
use it.

\textbf{Tortoise:} He has said \emph{today I use it} three times since I
arrived. Each time he shuts the drawer again.

\textbf{Crab:} This time it stays open. Sit, Achilles --- thaw first.

\emph{{[}The fire has been built high against the cold, and it spits and
shifts as the fresh wood catches; the snow that came in on Achilles's
boots goes dark, then to water, then to nothing on the hearth-stone.
Beyond the one tall window the afternoon has gone the colour of pewter
and is failing early, as it does in the deep of the year, so that the
newly-lit lamp already does the greater part of the work. In the far
corner, on the sill, something keeps its own counsel, curled ---
settling deeper, it may be, or it may be about to stretch.{]}}

\textbf{Achilles:} \emph{{[}warming his hands, then nodding at the
corner{]}} And who's your friend on the sill? Settling down for the
night, or waking up into it?

\textbf{Cat:} \emph{{[}without troubling to open more than one eye{]}} A
little of each. More of the one than the other --- and the other than
that --- depending how you ask.

\textbf{Achilles:} That is no sort of answer.

\textbf{Tortoise:} It is the exact size of your question. Ask it
something baggy, it hands you something baggy back. Drink your tea.
You'll grow used to the manners.

\emph{{[}Achilles takes the offered cup, folds both hands round it, and
lets the fire do its work while the light in the window fails by the
minute.{]}}

\begin{center}\rule{0.5\linewidth}{0.5pt}\end{center}

\hypertarget{first-challenge-the-look-that-forces}{%
\section{\texorpdfstring{First Challenge --- \emph{The Look That
Forces}}{First Challenge --- The Look That Forces}}\label{first-challenge-the-look-that-forces}}

\emph{{[}The Crab's parlour in the deep of winter, the short afternoon
sinking toward an early dusk, and the lamp lit already to meet it. It is
a room that shuts and covers: winter shutters folded back in their
reveals, waiting to be drawn; a sideboard of lidded jars and stoppered
bottles; cloths laid over whatever the cold would spoil; a drawer or two
that keep their contents to themselves. The one tall window has not yet
decided what it is --- half a view still of the white lane and the snow
sifting down, and half a ghost of the warm room laid over it, the two
neither one thing nor the other; and at its cold lower corners the glass
has begun to fog with the breath of the room, a soft grey bloom that has
not settled what it will be. On the scullery wall, where the lamplight
reaches only at its edge, hangs a print the Crab has kept these many
years --- Escher's \textup{Eye}, a single vast eye drawn so close that
every lash stands up like a rafter, and far down in the dark of its
pupil something small and steady looking back out. It was on that wall
before any of them came in from the cold, and it will be there after
they have gone.{]}}

\textbf{Achilles:} Then let me have it. All this while I have handled
things but never once \emph{looked} at one while it was making up its
mind. Today I look. What do I look at?

\textbf{Crab:} \emph{{[}bringing a shallow drawer to the table{]}}
These. Little stones, and they come in pairs --- each pair cut of one
piece, so the two are never quite two apart. A restless stone hasn't
settled what it is; it is both its notes at once, the way your friend on
the sill is neither asleep nor awake. \emph{{[}He sets one small stone
apart.{]}} But a look spends a stone. Once you have looked, it is
decided, and it stays decided --- so each showing wants a fresh pair,
and I have a drawerful.

\textbf{Crab:} And this is how you look. \emph{{[}He sets a squat brass
instrument on the table: an eyepiece at the top, a fine needle along a
scale on its face, a stoppered flask sunk in the base.{]}} The viewer.
From the same bench in Delft as my music boxes --- the daughter grinds
glass now, as well as cutting reeds. You look \emph{through} it, at the
stone, never \emph{at} it: put your eye to its works to catch how it
does the trick and the needle jams and it is spoilt. I have never once
looked inside, and it costs me. Look through it at a restless stone and
it forces one of the two notes on you --- which one, nobody can say
beforehand --- swings the needle to that side, and breathes once into
the flask below.

\textbf{Crab:} The answer it forces is cut here. \emph{{[}a slate, and a
stylus{]}} One notch a look, and there it stays: a plain record, settled
and done. And every look I have ever taken is in this. \emph{{[}a fat
ledger{]}} Eleven years of them.

\emph{{[}The viewer sits square in the lamplight, all brass and ground
glass, as finely made as anything from that Delft bench --- and, like
the best of that bench's work, it does not care to be studied: left
alone it keeps up the faintest tick somewhere inside its brass, a small
live sound at the very edge of hearing, that falls silent the moment
anyone leans in close to catch where it comes from, and takes itself up
again when they sit back. The Crab draws a chair up for Achilles and
another for himself, and keeps his own gaze at a careful middle
distance. Achilles settles in and bends to the eyepiece.{]}}

\textbf{Tortoise:} Go on, Achilles. Look at it.

\emph{{[}The stone lies restless --- neither note, both notes, trembling
at nothing. Achilles sets his eye to the viewer and looks through it at
the stone. The needle, which has hovered at the middle of its scale,
jumps --- and swings hard over to the high stop, and holds there; it
might as well have gone the other way, and no one in the room could have
said beforehand which. A single notch is scratched on the slate. And the
stone that trembled a moment ago lies quite still now: it has an answer,
and it is spent.{]}}

\textbf{Achilles:} High. \emph{{[}He looks again.{]}} High again --- and
again, high. It won't shift.

\textbf{Crab:} It won't. You have spent it. Look till the small hours
and it says high; but the both-at-once you looked \emph{through} is
gone. You didn't find it high---

\textbf{Achilles:} \ldots Oh. The look didn't read the stone, it
\emph{made} it. The stone had no answer until I looked; my looking is
what settled it. \emph{{[}He takes a fresh stone, looks; a notch is cut.
Another fresh one, looks; another notch.{]}} High. Low. Low. High. Low.
A fresh stone each time, and a row of notches with no more pattern in it
than a tossed coin.

\emph{{[}At the window one of the fogged patches, chilled past bearing,
springs of a sudden into frost --- a whole white fern flung out across
the glass in the space of a breath, from no cause that anyone could have
named beforehand, and fixed there now for good. A moment since it was
undecided grey; it will not be grey again tonight.{]}}

\textbf{Tortoise:} One look, one answer, and a restless thing is put to
rest. That is the whole of it.

\textbf{Achilles:} The whole of it --- except for what nags. Suppose I
look, the stone decides, and I \emph{don't} read the slate. I have spent
the stone; it has its answer; but I don't know it. It isn't the restless
both-at-once --- that's gone. And it isn't a plain high or low --- I
never read it. So what is a stone that has been looked at, and never
read?

\begin{center}\rule{0.5\linewidth}{0.5pt}\end{center}

\hypertarget{second-challenge-a-stone-thats-been-looked-at}{%
\section{\texorpdfstring{Second Challenge --- \emph{A Stone That's Been
Looked
At}}{Second Challenge --- A Stone That's Been Looked At}}\label{second-challenge-a-stone-thats-been-looked-at}}

\textbf{Tortoise:} Then make one, and we'll look at it --- the stone,
mind, not the slate. Mr Crab, cover the dial and the slate before he
looks. He is to spend a stone without learning its answer.

\textbf{Crab:} \emph{{[}draping one cloth over the needle's dial,
another over the slate{]}} Covered. Look through, Achilles --- you'll
see nothing come of it.

\emph{{[}The stone lies restless, both at once. Achilles sets his eye to
the viewer and looks through. Behind the cloth the needle swings and a
notch is cut --- but nobody sees to which side it swung, nobody reads
the mark. When he lifts his eye the stone lies still: spent, settled,
decided. Only \emph{which way} it decided, no one in the room knows. It
is a stone that has been looked at and never read.{]}}

\textbf{Achilles:} So it has an answer, and I can't get at it. It isn't
restless any longer --- but it isn't a plain high or low to me either.
What have I got?

\textbf{Tortoise:} A black glass to show it on, is what we want. Mr Crab
--- the slab?

\textbf{Crab:} \emph{{[}drawing a polished black slab into the
lamplight{]}} My scrying-slab; I keep it for its looks. Set anything on
it and it throws back a reflection so clean you would swear there were
two of the thing. \emph{{[}wry{]}} I begin to suspect I shall regret
lending it to a demonstration.

\emph{{[}The window has given up the last of its plain view. Where an
hour since there was a lane, there is now a muddle of both worlds at
once --- the snow-light and the lamp-flame, the frost on the pane and
the reflected room, laid one across another until a stranger looking in
could not have told which things stood in the parlour and which only
seemed to. In the glass the lamp burns beside its own reflected self;
and for a moment Achilles, glancing up, cannot say which is the flame
and which the window's.{]}}

\emph{{[}The Tortoise sets a fresh, restless stone at the edge of the
slab. It trembles at nothing --- both notes at once --- and in the black
its reflection trembles with it, the two locked step for step, so close
you could not slip a hair between their timing. She rocks the stone a
hair with one fingertip; the reflection rocks the same hair, the same
instant, the same distance. A thing, and its reflection, moving as one
thing.{]}}

\textbf{Tortoise:} A thing and its reflection. While the thing is
restless, the two are tight --- one motion, twice over.

\emph{{[}She takes the spent-and-unread stone and sets it where the
restless one stood. It lies still. But when she rocks it, its reflection
answers slack: a beat late, drifting a little wide, no longer keeping
its partner's exact time. The same stone, the same twin thrown back in
the glass; only the coupling between them has gone loose.{]}}

\textbf{Tortoise:} And a thing that has been looked at --- still a thing
and its reflection, but slack. That looseness is the whole of what your
look left you: not a note, not the other note, but a thing gone loose
from its own mirror.

\textbf{Achilles:} So my looked-at stone is --- half the time high, half
the time low, and I simply can't say which. A weighted list, not one
note and not the other. \emph{{[}working it through{]}} And the restless
one is the tight end of that: one whole possibility, its reflection in
perfect step. Every looseness between the two is a look taken and not
read.

\emph{{[}Beyond the glass the snow is at its slow work still. The flakes
coming down now are sharp six-pointed stars, each one particular and cut
clean; where they have lain a while on the sill they have softened and
slumped, arm run into arm, until they are only a bank of rounded grains
--- the same snow, every edge of it gone, keeping nothing of the figure
it fell in.{]}}

\textbf{Crab:} And each of these throws its own reflection, and my slab
is filling up with twins. Might we clear a few off before I lose which
is a stone and which is the ghost of one?

\textbf{Achilles:} Clear them, by all means --- but that is my next
trouble exactly. If every stone is a thing and its reflection now, and a
pair is two of them\ldots{} suppose I care about only \emph{one} half of
a shared pair. The far half's business is the far half's. How do I get
an honest account of just my half, and cover the rest?

\begin{center}\rule{0.5\linewidth}{0.5pt}\end{center}

\hypertarget{third-challenge-cover-the-other-half}{%
\section{\texorpdfstring{Third Challenge --- \emph{Cover the Other
Half}}{Third Challenge --- Cover the Other Half}}\label{third-challenge-cover-the-other-half}}

\emph{{[}The dusk has settled the quarrel the window was having with
itself: it is nearly all mirror now, a clean dark sheet with the whole
room stood complete and orderly a pane's-depth inside it. The Crab,
passing, feeds the fire a single log and sets the guard back; the tea
has gone tepid in the pot and no one minds. Outside, past the reflected
room, the snow has thickened to a steady, straight-down fall.{]}}

\textbf{Tortoise:} You cover a half by folding it into its own
reflection --- a new move, and a plain one. To cover a thing is not to
look at it and not read it; it is to refuse it any account at all, to
let it go. On the slab, letting a thing go is folding it down to meet
its reflection: they touch, and cancel, and there is an end of it.

\emph{{[}Here is a plain pair --- two stones with nothing between them,
cut from different pieces --- each with its reflection standing beneath
it in the black. The near stone is Achilles's; the far one is no
business of his. The Tortoise tips the near stone down toward the glass;
its reflection rises to meet it, mirror-fashion; they touch at the
surface --- and where thing met twin there is now neither, only a bare
patch of black slab. The near stone is not spent, not looked at, not
read: simply gone, folded into its own mirror. The far stone, and its
reflection, stand exactly where they stood, not stirred by a hair.{]}}

\emph{{[}On the roof of the house across the lane the afternoon's snow
has crept at last to its own weight, and the near slope lets go all
together --- a whole clean sheet of it sliding off into the dark below
and leaving the slates bare --- while the ridge above and the far slope
hold their white exactly as it lay, not a crystal shifted. What went,
went wholly; what stayed, stayed untouched.{]}}

\textbf{Tortoise:} That fold is ``cover the other half.'' I covered the
near one --- mine --- and the far one never felt it. Cover an unrelated
half and you lose nothing you wanted, and change nothing you kept.

\textbf{Achilles:} Because the two had nothing to do with each other.
\emph{{[}nods{]}} But that is the dull case. Do it on a pair truly
\emph{of one piece} --- the way your stones are cut --- and cover the
\emph{near} half, the one in my hand. What is left of the far one, that
I never touched?

\begin{center}\rule{0.5\linewidth}{0.5pt}\end{center}

\hypertarget{fourth-challenge-what-the-twin-cannot-tell}{%
\section{\texorpdfstring{Fourth Challenge --- \emph{What the Twin Cannot
Tell}}{Fourth Challenge --- What the Twin Cannot Tell}}\label{fourth-challenge-what-the-twin-cannot-tell}}

\textbf{Achilles:} \emph{{[}to the Cat, half to himself{]}} You on the
sill --- if I work this cleverly enough, will my message get across?

\textbf{Cat:} \emph{{[}one eye barely open{]}} It will, and it won't.
And the won't has the better of it.

\textbf{Tortoise:} Loose question, loose answer --- leave it loose while
it's on the sill. A plain question is a look, remember; and we are not
looking \emph{there} yet.

\emph{{[}The window is almost wholly a mirror now, the whole lit room
hung in it --- the lamp, the bowed heads, the table and its business ---
with only a last dull grey low down where the lane had been, and the
snow beyond gone to a blue-dark nothing. The Crab turns the lamp up a
little against the dark; the fire has burned down to a steady,
companionable red. On the sill the curled shape is a warm hollow and a
smooth one at once, an eye's-worth of lamplight caught in it and the
rest given over to the dark --- settling, or waking, there is no telling
which.{]}}

\textbf{Crab:} A pair of one piece, then. \emph{{[}He sets one fresh
stone by Achilles's hand and carries its partner to the far end of the
table, under a cloth.{]}} The near one is yours to bully. The far one
goes under here, and nobody looks at it --- not you, not me. Those are
the rules.

\textbf{Achilles:} Right. Now I send a message across this table with no
more than my own half. Watch. If I \emph{look} at my half, the pair is
cut of one piece, so the far half must fall the same way --- a twitch I
could send. If I \emph{fold} my half away instead, that's a different
twitch. If I \emph{leave it be}, a third. Three things I can do, three
states of the far stone, and a code between them. \emph{{[}He reaches
for the cloth.{]}}

\textbf{Crab:} You do not look at the far one. Rules.

\textbf{Tortoise:} You need not, and you must not. That is what the
ledger is for. Every one of those things you might do to a near half, I
have done a hundred times over the years, on a hundred pairs, and
written down what the far halves said. Read them off --- by what was
\emph{done}, not by what fell.

\textbf{Achilles:} \emph{{[}at the ledger{]}} Then --- every page where
the near half was looked at, whatever it came up. The far halves read:
high, low, low, high, high, low. A coin. \emph{{[}turning pages{]}}
Every page where the near half was folded away and let go. The far
halves: low, high, high, low, high. The same coin. \emph{{[}more
pages{]}} Every page where nothing at all was done to the near half. A
coin again. Three different things done here, and the far half a fair
coin under every one of them. No difference for her to read.

\emph{{[}The ledger's pages have gone soft as cloth at the corners,
thumbed through eleven winters of such evenings. At the table's far end
the kept stone lies dumb under its cloth where the Crab set it, a small
covered weight that no one has touched and no one means to. Achilles
reads on, chin propped in one hand, warming to the chase.{]}}

\textbf{Achilles:} But wait --- let me \emph{sort} them. Only the pages
where my near half came up \emph{high}. Surely then the far halves---

\textbf{Crab:} ---all read high, yes. Cut of one piece; they always
agree. But you did not \emph{pick} high, Achilles. The viewer did, at
random --- you told me so yourself an hour ago. And your partner across
the room has no way on earth to know which pages are hers, the high pile
or the low, unless you write and tell her by post. Sorting a thing you
could not choose is no message.

\textbf{Achilles:} \ldots Oh. There's nothing I can do here that reaches
there. The halves agree --- but the agreeing is no use to me, because I
cannot steer it and she cannot see it. Whatever I \emph{do}, her half is
the same grey coin.

\emph{{[}Outside, the wind that has been gathering all afternoon finds
its voice, and flings a fistful of snow against the panes, and then
another, and leans on the old house until the timbers creak. Inside, the
lamp-flame stands straight and unwavering and does not so much as lean;
the surface of the tea lies still in the cup; nothing in the sealed warm
room takes the least notice of all that battering at the glass.{]}}

\textbf{Tortoise:} \emph{{[}a length of string from her shell{]}} One
picture, and you'll see \emph{why} it can't. This is a string --- the
plainest thing I know that goes round a bend and stays all one piece,
which is just what a pair cut of one piece is. Let me lay your pair on
the slab, halves and reflections and all.

\emph{{[}On the black she lays the string from the far stone across to
the near stone --- that is the pair, joined of one piece --- and in the
glass its own reflection runs beneath it, from the far stone's twin
across to the near stone's twin. Then she folds the near stone down into
its reflection, the same cover as before; and where the near stone meets
its twin, the string above joins the string below. Far stone, across to
near, round the fold, and back beneath to the far stone's twin: one
string and its mirror, with the near half pinched away in the
middle.{]}}

\emph{{[}She sets one finger on the far stone to hold it, and with her
other hand draws the slack out of the middle --- toward the far stone.
The two bends slide together; the fold between them closes; the near
half and its twin are drawn in and swallowed; and the string plunges
straight down into the glass at the far stone's own foot, to meet its
own reflection there. What was far-across-near-fold-near-back is now one
clean bend: the far stone, joined to nothing but its own reflection in
the black. Beside it she chalks the one word: \emph{half}.{]}}

\textbf{Achilles:} The far half, bent round to its own reflection, and a
half beside it. That is your slack stone from before --- the thing loose
from its own mirror, the flattest grey there is. A fair coin, every way
of asking at once.

\textbf{Tortoise:} One pull. And whatever you had done to the near half
--- any box you please on that middle leg --- the fold swallows it
before the pull even begins. It never reaches the far stone. She was
always going to come out that same grey.

\textbf{Achilles:} But your fold is \emph{clean}. A real look isn't ---
mine swung a needle, breathed into a flask, left a leak in the room.
Fold a clean half away and nothing crosses, granted. But a real look
drags all that plumbing behind it. Does the leak get out? Could the
\emph{leak} carry across what I couldn't?

\begin{center}\rule{0.5\linewidth}{0.5pt}\end{center}

\hypertarget{fifth-challenge-a-look-with-the-lights-on}{%
\section{\texorpdfstring{Fifth Challenge --- \emph{A Look with the
Lights
On}}{Fifth Challenge --- A Look with the Lights On}}\label{fifth-challenge-a-look-with-the-lights-on}}

\textbf{Crab:} Then we lay it open --- the look itself, mind, not the
trick that does it. \emph{{[}He unclasps the viewer's case and folds it
back, keeping his eye carefully off the works.{]}} Here is where a real
look happens. The near stone goes in here. The needle here, that swings
to the answer. And here, sunk in the base, the flask --- the vent; every
look breathes once into it, and then it is stoppered and set by. Three
things, bound into one by the works in the instant of looking.
\emph{{[}aching{]}} I could tell you \emph{how} they are bound, if I put
my eye to the works. I never have, and I never shall; that is what this
device costs me, and I pay it gladly --- or nearly gladly.

\emph{{[}The opened viewer lies on the table like a watch with its back
off, laid open to the lamp --- and it is, uncovered so, a little less of
a marvel and a little more of an ordinary clever thing, as such things
always are once one can see round behind them. The Crab keeps his eyes
resolutely on the far wall while his claws do the work, the way a man
carries a brimming cup without watching his own feet. It has cost him
something to lay it bare, and he does not pretend that it hasn't.{]}}

\emph{{[}On the black slab he stands them out for the look, each with
its reflection beneath it: the near stone and its twin, the needle and
its twin, the flask and its twin. The far stone keeps to its cloth at
the end of the table, its own reflection with it. The works will bind
the near three; the far one they never touch. Off the black, beside the
slab, the Crab lays the slate.{]}}

\emph{{[}The needle and its reflection stand together, two pointers,
both quivering toward an answer. The Tortoise brings them together and
pins the needle down onto its own reflection, at the one point where the
two meet: where there were two quivering pointers there is now a single
still one. From that pinning a notch is cut --- but on the slate, which
she keeps off the black glass, so that this one mark alone throws no
twin of itself. Two pointers in; one settled mark out; and the mark
stands single, unpartnered, as the first notch of the evening did.{]}}

\emph{{[}She folds the flask down into its reflection --- the same
cancelling fold as the stones --- and it is gone: breathed into,
stoppered, let go.{]}}

\textbf{Tortoise:} Now the leak. The works tie the near stone, the
needle and the flask into one --- a box across the three of them, and
its reflection a box across the three twins. Then we let them all go:
the flask folded away, the near stone folded away, the needle read to
its one notch and finished --- every output of the works folded down
into its own reflection. Watch what that does to the far half.

\emph{{[}The string runs from the far stone into the works, through the
box, and out the far side; and every leg it comes out on is folded into
its reflection. The Tortoise sets a finger on the far stone and draws
the slack in toward it. The box, coming out, meets its own reflection
under the fold --- arrives upon itself, and cancels --- and the whole
apparatus is swallowed as the fold closes. The string plunges straight
into the glass at the far stone's foot, to its own reflection: one clean
bend, the plain half beside it, exactly as for the clean fold. The
needle's notch stands on the slate; the flask sits stoppered by the
base. But the far stone is the same flat coin it was before --- the
works reached it not at all.{]}}

\emph{{[}Up the chimney the fire is pouring out everything it has ---
smoke, and the afternoon's stored heat, and a low steady roar in the
flue --- away into the freezing dark above the roofs. And for all of it
the night is not altered by a breath: the snow lies along the roofs the
length of the lane no thicker and no thinner for the smoke going up past
it, and the dark takes the roar and gives back nothing. All that
pouring-out, and none of it reaches across.{]}}

\textbf{Achilles:} Nothing. All that plumbing --- needle, flask, leak
--- and the far half is the same grey coin.

\textbf{Tortoise:} Because discarding after the works is discarding
straight off. Let a thing go after all its machinery, or never build the
machinery at all --- it comes to one and the same. The works meet their
own reflection under the fold and cancel; the far half never hears them.
A real look, laid bare, leaks nothing across.

\emph{{[}As she speaks the last grey light gives out in the window, and
the room drops to the lamp alone. In that sudden evenness the print on
the scullery wall seems to stand a half-step out from its wall, the
great eye darkening as the shadows gather into it.{]}}

\emph{{[}The Tortoise, seeing the look itself now laid out on the slab
as one more thing among the stones, glances up at the print on the
scullery wall --- Escher's \textup{Eye}, a single vast eye with
something small looking back out of its own pupil.{]}}

\textbf{Tortoise:} The eye that looks --- set down as one more thing to
be looked at. \emph{{[}She turns back to the slab.{]}}

\textbf{Crab:} \emph{{[}looking over the spent stones, the row of
notches, the near-empty drawer{]}} Eleven years in that drawer, and the
better part of it gone in an afternoon. The daylight with it --- it is
the lamp now, and no more. \emph{{[}He gathers the stoppered flask and
sets it by.{]}}

\emph{{[}The room has drawn in to the size of the lamplight now: the
fire fallen low and red, the drawer all but empty, the little spent
stones gathered in a dull heap, the slate set off to one side of the
black. It is the hour when a long day's work is all but done, and the
doing of it has gone soft and unhurried. Someone's tea has been refilled
and let go tepid again, and nobody minds it.{]}}

\textbf{Tortoise:} \emph{{[}quietly{]}} Let them be spent; they were
made for it. Restless, every one --- and every one put to rest, a notch
apiece, and still.

\textbf{Crab:} There is warm tea yet, if anyone---

\textbf{Achilles:} \emph{{[}not quite hearing him, looking round the
settled table --- the spent stones, the folded halves, the one notch
alone on the slate --- and then to the corner{]}} In a moment. There is
one thing here that no one has looked at all evening. \emph{{[}nodding
to the sill{]}} Your friend there. Restless as any stone, both at once,
and nobody has spent it.

\emph{{[}The sill in the corner keeps its small mystery to the last: a
dent in the cushion that is also, looked at again, quite smooth; a
gathered warmth with the cold of the glass just beyond it; the gleam of
a half-open eye that might be a bead of lamplight caught and held, or
might be a look returned. Asleep, or awake --- the room has let it be
both all evening, and has asked nothing of it.{]}}

\textbf{Tortoise:} \emph{{[}softly{]}} Mind how you ask it, then. All
evening you have asked it loose things and had loose things back. A
plain question is a look --- and a look, tonight, we have seen what it
does.

\textbf{Achilles:} I have. \emph{{[}to the Cat, gently{]}} I have looked
at everything else on this table. Now I look at you. Cat --- asleep, or
awake?

\textbf{Cat:} Awake.

\emph{{[}The Cat settles to its one answer; and as the rest of it goes
from the sill --- the curled warmth, the folded paws, the one half-open
eye --- the last of it to fade is the grin, which hangs a moment over
the empty place, and then is gone too.{]}}

\emph{{[}The corner is only a windowsill now, and a cushion, and the
cold coming off the glass. The room settles the last of the way down
into its quiet: the fire sunk to a bed of embers, the lamp steady, and
on the table where the work left them the spent stones in a dull heap,
the kept stone still under its cloth at the far end, the viewer open
with its needle stilled, the flask stoppered and set by, the slate off
to one side. In the black window the whole warm room stands doubled and
still, and beyond it the snow has stopped falling at last. On the
scullery wall the great eye has gone back into the dark of its frame, a
plain rectangle with a little lamplight on its grain, part of the room's
stillness now. Only the lamp, and the embers, and no one left in the
room with a question to ask of anything.{]}}

\hypertarget{chapter-9-now-you-look}{%
\chapter{Chapter 9: Now You Look}\label{chapter-9-now-you-look}}

\emph{--- Lullaby for a Cat ---}

\emph{Part III: Pictures Become Proofs}

\emph{Escher: ``Eye'' --- a single vast eye, with something looking back
from inside its own pupil}

\begin{center}\rule{0.5\linewidth}{0.5pt}\end{center}

When you finally look at something, what has the look done to it --- and
could anyone far away ever tell that you looked?

For eight chapters we have kept the lights off. We prepared states, slid
boxes round bends, cooked tangles down to crossings, and every time a
question was asked of a wire we quietly kept only the answer we wanted
--- the postselected cap, the outcome that happened to suit us. No real
device is so obliging. A real device forces one of several answers on
you, at random, and destroys the rest; and the forcing is not a footnote
to the physics but a process in its own right, a thing that has to be
drawn. The last chapter promised as much and then closed the book for
the night. Now we open it and look.

Two things go wrong at once, and this chapter is about both. The first
is near at hand: a look leaves the thing looked at changed, and if we
did not read the outcome we are left holding something that is neither
the restless superposition we started with nor a plain definite answer
--- something new, for which a ket is the wrong kind of object. We shall
build the right kind. The second is across the room: if a look here
changed a thing here, and that thing was one half of a pair shared with
someone far away, did the look reach \emph{there}? Could a message be
sent, faster than any letter, by choosing what to do to one's own half?
The answer is no, exactly no, and the reason is one of the most quietly
astonishing facts in the theory. Half of a shared pair, examined on its
own, is the same featureless grey no matter what is done to the other
half --- and drawn in the right notation, that fact is a single tug on a
bent wire.

To get there we need three new marks on the page, and one honest
reckoning with the number in front of a cup. Let us make the marks.

\begin{center}\rule{0.5\linewidth}{0.5pt}\end{center}

\hypertarget{the-look-that-forces}{%
\section{§9.1 The Look That Forces}\label{the-look-that-forces}}

Achilles has never once watched a wire. Let him look now, at the
plainest thing there is: a single qubit, and the question ``are you
\(\ket0\) or \(\ket1\)?''

\textbf{On the blackboard.} To ask that question we use two
\textbf{projectors}, \(P_0=\ket0\bra0\) and \(P_1=\ket1\bra1\) --- each
one a box that keeps the part of a state lying along its own axis and
throws the rest away. Let the qubit be the unit vector
\(\ket\psi=\psi_0\ket0+\psi_1\ket1\). The \textbf{Born rule}, which is
the whole physics of a look, says two things. The look returns outcome
\(i\) with probability

\[p_i \;=\; \bra\psi P_i\ket\psi \;=\; |\psi_i|^2 ,\]

and, having returned \(i\), it leaves the qubit in the state
\(P_i\ket\psi/\sqrt{p_i}=\ket i\). Which outcome comes is not fixed by
anything: it is a genuine coin, weighted by the \(|\psi_i|^2\). Ask
\(\ket+ = \tfrac1{\sqrt2}(\ket0+\ket1)\) and the two outcomes each have
probability \(\tfrac12\); ask \(\ket0\) and the answer is \(0\) with
certainty, \(p_0=1\). And once the look has spoken, the question is
settled: measuring a second time returns the same \(i\), because
\(P_i\ket i=\ket i\). The look did not \emph{read off} a bit the qubit
was already carrying. It \emph{made} one, and destroyed the alternative.

What comes out is a record --- a definite value, already decided, no
longer able to interfere with anything. It is a new kind of quantity,
and it deserves a new kind of wire.

\textbf{In the sketchbook.} So far every wire has been a quantum wire,
drawn as a single stroke. A definite outcome is different: its value is
settled, it carries no phase, it cannot be put into superposition with
anything. We draw it as a \textbf{classical wire} --- a single line, but
stroked \emph{double}, to keep it distinct on the page from a quantum
wire.

\begin{center}
\begin{tikzpicture}[sd]
  \draw[cwire] (0,0) -- (3,0);
\end{tikzpicture}
\end{center}

One system, one wire; the double stroke is a reminder, not a second
line. Now the look itself. To draw it we lean on a fact we have owned
since Chapter 6. There we saw that a state has a \textbf{conjugate},
\(\ket{\psi^{*}}\) --- the same ket with every amplitude
complex-conjugated (§6.2), a genuinely different state when \(\psi\) is
not real. Feed a green \(Z\)-spider two legs --- on one, the state
\(\ket\psi\); on the other, its conjugate \(\ket{\psi^{*}}\), which is
an ordinary state box, a ket running rightward, and emphatically
\emph{not} the effect \(\bra\psi\). Let the spider have one classical
leg out. This is the \textbf{read-spider},

\[Z(2\!\to\!1)\;=\;\sum_i \ket i\bra{ii}\;=\;\ket0\bra{00}+\ket1\bra{11},\]

the very \(Z\)-spider of Chapter 8, put to a new use. On
\(\ket\psi\otimes\ket{\psi^{*}}\) it keeps only the branches where the
two legs agree, and passes their product to the classical wire:

\[\sum_i \psi_i\,\overline{\psi_i}\,\ket i \;=\; \sum_i |\psi_i|^2\,\ket i \;=\; (p_0,\,p_1).\]

The classical wire comes out carrying the Born probabilities, exactly
--- for \(\ket+\) the vector \((\tfrac12,\tfrac12)\), for \(\ket0\) the
vector \((1,0)\).

\begin{center}
\begin{tikzpicture}[sd,
  zsp/.style={circle,draw,line width=0.9pt,fill=green!55,minimum size=1.4em,inner sep=0pt}]
  \node[sdstate] at (0,0.9) {$\psi$};
  \node[sdstate] at (0,-0.9) {$\psi^{*}$};
  \draw[wire] (0,0.9) -- (2.1,0.9) -- (2.8,0);
  \draw[wire] (0,-0.9) -- (2.1,-0.9) -- (2.8,0);
  \node[zsp] at (2.8,0) {};
  \draw[cwire] (2.8,0) -- (5,0);
\end{tikzpicture}
\end{center}

\begin{quote}
\textbf{Result 9.1 (the look reads the Born rule).} The read-spider
\(Z(2\!\to\!1)\), fed a state on one leg and its conjugate on the other,
outputs on its classical wire the vector of Born probabilities
\(\bigl(|\psi_0|^2,|\psi_1|^2\bigr)\) --- exactly, scalar \(1\). See
\texttt{Code/Ch9/Measurement.py}.
\end{quote}

\textbf{\emph{Why this definition.}} A probability is an amplitude
multiplied by its own conjugate, so the look must be handed \emph{two}
copies of the state --- itself and its mirror --- for the spider to
multiply; and the outcome must ride a wire of its own, because it is no
longer a thing that can interfere, only a thing that has happened.

A word on scalars before we go on. From Chapter 8 we allowed ourselves
to state a picture equality up to a non-zero number. That licence is
withdrawn for the rest of this chapter wherever the number \emph{means}
something --- a probability, a trace, a weight in a mixture. Every
equation below is exact, the scalar written in, and we say so each time.

\begin{quote}
\textbf{Feynman's Notebook.}

The idea that a look leaves a \emph{mixture}, and that the mixture is
bookkept by a matrix rather than a vector, goes back to John von
Neumann. Wojciech Zurek's review argues that in a purely unitary account
``there is no room for a `real collapse','' and traces that conclusion
to von Neumann (1932). Zurek's own subject is \emph{decoherence}: an
apparatus that couples to a qubit is itself watched by an unruly
environment, which ``monitors'' it, selects a preferred set of
\textbf{pointer states}, and enforces the effective --- Zurek's phrase
--- ``collapse of the wavepacket.'' The coupling that starts it all is a
controlled-NOT, ``known as a `bit by bit measurement' (Zurek, 1981;
1983)'': the apparatus copies the system's bit the way the green dot
copies a wire. On this telling the classical world is not put in by
hand; it is what survives the environment's looking on our behalf.
\end{quote}

\begin{center}\rule{0.5\linewidth}{0.5pt}\end{center}

\hypertarget{a-state-thats-been-looked-at}{%
\section{§9.2 A State That's Been Looked
At}\label{a-state-thats-been-looked-at}}

Achilles measured the qubit but --- suppose --- did not read the slate.
The outcome is settled, but he does not know it. What is he holding? Not
\(\ket0\); not \(\ket1\); not the \(\ket+\) he began with, whose phase
is gone. It is a thing with no ket, and Chapter 6 promised it to us:
back in §6.2 we said there was no object that could be called ``the
state of one wire,'' and that this chapter would build one, and that it
would not be a ket. Here it is.

\textbf{On the blackboard.} A \textbf{density matrix} is a weighted list
of possibilities,

\[\rho \;=\; \sum_i p_i\,\ket{\psi_i}\bra{\psi_i}, \qquad p_i\ge 0,\ \sum_i p_i = 1 .\]

Achilles's looked-at-unread \(\ket+\) is: with probability \(\tfrac12\)
it is \(\ket0\), with probability \(\tfrac12\) it is \(\ket1\). So

\[\rho \;=\; \tfrac12\ket0\bra0 + \tfrac12\ket1\bra1 \;=\; \tfrac12 I ,\]

the \textbf{maximally mixed state} --- the flat grey that knows nothing,
the same in every basis. Compare it with the state he \emph{started}
from: the pure
\(\ket+\bra+ = \tfrac12\bigl(\begin{smallmatrix}1&1\\1&1\end{smallmatrix}\bigr)\),
whose off-diagonal \(\tfrac12\)'s are exactly the phase the look
destroyed. Both have trace \(1\); that is the density matrix's one law,
\(\operatorname{tr}\rho=1\), the total probability. A \textbf{pure}
state is the special case \(\rho=\ket\psi\bra\psi\) --- one possibility,
weight one, a matrix of rank one. Everything else is a genuine mixture,
and no ket describes it, because a ket has no room to record that a
phase was lost.

\textbf{In the sketchbook.} The picture of a density matrix is
\textbf{doubling}. A quantum system is drawn, from now on, as \emph{two}
lines --- its wire, and a second wire carrying its \textbf{conjugate}
(mirror) copy. A pure state \(\ket\psi\) is its two copies side by side,
matched: \(\ket\psi\) above and \(\ket{\psi^{*}}\) below, exactly the
pair the read-spider ate. A general \(\rho\) is a single box straddling
the two lines, whose copies no longer come apart:

\[\text{pure }\ket\psi:\quad
\begin{tikzpicture}[sd]
  \draw[wire] (0,0.6)--(1.8,0.6);
  \draw[wire] (0,-0.6)--(1.8,-0.6);
  \node[sdstate] at (0,0.6) {$\psi$};
  \node[sdstate] at (0,-0.6) {$\psi^{*}$};
\end{tikzpicture}
\qquad\qquad
\text{mixed }\rho:\quad
\begin{tikzpicture}[sd]
  \draw[wire] (0,0.6)--(1.8,0.6);
  \draw[wire] (0,-0.6)--(1.8,-0.6);
  \node[sdstate, minimum height=2.6em] at (0,0) {$\rho$};
\end{tikzpicture}\]

The second copy is the \emph{conjugate}, and it matters which. When a
unitary \(U\) evolves the state, \(\rho\mapsto U\rho U^{\dagger}\); in
doubled form \(U\) acts on the top (ket) copy and its conjugate
\(\bar U\) on the bottom (mirror) copy, because the mirror carries
conjugated amplitudes and a conjugated amplitude is turned by the
conjugated matrix. We draw this doubled process as a box on both lines
and call it \(U\otimes\bar U\):

\[\begin{tikzpicture}[sd]
  \draw[wire] (0,0.7)--(3,0.7);
  \draw[wire] (0,-0.7)--(3,-0.7);
  \node[sdbox] at (1.5,0.7) {$U$};
  \node[sdbox] at (1.5,-0.7) {$\bar U$};
\end{tikzpicture}\]

Note that this is neither the transpose \(U^{T}\) (the half-turn of
Chapter 6) nor the dagger \(U^{\dagger}\) (the left-to-right reflection
of Chapter 3): it is the plain conjugate, no wires bent, no reflection.
Doubling is what makes the scalar in a picture a real, non-negative,
\emph{meaningful} number: an overall phase \(e^{i\theta}\) on the top
copy is cancelled by \(e^{-i\theta}\) on the mirror, so only magnitudes
survive.

\textbf{\emph{Why this definition.}} Why carry a whole second copy of
every wire? Because a probability lives in the product of an amplitude
with its conjugate, and only an object that holds both --- the wire
\emph{and} its mirror --- can record what a look leaves. A ket holds
one; the doubled wire holds both, which is exactly enough to hold a
mixture. See \texttt{Code/Ch9/DensityMatrix.py}.

\begin{quote}
\textbf{Dirac's Blackboard.}

\emph{The state that is not a ket.} Chapter 6 asked for ``the state of
the top wire'' of a Bell pair and had nothing to hand back. Now we do:
it is the reduced density matrix, and for \(\ket{\Phi^+}\) it will turn
out to be \(\tfrac12 I\). Three facts pin \(\rho\) down. It is
\textbf{Hermitian}, \(\rho=\rho^{\dagger}\), since each
\(\ket{\psi_i}\bra{\psi_i}\) is. It is \textbf{positive},
\(\bra\phi\rho\ket\phi=\sum_i p_i|\braket{\phi|\psi_i}|^2\ge0\) for
every \(\ket\phi\). And it has \textbf{unit trace},
\(\operatorname{tr}\rho=\sum_i p_i=1\). Purity is measured by
\(\operatorname{tr}(\rho^2)\): equal to \(1\) for a pure
\(\ket\psi\bra\psi\), and \(\tfrac1d\) --- here \(\tfrac12\) --- for the
maximally mixed \(\tfrac1d I\), the flattest grey a \(d\)-dimensional
system allows. No column of amplitudes carries this; a ket has
\(\operatorname{tr}(\rho^2)=1\) always. The mixture is a new object, and
it needed the whole matrix to say it.
\end{quote}

\begin{center}\rule{0.5\linewidth}{0.5pt}\end{center}

\hypertarget{covering-a-half-partial-trace}{%
\section{§9.3 Covering a Half --- Partial
Trace}\label{covering-a-half-partial-trace}}

Achilles holds one half of a shared pair and wants its honest
description, the far half being none of his business. He must
\textbf{cover} the other half --- not measure it, not read it, simply
refuse to account for it. What is left?

\textbf{On the blackboard.} The description of the half you keep is the
\textbf{partial trace} over the half you discard. For a joint density
matrix \(\rho_{AB}\) on two systems, the marginal on \(A\) is

\[\rho_A \;=\; \operatorname{tr}_B\rho_{AB} \;=\; \sum_c (I\otimes\bra c)\,\rho_{AB}\,(I\otimes\ket c),\]

summing the effect \(\bra c\) over a basis of \(B\) on both sides. If
the pair is a plain product \(\rho_A\otimes\rho_B\), this returns
\(\rho_A\,\operatorname{tr}\rho_B=\rho_A\): covering an unrelated half
changes nothing, as it should. The interest is entirely in what it does
to a pair that is \emph{not} a product --- which is the next section's
business.

\textbf{In the sketchbook.} To cover a system is to \textbf{discard} it.
We draw the discard as a new mark: a wire that ends not in a box or a
bend but in the \textbf{ground symbol} --- a short stem closing in a few
shrinking bars, the sign that a wire is let go into the environment and
never asked another question.

\begin{center}
\begin{tikzpicture}[sd]
  \draw[wire] (0,0) -- (1.5,0);
  \sdground{1.5}{0}
\end{tikzpicture}
\end{center}

What does discarding \emph{do}, in the doubled picture? A system is now
two lines, its wire and its mirror. To let it go is to fold the two
copies together --- to join them with a \textbf{cap}, the effect
\(\bra{00}+\bra{11}\) of Chapter 6. The ground symbol on a system is
shorthand for exactly this cap on its two copies:

\[\begin{tikzpicture}[sd]
  \draw[wire] (0,0) -- (1.5,0);
  \sdground{1.5}{0}
\end{tikzpicture}
\;\;=\;\;
\begin{tikzpicture}[sd]
  \draw[wire] (0,0.6)--(1.6,0.6);
  \draw[wire] (0,-0.6)--(1.6,-0.6);
  \sdcap{1.6}{0.6}{-0.6}
\end{tikzpicture}\]

the left drawn as one line only because the doubling is understood;
unfold it and the two legs of the cap reappear. As always the cap is
unnormalized, \(\bra{00}+\bra{11}\); the scalar is load-bearing here and
we shall watch it. Cap a system to its mirror and you are computing a
trace. Applied to a doubled \(\rho\) --- the box on two lines --- the
cap contracts the two copies and returns \(\operatorname{tr}\rho\).

Two small closed pictures make the scalar honest. The cup joining a
system to its mirror, \(\sum_i\ket{ii}\), is the doubled form of the
identity matrix (Chapter 6, §6.5). Cap it --- discard it --- and you get
a number,
\(\sum_i\bra{ii}\sum_j\ket{jj}=\sum_i 1 = 2 = \operatorname{tr} I\). Put
a \(\tfrac12\) in front and the number is \(1\): the mirror-cup, scaled
by \(\tfrac12\), is the maximally mixed state \(\tfrac12 I\), and its
trace is \(1\) as a density matrix's must be. This is the image we will
need in a moment:

\[\tfrac12 I \;=\; \tfrac12\;
\begin{tikzpicture}[sd]
  \sdcup{0}{0.6}{-0.6}
  \draw[wire] (0,0.6)--(1.8,0.6);
  \draw[wire] (0,-0.6)--(1.8,-0.6);
\end{tikzpicture}\]

--- a system cupped to its own mirror copy, with the number \(\tfrac1d\)
beside it. One convention, to fix the pictures that follow: we draw a
system beside its mirror copy, and when a system is to be discarded we
keep its two copies side by side, so that its cap is small and its
neighbours are undisturbed.

\begin{quote}
\textbf{Result 9.2 (discard is partial trace).} For every joint state
\(\rho_{AB}\), capping the two copies of the \(B\)-system in the doubled
picture equals the partial trace \(\operatorname{tr}_B\rho_{AB}=\rho_A\)
--- exactly, scalar included.
\end{quote}

\emph{Proof idea.} The cap says ``join \(B\) to its mirror and sum'';
summing a matrix's diagonal over one factor is precisely tracing that
factor out.

\emph{Proof (blackboard).} Write
\(\rho_{AB}=\sum \rho_{ab,a'b'}\,\ket{ab}\bra{a'b'}\). In doubled form
this is a vector with a ket copy \(\ket{ab}\) and a mirror copy
\(\ket{a'b'}\) on each term. The cap on the \(B\)-copies applies
\(\sum_c\bra{cc}\) to the pair \((b,b')\), keeping only \(b=b'=c\) and
summing:

\[\sum_{a,a',c}\rho_{ac,a'c}\,\ket a\bra{a'}=\operatorname{tr}_B\rho_{AB}.\]

Every ``every'' here is checked on a random \(\rho_{AB}\), and the two
scalars are confirmed with it --- the doubled bare cup discards to \(I\)
(a closed value of \(2\)), the maximally mixed \(\tfrac12\)-cup to trace
\(1\). See \texttt{Code/Ch9/PartialTrace.py}. \(\;\blacksquare\)

\textbf{\emph{Why this definition.}} Discarding must be the trace and
nothing else --- the one effect that asks a system no question, only
lets it go --- and drawing it as a cap is what lets a discarded half
inherit \emph{yanking}. That inheritance is the whole reason the
chapter's theorem is one move rather than a page.

\begin{quote}
\textbf{Coecke's Sketchbook.}

\emph{The environment structure.} The ground symbol is not an ornament;
it is an \emph{axiom}. Coecke and Perdrix (2010) equip a theory with a
chosen discard \(\top_A\colon A\to I\) on every system --- for density
matrices, exactly the trace \(\rho\mapsto\operatorname{tr}\rho\) --- and
read it, in their words, as ``this system being part of the
environment.'' What they require of it is a single law (their Eq. 10):
two pure processes are the same up to a phase precisely when they look
the same after the discard, \(\top\circ f=\top\circ g\). Take \(g\) the
identity and the law says something we will use directly --- a process
leaves the discard untouched, \(\top\circ f=\top\), exactly when it is
trace-preserving. This book calls that property \textbf{causality}; the
word is ours, not the paper's, but the fact is theirs, and it is what
forbids a signal.
\end{quote}

\begin{center}\rule{0.5\linewidth}{0.5pt}\end{center}

\hypertarget{the-rosetta-table-7}{%
\section{§9.4 The Rosetta Table}\label{the-rosetta-table-7}}

The rows this chapter adds to the running dictionary, numbered on from
Chapter 8's sixty-four. Three introduce new picture elements --- the
classical wire, doubling, and the discard; the rest are constructions
built from them. The last three name what §9.6 will prove --- a channel,
its dilation, and the causality that makes the master hold; they are
gathered here so the chapter's additions sit on one page.

\begingroup\renewcommand{\arraystretch}{1.5}

\begin{longtable}[]{@{}
  >{\raggedright\arraybackslash}p{(\columnwidth - 4\tabcolsep) * \real{0.2727}}
  >{\raggedright\arraybackslash}p{(\columnwidth - 4\tabcolsep) * \real{0.3636}}
  >{\raggedright\arraybackslash}p{(\columnwidth - 4\tabcolsep) * \real{0.3636}}@{}}
\toprule\noalign{}
\begin{minipage}[b]{\linewidth}\raggedright
\#
\end{minipage} & \begin{minipage}[b]{\linewidth}\raggedright
Algebra (Dirac)
\end{minipage} & \begin{minipage}[b]{\linewidth}\raggedright
Picture (Coecke)
\end{minipage} \\
\midrule\noalign{}
\endhead
\bottomrule\noalign{}
\endlastfoot
65 & a \textbf{classical wire}: a system whose value is already definite
(an outcome) & a single wire drawn with a double stroke (\texttt{cwire})
--- \emph{new element} \\
66 & \textbf{doubling}: a system carries \(\rho\); a pure \(\ket\psi\)
becomes \(\ket\psi\otimes\ket{\psi^{*}}\) & a quantum wire drawn as two
lines --- the wire and its conjugate (mirror) copy --- \emph{new
element} \\
67 & \textbf{discard} \(\top_A\colon\rho\mapsto\operatorname{tr}\rho\)
(the environment structure) & a \textbf{ground symbol} closing a wire
\(=\) a cap joining a doubled system's two copies --- \emph{new
element} \\
68 & projective measurement; Born rule \(p_i=\bra\psi P_i\ket\psi\) &
the \textbf{read-spider}: a phase-free \(Z\)-spider on a doubled system
emitting one classical wire, \(\sum_i\ket i\bra{ii}\) \\
69 & maximally mixed state \(\tfrac1d I\) & a system cupped to its own
mirror copy, with the number \(\tfrac1d\) beside it \\
70 & channel \(\Lambda(\rho)=\sum_i K_i\rho K_i^{\dagger}\) (CPTP) & a
box on doubled wires --- the doubled process \\
71 & Stinespring: \(\Lambda=\operatorname{discard}_E\circ V\), \(V\) an
isometry \(A\to B\otimes E\) & a doubled unitary \(V=U\otimes\bar U\)
with the environment leg grounded \\
72 & causality:
\(\operatorname{tr}(\Lambda\rho)=\operatorname{tr}\rho\),
i.e.~\(\top\circ\Lambda=\top\) & a unitary capped by a discard yanks to
a bare discard \\
\end{longtable}

\endgroup

Rows 65--69 are exact by construction. Rows 70--72 are the subject of
§9.6, and each is exact there, scalar included.

\begin{center}\rule{0.5\linewidth}{0.5pt}\end{center}

\hypertarget{the-theorem-twice-8}{%
\section{§9.5 The Theorem, Twice}\label{the-theorem-twice-8}}

Now the pair. Let Achilles and a partner across the room share the Bell
state \(\ket{\Phi^+}=\tfrac1{\sqrt2}(\ket{00}+\ket{11})\), the top wire
his, the bottom wire hers. He wants to send a message by acting on his
half --- flip it, measure it, leave it alone --- and have her half
twitch in reply. Take the simplest scheme: he does nothing, or he
applies the bit-flip \(X\). Can she tell which, looking only at her own
half? Her half's honest description is the marginal,
\(\operatorname{tr}\) over \emph{his} wire --- and here is the theorem
that dashes his hopes.

\begin{quote}
\textbf{Result 9.3 (half a Bell pair is maximally mixed;
no-signalling).} Tracing out either half of \(\ket{\Phi^+}\) leaves the
other maximally mixed,
\[\operatorname{tr}_2\,\ket{\Phi^+}\bra{\Phi^+} \;=\; \tfrac12 I,\]
exactly, the \(\tfrac12\) included. Moreover, for \textbf{any}
trace-preserving channel \(\Lambda\) applied to the traced-out half, the
surviving marginal is unchanged:
\[\operatorname{tr}_2\bigl((I\otimes\Lambda)\,\ket{\Phi^+}\bra{\Phi^+}\bigr)\;=\;\tfrac12 I .\]
The far half is \(\tfrac12 I\) no matter what is done to the near half.
No message crosses.
\end{quote}

Here \(\operatorname{tr}_2\) traces the \emph{bottom} wire --- the near,
acted-on half --- and returns the far half's description.

\emph{Proof idea.} Her half, asked anything, is a fair coin (Chapter 6,
§6.2), in every basis at once --- and that flat grey is \(\tfrac12 I\).
Whatever he does is a trace-preserving process, and a trace-preserving
process is, by definition, one the discard cannot feel. The picture will
make the second sentence a single tug.

\hypertarget{diracs-blackboard-8}{%
\subsection{Dirac's Blackboard}\label{diracs-blackboard-8}}

Write \(\ket{\Phi^+}\bra{\Phi^+}=\tfrac12\sum_{a,b}\ket{aa}\bra{bb}\).
Trace the second wire,
\(\operatorname{tr}_2(\ket{aa}\bra{bb})=\ket a\bra b\,\braket{b|a}=\delta_{ab}\ket a\bra b\):

\[\operatorname{tr}_2\ket{\Phi^+}\bra{\Phi^+} \;=\; \tfrac12\sum_{a,b}\delta_{ab}\,\ket a\bra b \;=\; \tfrac12\sum_a \ket a\bra a \;=\; \tfrac12 I .\]

Flipping his half with \(X\) gives \((I\otimes X)\ket{\Phi^+}\), whose
marginal recomputes the same way to \(\tfrac12 I\); doing nothing gives
\(\tfrac12 I\) as well. For a general channel
\(\Lambda(\rho)=\sum_i K_i\rho K_i^{\dagger}\) with
\(\sum_i K_i^{\dagger}K_i=I\), the marginal is

\[\operatorname{tr}_2\bigl((I\otimes\Lambda)\ket{\Phi^+}\bra{\Phi^+}\bigr)
=\sum_i \operatorname{tr}_2\bigl((I\otimes K_i)\ket{\Phi^+}\bra{\Phi^+}(I\otimes K_i^{\dagger})\bigr),\]

and cyclicity of the (partial) trace gathers \(K_i^{\dagger}K_i\) under
it:

\[\sum_i\operatorname{tr}_2\bigl(\ket{\Phi^+}\bra{\Phi^+}(I\otimes K_i^{\dagger}K_i)\bigr)=\operatorname{tr}_2\bigl(\ket{\Phi^+}\bra{\Phi^+}(I\otimes I)\bigr)=\tfrac12 I .\]

The step that saves it is \(\sum_i K_i^{\dagger}K_i=I\) ---
trace-preservation, and nothing else. See
\texttt{Code/Ch9/NoSignalling.py}, which runs a random non-real Kraus
channel.

\hypertarget{coeckes-sketchbook-3}{%
\subsection{Coecke's Sketchbook}\label{coeckes-sketchbook-3}}

Draw \(\ket{\Phi^+}\bra{\Phi^+}\) doubled. Since
\(\ket{\Phi^+}=\tfrac1{\sqrt2}\text{cup}\), its doubled form is
\(\tfrac1{\sqrt2}\) times \(\tfrac1{\sqrt2}\), one factor for the ket
copy and one for the mirror:

\[\ket{\Phi^+}\bra{\Phi^+}\;=\;\tfrac12\;\text{cup}(S,M)\otimes\text{cup}(\bar S,\bar M),\]

with the far half's two copies \(S,\bar S\) outermost and the near
half's copies \(M,\bar M\) innermost. To trace the near half, cap \(M\)
to \(\bar M\) (Result 9.2). What stands on the page is a cup, a cap, a
cup, in a column --- an \(S\)-bend with the near half's copies pinched
in the middle:

\[\tfrac12\;
\begin{tikzpicture}[sd]
  \sdcup{0}{4.5}{3}
  \sdcup{0}{1.5}{0}
  \draw[wire] (0,4.5) -- (3,4.5);
  \draw[wire] (0,3) -- (2,3);
  \draw[wire] (0,1.5) -- (2,1.5);
  \draw[wire] (0,0) -- (3,0);
  \sdcap{2}{3}{1.5}
  \node[sdlabel,anchor=east] at (-0.2,4.5) {$S$};
  \node[sdlabel,anchor=east] at (-0.2,3) {$M$};
  \node[sdlabel,anchor=east] at (-0.2,1.5) {$\bar M$};
  \node[sdlabel,anchor=east] at (-0.2,0) {$\bar S$};
\end{tikzpicture}
\;\;\overset{\text{yank}}{=}\;\;
\tfrac12\;
\begin{tikzpicture}[sd]
  \sdcup{0}{1.5}{0}
  \draw[wire] (0,1.5) -- (3,1.5);
  \draw[wire] (0,0) -- (3,0);
  \node[sdlabel,anchor=east] at (-0.2,1.5) {$S$};
  \node[sdlabel,anchor=east] at (-0.2,0) {$\bar S$};
\end{tikzpicture}
\;\;=\;\;\tfrac12 I .\]

\textbf{One yank (Theorem 6.2).} The cup on \((S,M)\), the cap on
\((M,\bar M)\) and the cup on \((\bar M,\bar S)\) pull straight into a
single cup on \((S,\bar S)\). What remains is the far half cupped to its
own mirror, times \(\tfrac12\) --- which is the maximally-mixed picture
of §9.3, \(\tfrac12 I\). The scalar is carried, not dropped, at every
step.

And the corollary is one more tug. Before the discard, let the near half
pass through the channel --- a doubled box \(\Lambda\) on
\((M,\bar M)\). Being trace-preserving is, by the environment-structure
law (causality, proved as Result 9.4 in §9.6), exactly the statement
that the discard slides over it and off: capping after \(\Lambda\) is
capping alone. So the box vanishes under the cap before the yank even
begins, and the far half is \(\tfrac12 I\) as before, whatever
\(\Lambda\) was.

\hypertarget{the-verdict-8}{%
\subsection{The verdict}\label{the-verdict-8}}

The picture wins, on both counts, and we should say so plainly, because
it is the rarer verdict. The headline is a \emph{single yank} with a
\(\tfrac12\) written beside it, and the reason it is one yank is
visible: the maximally mixed state \emph{is} a system cupped to its
mirror, and half a Bell pair, discarded, straightens into exactly that.
The blackboard's partial trace is short too --- four lines --- but it
shows only \emph{that} the answer is \(\tfrac12 I\), never \emph{why}.
On the corollary the gap widens: the picture removes the channel with
one causality slide, while the blackboard must summon
\(\sum_i K_i^{\dagger}K_i=I\) and the cyclicity of the trace. What the
picture proves, though, is an equality of processes and no more: that a
certain marginal does not move. It is not a claim about bandwidth, nor a
security guarantee --- only that there is nothing here that reaches
there.

\begin{center}\rule{0.5\linewidth}{0.5pt}\end{center}

\hypertarget{the-unmoved-half}{%
\section{§9.6 The Unmoved Half}\label{the-unmoved-half}}

A clean look was too kind. A real measurement has a needle that swings,
a flask that vents, a leak into the room; drawn honestly, it is not a
bare cap but a whole apparatus. Does the leak carry a signal the clean
cap did not?

\textbf{On the blackboard.} A real, irreversible look is a
\textbf{channel} --- a completely positive, trace-preserving map. Drawn,
it is a box on the doubled wires (Rosetta row 70), and bending such a
box into a state is the process--state duality Chapter 6 named for Choi
and Jamiołkowski --- the ``grander kind of process'' it promised, now in
hand. Two classical facts, which we state and cite rather than prove,
tell us what every channel is made of. First, the \textbf{Kraus form}:
every channel can be written
\(\Lambda(\rho)=\sum_i K_i\rho K_i^{\dagger}\) with
\(\sum_i K_i^{\dagger}K_i=I\). Second, \textbf{Stinespring's dilation}:
every channel is a single unitary on the system together with a fresh
apparatus and environment, followed by discarding the environment ---
\(\Lambda(\rho)=\operatorname{tr}_E\bigl(V\rho V^{\dagger}\bigr)\) for
an isometry \(V\colon A\to B\otimes E\), \(V^{\dagger}V=I\). A leaky
look is a clean, reversible turn of a \emph{larger} system, with the
part we cannot keep let go. We take both facts, with their history, from
the eighth chapter of Nielsen and Chuang. Trace-preservation is the
causality condition
\(\operatorname{tr}(\Lambda\rho)=\operatorname{tr}\rho\): the discard
cannot feel the channel.

\textbf{In the sketchbook.} Causality, drawn, is Chapter 6 coming back
for its debt. Dilate the near half through the isometry, doubled as
\(V\otimes\bar V\) on the system's two copies, and then discard the
whole output --- cap the top copy to the bottom. Slide \(\bar V\) round
the cap (Theorem 6.4, the transpose trick in its cap form): a box
carried round a bend arrives transposed, so \(\bar V\) arrives as
\(\bar V^{T}=V^{\dagger}\), meeting \(V\). But \(V^{\dagger}V=I\), and
the cap is left bare.

\[\begin{tikzpicture}[sd]
  \draw[wire] (0,0.7)--(2.4,0.7);
  \draw[wire] (0,-0.7)--(2.4,-0.7);
  \node[sdbox] at (1.2,0.7) {$V$};
  \node[sdbox] at (1.2,-0.7) {$\bar V$};
  \sdcap{2.4}{0.7}{-0.7}
\end{tikzpicture}
\;\;\overset{\text{slide}}{=}\;\;
\begin{tikzpicture}[sd]
  \draw[wire] (0,0.7)--(2.4,0.7);
  \draw[wire] (0,-0.7)--(2.4,-0.7);
  \node[sdbox] at (1.2,0.7) {$V^{\dagger}V$};
  \sdcap{2.4}{0.7}{-0.7}
\end{tikzpicture}
\;\;=\;\;
\begin{tikzpicture}[sd]
  \draw[wire] (0,0.7)--(2,0.7);
  \draw[wire] (0,-0.7)--(2,-0.7);
  \sdcap{2}{0.7}{-0.7}
\end{tikzpicture}\]

\begin{quote}
\textbf{Result 9.4 (causality).} Discarding after an isometry is
discarding: a doubled isometry capped by a discard yanks to a bare
discard, because \(V^{\dagger}V=I\). For a general trace-preserving
channel this equality \emph{is} the definition of trace-preserving; via
Stinespring it reduces to the isometry case just drawn.
\end{quote}

\begin{quote}
\textbf{Feynman's Notebook.}

\emph{Where the channel picture was pinned down.} That every noisy
process is an operator-sum,
\(\Lambda(\rho)=\sum_i K_i\rho K_i^{\dagger}\), and that every such
process is a clean unitary on a larger system with a part discarded, are
the Kraus decomposition and Stinespring's dilation; we lean for both on
the eighth chapter of Nielsen and Chuang, where they are proved in full
and their history set out. The doubling, the discard and the environment
structure we draw to make them one picture are newer. Bob Coecke and
Simon Perdrix set them down in 2010 (arXiv:1004.1598), reading the trace
as a system passing quietly into its environment and axiomatising that
reading in their Equation 10 --- the law this chapter has been calling
causality.
\end{quote}

One more small equality, and it is the hinge of the master. A look does
not only leave the apparatus; it also writes an outcome, and sometimes
we do not read it. Reading with the read-spider and then discarding ---
\emph{forgetting} --- the classical outcome is the same as never having
read at all:

\[\Bigl(\textstyle\sum_j\bra j\Bigr)\circ\Bigl(\sum_i\ket i\bra{ii}\Bigr)\;=\;\sum_i\bra{ii}\;=\;\text{cap}.\]

\begin{quote}
\textbf{Result 9.5 (reading, then forgetting, is discarding).} The
read-spider followed by discarding its classical wire equals the plain
discard (cap) of the measured system --- exactly. See
\texttt{Code/Ch9/Stinespring.py}.
\end{quote}

Now the master. It is the no-signalling scene with the apparatus laid
bare, drawn in unfolded doubled form so that every line shows. On the
left, the doubled Bell state \(\ket{\Phi^+}\bra{\Phi^+}\) --- two cups
with the factor \(\tfrac12\) written beside them (cup
\(=\ket{00}+\ket{11}\), as always) --- ties the far half \(S\) to the
near half \(M\). The near half passes through the dilation box
\(V=U\otimes\bar U\), drawn as one box because its two halves --- \(U\)
on the kets, \(\bar U\) on the mirrors --- are wired the same way and
only too far apart to draw apart; into that box also go the doubled
pointer \(P\) (a fresh \(\ket0\)) and the doubled environment \(E\)
(another \(\ket0\)). Out the other side: the near half \(M\) is
discarded (a cap on its two copies), the environment \(E\) is discarded
(another cap), and the pointer \(P\) is read by the read-spider, which
emits the single classical outcome \(c\). The far half runs straight
across, untouched, above and below everything.

\begin{center}
\begin{tikzpicture}[sd,
  zsp/.style={circle,draw,line width=0.9pt,fill=green!55,minimum size=1.35em,inner sep=0pt},
  ket/.style={sdstate, minimum height=1.1em, minimum width=1.4em, inner sep=1pt, font=\footnotesize}]
  % far half, straight across
  \draw[wire] (0,11) -- (13,11);
  \draw[wire] (0,-0.5) -- (13,-0.5);
  % the factor 1/2 of the doubled Bell state (cup = |00>+|11>, so |Phi+><Phi+| = 1/2 cup cup)
  \node[sdlabel, font=\normalsize] at (-4,5.25) {$\tfrac12$};
  % Bell cups on the left: (S,M) small, (S-bar,M-bar) tall
  \sdcup{0}{11}{9}
  \sdcup{0}{7.5}{-0.5}
  % near half into V
  \draw[wire] (0,9) -- (3,9);
  \draw[wire] (0,7.5) -- (3,7.5);
  % pointer and environment fresh states into V
  \node[ket] at (1.2,6) {$0$};
  \node[ket] at (1.2,4.5) {$0$};
  \node[ket] at (1.2,3) {$0$};
  \node[ket] at (1.2,1.5) {$0$};
  \draw[wire] (1.2,6) -- (3,6);
  \draw[wire] (1.2,4.5) -- (3,4.5);
  \draw[wire] (1.2,3) -- (3,3);
  \draw[wire] (1.2,1.5) -- (3,1.5);
  % the dilation box
  \node[sdbox, minimum height=8.8em, minimum width=4em] at (5,5.25) {$U\otimes\bar U$};
  % outputs
  \draw[wire] (7,9) -- (9,9);
  \draw[wire] (7,7.5) -- (9,7.5);
  \sdcap{9}{9}{7.5}
  \draw[wire] (7,6) -- (8.5,6);
  \draw[wire] (7,4.5) -- (8.5,4.5);
  \draw[wire] (8.5,6) -- (9.3,5.25);
  \draw[wire] (8.5,4.5) -- (9.3,5.25);
  \node[zsp] at (9.3,5.25) {};
  \draw[cwire] (9.3,5.25) -- (13,5.25);
  \draw[wire] (7,3) -- (9,3);
  \draw[wire] (7,1.5) -- (9,1.5);
  \sdcap{9}{3}{1.5}
  % system labels
  \node[sdlabel,anchor=west] at (13.1,11) {$S$};
  \node[sdlabel,anchor=west] at (13.1,-0.5) {$\bar S$};
  \node[sdlabel,anchor=west] at (13.1,5.25) {$c$};
  \node[sdlabel,anchor=south] at (2.1,9.05) {$M$};
  \node[sdlabel,anchor=south] at (2.1,7.55) {$\bar M$};
  \node[sdlabel,anchor=east] at (0.4,6) {$P$};
  \node[sdlabel,anchor=east] at (0.4,4.5) {$\bar P$};
  \node[sdlabel,anchor=east] at (0.4,3) {$E$};
  \node[sdlabel,anchor=east] at (0.4,1.5) {$\bar E$};
\end{tikzpicture}
\end{center}

The lines are the far half and its mirror, \(S\) and \(\bar S\); the
near half and its mirror, \(M\) and \(\bar M\); the pointer and its
mirror, \(P\) and \(\bar P\); the environment and its mirror, \(E\) and
\(\bar E\); and the classical outcome \(c\) threaded out of the
read-spider. We call this diagram \textbf{The Unmoved Half}.

\begin{quote}
\textbf{Result 9.6 (the unmoved half).} In The Unmoved Half, \emph{avert
your eyes from the outcome} --- trace the classical wire \(c\) --- and
the far half evaluates to \(\tfrac12 I\), for \textbf{every} dilation
\(V\), independent of which measurement the near party chose. The far
half never moves.
\end{quote}

\emph{Proof.} Forget the outcome: by Result 9.5 the read-spider with
\(c\) discarded becomes a plain cap on \((P,\bar P)\), so now \emph{all}
of the box's outputs --- \(M\), \(P\), \(E\) and their mirrors --- are
discarded. By causality (Result 9.4) the whole discard slides over
\(V=U\otimes\bar U\) and off, and the fresh \(\ket0\) pointer and
environment each discard to \(\braket{0|0}=1\). What is left is the
doubled Bell cup with its near half capped --- the very picture of §9.5
--- which yanks to \(\tfrac12 I\). Nothing about \(U\) survived the
discard, so the far half is the same \(\tfrac12 I\) for every apparatus
and every choice of look. See \texttt{Code/Ch9/Stinespring.py}, which
builds the diagram, tallies its edges, and checks the marginal for a
random non-real \(V\). \(\;\blacksquare\)

Be careful with the converse, and honest about it. Do \emph{not} read
this as ``the far half is independent of the outcome.'' It is not: leave
\(c\) open --- condition on a particular reading --- and the far half is
thrown into the state correlated with that outcome, exactly as a shared
pair must be. What is independent of anything the near party does is
only the far half \emph{averaged over the outcomes it cannot see}. The
signal that could have crossed is drowned, precisely and completely, in
the party's own ignorance of which outcome fell.

The Tortoise, drawing the apparatus into the picture --- the look itself
now a thing on the page --- glances up at \emph{Eye} on the scullery
wall, the great eye with something watching back from inside its own
pupil, and remarks that the eye which looks is also a thing looked at.

\begin{quote}
\textbf{Dirac's Blackboard.}

\emph{Half a Bell pair, and Kraus in full.} The headline in coordinates:
\(\ket{\Phi^+}\bra{\Phi^+}=\tfrac12\bigl(\ket{00}\bra{00}+\ket{00}\bra{11}+\ket{11}\bra{00}+\ket{11}\bra{11}\bigr)\),
and \(\operatorname{tr}_2\) keeps the blocks whose second labels agree,
\(\tfrac12(\ket0\bra0+\ket1\bra1)=\tfrac12 I\). The off-diagonal
\(\ket{00}\bra{11}\) and \(\ket{11}\bra{00}\) --- the coherence, the
part that carried the correlation --- are killed by the partial trace,
not because they are small but because they are off the diagonal of the
traced factor. For the leaky look, the Kraus operators of a Stinespring
dilation are read straight off: \(K_i=(\,I_B\otimes\bra{e_i}\,)V\), one
for each environment basis vector \(\ket{e_i}\), and
\(\sum_i K_i^{\dagger}K_i=V^{\dagger}\bigl(I_B\otimes\sum_i\ket{e_i}\bra{e_i}\bigr)V=V^{\dagger}V=I\).
Trace-preservation is \(V^{\dagger}V=I\), wearing coordinates.
\end{quote}

\begin{center}\rule{0.5\linewidth}{0.5pt}\end{center}

\hypertarget{now-the-other-half}{%
\section{§9.7 Now the Other Half}\label{now-the-other-half}}

We came to look, and looking cost us the comfort of the ket. A qubit,
measured and unread, is not a ket but a mixture, and to draw a mixture
every wire had to be doubled, paired with its own conjugate. A part of a
pair, honestly described, is the whole with the rest discarded --- a cap
folding a system into its mirror, which is the partial trace. Half of a
Bell pair, so discarded, is \(\tfrac12 I\); and it stays \(\tfrac12 I\)
through any channel on the other half, because a trace-preserving
process is one the discard slides over and off. That is no-signalling,
and it survived the apparatus being laid bare: a real, leaky,
Stinespring-dilated look, its environment vented and its outcome
forgotten, leaves the far half exactly as grey as before. The whole of
that argument was one yank and one slide.

But we have been throwing outcomes away. A cap, since Chapter 6, has
been a \emph{postselected} answer --- the one reading out of several
that happened to suit us --- and Chapter 6 warned that the other
readings were not failures but corrections waiting to be applied. Now we
have the tools to apply them: a classical wire to carry the outcome that
fell, and a box, chosen by that wire, to set the far half right. Thread
a cup, keep the outcome, send it across, and turn the far half by the
correction it names, and the postselected cap of Chapter 6 becomes a
protocol that works every time. That is teleportation, and it is the
next chapter.

\hypertarget{out-of-order-canon}{%
\chapter{Out-of-Order Canon}\label{out-of-order-canon}}

\begin{center}\rule{0.5\linewidth}{0.5pt}\end{center}

\emph{雪晴や}\\
\emph{文より先に}\\
\emph{届く影}

\emph{yukibare ya}\\
\emph{fumi yori saki ni}\\
\emph{todoku kage}

\emph{Clear after the snow ---}\\
\emph{before the letter itself}\\
\emph{its shadow arrives.}

\begin{center}\rule{0.5\linewidth}{0.5pt}\end{center}

\emph{{[}The Crab's parlour, the next morning. Last night's black slab
is still out along the table, wiped clean and bare; beside it, off the
black, stands something wrapped in chamois. The Tortoise is already at
the table with her tea. Achilles comes in from the cold with a handful
of letters.{]}}

\emph{{[}It is a hard, bright, windless morning, and the parlour is full
of a white light that comes from everywhere at once and throws no
shadows. The one tall window at the near end of the long table, which
last night was a black glass giving the lamplit room back to itself, has
gone white in the dark: every pane is furred from edge to edge with
frost, fern grown over fern, so thick that the lane beyond is only a
brightness. The sun is up, low and white on the snow outside, but it is
still round the side of the house and has not come to this glass. At eye
height in one pane there is a small round clear place the size of a
breath, where someone stood at first light to watch for the post, and
already its rim is feathering over again. Last night's lamp stands cold
on the dresser, with the viewer beside it shut in its case; the fire has
been built up high, the kettle murmurs on the hob, and a plate of toast
is keeping warm on the hearth. Down the middle of the table the black
slab runs from the window end to the far end, bare. On the wall facing
the slab hangs a print the Crab has had so long that he has stopped
seeing it --- Escher's \textup{Magic Mirror}: a mirror standing upright,
and creatures stepping out of it. The frost-light falls on it evenly,
without favour, as it falls on everything this morning; it was on that
wall long before anyone came in out of the cold, and it will be there
after. The Tortoise has one of last night's quilts round her shoulders,
and the door, as Achilles pushes it to behind him, lets in a last breath
of air so cold it smells of iron.{]}}

\textbf{Achilles:} Good morning! I met your postman at the gate, Mr
Crab, blowing on his fingers, and said I'd carry the last few yards for
him. Three for you. One of them's from Ostend.

\textbf{Crab:} My cousin. He writes every quarter-day to ask for the
loan of my tuning key, and every quarter-day I tell him no. \emph{{[}He
takes the letters and props them against the teapot.{]}} Come to the
fire, come in. You've brought the whole frost in on your coat.

\textbf{Achilles:} It's the white sort --- every twig in the lane has
grown a fur of it overnight, and the hedges crackle when you brush them.
I can't feel my hands at all.

\emph{{[}He peels off a pair of knitted mittens gone stiff as card with
rime and lays them side by side on the fender to thaw, and at once they
begin to steam. The Crab pushes a chair in close to the fire for him;
Achilles sits on the edge of it and holds his bare hands out to the
blaze, flexing them, as if to make sure they are still his.{]}}

\textbf{Tortoise:} I never went home. The snow had the lane before we'd
finished last night, and Mr Crab put me in the little room over the
stairs with three quilts and a hot brick.

\textbf{Crab:} And she was up before the light, following the frost on
the windowpane with one claw, as if it might need her help.

\textbf{Tortoise:} It doesn't. It is a very good draughtsman. I was only
taking lessons.

\textbf{Achilles:} \emph{{[}folding both hands round the cup he is
given{]}} Ah --- that's better. That's much better. Your postman told me
nothing he carries is ever the worse for the road, except eggs.

\textbf{Crab:} Eggs, and my cousin's temper. Mind the slab, by the way;
I've left it out from last night, and it will show you two of everything
you set on it. \emph{{[}He goes to the dresser drawer, looks into it a
moment, and brings out what is left there: the last two pairs of little
stones.{]}} The last of them. I didn't like to leave them alone in a
drawer on a morning like this.

\emph{{[}He lays the first pair at the near end of the slab. One stone
he covers at once with a small square of cloth, and it stays there,
hidden, its reflection hidden with it under the cloth's edge; its twin
he sets beside it, bare. The second pair he parts: one stone he leaves
at the near end, by the others; the other he carries down the whole
length of the black and sets at the far end, alone. Each stone, as it
touches down, finds its reflection waiting for it underneath.{]}}

\textbf{Crab:} The one under the cloth is my keepsake. Nobody touches it
--- not today, not ever.

\textbf{Achilles:} And why send that one all the way down there by
itself?

\textbf{Crab:} Because it is a long table, and I want it a long way off.
You'll see.

\emph{{[}The stones sit where he left them down the length of the black:
at the near end the small square of cloth, with the keepsake under it;
the keepsake's twin, bare beside it; the near half of the parted pair
beside them; and far away at the other end, by itself, its partner. The
white window-light lies on the slab like milk on slate. Achilles, his
cup held up under his chin for the warmth of it, keeps looking down the
table at the lone stone at the far end, as though it might be
lonely.{]}}

\begin{center}\rule{0.5\linewidth}{0.5pt}\end{center}

\hypertarget{first-challenge-four-ways-to-answer}{%
\section{\texorpdfstring{First Challenge --- \emph{Four Ways to
Answer}}{First Challenge --- Four Ways to Answer}}\label{first-challenge-four-ways-to-answer}}

\textbf{Crab:} \emph{{[}drawing the wrapped thing nearer and folding
back the chamois{]}} Now. I promised myself I'd bring this out the
morning after we looked, and it is the morning after. The pairing-glass
--- from the same bench in Delft; the daughter ground the lens and her
father made the brass. It has two cups side by side under one lens, and
it reads two stones together, never one.

\emph{{[}It stands beside the slab, off the black, on small feet: two
shallow cups in a brass cradle, one lens across both, and under each cup
a little tilting tray. In the left tray lies a brass token, in the right
a copper one, each face-down, blank side up. A brass dish waits beneath
the trays' lips.{]}}

\textbf{Crab:} When both cups are full it chimes and tips its trays, and
a token slides from each into the dish and lands face-up or face-down.
That is its whole answer.

\textbf{Achilles:} Then the brass token is the left stone's answer and
the copper token is the right one's. Two stones, two tokens --- a fact
apiece.

\textbf{Crab:} That is just what I thought, the day it came. It isn't
so. Neither token belongs to either stone; both of them tell you only
how the two stones stand \emph{together}.

\textbf{Achilles:} Then let me see each stone by itself. Put your eye to
the lens while there's only the one in.

\textbf{Crab:} \emph{{[}with feeling{]}} Oh, I have wanted to, since the
day it came. But it will not read one stone alone --- leave one in by
itself till Easter and it only waits. And put your eye to the lens to
see a single stone, and the glass jams and reads nothing, and the stone
you peered at is spent for nothing. I have never once done it, and it
costs me every time I look at the thing.

\emph{{[}In the white light the glass is a lovely thing --- brass worked
as fine as a watch-case, a lens as clear as a drop of water, the little
trays set under the cups like a pair of cupped hands --- and it stands
so finely on its small feet that a heavy step on the boards makes the
trays shiver. Left to itself it keeps up a faint, thin singing somewhere
in its brass, like a wet finger going round the rim of a glass in
another room. When Achilles leans towards it the singing stops, and the
lens is only a round of glass; when he sits back, it takes up again.
Worn almost smooth on the foot of the cradle is a maker's stamp: two
sets of initials, struck one over the other.{]}}

\textbf{Tortoise:} Then what \emph{does} it say, when it speaks? You
have seen it used.

\textbf{Crab:} Not by me. By two of my correspondents, Alice and Bob.
They live a long way apart, and for years now the workshop's own carrier
--- slow, careful, once a season --- has brought them pairs, a half of
each. \emph{{[}nodding down the slab{]}} That is why I put that stone so
far off: let this end be Alice's, and that end Bob's.

\textbf{Crab:} Alice has a glass like this one. And they copy me every
letter that passes between them, and I keep them in here. \emph{{[}a fat
ledger, which he opens by the dish{]}} One line to a letter. Here is
what Alice's glass has dropped, letter by letter. Watch the tokens.

\emph{{[}The Crab lifts the two tokens out of their trays into the dish,
both still blank side up. He sets a claw at the first line of the ledger
and turns the brass token over: it rocks onto its back and shows its
stamp, face-up, while the copper token lies exactly as it was. One token
has changed; the other has held still.{]}}

\textbf{Achilles:} Brass up, copper down.

\emph{{[}At the next line the Crab turns the brass back down and the
copper over: now it is the copper that shows its stamp, and the brass
lies blank. At the line after, he turns the brass up again beside it,
and both lie stamped. At the next, both go down together, and the dish
looks as it did at the start.{]}}

\textbf{Achilles:} Copper up. Both up. Both down. \emph{{[}The claw goes
on down the page, and the tokens go on turning under it, one or both or
neither at each line.{]}} Brass up. Both down. Copper up. Copper up.
Both up. Brass up\ldots{}

\textbf{Achilles:} \emph{{[}ticking the margin of the Crab's newspaper
with a pencil stub{]}} That's twenty letters: five of the one, six of
another, four and five of the rest. Four ways the tokens can lie, and
every one of them turns up about a quarter of the time. Never a fifth
way, and never two at once.

\emph{{[}The Crab, pleased with him, refills his cup without being
asked, fetches the toast from the hearth and sets the plate at his
elbow. Achilles takes a slice and eats it one-handed, his eyes still on
the dish.{]}}

\textbf{Crab:} \emph{{[}turning both tokens blank side up again and
tipping them back into their trays{]}} And each way means something.
Both down: the two stones stand together just as they were cut. Brass
up: as though one of them had been flipped --- turned over, top to
bottom, like this. \emph{{[}His empty claw turns over in the air, about
a level line.{]}} Copper up: as though one had been spun --- turned half
round where it stands. \emph{{[}His claw turns half round on the spot,
about its upright.{]}} Both up: both.

\textbf{Achilles:} One of them. Which one?

\textbf{Crab:} It never says.

\textbf{Achilles:} Because it can't see. But one of them \emph{was}
flipped --- Alice's or Bob's. That's a fact about one stone, and your
glass is only too clumsy to report it.

\emph{{[}The breath-sized clear place in the window has feathered shut
while nobody was watching it. The pane is white now from edge to edge,
and whatever the lane is doing this morning, the parlour cannot see it:
only the brightness comes through, the same from every pane.{]}}

\textbf{Tortoise:} May I draw it? Not on the black, Mr Crab --- I'll not
lay string on your glass, where everything comes out twice. \emph{{[}She
spreads a small drawing-cloth on the table beside the slab, running the
length of it.{]}} This is only a cloth to draw on. Tell us the tale, and
I'll draw.

\textbf{Crab:} Alice calculates; she has a column for everything. Bob
would rather draw.

\textbf{Tortoise:} Then I'll draw for Bob. \emph{{[}She takes a length
of string from her shell.{]}} This is my string, as in the autumn: a
pencil that can go round corners. A bend in it is a pair cut of one
piece, with an end at each stone. Let Bob draw curves, then: the one
bend between Alice's half and his.

\emph{{[}She lays the string on the cloth in one long bend: one end laid
down opposite the near end of the slab, where Alice's stone stands; the
string carried out along the cloth and round; its other end laid down
opposite the far end, where Bob's stone waits alone.{]}}

\textbf{Crab:} And Alice has a stone she wants Bob to have --- as it
might be that one there, my keepsake's twin. But stones don't survive
the ordinary post; the road looks at everything on it, a thousand times
between her gate and his. Letters survive it, being settled already.

\textbf{Tortoise:} Let Bob draw one more line: the stone Alice wants to
send, coming in to her. \emph{{[}She lays a second length on the cloth,
a straight run coming in from the left towards Alice's end of the bend
and stopping short of it.{]}} And now her glass. She sets the stone she
wants to send in one cup, and her own half of the pair in the other, and
it reads them together.

\textbf{Tortoise:} \emph{{[}bringing out a little pouch{]}} These are
some of Mr Crab's rings from the scullery beam, red and green, each with
its bead set exactly halfway round. A red ring on a string is a flip,
and a green ring is a spin. Let Bob draw every way that reading could
end.

\emph{{[}She lays four short bends in a row on the cloth, each a length
of string turned round on itself once. The first she leaves plain. Onto
one leg of the second she threads a red ring and pushes it down until it
sits close in by the bend; the bend itself does not move, only the ring
rides down to it. Onto the third, a green ring, in the same place. Onto
the fourth, a red ring close in by the bend and a green ring further out
on the same leg.{]}}

\textbf{Achilles:} A plain bend, a flip, a spin, and both. The four ways
the tokens lie. \emph{{[}He puts a finger by the second bend.{]}} And
there's my fact: the red ring is on one leg. Whichever stone that leg
goes to --- that's the one that was flipped.

\textbf{Tortoise:} Is it? \emph{{[}She puts a claw to the red ring.{]}}

\emph{{[}The red ring sits on one leg of the second bend, close in, its
bead halfway round. The Tortoise slides it along the string, into the
bend, round it, and out along the other leg. The string does not move
and neither end stirs; only the ring travels. It comes to rest on the
other leg, the same distance from the bend --- still red, its bead still
exactly halfway round. Nothing about it has changed but the leg it sits
on.{]}}

\textbf{Achilles:} It came round exactly as it went. Your bead in the
autumn came round end for end.

\textbf{Tortoise:} That bead was lopsided. A half-turn has nothing
lopsided about it to turn: one ring comes round just as it went. Two
together come round the other way about --- but one comes round
unchanged.

\emph{{[}On the hob the kettle has begun to sing, and its note lies so
close to the glass's thin singing that the two make between them a slow
throb, swelling and fading, that belongs to neither of them. Nobody in
the room could say which of the two is the higher. The Crab eases the
kettle a little off the heat; the throb slows and lengthens, and the
room goes on keeping it.{]}}

\textbf{Achilles:} Then the flip isn't on Alice's leg or on Bob's. It's
on the bend --- it slides to whichever leg you please. So the glass
isn't clumsy at all; there was never a fact about one stone for it to
report.

\textbf{Crab:} I told you so. Both tokens, about the two together.

\textbf{Achilles:} Four answers about how two stones stand together. So
if Alice reads \emph{her} stone against the one she wants to send, what
does Bob's stone, a whole house away, know about it?

\begin{center}\rule{0.5\linewidth}{0.5pt}\end{center}

\hypertarget{second-challenge-two-bits-early}{%
\section{\texorpdfstring{Second Challenge --- \emph{Two Bits
Early}}{Second Challenge --- Two Bits Early}}\label{second-challenge-two-bits-early}}

\textbf{Tortoise:} Let us put the tale together and see.

\emph{{[}She takes up her longest string and lays the whole tale on the
cloth in one piece. In from the left it comes, the message; at Alice's
end it turns round in a bend, and onto its incoming leg she moves the
two rings from her fourth little bend, the red close in by the bend and
the green further out. From there the string runs back towards the left,
round a second bend facing the other way --- the pair between Alice and
Bob --- and out along the far leg to Bob's end, where it stops opposite
his stone. Laid flat, it makes an S: in with the message, round Alice's
reading, back, round the pair, and out to Bob.{]}}

\textbf{Achilles:} Let me work it through. The stone Alice sends goes
into her glass with her half of the pair, and both are read there. Bob's
stone is the other half of a pair she's just used up. So Bob has lost
his partner and never had a thing to do with the message. Nothing
reaches Bob until the letter does.

\emph{{[}The fire has begun on the window from below. On the lowest
panes the ferns have gone glassy at their tips, half frost and half
water and each tangled into the next, and through them the lane comes in
as a wavering blur of white and dark that might be hedges or might be
anything. Above, the frost holds, white and whole.{]}}

\textbf{Tortoise:} Take the message end. I'll take Bob's. Pull.

\emph{{[}Achilles takes the loose end on the left; the Tortoise takes
the loose end at Bob's. They draw apart. The S tightens; the bend at
Alice's end and the bend of the pair slide towards each other, meet, and
run out through one another into nothing. The rings ride with the
string, never lifted and never passing each other. What was in, round,
back, round and out is now one taut line from Achilles's hand to the
Tortoise's --- from the message straight to Bob's end --- and on it,
still, a green ring nearer the message and a red ring nearer Bob. The
string is no longer and no shorter than it was; only the bends are
gone.{]}}

\textbf{Achilles:} It's at Bob's end. The message runs straight through
to him --- and he hasn't had a letter. Alice has only just read.

\textbf{Tortoise:} And what is on it?

\textbf{Achilles:} Her rings. A spin and a flip. It's there, but turned
--- spun and flipped, the way her tokens fell.

\textbf{Crab:} \emph{{[}half out of his chair, a claw reaching towards
the far end of the slab{]}} Then Bob's stone --- this very minute --- if
I only---

\textbf{Tortoise:} Nothing has been sent on this table, Mr Crab. And if
it had, one look would show you a coin, and spend your last pair.

\textbf{Crab:} \emph{{[}sitting back{]}} A coin. Yes. \emph{{[}He turns
pages in the ledger.{]}} Bob was no better at waiting than I am. Here
--- many a morning he peeked at his stone through his viewer before the
letter came.

\textbf{Achilles:} \emph{{[}reading down the column{]}} High. Low. Low.
High. Low. High\ldots{} And here, where Alice had sent something quite
different --- high, low, high, low. A coin toss, every time, whatever
Alice had sent.

\textbf{Tortoise:} Which bend was she on, Achilles --- the plain one,
the flip, the spin, or both? Bob can't know. Four ways turned, and all
four together look like a coin.

\emph{{[}Against the teapot, with the white window behind it, the letter
from Ostend has gone half transparent in the glare, and the cousin's
hand shows through envelope and folded sheet alike --- the whole of it
at once. But every fold has laid the writing down a different way: some
of it the right way round, some mirror-wise, some upside down in its
mirror, and some simply standing on its head; and laid one over another
the lines make only an even grey lattice that no one could read a word
of. It is all there. It says nothing at all.{]}}

\textbf{Achilles:} \ldots Oh. It's there before the letter --- and it's
no use to him.

\emph{{[}He sits back in his chair and blows on his tea, and for a
moment nobody says anything. The kettle murmurs on the hob; the Tortoise
draws her quilt up round her.{]}}

\textbf{Achilles:} Then what does the letter carry? Not the stone ---
two tokens?

\begin{center}\rule{0.5\linewidth}{0.5pt}\end{center}

\hypertarget{third-challenge-the-letter}{%
\section{\texorpdfstring{Third Challenge --- \emph{The
Letter}}{Third Challenge --- The Letter}}\label{third-challenge-the-letter}}

\emph{{[}The lower half of the window has given in to the fire. Its
small panes are clear now, clean glass from bar to bar, each with its
own sharp little picture of the lane --- the white road, the hedge-tops
heavy with rime, the stable roof across the way --- while the upper
panes still keep their ferns, and the frost stops exactly at the wood of
the bars, as tidy as a border in a book. The Crab tops up the pot from
the kettle, sets it back by the letters, and tucks its cosy round
it.{]}}

\textbf{Crab:} Two tokens, just as they fell; that's all Alice ever
posts. \emph{{[}He shakes his cousin's letter out of its envelope and
holds the empty envelope up.{]}} This will do for a letter. The tokens
go in exactly as they lay, the flap goes down, and the road can look at
them all it likes --- they're settled already, and looking can't
unsettle them.

\textbf{Crab:} And she writes the same two words on the back of every
one, in case Bob forgets. \emph{{[}He writes across the back of the
envelope: brass first.{]}}

\emph{{[}The Crab tips the two tokens out of their trays into the dish;
they land as they lay, both blank side up.{]}}

\textbf{Crab:} Here's one of hers from the autumn. \emph{{[}At the
ledger's line he turns the brass token face-up in the dish, and then the
copper: both lie stamped. He slides them into the envelope, faces as
they are, and folds the flap down over them.{]}} Postman?

\textbf{Achilles:} \emph{{[}taking it{]}} The long way round, I suppose
--- not over the black.

\emph{{[}He carries the envelope round the edge of the table: past the
window end, down the long side, never once over the slab, and lays it on
the cloth at Bob's end. Inside, the tokens have not stirred: brass up,
copper up, just as they were laid.{]}}

\textbf{Achilles:} \emph{{[}opening it{]}} Brass up: flip. Copper up:
spin. And on the back, brass first.

\textbf{Tortoise:} Let Bob draw corrections, then, just as the letter
says.

\emph{{[}The straight string still lies as the pull left it: from the
message end, a green ring, then a red. Now the Tortoise threads two more
on from Bob's end. First a red ring, which she slides up the string
until it sits a little way below Alice's red; then a green ring, which
she leaves nearer Bob's hand. Along the whole string, from the message
to Bob: green, red, red, green --- Alice's two and Bob's two, facing
each other across the middle.{]}}

\textbf{Tortoise:} Let Bob draw kin to kin.

\emph{{[}She slides Alice's red ring down the string towards Bob's red.
The string holds still; only the ring travels, and there is nothing in
its way --- no green between the two reds. They touch, and run together
into one ring: two beads, each halfway round, come to rest as one bead
gone all the way round, back at the top, where no turn at all sits. A
red ring with no turn on it, one string in and one out, is only string;
the Tortoise lifts it off. Now Alice's green has a clear run down to
Bob's green. She slides it; they meet; the two halves make a whole; she
lifts that off too. The string lies bare from the message end to Bob's
end --- nothing on it at all, and no longer or shorter than before.{]}}

\textbf{Tortoise:} A whole turn is no turn.

\emph{{[}Across the lane the weathercock on the stable roof, which froze
in the night with its tail to the weather, is loosed by the first of the
day's warmth and caught by the morning's first small breeze. It goes
round with a creak that carries clear over the snow --- half round ---
and stops facing exactly as it faced yesterday, as though it had never
been turned at all.{]}}

\textbf{Achilles:} Straight from the message to Bob, and nothing on it.
He has it exactly.

\textbf{Achilles:} But brass first --- must it be? A flip and a spin are
a flip and a spin, whichever you do first.

\textbf{Tortoise:} His stone could never tell the difference. But on the
string, only this way round does each ring meet its own kind; spin
first, and a green would have to get past a red to reach its partner.

\emph{{[}Achilles brings the envelope back round the table's edge to the
near end, and the Crab tips the two tokens out of it into the dish. They
slide out just as they went in, brass up, copper up; only their place
has changed.{]}}

\textbf{Crab:} And the other letters? \emph{{[}He lays out three more
ledger lines in the dish, one after another, and at each the Tortoise
re-threads the string: Alice's rings as the tokens fell, Bob's as the
letter says.{]}}

\emph{{[}Brass alone: one red ring left on the string by Alice's
reading, one red put on by Bob; they slide together, make their whole
turn, and come off. Copper alone: green to green, and off. Both down:
nothing on the string from Alice, nothing for Bob to put on, and the
string bare from the start. Every pattern, and every time the string
lies clean from the message to Bob.{]}}

\textbf{Crab:} \emph{{[}turning a page{]}} And here are the mornings
Alice set her stone herself, some way of her own choosing, and wrote
afterwards to tell Bob how. When he had turned his stone as the letter
said, he looked at it that way through his viewer. It matched. Every
time, whatever the tokens were.

\textbf{Crab:} \emph{{[}closing the ledger with both claws{]}} I am so
glad. I always liked Bob. And the letter went by the post, and the stone
never had to.

\emph{{[}The Crab turns both tokens blank side up --- the last letter
left them so already --- and tips them back into the glass's trays, one
to each, where they lay when the glass was unwrapped.{]}}

\emph{{[}He blows his nose, rather loudly, into a large spotted
handkerchief, and the Tortoise passes him the toast. Achilles laughs
under his breath and helps himself to the last slice before anyone can
stop him.{]}}

\textbf{Achilles:} Two tokens to carry one stone's worth. Could it go
the other way --- one stone to carry two tokens' worth?

\begin{center}\rule{0.5\linewidth}{0.5pt}\end{center}

\hypertarget{fourth-challenge-the-other-way-round}{%
\section{\texorpdfstring{Fourth Challenge --- \emph{The Other Way
Round}}{Fourth Challenge --- The Other Way Round}}\label{fourth-challenge-the-other-way-round}}

\textbf{Crab:} It can, and they did --- last spring, when Alice had two
tokens' worth of news and nothing but a stone to send it by. \emph{{[}He
tips the tokens out of their trays into the dish, both blank side
up.{]}} She set her news out first, by hand.

\emph{{[}Both tokens lie blank. The Crab, for Alice, turns the copper
one over: it shows its stamp, and the brass lies blank beside it.
Nothing has fallen from any glass; this time the tokens were chosen.{]}}

\textbf{Crab:} Then she turned her half of the pair as her tokens said
--- flip if the brass was up, then spin if the copper was. So: no flip,
and a spin. And she sent Bob the \emph{stone}, by the workshop's
carrier, slow and careful.

\emph{{[}Through the clear lower panes the milkman can be seen in the
lane, on foot this morning because of the snow, stopping at the house
opposite to look at the card in its front window: a square of pasteboard
that the house turns over, or turns round, or both, or leaves as it is,
to tell him the whole of what it wants without a word written. He reads
it at a glance, by the way it was agreed between them long ago, and
takes his cans up the path.{]}}

\textbf{Achilles:} And Bob?

\textbf{Crab:} Bob set her stone and his own in his glass, and it read
them together. \emph{{[}He finds Bob's line in the ledger, reads it, and
does not touch the tokens: they already lie as Bob's glass dropped
them.{]}} Copper up, brass down. The very tokens she chose.
\emph{{[}turning pages{]}} And the next letter, and the next ---
whatever she set, his glass dropped the same. Never once otherwise.

\textbf{Achilles:} Flip if the brass is up, then spin if the copper is.
That's exactly what Bob did with the letter. \ldots Oh. It's the
letter's own table. The half-turns Bob used to mend, Alice uses to
write.

\emph{{[}As he says it a bright fleck of sun, thrown in off the
milkman's cans as he comes back down the path opposite, slides across
the wall and over the glass of the print, and for a breath the upright
mirror in it is lit as if from within. Then the fleck has gone on, and
the wall is only white again.{]}}

\emph{{[}Achilles glances up at the print on the wall facing the slab
--- Escher's \textup{Magic Mirror}: a mirror standing upright, and
creatures stepping out of it.{]}}

\textbf{Achilles:} Stepping out of the looking-glass, as solid as
anything that went in.

\textbf{Tortoise:} \emph{{[}taking up a length of string{]}} I'll draw
this one myself. A loop of string has no loose end, so nothing comes out
of it but a yes or a no.

\emph{{[}She lays a bend on the cloth --- the pair, Alice's leg along
the top and Bob's along the bottom. On Alice's leg she threads Alice's
writing: no red, for the brass lay down, and a green ring for the
copper. Further along the same leg she threads the answer Bob's glass is
being asked about --- copper up, brass down: one green ring. Then she
brings the far end round in a second bend, back to meet Bob's leg, and
ties it off. A closed loop of string, no loose end anywhere, with two
green rings side by side on it.{]}}

\textbf{Tortoise:} Kin to kin.

\emph{{[}She slides the answer's green ring along the string to Alice's
green. They touch; the two half-turns make a whole; she lifts the ring
off. The loop lies bare, closed, nothing on it anywhere.{]}}

\textbf{Tortoise:} A clean loop: a sure yes. Now ask Bob's glass about
brass up instead.

\emph{{[}She threads Alice's green back on, and in place of the answer's
green she puts a red ring, close in by the far bend. Then she slides the
red round the loop --- round the far bend, along Bob's leg, round the
near bend and up Alice's leg --- until it sits beside the green. It is
still red, its bead still halfway round, with no red anywhere to meet
it; and the green has no green to meet. Nothing comes off.{]}}

\textbf{Tortoise:} A loop with a ring left over is a no: that answer
never comes. Whatever Alice writes, only her own answer can.

\emph{{[}Her quilt has slipped off one shoulder while she worked; the
Crab, passing behind her chair to see the loop better, puts it back
without a word, and she pats his claw. The fire snaps and settles. Above
them the frost on the upper panes has thinned to a grey lace, with blue
sky showing through the holes in it.{]}}

\textbf{Achilles:} \emph{{[}looking from one bend of the loop to the
other{]}} And the two bends. That one is Alice writing, and this one is
Bob's glass reading --- the same bend with the same rings, only facing
back the other way.

\textbf{Crab:} Don't think it came for nothing, mind. The carrier took
Bob his half of the pair long before; one stone travelled then and one
now. Two stones for two tokens' worth.

\textbf{Achilles:} And if the stone Alice sends is \emph{itself} half of
a pair?

\begin{center}\rule{0.5\linewidth}{0.5pt}\end{center}

\hypertarget{fifth-challenge-the-mended-junction}{%
\section{\texorpdfstring{Fifth Challenge --- \emph{The Mended
Junction}}{Fifth Challenge --- The Mended Junction}}\label{fifth-challenge-the-mended-junction}}

\textbf{Crab:} \emph{{[}a long look at the near end of the slab: the
cloth, and the bare stone beside it{]}} Then it's this one. My
keepsake's twin --- half of a pair, and the other half under that cloth.
\emph{{[}He straightens.{]}} We'll do it for real, on the table, with
the last of the drawer. This end is Alice's. You shall be Bob, Madam
Tortoise, down at his end, and do exactly what Bob would do. And
Achilles knows the road.

\textbf{Tortoise:} \emph{{[}going down the long side of the table to the
far end and settling there beside the lone stone, the drawing-cloth
running along the slab beside her{]}} I'm at home.

\emph{{[}The Crab has Achilles carry her cup down to her the long way
round, by the table's edge, and she sets it at her elbow on the table
beside her cloth, well clear of the black. Down there at the far end,
with her quilt round her and her tea beside her, she looks very small
and entirely comfortable.{]}}

\emph{{[}The tokens still lie in the dish as Alice set them: copper up,
brass down. The Crab turns the copper one back onto its face, so that
both lie blank side up again, and tips them into the glass's trays,
brass to the left, copper to the right, where they lay at the start of
the morning. The envelope, empty, lies ready beside the dish.{]}}

\textbf{Crab:} \emph{{[}putting a claw to the message stone, and
hesitating{]}} The message, and Alice's half. Into the glass.

\emph{{[}At the near end the message stone and Alice's stone lie on the
black, each with its reflection under it. The Crab lifts the message
stone and sets it in the left cup of the glass; as it leaves the black
its reflection sinks away beneath it and is gone. The glass waits. He
lifts Alice's stone into the right cup, and its reflection goes the same
way. Under the one lens the two sit side by side, and the moment the
second settles the glass chimes, clear and small. Both trays tip at
once; the tokens slide off and ring down into the dish --- the brass
face-up, the copper face-up --- and lie there, off the black, casting
nothing. In the cups the two stones sit still: read together, and spent.
At the far end Bob's stone has not moved, and nobody has touched it;
nothing about it shows anything at all.{]}}

\textbf{Tortoise:} The letter never went, and then the glass chimed.

\textbf{Achilles:} It hasn't gone --- I haven't touched it. And the
glass chimed first; we all heard it. You've said it back to front, or
inside out, or both --- and you can't tell which, down there.

\textbf{Tortoise:} Not from here.

\textbf{Crab:} \emph{{[}very quietly, looking down the length of the
slab{]}} It's there now, isn't it. Down at that end, this very minute
--- my keepsake's partner already, only turned. And I cannot so much as
glance at it.

\textbf{Tortoise:} One glance and you'd see a coin, and have a line in
your ledger where you had a pair.

\textbf{Crab:} \emph{{[}folding his claws together, and keeping them
folded{]}} The only pair ever made by post in this house. I shan't look.

\emph{{[}The glass stands at the near end with its work done. The small
singing in its brass has not come back since the chime: it is a handsome
thing of brass and ground glass with its trays tipped empty, no more,
and the room is a little plainer for it. Nobody looks down the table.
The Crab keeps his eyes on the fire, and his claws stay folded.{]}}

\textbf{Achilles:} Then the post had better come.

\emph{{[}The tokens lie in the dish where they fell, brass up, copper
up. Achilles lifts them out one at a time by their edges, keeping each
on exactly the face it fell on, and lays them in the envelope side by
side, above the two words on its back. He folds the flap. Nothing about
them has changed but where they are.{]}}

\textbf{Achilles:} The glass chimed, and then the letter went.

\emph{{[}He carries it the long way round: along the table's edge, round
the window end, down the far side, never once over the black. At Bob's
end he lays it on the cloth in front of the Tortoise.{]}}

\textbf{Achilles:} \emph{{[}opening it{]}} Flip. Spin.

\emph{{[}Bob's stone sits at the far end as it has sat all morning, its
reflection beneath it. The Tortoise keeps her eyes on the two tokens in
the open envelope and puts out a claw to the stone by feel. She turns it
over, top to bottom, about a level line through its middle: what lay
uppermost goes down against the black, and what lay against the black
comes up. Beneath it the reflection turns over too, mirror-fashion, and
meets it again face to face at the surface. The stone has not left its
spot and no one has looked at it; only which way up it lies has
changed.{]}}

\emph{{[}Then, still by feel, she turns it half round where it stands,
about its upright: the side that faced the window now faces the room,
and its top stays its top. Its reflection turns the same half round
beneath it. The stone is where it was, and still unread; only the way it
faces has changed. She tips the two tokens out of the envelope onto the
cloth beside her, faces as they fell, brass up, copper up; the envelope
is empty now.{]}}

\textbf{Tortoise:} The glass chimed, and then the letter went.

\emph{{[}At the window the last ferns on the upper panes have begun to
burn gold along their top edges: the sun, coming round the corner of the
house, is almost on the glass. The whole room lightens by a shade, as if
someone had opened a door somewhere.{]}}

\textbf{Tortoise:} Let Bob draw end to end. \emph{{[}She lays her string
straight along the drawing-cloth, from opposite the keepsake under its
cloth at the near end, the whole length of the slab, to opposite Bob's
stone at the far: one line, no bend in it, nothing on it.{]}}

\textbf{Achilles:} Your keepsake, and Bob's stone. One piece now --- and
they were never in the same hand, never even at the same end of the
room. Nothing crossed the room but a letter. And the answer was waiting
there before it came.

\emph{{[}Out in the garden, between the house and the lane, the old pear
against the wall stands up black out of the snow, and among its own
boughs is one that came from a tree clear across the town, set into it
by some grafter in another winter. The join grew over long ago. Nobody
could find now where the one tree stops and the other begins; in the
spring the whole of it comes into leaf together.{]}}

\textbf{Crab:} \emph{{[}opening the ledger once more, at the back{]}}
Alice and Bob made pairs like that for their friends, now and then, the
same way. Whatever tokens fell, once the letter was obeyed, every one of
them read afterwards as one piece. The winter the letters went astray in
the snow, the pairs read like two strangers' coins.

\textbf{Crab:} \emph{{[}closing it{]}} And these two I shall never read.
The last pair in the house, and it never met. \emph{{[}He goes to the
dresser and looks into the drawer, which is empty, and shuts it.{]}}

\textbf{Tortoise:} \emph{{[}coming back up the long side of the
table{]}} Then you have a thing worth more than a line in a book, Mr
Crab. Most people only have the book.

\emph{{[}The sun has come round the corner of the house at last and
stands full on the window, low and gold, and the ferns on the upper
panes darken at their edges and begin to gather themselves into beads of
water, pane after pane, until the lane shows through the whole window,
dazzling. Warmth comes in with the light, and for the first time this
morning the room has a sunny side and a cool one.{]}}

\textbf{Crab:} There's tea left. Hot, too --- look, the sun's got round
to the window, and your frost is running, Madam Tortoise.

\textbf{Tortoise:} It will draw it again tonight. I'll take more
lessons.

\emph{{[}The Crab pours; Achilles, who has been standing at Bob's end
with the empty envelope in his hand, comes back to his chair and takes
his cup. For a while nobody reaches for anything on the table.{]}}

\textbf{Achilles:} It was there before the letter came.

\textbf{Tortoise:} And no use to anyone until it did.

\textbf{Achilles:} \emph{{[}turning the empty envelope over, to the two
words on its back{]}} Then what was it I carried?

\emph{{[}Nobody answers him, and for the moment nobody needs to. The sun
lies along the near end of the room, low and warm, and the window is
clear from top to bottom, a few beads of water still standing on the
glass; beyond it the lane is white and quiet, and the postman's
footprints go away down it towards the next gate. On the black the
keepsake keeps to its cloth at the near end, and far down at the other
end Bob's stone sits alone; neither has been read, and neither will be.
Beside the slab the pairing-glass stands silent, the two stones that
went into it still in its cups, off the black; on the cloth at the far
end the brass token and the copper lie out of their envelope, off the
black. The mittens on the fender have dried. On the wall facing the slab
the print hangs in the ordinary daylight --- a mirror, and its
creatures, in a plain frame --- one more thing in a sunny room. Achilles
sits with the empty envelope on his knee, the Crab fills the Tortoise's
cup and then his own, and the kettle, set back from the fire, ticks as
it cools; the frost, which will be back tonight, has left the window to
the sun.{]}}

\hypertarget{chapter-10-teleportation}{%
\chapter{Chapter 10: Teleportation}\label{chapter-10-teleportation}}

\emph{--- Out-of-Order Canon ---}

\emph{Part III: Pictures Become Proofs}

\emph{Escher: ``Magic Mirror'' --- a mirror standing upright, and
creatures stepping out of it}

\begin{center}\rule{0.5\linewidth}{0.5pt}\end{center}

Can a thing that nobody is allowed to look at be sent to someone far
away, when the only thing the post will carry is a short letter?

Suppose you hold a single qubit in a state you do not know, and a friend
in another city needs it. You cannot look at it and write down what you
saw. Chapter 9 taught us what a look does: it forces one answer out of
several, at random, and destroys the rest, so what you wrote down would
be a coin toss and not the qubit. Nor can you put the qubit itself in
the post. Say the road is rough, and anything quantum that travels it is
looked at a thousand times on the way. What you can send is a letter: a
few definite marks on paper, which no amount of looking can spoil,
because they are settled already. And the qubit is not a few marks. It
is a pair of complex numbers, a point anywhere on a sphere; to write it
out exactly would take digits without end.

So the task looks impossible, and it is exactly what this chapter does.
It needs one extra thing, arranged long before: a shared pair, one half
each, the bent wire of Chapter 6. With that in hand, a letter carrying
\emph{two bits} is enough.

The book has owed this chapter for some time, and it owes it three
promises. In Chapter 6 a cap was a postselected outcome, ``the answer
`\(\Phi^+\)' to a question with four possible answers'', and we wrote:
``Chapter 10 will show what to do about the other three, and the result
will be teleportation.'' At the end of the same chapter, joining two
pairs at a junction worked one time in four, and we said that each of
the other answers ``leaves Alice and Bob with a pair that is a known
one-box correction away from \(\ket{\Phi^+}\) \ldots{} That is Chapter
10.'' And Chapter 9 closed on the recipe itself: thread a cup, keep the
outcome, send it across, turn the far half by the correction it names,
``and the postselected cap of Chapter 6 becomes a protocol that works
every time.''

There is one strange fact to watch for. The moment the question is asked
at the near end, the message is \emph{already} at the far end, before
any letter has been written. It arrives turned, though, by one of four
half-turns, and nobody at the far end can tell which. The letter comes
afterwards and names the half-turn to undo. The answer arrives before
its cue, and the cue sets it right. We shall see why it cannot be used
early, and why, once the letter arrives, the picture proof is a wire
pulled straight and two pairs of rings lifted off.

\begin{center}\rule{0.5\linewidth}{0.5pt}\end{center}

\hypertarget{four-answers-to-one-question}{%
\section{§10.1 Four Answers to One
Question}\label{four-answers-to-one-question}}

Chapter 6's cap asked two wires one question --- ``are you
\(\ket{\Phi^+}\)?'' --- and kept only the answer yes. A real device
asked about two wires must say \emph{something} every time. What are the
other things it can say?

\textbf{On the blackboard.} We need two of the Pauli matrices of Chapter
3: the bit-flip
\(X=\bigl(\begin{smallmatrix}0&1\\1&0\end{smallmatrix}\bigr)\) and the
phase-flip
\(Z=\bigl(\begin{smallmatrix}1&0\\0&-1\end{smallmatrix}\bigr)\). Each is
its own inverse, \(X^2=Z^2=I\), and each is real and symmetric. For a
bit \(m\in\{0,1\}\) write \(X^m\) and \(Z^m\) for ``do it if \(m=1\)'':
\(X^0=Z^0=I\), \(X^1=X\), \(Z^1=Z\). The two do not commute; one line of
multiplication gives \(XZ=-ZX\).

Chapter 6 built \(\ket{\Phi^+}\) as
\(\mathrm{CNOT}\,(H\otimes I)\ket{00}\). Feed the same little circuit
the other three inputs. For two bits \(m_1,m_2\) define

\[\ket{B_{m_1m_2}} \;=\; \mathrm{CNOT}\,(H\otimes I)\,\ket{m_1m_2}.\]

Take \(\ket{11}\). The Hadamard turns the top \(\ket1\) into \(\ket-\),
giving \(\tfrac1{\sqrt2}(\ket{01}-\ket{11})\); the CNOT flips the bottom
wire where the top reads \(1\), giving
\(\tfrac1{\sqrt2}(\ket{01}-\ket{10})\). The four together:

\[\begin{aligned}
\ket{B_{00}} &= \tfrac1{\sqrt2}\bigl(\ket{00}+\ket{11}\bigr)=\ket{\Phi^+}, &\qquad
\ket{B_{01}} &= \tfrac1{\sqrt2}\bigl(\ket{01}+\ket{10}\bigr)=\ket{\Psi^+},\\
\ket{B_{10}} &= \tfrac1{\sqrt2}\bigl(\ket{00}-\ket{11}\bigr)=\ket{\Phi^-}, &\qquad
\ket{B_{11}} &= \tfrac1{\sqrt2}\bigl(\ket{01}-\ket{10}\bigr)=\ket{\Psi^-}.
\end{aligned}\]

These are the four \textbf{Bell states}. Now look at them against
\(\ket{\Phi^+}\). \(\ket{\Psi^+}\) is \(\ket{\Phi^+}\) with the top bit
flipped; \(\ket{\Phi^-}\) is \(\ket{\Phi^+}\) with a sign on the
\(\ket{11}\) branch, a phase-flip on the top; and \(\ket{\Psi^-}\) is
both. Check the last:
\((ZX\otimes I)\tfrac1{\sqrt2}(\ket{00}+\ket{11})=\tfrac1{\sqrt2}(Z\ket1\otimes\ket0+Z\ket0\otimes\ket1)=\tfrac1{\sqrt2}(\ket{01}-\ket{10})\).
Each of the four is the first one with a Pauli on one leg.

\begin{quote}
\textbf{Result 10.1 (the Bell basis).} For bits \(m_1,m_2\), exactly,
\[\begin{aligned}\ket{B_{m_1m_2}} \;&=\; \mathrm{CNOT}\,(H\otimes I)\ket{m_1m_2}\\ &=\; \bigl(Z^{m_1}X^{m_2}\otimes I\bigr)\ket{\Phi^+} \;=\; \bigl(I\otimes X^{m_2}Z^{m_1}\bigr)\ket{\Phi^+}.\end{aligned}\]
The four are orthonormal and complete: \(\braket{B_j|B_k}=\delta_{jk}\)
and \(\sum_k\ket{B_k}\bra{B_k}=I\).
\end{quote}

\emph{Proof idea.} The circuit is a unitary, so it carries the four
orthonormal inputs \(\ket{m_1m_2}\) to four orthonormal outputs, and
four orthonormal vectors in a four-dimensional space leave no room for a
fifth. The Pauli forms are the table read off, and the bottom-leg form
is Chapter 6's slide.

\emph{Proof.} The first equality is the definition, and the second is
the table above, checked in all four cases. For the third, slide the box
off the top leg of \(\ket{\Phi^+}=\tfrac1{\sqrt2}\,\)cup (Theorem 6.4):
it arrives on the bottom leg transposed, and
\((Z^{m_1}X^{m_2})^{T}=(X^{m_2})^{T}(Z^{m_1})^{T}=X^{m_2}Z^{m_1}\),
since \(X\) and \(Z\) are symmetric. Orthonormality:
\(\mathrm{CNOT}\,(H\otimes I)\) is unitary. Completeness: four
orthonormal vectors in a four-dimensional space are a basis.
\(\;\blacksquare\) See \texttt{Code/Ch10/FourAnswers.py}.

A \textbf{Bell reading} is the question ``which of the four?''. Its four
answers are the four effects that are the mirror images of the Bell
states (Chapter 3: the mirror is the dagger, the conjugate transpose).
Taking the dagger of \((Z^{m_1}X^{m_2}\otimes I)\ket{\Phi^+}\) reverses
the order of the product, and the Paulis are their own daggers, so

\[\bra{B_{m_1m_2}} \;=\; \bra{\Phi^+}\bigl(X^{m_2}Z^{m_1}\otimes I\bigr) \;=\; \bra{m_1m_2}\,(H\otimes I)\,\mathrm{CNOT}.\]

The second form says how to build the reading: run the circuit backwards
(CNOT, then \(H\) on the top wire) and then look at both wires with the
plain look of Chapter 9. The bit \(m_1\) is read off the Hadamarded top
wire, \(m_2\) off the bottom. By the Born rule (§9.1), a two-wire state
\(\ket\Omega\) gives answer \(m_1m_2\) with probability
\(|\braket{B_{m_1m_2}|\Omega}|^2\). Chapter 6's cap was
\(\sqrt2\,\bra{B_{00}}\), the first of these four, and nothing more.

\textbf{\emph{Why this definition.}} Why define the Bell states by a
circuit, when we could simply list four vectors? Because a definition
should say how to build the thing, and run backwards this circuit
\emph{is} the reading. And the four vectors it produces turn out to be
Chapter 6's one bent wire with a Pauli on one leg. That is what will
make every one of the ``wrong'' answers correctable.

\textbf{In the sketchbook.} Chapter 8 gave us the two Paulis as dots. A
red dot with one leg in, one leg out and phase \(\pi\) is exactly \(X\),
and a green one is exactly \(Z\) (§8.2, Dirac's Blackboard). We now need
a small new mark: a \textbf{ring whose phase is a bit times \(\pi\)}.
Write \(m\pi\) beside a green or red dot on a wire, with \(m\) a bit
left open. For \(m=0\) the phase is \(0\), and a phase-free dot with one
leg in and one out is a bare wire (§8.1); for \(m=1\) it is the
half-turn. So one drawing stands for two, exactly:

\[\begin{tikzpicture}[sd, x=1.4em, y=1.4em,
  zsp/.style={circle,draw,line width=0.9pt,fill=green!55,minimum size=1.1em,inner sep=0pt},
  ph/.style={font=\scriptsize,inner sep=1pt}]
  \draw[wire] (0,0) -- (2.6,0);
  \node[zsp,label={[ph]above:{$m\pi$}}] at (1.3,0) {};
\end{tikzpicture}
\;=\;Z^{m},
\qquad\qquad
\begin{tikzpicture}[sd, x=1.4em, y=1.4em,
  xsp/.style={circle,draw,line width=0.9pt,fill=red!75,minimum size=1.1em,inner sep=0pt},
  ph/.style={font=\scriptsize,inner sep=1pt}]
  \draw[wire] (0,0) -- (2.6,0);
  \node[xsp,label={[ph]above:{$m\pi$}}] at (1.3,0) {};
\end{tikzpicture}
\;=\;X^{m}.\]

A picture with two such rings stands for four pictures. Labelling a
phase by an unknown bit is how van de Wetering's survey draws
teleportation (§5.4), and it is what lets one drawing carry all four
answers of a Bell reading.

With the ring in hand, Result 10.1 draws itself. The cup is, as always,
unnormalized, cup \(=\ket{00}+\ket{11}\), so
\(\ket{\Phi^+}=\tfrac1{\sqrt2}\,\)cup; and a Bell state is that cup with
a red ring \(m_2\pi\) and then a green ring \(m_1\pi\) on its top leg,
read outward from the bend:

\[\ket{B_{m_1m_2}}\;=\;\tfrac1{\sqrt2}\;
\begin{tikzpicture}[sd, x=1.4em, y=1.4em,
  zsp/.style={circle,draw,line width=0.9pt,fill=green!55,minimum size=1.1em,inner sep=0pt},
  xsp/.style={circle,draw,line width=0.9pt,fill=red!75,minimum size=1.1em,inner sep=0pt},
  ph/.style={font=\scriptsize,inner sep=1pt}]
  \sdcup{0}{1.5}{0}
  \draw[wire] (0,1.5) -- (4.6,1.5);
  \draw[wire] (0,0) -- (4.6,0);
  \node[xsp,label={[ph]above:{$m_2\pi$}}] at (1.4,1.5) {};
  \node[zsp,label={[ph]above:{$m_1\pi$}}] at (3.0,1.5) {};
\end{tikzpicture}
\qquad
\bra{B_{m_1m_2}}\;=\;\tfrac1{\sqrt2}\;
\begin{tikzpicture}[sd, x=1.4em, y=1.4em,
  zsp/.style={circle,draw,line width=0.9pt,fill=green!55,minimum size=1.1em,inner sep=0pt},
  xsp/.style={circle,draw,line width=0.9pt,fill=red!75,minimum size=1.1em,inner sep=0pt},
  ph/.style={font=\scriptsize,inner sep=1pt}]
  \draw[wire] (0,1.5) -- (4.6,1.5);
  \draw[wire] (0,0) -- (4.6,0);
  \sdcap{4.6}{1.5}{0}
  \node[zsp,label={[ph]above:{$m_1\pi$}}] at (1.6,1.5) {};
  \node[xsp,label={[ph]above:{$m_2\pi$}}] at (3.2,1.5) {};
\end{tikzpicture}\]

The effect on the right is the mirror image of the state on the left:
reflected left to right, so the cup becomes a cap and the order of the
rings is reversed, and every phase conjugated, which changes nothing
here, because \(e^{-i\pi}=e^{i\pi}=-1\). A half-turn ring is its own
mirror image.

\textbf{\emph{Why this definition.}} Why label a phase by a bit rather
than draw four pictures? Because then a proof drawn once is four proofs.
The only fact about the label that any rewrite will ever use is that two
of them add to \(2m\pi\), a whole turn, which is no turn at all.

\hypertarget{the-reading-drawn-from-its-circuit}{%
\subsection{The reading, drawn from its
circuit}\label{the-reading-drawn-from-its-circuit}}

We have \emph{asserted} that the reading's effect is a cap with two
rings. The reading we can build is the circuit
\(\bra{m_1m_2}(H\otimes I)\,\mathrm{CNOT}\). The claim that the two
agree is a real debt, and we pay it with Chapter 8's rules, keeping
every number.

Three translations, all from Chapter 8. The CNOT is a green dot on the
top wire tied by a cord to a red dot on the bottom, and the bare picture
is \(\mathrm{CNOT}/\sqrt2\) (row 64); so \(\mathrm{CNOT}\) is \(\sqrt2\)
times that picture. A red dot with one leg and phase \(0\) is the state
\(\sqrt2\,\ket0\), and with phase \(\pi\) it is \(\sqrt2\,\ket1\) (§8.2,
§8.3); its mirror image, a red dot with one leg coming \emph{in}, is
\(\sqrt2\,\bra{m}\) when its phase is \(m\pi\). So \(\bra m\) is
\(\tfrac1{\sqrt2}\) times a one-legged red effect. And \(H\) is the
yellow box. The numbers in front multiply to
\(\sqrt2\cdot\tfrac1{\sqrt2}\cdot\tfrac1{\sqrt2}=\tfrac1{\sqrt2}\):

\[\bra{m_1m_2}(H\otimes I)\,\mathrm{CNOT}\;=\;\tfrac1{\sqrt2}\;
\begin{tikzpicture}[sd, x=1.4em, y=1.4em,
  zsp/.style={circle,draw,line width=0.9pt,fill=green!55,minimum size=1.1em,inner sep=0pt},
  xsp/.style={circle,draw,line width=0.9pt,fill=red!75,minimum size=1.1em,inner sep=0pt},
  hbox/.style={draw,line width=0.9pt,fill=yellow!85,minimum size=0.9em,inner sep=1pt},
  ph/.style={font=\scriptsize,inner sep=1pt}]
  \draw[wire] (0,1.5) -- (4.4,1.5);
  \draw[wire] (0,0) -- (4.4,0);
  \draw[wire] (1.2,1.5) -- (1.2,0);
  \node[zsp] at (1.2,1.5) {};
  \node[xsp] at (1.2,0) {};
  \node[hbox] at (2.8,1.5) {};
  \node[xsp,label={[ph]above:{$m_1\pi$}}] at (4.4,1.5) {};
  \node[xsp,label={[ph]below:{$m_2\pi$}}] at (4.4,0) {};
\end{tikzpicture}
\;\overset{\text{colour change}}{=}\;\tfrac1{\sqrt2}\;
\begin{tikzpicture}[sd, x=1.4em, y=1.4em,
  zsp/.style={circle,draw,line width=0.9pt,fill=green!55,minimum size=1.1em,inner sep=0pt},
  xsp/.style={circle,draw,line width=0.9pt,fill=red!75,minimum size=1.1em,inner sep=0pt},
  ph/.style={font=\scriptsize,inner sep=1pt}]
  \draw[wire] (0,1.5) -- (3.0,1.5);
  \draw[wire] (0,0) -- (3.0,0);
  \draw[wire] (1.2,1.5) -- (1.2,0);
  \node[zsp] at (1.2,1.5) {};
  \node[xsp] at (1.2,0) {};
  \node[zsp,label={[ph]above:{$m_1\pi$}}] at (3.0,1.5) {};
  \node[xsp,label={[ph]below:{$m_2\pi$}}] at (3.0,0) {};
\end{tikzpicture}\]

The first step is colour change (Theorem 8.2): a red effect with a
yellow box on its one leg is a green effect of the same phase. Now each
wire carries two dots of one colour in a row, and fusion (Theorem 8.1)
merges them, adding the phases, \(0+m_1\pi\) on top and \(0+m_2\pi\)
below:

\[\;\overset{\text{fusion}}{=}\;\tfrac1{\sqrt2}\;
\begin{tikzpicture}[sd, x=1.4em, y=1.4em,
  zsp/.style={circle,draw,line width=0.9pt,fill=green!55,minimum size=1.1em,inner sep=0pt},
  xsp/.style={circle,draw,line width=0.9pt,fill=red!75,minimum size=1.1em,inner sep=0pt},
  ph/.style={font=\scriptsize,inner sep=1pt}]
  \draw[wire] (0,1.5) -- (1.4,1.5);
  \draw[wire] (0,0) -- (1.4,0);
  \draw[wire] (1.4,1.5) -- (1.4,0);
  \node[zsp,label={[ph]above:{$m_1\pi$}}] at (1.4,1.5) {};
  \node[xsp,label={[ph]below:{$m_2\pi$}}] at (1.4,0) {};
\end{tikzpicture}
\;\overset{\text{connectivity}}{=}\;\tfrac1{\sqrt2}\;
\begin{tikzpicture}[sd, x=1.4em, y=1.4em,
  zsp/.style={circle,draw,line width=0.9pt,fill=green!55,minimum size=1.1em,inner sep=0pt},
  xsp/.style={circle,draw,line width=0.9pt,fill=red!75,minimum size=1.1em,inner sep=0pt},
  ph/.style={font=\scriptsize,inner sep=1pt}]
  \draw[wire] (0,1.5) -- (4.6,1.5);
  \draw[wire] (0,0) -- (4.6,0);
  \sdcap{4.6}{1.5}{0}
  \node[zsp,label={[ph]above:{$m_1\pi$}}] at (1.6,1.5) {};
  \node[xsp,label={[ph]above:{$m_2\pi$}}] at (3.2,1.5) {};
\end{tikzpicture}\]

What remains is a green dot with one leg in and one out, joined by the
cord to a red dot with two legs in and none out. That red dot is
\(\bra{++}+e^{im_2\pi}\bra{--}\), which is \(\bra{00}+\bra{11}\) for
\(m_2=0\) and \(\bra{01}+\bra{10}\) for \(m_2=1\): a cap, carrying a red
ring \(m_2\pi\) on one leg. Since only the connectivity matters (§8.5),
the whole is a single wire running in at the top, through a green ring
and a red ring, and back out at the bottom: the cap of the reading.

\begin{quote}
\textbf{Result 10.2 (the reading is a ringed cap).} For each of the four
outcomes, exactly,
\[\bra{m_1m_2}(H\otimes I)\,\mathrm{CNOT} \;=\; \bra{B_{m_1m_2}} \;=\; \tfrac1{\sqrt2}\,\text{cap}\,\bigl(X^{m_2}Z^{m_1}\otimes I\bigr),\]
and the picture derivation above uses only colour change and fusion,
both exact, with the scalar \(\tfrac1{\sqrt2}\) carried from the two
translations. See \texttt{Code/Ch10/FourAnswers.py}, which checks the
circuit form against the ringed cap for all four outcomes, before and
after reducing it by the rules.
\end{quote}

\vskip 0pt plus 4\baselineskip\penalty-100\vskip 0pt plus
-4\baselineskip

\begin{quote}
\textbf{Coecke's Sketchbook.}\nopagebreak

\emph{Rings round a bend.} A single half-turn ring slides round a bend
and arrives unchanged: sliding transposes (Theorem 6.4), and
\(X^{T}=X\), \(Z^{T}=Z\). Two rings in a row do something that is easy
to miss. The transpose of a product reverses it,
\((BA)^{T}=A^{T}B^{T}\), so a chain slid round a bend arrives in the
\emph{opposite order}. That is Result 10.1's third form: red then green
on the top leg is green then red on the bottom leg. Abramsky and Coecke
noticed the same thing in the abstract, remarking on a ``seemingly
`acausal'\,'' inversion of the order in which things are applied along a
bent wire (their §3.3). The order matters here because a red half-turn
and a green one do not commute: \(XZ=-ZX\). Passing one through the
other costs a sign (Theorem 8.4 with \(\alpha=\pi\)). Every rewrite in
this chapter is planned so that no ring ever has to pass a ring of the
other colour.
\end{quote}

\begin{center}\rule{0.5\linewidth}{0.5pt}\end{center}

\hypertarget{two-bits-early}{%
\section{§10.2 Two Bits Early}\label{two-bits-early}}

Now the scene. Alice holds a qubit on wire \(1\) in a state
\(\ket\psi=\alpha\ket0+\beta\ket1\) that she does not know. She also
holds wire \(a\), one half of a pair \(\ket{\Phi^+}\); the other half,
wire \(c\), is with Bob, far away. The wires are drawn top to bottom
\(1,a,c\). Alice performs a Bell reading on her two wires, \(1\) and
\(a\). What happens to Bob's?

\textbf{On the blackboard.} Write the three wires out:

\[\ket\psi\otimes\ket{\Phi^+} \;=\; \tfrac1{\sqrt2}\bigl(\alpha\ket{000}+\alpha\ket{011}+\beta\ket{100}+\beta\ket{111}\bigr).\]

To see what the reading does we must write Alice's two wires in the Bell
basis. Reading the table of §10.1 backwards,

\[\begin{aligned}
\ket{00}&=\tfrac1{\sqrt2}\bigl(\ket{\Phi^+}+\ket{\Phi^-}\bigr), &\qquad
\ket{11}&=\tfrac1{\sqrt2}\bigl(\ket{\Phi^+}-\ket{\Phi^-}\bigr),\\
\ket{01}&=\tfrac1{\sqrt2}\bigl(\ket{\Psi^+}+\ket{\Psi^-}\bigr), &\qquad
\ket{10}&=\tfrac1{\sqrt2}\bigl(\ket{\Psi^+}-\ket{\Psi^-}\bigr).
\end{aligned}\]

Substitute into the four terms, which become eight, and gather them by
Bell state (the sidebar below writes every term):

\[\begin{aligned}
\ket\psi\otimes\ket{\Phi^+}\;=\;\tfrac12\Bigl[\;&\ket{\Phi^+}\bigl(\alpha\ket0+\beta\ket1\bigr)+\ket{\Psi^+}\bigl(\alpha\ket1+\beta\ket0\bigr)\\
+\;&\ket{\Phi^-}\bigl(\alpha\ket0-\beta\ket1\bigr)+\ket{\Psi^-}\bigl(\alpha\ket1-\beta\ket0\bigr)\Bigr].
\end{aligned}\]

Look at what stands next to each Bell state. Next to \(\ket{\Phi^+}\) is
\(\ket\psi\) itself. Next to \(\ket{\Psi^+}\) is \(X\ket\psi\), flipped;
next to \(\ket{\Phi^-}\) is \(Z\ket\psi\), phase-flipped; next to
\(\ket{\Psi^-}\) is \(XZ\ket\psi=X(\alpha\ket0-\beta\ket1)\), both. The
message is on Bob's wire in every branch, turned by the Paulis that the
branch's own Bell state carries.

\begin{quote}
\textbf{Result 10.3 (the message arrives turned).} For every one-qubit
state \(\ket\psi\), exactly,
\[\ket\psi_1\otimes\ket{\Phi^+}_{ac}\;=\;\tfrac12\sum_{m_1,m_2}\ket{B_{m_1m_2}}_{1a}\otimes X^{m_2}Z^{m_1}\ket\psi_c .\]
So Alice's reading gives each of the four outcomes with probability
\(\tfrac14\), whatever \(\ket\psi\) is, and outcome \(m_1m_2\) leaves
Bob's wire in \(X^{m_2}Z^{m_1}\ket\psi\).
\end{quote}

\emph{Proof idea.} A Bell state is a bent wire carrying the outcome's
rings. Asking whether Alice's wires stand in that Bell state bends the
message wire into Bob's, rings and all.

\emph{Proof (blackboard).} The calculation above, with \(\alpha\) and
\(\beta\) left free, is the whole proof: it holds for every \(\ket\psi\)
at once. Each branch has amplitude \(\tfrac12\) times a unit vector,
since Paulis preserve length, so its probability is \(\tfrac14\).
\(\;\blacksquare\) See \texttt{Code/Ch10/TwoBitsEarly.py}.

\textbf{In the sketchbook.} On the branch where the reading says
\(m_1m_2\), draw the scene. The input comes in on the top wire. Alice's
answer is the ringed cap of Result 10.2, closing the top wire onto the
middle one. Bob's pair is a cup on the middle and bottom wires. Two
factors of \(\tfrac1{\sqrt2}\), one from the reading and one from the
pair, make the \(\tfrac12\) in front. Bob's wire runs out at the bottom.
The wire makes an S, and Theorem 6.2 pulls it straight:

\[\tfrac12\;
\begin{tikzpicture}[sd, x=1.4em, y=1.4em,
  zsp/.style={circle,draw,line width=0.9pt,fill=green!55,minimum size=1.1em,inner sep=0pt},
  xsp/.style={circle,draw,line width=0.9pt,fill=red!75,minimum size=1.1em,inner sep=0pt},
  ph/.style={font=\scriptsize,inner sep=1pt}]
  \draw[wire] (0,3) -- (4.8,3);
  \sdcap{4.8}{3}{1.5}
  \draw[wire] (1.2,1.5) -- (4.8,1.5);
  \sdcup{1.2}{1.5}{0}
  \draw[wire] (1.2,0) -- (6.4,0);
  \node[zsp,label={[ph]above:{$m_1\pi$}}] at (1.6,3) {};
  \node[xsp,label={[ph]above:{$m_2\pi$}}] at (3.2,3) {};
  \node[sdlabel,anchor=east] at (-0.2,3) {$1$};
  \node[sdlabel,anchor=south] at (3.0,1.55) {$a$};
  \node[sdlabel,anchor=west] at (6.6,0) {$c$};
\end{tikzpicture}
\;\;\overset{\text{yank}}{=}\;\;
\tfrac12\;
\begin{tikzpicture}[sd, x=1.4em, y=1.4em,
  zsp/.style={circle,draw,line width=0.9pt,fill=green!55,minimum size=1.1em,inner sep=0pt},
  xsp/.style={circle,draw,line width=0.9pt,fill=red!75,minimum size=1.1em,inner sep=0pt},
  ph/.style={font=\scriptsize,inner sep=1pt}]
  \draw[wire] (0,0) -- (4.6,0);
  \node[zsp,label={[ph]above:{$m_1\pi$}}] at (1.4,0) {};
  \node[xsp,label={[ph]above:{$m_2\pi$}}] at (3.0,0) {};
\end{tikzpicture}\]

That is Result 10.3 for all four branches at once: the message comes out
on Bob's wire, times \(\tfrac12\), still wearing Alice's rings,
\(X^{m_2}Z^{m_1}\ket\psi\). It is there \emph{already}, at the instant
of the reading, and no letter has been sent.

It is also no use to Bob. He does not know which branch he is on, and
the four branches are equally likely. What he holds, honestly described,
is the average.

\begin{quote}
\textbf{Result 10.4 (nothing before the letter).} Before the outcome
reaches him, Bob's wire is in the maximally mixed state \(\tfrac12 I\),
for every input, pure or mixed. Nothing Alice sends changes what Bob
holds until the letter arrives.
\end{quote}

\emph{Proof (blackboard).} The input may as well be a density matrix
\(\rho=\bigl(\begin{smallmatrix}r&s\\t&u\end{smallmatrix}\bigr)\), with
\(r+u=\operatorname{tr}\rho=1\); the pure case is
\(\rho=\ket\psi\bra\psi\). By Result 10.3 and linearity, branch
\(m_1m_2\) leaves Bob with \(\tfrac14\,P\rho P^{\dagger}\) where
\(P=X^{m_2}Z^{m_1}\), so Bob's average is
\(\tfrac14\bigl(\rho+Z\rho Z+X\rho X+XZ\rho ZX\bigr)\). Now \(Z\rho Z\)
flips the signs of \(s\) and \(t\), so
\(\rho+Z\rho Z=\mathrm{diag}(2r,2u)\); and \(X(\cdot)X\) exchanges the
two diagonal entries, so the sum of all four is
\(\mathrm{diag}(2r+2u,\,2r+2u)=2I\). A quarter of that is
\(\tfrac12 I\). \(\;\blacksquare\)

\emph{Proof (sketchbook).} Not knowing the outcome is forgetting it. The
reading is the circuit \((H\otimes I)\,\mathrm{CNOT}\), which is
unitary, followed by two of Chapter 9's looks. Draw the reading on
doubled wires, as every real process must be drawn since Chapter 9: a
box \(M\) that takes Alice's two wires and their two mirror copies, four
lines in, and emits two classical wires, one per bit. To forget the
outcome is to end each classical wire in the ground symbol. A look
followed by forgetting its outcome is a discard (Result 9.5), and a
discard after a unitary is a discard (Result 9.4, causality). So Alice's
reading with its record forgotten is simply the discard of both her
wires, a cap on each wire and its mirror:

\[\begin{tikzpicture}[sd]
  \draw[wire] (0,3.6) -- (1.6,3.6);
  \draw[wire] (0,2.4) -- (1.6,2.4);
  \draw[wire] (0,1.2) -- (1.6,1.2);
  \draw[wire] (0,0) -- (1.6,0);
  \draw[cwire] (2,2.4) -- (5.0,2.4);
  \draw[cwire] (2,1.2) -- (5.0,1.2);
  \sdground{5.0}{2.4}
  \sdground{5.0}{1.2}
  \node[sdbox, minimum height=4.9em, minimum width=2em] at (2.2,1.8) {$M$};
  \node[sdlabel,anchor=east] at (-0.2,3.6) {$1$};
  \node[sdlabel,anchor=east] at (-0.2,2.4) {$\bar 1$};
  \node[sdlabel,anchor=east] at (-0.2,1.2) {$a$};
  \node[sdlabel,anchor=east] at (-0.2,0) {$\bar a$};
\end{tikzpicture}
\;\;=\;\;
\begin{tikzpicture}[sd]
  \draw[wire] (0,3.6) -- (1.8,3.6);
  \draw[wire] (0,2.4) -- (1.8,2.4);
  \draw[wire] (0,1.2) -- (1.8,1.2);
  \draw[wire] (0,0) -- (1.8,0);
  \sdcap{1.8}{3.6}{2.4}
  \sdcap{1.8}{1.2}{0}
\end{tikzpicture}\]

Discarding the input gives its trace, \(1\). What is left is a Bell pair
with Alice's half discarded, and Result 9.3 says Bob's half is then
\(\tfrac12 I\). \(\;\blacksquare\) See
\texttt{Code/Ch10/TwoBitsEarly.py}.

So the message has arrived and cannot be read. Bob's wire, before the
letter, is the same featureless grey it was before Alice did anything;
Chapter 9's no-signalling theorem is at work. Teleportation sends
nothing faster than light. The message is there, hidden in a four-way
average, and only two bits of news can pull it out.

\begin{quote}
\textbf{Dirac's Blackboard.}

\emph{The eight terms, and what the outcomes say.} Substituting the Bell
basis into
\(\sqrt2\,\ket\psi\ket{\Phi^+}=\alpha\ket{00}\ket0+\alpha\ket{01}\ket1+\beta\ket{10}\ket0+\beta\ket{11}\ket1\)
gives, times \(\tfrac1{\sqrt2}\),
\[\begin{aligned}&\alpha\ket{\Phi^+}\ket0+\alpha\ket{\Phi^-}\ket0+\alpha\ket{\Psi^+}\ket1+\alpha\ket{\Psi^-}\ket1\\ {}+{}&\beta\ket{\Psi^+}\ket0-\beta\ket{\Psi^-}\ket0+\beta\ket{\Phi^+}\ket1-\beta\ket{\Phi^-}\ket1 .\end{aligned}\]
Gathered by Bell state, these are the four brackets of §10.2, with
\(\tfrac1{\sqrt2}\cdot\tfrac1{\sqrt2}=\tfrac12\) in front. The four
probabilities are
\(\bigl\|\tfrac12 X^{m_2}Z^{m_1}\psi\bigr\|^2=\tfrac14\bigl(|\alpha|^2+|\beta|^2\bigr)=\tfrac14\),
and they sum to \(1\). Notice what does \emph{not} appear: \(\alpha\)
and \(\beta\). The outcome is a fair four-way toss whatever the message
was, so the two bits Alice will write in her letter say nothing whatever
about \(\ket\psi\). The letter carries information \emph{about the
turn}, and none about the thing turned.
\end{quote}

\begin{center}\rule{0.5\linewidth}{0.5pt}\end{center}

\hypertarget{the-letter-and-the-half-turns}{%
\section{§10.3 The Letter and the
Half-Turns}\label{the-letter-and-the-half-turns}}

Alice writes the two bits \(m_1m_2\) in a letter and posts it. When it
arrives, Bob knows he holds \(X^{m_2}Z^{m_1}\ket\psi\). To undo a
product, undo its factors in the reverse order: the inverse of
\(X^{m_2}Z^{m_1}\) is \(Z^{m_1}X^{m_2}\), since each factor is its own
inverse. So Bob applies \(X^{m_2}\) \textbf{first}, and then
\(Z^{m_1}\).

\textbf{On the blackboard.} Take the hardest case, outcome \(11\). Bob
holds \(\tfrac12\,XZ\ket\psi=\tfrac12(\alpha\ket1-\beta\ket0)\). The
flip first: \(\tfrac12(\alpha\ket0-\beta\ket1)\). Then the phase-flip:
\(\tfrac12(\alpha\ket0+\beta\ket1)=\tfrac12\ket\psi\). The message,
exactly, times the branch's \(\tfrac12\). Now do it in the other order.
The phase-flip first:
\(Z(\alpha\ket1-\beta\ket0)=-\alpha\ket1-\beta\ket0\); then the flip:
\(-\alpha\ket0-\beta\ket1\). That is \(-\tfrac12\ket\psi\): the message
with a minus sign in front. The other three outcomes need at most one
Pauli, so order cannot matter for them.

\textbf{In the sketchbook.} Bob's correction is a box whose content is
chosen by what the letter says. This needs one more new mark. In Chapter
9 classical wires only ever \emph{left} something: the look emitted one.
Now a classical wire runs \textbf{into a box}. A box with classical
wires running into it is a box chosen by their values: on the branch
where the wires carry \(m_1m_2\), the box is the process named for
\(m_1m_2\). We let the classical wires run on through the box and out
the other side, the record kept, since reading a definite value does not
use it up. Here is Bob's correction box \(C\), fed by the two bits of
the letter, a classical wire \(x\) carrying the flip bit \(m_2\) and a
classical wire \(z\) carrying the phase bit \(m_1\). Bob's wire \(c\) is
drawn once, its mirror copy understood, as Chapter 9 drew the ground
symbol on a single line:

\[\begin{tikzpicture}[sd]
  \draw[cwire] (0,2.4) -- (6,2.4);
  \draw[cwire] (0,1.2) -- (6,1.2);
  \draw[wire] (0,0) -- (6,0);
  \node[sdbox, minimum height=4.1em, minimum width=2em] at (3,1.2) {$C$};
  \node[sdlabel,anchor=east] at (-0.2,2.4) {$x$};
  \node[sdlabel,anchor=east] at (-0.2,1.2) {$z$};
  \node[sdlabel,anchor=east] at (-0.2,0) {$c$};
\end{tikzpicture}
\qquad\xrightarrow{\;\text{on the branch } m_1m_2\;}\qquad
\begin{tikzpicture}[sd, x=1.4em, y=1.4em,
  zsp/.style={circle,draw,line width=0.9pt,fill=green!55,minimum size=1.1em,inner sep=0pt},
  xsp/.style={circle,draw,line width=0.9pt,fill=red!75,minimum size=1.1em,inner sep=0pt},
  ph/.style={font=\scriptsize,inner sep=1pt}]
  \draw[wire] (0,0) -- (4.6,0);
  \node[xsp,label={[ph]above:{$m_2\pi$}}] at (1.4,0) {};
  \node[zsp,label={[ph]above:{$m_1\pi$}}] at (3.0,0) {};
\end{tikzpicture}\]

On the blackboard the box is a sum over the branches. Drawn on doubled
wires, as a real process must be (Chapter 9), with the classical record
passed through,

\[C\;=\;\sum_{m_1,m_2}\ket{m_2m_1}\bra{m_2m_1}_{xz}\otimes\bigl(Z^{m_1}X^{m_2}\bigr)_c\otimes\bigl(\overline{Z^{m_1}X^{m_2}}\bigr)_{\bar c},\]

where the bar is the complex conjugate, which for the real Paulis
changes nothing. The two new marks meet here. The bit on the classical
wire is the bit in the ring's label: \emph{classical control} is the
bit-labelled ring with its bit supplied by a wire.

\textbf{\emph{Why this definition.}} Why not simply draw four pictures,
one per outcome? Because Bob does not choose his box; the letter does. A
classical wire into a box says so on the page. It records that the
correction depends on a value that was settled somewhere else, and has
to travel to get here.

\begin{quote}
\textbf{Result 10.5 (teleportation, all four outcomes, exact).} Let
wires \(1,a,c\) run top to bottom, let
\(\ket{\Phi^+}=\tfrac1{\sqrt2}\,\)cup on \((a,c)\), and let Alice's Bell
reading on \((1,a)\) have effects
\(\bra{B_{m_1m_2}}=\bra{\Phi^+}(X^{m_2}Z^{m_1}\otimes I)\). Then for
\textbf{every} one-qubit \(\ket\psi\) and \textbf{each} outcome
\(m_1m_2\),
\[\bigl(Z^{m_1}X^{m_2}\bigr)_c\,\bigl(\bra{B_{m_1m_2}}_{1a}\otimes I_c\bigr)\bigl(\ket\psi_1\otimes\ket{\Phi^+}_{ac}\bigr)\;=\;\tfrac12\,\ket\psi_c ,\]
exactly, the scalar included. Each outcome has probability
\(\bigl\|\tfrac12\psi\bigr\|^2=\tfrac14\), and after the correction
(\(X^{m_2}\) first, then \(Z^{m_1}\)) Bob holds exactly \(\ket\psi\). As
a channel: read, send both bits, correct and forget the bits, and the
result is the identity, \(\rho\mapsto\rho\), each branch contributing
exactly \(\tfrac14\rho\).
\end{quote}

The statement matches Abramsky and Coecke's account (their §2.1): four
outcomes, four unitary corrections on Bob's qubit, and at the end Bob's
qubit in the state the input began in. §10.6 proves it twice. See
\texttt{Code/Ch10/Teleportation.py}, which evaluates both sides on a
non-real \(\ket\psi\) and a non-real \(\rho\), and also checks that the
wrong order gives \(-\tfrac12\psi\) on outcome \(11\).

\begin{quote}
\textbf{Dirac's Blackboard.}

\emph{Why the order matters exactly and not physically.} With the
corrections in the wrong order, outcome \(11\) ends in
\(-\tfrac12\ket\psi\) rather than \(\tfrac12\ket\psi\). The minus sign
is a global phase, and no experiment can see it: in the doubled picture
of Chapter 9 an overall phase \(e^{i\theta}\) on the ket copy meets
\(e^{-i\theta}\) on the mirror copy, and \((-1)\cdot(-1)=1\). So
physically the two orders do the same thing. But Result 10.5 is an
equation between vectors, and it holds with \(+\tfrac12\) only when the
flip comes first. Abramsky and Coecke's list of corrections shows the
same sign: of their four, three are their own inverses, and the fourth
is minus its own inverse (their §2.1). We fix the order once,
\(Z^{m_1}X^{m_2}\), and keep it for the rest of the chapter, so that
every equation is exact as written.
\end{quote}

\begin{quote}
\textbf{Feynman's Notebook.}

\emph{Two channels.} The protocol was published in 1993 by C. H.
Bennett, G. Brassard, C. Crépeau, R. Jozsa, A. Peres and W. K. Wootters,
in \emph{Physical Review Letters} 70, 1895--1899, under the title
``Teleporting an unknown quantum state via dual classical and
Einstein-Podolsky-Rosen channels'' (Abramsky and Coecke's reference
{[}5{]}). The title names both ingredients: a classical channel, which
is the letter, and an Einstein--Podolsky--Rosen channel, which is the
shared pair. Abramsky and Coecke, opening their 2004 paper with
teleportation, dwell on what the count means. The unknown qubit is an
arbitrary pair of complex numbers, and yet only two classical bits
cross, which they call ``no mean feat''. Since a continuous quantity has
been moved and only two bits have travelled, they argue that there must
be a quantum flow of information besides the classical one. In the
sketchbook that flow is the wire, bent at the source of the pair and
pulled straight by the reading.
\end{quote}

\begin{center}\rule{0.5\linewidth}{0.5pt}\end{center}

\hypertarget{the-other-way-round-superdense-coding}{%
\section{§10.4 The Other Way Round: Superdense
Coding}\label{the-other-way-round-superdense-coding}}

Two bits in a letter moved one qubit. Can it go the other way: can one
qubit sent carry two bits?

\textbf{On the blackboard.} Again Alice and Bob share \(\ket{\Phi^+}\),
Alice's half on wire \(a\) and Bob's on wire \(c\). This time Alice has
two bits \(k_1k_2\) to send. She turns her half by \(Z^{k_1}X^{k_2}\)
(the flip first, then the phase-flip, as in Bob's correction table) and
posts the \emph{qubit} to Bob. Bob now holds both halves, and he makes a
Bell reading on them.

Take \(k_1k_2=10\). Alice applies \(Z\), and
\((Z\otimes I)\ket{\Phi^+}=\ket{\Phi^-}=\ket{B_{10}}\). Bob's reading is
certain to answer \(10\). In general there is nothing to calculate: by
Result 10.1, \((Z^{k_1}X^{k_2}\otimes I)\ket{\Phi^+}\) \emph{is}
\(\ket{B_{k_1k_2}}\), and the four Bell states are orthonormal, so the
reading tells them apart without error. The amplitude that Bob reads
\(j=j_1j_2\) when Alice wrote \(k=k_1k_2\) is

\[\bra{B_j}\bigl(Z^{k_1}X^{k_2}\otimes I\bigr)\ket{\Phi^+}\;=\;\braket{B_j|B_k}\;=\;\delta_{jk}.\]

There is a second way to see the same number, and it is the one the
picture uses. Write the reading's effect in its ring form: the amplitude
is
\(\bra{\Phi^+}\bigl(X^{j_2}Z^{j_1}Z^{k_1}X^{k_2}\otimes I\bigr)\ket{\Phi^+}\),
a Bell state and a Bell effect with one box between them on the top
wire. That shape has a name.

\textbf{In the sketchbook.} Close a wire on itself: a cup, and then a
cap on the same two wires. Nothing comes in and nothing goes out, so it
is a number. We call it a \textbf{closed loop}. Put a box \(A\) on it,
and the loop is the \textbf{trace} of \(A\):

\[\begin{tikzpicture}[sd]
  \sdcup{0}{1.5}{0}
  \draw[wire] (0,1.5) -- (3.4,1.5);
  \draw[wire] (0,0) -- (3.4,0);
  \sdcap{3.4}{1.5}{0}
  \node[sdbox] at (1.7,1.5) {$A$};
\end{tikzpicture}
\;=\;\operatorname{tr}A,
\qquad\qquad
\begin{tikzpicture}[sd]
  \sdcup{0}{1.5}{0}
  \draw[wire] (0,1.5) -- (2,1.5);
  \draw[wire] (0,0) -- (2,0);
  \sdcap{2}{1.5}{0}
\end{tikzpicture}
\;=\;2 .\]

\begin{quote}
\textbf{Result 10.6 (a loop is a trace).} For every \(2\times2\) matrix
\(A\), \(\;\text{cap}\,(A\otimes I)\,\text{cup}=\operatorname{tr}A\),
exactly. In particular the bare loop is \(\operatorname{tr}I=2\).
\end{quote}

\emph{Proof.}
\(\text{cap}\,(A\otimes I)\,\text{cup}=\sum_{i,k}\bra{ii}\bigl(A\ket k\otimes\ket k\bigr)=\sum_{i,k}\bra iA\ket k\braket{i|k}=\sum_i A_{ii}\).
The bare loop was met in §9.3, where it was a system's mirror-cup
discarded. \(\;\blacksquare\)

\textbf{\emph{Why this definition.}} Why give a name to a picture that
is only a number? Because the loop is the shape of every question of the
form ``does the reading return what was prepared?''. A preparation is a
cup, a reading is a cap, and whatever was done in between sits on the
loop. Naming it the trace turns sixteen amplitudes into one fact about
the traces of four matrices.

The loop is where the unnormalized cup pays us back. Because cup and cap
carry no \(\tfrac1{\sqrt2}\), a loop is a plain trace, and the Bell
state and Bell effect contribute
\(\tfrac1{\sqrt2}\cdot\tfrac1{\sqrt2}=\tfrac12\), written in front. Now
draw dense coding. The cup is the shared pair. On its top leg, reading
outward from the bend, come Alice's rings, red \(k_2\pi\) then green
\(k_1\pi\), and then the reading's rings, green \(j_1\pi\) then red
\(j_2\pi\), and the cap closes the loop:

\[\bra{B_j}\bigl(Z^{k_1}X^{k_2}\otimes I\bigr)\ket{\Phi^+}\;=\;\tfrac12\;
\begin{tikzpicture}[sd, x=1.4em, y=1.4em,
  zsp/.style={circle,draw,line width=0.9pt,fill=green!55,minimum size=1.1em,inner sep=0pt},
  xsp/.style={circle,draw,line width=0.9pt,fill=red!75,minimum size=1.1em,inner sep=0pt},
  ph/.style={font=\scriptsize,inner sep=1pt}]
  \sdcup{0}{1.5}{0}
  \draw[wire] (0,1.5) -- (7.6,1.5);
  \draw[wire] (0,0) -- (7.6,0);
  \sdcap{7.6}{1.5}{0}
  \node[xsp,label={[ph]above:{$k_2\pi$}}] at (1.4,1.5) {};
  \node[zsp,label={[ph]above:{$k_1\pi$}}] at (3.0,1.5) {};
  \node[zsp,label={[ph]above:{$j_1\pi$}}] at (4.6,1.5) {};
  \node[xsp,label={[ph]above:{$j_2\pi$}}] at (6.2,1.5) {};
\end{tikzpicture}\]

In the middle, green meets green, and fusion adds them to
\((j_1+k_1)\pi\). The last red ring is a single ring, so it slides round
the cap, along the bottom and round the cup, arriving unchanged (Theorem
6.4, twice), until it sits next to the first red ring; fusion adds those
to \((j_2+k_2)\pi\). What is left is a loop carrying one green ring
\((j_1+k_1)\pi\) and one red ring \((j_2+k_2)\pi\). If \(j=k\), both
phases are whole turns, both rings are bare wire (§8.1), and the value
is \(\tfrac12\) times the bare loop, \(\tfrac12\cdot2=1\). If
\(j\neq k\), a ring is left over: a red one, a green one, or one of
each. By Result 10.6 the loop is then the trace of \(X\), of \(Z\) or of
\(ZX\), and all three traces are \(0\). So the read label matches the
sent one with certainty.

\begin{quote}
\textbf{Result 10.7 (superdense coding).} For all sixteen pairs of
two-bit labels \(j,k\), exactly,
\[\bra{B_j}\bigl(Z^{k_1}X^{k_2}\otimes I\bigr)\ket{\Phi^+}\;=\;\tfrac12\operatorname{tr}\bigl(X^{j_2}Z^{j_1}Z^{k_1}X^{k_2}\bigr)\;=\;\delta_{jk}.\]
Alice's codebook, the half-turn \(Z^{k_1}X^{k_2}\) for message
\(k_1k_2\), is Bob's correction table in teleportation, entry for entry.
See \texttt{Code/Ch10/DenseCoding.py}.
\end{quote}

Here is what ``the mirror image'' means, and it needs saying precisely,
because in this book the mirror is the dagger. The same three pieces
appear in both protocols: one shared pair, the Bell reading, and the
four half-turns \(Z^{m_1}X^{m_2}\). In \textbf{teleportation} Alice
\emph{reads} and Bob \emph{turns}: one pair and two bits sent move one
qubit. In \textbf{superdense coding} Alice \emph{turns} and Bob
\emph{reads}: one pair and one qubit sent move two bits. The classical
and quantum channels have changed places. And the pieces themselves go
into each other in the mirror. Alice's four readings in teleportation
are the Bell effects \(\bra{B_k}\); reflected, they are exactly her four
preparations in dense coding, the Bell states
\(\ket{B_k}=\bra{B_k}^{\dagger}\), each a cup carrying the same rings,
because a half-turn ring is its own mirror image. The dagger has
exchanged the two ends of the pieces. It has \emph{not} reversed the
protocol. The dagger of teleportation's corrected branch,
\(\tfrac12 I\), is \(\tfrac12 I\) again: a qubit going in and a qubit
coming out, teleportation run backwards, which is not dense coding.

And the honest accounting. Two bits per qubit sent sounds like something
for nothing, and it is not. The pair had to be shared beforehand; one of
its halves travelled to Bob at some earlier time. Counting that half,
two qubits carried two bits.

In the dialogue, as this sinks in, Achilles looks up at the print on the
parlour wall, \emph{Magic Mirror}: a mirror standing upright, and
creatures stepping out of it.

\begin{quote}
\textbf{Coecke's Sketchbook.}

\emph{A mirror is not a half-turn.} The reading is the dagger of the
preparation, not its transpose. With the real Paulis the two cannot be
told apart, so here is a test that can. Pick a unitary \(W\) that is
neither real nor symmetric, and use the twisted boxes
\(U_k=Z^{k_1}X^{k_2}W\) in place of the Paulis, with states
\(\ket{E_k}=(U_k\otimes I)\ket{\Phi^+}\). The reading's effects are
\(\bra{E_k}=\ket{E_k}^{\dagger}=\bra{\Phi^+}(U_k^{\dagger}\otimes I)\),
which is not \(\ket{E_k}^{T}\). The yank of §10.2 then delivers
\(\tfrac12U_k^{\dagger}\ket\psi\), for every unitary \(U_k\), so the
correction is \(U_k\) itself: not \(U_k^{T}\), not \(\bar U_k\), not
\(U_k^{\dagger}\). Put the box on the \emph{other} leg instead,
\(\ket{E'_k}=(I\otimes U_k)\ket{\Phi^+}\), and sliding moves it to the
top leg as \(U_k^{T}\), so the correction becomes \(U_k^{T}\). The
transpose enters by sliding, and only by sliding, and Bob's output is
always \(\ket\psi\), never its conjugate. The companion code runs this
with a random \(W\).
\end{quote}

\begin{quote}
\textbf{Feynman's Notebook.}\nopagebreak

\emph{The mirror came first.} The other way round was published a year
before teleportation: Bennett and Wiesner, ``Communication via one- and
two-particle operators on Einstein-Podolsky-Rosen states'',
\emph{Physical Review Letters} 69, 2881--2884 (1992). Read with this
chapter's eyes, the title names the two halves of the scheme: the sender
acts on one particle, her half of the pair, and the receiver on two
together, with a joint reading. That gloss is ours; the paper itself
should be consulted for its own framing. The name \emph{superdense
coding} is the one the protocol goes by now. The two papers share an
author and a resource. What one of them spends the shared pair on,
moving a qubit with two bits, the other spends it on in reverse, moving
two bits with a qubit.
\end{quote}

\begin{center}\rule{0.5\linewidth}{0.5pt}\end{center}

\hypertarget{the-rosetta-table-8}{%
\section{§10.5 The Rosetta Table}\label{the-rosetta-table-8}}

The rows this chapter adds to the running dictionary, numbered on from
Chapter 9's seventy-two. Three introduce new picture elements: the
bit-labelled phase (75), the classical wire into a box (76) and the
closed loop (77). The rest are protocols built from them.

\begingroup\renewcommand{\arraystretch}{1.5}

\begin{longtable}[]{@{}
  >{\raggedright\arraybackslash}p{(\columnwidth - 4\tabcolsep) * \real{0.2727}}
  >{\raggedright\arraybackslash}p{(\columnwidth - 4\tabcolsep) * \real{0.3636}}
  >{\raggedright\arraybackslash}p{(\columnwidth - 4\tabcolsep) * \real{0.3636}}@{}}
\toprule\noalign{}
\begin{minipage}[b]{\linewidth}\raggedright
\#
\end{minipage} & \begin{minipage}[b]{\linewidth}\raggedright
Algebra (Dirac)
\end{minipage} & \begin{minipage}[b]{\linewidth}\raggedright
Picture (Coecke)
\end{minipage} \\
\midrule\noalign{}
\endhead
\bottomrule\noalign{}
\endlastfoot
73 & Bell states
\(\ket{B_{m_1m_2}}=(Z^{m_1}X^{m_2}\otimes I)\ket{\Phi^+}\) & a cup,
\(\times\tfrac1{\sqrt2}\), with a red ring (\(m_2\pi\)) then a green
ring (\(m_1\pi\)) on its top leg \\
74 & the Bell reading: effects
\(\bra{B_{m_1m_2}}=\bra{\Phi^+}(X^{m_2}Z^{m_1}\otimes I)\), outcome two
bits & the mirrored cap with the rings on its incoming leg; as a box
\(M\) on doubled wires, two classical wires out \\
75 & a Pauli chosen by a bit: \(Z^m\), \(X^m\), \(m\in\{0,1\}\) & a
spider whose phase is a bit times \(\pi\): one picture for all the cases
--- \emph{new element} \\
76 & classically controlled correction: apply \(X^{m_2}\), then
\(Z^{m_1}\) & a classical wire running into a box (classical control)
--- \emph{new element} \\
77 & \(\operatorname{tr}A=\text{cap}\,(A\otimes I)\,\text{cup}\);
\(\operatorname{tr}I=2\) & a closed loop through a box \(A\) is its
trace; the bare loop is the number \(2\) (§9.3) --- \emph{new
element} \\
78 & teleportation:
\((Z^{m_1}X^{m_2})(\bra{B_{m_1m_2}}\otimes I)(\psi\otimes\Phi^+)=\tfrac12\psi\),
\(p=\tfrac14\) & the S-bend yanks; red meets red, green meets green; a
bare wire \(\times\tfrac12\) \\
79 & superdense coding:
\(\bra{B_j}(Z^{k_1}X^{k_2}\otimes I)\ket{\Phi^+}=\delta_{jk}\) & the
loop with Alice's rings and the reading's rings: clean (value \(1\)) iff
\(j=k\), otherwise a red ring, a green ring or one of each is left,
value \(0\) \\
80 & swapping corrected:
\(\tfrac14\sum_k\ket k\!\bra k\otimes\ket{\Phi^+}\!\bra{\Phi^+}_{dc}\) &
The Mended Junction: two doubled cups, the reading box, two classical
wires, the correction box; one doubled cup \((d,c)\) on every outcome \\
\end{longtable}

\endgroup

Every row is an exact equality, scalar included; none of them needs the
``up to a non-zero scalar'' of Chapter 8. Rows 78 and 80 are proved in
§10.6 and §10.7.

\begin{center}\rule{0.5\linewidth}{0.5pt}\end{center}

\hypertarget{the-theorem-twice-9}{%
\section{§10.6 The Theorem, Twice}\label{the-theorem-twice-9}}

\begin{quote}
\textbf{Result 10.5 (teleportation, all four outcomes, exact),
restated.} For every one-qubit \(\ket\psi\) and each outcome \(m_1m_2\)
of Alice's Bell reading on wires \(1,a\),
\[\bigl(Z^{m_1}X^{m_2}\bigr)_c\,\bigl(\bra{B_{m_1m_2}}_{1a}\otimes I_c\bigr)\bigl(\ket\psi_1\otimes\ket{\Phi^+}_{ac}\bigr)\;=\;\tfrac12\,\ket\psi_c ,\]
exactly; each outcome has probability \(\tfrac14\); as a channel, the
protocol is the identity.
\end{quote}

\hypertarget{diracs-blackboard-9}{%
\subsection{Dirac's Blackboard}\label{diracs-blackboard-9}}

Let \(\ket\psi=\alpha\ket0+\beta\ket1\) with \(|\alpha|^2+|\beta|^2=1\).

\textbf{Step 1: expand.}
\(\ket\psi\otimes\ket{\Phi^+}=\tfrac1{\sqrt2}\bigl(\alpha\ket{00}\ket0+\alpha\ket{01}\ket1+\beta\ket{10}\ket0+\beta\ket{11}\ket1\bigr)\),
Alice's two wires written first.

\textbf{Step 2: change basis on Alice's wires.} Substitute
\(\ket{00},\ket{01},\ket{10},\ket{11}\) in terms of the Bell basis
(§10.2), which gives eight terms (the Dirac's Blackboard of §10.2), and
regroup:

\[\begin{aligned}
\ket\psi\otimes\ket{\Phi^+}=\tfrac12\Bigl[\;&\ket{B_{00}}(\alpha\ket0+\beta\ket1)+\ket{B_{01}}(\alpha\ket1+\beta\ket0)\\
+\;&\ket{B_{10}}(\alpha\ket0-\beta\ket1)+\ket{B_{11}}(\alpha\ket1-\beta\ket0)\Bigr].
\end{aligned}\]

\textbf{Step 3: read.} The effect \(\bra{B_{m_1m_2}}\) on Alice's wires
keeps exactly one bracket, because the Bell states are orthonormal
(Result 10.1). Branch \(m_1m_2\) leaves Bob's wire with \(\tfrac12\)
times its bracket.

\textbf{Step 4: correct, all four branches.}

\[\begin{aligned}
00:&\quad \tfrac12(\alpha\ket0+\beta\ket1) &&\xrightarrow{\;\text{nothing}\;}\;\tfrac12(\alpha\ket0+\beta\ket1),\\
01:&\quad \tfrac12(\alpha\ket1+\beta\ket0) &&\xrightarrow{\;X\;}\;\tfrac12(\alpha\ket0+\beta\ket1),\\
10:&\quad \tfrac12(\alpha\ket0-\beta\ket1) &&\xrightarrow{\;Z\;}\;\tfrac12(\alpha\ket0+\beta\ket1),\\
11:&\quad \tfrac12(\alpha\ket1-\beta\ket0) &&\xrightarrow{\;X\;}\;\tfrac12(\alpha\ket0-\beta\ket1)\;\xrightarrow{\;Z\;}\;\tfrac12(\alpha\ket0+\beta\ket1).
\end{aligned}\]

Every branch ends in \(\tfrac12\ket\psi\), exactly.

\textbf{Step 5: probabilities and the channel.} Each branch has squared
length \(\tfrac14(|\alpha|^2+|\beta|^2)=\tfrac14\), and four quarters
make \(1\). For a mixed input
\(\rho=\sum_i p_i\ket{\psi_i}\bra{\psi_i}\) apply Steps 1--4 to each
\(\ket{\psi_i}\). Branch \(m_1m_2\) then contributes
\(\sum_ip_i\cdot\tfrac12\ket{\psi_i}\cdot\tfrac12\bra{\psi_i}=\tfrac14\rho\),
and the four branches, with the bits forgotten, sum to \(\rho\). The
protocol is the identity channel. \(\;\blacksquare\)

The page is honest and complete, and it shows every amplitude. What it
does not show is \emph{why} the four brackets are the message, turned.
That is a fact about the Bell basis, visible only in the regrouping.

\hypertarget{coeckes-sketchbook-4}{%
\subsection{Coecke's Sketchbook}\label{coeckes-sketchbook-4}}

Fix a branch \(m_1m_2\); the bit-labelled rings will carry all four at
once. Every rule used below is exact.

\textbf{Step 0: substitution.} Alice's reading is the ringed cap,
\(\tfrac1{\sqrt2}\) in front (Result 10.2). Bob's pair is a cup,
\(\tfrac1{\sqrt2}\) in front (cup \(=\ket{00}+\ket{11}\), unnormalized).
Bob's correction box, on this branch, is red \(m_2\pi\) then green
\(m_1\pi\) (row 76). The two \(\tfrac1{\sqrt2}\) make the \(\tfrac12\).

\textbf{Step 1: yanking (Theorem 6.2).} The cap and the cup form an S;
pull it straight. The rings go with the wire, in the order the wire
meets them:

\[\tfrac12\;
\begin{tikzpicture}[sd, x=1.4em, y=1.4em,
  zsp/.style={circle,draw,line width=0.9pt,fill=green!55,minimum size=1.1em,inner sep=0pt},
  xsp/.style={circle,draw,line width=0.9pt,fill=red!75,minimum size=1.1em,inner sep=0pt},
  ph/.style={font=\scriptsize,inner sep=1pt}]
  \draw[wire] (0,3) -- (4.8,3);
  \sdcap{4.8}{3}{1.5}
  \draw[wire] (1.2,1.5) -- (4.8,1.5);
  \sdcup{1.2}{1.5}{0}
  \draw[wire] (1.2,0) -- (9.4,0);
  \node[zsp,label={[ph]above:{$m_1\pi$}}] at (1.6,3) {};
  \node[xsp,label={[ph]above:{$m_2\pi$}}] at (3.2,3) {};
  \node[xsp,label={[ph]below:{$m_2\pi$}}] at (6.4,0) {};
  \node[zsp,label={[ph]below:{$m_1\pi$}}] at (8.0,0) {};
\end{tikzpicture}
\;\;\overset{\text{yank}}{=}\;\;
\tfrac12\;
\begin{tikzpicture}[sd, x=1.4em, y=1.4em,
  zsp/.style={circle,draw,line width=0.9pt,fill=green!55,minimum size=1.1em,inner sep=0pt},
  xsp/.style={circle,draw,line width=0.9pt,fill=red!75,minimum size=1.1em,inner sep=0pt},
  ph/.style={font=\scriptsize,inner sep=1pt}]
  \draw[wire] (0,0) -- (7.6,0);
  \node[zsp,label={[ph]above:{$m_1\pi$}}] at (1.4,0) {};
  \node[xsp,label={[ph]above:{$m_2\pi$}}] at (3.0,0) {};
  \node[xsp,label={[ph]above:{$m_2\pi$}}] at (4.6,0) {};
  \node[zsp,label={[ph]above:{$m_1\pi$}}] at (6.2,0) {};
\end{tikzpicture}\]

\textbf{Step 2: fusion (Theorem 8.1), red meets red.} The two red rings
are neighbours; they fuse, and their phases add to \(2m_2\pi\).

\textbf{Step 3: the identity rule (§8.1).} \(2m_2\pi\) is \(0\) or
\(2\pi\), a whole turn either way, so the fused red dot is a phase-free
dot with one leg in and one out: a bare wire. Remove it. Now the two
green rings are neighbours.

\textbf{Step 4: fusion, green meets green}, to \(2m_1\pi\). \textbf{Step
5: the identity rule} again.

\[\overset{\text{fusion}}{=}\;\tfrac12\;
\begin{tikzpicture}[sd, x=1.4em, y=1.4em,
  zsp/.style={circle,draw,line width=0.9pt,fill=green!55,minimum size=1.1em,inner sep=0pt},
  xsp/.style={circle,draw,line width=0.9pt,fill=red!75,minimum size=1.1em,inner sep=0pt},
  ph/.style={font=\scriptsize,inner sep=1pt}]
  \draw[wire] (0,0) -- (6.0,0);
  \node[zsp,label={[ph]above:{$m_1\pi$}}] at (1.4,0) {};
  \node[xsp,label={[ph]above:{$2m_2\pi$}}] at (3.0,0) {};
  \node[zsp,label={[ph]above:{$m_1\pi$}}] at (4.6,0) {};
\end{tikzpicture}
\;\overset{\text{id}}{=}\;\tfrac12\;
\begin{tikzpicture}[sd, x=1.4em, y=1.4em,
  zsp/.style={circle,draw,line width=0.9pt,fill=green!55,minimum size=1.1em,inner sep=0pt},
  ph/.style={font=\scriptsize,inner sep=1pt}]
  \draw[wire] (0,0) -- (4.4,0);
  \node[zsp,label={[ph]above:{$m_1\pi$}}] at (1.4,0) {};
  \node[zsp,label={[ph]above:{$m_1\pi$}}] at (3.0,0) {};
\end{tikzpicture}\]

\[\overset{\text{fusion}}{=}\;\tfrac12\;
\begin{tikzpicture}[sd, x=1.4em, y=1.4em,
  zsp/.style={circle,draw,line width=0.9pt,fill=green!55,minimum size=1.1em,inner sep=0pt},
  ph/.style={font=\scriptsize,inner sep=1pt}]
  \draw[wire] (0,0) -- (2.8,0);
  \node[zsp,label={[ph]above:{$2m_1\pi$}}] at (1.4,0) {};
\end{tikzpicture}
\;\overset{\text{id}}{=}\;\tfrac12\;
\begin{tikzpicture}[sd]
  \draw[wire] (0,0) -- (2.4,0);
\end{tikzpicture}\]

A bare wire, times \(\tfrac12\): Bob's output is \(\tfrac12\ket\psi\) on
every branch. Van de Wetering's ZX proof of teleportation ends the same
way, with the remark that \(2b\pi\) is either \(2\pi\) or \(0\), both
zero modulo \(2\pi\) (his §5.4). The probability is the square of the
drawn scalar, \(\bigl(\tfrac12\bigr)^2=\tfrac14\); in the doubled
picture it is simply the number written in front,
\(\tfrac12\cdot\tfrac12\). \(\;\blacksquare\)

Nothing had to slide and nothing had to pass anything else. The order of
Bob's corrections, flip first and phase-flip second, is exactly what
makes like rings neighbours. In the other order the straight wire would
read green, red, green, red, and the middle pair could meet only by
passing a green half-turn through a red one, at the cost of a sign.

\hypertarget{the-verdict-9}{%
\subsection{The verdict}\label{the-verdict-9}}

The picture wins, and it wins on both counts: it is shorter, and it
shows why. One yank and two fusions settle all four outcomes at once,
and the drawing makes the reason plain. The wire simply runs from
Alice's input to Bob's output, and the only things ever on it are the
reading's half-turns and the letter's half-turns, which meet their own
kind and vanish. The blackboard needs a page and an eight-term
regrouping, and at the end we know \emph{that} each bracket is the
message turned, not why.

But the picture has a debt, and we name it. That Alice's reading
\emph{is} a cap carrying those two rings is a translation: Result 10.2,
a four-effect check on the blackboard or a derivation by colour change
and fusion from the CNOT and the Hadamard. The \(\tfrac14\) is a scalar
the picture carries rather than explains. And what the picture proves is
an equality of processes, no more. It makes no claim about speed or
about security, and it sends nothing faster than light; the letter is
needed, as §10.2 showed.

\begin{center}\rule{0.5\linewidth}{0.5pt}\end{center}

\hypertarget{the-mended-junction}{%
\section{§10.7 The Mended Junction}\label{the-mended-junction}}

Chapter 6 ended with a junction. Two pairs, a cup on wires \(d,a\) and a
cup on wires \(b,c\); a cap on the inner wires \(a,b\); and the outer
wires \(d\) and \(c\), which had never met, pulled into a pair by a
single yank. The cap was a postselected answer, one in four, and each of
the other three, Chapter 6 said, left ``a pair that is a known one-box
correction away'' from the one we wanted. (The history of that junction,
proposed and then performed, is told in Chapter 6's Feynman's Notebook,
§6.8; it is not repeated here.) This chapter adds the two things Chapter
6 lacked: the three other answers, each corrected, and the record of
which answer came, carried as classical wires and drawn in the doubled
form of Chapter 9.

The junction is teleportation with a particular message. The wire \(a\)
that Alice reads together with her half \(b\) of the shared pair is
itself half of another pair, \((d,a)\). So the junction teleports
\emph{half a pair}. And Theorem 6.5 said a pair is a bent wire, the
identity bent into a state. Teleporting half of it is teleporting the
identity; this is the headline, bent.

\textbf{On the blackboard.} Wires top to bottom \(d,a,b,c\); the state
is \(\ket{\Phi^+}_{da}\otimes\ket{\Phi^+}_{bc}\); Alice reads \((a,b)\);
Bob corrects \(c\). On branch \(m_1m_2\),

\[\bigl(I_d\otimes\bra{B_{m_1m_2}}_{ab}\otimes(Z^{m_1}X^{m_2})_c\bigr)\bigl(\ket{\Phi^+}_{da}\otimes\ket{\Phi^+}_{bc}\bigr)\;=\;\tfrac12\ket{\Phi^+}_{dc},\]

and without the correction it is \(\tfrac12\ket{B_{m_1m_2}}_{dc}\): the
pair left at the far ends is \emph{the very Bell state the reading
named}. Both follow from Result 10.5 applied to the wire \(a\), with
\(d\) carried along.

\textbf{In the sketchbook.} On one branch, with three factors of
\(\tfrac1{\sqrt2}\) in front (two pairs and the reading), the wire runs
out of \(d\), round the first cup, along \(a\) through the reading's
rings, round the cap, back along \(b\), round the second cup, and out
along \(c\) through Bob's rings. Yank it:

\[\tfrac1{2\sqrt2}\;
\begin{tikzpicture}[sd, x=1.4em, y=1.4em,
  zsp/.style={circle,draw,line width=0.9pt,fill=green!55,minimum size=1.1em,inner sep=0pt},
  xsp/.style={circle,draw,line width=0.9pt,fill=red!75,minimum size=1.1em,inner sep=0pt},
  ph/.style={font=\scriptsize,inner sep=1pt}]
  \sdcup{0}{4.5}{3}
  \sdcup{0}{1.5}{0}
  \draw[wire] (0,4.5) -- (8.6,4.5);
  \draw[wire] (0,3) -- (4.4,3);
  \draw[wire] (0,1.5) -- (4.4,1.5);
  \draw[wire] (0,0) -- (8.6,0);
  \sdcap{4.4}{3}{1.5}
  \node[zsp,label={[ph]above:{$m_1\pi$}}] at (1.2,3) {};
  \node[xsp,label={[ph]above:{$m_2\pi$}}] at (2.8,3) {};
  \node[xsp,label={[ph]below:{$m_2\pi$}}] at (5.6,0) {};
  \node[zsp,label={[ph]below:{$m_1\pi$}}] at (7.2,0) {};
  \node[sdlabel,anchor=west] at (8.8,4.5) {$d$};
  \node[sdlabel,anchor=west] at (8.8,0) {$c$};
\end{tikzpicture}\]

\[\overset{\text{yank}}{=}\;
\tfrac1{2\sqrt2}\;
\begin{tikzpicture}[sd, x=1.4em, y=1.4em,
  zsp/.style={circle,draw,line width=0.9pt,fill=green!55,minimum size=1.1em,inner sep=0pt},
  xsp/.style={circle,draw,line width=0.9pt,fill=red!75,minimum size=1.1em,inner sep=0pt},
  ph/.style={font=\scriptsize,inner sep=1pt}]
  \sdcup{0}{1.8}{0}
  \draw[wire] (0,1.8) -- (7.2,1.8);
  \draw[wire] (0,0) -- (7.2,0);
  \node[zsp,label={[ph]below:{$m_1\pi$}}] at (1.2,0) {};
  \node[xsp,label={[ph]below:{$m_2\pi$}}] at (2.8,0) {};
  \node[xsp,label={[ph]below:{$m_2\pi$}}] at (4.4,0) {};
  \node[zsp,label={[ph]below:{$m_1\pi$}}] at (6.0,0) {};
\end{tikzpicture}\]

Then red meets red and green meets green, exactly as in §10.6, and a
bare cup on \((d,c)\) is left: \(\tfrac1{2\sqrt2}\,\)cup
\(=\tfrac12\ket{\Phi^+}\). Stop before Bob's rings and the cup still
carries the reading's green and red, which is
\(\tfrac12\ket{B_{m_1m_2}}\) in its bottom-leg form (Result 10.1).

Now the whole thing, every branch together, with the record kept. Two
boxes on doubled wires do the work. The \textbf{reading box} \(M\) takes
Alice's two doubled wires, four lines, and emits two classical wires,
\(x\) carrying \(m_2\) and \(z\) carrying \(m_1\); on the blackboard it
is
\(M=\sum_{m_1,m_2}\ket{m_2m_1}_{xz}\bigl(\bra{B_{m_1m_2}}\otimes\bra{B_{m_1m_2}}\bigr)\),
the Bell effect on the kets and its mirror on the mirror copies (the
Bell states are real, so the conjugate changes nothing). It is the
circuit of §10.1 followed by two of Chapter 9's looks. The
\textbf{correction box} \(C\) is the classically controlled box of
§10.3, on Bob's doubled wire, with \(x\) and \(z\) running in and on
through it. The two doubled pairs are \(\tfrac12\,\)cup\(\otimes\)cup
each (§9.5), so \(\tfrac14\) is written in front of the four cups; as
ever, cup \(=\ket{00}+\ket{11}\).

\begin{center}
\begin{tikzpicture}[sd]
  \node[sdlabel, font=\normalsize] at (-3.4,5.25) {$\tfrac14$};
  \sdcup{0}{10.5}{6}
  \sdcup{0}{9}{7.5}
  \sdcup{0}{4.5}{0}
  \sdcup{0}{3}{1.5}
  \draw[wire] (0,10.5) -- (14,10.5);
  \draw[wire] (0,9) -- (14,9);
  \draw[wire] (0,7.5) -- (4,7.5);
  \draw[wire] (0,6) -- (4,6);
  \draw[wire] (0,4.5) -- (4,4.5);
  \draw[wire] (0,3) -- (4,3);
  \draw[cwire] (5.2,4.5) -- (14,4.5);
  \draw[cwire] (5.2,3) -- (14,3);
  \draw[wire] (0,1.5) -- (14,1.5);
  \draw[wire] (0,0) -- (14,0);
  \node[sdbox, minimum height=6.3em, minimum width=2.4em] at (4.6,5.25) {$M$};
  \node[sdbox, minimum height=6.3em, minimum width=2.4em] at (9.6,2.25) {$C$};
  \node[sdlabel,anchor=west] at (14.2,10.5) {$\bar d$};
  \node[sdlabel,anchor=west] at (14.2,9) {$d$};
  \node[sdlabel,anchor=west] at (14.2,4.5) {$x$};
  \node[sdlabel,anchor=west] at (14.2,3) {$z$};
  \node[sdlabel,anchor=west] at (14.2,1.5) {$c$};
  \node[sdlabel,anchor=west] at (14.2,0) {$\bar c$};
  \node[sdlabel,anchor=south] at (2.2,7.55) {$a$};
  \node[sdlabel,anchor=south] at (2.2,6.05) {$\bar a$};
  \node[sdlabel,anchor=south] at (2.2,4.55) {$\bar b$};
  \node[sdlabel,anchor=south] at (2.2,3.05) {$b$};
\end{tikzpicture}
\end{center}

The lines are the keepsake half \(d\) and its mirror \(\bar d\), running
straight out; the message half \(a\) and its mirror \(\bar a\), from the
first pair into \(M\); Alice's half \(b\) and its mirror \(\bar b\),
from the second pair into \(M\); Bob's half \(c\) and its mirror
\(\bar c\), through \(C\) and out; and the two classical wires, \(x\)
for the flip bit and \(z\) for the phase bit, from \(M\) through \(C\)
and out, the record kept. The mirror pairs are drawn so that their cups
nest and no line crosses another. We call this diagram \textbf{The
Mended Junction}.

\begin{quote}
\textbf{Result 10.8 (the mended junction).} The Mended Junction
evaluates, exactly, to
\[\tfrac14\sum_{m_1,m_2}\ket{m_2m_1}\!\bra{m_2m_1}_{xz}\;\otimes\;\ket{\Phi^+}\!\bra{\Phi^+}_{dc}.\]
The record is a fair four-way toss; on every one of the four outcomes
the pair \((d,c)\) is exactly \(\ket{\Phi^+}\); and the two are
independent. Forget the record and \((d,c)\) is exactly
\(\ket{\Phi^+}\!\bra{\Phi^+}\). Remove the correction box and forget the
record, and \((d,c)\) is \(\tfrac14 I\), no pair at all.
\end{quote}

\emph{Proof idea.} Fix the record's value and the diagram falls apart
into the one-branch picture above and its mirror image. Each yanks to a
cup on \((d,c)\), and what is left to do is count scalars.

\emph{Proof.} On the branch \(m_1m_2\) the classical wires carry a
definite value, \(M\) becomes the Bell effect on the kets and on the
mirrors, and \(C\) becomes the two rings on \(c\) and their conjugates
on \(\bar c\). The doubled diagram is the pure one-branch diagram beside
its mirror copy. Count the numbers: \(\tfrac14\) from the drawn pairs;
\(\tfrac1{\sqrt2}\cdot\tfrac1{\sqrt2}=\tfrac12\) from the normalized
Bell effect on the two copies. The bare cups and ringed caps yank, and
fuse, to cup \(\otimes\) cup on \((d,c)\) and \((\bar d,\bar c)\), which
is \(2\,\ket{\Phi^+}\!\bra{\Phi^+}\). So each branch contributes
\(\tfrac14\cdot\tfrac12\cdot2=\tfrac14\,\ket{\Phi^+}\!\bra{\Phi^+}\),
tagged by its record \(\ket{m_2m_1}\!\bra{m_2m_1}\); summing over the
four gives the formula. Forgetting the record adds the four quarters,
giving \(\ket{\Phi^+}\!\bra{\Phi^+}\). Without the correction each
branch leaves \(\tfrac14\ket{B_{m_1m_2}}\!\bra{B_{m_1m_2}}\), and by the
completeness of Result 10.1 the four sum to \(\tfrac14 I\).
\(\;\blacksquare\) See \texttt{Code/Ch10/MendedJunction.py}, which
builds the diagram in doubled form with its classical record, tallies
its lines, and checks the value exactly.

So the promise of Chapter 6 is kept. The junction no longer works one
time in four: it works every time, provided the letter is obeyed. The
keepsake half \(d\) and Bob's half \(c\) come out as one pair, although
they were never in the same hand, and no quantum thing crossed from
Alice to Bob; only the two bits did. The uncorrected branches were not
failures. Each was a named Bell state, one known turn away from the pair
we wanted. And forgetting the record without correcting leaves
\(\tfrac14 I\), the grey of two uncorrelated coins, which is Result 10.4
again, one level up.

\begin{quote}
\textbf{Coecke's Sketchbook.}

\emph{Keeping the scalars.} Abramsky and Coecke's own diagrammatic
proofs of teleportation and of entanglement swapping say, each time,
``We ignore the scalars -- which cancel out against each other -- in
this proof'' (their Appendix). In those proofs the scalars are set
aside. They do keep them in the statement: their teleportation sends the
input to the output along each of the four branches with a weight
\(s^{\dagger}s\), where \(2s^{\dagger}s=1\), which is our \(\tfrac12\),
``corresponding to the probability amplitude for that branch'' (their
§9.1). This book writes the \(\tfrac12\) and the \(\tfrac14\) into the
pictures, because here they are the whole point. They are the
probabilities. A picture that drops them proves that the protocol is
right on each branch; a picture that keeps them proves also that the
branches happen a quarter of the time each, and add up to certainty.
\end{quote}

\begin{quote}
\textbf{Feynman's Notebook.}

\emph{Teleporting a gate.} Abramsky and Coecke treat two cousins of
teleportation alongside it. One is entanglement swapping, whose history
this book told in Chapter 6. The other is logic-gate teleportation,
which they credit to Gottesman and Chuang, ``Quantum teleportation is a
universal computational primitive'' (arXiv:quant-ph/9908010, 1999;
Abramsky and Coecke's reference {[}12{]}), published in \emph{Nature}
402, 390--393 (1999) under a different title. Replace the shared pair by
some other entangled state, and the qubit that arrives is not the input
but a function of it, fixed by the state. Done with pairs of qubits in
place of single ones, Abramsky and Coecke report (their §2.2), this
gives a universal computational primitive, one that is moreover
fault-tolerant. In the sketchbook the idea is one bend. Put a box on the
cup, and by Theorem 6.5 the box rides along the wire and is applied to
whatever is teleported through it; the corrections must then be chosen
to suit the box.
\end{quote}

\begin{center}\rule{0.5\linewidth}{0.5pt}\end{center}

\hypertarget{what-travelled}{%
\section{§10.8 What Travelled}\label{what-travelled}}

Let us gather it up. A Bell reading has four answers, the four Bell
states: Chapter 6's bent wire, plain or carrying a flip, a phase-flip,
or both, and one drawing with two bit-labelled rings holds all four
(§10.1). The reading's effect is a cap with those rings, a translation
we paid for with colour change and fusion. Read half of a pair against a
message and the message is at the far end at once, turned by the
reading's rings, and useless there, a flat grey, until the outcome
arrives (§10.2). Send the outcome, two bits, and the far end turns the
message back, flip first and phase-flip second; in the picture the wire
runs straight through, like rings meet and vanish, and \(\tfrac12\) is
left, whose square is the \(\tfrac14\) of each branch (§10.3, §10.6).
Run the same pieces the other way round and one qubit carries two bits,
the correction table serving as the codebook and the preparations as the
readings' mirror images (§10.4). And teleport half of a pair, and
Chapter 6's junction is mended on every outcome (§10.7).

What travelled? Not the qubit: nothing quantum went from Alice to Bob.
Not a description of it either, since the two bits are a fair toss
whatever the message was. What travelled was \emph{which turn}, and the
rest was already there, in a wire that was bent long ago.

One thing more has happened, and the dialogue lets it pass in silence.
When the reading was done, Alice's message was gone from her end. Her
wire gave up its state in the reading; as Abramsky and Coecke put it,
the information in the source ``has been `destroyed' in transferring it
to the target''. Teleportation \emph{moves} a state. At no moment in the
whole procedure were there two copies of \(\ket\psi\), one at each end.

Was that bad luck, or bad design, or is it a law? Could a cleverer
protocol have left Alice her message and sent Bob a second one? Could
anything, anywhere, take an unknown qubit and make two? The next chapter
asks, and the answer is written into the very linearity that made every
picture in this one work.

\hypertarget{gigue-that-may-not-be-repeated}{%
\chapter{Gigue That May Not Be
Repeated}\label{gigue-that-may-not-be-repeated}}

\begin{center}\rule{0.5\linewidth}{0.5pt}\end{center}

\emph{落ちし玉}\\
\emph{二度は落ちざる}\\
\emph{氷柱かな}

\emph{ochishi tama}\\
\emph{nido wa ochizaru}\\
\emph{tsurara kana}

\emph{A bead once fallen}\\
\emph{does not fall a second time ---}\\
\emph{the eaves' icicle.}

\begin{center}\rule{0.5\linewidth}{0.5pt}\end{center}

\emph{{[}The Crab's parlour, the same day, the hour after the midday
meal. The black slab still lies along the table, wiped and bare. Beside
it, off the black, stand the viewer, fetched down from the dresser, and
a little wooden tray with three pewter counters stacked in it. The
Tortoise's drawing-cloth runs beside the slab, her string and rings on
it. Achilles stands at the window.{]}}

\emph{{[}The thaw has come at noon, as it will on the brightest of
deep-winter days. The sun, low even at its highest, stands full on the
house front, and the parlour is lit from its near end, through the one
tall window, while the far end keeps a cooler, quieter light of its own.
The frost that furred the glass at breakfast has gone from it
altogether; the small panes are clear, and every one of them, blown by
hand long before the Crab's time, bends the lane a little differently
--- the gate leans in one, the hedge swims in the next, and no two show
quite the same lane. Along the eaves above the window hangs a fringe of
icicles, and the sun has set them all going: a drop here, two there, a
long wait, then three at once, at no pace anyone could follow. Each drop
falls once, past the glass, and is gone into the snow on the sill. The
lace curtain is drawn back to one side of the window and hangs there in
loose folds. On the long wall the fire is built up high, the kettle
stands at the side of the hob, and the dresser holds the shut ledger,
the cake-tin, and the Crab's cups --- he buys them one at a time,
wherever he happens to be, so that there are no two of a pattern on the
shelf. Down the middle of the table the black slab runs from the window
to the far end, bare. In the tray beside it the pewter counters sit in
their stack: the top one new and bright, the middle one dulled with
handling, the bottom one with a small nick in its rim. On the end wall
beyond the far end of the table, facing the window down the whole length
of the room, hangs a small print in a plain frame --- Escher's
\textup{Drawing Hands}: two hands, each drawing the other. The noon
light reaches it only at second hand, pale, off the ceiling. It has hung
there since long before the slab was first laid out on that table, and
it will hang there after. The Tortoise has kept this morning's quilt
round her shoulders, and the room smells of the pudding and, faintly, of
wet wool.{]}}

\textbf{Achilles:} Will it draw the very same pattern tonight, do you
think, Madam Tortoise --- line for line, just as it lay on the pane this
morning?

\textbf{Tortoise:} Not in all the winters I have watched it. It draws a
new one every night and never looks back at the old.

\textbf{Crab:} \emph{{[}stacking the last plates{]}} The eaves are no
better. They've dripped since the pudding, and I've stopped hearing
them.

\textbf{Achilles:} The whole lane is running with it. I walked down to
the gate after lunch to look at the thaw and brought half of it back in
my boots.

\textbf{Crab:} Then put them by the fender and your feet by the fire,
and let the rug forgive you. \emph{{[}He carries the plates out and
comes back drying his claws.{]}} This morning's keepsake pair is in the
drawer, unread, and there it stays; the pairing-glass is on its shelf.
The table is clean for once.

\emph{{[}Achilles sits down on the fender stool and works his boots off,
and stands them on the hearth, where they begin at once to steam and to
leak a small dark pool of thaw onto the stone. He puts his stockinged
feet out to the fire and wriggles his toes at it. The Crab, passing
behind him, drops a dry towel over his knees without a word.{]}}

\textbf{Tortoise:} Clean, and waiting for something. You fetched the
viewer down, and you've looked at the clock four times since the soup.

\textbf{Crab:} Because the midday post is late, and the midday post is
bringing me a parcel. The old man at the workshop wrote to say so.

\emph{{[}A knock at the door. The Crab is out of his chair before it has
finished, and comes back with a squarish parcel in brown paper, stiff
with frost at the corners, straw showing at a torn seam.{]}}

\textbf{Crab:} From Delft. The postman made me promise not to drop it on
his foot.

\textbf{Achilles:} Is it heavy?

\textbf{Crab:} Heavier than anything he has ever sent me. \emph{{[}He
cuts the string and folds back the paper on a nest of straw. From the
top of the straw he lifts a small box and opens it.{]}} And, before
anything, a fresh pair --- cut of one piece, as always. He never sends a
thing without a pair to try it on.

\emph{{[}He sets one stone at the window end of the slab. Its reflection
is there at once in the black, beneath it, meeting it at the surface.
The other he carries down the whole length of the table and sets at the
far end, and its reflection rises to meet it there. Two stones, a long
table apart, each doubled in the glass; neither has been looked at.{]}}

\textbf{Crab:} One here, one there. The rest of the parcel we'll come
to.

\emph{{[}The two stones sit a table's length apart. The near one, at the
window end, has the sun full on it and is plain to see from anywhere in
the room; the far one, down by the end wall, sits in the cooler light
there, and from the fire it looks smaller than it is. Achilles, his feet
still to the blaze, keeps turning his head from the one to the other, as
though they were having a conversation and he was missing it.{]}}

\begin{center}\rule{0.5\linewidth}{0.5pt}\end{center}

\hypertarget{first-challenge-the-press-copies-its-own}{%
\section{\texorpdfstring{First Challenge --- \emph{The Press Copies Its
Own}}{First Challenge --- The Press Copies Its Own}}\label{first-challenge-the-press-copies-its-own}}

\emph{{[}The Crab lifts out of the straw a small iron press: a hopper
with a lid on top, a lever at the side, a hood with a shutter over the
middle, and two little chutes below, one to the left and one to the
right. He sets it on the black at the far end, beside the far stone.
Under it, upside down in the glass, stands the press's reflection,
hopper to hopper, chutes to chutes.{]}}

\textbf{Crab:} The copying press. The old man has been building it since
the summer, and I have been asking him for one since before I could see
over his bench. Put a stone in this hopper, pull the lever once, and two
stones drop from these chutes.

\textbf{Achilles:} Two from one? Where does the second come from?

\textbf{Crab:} From in here. \emph{{[}He taps the housing.{]}} It comes
loaded with blanks --- raw stones, settled low at the workshop's bench
and sealed inside --- and it eats one on every strike.

\textbf{Achilles:} So the stone goes in, and it comes out with a twin.

\textbf{Crab:} Two come out. Which of them, if either, is the one that
went in, nobody can say.

\emph{{[}In the noon light the press is a lovely small thing: blacked
iron polished like a good stove-front, a hopper lid of brass, a lever
with a knob of turned boxwood worn pale by the maker's hand. It has come
all the way from Delft in straw, and still it stands so finely on its
four small feet that a cart going by in the lane sets a faint shiver in
it. Under the lever, struck small into the iron, is the workshop's sign,
a little anchor, and beside it the old man's initials, their edges gone
soft. And at the seam of the hood, seen only from the side of the eye,
there is now and then a faint warm glow, as of a coal banked somewhere
in the works; looked at straight, it is only a join in blacked iron.{]}}

\emph{{[}He lifts the hopper's lid: inside, a small dark cup, empty. He
lets the lid down. Then he sets a claw to the shutter on the hood and
slides it open. Behind it there is only dark iron --- no stone, no
light, nothing turning. He slides it shut again. Nothing has moved; the
press stands as it stood.{]}}

\textbf{Crab:} Open, while it idles, it shows you dark works and nothing
else. Open it while it strikes, and it jams, and the stone in the hopper
is spent for nothing. So it is worked blind, and I shall never see how
it does it. \emph{{[}He looks at the shutter a moment longer than he
needs to.{]}} I knew that when I asked for it.

\emph{{[}When at last he looks away, the glow at the seam is back, just
where the eye cannot quite get at it. The Tortoise, who has been
watching the Crab rather than the press, hands him his cup, and he
drinks from it without noticing that he has.{]}}

\textbf{Achilles:} All this morning we moved a stone from one end of the
room to the other, and not once were there two of it. And here is a
machine that makes two. So it can be done after all.

\textbf{Crab:} Read what it does before you say so. \emph{{[}From the
straw he takes a slim, cloth-bound book.{]}} The maker's proving-book.
Every pressing made at the workshop, one line each: how the stone was
settled beforehand, and how its two struck stones read afterwards.

\textbf{Tortoise:} Settled --- not looked at?

\textbf{Crab:} Settled. That is the workshop's bench making a stone a
certain way, not reading it; it spends nothing.

\textbf{Crab:} \emph{{[}laying a hand on the viewer{]}} And they read
the struck stones through a viewer like mine. \emph{{[}He turns its
collar a quarter round and back.{]}} This collar I never showed you last
night: upright, it asks a stone the plain question, high or low; turned
crosswise, it asks the crosswise question, and the needle swings to the
same two marks. Either way, a reading spends the stone.

\textbf{Crab:} \emph{{[}a claw on the little tray{]}} And these counters
keep the readings. Whoever reads a stone lays one on the table's edge,
face-up for high and face-down for low; stacked in the tray, they mean
nothing at all.

\textbf{Crab:} \emph{{[}taking a letter from the straw and setting it,
unopened, by the book{]}} There's a letter from him too, mostly about a
clerk called Eve. He says I'm to hear it before I use the press for
anything clever.

\textbf{Achilles:} Then the book first. What did the plain stones do?

\textbf{Crab:} \emph{{[}opening the book and running a claw down the
first page{]}} Settled plain-high, pressed, and both struck stones read
high. Settled plain-low: both low. High and high; low and low; high and
high --- all the way down the page and over the next.

\textbf{Achilles:} One stone in, two stones out: that's all a copy is.
And now the crosswise ones. A crosswise stone is half high and half low,
near enough, and a press that copies high and copies low will copy
anything in between. So out come two crosswise stones, and each, read at
cross, will say what its original would have said, every time.

\textbf{Crab:} \emph{{[}turning a page{]}} Settled crosswise one way,
pressed, and the struck pair read at cross. The first struck stone: high
mark; low mark; low; high; low. The second, just the same kind of
scatter. A coin toss, each of them.

\emph{{[}Outside, the sun has got well onto the eaves, and the drip has
quickened into a patter on the sill: a light, quick, running figure that
seems each moment about to come round to its beginning again, and never
quite does. Achilles, by the fire, taps it on his knee for a bar or two
and loses it.{]}}

\textbf{Achilles:} Each? But a stone settled crosswise, and asked the
crosswise question, gives the same answer every time.

\textbf{Crab:} Its struck stones don't. And yet look along the lines,
not down them. Whenever the first said high, the second said high;
whenever low, low --- not one line where they differed. And where the
stone went in settled crosswise the other way, the two always disagreed,
every line of it.

\textbf{Achilles:} What did they do when they were asked the plain
question?

\textbf{Crab:} There are some pages of that at the back. \emph{{[}He
finds them.{]}} Settled crosswise, either way, and the struck pair read
plain: both high or both low, half the time each. Never one of each.

\textbf{Achilles:} Not a single high-and-low. Two honest crosswise
stones, read plain, would be two separate tosses, one high and one low
as often as not. It copied the high part into both-high and the low part
into both-low, and there's nowhere left for one-of-each to come from.
\ldots Oh. It doesn't copy the crosswise stone. It makes a \emph{pair}
of it.

\textbf{Tortoise:} Let me show you on the cloth, where nothing gets
spent. \emph{{[}She takes from her shell a green ring with one cord
running in and two running out.{]}} You know this from Mr Crab's beam: a
green ring made to fork, one cord in and two out, and it does to its own
kind what his press does. \emph{{[}She lays beside it two short cords,
each ending in a single-legged ring, one red, one green.{]}} And these
are points: a red point is a stone settled plain, and a green point is a
stone settled crosswise.

\emph{{[}The fork lies on the cloth, one cord in from the left, two out
to the right. The Tortoise threads the red point's cord onto the fork's
in-cord and slides the red ring up it. The red ring meets the green and
does not stop: the green ring opens round it and is gone. On each of the
two out-cords a red ring rides out to the end and comes to rest there.
Where there was one red point and a green fork, there are now two red
points, each the very twin of the one that went in; the cords are no
longer and no shorter, and only the green ring has vanished.{]}}

\textbf{Achilles:} Two, and both red. That's the plain page of the book.

\emph{{[}She sets out a fresh fork and threads the green point onto it
instead. Green meets green: the two rings slide into each other, as like
rings always do, and are one ring. The point's own cord is swallowed
with it, and what is left is a single green ring with just the two
out-cords hanging from it. A ring with two cords and no turn on it is
only a bend; she lifts the ring off, and on the cloth lies one length of
string, bent round on itself, both its ends running out to the right.
Nothing has slid out along either cord. Where two green points might
have lain, there lies one bend --- two ends, but one piece.{]}}

\textbf{Tortoise:} Red goes through as two. Green goes in, and one bend
comes out: a pair cut of one piece, not two copies of anything.

\emph{{[}Somewhere above them the thaw loosens the drift that has lain
along the roof-ridge since the night, and it goes all at once, down both
slopes together: a soft thump into the lane at the front of the house
and another into the garden at the back, so exactly together that they
make one sound, heard from two sides. Achilles looks up at the ceiling,
and then at the window, where a little loose powder is still sifting
past the glass.{]}}

\textbf{Achilles:} So this press copies only its own kind. But a better
press, one cleverer inside --- might that copy anything?

\begin{center}\rule{0.5\linewidth}{0.5pt}\end{center}

\hypertarget{second-challenge-one-loop-two-loops}{%
\section{\texorpdfstring{Second Challenge --- \emph{One Loop, Two
Loops}}{Second Challenge --- One Loop, Two Loops}}\label{second-challenge-one-loop-two-loops}}

\textbf{Crab:} That is exactly what I mean to ask the old man. He swears
this one loses nothing --- whatever goes in is all accounted for in what
comes out --- but he could build round it, surely, with better works.

\textbf{Tortoise:} Then let me draw any press whatever, as clever as he
pleases. \emph{{[}She takes a small card from her shell and threads it
on a length of string.{]}} This card is the press --- its works, its
blanks, whatever it keeps inside, all folded into one. \emph{{[}She lays
a second card on the first, traces it through, turns the tracing over,
and threads that on too.{]}} And this is its mirror card, traced from it
back to front: the same press, run the other way about.

\textbf{Tortoise:} One more thing to lay out. A loop of string with a
stone's point at one end and, closing it, the mirror of another stone's
point, is how alike the two stones are: a whole loop if they are the
very same, nothing at all if they are wholly unalike, and something
between for the rest.

\emph{{[}She makes the loop. At one end of the string she fixes a red
point --- a stone settled plain-low. From it the string runs through the
card, then through the mirror card, and on to the mirror of a green
point, a stone settled crosswise, where she brings it round and ties it
back to the start. A closed loop, with two points on it and the two
cards between them.{]}}

\textbf{Tortoise:} That is how alike the two stones are before any
pressing --- the press and its mirror both in the way. A fair press that
loses nothing always meets its own mirror like this.

\emph{{[}She slides the card along the string towards the mirror card.
The string does not move; only the card travels. The two cards touch,
face to face, and she lifts them off the string together, as one. The
loop lies closed as before, the red point and the mirrored green point
still on it, the string no shorter than it was; only the cards are
gone.{]}}

\textbf{Tortoise:} One loop. The press has left no mark on it at all:
its likeness is just what it was before.

\emph{{[}Then, beside it, she lays out the same two stones after a
faithful pressing: each stone now two copies. The first copy of the
plain stone she closes on the first copy of the crosswise one, in a loop
of its own; the second copy on the second, in another. Two separate
loops lie side by side, not touching, each with its red point and its
mirrored green point. A little way off she lays a small plain loop.{]}}

\textbf{Tortoise:} And that little one is whatever the works kept for
themselves. At best it is worth a whole, and never more.

\textbf{Tortoise:} One loop is the likeness; two loops side by side are
the likeness taken twice over, and multiplied.

\textbf{Crab:} Say it again, slowly, so I have it exactly. No --- let
me. One loop is how alike they were, and two loops side by side are that
same likeness laid down twice and multiplied --- didn't you say?

\textbf{Achilles:} \emph{{[}laughing{]}} Close, Mr Crab.

\textbf{Crab:} Bother. I had it word for word in my head, and it came
out of my mouth another shape.

\textbf{Tortoise:} Keep yours, Mr Crab. It came out a new one, and a
better one than mine.

\emph{{[}The Crab, mollified, laughs under his breath and fills all
three cups from the pot, and the Tortoise folds both claws round hers.
Outside, the drip has settled: each icicle along the eaves has found a
pace of its own and keeps to it --- one quick, one slow, one that waits
and waits and then lets go --- and none of them is in step with
another.{]}}

\textbf{Achilles:} Then let me put numbers on your loops. A plain-low
stone and a crosswise one: ask the crosswise stone the plain question
and it says low half the time, so call them half alike. Before the
press, one loop: a half. After a faithful pressing, two loops: a half of
a half, a quarter --- and the works' little loop can only make that
smaller. A half against a quarter.

\textbf{Tortoise:} And the press, which loses nothing, must keep the
half it was given.

\textbf{Crab:} Then it can't be done for those two. But for others ---
any two at all?

\textbf{Tortoise:} One loop equals two only when the loop is worth all
or nothing: stones the very same, or wholly unalike. For any two
between, there is no press.

\emph{{[}The sun lies full on the lace curtain where it hangs gathered
at the side of the window. Where the lace lies single, the lane shows
through it dimmed; where a fold lies over it, dimmer again, by the same
share once more; and in the thick bunch at the hook no lane comes
through at all. Only that dark bunch, and the bare glass beside the
curtain, would look no different for one more fold.{]}}

\textbf{Crab:} Not even if the old man made it cleverer inside?

\textbf{Tortoise:} The cleverness lives in the card, and the card came
off.

\textbf{Crab:} \emph{{[}sitting back, and looking at the letter by the
proving-book{]}} Then Eve's press --- the letter says she had one ---
never copied a thing either. So what was she doing with it all those
years?

\begin{center}\rule{0.5\linewidth}{0.5pt}\end{center}

\hypertarget{third-challenge-eves-tap}{%
\section{\texorpdfstring{Third Challenge --- \emph{Eve's
Tap}}{Third Challenge --- Eve's Tap}}\label{third-challenge-eves-tap}}

\textbf{Crab:} \emph{{[}breaking the seal and unfolding the letter{]}}
Let us find out. He writes a small hand. \emph{{[}reading{]}} ``You will
remember Eve, who clerked for the carrier --- the slow, careful one that
takes our stones about the country once a season. She had a press like
yours. For three winters she pressed every stone that passed through her
hands, kept one of the struck stones in her drawer, and sent the other
on to the customer.''

\textbf{Achilles:} A thief! And nobody noticed?

\textbf{Crab:} \emph{{[}reading on{]}} ``The customers whose stones were
settled plain noticed nothing: they came as good as ever. But every
stone that was sent settled crosswise, either way, arrived looked-at.
Asked the crosswise question, it answered like a coin.'' \emph{{[}He
lowers the letter.{]}} Looked-at, every one. People paid for those
stones, Achilles, and waited a season for them.

\textbf{Achilles:} Wait, though --- why were the plain ones spared?
\emph{{[}He thinks.{]}} Nothing in them to spoil: they were settled
already. A press that copies high and low leaves high and low alone, and
looks at everything in between.

\emph{{[}Through the clear panes the lane lies in the full thaw. The
morning's footprints go along it still, every one still there, a hollow
wherever a foot came down; but the thaw has gone softly into each of
them and taken out the tread and the nail-heads, so that whose boots
they were can no longer be told from any one of them.{]}}

\textbf{Tortoise:} Let me draw it as your slab would show it, Mr Crab,
with everything doubled. \emph{{[}She lays a second length of string
beside the first on the cloth, and runs it the other way.{]}} This is
the mirror string: it carries the reflection of whatever I lay on the
first. And to let a thing go, as we did last night, I fold its end over
onto its own mirror twin.

\emph{{[}On the first string she threads a fork: the stone's cord coming
in, two copies going out. On the mirror string, facing it, she threads
the fork's reflection, the stone's reflection coming in, the two copies'
reflections going out. Then she takes the second copy's end and folds it
over onto the second copy's reflection: they meet end to end, and the
two forks are now joined, fork to mirror fork, through that single
folded string. She slides the forks together along it. Green meets
green; the two rings run into each other and are one. What lies on the
cloth is a single green ring with four cords: the stone's and its
reflection's coming in, the kept copy's and its reflection's going out.
What has changed: the second copy and its reflection are gone into the
fold, and the two forks have become one ring. What has held: a stone and
its reflection still come in, and one stone and its reflection still go
out.{]}}

\textbf{Tortoise:} Stone and reflection in, stone and reflection out,
and one ring holding all four together.

\textbf{Achilles:} I've seen that. Last night --- the stone I looked at
with the needle hidden behind the cloth, and never read. A look that
nobody read.

\textbf{Tortoise:} Laid out in string, it is that, exactly.

\textbf{Crab:} So her press copied nothing. It looked, and kept the look
in her drawer where no one could read it, and sent my customers the
leavings.

\emph{{[}He sits a moment with the letter on his knee, looking at
nothing; Achilles leans forward from the fender and puts another piece
of coal on the fire for him. On the sill below the window each icicle
has worn its own small round pit in the snow, a neat row of them under
the eaves, and every drop that falls now falls into its own icicle's pit
and no other.{]}}

\textbf{Achilles:} And her blanks? The press eats one a strike. Couldn't
she have made her own, out of stones she had two of?

\begin{center}\rule{0.5\linewidth}{0.5pt}\end{center}

\hypertarget{fourth-challenge-eves-blanks}{%
\section{\texorpdfstring{Fourth Challenge --- \emph{Eve's
Blanks}}{Fourth Challenge --- Eve's Blanks}}\label{fourth-challenge-eves-blanks}}

\textbf{Crab:} She tried. \emph{{[}turning the sheet over{]}} ``Blanks
cost money, and hers ate one a strike. Now and then a customer sent two
stones settled alike by the same hand --- two of a kind, sent together
--- and she never knew how they had been settled. So she built a frame
to turn one of the two into a blank, and leave the other as it was. Her
frame had a tell: a small stone inside, which shifted whenever it
worked.''

\textbf{Achilles:} And it worked?

\textbf{Crab:} ``It always gave her a blank.'' \emph{{[}reading on{]}}
``When she was found out, we bought the frame from her and tried it at
the bench: two stones settled alike, the frame run, and the tell read at
the setting the stones had been given. The tell always came out holding
the stone she had blanked. It read exactly as that stone would have
read, every settling, every trial.''

\textbf{Achilles:} Then she did delete it; the tell is just a bit of
ballast that shook when it ran. A blank is a blank.

\textbf{Tortoise:} Watch where the stone goes. \emph{{[}She lays two
strings side by side on the cloth.{]}} The upper line is her second
stone. The lower is the tell, which starts settled low, as a blank does.
And between them, a knot you know from Mr Crab's beam: three cords from
line to line, the middle one tied the other way about.

\emph{{[}She ties it: three short cords across from the upper line to
the lower, the first and last tied one way, the middle the other. The
upper line runs in from the second stone at the left; the lower runs in
from the tell's settled start. She takes both right-hand ends in one
claw and both left-hand ends in the other and pulls. The three cords
draw tight, slide into one another, and give way; the knot comes undone.
What is left is two lines and a single crossing. The line that came in
from the stone now runs down across the middle and out along the bottom,
into the tell's place; the line that came in from the tell's settled
start runs up and out along the top, where the second stone's line used
to leave. Nothing is on either line that was not there before, and
nothing is missing; only which one ends where has changed.{]}}

\textbf{Achilles:} A crossing. The stone went down into the tell, and
the tell's own start came up to take its place --- that's her blank. Two
stones in, one blank out, and nothing gone: that's all her un-copying
was.

\emph{{[}Across the lane a man in a greatcoat is clearing the path to
his door. It comes clear only a shovelful at a time, and every shovelful
goes onto the bank beside it, which grows by just what the path has
lost. When he straightens at his door there is a clean path and a taller
bank, and not a flake fewer in his garden.{]}}

\textbf{Crab:} ``Nothing came free,'' he says. He has underlined it.

\textbf{Tortoise:} She could move the stone. She could not make it be
nowhere.

\emph{{[}The icicles have gone thin and glassy with the afternoon, each
tip a bead of light, and the drops come slower now, as though each were
being considered before it fell. The sun is lower and reaches further
in: its patch has crept along the floorboards beside the table, halfway
down the room towards the end wall. Achilles draws his feet back from
the fire and finds his stockings dry.{]}}

\textbf{Crab:} \emph{{[}folding the letter, and looking down the table
at the press, and at the fresh stones, one at each end{]}} Then let us
neither copy nor delete. Let us \emph{use} it. The post to Delft takes a
week in this weather. Say that end of the table is Delft. If the far
stone there could be copied, could it not carry word faster than the
post?

\begin{center}\rule{0.5\linewidth}{0.5pt}\end{center}

\hypertarget{fifth-challenge-the-forbidden-telegraph}{%
\section{\texorpdfstring{Fifth Challenge --- \emph{The Forbidden
Telegraph}}{Fifth Challenge --- The Forbidden Telegraph}}\label{fifth-challenge-the-forbidden-telegraph}}

\textbf{Tortoise:} Suppose the press were perfect, and made true copies
of anything. Achilles, you are here at the window with the near stone,
and you choose how to read it. What does Delft see?

\textbf{Achilles:} Read the near stone plain, and the far one settles
plain, high or low; perfect copies of that are two alike, and read plain
they always agree. Read the near stone crosswise, and the far one
settles crosswise; perfect copies of a crosswise stone, read plain, are
two separate coin tosses, and they agree only half the time. So Delft
presses, reads both copies, and knows: all agreeing, I read plain; half
agreeing, I read crosswise. A telegraph!

\textbf{Tortoise:} But this press is not perfect. Run it, Mr Crab.

\textbf{Crab:} \emph{{[}a long breath{]}} Once, then. The only strike of
the afternoon. Achilles, read yours crosswise.

\emph{{[}Achilles gets up from the fender, sets his cup down on the
hearth, and wipes his palms on his knees. The Tortoise pushes the quilt
back off one shoulder to have a claw free, moves her chair an inch
nearer the table, and then sits very still.{]}}

\emph{{[}Achilles takes the viewer by its collar. It stands upright, set
for the plain question; he turns it a quarter round until it lies
crosswise, and the needle, resting at the middle of its dial, does not
stir. Only the question has changed. He carries the viewer along the
table's edge to the window end, never over the black, and looks through
it at the near stone.{]}}

\emph{{[}The needle leaves the middle of the dial and swings hard over
to the high mark, and holds there; it might as well have gone the other
way. The near stone, which a moment ago could have answered either way,
lies still now, spent, its reflection still beneath it. Achilles takes
one counter off the stack in the tray and lays it on the table's edge
opposite the near stone, face-up. The other two stay stacked in the
tray, meaning nothing yet.{]}}

\textbf{Achilles:} It swung to the high mark, at cross. Face-up. And
I've told nobody.

\emph{{[}At the far end the Crab lifts the hopper's lid, takes up the
far stone, and sets it in the hopper's cup. As the stone leaves the
black its reflection sinks from the surface and goes down, into the
hopper of the press's reflection. The Crab lets the lid fall. The far
stone is inside now; the place where it stood is bare black.{]}}

\textbf{Crab:} Shutter shut. \emph{{[}He tries it with a claw: shut.{]}}
And eyes on the fire.

\emph{{[}He turns his head to the grate and keeps it there, and feels
for the lever, and pulls it once, blind. A thump in the housing. From
the left chute a stone drops onto the black, and the instant it lands
its reflection is there beneath it. From the right chute drops another,
and its reflection meets it. The hopper is empty. Where one stone went
in, two stand on the black below the chutes, each doubled in the glass,
and nothing about either says which, if either, was the stone he put
in.{]}}

\emph{{[}The patch of sun, which has been creeping down the room all
afternoon, reaches the end wall just then and lies full across the glass
of the print, so that the frame seems for a breath to stand nearer the
table than the wall.{]}}

\emph{{[}The Crab turns back from the fire, and his eye goes past the
two stones to the print on the wall --- Escher's
\textup{Drawing Hands}.{]}}

\textbf{Crab:} Two hands, each drawing the other. And neither of them is
the one I put in.

\textbf{Tortoise:} \emph{{[}taking the viewer{]}} Now Delft reads.
\emph{{[}She turns the collar back until it stands upright, and carries
the viewer along the edge of the table, the long way round, to the far
end.{]}}

\emph{{[}She looks through it at the first struck stone. The needle
swings to the low mark and holds. The stone is spent, and still. She
takes a counter from the tray and lays it on the edge opposite,
face-down. She moves the viewer to the second struck stone and looks.
The needle goes to the low mark again, and holds. She lays the last
counter beside the other, face-down. The tray is empty. At the window
end one counter lies face-up; at the far end two lie face-down, side by
side.{]}}

\textbf{Tortoise:} Low on the first; low on the second as well.

\textbf{Achilles:} They agree. But that's a single run --- agreeing is
what a perfect press would have shown me if I'd read plain.

\textbf{Crab:} \emph{{[}opening the proving-book at its last pages{]}}
Then he has tried it for us, many times. ``Pressings of one half of a
fresh pair, the other half read first, at either setting.''
\emph{{[}reading down{]}} The struck pair always agreed, high with high
or low with low, half the time each. When the near half was read plain
--- the same. When it was read crosswise --- the same. Page after page,
you could not tell one setting's page from the other's.

\textbf{Crab:} \emph{{[}a finger at a margin{]}} And where the near half
was read plain, the struck pair matched the near reading too. But that
match lives in the near end's column, and Delft only ever gets that
column by post.

\textbf{Achilles:} So Delft sees the same whatever I do here.

\textbf{Tortoise:} Watch the counters' cords. \emph{{[}She lays it out
on the cloth: the bend of the pair, the fork on its far leg, and a short
cord for each counter.{]}} A plain reading I draw as a green ring
gathering a stone and its reflection; a crosswise one, as a red ring.

\emph{{[}First the plain setting. The near reading's ring is green, and
so is everything at the far end: the fork, its mirror, the two far
readings. Green meets green along the cords, and they slide together
into a single green ring, with all three counters' cords hanging from
it. Pull any one of the three, and the other two come with it.{]}}

\emph{{[}Then she takes off the near green ring and puts a red one in
its place, for the crosswise reading. The red ring and the far green
ring are joined by two cords, the stone's and its reflection's. She
draws the red ring away from the green. The two cords between them slip
loose together and fall away, as a doubled link between the two colours
always does; the red ring lies apart with the near counter's cord
hanging from it, tied to nothing at the far end. The far green ring
still holds both far cords together. What changed: the near cord is free
of the far ones. What held: the two far cords still hang as one.{]}}

\textbf{Achilles:} At cross, my counter isn't tied to theirs at all. And
theirs are still tied to each other, so they agree --- but about nothing
I did.

\emph{{[}Through the window at Achilles's end, the two chimneys of the
house across the lane are smoking side by side. Their smoke leans the
same way, bends at the same height, and frays out together over the
stable roof; it is the one small wind on both.{]}}

\textbf{Tortoise:} No cleverness would mend it: to copy every stone, a
press would first have to cut every bend in two, and a bend in two
pieces is no pair at all.

\textbf{Achilles:} It copies its own kind and makes a pair of every
other; and a press that copied every kind would beat the post to Delft.
So there's no such press.

\emph{{[}Outside, the drops come slower still. The sun has slid along
the eaves towards the corner of the house, and each icicle holds its
bead at the tip a long while, bright, before it lets it go.{]}}

\textbf{Crab:} \emph{{[}quietly, closing the proving-book{]}} The post
to Delft, then. I shall write tonight, and thank him, and tell him I
worked it blind.

\textbf{Tortoise:} And tell him you never once looked. He'll know what
that cost you.

\emph{{[}The Crab slides the press's shutter home a last time, though it
is shut already, and lays his claw flat on the hood. On the black lie
the spent stones, each with its reflection; on the table's edge, the
laid counters; the fresh pair is gone.{]}}

\emph{{[}The patch of sun climbs off the end wall and goes out of the
room. Along the eaves the last drops let go, one icicle after another,
each in its own time, and nothing comes after them. The Crab keeps his
claw where it is a moment longer, feeling the iron warm under it.{]}}

\textbf{Crab:} Listen --- the eaves have stopped. The sun's gone off
them.

\textbf{Achilles:} So it has. It's colder already by the window.

\textbf{Crab:} Then come away from it, both of you, and I'll put the
kettle back on. There's the end of the seed cake.

\emph{{[}The Tortoise gets up, stiff from sitting, and comes round the
table's end to the fire with her quilt gathered about her. The Crab cuts
the end of the seed cake into three pieces, none of them the size of
another, and gives Achilles the largest.{]}}

\emph{{[}The Crab pours. Achilles takes his cup to the fire and stands
with it, looking back at the counters on the edge of the table.{]}}

\textbf{Achilles:} Not this press, and not a cleverer one. But anything
at all --- anywhere --- that could take a stone nobody knows and make
two of it?

\textbf{Tortoise:} Ask the frost tonight to draw this morning's pane
again, line for line, and see what it gives you.

\emph{{[}Nobody takes her up on it; there is cake to be eaten first. The
sun is off the house now, and the parlour's light has gone even and
grey-white, the light of a winter afternoon that has begun, without any
fuss, to think about the evening. Along the eaves every icicle holds one
last bead at its tip, and the beads have stopped falling: they are
stiffening where they hang, and the little pits in the snow on the sill
are glazing over. At the very bottom corner of the lowest pane, where
the glass is coldest, a first faint grey of frost has started, too fine
yet to have any pattern at all. Down the table the stones lie where the
afternoon left them: one at the window end, with the bright counter on
the edge opposite; two at the far end below the chutes, with the dulled
counter and the nicked one side by side on the edge before them; the
tray between is empty. The press stands at the far end keeping its works
to itself, and at the seam of its hood, just where the eye cannot quite
get at it, the faint glow is back; nobody turns to look. In the dresser
drawer this morning's pair keeps its own counsel. On the end wall the
print has gone grey with the rest of the wall now the sun is off it ---
two hands, each drawing the other, in a plain frame --- and nobody has
glanced at it since. The three of them sit close to the fire with their
odd cups: Achilles in his stockinged feet, his boots dried stiff on the
hearth beside him; the Crab with crumbs of seed cake down his front; the
Tortoise with the quilt up to her chin. The kettle, pulled to the side
of the hob, sings itself slowly down to nothing.{]}}

\hypertarget{chapter-11-copying-is-forbidden}{%
\chapter{Chapter 11: Copying Is
Forbidden}\label{chapter-11-copying-is-forbidden}}

\emph{--- Gigue That May Not Be Repeated ---}

\emph{Part III: Pictures Become Proofs}

\emph{Escher: ``Drawing Hands'' --- two hands, each drawing the other}

\begin{center}\rule{0.5\linewidth}{0.5pt}\end{center}

Can a copy be made of something that nobody has looked at?

Every copier we know begins with a look. A photocopier shines a light on
the page and records what comes back. A scribe reads a line before
writing it out again. A computer copies a bit by finding out whether it
is \(0\) or \(1\) and writing the same thing somewhere else. Copying is
so cheap in the everyday world that we hardly notice that it always
starts by looking. Chapter 9 told us what a look does to a qubit: it
forces one answer out of several, at random, and destroys the rest. So a
copier of qubits, if there is one, would have to work blind. It would
take a state it has not examined, and cannot examine without spoiling,
and hand back two of it.

If such a machine existed, a great deal would follow. You could make a
thousand copies of an unknown state and ask each copy a different
question, until the state had no secrets left. And you could, as the end
of this chapter will show, send a message faster than the post. Chapter
10 ended on exactly this question and left it open: ``Could anything,
anywhere, take an unknown qubit and make two? The next chapter asks, and
the answer is written into the very linearity that made every picture in
this one work.'' In teleportation the message was gone from Alice's end
the moment it was read, and at no moment were there two of it. Was that
a quirk of one clever protocol, or a law?

It is a law, and one of the shortest in the book. On the blackboard it
is a single line, in which a number would have to equal its own square.
In the sketchbook it is a count of closed pictures: one on one side, two
on the other. But the law has a positive side, and the pictures see it
first. Some things \emph{can} be copied: the things that were settled
already. The machine the Crab unpacked this afternoon copies exactly
those, and turns everything else into a pair. From there the chapter
follows a copy thrown away (it turns out to be a look that nobody read)
and a copy that someone tries to unmake (it can only be moved). At the
end we meet the only telegraph a copier could ever build, one that no
linear picture can draw.

\begin{center}\rule{0.5\linewidth}{0.5pt}\end{center}

\hypertarget{the-press-copies-its-own}{%
\section{§11.1 The Press Copies Its
Own}\label{the-press-copies-its-own}}

The press in the dialogue is a copier with a magazine of blanks. Settle
a stone plain-high, press it, and two high stones drop out. Settle it
plain-low, and two low ones. The proving-book says so, line after line.
What does the same press do to a stone settled crosswise?

\textbf{On the blackboard.} Chapter 8 met the copier (§8.3, row 59), and
we now give it a letter of its own. The \textbf{copy map} is

\[\delta \;=\; \ket{00}\bra0+\ket{11}\bra1 ,\]

the phase-free green dot with one leg in and two out, \(Z(1,2,0)\) of
§8.1; the Dirac's Blackboard of §8.2 names it as a \(4\times2\) matrix.
On the two plain states it does what its name says, exactly:

\[\delta\ket0=\ket{00}=\ket0\otimes\ket0,\qquad \delta\ket1=\ket{11}=\ket1\otimes\ket1 .\]

Now the crosswise stone, \(\ket+=\tfrac1{\sqrt2}(\ket0+\ket1)\). The
copier is linear, so it acts on each branch separately:

\[\delta\ket+ \;=\; \tfrac1{\sqrt2}\bigl(\ket{00}+\ket{11}\bigr) \;=\; \ket{\Phi^+}.\]

Two copies of \(\ket+\) would be
\(\ket+\otimes\ket+=\tfrac12\bigl(\ket{00}+\ket{01}+\ket{10}+\ket{11}\bigr)\).
The copier has not made two crosswise stones. It has made the Bell pair
of Chapter 6: an entangled state, neither half of which is \(\ket+\),
and each half of which, asked anything on its own, answers like a fair
coin (§6.2). In the same way
\(\delta\ket-=\ket{\Phi^-}=\tfrac1{\sqrt2}(\ket{00}-\ket{11})\). For a
general state \(\ket\psi=\alpha\ket0+\beta\ket1\) the two things to
compare are

\[\delta\ket\psi=\alpha\ket{00}+\beta\ket{11},\qquad
\ket\psi\otimes\ket\psi=\alpha^2\ket{00}+\alpha\beta\ket{01}+\alpha\beta\ket{10}+\beta^2\ket{11},\]

and they agree only when \(\alpha\beta=0\).

The press has a magazine of blanks, and \(\delta\) seems to have none.
Where did the blank go? Put a blank \(\ket0\) on a second wire and run a
CNOT from the stone to the blank. Then
\(\mathrm{CNOT}\,(\ket\psi\otimes\ket0)=\alpha\ket{00}+\beta\ket{11}=\delta\ket\psi\)
for every \(\ket\psi\), so, exactly,

\[\mathrm{CNOT}\,(I\otimes\ket0)\;=\;\delta .\]

The copier is a CNOT with its blank already loaded.

Is the press simply badly made? Could a cleverer one, with works inside
that are allowed to change as it runs, copy the plain stones \emph{and}
everything else? To ask that properly we need to say what copying is.

We say a process \textbf{clones} a state \(\ket\psi\) if, started on
\(\ket\psi\), a blank \(\ket0\), and its own apparatus in a fixed
starting state \(\ket A\), it ends in

\[\ket\psi\otimes\ket\psi\otimes\ket{A_{\psi}},\]

for some state \(\ket{A_{\psi}}\) of the apparatus, which is allowed to
depend on \(\ket\psi\). The apparatus may be of any size.

\textbf{\emph{Why this definition.}} Why insist that the two copies come
out as a product, \(\ket\psi\otimes\ket\psi\), while the apparatus is
allowed to do anything at all? Because a copy is something you can walk
away with. Each copy, taken on its own, must be \(\ket\psi\), and the
copies must not be tied to each other or to the machine; otherwise they
are parts of one larger thing, as \(\delta\ket+\) is. The apparatus is
left free so that the definition rules out no cleverness inside the
machine. (§11.3 relaxes the demand on the copies too, and finds that the
relaxed demand fails as well.)

Here is the first answer. It needs nothing but linearity.

\begin{quote}
\textbf{Result 11.1 (no-cloning from linearity).} Let \(L\) be any
linear process on a stone, a blank that starts in \(\ket0\), and an
apparatus of any size that starts in a fixed state \(\ket A\). Suppose
\(L\) clones the two plain states:
\[L\bigl(\ket0\ket0\ket A\bigr)=\ket0\ket0\ket{A_0},\qquad L\bigl(\ket1\ket0\ket A\bigr)=\ket1\ket1\ket{A_1}.\]
Then for every \(\ket\psi=\alpha\ket0+\beta\ket1\),
\[L\bigl(\ket\psi\ket0\ket A\bigr)\;=\;\alpha\ket{00}\ket{A_0}+\beta\ket{11}\ket{A_1},\]
which has \textbf{no \(\ket{01}\) part}, while a clone
\(\ket\psi\ket\psi\ket{A_{\psi}}\) has \(\ket{01}\) part
\(\alpha\beta\ket{A_{\psi}}\). So \(L\) clones \(\ket\psi\) only if
\(\alpha\beta=0\). With no apparatus, \(L\) with its blank loaded is
exactly \(\delta\).
\end{quote}

By the \textbf{\(\ket{01}\) part} of a state on stone, copy and
apparatus we mean what is left when the effect \(\bra{01}\) is applied
to the first two wires, \((\bra{01}\otimes I)\), and the apparatus is
left open.

\emph{Proof idea.} A linear machine that copies the two plain states has
no choice about anything else: it treats a superposition branch by
branch, so every branch it makes has stone and copy agreeing. A true
copy of a superposition needs branches in which they disagree.

\emph{Proof.} Linearity gives
\(L(\ket\psi\ket0\ket A)=\alpha L(\ket0\ket0\ket A)+\beta L(\ket1\ket0\ket A)\),
which is the displayed state. Apply \(\bra{01}\otimes I\): both terms
die, since \(\braket{01|00}=\braket{01|11}=0\). For the clone,
\((\bra{01}\otimes I)\bigl(\ket\psi\ket\psi\ket{A_{\psi}}\bigr)=\braket{0|\psi}\braket{1|\psi}\ket{A_{\psi}}=\alpha\beta\ket{A_{\psi}}\),
and \(\ket{A_{\psi}}\) is a unit vector, so this is zero only when
\(\alpha\beta=0\). Without an apparatus, a linear map from the stone to
two wires is fixed by where it sends \(\ket0\) and \(\ket1\), and
\(\ket{00}\), \(\ket{11}\) are where \(\delta\) sends them.
\(\;\blacksquare\) See \texttt{Code/Ch11/OwnBasis.py}, which also runs
the three-wire basis cloner
\(\mathrm{CNOT}_{0\to2}\,\mathrm{CNOT}_{0\to1}\) and finds its
\(\ket{01}\) part empty on a non-real \(\ket\psi\).

\vskip 0pt plus 4\baselineskip\penalty-100\vskip 0pt plus
-4\baselineskip

\begin{quote}
\textbf{Dirac's Blackboard.}\nopagebreak

\emph{The missing \(\ket{01}\), and what the apparatus knows.} Take
\(\ket\psi=\ket+\), so \(\alpha=\beta=\tfrac1{\sqrt2}\). A clone would
be
\(\ket{++}\ket{A_+}=\tfrac12\bigl(\ket{00}+\ket{01}+\ket{10}+\ket{11}\bigr)\ket{A_+}\),
with amplitude \(\tfrac12\) on ``stone low, copy high''. The linear
machine gives
\(\tfrac1{\sqrt2}\bigl(\ket{00}\ket{A_0}+\ket{11}\ket{A_1}\bigr)\):
amplitude \(0\) there, whatever the apparatus does. The apparatus cannot
repair the gap, because it only rides along on each branch; it cannot
create a branch. It can do something else, though. If \(\ket{A_0}\) and
\(\ket{A_1}\) differ, the apparatus holds a record of which branch it is
on. When they are orthogonal, the record is perfect: the machine has in
effect \emph{looked} at the stone, in the plain way, and written down
the answer. That remark will come back in §11.3, where copying and
looking turn out to be the same act.
\end{quote}

\textbf{In the sketchbook.} Chapter 8's points carried a \(\sqrt2\). A
red dot with one leg and phase \(0\) is \(\sqrt2\,\ket0\), and with
phase \(\pi\) it is \(\sqrt2\,\ket1\) (§8.2, §8.3). A green dot with one
leg and phase \(0\) is \(\ket0+\ket1=\sqrt2\,\ket+\), and with phase
\(\pi\) it is \(\sqrt2\,\ket-\). With Chapter 10's bit-labelled phase
(row 75) one drawing covers both cases of each colour: for a bit \(a\),
the red point \(a\pi\) is \(\sqrt2\,\ket a\), and the green point
\(a\pi\) is \(\sqrt2\) times \(\ket+\) or \(\ket-\). So a normalized
state is a point with \(\tfrac1{\sqrt2}\) written in front, and we write
it in every time.

Feed a red point to the copier. The copy rule (Theorem 8.3) sends it out
as two red points, with a factor \(\tfrac1{\sqrt2}\):

\[\tfrac1{\sqrt2}\;
\begin{tikzpicture}[sd, x=1.4em, y=1.4em,
  zsp/.style={circle,draw,line width=0.9pt,fill=green!55,minimum size=1.1em,inner sep=0pt},
  xsp/.style={circle,draw,line width=0.9pt,fill=red!75,minimum size=1.1em,inner sep=0pt},
  ph/.style={font=\scriptsize,inner sep=1pt}]
  \draw[wire] (0,0) -- (1.8,0);
  \draw[wire] (1.8,0) -- (2.8,0.7) -- (3.8,0.7);
  \draw[wire] (1.8,0) -- (2.8,-0.7) -- (3.8,-0.7);
  \node[xsp,label={[ph]above:{$a\pi$}}] at (0,0) {};
  \node[zsp] at (1.8,0) {};
\end{tikzpicture}
\;\;\overset{\text{copy}}{=}\;\;
\tfrac12\;
\begin{tikzpicture}[sd, x=1.4em, y=1.4em,
  xsp/.style={circle,draw,line width=0.9pt,fill=red!75,minimum size=1.1em,inner sep=0pt},
  ph/.style={font=\scriptsize,inner sep=1pt}]
  \draw[wire] (0,0.7) -- (1.8,0.7);
  \draw[wire] (0,-0.7) -- (1.8,-0.7);
  \node[xsp,label={[ph]above:{$a\pi$}}] at (0,0.7) {};
  \node[xsp,label={[ph]below:{$a\pi$}}] at (0,-0.7) {};
\end{tikzpicture}\]

The numbers in front work out exactly: \(\tfrac1{\sqrt2}\) for the point
that went in, times the rule's \(\tfrac1{\sqrt2}\), is
\(\tfrac12=\tfrac1{\sqrt2}\cdot\tfrac1{\sqrt2}\), one normalizing factor
for each point that comes out. So \(\delta\ket a=\ket a\otimes\ket a\),
exactly. Chapter 8 stated the copy rule ``up to \(\sqrt2\)'' (row 59)
only because its points were unnormalized. With normalized states the
rule is exact.

Now feed it a green point. There is nothing to copy here. A green point
and a green dot joined by a wire are two dots of one colour sharing a
leg, and fusion (Theorem 8.1) makes them one:

\[\tfrac1{\sqrt2}\;
\begin{tikzpicture}[sd, baseline={([yshift=-0.25em]0,0)}, x=1.4em, y=1.4em,
  zsp/.style={circle,draw,line width=0.9pt,fill=green!55,minimum size=1.1em,inner sep=0pt},
  ph/.style={font=\scriptsize,inner sep=1pt}]
  \draw[wire] (0,0) -- (1.8,0);
  \draw[wire] (1.8,0) -- (2.8,0.7) -- (3.8,0.7);
  \draw[wire] (1.8,0) -- (2.8,-0.7) -- (3.8,-0.7);
  \node[zsp,label={[ph]above:{$a\pi$}}] at (0,0) {};
  \node[zsp] at (1.8,0) {};
\end{tikzpicture}
\;\;\overset{\text{fusion}}{=}\;\;
\tfrac1{\sqrt2}\;
\begin{tikzpicture}[sd, baseline={([yshift=-0.25em]0,0)}, x=1.4em, y=1.4em,
  zsp/.style={circle,draw,line width=0.9pt,fill=green!55,minimum size=1.1em,inner sep=0pt},
  ph/.style={font=\scriptsize,inner sep=1pt}]
  \draw[wire] (0,0) -- (1,0.7) -- (2.2,0.7);
  \draw[wire] (0,0) -- (1,-0.7) -- (2.2,-0.7);
  \node[zsp,label={[ph]left:{$a\pi$}}] at (0,0) {};
\end{tikzpicture}
\;\;=\;\;
\tfrac1{\sqrt2}\;
\begin{tikzpicture}[sd, baseline={([yshift=-0.25em]0,0)}, x=1.4em, y=1.4em,
  zsp/.style={circle,draw,line width=0.9pt,fill=green!55,minimum size=1.1em,inner sep=0pt},
  ph/.style={font=\scriptsize,inner sep=1pt}]
  \sdcup{0}{0.7}{-0.7}
  \draw[wire] (0,0.7) -- (2.6,0.7);
  \draw[wire] (0,-0.7) -- (2.6,-0.7);
  \node[zsp,label={[ph]above:{$a\pi$}}] at (1.3,0.7) {};
\end{tikzpicture}\]

The middle picture is a green dot with no legs in and two out, phase
\(a\pi\): by §8.1's formula,
\(\ket{00}+(-1)^a\ket{11}=(Z^a\otimes I)\,\)cup. Read fusion backwards
to pull the phase out onto the top leg as a one-in, one-out dot; what is
left, the phase-free green dot with two legs out, is
\(\ket{00}+\ket{11}\), the cup itself. That is the right-hand picture.
The cup is unnormalized, as always, cup \(=\ket{00}+\ket{11}\), so
\(\ket{\Phi^+}=\tfrac1{\sqrt2}\,\)cup, and the equation says
\(\delta\ket\pm=\ket{\Phi^{\pm}}\) exactly, with the drawn
\(\tfrac1{\sqrt2}\) doing real work. The point does not come out
doubled. It melts into the copier, and what leaves is one bent wire: a
pair, not two copies.

The blank can be drawn going in too. A CNOT is a green dot tied by a
cord to a red dot, and the bare picture is \(\mathrm{CNOT}/\sqrt2\) (row
64); the blank \(\ket0\) is \(\tfrac1{\sqrt2}\) times a red point. The
two numbers cancel:

\[\mathrm{CNOT}\,(I\otimes\ket0)\;=\;\sqrt2\cdot\tfrac1{\sqrt2}\;
\begin{tikzpicture}[sd, x=1.4em, y=1.4em,
  zsp/.style={circle,draw,line width=0.9pt,fill=green!55,minimum size=1.1em,inner sep=0pt},
  xsp/.style={circle,draw,line width=0.9pt,fill=red!75,minimum size=1.1em,inner sep=0pt}]
  \draw[wire] (-0.4,1.2) -- (3,1.2);
  \draw[wire] (0.2,0) -- (3,0);
  \draw[wire] (1.6,1.2) -- (1.6,0);
  \node[zsp] at (1.6,1.2) {};
  \node[xsp] at (1.6,0) {};
  \node[xsp] at (0.2,0) {};
\end{tikzpicture}
\;\overset{\text{fusion}}{=}\;
\begin{tikzpicture}[sd, x=1.4em, y=1.4em,
  zsp/.style={circle,draw,line width=0.9pt,fill=green!55,minimum size=1.1em,inner sep=0pt},
  xsp/.style={circle,draw,line width=0.9pt,fill=red!75,minimum size=1.1em,inner sep=0pt}]
  \draw[wire] (0,1.2) -- (3,1.2);
  \draw[wire] (1.6,1.2) -- (1.6,0) -- (3,0);
  \node[zsp] at (1.6,1.2) {};
  \node[xsp] at (1.6,0) {};
\end{tikzpicture}
\;\overset{\text{id}}{=}\;
\begin{tikzpicture}[sd, x=1.4em, y=1.4em,
  zsp/.style={circle,draw,line width=0.9pt,fill=green!55,minimum size=1.1em,inner sep=0pt}]
  \draw[wire] (0,1.2) -- (3,1.2);
  \draw[wire] (1.6,1.2) to[out=-90,in=180] (3,0);
  \node[zsp] at (1.6,1.2) {};
\end{tikzpicture}
\;=\;\delta .\]

The blank's red point fuses into the red target of the CNOT (same
colour, shared leg). What remains of the red dot has one leg in and one
out and no phase, so it is a bare wire (§8.1). What is left is the green
copier. The blank has no line of its own: it is part of the copier.

\begin{quote}
\textbf{Result 11.2 (the copier copies its own basis).} Exactly, for
each bit \(a\): \[\delta\ket a=\ket a\otimes\ket a,\]
\[\delta\ket\pm=\ket{\Phi^{\pm}}=\tfrac1{\sqrt2}\,(Z^a\otimes I)\,\text{cup},\]
with \(\ket+\) for \(a=0\) and \(\ket-\) for \(a=1\). In pictures, red
points go through the green copier as two red points, and green points
fuse with it into a single cup. For \(\ket\psi=\alpha\ket0+\beta\ket1\),
the output \(\delta\ket\psi\) is a product state only if
\(\alpha\beta=0\), that is, only if \(\ket\psi\) is a plain state up to
a phase.
\end{quote}

\emph{Proof.} The two pictures above, each step exact (copy with the
scalar kept, and fusion). For the last sentence, write
\(\delta\ket\psi=\alpha\ket{00}+0\ket{01}+0\ket{10}+\beta\ket{11}\) as
its square of coefficients,
\(\bigl(\begin{smallmatrix}\alpha&0\\0&\beta\end{smallmatrix}\bigr)\).
The product test of Chapter 6 (Theorem 6.1) says it is a product exactly
when the determinant, \(\alpha\beta\), vanishes. \(\;\blacksquare\) See
\texttt{Code/Ch11/OwnBasis.py}.

\textbf{\emph{Why this definition.}} Why take \(\delta\), a map with no
apparatus and one fixed basis, as \emph{the} copier, when we are asking
about all copiers? Because Result 11.1 says every linear machine that
copies the plain states, run with no apparatus, \emph{is} \(\delta\),
and Result 11.2 shows what \(\delta\) does with everything else. The
question is not whether \(\delta\) is well made; it is whether any
machine could copy more than one basis. The next section says no.

\vskip 0pt plus 4\baselineskip\penalty-100\vskip 0pt plus
-4\baselineskip

\begin{quote}
\textbf{Coecke's Sketchbook.}\nopagebreak

\emph{Red copies, green fuses.} The green copier copies exactly the red
points, phase \(0\) or \(\pi\), and nothing else. A green point is not
copied. It fuses, and what comes out is a cup. Van de Wetering's survey
says the same of the rule itself: the state-copy rules ``only hold when
the spider being copied has a phase of \(0\) or \(\pi\)'' (§4.2). Give
the red point any other turn and the blackboard agrees:
\(\tfrac1{\sqrt2}X(0,1,\theta)\) has
\(\alpha\beta=\tfrac14(1-e^{2i\theta})\), which vanishes only for
\(\theta\in\{0,\pi\}\). Coecke and Duncan read no-cloning this way
round, as a positive statement: quantum states may be copied ``if they
are known to lie in a given basis'' (their §1). A copier \emph{is} a
choice of basis, and the states it copies are the ones already settled
in it. Each colour copies the other's points and swallows its own. That
is complementarity, and Chapter 8's Hopf law (row 62) is its other face:
van de Wetering (§4.5) reads Hopf as saying that the green and red bases
are complementary.
\end{quote}

\begin{center}\rule{0.5\linewidth}{0.5pt}\end{center}

\hypertarget{one-loop-two-loops}{%
\section{§11.2 One Loop, Two Loops}\label{one-loop-two-loops}}

The Crab's next hope is the natural one. Perhaps a cleverer press, with
more elaborate works, could copy more. It need not copy everything.
Suppose it only has to copy two stones: one settled plain-low,
\(\ket0\), and one settled crosswise, \(\ket+\).

\textbf{On the blackboard.} Linearity alone does not forbid this. The
two states \(\ket0\) and \(\ket+\) are linearly independent, so there is
a linear map \(L\) from one wire to two with

\[L\ket0=\ket{00},\qquad L\ket+=\ket{++}.\]

It is fixed on the basis \(\ket0,\ket1\) by
\(\ket1=\sqrt2\,\ket+-\ket0\), giving
\(L\ket1=\sqrt2\,\ket{++}-\ket{00}\). Result 11.1 does not apply,
because \(L\) does not copy \(\ket1\). But look at the length of
\(L\ket1\):

\[\bigl\lVert L\ket1\bigr\rVert^2 \;=\; 2+1-2\sqrt2\,\braket{00|++} \;=\; 3-\sqrt2 \;\approx\; 1.59 .\]

\(L\) stretches a state out of the unit sphere, so it cannot be a
physical process (Theorem 3.2). Here is what went wrong, stated in a
form that will generalize. Before the press, the two stones had
amplitude \(\braket{0|+}=\tfrac1{\sqrt2}\) for each other. After a
faithful copy, the pairs would have
\(\braket{00|++}=\braket{0|+}^2=\tfrac12\). A unitary box preserves
amplitudes between states (Theorem 3.2), and
\(\tfrac1{\sqrt2}\neq\tfrac12\). In probabilities, squaring both, it is
a half against a quarter.

Call the amplitude \(\braket{\varphi|\psi}\) the \textbf{likeness} of
the two states: it is \(1\) when they are the same state, \(0\) when
they are orthogonal, and in between otherwise. Copying multiplies a
likeness by itself, while every unitary process keeps it fixed.

\textbf{In the sketchbook.} Put the state \(\psi\) at the left end of a
wire and the effect \(\varphi\) at the right end: the closed diagram of
Chapter 2, whose value is the amplitude (§2.2). The effect labelled
\(\varphi\) is the \emph{mirror image} of the state \(\varphi\), so its
entries are conjugated. We call this closed picture a \textbf{likeness}.
The dialogue's string makes it a loop, and it is one. Put the state
\(\psi\) and the effect \(\varphi\) on Chapter 10's closed loop, as the
box \(\ket\psi\bra\varphi\); the loop is the trace of that box (Result
10.6), which is \(\braket{\varphi|\psi}\), and yanking straightens it
into the segment below. Two likenesses stacked, with nothing joining
them, are the product of their numbers (Proposition 4.1):

\[\braket{\varphi|\psi}\;=\;
\begin{tikzpicture}[sd]
  \draw[wire] (0,0) -- (3.6,0);
  \node[sdstate, text height=1.5ex, text depth=0.5ex] at (0,0) {$\psi$};
  \node[sdeffect, text height=1.5ex, text depth=0.5ex] at (3.6,0) {$\varphi$};
\end{tikzpicture}
\qquad\qquad
\braket{\varphi|\psi}^2\;=\;
\begin{tikzpicture}[sd]
  \draw[wire] (0,1.1) -- (3.6,1.1);
  \draw[wire] (0,-1.1) -- (3.6,-1.1);
  \node[sdstate, text height=1.5ex, text depth=0.5ex] at (0,1.1) {$\psi$};
  \node[sdeffect, text height=1.5ex, text depth=0.5ex] at (3.6,1.1) {$\varphi$};
  \node[sdstate, text height=1.5ex, text depth=0.5ex] at (0,-1.1) {$\psi$};
  \node[sdeffect, text height=1.5ex, text depth=0.5ex] at (3.6,-1.1) {$\varphi$};
\end{tikzpicture}\]

One loop is the likeness of the stones. Two separate loops are the
likeness of the pairs, the same number taken twice over and multiplied.

\textbf{\emph{Why this definition.}} Why measure likeness by the
amplitude \(\braket{\varphi|\psi}\) and not by the probability
\(\lvert\braket{\varphi|\psi}\rvert^2\)? Both would serve for the
argument, but the amplitude is the thing a closed picture \emph{is}. A
process that meets its own mirror image and vanishes leaves the
amplitude untouched, phase and all, so the amplitude is what the
pictures can carry through a proof unchanged.

The worked instance is the copier itself. The mirror image of the green
fork is a green dot with two legs in and one out; it is
\(\delta^{\dagger}\), and since \(\delta\) is real no conjugation shows.
Put \(\delta\) and then \(\delta^{\dagger}\) on a wire. The two dots are
joined by \emph{two} wires, and fusion applies to spiders ``connected by
one or more wires'' (van de Wetering, §4.1):

\[\begin{tikzpicture}[sd, baseline={([yshift=-0.25em]0,0)}, x=1.4em, y=1.4em,
  zsp/.style={circle,draw,line width=0.9pt,fill=green!55,minimum size=1.1em,inner sep=0pt}]
  \draw[wire] (0,0) -- (1.2,0);
  \draw[wire] (1.2,0) to[out=50,in=130] (3.6,0);
  \draw[wire] (1.2,0) to[out=-50,in=-130] (3.6,0);
  \draw[wire] (3.6,0) -- (4.8,0);
  \node[zsp] at (1.2,0) {};
  \node[zsp] at (3.6,0) {};
\end{tikzpicture}
\;\;\overset{\text{fusion}}{=}\;\;
\begin{tikzpicture}[sd, baseline={([yshift=-0.25em]0,0)}, x=1.4em, y=1.4em,
  zsp/.style={circle,draw,line width=0.9pt,fill=green!55,minimum size=1.1em,inner sep=0pt}]
  \draw[wire] (0,0) -- (2.8,0);
  \draw[wire] (1.4,0) .. controls (0.4,1.6) and (2.4,1.6) .. (1.4,0);
  \node[zsp] at (1.4,0) {};
\end{tikzpicture}
\;\;\overset{\text{self-loop}}{=}\;\;
\begin{tikzpicture}[sd, baseline={([yshift=-0.25em]0,0)}, x=1.4em, y=1.4em]
  \draw[wire] (0,0) -- (2.4,0);
\end{tikzpicture}\]

After fusing along one wire, the other wire runs from the dot back to
itself: a \textbf{self-loop}. Identity removal, van de Wetering notes in
the same section, also gets rid of self-loops on spiders: the loop is
erased, and no number appears. (On the blackboard, closing two legs of a
green dot with a wire only forces two of its all-agree labels to agree,
which they already do.) The result is a bare wire, exactly. On the
blackboard, in one line:
\(\delta^{\dagger}\delta=\bigl(\sum_i\ket i\bra{ii}\bigr)\bigl(\sum_j\ket{jj}\bra j\bigr)=\sum_i\ket i\bra i=I\).
The copier meets its mirror and vanishes, as a unitary box does in
Chapter 3; so the copier preserves likeness,
\(\bra\varphi\delta^{\dagger}\delta\ket\psi=\braket{\varphi|\psi}\). It
is an \emph{isometry}, not a unitary: in the other order,
\(\delta\delta^{\dagger}\) is not the identity on two wires.

So the question is now exact. A faithful press would have to turn one
loop into two, and a likeness into its own square, while keeping it
fixed. With the apparatus kept, here is the law.

\begin{quote}
\textbf{Result 11.3 (no-cloning, unitarity form).} Let \(U\) be any
unitary on a stone, a blank and an apparatus of any size, with the blank
starting in \(\ket0\) and the apparatus in a fixed state \(\ket A\).
Suppose \(U\) clones two states,
\[U\bigl(\ket\psi\ket0\ket A\bigr)=\ket\psi\ket\psi\ket{A_{\psi}},\qquad U\bigl(\ket\varphi\ket0\ket A\bigr)=\ket\varphi\ket\varphi\ket{A_{\varphi}}.\]
Put \(s=\braket{\varphi|\psi}\) and \(a=\braket{A_{\varphi}|A_{\psi}}\).
Then \[s \;=\; s^2\,a ,\] and since \(\lvert a\rvert\le1\),
\textbf{either \(s=0\) or \(\lvert s\rvert=1\)}: a unitary can clone two
states only if they are orthogonal or equal up to a phase. So no unitary
clones every state; for instance \(\ket0\) and \(\ket+\) have
\(s=\tfrac1{\sqrt2}\).
\end{quote}

\emph{Proof idea.} Count the loops. Before the machine, the likeness of
the two inputs is one loop (the blank's loop and the apparatus's loop
are each worth \(1\)). The machine and its mirror vanish, so the
likeness is unchanged. After a faithful copy the same likeness is two
loops and the apparatus's loop, multiplied. One loop can equal two loops
only if the loop is worth nothing or has size one.

The proof is given twice, in full, in §11.6. See
\texttt{Code/Ch11/NoCloning.py}, which runs the unitarity identity for a
random \(8\times8\) unitary and non-real states, and builds the one-loop
and two-loop pictures, with values exactly \(s\) and \(s^2\).

\vskip 0pt plus 4\baselineskip\penalty-100\vskip 0pt plus
-4\baselineskip

\begin{quote}
\textbf{Dirac's Blackboard.}\nopagebreak

\emph{Why the apparatus cannot help.} The apparatus enters the equation
once, as the factor \(a=\braket{A_{\varphi}|A_{\psi}}\) on the right of
\(s=s^2a\). Both \(\ket{A_{\varphi}}\) and \(\ket{A_{\psi}}\) are unit
vectors, so by the Cauchy--Schwarz inequality \(\lvert a\rvert\le1\):
the apparatus can only shrink the right-hand side, never enlarge it.
Take moduli:
\(\lvert s\rvert=\lvert s\rvert^2\lvert a\rvert\le\lvert s\rvert^2\). If
\(s\neq0\), divide by \(\lvert s\rvert\) to get \(\lvert s\rvert\ge1\);
and \(\lvert s\rvert\le1\), again by Cauchy--Schwarz, so
\(\lvert s\rvert=1\). Equality in Cauchy--Schwarz happens only when one
vector is a multiple of the other, so
\(\ket\varphi=e^{i\theta}\ket\psi\): a phase on the whole state, which
no question can detect. With no apparatus at all, \(a=1\) and the law
reads \(s=s^2\). An apparatus whose final states are orthogonal,
\(a=0\), is worse off still: then \(s=0\) is forced. Its works have
recorded which stone went in, which is a look.
\end{quote}

\vskip 0pt plus 4\baselineskip\penalty-100\vskip 0pt plus
-4\baselineskip

\begin{quote}
\textbf{Feynman's Notebook.}\nopagebreak

\emph{Three statements of one law.} Both of our sources credit the
no-cloning theorem to two papers of 1982. One is W. K. Wootters and W.
H. Zurek, ``A single quantum cannot be cloned'', \emph{Nature} 299,
802--803. The other is D. Dieks, ``Communication by EPR devices'',
\emph{Physics Letters A} 92, 271--272. We have had neither paper in
hand, so we say nothing about how either one argued. The statements
below come from our sources. In their preprint of November 1995, Howard
Barnum, Carlton M. Caves, Christopher A. Fuchs, Richard Jozsa and
Benjamin Schumacher describe the theorem as the fact that there are ``no
physical means by which an unknown pure quantum state can be reproduced
or copied''. Arun Kumar Pati and Samuel L. Braunstein, writing about
photons, put the unitary form in one clause: two non-orthogonal
photon-states ``cannot be copied perfectly by a unitary process''.
Result 11.3 is that clause with an apparatus kept. The derivation of
\(s=s^2a\) given here is this book's own.
\end{quote}

\begin{center}\rule{0.5\linewidth}{0.5pt}\end{center}

\hypertarget{copying-is-looking}{%
\section{§11.3 Copying Is Looking}\label{copying-is-looking}}

The maker's letter tells of Eve, a clerk at a carrier, who had a press
like the Crab's. She pressed every stone that passed through her hands,
kept one struck stone, and sent the other on. Her customers noticed
nothing when their stones had been settled plain. Stones settled
crosswise, though, arrived ``looked-at'': at the cross setting they read
as a coin toss. What exactly did her tap do?

\textbf{On the blackboard.} Take the two cases in turn. A plain stone
\(\ket0\) goes through the press as \(\ket{00}\); Eve keeps one, and the
customer receives \(\ket0\), untouched. A crosswise stone \(\ket+\) goes
through as \(\ket{\Phi^+}\), and the customer receives half of a Bell
pair. By Result 9.3 that half, on its own, is \(\tfrac12 I\), the flat
grey that reads as a fair coin at every setting. The press did to
\(\ket+\) what a look does.

To see the general case, let the stone carry any density matrix
\(\rho=\sum_{i,j}\rho_{ij}\ket i\bra j\) (Chapter 9). The copier sends
it to
\(\delta\rho\,\delta^{\dagger}=\sum_{i,j}\rho_{ij}\ket{ii}\bra{jj}\).
Discard either copy, which is the partial trace (Result 9.2). The
partial trace keeps only the terms whose discarded labels agree,
\(i=j\):

\[\operatorname{tr}_2\bigl(\delta\rho\,\delta^{\dagger}\bigr)\;=\;\operatorname{tr}_1\bigl(\delta\rho\,\delta^{\dagger}\bigr)\;=\;\sum_i\rho_{ii}\,\ket i\bra i\;=\;\sum_i P_i\,\rho\,P_i ,\]

with \(P_i=\ket i\bra i\). The diagonal survives; the off-diagonal
entries, where the phase lives, are gone. This is the state of §9.2: the
stone after a plain look whose outcome nobody read.

\begin{quote}
\textbf{Result 11.4 (copying is looking).} For every density matrix
\(\rho\), both marginals of \(\delta\rho\,\delta^{\dagger}\) equal
\(\sum_i P_i\rho P_i\), the state after a plain look with the outcome
unread. They equal \(\rho\) exactly when \(\rho\) is diagonal in the
plain basis.
\end{quote}

\emph{Proof.} The calculation above, for each marginal. The last
sentence holds because \(\sum_iP_i\rho P_i\) is \(\rho\) with its
off-diagonal entries set to zero. \(\;\blacksquare\) See
\texttt{Code/Ch11/NoBroadcasting.py}, which runs a non-real,
non-diagonal \(\rho\).

\textbf{In the sketchbook.} Draw the tap on doubled wires, as every real
process must be drawn since Chapter 9. The stone is two lines, its wire
\(e\) and its mirror copy \(\bar e\) (row 66). The copier acts on \(e\),
and its conjugate acts on \(\bar e\); since \(\delta\) is real, its
conjugate is the green fork again. Keep one copy, \(e_1\) with its
mirror \(\bar e_1\). Discard the other: a cap on \(e_2\) and
\(\bar e_2\) (row 67). The cap is a bent wire, so the two copiers are
now joined by a single wire, and fusion (Theorem 8.1) makes them one
dot:

\[\begin{tikzpicture}[sd,
  zsp/.style={circle,draw,line width=0.9pt,fill=green!55,minimum size=1.3em,inner sep=0pt}]
  \draw[wire] (0,1.7) -- (2,1.7);
  \draw[wire] (0,-1.7) -- (2,-1.7);
  \draw[wire] (2,1.7) -- (3,2.5) -- (5.8,2.5);
  \draw[wire] (2,1.7) -- (3,0.8) -- (4.2,0.8);
  \draw[wire] (2,-1.7) -- (3,-0.8) -- (4.2,-0.8);
  \draw[wire] (2,-1.7) -- (3,-2.5) -- (5.8,-2.5);
  \sdcap{4.2}{0.8}{-0.8}
  \node[zsp] at (2,1.7) {};
  \node[zsp] at (2,-1.7) {};
  \node[sdlabel,anchor=east] at (-0.2,1.7) {$e$};
  \node[sdlabel,anchor=east] at (-0.2,-1.7) {$\bar e$};
  \node[sdlabel,anchor=west] at (6.0,2.5) {$e_1$};
  \node[sdlabel,anchor=west] at (6.0,-2.5) {$\bar e_1$};
  \node[sdlabel,anchor=south] at (3.6,0.85) {$e_2$};
  \node[sdlabel,anchor=north] at (3.6,-0.85) {$\bar e_2$};
\end{tikzpicture}
\;\;\overset{\text{fusion}}{=}\;\;
\begin{tikzpicture}[sd,
  zsp/.style={circle,draw,line width=0.9pt,fill=green!55,minimum size=1.3em,inner sep=0pt}]
  \draw[wire] (0,1.2) -- (1.6,0);
  \draw[wire] (0,-1.2) -- (1.6,0);
  \draw[wire] (1.6,0) -- (3.2,1.2);
  \draw[wire] (1.6,0) -- (3.2,-1.2);
  \node[zsp] at (1.6,0) {};
  \node[sdlabel,anchor=east] at (-0.2,1.2) {$e$};
  \node[sdlabel,anchor=east] at (-0.2,-1.2) {$\bar e$};
  \node[sdlabel,anchor=west] at (3.4,1.2) {$e_1$};
  \node[sdlabel,anchor=west] at (3.4,-1.2) {$\bar e_1$};
\end{tikzpicture}\]

The result is a phase-free green dot with four legs, two in and two out,
\(Z(2,2,0)=\sum_i\ket{ii}\bra{ii}\), exactly. On doubled wires it sends
\(\rho\) to \(\sum_i\rho_{ii}\ket i\bra i\): the blackboard's answer.

Now draw a look that leaves the stone in place. Chapter 9's read-spider
takes the doubled stone and emits a classical wire (row 68). Its mirror
image is a green dot that takes a classical wire in and puts a doubled
stone out, \(\sum_i\ket{ii}\bra i\): it sets a stone to whatever the
record says. Read the stone, then set it as read, and throw the record
away between the two. The classical wire joins two green dots, and it is
a wire like any other to the rules; the double stroke is only a reminder
(§9.1). So fusion makes them one:

\[\begin{tikzpicture}[sd,
  zsp/.style={circle,draw,line width=0.9pt,fill=green!55,minimum size=1.3em,inner sep=0pt}]
  \draw[wire] (0,1.2) -- (1.6,0);
  \draw[wire] (0,-1.2) -- (1.6,0);
  \draw[cwire] (1.6,0) -- (4.2,0);
  \draw[wire] (4.2,0) -- (5.8,1.2);
  \draw[wire] (4.2,0) -- (5.8,-1.2);
  \node[zsp] at (1.6,0) {};
  \node[zsp] at (4.2,0) {};
  \node[sdlabel,anchor=east] at (-0.2,1.2) {$e$};
  \node[sdlabel,anchor=east] at (-0.2,-1.2) {$\bar e$};
  \node[sdlabel,anchor=west] at (6.0,1.2) {$e_1$};
  \node[sdlabel,anchor=west] at (6.0,-1.2) {$\bar e_1$};
\end{tikzpicture}
\;\;\overset{\text{fusion}}{=}\;\;
\begin{tikzpicture}[sd,
  zsp/.style={circle,draw,line width=0.9pt,fill=green!55,minimum size=1.3em,inner sep=0pt}]
  \draw[wire] (0,1.2) -- (1.6,0);
  \draw[wire] (0,-1.2) -- (1.6,0);
  \draw[wire] (1.6,0) -- (3.2,1.2);
  \draw[wire] (1.6,0) -- (3.2,-1.2);
  \node[zsp] at (1.6,0) {};
\end{tikzpicture}\]

It is the same four-legged dot. Copying and throwing one copy away is,
as a picture, a look that nobody read. That is Eve's tap. It left plain
stones alone because a plain stone is already settled, and a look in its
own basis finds nothing to spoil. It turned crosswise stones grey
because a look at the plain setting destroys the crosswise phase. (A
small check: cap the output as well, discarding the stone after the
look, and fusion leaves a green dot with the two input legs only,
\(\sum_i\bra{ii}\), which is the cap. The tapped stone, thrown away, is
a stone thrown away, as causality requires; Result 9.4.)

Barnum and his colleagues relaxed the demand on a copier in exactly the
direction that Eve's customers would care about. The copies need not be
independent of each other; each one, examined alone, need only look
right.

We say a process \textbf{broadcasts} a set of states
\(\{\rho_0,\rho_1\}\) on a system \(A\) if, fed \(\rho_s\) together with
a second system \(B\) in a fixed standard state, it produces a joint
state on \(A\) and \(B\) whose \textbf{two marginals are both
\(\rho_s\)}, for \(s=0\) and for \(s=1\). \textbf{Cloning} is the
special case where the joint state is \(\rho_s\otimes\rho_s\).

\textbf{\emph{Why this definition.}} Why relax cloning to broadcasting?
Because a customer who receives one stone can only examine that stone. A
tap that produced correlated copies, each of which looked perfect alone,
would be invisible to every customer. For a \emph{pure} state the
relaxation buys nothing: Barnum et al.~point out that the only way to
broadcast a pure \(\ket\psi\) is to produce \(\ket\psi\otimes\ket\psi\).
For mixed states it could buy a great deal, and that is the question.

\begin{quote}
\textbf{Result 11.5 (no-broadcasting; Barnum, Caves, Fuchs, Jozsa,
Schumacher; the converse stated, not proved).} A pair of states
\(\{\rho_0,\rho_1\}\) can be broadcast \textbf{if and only if}
\(\rho_0\rho_1=\rho_1\rho_0\). The pair can be cloned if and only if
\(\rho_0\) and \(\rho_1\) are identical or orthogonal
(\(\rho_0\rho_1=0\)).
\end{quote}

\emph{Proof of the easy direction (theirs too).} If \(\rho_0\) and
\(\rho_1\) commute, they are diagonal in one common orthonormal basis
\(\ket b\): \(\rho_s=\sum_b\lambda_{sb}\ket b\bra b\). Use the copier of
\emph{that} basis, \(\delta_{\mathrm{B}}=\sum_b\ket{bb}\bra b\), with
the second system's standard state as its blank. It gives
\(\delta_{\mathrm{B}}\rho_s\delta_{\mathrm{B}}^{\dagger}=\sum_b\lambda_{sb}\ket{bb}\bra{bb}\),
and by Result 11.4, in the basis \(\ket b\), each marginal is \(\rho_s\)
with its off-diagonal entries removed. It has none, so each marginal is
\(\rho_s\). This is broadcasting and not cloning. The state
\(\sum_b\lambda_{sb}\ket{bb}\bra{bb}\) equals
\(\rho_s\otimes\rho_s=\sum_{b,c}\lambda_{sb}\lambda_{sc}\ket{bc}\bra{bc}\)
only if every \(\lambda_{sb}\lambda_{sc}\) with \(b\neq c\) vanishes and
every \(\lambda_{sb}^2=\lambda_{sb}\), that is, only if \(\rho_s\) is
pure. The copier makes correlated copies: broadcast, but not cloned.
\emph{The converse}, that broadcasting forces the two states to commute,
is the hard half. Barnum et al.~prove it with the fidelity between
density matrices, and the proof is beyond this chapter; we cite it. See
\texttt{Code/Ch11/NoBroadcasting.py}, which checks the easy half and the
correlated state, and states the theorem without testing it.

Eve wanted to tap stones settled plain and stones settled crosswise, say
\(\ket0\bra0\) and \(\ket+\bra+\). Those two do not commute. So no
press, however cunning its works, could have passed both on unmarked.
Hers marked exactly the ones the theorem says must be marked.

\vskip 0pt plus 4\baselineskip\penalty-100\vskip 0pt plus
-4\baselineskip

\begin{quote}
\textbf{Feynman's Notebook.}\nopagebreak

\emph{Broadcasting.} The theorem comes from ``Noncommuting mixed states
cannot be broadcast'', by Barnum, Caves, Fuchs, Jozsa and Schumacher
(arXiv:quant-ph/9511010, posted in November 1995 and marked as submitted
to \emph{Physical Review Letters}), which gives its address as the
Center for Advanced Studies of the University of New Mexico, in
Albuquerque. They set the problem as a question about secrets. System A
is ``secretly prepared'' in one of two states, a second system waits in
a standard state, and the question is whether any physical process can
leave both systems looking, separately, like the secret one. Their most
general process is the one of Chapter 9: a unitary together with an
auxiliary system that is afterwards ignored. They also say why cloning
had to fail. If nonorthogonal states could be cloned, one could measure
enough copies to tell the states apart as reliably as one liked, which
quantum measurements cannot do. Any broadcasting process, they add, can
be seen as creating distinguishability ``ex nihilo''. Only for commuting
states does it create none.
\end{quote}

\begin{center}\rule{0.5\linewidth}{0.5pt}\end{center}

\hypertarget{no-blank-for-free}{%
\section{§11.4 No Blank for Free}\label{no-blank-for-free}}

The letter goes on. Blanks cost money, and Eve's press ate one on every
strike. So she tried to make her own. Given two stones that the same
unknown hand had settled alike, could she turn one of them into a blank
and leave the other as it was?

Pati and Braunstein asked the question in this form, for two photons of
the same unknown polarization. Their \textbf{deleting machine} acts on
two copies of a state and an apparatus in a ready state \(\ket A\), and
is required to satisfy

\[\ket\psi\ket\psi\ket A\;\longmapsto\;\ket\psi\ket\Sigma\ket{A_{\psi}}\]

for every \(\ket\psi\). Here \(\ket\Sigma\) is a fixed standard state,
the blank, and the apparatus may end in a state \(\ket{A_{\psi}}\) that
depends on \(\ket\psi\). Eve's frame has such an apparatus: the letter
calls it the tell.

\textbf{\emph{Why this definition.}} Why allow the apparatus to depend
on \(\ket\psi\)? For the same reason as in §11.1: the definition should
forbid no cleverness inside the machine. It also leaves a loophole, and
the loophole is the whole story.

\textbf{An example that works too well.} Let the tell be one more qubit
starting in \(\ket0\), and tie the second stone to the tell with Chapter
8's three CNOTs. By Theorem 8.8 they are a swap, exactly, so

\[\ket\psi\ket\psi\ket0\;\longmapsto\;\ket\psi\ket0\ket\psi\]

for every \(\ket\psi\). The second stone has become the blank
\(\ket\Sigma=\ket0\), and the first is untouched. It is a deleting
machine by the letter of the definition, and it deleted nothing. The
tell's final state is \(\ket{A_{\psi}}=\ket\psi\): the stone has been
moved into the machine. In the sketchbook the knot is pulled straight.
The three spider-CNOTs, drawn bare, carry \(\tfrac1{2\sqrt2}\) (Theorem
8.8), so \(2\sqrt2\) stands in front:

\[2\sqrt2\;
\begin{tikzpicture}[sd, y=1.35em,
  zsp/.style={circle,draw,line width=0.9pt,fill=green!55,minimum size=1.1em,inner sep=0pt},
  xsp/.style={circle,draw,line width=0.9pt,fill=red!75,minimum size=1.1em,inner sep=0pt}]
  \draw[wire] (0,2.6) -- (7.2,2.6);
  \draw[wire] (0,1.1) -- (7.2,1.1);
  \draw[wire] (0,-0.4) -- (7.2,-0.4);
  \draw[wire] (2.4,1.1) -- (2.4,-0.4);
  \draw[wire] (3.9,1.1) -- (3.9,-0.4);
  \draw[wire] (5.4,1.1) -- (5.4,-0.4);
  \node[zsp] at (2.4,1.1) {}; \node[xsp] at (2.4,-0.4) {};
  \node[xsp] at (3.9,1.1) {}; \node[zsp] at (3.9,-0.4) {};
  \node[zsp] at (5.4,1.1) {}; \node[xsp] at (5.4,-0.4) {};
  \node[sdstate, text height=1.5ex, text depth=0.5ex] at (0,2.6) {$\psi$};
  \node[sdstate, text height=1.5ex, text depth=0.5ex] at (0,1.1) {$\psi$};
  \node[sdstate, text height=1.5ex, text depth=0.5ex] at (0,-0.4) {$0$};
\end{tikzpicture}
\;\;\overset{\text{Thm 8.8}}{=}\;\;
\begin{tikzpicture}[sd, y=1.35em]
  \draw[wire] (0,2.6) -- (4.4,2.6);
  \draw[wire] (0,1.1) -- (1.2,1.1) to[out=0,in=180] (3.2,-0.4) -- (4.4,-0.4);
  \draw[wire] (0,-0.4) -- (1.2,-0.4) to[out=0,in=180] (3.2,1.1) -- (4.4,1.1);
  \node[sdstate, text height=1.5ex, text depth=0.5ex] at (0,2.6) {$\psi$};
  \node[sdstate, text height=1.5ex, text depth=0.5ex] at (0,1.1) {$\psi$};
  \node[sdstate, text height=1.5ex, text depth=0.5ex] at (0,-0.4) {$0$};
\end{tikzpicture}
\;\;=\;\;
\begin{tikzpicture}[sd, y=1.35em]
  \draw[wire] (0,2.6) -- (2.4,2.6);
  \draw[wire] (0,1.1) -- (2.4,1.1);
  \draw[wire] (0,-0.4) -- (2.4,-0.4);
  \node[sdstate, text height=1.5ex, text depth=0.5ex] at (0,2.6) {$\psi$};
  \node[sdstate, text height=1.5ex, text depth=0.5ex] at (0,1.1) {$0$};
  \node[sdstate, text height=1.5ex, text depth=0.5ex] at (0,-0.4) {$\psi$};
\end{tikzpicture}\]

Pati and Braunstein saw this coming. Swapping the second copy into the
apparatus solves the defining equation, they note, but it is only
erasure in the ordinary sense, with the extra copy playing no part; so
they ``explicitly exclude swapping'' as a form of deleting. Then they
show that linearity leaves nothing else. Here is their argument, in
outline. Write \(\ket\psi=\alpha\ket H+\beta\ket V\) in the horizontal
and vertical polarizations, and let the machine act on \(\ket H\ket H\)
and \(\ket V\ket V\) as required, with final tells \(\ket{A_H}\) and
\(\ket{A_V}\). On the mixed input
\(\tfrac1{\sqrt2}(\ket H\ket V+\ket V\ket H)\) it may produce any fixed
state at all. Linearity then makes the output on \(\ket\psi\ket\psi\) a
\emph{quadratic} polynomial in \(\alpha\) and \(\beta\), and it must
equal the target \((\alpha\ket H+\beta\ket V)\ket\Sigma\ket{A_{\psi}}\)
for every \(\alpha\) and \(\beta\). Since the output on the mixed input
does not depend on them, they conclude that \(\ket{A_{\psi}}\) must be
linear in them: \(\ket{A_{\psi}}=\alpha\ket{A_H}+\beta\ket{A_V}\).
Normalization, for every \(\alpha\) and \(\beta\), then forces
\(\ket{A_H}\) and \(\ket{A_V}\) to be orthogonal. So the tell carries
\(\alpha\) and \(\beta\) faithfully in a two-dimensional part of itself:
``merely swapping'', in their words. Their machine is only required to
be linear. They remark that if it were known to be unitary, it ``might
work like the time-reverse of cloning''.

For unitary machines there is a one-line version in the likeness of
§11.2, and it is the headline's argument run the other way.

\begin{quote}
\textbf{Result 11.6 (no-deleting, loop form).} Let \(U\) be a unitary
with \(U(\ket\psi\ket\psi\ket A)=\ket\psi\ket\Sigma\ket{A_{\psi}}\) and
\(U(\ket\varphi\ket\varphi\ket A)=\ket\varphi\ket\Sigma\ket{A_{\varphi}}\),
and let \(s=\braket{\varphi|\psi}\neq0\). Then
\[\braket{A_{\varphi}|A_{\psi}}\;=\;s .\] The tell's two final states
are exactly as alike as the stones were: whatever was deleted from the
second stone is now in the tell.
\end{quote}

\emph{Proof.} Compare likenesses before and after, as in §11.2. Before:
\(\braket{\varphi\varphi A|\psi\psi A}=s\cdot s\cdot1=s^2\), two loops
and the tell's starting loop. After:
\(\braket{\varphi\Sigma A_{\varphi}|\psi\Sigma A_{\psi}}=s\cdot1\cdot\braket{A_{\varphi}|A_{\psi}}\),
one loop, the blank's loop, and the tell's final loop. A unitary box
meets its mirror and vanishes, so the two are equal:
\(s^2=s\,\braket{A_{\varphi}|A_{\psi}}\). Divide by \(s\neq0\).
\(\;\blacksquare\) See \texttt{Code/Ch11/NoDeleting.py}, which also
shows that a would-be deleter whose tell ignores the stone
(\(\ket{A_H}=\ket{A_V}\)), extended by linearity, sends some states to
vectors that are not of unit length.

In the picture it is a matter of counting loops. Two loops went in. One
comes out on the stones, so the missing loop must come out on the tell.
The three-cord knot shows how that happens for every state at once, and
Pati and Braunstein's argument above shows there is no other way: the
tell \emph{becomes} the stone. Result 11.6 is this book's formulation
for unitary machines; Pati and Braunstein's theorem is stronger in one
respect, since it asks only for linearity.

\vskip 0pt plus 4\baselineskip\penalty-100\vskip 0pt plus
-4\baselineskip

\begin{quote}
\textbf{Feynman's Notebook.}\nopagebreak

\emph{Uncopying.} Arun Kumar Pati and Samuel L. Braunstein, of the
University of Wales at Bangor, turned the cloning question round. Given
two photons in the same unknown polarization, can one of them be reset
to a standard state while the other is left alone? Their motive was
practical. A blank made that way could be copied onto, approximately by
a deterministic cloner or exactly by a probabilistic one, and a quantum
computer could store new information in an already computed state by
deleting the old. Classically, they note, one copy can be deleted
against the others in a perfectly reversible way, while erasing a single
copy has an energy cost (Landauer's principle). What they proved is that
the quantum case is different: this is their ``quantum no-deleting''
principle, published in \emph{Nature} 404, 164 (2000). They call it
``complementary in spirit'' to no-cloning. They also suggest that it
gives quantum files a kind of intrinsic security: no one can obliterate
one copy of an unknown file from a collection of several copies.
\end{quote}

\begin{center}\rule{0.5\linewidth}{0.5pt}\end{center}

\hypertarget{the-rosetta-table-9}{%
\section{§11.5 The Rosetta Table}\label{the-rosetta-table-9}}

The rows this chapter adds to the running dictionary, numbered on from
Chapter 10's eighty. One introduces a new picture element, the cut (87),
drawn in §11.7. The rest are new uses of elements we already had: the
copier of row 59, the cup, the doubled wire, the read-spider and the
closed picture.

\begingroup\renewcommand{\arraystretch}{1.5}

\begin{longtable}[]{@{}
  >{\raggedright\arraybackslash}p{(\columnwidth - 4\tabcolsep) * \real{0.2727}}
  >{\raggedright\arraybackslash}p{(\columnwidth - 4\tabcolsep) * \real{0.3636}}
  >{\raggedright\arraybackslash}p{(\columnwidth - 4\tabcolsep) * \real{0.3636}}@{}}
\toprule\noalign{}
\begin{minipage}[b]{\linewidth}\raggedright
\#
\end{minipage} & \begin{minipage}[b]{\linewidth}\raggedright
Algebra (Dirac)
\end{minipage} & \begin{minipage}[b]{\linewidth}\raggedright
Picture (Coecke)
\end{minipage} \\
\midrule\noalign{}
\endhead
\bottomrule\noalign{}
\endlastfoot
81 & the copy map \(\delta=\ket{00}\bra0+\ket{11}\bra1\);
\(\delta\ket0=\ket{00}\), \(\delta\ket1=\ket{11}\) (normalized, exact);
the only linear map copying \(\ket0,\ket1\) & the green dot with one leg
in and two out; a red point in comes out as two red points (row 59,
exact once the points carry their \(\tfrac1{\sqrt2}\)) \\
82 & \(\delta\ket\pm=\ket{\Phi^{\pm}}\);
\(\delta(\alpha\ket0+\beta\ket1)=\alpha\ket{00}+\beta\ket{11}\neq\ket\psi\ket\psi\)
unless \(\alpha\beta=0\) & a \textbf{green} point into the green copier
\textbf{fuses} into one cup \(\times\tfrac1{\sqrt2}\) (a half-turn point
gives the cup carrying it): a pair, not two copies;
\textbf{complementarity}: red points copy, green points fuse \\
83 & \(\delta^{\dagger}\delta=I\); the likeness
\(\langle\varphi\vert\psi\rangle\);
\(\langle\varphi\varphi\vert\psi\psi\rangle=\langle\varphi\vert\psi\rangle^2\)
& the copier meets its mirror and they leave a bare wire (fusion along
two wires, then self-loop removal; van de Wetering §4.1); a likeness is
a closed picture, and two likenesses side by side multiply \\
84 & \textbf{no-cloning:}
\(s=s^2\langle A_{\varphi}\vert A_{\psi}\rangle\Rightarrow s=0\) or
\(\lvert s\rvert=1\) & one loop \(=\) two loops only if the loop is all
or nothing \\
85 & copy, then discard one copy \(=\) dephasing \(\sum_iP_i\rho P_i\);
\textbf{broadcasting} \(\{\rho_0,\rho_1\}\) possible iff
\(\rho_0\rho_1=\rho_1\rho_0\) (Barnum et al., stated) & the doubled
copier with one output capped onto its mirror fuses into a four-legged
green dot: a look, the stone left as looked at, the record not read \\
86 & \textbf{no-deleting:}
\(\langle A_{\varphi}\vert A_{\psi}\rangle=\langle\varphi\vert\psi\rangle\);
the copy is moved into the apparatus (Pati--Braunstein) & three CNOT
cords between copy and apparatus pull into a crossing:
\(\psi\,\psi\,0\mapsto\psi\,0\,\psi\) \\
87 & rank across a cut \(\le2^k\);
\(\dim\operatorname{span}\{\psi\otimes\psi\}=3>2\); a universal copier
forces rank \(1\) & \textbf{the cut}: a thin, densely dotted grey curve
across a picture, not a wire; rank \(1\) across a cut means the picture
falls apart there, \textbf{disconnected} --- \emph{new element} \\
88 & the Forbidden Telegraph: \(p=\tfrac14\) whenever \(z_1=z_2\)
(cross); \(p=\tfrac12\) whenever \(m=z_1=z_2\) (plain); \(p=0\)
otherwise; far marginal the same & the master: at the cross setting,
fusion and Hopf leave the near record on its own \\
\end{longtable}

\endgroup

Every row is exact, scalar included; none needs Chapter 8's ``up to a
non-zero scalar''. Row 85's second half and the ideal cloner of row 88
are the only entries the pictures do not prove. The first is cited; the
second, as §11.7 explains, has no picture at all.

\begin{center}\rule{0.5\linewidth}{0.5pt}\end{center}

\hypertarget{the-theorem-twice-10}{%
\section{§11.6 The Theorem, Twice}\label{the-theorem-twice-10}}

\begin{quote}
\textbf{Result 11.3 (no-cloning, unitarity form), restated.} Let \(U\)
be any unitary on stone \(\otimes\) blank \(\otimes\) apparatus, the
apparatus of any size, with fixed starting states \(\ket0\) and
\(\ket A\). If \(U(\ket\psi\ket0\ket A)=\ket\psi\ket\psi\ket{A_{\psi}}\)
and
\(U(\ket\varphi\ket0\ket A)=\ket\varphi\ket\varphi\ket{A_{\varphi}}\),
then, with \(s=\braket{\varphi|\psi}\) and
\(a=\braket{A_{\varphi}|A_{\psi}}\),
\[s=s^2a,\qquad\text{so}\qquad s=0\ \text{ or }\ \lvert s\rvert=1 .\]
\end{quote}

\hypertarget{diracs-blackboard-10}{%
\subsection{Dirac's Blackboard}\label{diracs-blackboard-10}}

\textbf{Step 1: the likeness of the inputs.} Inner products of product
states multiply (Proposition 4.1), and \(\braket{0|0}=\braket{A|A}=1\):

\[\bigl\langle\varphi\,0\,A\,\big|\,\psi\,0\,A\bigr\rangle\;=\;\braket{\varphi|\psi}\,\braket{0|0}\,\braket{A|A}\;=\;s .\]

\textbf{Step 2: unitarity.} For any vectors \(x,y\),
\(\;\langle Ux,Uy\rangle=\bra xU^{\dagger}U\ket y=\braket{x|y}\),
because \(U^{\dagger}U=I\) (Theorem 3.2, whose proof works for a box on
any number of wires). With \(x=\varphi\,0\,A\) and \(y=\psi\,0\,A\):

\[\bigl\langle U(\varphi\,0\,A),\,U(\psi\,0\,A)\bigr\rangle\;=\;s .\]

\textbf{Step 3: substitute the cloning hypothesis}, and multiply again
by Proposition 4.1:

\[\bigl\langle\varphi\,\varphi\,A_{\varphi}\,\big|\,\psi\,\psi\,A_{\psi}\bigr\rangle\;=\;\braket{\varphi|\psi}\,\braket{\varphi|\psi}\,\braket{A_{\varphi}|A_{\psi}}\;=\;s^2a .\]

\textbf{Step 4: equate.} Steps 2 and 3 compute the same inner product,
so \(s=s^2a\).

\textbf{Step 5: take moduli.}
\(\lvert s\rvert=\lvert s\rvert^2\lvert a\rvert\le\lvert s\rvert^2\),
since \(\lvert a\rvert\le1\) by Cauchy--Schwarz for the unit vectors
\(\ket{A_{\varphi}},\ket{A_{\psi}}\). If \(s\neq0\), then
\(\lvert s\rvert\ge1\); and \(\lvert s\rvert\le1\) by Cauchy--Schwarz
again, so \(\lvert s\rvert=1\), which for unit vectors means
\(\ket\varphi=e^{i\theta}\ket\psi\).

\textbf{Step 6: the example.} \(\ket0\) and \(\ket+\) have
\(s=\tfrac1{\sqrt2}\), and \(0<\tfrac1{\sqrt2}<1\). No unitary clones
both, so none clones every state. \(\;\blacksquare\)

The whole argument is Step 4, a single line. Everything else is
bookkeeping and one inequality.

\hypertarget{coeckes-sketchbook-5}{%
\subsection{Coecke's Sketchbook}\label{coeckes-sketchbook-5}}

We work with a general unitary box \(U\) on three wires, stone, blank
and apparatus, top to bottom. The copier \(\delta\) is one example of
such a box (with the blank loaded and no apparatus), but a proof drawn
for \(\delta\) alone would prove nothing about every \(U\). Build one
closed picture. The states \(\psi\), \(0\), \(A\) enter on the left.
\(U\) acts, then its mirror image \(U^{\dagger}\), and the effects
\(\varphi\), \(0\), \(A\) close the three wires on the right. By the
mirror of a row (Theorem 3.1), the right half, \(U^{\dagger}\) followed
by the three effects, is the mirror image of the left half with
\(\varphi\) in place of \(\psi\). We evaluate the picture in two ways.

\textbf{First way, Step 1: a box meets its mirror (§3.3).} \(U\) is
unitary, so \(U\) followed by \(U^{\dagger}\) is three bare wires. The
rule is licensed by Theorem 3.2 and by nothing else.

\textbf{Step 2: stacked closed pictures multiply (Proposition 4.1).}
Three separate likenesses remain. The blank's likeness and the
apparatus's likeness are each \(1\) (Theorem 2.1).

\[\begin{tikzpicture}[sd, y=1.4em]
  \draw[wire] (0,2.4) -- (8.4,2.4);
  \draw[wire] (0,1.2) -- (8.4,1.2);
  \draw[wire] (0,0) -- (8.4,0);
  \node[sdstate, text height=1.5ex, text depth=0.5ex] at (0,2.4) {$\psi$};
  \node[sdstate, text height=1.5ex, text depth=0.5ex] at (0,1.2) {$0$};
  \node[sdstate, text height=1.5ex, text depth=0.5ex] at (0,0) {$A$};
  \node[sdbox, minimum height=5.4em] at (2.8,1.2) {$U$};
  \node[sdbox, minimum height=5.4em] at (5.6,1.2) {$U^{\dagger}$};
  \node[sdeffect, text height=1.5ex, text depth=0.5ex] at (8.4,2.4) {$\varphi$};
  \node[sdeffect, text height=1.5ex, text depth=0.5ex] at (8.4,1.2) {$0$};
  \node[sdeffect, text height=1.5ex, text depth=0.5ex] at (8.4,0) {$A$};
\end{tikzpicture}
\;\;\overset{\text{mirror}}{=}\;\;
\begin{tikzpicture}[sd, y=1.4em]
  \draw[wire] (0,2.4) -- (3.4,2.4);
  \draw[wire] (0,1.2) -- (3.4,1.2);
  \draw[wire] (0,0) -- (3.4,0);
  \node[sdstate, text height=1.5ex, text depth=0.5ex] at (0,2.4) {$\psi$};
  \node[sdstate, text height=1.5ex, text depth=0.5ex] at (0,1.2) {$0$};
  \node[sdstate, text height=1.5ex, text depth=0.5ex] at (0,0) {$A$};
  \node[sdeffect, text height=1.5ex, text depth=0.5ex] at (3.4,2.4) {$\varphi$};
  \node[sdeffect, text height=1.5ex, text depth=0.5ex] at (3.4,1.2) {$0$};
  \node[sdeffect, text height=1.5ex, text depth=0.5ex] at (3.4,0) {$A$};
\end{tikzpicture}
\;\;\overset{\text{stack}}{=}\;\;s\cdot1\cdot1 .\]

\textbf{Second way, Step 1: the hypothesis.} The left half, the states
\(\psi,0,A\) followed by \(U\), is the state \(\psi,\psi,A_{\psi}\).

\textbf{Step 2: its mirror image (Theorem 3.1).} The right half,
\(U^{\dagger}\) followed by the effects \(\varphi,0,A\), is the mirror
image of the hypothesis for \(\varphi\): the effects
\(\varphi,\varphi,A_{\varphi}\).

\textbf{Step 3: stacked closed pictures multiply (Proposition 4.1).} Now
the three likenesses are \(\psi\) against \(\varphi\), \(\psi\) against
\(\varphi\) again, and \(A_{\psi}\) against \(A_{\varphi}\):

\[\begin{tikzpicture}[sd, y=1.4em]
  \draw[wire] (0,2.4) -- (8.4,2.4);
  \draw[wire] (0,1.2) -- (8.4,1.2);
  \draw[wire] (0,0) -- (8.4,0);
  \node[sdstate, text height=1.5ex, text depth=0.5ex] at (0,2.4) {$\psi$};
  \node[sdstate, text height=1.5ex, text depth=0.5ex] at (0,1.2) {$0$};
  \node[sdstate, text height=1.5ex, text depth=0.5ex] at (0,0) {$A$};
  \node[sdbox, minimum height=5.4em] at (2.8,1.2) {$U$};
  \node[sdbox, minimum height=5.4em] at (5.6,1.2) {$U^{\dagger}$};
  \node[sdeffect, text height=1.5ex, text depth=0.5ex] at (8.4,2.4) {$\varphi$};
  \node[sdeffect, text height=1.5ex, text depth=0.5ex] at (8.4,1.2) {$0$};
  \node[sdeffect, text height=1.5ex, text depth=0.5ex] at (8.4,0) {$A$};
\end{tikzpicture}
\;\;\overset{\text{hyp.}}{=}\;\;
\begin{tikzpicture}[sd, y=1.4em]
  \draw[wire] (0,2.4) -- (3.6,2.4);
  \draw[wire] (0,1.2) -- (3.6,1.2);
  \draw[wire] (0,0) -- (3.6,0);
  \node[sdstate, text height=1.5ex, text depth=0.5ex] at (0,2.4) {$\psi$};
  \node[sdstate, text height=1.5ex, text depth=0.5ex] at (0,1.2) {$\psi$};
  \node[sdstate, text height=1.5ex, text depth=0.8ex, inner xsep=4.5pt] at (0,0) {$A_{\psi}$};
  \node[sdeffect, text height=1.5ex, text depth=0.5ex] at (3.6,2.4) {$\varphi$};
  \node[sdeffect, text height=1.5ex, text depth=0.5ex] at (3.6,1.2) {$\varphi$};
  \node[sdeffect, text height=1.5ex, text depth=0.8ex, inner xsep=4.5pt] at (3.6,0) {$A_{\varphi}$};
\end{tikzpicture}
\;\;\overset{\text{stack}}{=}\;\;s\cdot s\cdot a .\]

The same picture has been evaluated twice, so \(s=s^2a\). The first way
leaves one loop that matters; the second leaves two, side by side. The
last step, from \(s=s^2a\) to ``\(s=0\) or \(\lvert s\rvert=1\)'', is a
fact about complex numbers, and the sketchbook borrows it from the
blackboard's Step 5. \(\;\blacksquare\)

The worked instance of the first way's Step 1 is §11.2's picture: when
\(U\) is the copier (with its blank inside, §11.1), the box meeting its
mirror is fusion along two wires and a self-loop removed. In that case
the second way fails at Step 1, because the hypothesis is false:
\(\delta\) sends \(\ket\psi\) to \(\alpha\ket{00}+\beta\ket{11}\), not
to \(\ket\psi\ket\psi\).

\hypertarget{the-verdict-10}{%
\subsection{The verdict}\label{the-verdict-10}}

The blackboard wins on length. Its proof is one line of inner products,
\(s=\langle U(\varphi0A),U(\psi0A)\rangle=s^2a\), and the rest is an
inequality. The sketchbook needs the box to meet its mirror, the mirror
of the hypothesis, and the rule that separate closed pictures multiply.
Then it hands the final step, a fact about numbers, back to the
blackboard.

The picture wins on insight, and it can say why. The input likeness is
one loop; the copies close into two \emph{separate} loops, because a
clone is a product and a product is a picture in two pieces. A closed
loop is a number, so a faithful copier would make a number equal its own
square. The picture also gives the positive reading that the blackboard
does not show. The copier copies its own basis: red points pass through
the green copier as two red points, and green points fuse into a cup, a
pair and not two copies. That is complementarity, and it is why the
forbidden machine fails exactly where it does. None of this is a bound
on how well an imperfect copier can do, a security claim, or a statement
about any machine outside the linear theory. A rewrite proves an
equality of processes, and here the equality is \(s=s^2a\).

\vskip 0pt plus 4\baselineskip\penalty-100\vskip 0pt plus
-4\baselineskip

\begin{quote}
\textbf{Coecke's Sketchbook.}\nopagebreak

\emph{What the loops do not say.} The one-loop, two-loop argument is
about \emph{perfect} copies: it shows that \(s=s^2a\) fails, and nothing
more. It does not say how close to a copy a real machine can come, or at
what cost. That question has a literature of its own, on approximate and
optimal cloning machines; Pati and Braunstein's closing paragraph points
to it, and suggests that one could try to build an approximate
\emph{deleting} machine by analogy. Nor does the picture know why
\(\lvert a\rvert\le1\); that is Cauchy--Schwarz, from the blackboard.
And it cannot tell a unitary box from any other by its outline. Chapter
3's warning holds here with full force: cancel a box against its mirror
only because the blackboard has shown that the box is unitary, and say
so. The copier is an isometry, so \(\delta^{\dagger}\delta\) cancels and
\(\delta\delta^{\dagger}\) does not. What the loops \emph{do} give, and
give at a glance, is the reason: one loop against two.
\end{quote}

\begin{center}\rule{0.5\linewidth}{0.5pt}\end{center}

\hypertarget{the-forbidden-telegraph}{%
\section{§11.7 The Forbidden Telegraph}\label{the-forbidden-telegraph}}

The Crab's last hope is the boldest. The post to Delft takes a week.
Suppose one half of a fresh pair stays here and the other goes to Delft,
and suppose the Delft end had a perfect copier. Could the near end then
send word by choosing how to read its half?

\textbf{On the blackboard, the machine that cannot exist.} Share
\(\ket{\Phi^+}\) between a near stone \(a\) and a far stone \(b\). The
near end reads \(a\) at a \textbf{setting}, plain (the \(\ket0,\ket1\)
basis) or cross (the \(\ket\pm\) basis), and gets a record \(m\). The
far end copies \(b\) and reads \emph{both} copies plain, getting records
\(z_1\) and \(z_2\). Regrouping the four terms,
\(\ket{\Phi^+}=\tfrac1{\sqrt2}(\ket{00}+\ket{11})=\tfrac1{\sqrt2}(\ket{++}+\ket{--})\).
So at plain the far stone is left in \(\ket0\) or \(\ket1\), and at
cross in \(\ket+\) or \(\ket-\), each half the time.

Now suppose, against everything so far, a machine \(K\) with
\(K\ket\psi=\ket\psi\otimes\ket\psi\) for every \(\ket\psi\), and apply
it on each branch. At plain the far end gets \(\ket{00}\) or
\(\ket{11}\), and its two records always agree. At cross it gets
\(\ket{++}\) or \(\ket{--}\), and each copy, read plain, is a fair coin
independent of the other, so the records agree only half the time. The
far end, with no letter, sees agreement with probability \(1\) or
\(\tfrac12\) according to the near setting. Run it on many pairs and the
near end is sending a message: a telegraph.

That is a blackboard computation of something that cannot be drawn.
\(K\) is not linear, and every picture in this book is a linear map. The
sketchbook cannot draw the forbidden machine; only the blackboard can
state it and watch it contradict itself. So we turn to the machine that
does exist.

\vskip 0pt plus 4\baselineskip\penalty-100\vskip 0pt plus
-4\baselineskip

\begin{quote}
\textbf{Dirac's Blackboard.}\nopagebreak

\emph{Why only a nonlinear machine could read the setting.} Before any
letter arrives, the far stone is \(\tfrac12 I\) whatever the near end
does (Result 9.3). But the same matrix can be written as two different
mixtures:
\[\tfrac12 I\;=\;\tfrac12\ket0\bra0+\tfrac12\ket1\bra1\;=\;\tfrac12\ket+\bra++\tfrac12\ket-\bra- .\]
The plain setting prepares the first mixture at the far end, and the
cross setting prepares the second. A linear machine acts on the matrix,
\(\tfrac12 I\), and cannot tell how it was mixed. \(K\) acts stone by
stone, and it sends the two mixtures to
\(\tfrac12(\ket{00}\bra{00}+\ket{11}\bra{11})\) and
\(\tfrac12(\ket{++}\bra{++}+\ket{--}\bra{--})\). These differ: the first
gives agreeing plain readings with probability \(1\), the second with
probability \(\tfrac12\). Telling apart two descriptions of the same
density matrix is exactly what linearity forbids. It is also exactly
what a telegraph would need.
\end{quote}

\textbf{The Forbidden Telegraph.} Now the real press. Draw the whole
scene in unfolded doubled form, so that every line shows, as in Chapter
9 §9.6 and Chapter 10 §10.7. On the left, the doubled pair
\(\ket{\Phi^+}\bra{\Phi^+}\): two cups, one joining the near stone \(a\)
to the far stone \(b\) and one joining their mirror copies \(\bar a\)
and \(\bar b\), with \(\tfrac12\) written beside them. As always, cup
\(=\ket{00}+\ket{11}\), so each normalized pair is
\(\tfrac1{\sqrt2}\,\)cup and the doubled pair is \(\tfrac12\) cup
\(\otimes\) cup. At the near end, \(a\) and \(\bar a\) run into Chapter
9's read-spider, which emits the classical near record \(m\). At the
cross setting each of \(a\) and \(\bar a\) first passes through a yellow
Hadamard box (row 56), which turns the plain reading into a crosswise
one; at the plain setting the two boxes are simply absent. At the far
end, \(b\) runs into the green copier, whose outputs are the first and
second struck stones \(b_1\) and \(b_2\). The mirror copy \(\bar b\)
runs into the copier's conjugate, which is the green copier again since
\(\delta\) is real, with outputs \(\bar b_1\) and \(\bar b_2\). Each
struck stone is read plain: one read-spider takes \(b_1\) with
\(\bar b_1\) and emits the classical record \(z_1\), and another takes
\(b_2\) with \(\bar b_2\) and emits \(z_2\). To bring each struck stone
next to its own mirror, \(\bar b_1\) and \(b_2\) cross once. A crossing
neither joins nor breaks a line.

\begin{center}
\begin{tikzpicture}[sd, x=1.3em, y=1.15em,
  zsp/.style={circle,draw,line width=0.9pt,fill=green!55,minimum size=1.2em,inner sep=0pt},
  hbox/.style={draw,line width=0.9pt,fill=yellow!85,minimum size=0.9em,inner sep=1pt}]
  \node[sdlabel, font=\normalsize, anchor=east] at (-3.9,8) {$\tfrac12$};
  % the doubled pair: outer cup (a-bar, b-bar), inner cup (a, b)
  \sdcup{0}{12}{4}
  \sdcup{0}{10}{8}
  % near end: a-bar, a through Hadamards into the near read-spider
  \draw[wire] (0,12) -- (5,12) -- (6.5,11);
  \draw[wire] (0,10) -- (5,10) -- (6.5,11);
  \node[hbox] at (3,12) {};
  \node[hbox] at (3,10) {};
  \draw[cwire] (6.5,11) -- (14,11);
  % far end: b and b-bar into the copier and its conjugate
  \draw[wire] (0,8) -- (4,8);
  \draw[wire] (0,4) -- (4,4);
  % the four struck lines; b2 and b-bar-1 cross once
  \draw[wire] (4,8) -- (5,8.8) -- (10,8.8) -- (11,8);
  \draw[wire] (4,8) -- (5,7.2) -- (6,7.2) to[out=0,in=180] (9,4.8) -- (10,4.8) -- (11,4);
  \draw[wire] (4,4) -- (5,4.8) -- (6,4.8) to[out=0,in=180] (9,7.2) -- (10,7.2) -- (11,8);
  \draw[wire] (4,4) -- (5,3.2) -- (10,3.2) -- (11,4);
  \draw[cwire] (11,8) -- (14,8);
  \draw[cwire] (11,4) -- (14,4);
  % spiders on top
  \node[zsp] at (6.5,11) {};
  \node[zsp] at (4,8) {};
  \node[zsp] at (4,4) {};
  \node[zsp] at (11,8) {};
  \node[zsp] at (11,4) {};
  % labels
  \node[sdlabel,anchor=south] at (1.5,12.05) {$\bar a$};
  \node[sdlabel,anchor=south] at (1.5,10.05) {$a$};
  \node[sdlabel,anchor=south] at (1.5,8.05) {$b$};
  \node[sdlabel,anchor=north] at (1.5,3.95) {$\bar b$};
  \node[sdlabel,anchor=south] at (7.5,8.85) {$b_1$};
  \node[sdlabel,anchor=north] at (5.6,7.15) {$b_2$};
  \node[sdlabel,anchor=south] at (5.6,4.85) {$\bar b_1$};
  \node[sdlabel,anchor=north] at (7.5,3.15) {$\bar b_2$};
  \node[sdlabel,anchor=west] at (14.2,11) {$m$};
  \node[sdlabel,anchor=west] at (14.2,8) {$z_1$};
  \node[sdlabel,anchor=west] at (14.2,4) {$z_2$};
\end{tikzpicture}
\end{center}

The lines are the near stone \(a\) and its mirror \(\bar a\), each
running from its cup through its Hadamard box into the near read-spider;
the far stone \(b\) and its mirror \(\bar b\), from the cups into the
copier and its conjugate; the struck stones \(b_1\) and \(b_2\) and
their mirrors \(\bar b_1\) and \(\bar b_2\), each a leg of a copier
running to a far read-spider; and the three classical records \(m\),
\(z_1\) and \(z_2\). The drawing is at the cross setting. We call this
diagram \textbf{The Forbidden Telegraph}.

\begin{quote}
\textbf{Result 11.7 (the Forbidden Telegraph).} The Forbidden Telegraph
evaluates, exactly, to the joint distribution of the three records:
\[\text{cross:}\quad p(m,z_1,z_2)=\tfrac14\,\delta_{z_1z_2},\]
\[\text{plain:}\quad p(m,z_1,z_2)=\tfrac12\,\delta_{mz_1}\delta_{z_1z_2}.\]
Here \(\delta_{z_1z_2}\) is the Kronecker delta, \(1\) when the two
records agree and \(0\) otherwise (likewise \(\delta_{mz_1}\)); it is
not the copier. At both settings the struck stones always agree,
\(z_1=z_2\), and the far records have the same distribution,
\(\sum_m p=\tfrac12\,\delta_{z_1z_2}\). \textbf{No signal crosses.} At
the cross setting the near record is a fair toss independent of the far
records.
\end{quote}

\emph{Proof, on the blackboard.} At plain, the near reading gives \(m\)
with probability \(\tfrac12\) and leaves \(b\) in \(\ket m\); the copier
makes \(\ket{mm}\), and both far records read \(m\). At cross, the near
reading gives \(\ket\pm\) with probability \(\tfrac12\) each and leaves
\(b\) in the same \(\ket\pm\); the copier makes \(\ket{\Phi^{\pm}}\)
(Result 11.2), whose two plain readings agree and are each a fair coin,
whichever sign came. Sum over \(m\) for the far marginal.
\(\;\blacksquare\)

\emph{Proof, in the sketchbook.} Only connectivity matters (§8.5), and a
cup is a bent wire, so the near read-spider is joined to the copier by
one wire and to the copier's conjugate by another. Every scalar below is
exact.

At the \textbf{plain} setting every dot is green and phase-free.
\textbf{Fusion} (Theorem 8.1) merges the near read-spider, the copier,
its conjugate and both far read-spiders into one green dot. The six
wires among them close two self-loops, which are \textbf{removed} (van
de Wetering §4.1). What is left is \(\tfrac12\) times a green dot with
three legs, \(m\), \(z_1\) and \(z_2\): \(\tfrac12\sum_i\ket{iii}\). All
three records agree.

At the \textbf{cross} setting, five steps.

\textbf{Step 1: colour change (Theorem 8.2).} Put two Hadamard boxes on
the record \(m\), which changes nothing, since \(H^2=I\). The record leg
is a leg like any other to the rules; the double stroke is only a
reminder (§9.1). Now the near read-spider wears a Hadamard on every leg,
so it is a red dot \(x\) of phase \(0\), with one Hadamard box left on
\(m\).

\textbf{Step 2: fusion (Theorem 8.1).} The copier, its conjugate and the
two far read-spiders are green and joined, so they fuse into one green
dot \(G\). The four wires among them contain one wire too many for a
tree, which leaves one self-loop on \(G\); remove it (van de Wetering
§4.1). \(G\) keeps the legs \(z_1\) and \(z_2\), and two wires to \(x\),
one through each cup.

\[\tfrac12\;
\begin{tikzpicture}[sd, x=1.2em, y=1.2em,
  zsp/.style={circle,draw,line width=0.9pt,fill=green!55,minimum size=1.2em,inner sep=0pt},
  xsp/.style={circle,draw,line width=0.9pt,fill=red!75,minimum size=1.2em,inner sep=0pt},
  hbox/.style={draw,line width=0.9pt,fill=yellow!85,minimum size=0.9em,inner sep=1pt}]
  \draw[wire] (0,1.6) to[out=-120,in=120] (0,-1.2);
  \draw[wire] (0,1.6) to[out=-60,in=60] (0,-1.2);
  \draw[wire] (0,1.6) -- (1.6,1.6);
  \draw[cwire] (1.6,1.6) -- (4.6,1.6);
  \node[hbox] at (1.6,1.6) {};
  \draw[cwire] (0,-1.2) -- (1.2,-0.6) -- (4.6,-0.6);
  \draw[cwire] (0,-1.2) -- (1.2,-1.8) -- (4.6,-1.8);
  \node[xsp] at (0,1.6) {};
  \node[zsp] at (0,-1.2) {};
  \node[sdlabel,anchor=east] at (-0.5,1.6) {$x$};
  \node[sdlabel,anchor=east] at (-0.5,-1.2) {$G$};
  \node[sdlabel,anchor=west] at (4.8,1.6) {$m$};
  \node[sdlabel,anchor=west] at (4.8,-0.6) {$z_1$};
  \node[sdlabel,anchor=west] at (4.8,-1.8) {$z_2$};
\end{tikzpicture}
\;\;\overset{\text{Hopf}}{=}\;\;
\tfrac14\;
\begin{tikzpicture}[sd, x=1.2em, y=1.2em,
  zsp/.style={circle,draw,line width=0.9pt,fill=green!55,minimum size=1.2em,inner sep=0pt},
  xsp/.style={circle,draw,line width=0.9pt,fill=red!75,minimum size=1.2em,inner sep=0pt},
  hbox/.style={draw,line width=0.9pt,fill=yellow!85,minimum size=0.9em,inner sep=1pt}]
  \draw[wire] (0,1.6) -- (1.6,1.6);
  \draw[cwire] (1.6,1.6) -- (4.6,1.6);
  \node[hbox] at (1.6,1.6) {};
  \draw[cwire] (0,-1.2) -- (1.2,-0.6) -- (4.6,-0.6);
  \draw[cwire] (0,-1.2) -- (1.2,-1.8) -- (4.6,-1.8);
  \node[xsp] at (0,1.6) {};
  \node[zsp] at (0,-1.2) {};
\end{tikzpicture}
\;\;\overset{\text{colour}}{=}\;\;
\tfrac14\;
\begin{tikzpicture}[sd, x=1.2em, y=1.2em,
  zsp/.style={circle,draw,line width=0.9pt,fill=green!55,minimum size=1.2em,inner sep=0pt}]
  \draw[cwire] (0,1.6) -- (3.4,1.6);
  \draw[cwire] (0,-1.2) -- (1.2,-0.6) -- (3.4,-0.6);
  \draw[cwire] (0,-1.2) -- (1.2,-1.8) -- (3.4,-1.8);
  \node[zsp] at (0,1.6) {};
  \node[zsp] at (0,-1.2) {};
\end{tikzpicture}\]

\textbf{Step 3: Hopf (Theorem 8.6).} \(G\) and \(x\) are a green dot and
a red dot joined by two wires. Unfuse the legs that take no part (fusion
read backwards), so that the shape is exactly Theorem 8.6's. Apply the
rule: the two wires fall away, with the factor \(\tfrac12\) of Theorem
8.6. Then fuse back. The red dot keeps only its leg toward \(m\), and
\(G\) keeps only \(z_1\) and \(z_2\).

\textbf{Step 4: colour change again.} A red point followed by a Hadamard
box is a green point, exactly: \(H\bigl(\sqrt2\ket0\bigr)=\ket0+\ket1\).

\textbf{Step 5: read off.} The result is
\(\tfrac14\,(\ket0+\ket1)_m\otimes(\ket{00}+\ket{11})_{z_1z_2}\): the
near record is a fair toss on its own, and the far records agree and are
a fair toss between them. \(\;\blacksquare\) See
\texttt{Code/Ch11/ForbiddenTelegraph.py}, which builds the doubled
diagram at both settings, counts its lines, checks both distributions
exactly, and derives the disconnection with the rules.

At the cross setting the near record has come off the picture
altogether. To say that precisely we need one new mark.

\textbf{In the sketchbook: the cut.} Draw a thin, densely dotted grey
curve across a picture, dividing it into two parts. We call it a
\textbf{cut}. It is not a wire and carries nothing; it is a line of
inspection, drawn over the picture to ask how the two sides are joined.
The two parts are joined only through the wires the curve crosses. (It
is drawn dotted and grey, never dashed: a dashed wire is Chapter 8's
Hadamard edge.) Here is the telegraph at the cross setting after the
Hopf step, with a cut drawn between the near record and the far records.
It crosses no wire at all:

\[\tfrac14\;
\begin{tikzpicture}[sd, x=1.2em, y=1.2em,
  zsp/.style={circle,draw,line width=0.9pt,fill=green!55,minimum size=1.2em,inner sep=0pt},
  xsp/.style={circle,draw,line width=0.9pt,fill=red!75,minimum size=1.2em,inner sep=0pt},
  hbox/.style={draw,line width=0.9pt,fill=yellow!85,minimum size=0.9em,inner sep=1pt}]
  \draw[wire] (0,1.6) -- (1.6,1.6);
  \draw[cwire] (1.6,1.6) -- (4.6,1.6);
  \node[hbox] at (1.6,1.6) {};
  \draw[cwire] (0,-1.2) -- (1.2,-0.6) -- (4.6,-0.6);
  \draw[cwire] (0,-1.2) -- (1.2,-1.8) -- (4.6,-1.8);
  \node[xsp] at (0,1.6) {};
  \node[zsp] at (0,-1.2) {};
  \draw[sdcut] (-0.8,0.75) -- (5.8,0.75);
  \node[sdlabel,anchor=west] at (4.8,1.6) {$m$};
  \node[sdlabel,anchor=west] at (4.8,-0.6) {$z_1$};
  \node[sdlabel,anchor=west] at (4.8,-1.8) {$z_2$};
\end{tikzpicture}\]

\textbf{On the blackboard: rank across a cut.} Bend all the open wires
on one side of the cut round to be inputs (Theorem 6.5), so that the
picture becomes a process from one side to the other. Its matrix has a
\textbf{rank}, the dimension of its image, and we call that the
picture's \textbf{rank across the cut}. For a state on two parts this is
its Schmidt rank. Every path from one side to the other runs through a
wire that the cut crosses, so the process factors through those wires.
If the cut crosses \(k\) qubit wires, the matrix factors through a space
of dimension \(2^k\), and so its rank is at most \(2^k\). A cut crossing
no wire gives rank at most \(2^0=1\): the picture is a product of its
two parts, and it falls apart there. We call such a picture
\textbf{disconnected} across the cut. Rank \(1\) across the cut in the
telegraph means \(p(m,z_1,z_2)=p(m)\,p(z_1,z_2)\): independence, at the
cross setting. At the plain setting the three records hang on one green
dot, and any cut separating \(m\) from \(z_1,z_2\) crosses a wire.

\textbf{\emph{Why this definition.}} Why draw a cut rather than simply
compute a rank? Because in the picture the rank can be read from where
the curve can pass. Rewriting moves wires, and a picture that has fallen
apart shows it by admitting a cut through empty paper. The picture gives
an upper bound for free. A lower bound, such as ``this cup really is
connected'', still needs the blackboard: the cup's rank is \(2\) because
its matrix is the identity.

As in Chapter 9, the independence is honest only on average. Across the
cross setting's cut, the near record is independent of the far ones. At
the plain setting it is not. Condition on \(m\) there, and the far
records follow it exactly, \(z_1=z_2=m\). What the far end cannot see is
the \emph{setting}, because it never sees \(m\). Averaged over the near
outcome, its records are \(\tfrac12\delta_{z_1z_2}\) either way. That is
Chapter 9's no-signalling (Result 9.3, and the master of §9.6), now with
a copier on the far wire.

In the dialogue, just after the press drops its two struck stones and
before they are read, the Crab glances at \emph{Drawing Hands}: two
hands, each drawing the other.

The cut also says why no better press could ever have helped the Crab.

\begin{quote}
\textbf{Result 11.8 (a copier cannot widen a wire; a universal copier
would disconnect every cup).} (i) For every linear process \(D\) with
one qubit wire in and two out, the picture \((I\otimes D)\,\)cup has
rank at most \(2\) across a cut crossing the wire into \(D\).
Equivalently, the image of \(D\) has dimension at most \(2\). (ii) The
clones \(\ket\psi\otimes\ket\psi\) span a space of dimension \(3\). So
no linear process copies every state. (iii) More generally, on a wire of
dimension \(r\) the clones span a space of dimension at least
\(r(r+1)/2\). A linear process copying every state of such a wire exists
only if \(r(r+1)/2\le r\), that is \(r\le1\), and on a wire of dimension
\(1\) every cup has rank \(1\) across every cut: disconnected. (iv) What
the copier actually does to a cup is exact: \((I\otimes\delta)\,\)cup
\(=\ket{000}+\ket{111}\), Chapter 7's three-legged node (row 46),
connected across every cut.
\end{quote}

\emph{Proof.} (i) The cut crosses one wire, so the rank is at most
\(2^1=2\); bending \((I\otimes D)\,\)cup back into \(D\) (Theorem 6.5)
shows the same number is the dimension of \(D\)'s image, which is
spanned by \(D\ket0\) and \(D\ket1\). (ii) The clones include
\(\ket{00}\), \(\ket{11}\) and \(\ket{++}\), and
\(\ket{++}-\tfrac12\ket{00}-\tfrac12\ket{11}=\tfrac12\bigl(\ket{01}+\ket{10}\bigr)\);
these three are independent. They also span everything, since every
clone
\(\alpha^2\ket{00}+\alpha\beta(\ket{01}+\ket{10})+\beta^2\ket{11}\) is a
combination of them. A universal copier would need an image of dimension
\(3\) through a wire that allows \(2\). (iii) For any two basis vectors
\(u,v\) of the wire,
\((u+v)\otimes(u+v)-u\otimes u-v\otimes v=u\otimes v+v\otimes u\). So
the span of the clones contains every \(e_i\otimes e_i\) and every
\(e_i\otimes e_j+e_j\otimes e_i\) with \(i<j\), which are
\(r+\tfrac12r(r-1)=\tfrac12r(r+1)\) independent vectors. The inequality
\(\tfrac12r(r+1)\le r\) gives \(r\le1\). A cup on a one-dimensional wire
is a single term, a product. (iv) The cup is a phase-free green dot with
two legs, and the copier fuses with it (Theorem 8.1). \(\;\blacksquare\)
See \texttt{Code/Ch11/ForbiddenTelegraph.py}, which checks (i) on random
\(D\), (ii) on random clones, and (iv) exactly.

\[\begin{tikzpicture}[sd]
  \sdcup{0}{2.6}{0}
  \draw[wire] (0,2.6) -- (6.4,2.6);
  \draw[wire] (0,0) -- (3.6,0);
  \draw[wire] (3.6,0.6) -- (6.4,0.6);
  \draw[wire] (3.6,-0.6) -- (6.4,-0.6);
  \node[sdbox, minimum height=2.3em] at (3.6,0) {$D$};
  \draw[sdcut] (7.6,1.6) -- (3.0,1.6) to[out=180,in=90] (1.9,0.4) -- (1.9,-1.6);
  \node[sdlabel,anchor=west] at (6.6,2.6) {$a$};
  \node[sdlabel,anchor=west] at (6.6,0.6) {$b_1$};
  \node[sdlabel,anchor=west] at (6.6,-0.6) {$b_2$};
\end{tikzpicture}
\;\;\text{rank}\le2
\qquad\qquad
\begin{tikzpicture}[sd,
  zsp/.style={circle,draw,line width=0.9pt,fill=green!55,minimum size=1.3em,inner sep=0pt}]
  \sdcup{0}{2.2}{0}
  \draw[wire] (0,2.2) -- (4.6,2.2);
  \draw[wire] (0,0) -- (2.2,0);
  \draw[wire] (2.2,0) -- (3.3,0.75) -- (4.6,0.75);
  \draw[wire] (2.2,0) -- (3.3,-0.75) -- (4.6,-0.75);
  \node[zsp] at (2.2,0) {};
\end{tikzpicture}
\;\;\overset{\text{fusion}}{=}\;\;
\begin{tikzpicture}[sd,
  zsp/.style={circle,draw,line width=0.9pt,fill=green!55,minimum size=1.3em,inner sep=0pt}]
  \draw[wire] (0,0.3) -- (1.4,1.6) -- (2.6,1.6);
  \draw[wire] (0,0.3) -- (1.4,0.3) -- (2.6,0.3);
  \draw[wire] (0,0.3) -- (1.4,-1.0) -- (2.6,-1.0);
  \node[zsp] at (0,0.3) {};
\end{tikzpicture}\]

On the left, any process \(D\) on one half of a cup, with the cut
crossing the single wire that feeds it: whatever \(D\) does, everything
on the far side of the cut came through one wire. On the right, the
green copier on one half of a cup. The cup is not copied. It grows a
leg, into the three-legged state of Chapter 7, and it stays connected:
any cut that separates one leg from the other two crosses one wire, and
the rank there is \(2\). This book's pictures are built on the connected
cup (Chapter 6). A universal copier could exist only in a world where
every cup came apart, so the universal copier and the connected cup
cannot live in the same picture. That is the precise sense in which a
universal copier would disconnect every cup. The argument is the book's
own; we attribute it to no one.

\vskip 0pt plus 4\baselineskip\penalty-100\vskip 0pt plus
-4\baselineskip

\begin{quote}
\textbf{Coecke's Sketchbook.}\nopagebreak

\emph{The cut is not a wire.} The dotted curve is drawn \emph{over} a
picture, never \emph{in} it. It carries no system, it is not counted
among the lines, and it changes no value; erase it and the picture means
what it meant before. It is a question put to the picture: how many
wires join this side to that? Its answer is an upper bound. If the curve
crosses \(k\) qubit wires, the rank across it is at most \(2^k\), but it
may be less. A product state drawn carelessly, as a single box with two
wires, admits no empty cut until it is rewritten into two pieces. So a
cut through empty paper \emph{proves} disconnection, while a cut that
must cross a wire proves nothing by itself about connection. For that
you go to the blackboard, or to a rewrite that ends in a spider no cut
can avoid, like the plain telegraph's single green dot. In this chapter
the cut crosses nothing at the cross setting, and that is the whole of
the no-telegraph theorem in one stroke of grey.
\end{quote}

\vskip 0pt plus 4\baselineskip\penalty-100\vskip 0pt plus
-4\baselineskip

\begin{quote}
\textbf{Feynman's Notebook.}\nopagebreak

\emph{Faster than the post.} The link between copying and signalling is
older than this chapter's small telegraph. Pati and Braunstein, citing
both 1982 papers, put it in one sentence: if an arbitrary state could be
cloned, then with non-local resources one could send signals faster than
light. They go on to ask whether their own no-deleting principle could
likewise be derived from the impossibility of faster-than-light
signalling, as no-cloning can, and they leave the question open. The
telegraph computed in this section, a copier on one half of a Bell pair,
read plain at the far end and at either setting at the near end, is this
book's own worked example of the remark. It is not taken from either
paper. Note which way the logic runs. No-signalling was proved in
Chapter 9 for every trace-preserving channel at the near end (Result
9.3): the far half stays \(\tfrac12 I\), and the real copier, being
linear, acts on that \(\tfrac12 I\) alone and inherits the law. Only a
nonlinear machine could break it, and the sketchbook has no way to draw
one.
\end{quote}

\begin{center}\rule{0.5\linewidth}{0.5pt}\end{center}

\hypertarget{what-may-be-passed-on}{%
\section{§11.8 What May Be Passed On}\label{what-may-be-passed-on}}

Let us gather it up. The green copier copies exactly the states of its
own basis. Fed a red point, it gives two red points; fed a green one, it
fuses into a cup and makes a pair. Linearity forbids any machine that
copies the plain states from doing better, apparatus or no apparatus
(§11.1). Unitarity forbids any machine at all from copying two states
unless they are orthogonal or the same, because a faithful copy would
turn one loop into two, and a likeness into its own square (§11.2,
proved twice in §11.6). A copy made and thrown away is a look that
nobody read, and even the gentler demand of broadcasting fails unless
the states commute (§11.3). A copy cannot be unmade either. The best a
deleter can do is move the copy into its own works, where it sits
exactly as unlike as the stones were (§11.4). And a perfect copier on
one half of a pair would be a telegraph. The real one is not, because at
the cross setting the near record falls away from the far ones, across a
cut that crosses nothing. The only copier that could do better exists in
a world where every cup has come apart (§11.7).

So Chapter 10's question has its answer. Teleportation moved the message
and did not copy it, and it could not have done otherwise. A quantum
state can be passed on, or answered, but never duplicated.

That fact has a use, and it is the next chapter. Barnum and his
colleagues closed their paper with one: in quantum cryptography over a
noisy channel, their theorem leaves an eavesdropper no means of taking
the signal meant for the legitimate receiver without disturbing what
that receiver gets. Think of a spy on the line. She cannot copy what
passes, so she cannot keep a copy and let the original go by untouched.
To learn anything she must look. A look is made at some setting, plain
or cross, and this chapter has shown what a look at the wrong setting
does: it turns a crosswise stone grey, as Eve's customers found. Chapter
12 makes those two settings into the two keys of a lock, built of light.
There the mark a spy's look leaves on the stones becomes an alarm.

\hypertarget{yesno-duet-with-eavesdropper}{%
\chapter{Yes--No Duet, with
Eavesdropper}\label{yesno-duet-with-eavesdropper}}

\begin{center}\rule{0.5\linewidth}{0.5pt}\end{center}

\emph{寒灯や}\\
\emph{二人の窓に}\\
\emph{人の息}

\emph{kantō ya}\\
\emph{futari no mado ni}\\
\emph{hito no iki}

\emph{Lamp on a cold night ---}\\
\emph{on the window of the two}\\
\emph{someone else's breath.}

\begin{center}\rule{0.5\linewidth}{0.5pt}\end{center}

\emph{{[}The Tortoise's parlour, early evening. On the table lies a
large black tray, and on it, set out beforehand, stands the Tortoise's
kit. In the middle is a small brass lantern with a lever, and a port on
either side. From its left port a slender glass rod, on low brass forks
a finger's breadth above the lacquer, runs to a clear crystal in a brass
collar at the tray's left end. From its right port another rod runs most
of the way along to a short brass sleeve, and from the sleeve a further
rod runs on to another crystal and collar at the right end. Beneath each
rod, in the lacquer, lies its reflection; beneath the lantern, the
sleeve and the crystals, theirs. Off the tray, on the cloth, a deep box
of draughtsmen stands open; a small brass bell sits on the table's
corner; beside the tray the Tortoise's drawing-cloth is spread, with her
string and rings on it. The door opens, and the Tortoise and Achilles
come in out of the dark.{]}}

\emph{{[}The Tortoise finds the matches on the dresser by feel and
lights the room lamp that stands high there, and the parlour comes up
out of the dark around it. It is a little after five and full night
already; the window on the long wall is black, and the cold stands off
it into the room. The stove, banked since morning, wakes when she opens
its damper, with a small roar behind its mica door and a first tick from
its iron. The table runs down the middle of the room, its left end
towards the stove and its right towards the dresser, and the kit lies
along it as she left it this morning: the long rod to the stove end has
a faint green in the thick of its glass; the one from the lantern to the
sleeve carries a single seed of air near its middle; the far rod, on to
the dresser end, is the palest of the three. The supper bell on the
table's corner is plain brass, its turned wooden handle worn pale by
forty years of suppers, and the draughts box is older still, its pieces
rubbed smooth at the rims. On the end wall beyond the dresser end,
beside the dresser, hangs a print in a narrow black frame --- Escher's
\textup{Three Worlds}, a pond: leaves on its surface, trees reflected in
it, a fish beneath. The room lamp reaches it from the side, and it takes
the light quietly, as it has every evening since long before the tray
was set out on that table, and as it will after. On the lower pane of
the black window, at the very bottom corner, the frost has made a start:
a few feathers, fine as hair, each going its own way.{]}}

\textbf{Achilles:} \emph{{[}stamping on the mat{]}} I can't feel my
toes. Every rut in your lane has a skin of ice on it, and I found them
all.

\textbf{Tortoise:} You would walk on the crown of the road. Boots by the
stove, and I'll see to the kettle.

\emph{{[}Achilles sits down on the fender and works his boots off, and
stands them by the stove, where they begin at once to steam. His gloves
he pulled off at the door by the fingertips, as one does with cold
hands, and they lie on the fender beside him inside out, with his
woollen cap, inside out too, dropped on top of them. The Tortoise fills
the kettle from the jug on the dresser and pushes it over the
firebox.{]}}

\textbf{Achilles:} \emph{{[}crossing to the window, rubbing his
hands{]}} It's begun already, look --- ferns, from the bottom corner up.
And nothing like this morning's. Not a frond of it in the same place.

\textbf{Tortoise:} It never is. It starts fresh each night, from
whatever the cold finds on the glass.

\emph{{[}Achilles puts a fingertip, warm from the rubbing, to the glass
where the first feathers have started, to feel how cold it is. They are
gone at once from under it, and a small round of wet bare glass is left
in their place, with the frost beginning again all round its edge.{]}}

\textbf{Achilles:} Mr Crab stood on his step with a candle to see us
off, and then came pattering down the path after us in his slippers.
\emph{{[}He takes a folded letter from his coat.{]}} This had come by
the last post, just as we were leaving. It's to him, from Alice ---
Bob's Alice. He wouldn't open it. He pressed it on me and said, ``It is
about light, and she has the glass for it.''

\textbf{Tortoise:} And what are we to do with it?

\textbf{Achilles:} Read it here, and tell him in the morning what it
says. He said he had done enough looking for one day.

\textbf{Tortoise:} \emph{{[}setting out two cups{]}} Poor Mr Crab. A
whole afternoon with his claws off the machine he has wanted since he
was small. We'll give him a good account.

\textbf{Achilles:} \emph{{[}turning to the table{]}} And you've set all
this out. When?

\textbf{Tortoise:} This morning, before I came over to his. I had a
feeling the day would end with light.

\textbf{Achilles:} And the supper bell on the corner --- is that a
promise?

\textbf{Tortoise:} Supper is later. The bell stays where it is till it's
wanted.

\emph{{[}The kettle comes to the boil over the firebox, and the Tortoise
makes the tea in a brown pot and fills both cups. Achilles stands with
his back to the stove and his cup in both hands until the feeling comes
back into his toes. The Tortoise keeps her shawl about her shoulders; at
the dresser end of the room it is still cold enough to see one's
breath.{]}}

\begin{center}\rule{0.5\linewidth}{0.5pt}\end{center}

\hypertarget{first-challenge-same-setting-same-answer}{%
\section{\texorpdfstring{First Challenge --- \emph{Same Setting, Same
Answer}}{First Challenge --- Same Setting, Same Answer}}\label{first-challenge-same-setting-same-answer}}

\textbf{Achilles:} \emph{{[}sitting, and breaking the seal{]}} Then the
letter first. She writes a round hand. \emph{{[}reading{]}} ``Dear Mr
Crab --- Bob and I have given up sending stones by the carrier; a season
is too long to wait for a word. The carrier has laid a glass line
between our towns instead, and we send glints down it, one at a time,
from a lamp of mine to a spar-key of Bob's.''

\textbf{Achilles:} \emph{{[}reading on{]}} ``What we want is a key: a
string of yeses and noes that only the two of us know, to lock our
letters with. The letters themselves go by the ordinary post, like this
one, which anyone may read but nobody can alter without its showing. I
want a lock made of light.''

\textbf{Achilles:} \emph{{[}looking up{]}} A lock made of light. What is
a glint, exactly?

\textbf{Tortoise:} The least light there is: a single spark, too small
to be split. \emph{{[}She lays a hand beside the lantern in the middle
of the tray.{]}} This is a glint-lamp. My great-uncle kept a light on
the north coast, and made it to signal to the next light along, one
glint at a time.

\textbf{Tortoise:} One press of this lever lets out a twin glint: one
glint from this port and one from that, of one piece, like Mr Crab's
pairs of stones. Nobody sees a glint in flight; it shows only where it
lands.

\textbf{Achilles:} And the tray? It's like his slab.

\textbf{Tortoise:} It's my grandmother's tea-tray, and tonight it does
the same work. Whatever stands on it stands again underneath, in the
lacquer; a glint has its mirror side, as his stones had, and the tray
keeps both in view.

\textbf{Tortoise:} \emph{{[}a claw along the glass{]}} These rods carry
a glint along them and out at the end, as a speaking-tube carries a
voice. This one runs to your end. That one runs to the brass sleeve,
which holds its end against the next rod's end, so that a glint passes
straight through as if the sleeve were not there, and on to my end.

\textbf{Tortoise:} And at each end, a spar-key: a crystal of Iceland
spar in a brass collar. Behind the crystal are two small frosted
windows, yes above and no below, facing whoever sits behind the key. A
glint that comes through lands in exactly one window, which flashes
once; and then the glint is spent.

\textbf{Achilles:} \emph{{[}bending to the collar at the left end{]}}
The collar turns.

\textbf{Tortoise:} It chooses the question. Turn it till you feel it
stop: at the upright stroke it asks the upright question, and at the
cross it asks the crosswise one. The stops make no sound, so a hand
cupped over the collar keeps its setting to itself.

\textbf{Achilles:} There's another mark past the cross --- a little
curl.

\textbf{Tortoise:} Another stop. We shan't want it yet. Sit at that end:
that key is yours, and its windows face you. Mine face me.

\emph{{[}Achilles carries his chair and his cup round to the stove end
and sits, and the stove's warmth finds his back at once. He sets the cup
on the corner of the stove to keep it hot. The Tortoise draws her own
chair up to the dresser end, under the room lamp, and settles into it
with the shawl pulled round her.{]}}

\textbf{Tortoise:} \emph{{[}drawing the draughts box to the middle of
the cloth{]}} The draughtsmen keep the score, and they stay off the
tray, where nothing is doubled. When your window flashes, lay a white
piece for yes or a black for no, next in your row along the table's edge
in front of your key.

\textbf{Tortoise:} Both upright, to begin. \emph{{[}She feels that her
collar stands at the stroke; Achilles, a hand on his, finds it there
too.{]}} Call it when you see it, both together.

\emph{{[}The tray stands as it is: the lamp shut and still, the rods
clear and empty from end to end, each with its reflection lying beneath
it in the black, and at either end a pair of frosted windows, dark. The
Tortoise reaches along the tray and presses the lever once. Nothing is
seen to leave the lamp, and nothing is seen to run along any rod; the
glass stays as clear as it was. In the same instant, at both ends of the
tray, one window flashes --- once --- and is dark again. The rods, the
lamp, the collars and all their reflections are exactly as they were.
What is different is only that two glints are gone, one spent at each
end, and that each of the two readers has seen a window flash that the
other has not.{]}}

\textbf{Achilles:} Yes.

\textbf{Tortoise:} Yes.

\emph{{[}The edge of the table in front of Achilles's key is bare cloth.
He takes a white piece from the box and sets it down there, off the
tray, close under his key. His window is already dark and the glint is
gone; the white piece stays where he put it, and nothing appears beneath
it, for it lies off the black. At her end the Tortoise lays a white
piece in front of her own key. Where a moment ago there was a flash that
only one of them could see, there is now a piece that anyone in the room
can see, and it will still be there in the morning.{]}}

\textbf{Tortoise:} Again. \emph{{[}The lever. A window flashes at each
end.{]}}

\textbf{Achilles:} No.

\textbf{Tortoise:} No.

\emph{{[}Two black pieces go down, one in each row, beside the
whites.{]}}

\textbf{Tortoise:} Again. \emph{{[}The lever; two flashes.{]}}

\textbf{Achilles:} Yes.

\textbf{Tortoise:} Yes.

\emph{{[}Two white pieces.{]}}

\textbf{Achilles:} In unison, every time. Well --- of course. They're
twins, cut of one piece. Each carries its yes or no out of the lamp, and
the key simply reads what's in it.

\emph{{[}Down at the bottom corner of the lower pane the frost is busy.
The first feathers have put out fronds, and the fronds fronds of their
own, one laid over another and crowding up the corner every which way,
faster than an eye can follow any one of them.{]}}

\textbf{Tortoise:} Then turn yours to the cross. I'll do the same with
mine.

\emph{{[}Each collar turns on to its other stop, and holds. The lever;
two flashes.{]}}

\textbf{Achilles:} No.

\textbf{Tortoise:} No.

\emph{{[}Two black pieces. The lever again.{]}}

\textbf{Achilles:} Yes.

\textbf{Tortoise:} Yes.

\emph{{[}Two white pieces.{]}}

\textbf{Achilles:} Crosswise too. So each glint carries two answers out
of the lamp, an upright one and a crosswise one, and its twin carries
the same two. And the two can't be strangers: a glint that is an upright
yes is a yes, and asked the crosswise question it ought still to lean
towards yes. Keep yours at the cross and I'll go back to the stroke, and
I say our calls will still agree more often than not --- three times in
four, I'd guess.

\textbf{Tortoise:} Three times in four. Let's hear it.

\emph{{[}Achilles turns his collar back to the upright stroke; the
Tortoise leaves hers at the cross. The lever.{]}}

\textbf{Achilles:} Yes.

\textbf{Tortoise:} No.

\emph{{[}He lays white; she lays black. The lever.{]}}

\textbf{Achilles:} Yes.

\textbf{Tortoise:} Yes.

\emph{{[}White, and white. The lever.{]}}

\textbf{Achilles:} No.

\textbf{Tortoise:} Yes.

\emph{{[}He lays black; she lays white.{]}}

\textbf{Achilles:} Once in three. But three rounds are nothing; toss a
coin three times and you can get anything.

\textbf{Tortoise:} Then ask someone who has tossed it more often. What
does Alice's lock-book say?

\textbf{Achilles:} \emph{{[}turning a page{]}} Here. ``I keep a
lock-book, a line for every glint I send, and Bob keeps one for every
glint he reads. These first pages are from our trials, before the
carrier routed the line through its relay house, when we compared
everything by post to learn how the line behaved.'' \emph{{[}reading
down{]}} ``Two hundred glints. Bob read ninety-six of them at my
setting, and every one came out as I sent it. The other hundred and four
he read at the other setting: yes on fifty-three and no on fifty-one,
about equally, near enough, and it came out the same whether I had sent
a yes or a no. I have added it all twice.''

\emph{{[}He reads it with his cup in his other hand, and forgets to
drink. The Tortoise leans over with the pot and tops the cup up to the
brim all the same, and he sets it back on the stove's corner without
looking.{]}}

\textbf{Achilles:} Then an upright yes, asked crosswise, doesn't lean at
all. Not a quarter of the way.

\textbf{Tortoise:} Let Bob draw circles, then, since Alice has done the
sums. \emph{{[}She moves to her drawing-cloth.{]}} These are my rings,
as at Mr Crab's. A green ring asks the upright question and a red ring
the crosswise one; each gathers a glint's cord and its mirror cord
together, with a short cord to one side, where the answer goes out. Run
the short cord in instead of out, and the same ring sends where it used
to read.

\emph{{[}On the cloth she lays a green ring at the left, its short cord
coming in from the cloth's left edge: Alice, sending upright. At the
right she lays another green ring, its short cord running out to the
right edge: Bob, reading upright. Between the two rings run two cords
side by side, the glint's and its mirror. She takes a ring in each claw
and slides them along the two cords towards each other. The cords
between them shorten and are gone; green meets green, and the two rings
run into one. A ring with one cord in and one cord out, and nothing else
on it, is no more than a joint, and she lifts it off. What lies on the
cloth now is a single cord, from the left edge to the right: Alice's
short cord running straight on into Bob's. The rings and the glint's
cords have gone; nothing has been added.{]}}

\textbf{Achilles:} Straight through. What she sends is what he reads.

\emph{{[}She lays the rings out again as before, but puts a red ring at
Bob's end: Alice sends upright, Bob reads crosswise. Green at the left,
red at the right, the glint's cord and its mirror running between them.
She draws the two rings apart. The two cords between them go slack
together, slip loose at both ends at once, and fall away onto the cloth.
Alice's short cord now ends in a green ring tied to nothing; Bob's short
cord ends in a red ring tied to nothing. What has changed: the two ends
no longer touch at all. What has held: each end still has its ring and
its short cord, and Bob still has an answer to read out --- only it is
joined to nothing Alice put in.{]}}

\textbf{Achilles:} Tied to nothing she put in. \emph{{[}He looks from
the cloth to his row, and to hers.{]}} \ldots Oh. There's no crosswise
answer in an upright glint, waiting to be read out. Asked crosswise, it
gives one there and then, and nothing in it comes from what she sent.

\emph{{[}At his elbow a drop of tea runs down the outside of his cup,
where it stands warming on the stove's corner, and falls onto the hot
iron. It does not hiss. It skates there, a bright bead, off one way and
then another, and is gone.{]}}

\textbf{Tortoise:} Then the rounds we asked different questions have
nothing to tell us about each other. We'll put them back.

\emph{{[}The rows as they stand: in front of Achilles, white, black,
white, black, white, and after them the three mixed rounds, white,
white, black. In front of the Tortoise, white, black, white, black,
white, and after them black, white, white. For five pieces the rows are
the same colour, piece for piece; after that they go their own ways. She
takes the first mixed round's piece from her row, and Achilles, at her
nod, the same round's piece from his, and they drop them back into the
box; the other two mixed rounds go the same way, a piece from each row
for each round. Nothing else on the table moves. Each row is shorter
now, and ends where the matching ended: white, black, white, black,
white at his end, and the same at hers, every piece still opposite its
own round's partner.{]}}

\textbf{Achilles:} So a key needs the settings to match. But if Alice
and Bob agree their settings by post beforehand, anyone who reads the
post can turn their own key to match, and read the line as well as Bob.
So how are they ever to match?

\begin{center}\rule{0.5\linewidth}{0.5pt}\end{center}

\hypertarget{second-challenge-say-only-the-setting}{%
\section{\texorpdfstring{Second Challenge --- \emph{Say Only the
Setting}}{Second Challenge --- Say Only the Setting}}\label{second-challenge-say-only-the-setting}}

\textbf{Tortoise:} Then let them say nothing beforehand. Choose in
secret, read, lay your piece where nobody sees it --- and only when the
glints are all spent, say which setting you chose, and nothing else.

\textbf{Tortoise:} \emph{{[}She lifts the lid from the draughts box and
stands it on its long edge in front of Achilles's row, a low wall
between his pieces and the rest of the table.{]}} Take a handful of each
colour behind that. \emph{{[}She draws the box itself to her own end and
sets it in front of her row, its deep side towards him.{]}}

\textbf{Tortoise:} Cup your hand over your collar and turn it by feel,
to whichever of its stops you like --- upright or cross. Don't tell me
which, and don't call.

\emph{{[}Achilles cups his hand over his collar and turns it; somewhere
under his palm it stops, without a sound, and not even he could have
said which stop from the feel of the room. At her end the Tortoise does
the same. The lever. A window flashes at each end, and neither says a
word. Achilles lays a piece behind his wall; the Tortoise lays one
behind the box. The rows each grow by one, but where in the first rounds
two calls hung in the air, there is nothing now, and nobody at the table
knows both colours: only the small sound of two draughtsmen set
down.{]}}

\emph{{[}Five times more: a hand over each collar, a turn felt and not
heard, the lever, a flash at each end, two pieces laid in silence behind
their screens.{]}}

\emph{{[}In the hush the room's own small sounds come forward: the stove
ticking as its iron swells, a coal settling in the firebox. On the lower
pane the frost has stopped crowding and found its figure. Every frond
now throws its side-shoots at the same few angles all along its length,
row on row, as if it had been ruled.{]}}

\textbf{Achilles:} \emph{{[}low{]}} It's very quiet.

\textbf{Tortoise:} Keys are quiet things. Now the settings, round by
round, and only the settings.

\textbf{Achilles:} First round: upright.

\textbf{Tortoise:} Upright. We keep it.

\textbf{Achilles:} Second: crosswise.

\textbf{Tortoise:} Upright. Strike it.

\emph{{[}Each reaches behind a screen, takes that round's piece from the
new end of his own row, and drops it into the box. The pieces after it
slide up to close the gap, so that at both ends every piece still stands
opposite its own round. What has gone is a round asked two different
questions; what stays is unseen by the other, as before.{]}}

\textbf{Achilles:} \emph{{[}stopping{]}} Wait --- we're saying the
settings aloud, the very thing I said a listener wanted. Somebody at the
keyhole ---

\textbf{Tortoise:} Then be the somebody. You heard me say ``upright''
for the first round, and you said it too. What do you know of my piece?

\textbf{Achilles:} Anyone with an ear to the keyhole has heard
``upright, upright'' and learned that we kept the round --- and not a
whisper of which piece we laid.

\emph{{[}Out in the lane a step goes by on the rime, crunching,
unhurried, and is gone towards the corner. Whose it was, nobody in the
house could say.{]}}

\textbf{Tortoise:} And the struck rounds he's welcome to; they're back
in the box. Third.

\textbf{Achilles:} Crosswise.

\textbf{Tortoise:} Crosswise. Keep.

\textbf{Achilles:} Fourth: upright.

\textbf{Tortoise:} Crosswise. Strike.

\emph{{[}Each takes that round's piece from behind the screen, drops it
into the box, and closes the gap.{]}}

\textbf{Achilles:} Fifth: upright.

\textbf{Tortoise:} Upright. Keep.

\textbf{Achilles:} Sixth: crosswise.

\textbf{Tortoise:} Upright. Strike.

\emph{{[}The last pieces of the run go back in the box.{]}}

\textbf{Tortoise:} Now look --- once, for the demonstration, and never
again. A key that has been shown is a key no longer. \emph{{[}She takes
down the lid and turns the box aside.{]}}

\emph{{[}In front of Achilles lie the five pieces from the first rounds
of the evening, and after them the kept pieces of the silent run: black,
white, black. In front of the Tortoise lie the same five, and after them
black, white, black. Neither saw the other's pieces laid; laid open now,
the two rows match piece for piece, right to the end, each piece
opposite its partner.{]}}

\textbf{Achilles:} The same, right to the end. And I didn't choose a
single one of them. Nor did you. The lamp did.

\emph{{[}The Tortoise laughs under her breath and sits back in her
chair. Then she gets up, fills the kettle again from the jug, and sets
it at the stove's cool edge, well back from the fire, where it will take
its own time. Achilles stretches his stockinged feet towards the stove
and wriggles his toes at it.{]}}

\textbf{Achilles:} \emph{{[}finding his place in the letter{]}} She says
so too. ``Of our two hundred trial glints, ninety-six had matching
settings --- about half, near enough --- and not one of those disagreed.
So each glint I send is half a letter of key, and I never choose a
letter of it myself.''

\textbf{Tortoise:} Let Bob draw only the rounds whose settings matched.

\emph{{[}On the cloth she lays two rounds, one above the other: above, a
green sending ring and a green reading ring with the glint's cord and
its mirror cord between them; below, the same in red. She slides each
pair of rings together, as before; each pair runs into one ring and
comes off, and each round is left as a single cord from the left edge to
the right --- upright above, crosswise below, and no difference to be
seen between the two cords. The rounds that mixed red with green she
does not lay at all.{]}}

\textbf{Achilles:} Two straight cords. A listener knows which colour
each one was, and nothing of what went along it.

\textbf{Achilles:} \emph{{[}turning the page{]}} There's more.
\emph{{[}reading{]}} ``But there is one thing you should know. The
carrier has since routed our line through its relay house, and the clerk
there is Eve --- the one with the press.''

\begin{center}\rule{0.5\linewidth}{0.5pt}\end{center}

\hypertarget{third-challenge-the-tell-glass}{%
\section{\texorpdfstring{Third Challenge --- \emph{The
Tell-Glass}}{Third Challenge --- The Tell-Glass}}\label{third-challenge-the-tell-glass}}

\textbf{Achilles:} Eve! The carrier's clerk from Mr Crab's letter this
afternoon --- the one who kept the struck stones in her drawer.

\textbf{Tortoise:} She has moved with the times. \emph{{[}From a small
wooden box by her chair she takes a block of glass the size of a sugar
lump, with a port at either end and another in its side; a short glass
rod; and a spar-key in its collar, like the others.{]}} This is a
tell-glass: what Eve would put on the line. A glint goes in at this end
and on out at that one, and at the same moment a second glint leaves by
this side port, carrying the passing glint's upright answer.

\textbf{Achilles:} Mr Crab's press --- in glass.

\textbf{Tortoise:} In glass, and quieter; nothing drops out of a chute.
\emph{{[}She sets the spar-key on the tray beside her own at the right
end, its windows turned towards her, and its reflection stands beneath
it in the black.{]}} And this is the tell's key. It stands at my end
with its windows facing me, so that only I see what it reads; the short
rod will carry the side glint to it.

\emph{{[}The right-hand rod and the rod that runs on to her key meet end
to end inside the brass sleeve: one line of glass, with one line beneath
it in the black. She loosens the sleeve and draws it off. The two rod
ends stand apart now, a finger's gap between them, each still on its
fork; nothing else on the tray stirs. Into the gap she sets the
tell-glass, the right-hand rod's end at its near port, the far rod's end
at its through-port, its side port turned towards her. The moment it
rests on its fork, a second tell-glass is standing beneath it, upside
down in the black, port to port. Where one line of glass ran through a
sleeve, there are now two rods meeting at a block of glass, and the
block has its reflection; the way from the lamp to her key is still open
end to end.{]}}

\emph{{[}From the side port she lays the short rod on its forks to the
tell's key, and turns that key's collar till it stops at the upright
stroke. Beneath the new rod, in the lacquer, its reflection lies.{]}}

\emph{{[}For a moment the room makes no sound at all: the stove has left
off its ticking, and nothing stirs at the kettle's spout. The room
lamp's flame stands perfectly still on the dresser, and the print on the
end wall hangs in its light as though it had only now been hung
there.{]}}

\emph{{[}Achilles, looking down at the tell-glass standing upside down
in the black, lifts his eyes from the tray to the print on the wall:
Escher's \textup{Three Worlds}.{]}}

\textbf{Achilles:} A pond: leaves on its surface, trees reflected in it,
a fish beneath. I have drunk a dozen pots of tea under that print and
never once looked past the leaves.

\textbf{Tortoise:} Then look now. Secret settings, hands over the
collars, as before --- but this time call your windows aloud, for the
demonstration. The tell's pieces I'll lay myself, and say nothing about
them.

\emph{{[}Achilles blows on his fingers before he cups them. The short
rod the Tortoise has just laid is newer glass than the rest, clear as
water, with none of the old rods' faint colour in it, and his eye keeps
going back to it.{]}}

\emph{{[}Hands over collars. The tell's key stands at the upright
stroke, its windows dark; the side rod is clear and empty, as it has
been since she laid it. The Tortoise presses the lever. A window flashes
at Achilles's end; a window flashes at her own key; and in the same
instant a window flashes at the tell's key as well, where no glint has
ever landed before. The lamp let out no more than it always does. The
rods, the tell-glass and every reflection are as they were. What is new
is one more flash at her end, from a glint the tell-glass sent off
sideways --- and the glint that went on to her own key arrived all the
same.{]}}

\textbf{Achilles:} Yes.

\textbf{Tortoise:} Yes.

\emph{{[}Two white pieces go down in the rows. Then the Tortoise takes a
white piece and, without a word, lays it on the table's edge in front of
the tell's key, beside her own row: the first piece of a new row.{]}}

\textbf{Tortoise:} Again. \emph{{[}Hands over collars; the lever; a
flash at each end, and one more at hers.{]}}

\textbf{Achilles:} No.

\textbf{Tortoise:} Yes.

\emph{{[}Black for him, white for her; and in silence, a white for the
tell.{]}}

\textbf{Tortoise:} Again.

\textbf{Achilles:} Yes.

\textbf{Tortoise:} No.

\emph{{[}White for him, black for her; for the tell, in silence,
black.{]}}

\textbf{Tortoise:} Again.

\textbf{Achilles:} Yes.

\textbf{Tortoise:} Yes.

\emph{{[}Two whites; in silence, a white for the tell.{]}}

\textbf{Tortoise:} Again.

\textbf{Achilles:} No.

\textbf{Tortoise:} No.

\emph{{[}Two blacks; in silence, a black for the tell.{]}}

\textbf{Tortoise:} Once more.

\textbf{Achilles:} No.

\textbf{Tortoise:} No.

\emph{{[}Two blacks; and for the tell, in silence, white.{]}}

\textbf{Tortoise:} Now the settings. Yours first, all of them.

\textbf{Achilles:} Upright, crosswise, crosswise, upright, upright,
crosswise.

\textbf{Tortoise:} Upright, upright, crosswise, crosswise, upright,
crosswise. The second round goes, and the fourth.

\emph{{[}From every row --- his, hers and the tell's --- the second
round's piece and the fourth round's go back into the box, and the gaps
are closed. The rows stand aligned as before, round opposite round, only
shorter: what is left of the run is the first and fifth rounds, asked
upright at both ends, and the third and sixth, asked crosswise at both
ends.{]}}

\textbf{Achilles:} The upright rounds, the first and the fifth: we
agreed both times, and the tell's pieces agree with us --- white, then
black. She has those, whole. \emph{{[}His finger moves along the
rows.{]}} The crosswise ones. In the third, I said yes and you said no
--- at the same setting. That hasn't happened all evening.

\textbf{Tortoise:} Not on an empty line.

\emph{{[}On the lower pane, the round that Achilles's fingertip melted
at the start of the evening has frozen over again, but not as the rest
has. The ferns that climbed up round it bend away at its edge, and
across the patch itself runs a single blunt frond, thick and crooked,
like nothing else on the glass.{]}}

\textbf{Achilles:} And in the sixth we agreed, both no, and the tell's
piece is white. On the crosswise rounds her pieces are no guide at all.

\textbf{Achilles:} \emph{{[}taking up the letter again{]}} ``Then the
carrier routed the line through its relay house, and the lock-book
changed. We had begun by then to check some of our kept rounds openly,
by post, before we used a key. Now some of the checks disagreed: every
one of them on a crosswise round, never an upright one --- about one
checked round in four, near enough. So we knew.''

\textbf{Achilles:} Then a cleverer Eve would have a glass that copies
the crosswise answer instead.

\textbf{Tortoise:} On which rounds? She hears the settings only when the
glints are long spent. Copy crosswise, and it's the upright rounds she
spoils; she can only guess, and every wrong guess spoils a round,
whichever way it goes.

\textbf{Achilles:} And the upright rounds came through untouched because
the tell-glass copies its own kind, as the press did: an upright glint
it copies and lets be, and a crosswise one it can only make a pair of.

\textbf{Tortoise:} Let Bob draw Eve where she sits, then. A forked ring
is the tell-glass on my cloth: a green ring with one cord in and two
out, one going on to Bob and one turning off to her.

\emph{{[}On the cloth she lays Alice's round again: a green sending ring
at the left, a red reading ring at Bob's end, the glint's cord and its
mirror cord running between them. Then she takes a green ring made to
fork, one cord in and two out, and threads it onto the glint's cord
partway along; its reflection, another fork, she threads onto the mirror
cord just opposite. From each fork one out-cord runs on to Bob's ring as
before, and the other turns off sideways, both side cords going together
to another ring, laid off to one side: Eve's. What has changed: Bob's
line now has a side road off it, and at its end a ring of its own. What
has held: Alice's ring, Bob's ring, and the line from one to the other,
which still runs through unbroken.{]}}

\textbf{Achilles:} A side road off his line, and Alice never sees it
from her end.

\emph{{[}The Tortoise rolls her shoulders under the shawl and flexes her
claws, stiff from the fine work. At the stove's cool edge the kettle has
begun a low hum, too low yet to call a note, and neither of them turns
to it.{]}}

\emph{{[}The Tortoise lifts the tell-glass from the tray, and its
reflection goes out of the black with it; she lays it back in her wooden
box. She slides the sleeve over the two rod ends again, and one line of
glass runs once more from the lamp to her key, with one line beneath it.
The side rod stays where it lies, its near end at nothing now, its
reflection beneath it; the tell's key stands idle beside her own, still
at the upright stroke. The tell's row stays on the cloth.{]}}

\textbf{Achilles:} But wait. \emph{{[}looking along the tray{]}} This
isn't Alice's way at all. Here the lamp sits in the middle and throws a
pair, one glint to each end; Alice sends from her end, and Bob reads at
his. Why should they come out the same?

\begin{center}\rule{0.5\linewidth}{0.5pt}\end{center}

\hypertarget{fourth-challenge-backwards-to-the-lamp}{%
\section{\texorpdfstring{Fourth Challenge --- \emph{Backwards to the
Lamp}}{Fourth Challenge --- Backwards to the Lamp}}\label{fourth-challenge-backwards-to-the-lamp}}

\textbf{Tortoise:} Let Bob draw clean through the lamp, and see. A bend
is how I draw the lamp's twin: one cord turned round on itself, a pair
cut of one piece, with an end at each key.

\emph{{[}On the cloth she lays the table as it stands tonight. In the
middle, where the lamp is, a bend: the glint's cord and its mirror cord,
each turned round on itself, so that both run out to the left and both
run out to the right. At the left, on the two cords, she threads a green
reading ring, Achilles's short cord running out of it to the left edge;
at the right, a green reading ring with her own short cord running out
to the right edge. Then she takes Achilles's short cord in one claw and
her own in the other, and pulls. Achilles's ring slides along its cords,
round the turn of the bend where the lamp stood, and comes out on the
other side, heading right; the bend runs out and is gone. What lies on
the cloth is one pair of cords from left to right, with a green ring at
each end: at the left, a ring whose short cord now comes in from
Achilles's side; at the right, a ring whose short cord goes out to hers.
What has changed: the lamp's bend is gone, and Achilles's ring has
turned from a reading ring into a sending one, his cord coming in where
it used to go out. What has held: every ring, every cord, and where each
short cord ends.{]}}

\textbf{Achilles:} A sending ring and a reading ring on one cord --- the
very picture from the start of the evening, Alice's lamp and Bob's key.
My reading at my end of the twin is her sending.

\textbf{Tortoise:} Your reading goes back into the lamp, and on out to
me.

\textbf{Achilles:} Then that's settled: same setting, same answer,
whichever way it's done.

\emph{{[}Achilles leans back, well satisfied, and reaches behind him to
spread his hands to the stove. The Tortoise watches him do it and says
nothing at all. On the window a second fern has started from the other
bottom corner of the lower pane, and is climbing towards the first.{]}}

\textbf{Tortoise:} Upright, both, first, and let's hear it. \emph{{[}Two
collars turn back from the cross to the upright stroke. The lever; a
flash at each end.{]}}

\textbf{Achilles:} Yes.

\textbf{Tortoise:} Yes.

\textbf{Tortoise:} And crosswise. \emph{{[}Two collars go to the cross.
The lever; a flash at each end.{]}}

\textbf{Achilles:} No.

\textbf{Tortoise:} No.

\textbf{Tortoise:} And now the stop you asked about. Past the cross, to
the curl. I'll turn mine to the curl as well.

\emph{{[}Each turns a collar on past the cross until it stops at the
little curl. The lever.{]}}

\textbf{Achilles:} Yes.

\textbf{Tortoise:} No.

\emph{{[}The lever.{]}}

\textbf{Achilles:} No.

\textbf{Tortoise:} Yes.

\emph{{[}The lever.{]}}

\textbf{Achilles:} No.

\textbf{Tortoise:} Yes.

\textbf{Achilles:} Every one wrong! The same setting --- the very same
curl at both ends --- and not once alike. \emph{{[}He looks along the
tray to the sleeve.{]}} The tell-glass is in its box. Is it the lamp? Or
has somebody put something on the line?

\textbf{Tortoise:} Nothing is on the line. Look at your collar in the
tray.

\emph{{[}Achilles bends down to the lacquer. His collar stands at the
curl, the little mark facing him and turning, as he looks at it straight
on, the way the sun goes round. In the black beneath, the collar's
reflection shows its curl turning the other way. He turns the collar
back to the cross: the cross in the tray is a cross, the same as the one
above. On to the stroke: the reflected stroke is just the stroke. He
turns it to the curl again, and the curl beneath turns against the curl
above. The stroke and the cross hold still in the mirror; only the curl
changes hands.{]}}

\textbf{Tortoise:} Round the bend, your question reaches my end as its
reflection. Upright and crosswise look the same in the lacquer, so
nobody notices. The curl does not: at the curl, my yes-window is your
no-window's reflection.

\textbf{Achilles:} \ldots Oh. Then if you'd read your windows the other
way about, we'd have agreed every time.

\emph{{[}On the fender his gloves lie where he dropped them when he came
in, pulled off by the fingertips and still inside out. The left one is,
to any eye, a right glove now. His woollen cap, inside out on top of
them, is only a cap.{]}}

\textbf{Tortoise:} Every time. A glint sent at the curl and read at the
curl comes through as it was sent; it's only the bend that hands on a
reflection. \emph{{[}She turns her collar back to the upright stroke;
Achilles turns his.{]}} Now suppose the calls agreed at upright and at
crosswise, every time --- no wrong note, ever. What could anyone else
have of such a pair?

\begin{center}\rule{0.5\linewidth}{0.5pt}\end{center}

\hypertarget{fifth-challenge-the-alarm}{%
\section{\texorpdfstring{Fifth Challenge --- \emph{The
Alarm}}{Fifth Challenge --- The Alarm}}\label{fifth-challenge-the-alarm}}

\textbf{Achilles:} I don't know. This evening Eve had half of every key
she touched.

\textbf{Tortoise:} This evening Eve had a glass on the line. Let's put
it back, and give her every advantage we can think of.

\emph{{[}At the stove's cool edge the kettle has found its note at last:
a thin, steady singing from the spout, which goes on under everything.
Neither of them turns round.{]}}

\emph{{[}She draws the sleeve off again and sets the tell-glass in its
place, the right-hand rod at its near port, the far rod at its
through-port, the side rod at its side port; its reflection stands
beneath it at once. The line to her key is whole as before, with the
tell-glass on it and the side road open.{]}}

\textbf{Tortoise:} This time we agree the setting aloud beforehand, and
I set the tell's key to the same stop: the strongest Eve there is, who
knows the question before she reads.

\textbf{Achilles:} Crosswise, then --- all of them.

\emph{{[}His collar and hers turn to the cross, and so does the
tell's.{]}}

\textbf{Tortoise:} \emph{{[}a claw on the bell at the table's corner{]}}
Now the bell. This is a checked round: we both call aloud, and if our
calls differ I ring it once, and if they agree it stays still. That is
all it does tonight.

\emph{{[}The lever. A flash at his end; two at hers.{]}}

\textbf{Achilles:} Yes.

\textbf{Tortoise:} No.

\emph{{[}The bell stands on the table's corner, silent, as it has stood
all evening. The Tortoise lifts it by its handle and rings it once; the
note goes round the room and fades. She sets it back where it stood.
Nothing on the tray has moved, nothing in the black; the lamp, the glass
and the rods are as they were. What has changed is only who knows: the
difference between the two calls, which a moment ago was a thing between
the two of them, is now a sound that anyone in the house could have
heard.{]}}

\emph{{[}He lays white; she lays black; and then, without a word, she
lays the tell's piece: black.{]}}

\emph{{[}Across the lane a dog, woken by the bell, barks twice and is
quiet; somewhere a door opens on the cold and shuts again. At the
stove's edge the kettle sings on, unheeded.{]}}

\textbf{Achilles:} And the tell's piece is black.

\textbf{Tortoise:} Another, crosswise still. \emph{{[}The lever.{]}}

\textbf{Achilles:} No.

\textbf{Tortoise:} No.

\emph{{[}The bell stays where it is, and does not sound. Two blacks go
down; and in silence, a white for the tell.{]}}

\textbf{Achilles:} No bell --- and the tell's piece is white.

\textbf{Tortoise:} At crosswise it always goes so: her piece is black
whenever the bell rings, and white whenever it keeps still.

\textbf{Achilles:} Then all she holds is whether she was caught --- and
nothing of the key.

\textbf{Tortoise:} Now upright, all of them, agreed aloud. \emph{{[}The
collars go back to the stroke, the tell's with them. The lever.{]}}

\textbf{Achilles:} No.

\textbf{Tortoise:} No.

\emph{{[}The bell is still. Two blacks; and in silence, a black for the
tell, the same as his.{]}}

\textbf{Achilles:} Silent again, and black like mine. At upright she has
the key, and the bell never hears her.

\textbf{Tortoise:} Which is why Alice and Bob check at both settings,
and say which only afterwards.

\emph{{[}Achilles takes up his cup from the stove's corner, finds it
stewed and bitter from standing, and drinks it anyway. The Tortoise
draws the shawl closer; the dresser end of the room has never quite
warmed.{]}}

\textbf{Achilles:} \emph{{[}turning to the last page{]}} She comes to
that. \emph{{[}reading{]}} ``Now we check about one kept round in three,
by post, and strike the checked ones from the key. If too many disagree,
we throw the whole key away and begin again. Bob, who draws, has found a
way to mend a key with only a few disagreements in it, and to shorten
it, so that what the clerk knows of it comes to next to nothing. I
cannot follow his proof yet. He says it is all in the drawing.''

\textbf{Tortoise:} Let Bob draw keys and checks apart, then.
\emph{{[}From her pouch she takes two small rings, one red and one
green, each with a bead set exactly halfway round.{]}} These are
half-turn rings, like this morning's. Threaded on a pair's bend, a red
one makes the pair disagree whenever both ends are asked upright, and a
green one whenever both are asked crosswise.

\emph{{[}On the cloth she lays a bend of string, one cord turned round
on itself, both its ends running out to the right: a pair, before anyone
reads it. On its upper leg she threads the red half-turn ring, and after
it the green, each bead sitting halfway round.{]}}

\textbf{Tortoise:} An upright check that never disagrees ---
\emph{{[}she slides the red ring along the leg and off its end{]}} ---
and a crosswise check that never disagrees. \emph{{[}The green ring
follows it off.{]}}

\emph{{[}The bend lies bare on the cloth. Before, two rings rode on its
upper leg, each a way the pair could go wrong; now there are none. The
bend itself has not moved: the same turn, the same two ends, one running
towards Achilles's side of the cloth and one towards hers.{]}}

\textbf{Tortoise:} Let Bob draw empty paper between the pair and anyone
else.

\emph{{[}She runs the edge of her claw down the cloth beside the bare
bend, from the top of the cloth to the bottom, between the bend and
everything else on the cloth. It crosses nothing: no cord, no ring. The
bend lies on one side of the line her claw made, whole, with both its
ends taken; on the other side there is only cloth.{]}}

\textbf{Tortoise:} Not an end left over for anyone.

\emph{{[}On the dresser shelf, below the room lamp, lies the
great-uncle's horseshoe magnet, its iron keeper laid across both its
ends, as it has lain for years. A steel thimble has rolled against the
keeper and lies there, held by nothing.{]}}

\textbf{Achilles:} Then that's the whole of it. Two calls at the same
setting sing the same; a listener has to look; a look at the wrong
setting is a wrong note, and the bell hears it; and a pair that never
sings a wrong note, upright or crosswise, has nothing to spare for
anyone else.

\emph{{[}On the lower pane the two ferns, climbing all evening from
their two corners, have met near the middle. Where they meet, each frond
runs into a gap in the other's, and neither grows on past the other.{]}}

\textbf{Tortoise:} The checks catch a spy. That a key with a few wrong
notes in it can still be made quite safe is Bob's proof, and not
tonight's.

\emph{{[}She lifts the tell-glass out of its place for the last time and
lays it in its box; its reflection goes out of the black with it. She
slides the sleeve back over the two rod ends, and shuts the lamp.{]}}

\textbf{Tortoise:} Leave the lamp to cool; it's warmer than it looks.
And the kettle has been singing to itself for some time --- did you hear
it?

\emph{{[}She lifts the kettle off the stove's edge with a fold of her
shawl round the handle, and the singing stops in the middle of its note.
She makes a second pot, and pours.{]}}

\textbf{Achilles:} I heard nothing all evening but yes and no, and a
bell. \emph{{[}He takes the cup she hands him, and goes with it to the
window.{]}} The ferns have reached the middle of the pane.

\textbf{Tortoise:} They'll have the whole of it by morning.

\textbf{Achilles:} Two people in different towns, with nothing between
them but a line anyone might tap, and the post. Could they really lock a
letter with nothing but light --- and know, afterwards, whether anyone
had laid a finger on it?

\textbf{Tortoise:} Bob says so, and Alice can't follow him yet. Drink
your tea, and tell me in the morning which of them you would trust with
your own letters.

\emph{{[}Achilles laughs, and drinks, and stays at the window with the
cup warm between his hands. The ferns stand at the middle of the lower
pane, where the two of them met; above them the glass is clear and
black, and past it the lane lies pale under the rime, the curtained
windows opposite lit, and no step in it. The stove has settled to a low
red behind its mica, and its ticking has slowed to one now and then.
Down the table the rods lie on their forks as the Tortoise left them ---
the long green one to the stove end, the one with the seed of air
running to the sleeve, the pale one on from the sleeve to her end, and
the short new one lying idle by the tell's key; before Achilles's key
his row runs along the table's edge, and before hers her own row and the
tell's lie side by side. The bell stands on its corner. On the fender
his gloves are dry, and still inside out. On the dresser shelf the
horseshoe magnet keeps its keeper, and the thimble lies where it rolled.
Behind the Tortoise's empty chair the print hangs on the end wall where
it has always hung, pale in its narrow frame; from the window it is too
far to make out more than the leaves on its surface. The Tortoise sits
down by the stove with her own cup, and the room is quiet but for the
slow ticking of the iron and the frost going on with its work on the
glass, too slowly for anyone to see.{]}}

\hypertarget{chapter-12-locks-made-of-light}{%
\chapter{Chapter 12: Locks Made of
Light}\label{chapter-12-locks-made-of-light}}

\emph{--- Yes--No Duet, with Eavesdropper ---}

\emph{Part III: Pictures Become Proofs}

\emph{Escher: ``Three Worlds'' --- a pond: leaves on its surface, trees
reflected in it, a fish beneath}

\begin{center}\rule{0.5\linewidth}{0.5pt}\end{center}

Can two people who share only a line that anyone may tap agree on a
secret, and know afterwards whether anyone listened?

Every lock needs a key, and every key has to travel. Two friends who
want to write to each other in a cipher must first agree on how to
scramble and unscramble their letters, and that agreement is itself a
message. If they are in the same room they can whisper it. If they live
in different towns it has to go by some carrier: a courier who might be
bribed, a letter that might be steamed open and sealed again, a wire
that might be tapped. The trouble is not only that a spy might read the
key. It is that a careful spy leaves no trace. A letter can be copied
and sent on unharmed, and a signal on a wire can be recorded without
weakening it by a detectable amount. In the everyday world, looking
costs nothing, so the two friends can never be sure they were alone.

The last chapter changed the terms of that bargain. A qubit cannot be
copied. A spy who wants to learn something from a qubit in transit has
to look at it, and a look is made at some setting. Chapter 11 ended
here: ``Think of a spy on the line. She cannot copy what passes, so she
cannot keep a copy and let the original go by untouched. To learn
anything she must look. A look is made at some setting, plain or cross,
and this chapter has shown what a look at the wrong setting does: it
turns a crosswise stone grey, as Eve's customers found. Chapter 12 makes
those two settings into the two keys of a lock, built of light. There
the mark a spy's look leaves on the stones becomes an alarm.''

This is that chapter. Its central fact is small enough to state on a
postcard: \emph{read a qubit at the setting it was sent with, and you
get back what was sent; read it at the other setting, and you get a coin
toss that knows nothing of what was sent.} From that one fact, Bennett
and Brassard built a way to share a key in 1984, and it is still the
first thing anyone learns about quantum cryptography. We will build it
in four steps. Two people choose their settings at random and afterwards
tell each other only the settings. A spy on the line is Chapter 11's
copier, and we will measure exactly what her look costs her. The same
key can be grown from a shared pair instead of a sent qubit, and a
single bend in a wire shows the two methods are one. Finally the checks
the two make become an alarm, and a pair that passes every check turns
out to have nothing left over for anyone else.

We will also be plain about where the pictures stop. Every rewrite in
this chapter proves an equality of processes. None of them proves that
the key is secure against every possible spy. That theorem exists, and
we will state it and say whose it is, but it is not drawn here.

\begin{center}\rule{0.5\linewidth}{0.5pt}\end{center}

\hypertarget{two-settings}{%
\section{§12.1 Two Settings}\label{two-settings}}

Chapter 11 left us two ways to read a qubit. Chapter 12 turns them into
the two keys of a lock.

\textbf{On the blackboard.} The two settings are the ones Chapter 11
used. The \textbf{plain} setting asks the question whose answers are
\(\ket0\) and \(\ket1\); the \textbf{cross} setting asks the question
whose answers are \(\ket+=\tfrac1{\sqrt2}(\ket0+\ket1)\) and
\(\ket-=\tfrac1{\sqrt2}(\ket0-\ket1)\). We label a setting by a bit:
\(0\) for plain, \(1\) for cross. The Hadamard box of Chapter 3 turns
one set of answers into the other, \(H\ket0=\ket+\) and
\(H\ket1=\ket-\), so one formula covers both settings. Write \(H^0=I\)
and \(H^1=H\). Then the answer \(x\) at setting \(r\) is the state

\[H^r\ket x,\qquad x,r\in\{0,1\}:\qquad \ket0,\ \ket1\ \ (r=0),\qquad \ket+,\ \ket-\ \ (r=1).\]

Alice, who wants to send a bit \(x\), picks a setting \(r\) and sends
\(H^r\ket x\). Bob, at the other end, picks a setting \(t\) and reads. A
reading at setting \(t\) has the two outcomes \(y=0\) and \(y=1\), with
effects the mirror images of the two states \(H^t\ket y\). The mirror
image of a state is its conjugate transpose, and \(H\) is real and
symmetric, so the effect for outcome \(y\) is \(\bra yH^t\). The
amplitude for Bob to read \(y\) is therefore

\[\bra y\,H^t H^r\,\ket x ,\]

and the probability is its squared modulus (§2.3).

Start with the case that matters most. Alice sends \(x=0\) at the plain
setting, \(\ket0\), and Bob reads at the cross setting. The two
amplitudes are \(\braket{+|0}=\tfrac1{\sqrt2}\) and
\(\braket{-|0}=\tfrac1{\sqrt2}\), so Bob reads \(0\) or \(1\) with
probability \(\tfrac12\) each. Had Alice sent \(x=1\), the amplitudes
would have been \(\braket{+|1}=\tfrac1{\sqrt2}\) and
\(\braket{-|1}=-\tfrac1{\sqrt2}\), and the probabilities again
\(\tfrac12\) each. Bob's reading at the wrong setting is a fair coin,
and it is the same fair coin whatever Alice sent. A reader who believes
that the qubit ``really'' carries an answer at both settings, and that
the cross question merely reveals the cross answer it was carrying all
along, has to explain why that answer is a fair coin whose odds do not
depend on the plain answer Alice put in. The simpler account is the true
one. A qubit sent at the plain setting has no cross answer waiting to be
read.

Two settings with this property have a name. Two orthonormal bases
\(\{\ket{u_0},\ket{u_1}\}\) and \(\{\ket{v_0},\ket{v_1}\}\) of a qubit
are \textbf{conjugate} if

\[\bigl\lvert\braket{u_i|v_j}\bigr\rvert^2=\tfrac12\qquad\text{for all } i,j .\]

In Bennett and Brassard's words, each vector of one basis has
equal-length projections onto all vectors of the other; they take the
word from Wiesner.

\textbf{\emph{Why this definition.}} Why ask for equal projections, and
not merely that the two bases differ? Because equal projections are
exactly what makes a reading at the wrong setting worthless. If
\(\lvert\braket{u_i|v_j}\rvert^2\) were \(0.6\) for one pair and \(0.4\)
for another, a reading in one basis would lean toward the state sent in
the other, and a reader at the wrong setting would learn something. At
exactly \(\tfrac12\) everywhere, the reading is a fair coin whatever was
sent.

\begin{quote}
\textbf{Result 12.1 (conjugate bases).} The plain and cross bases are
conjugate: \(\lvert\bra iH\ket j\rvert^2=\tfrac12\) for all bits
\(i,j\). The circular basis of the BB84 paper,
\(\ket{c_1}=\tfrac1{\sqrt2}(1,i)\) and
\(\ket{c_2}=\tfrac1{\sqrt2}(i,1)\), is conjugate to both.
\end{quote}

\emph{Proof.} The four entries of \(H\) are
\(\bra iH\ket j=(-1)^{ij}/\sqrt2\), each of squared modulus
\(\tfrac12\). For the circular basis,
\(\braket{0|c_1}=\tfrac1{\sqrt2}\), \(\braket{1|c_1}=\tfrac i{\sqrt2}\),
\(\braket{0|c_2}=\tfrac i{\sqrt2}\) and
\(\braket{1|c_2}=\tfrac1{\sqrt2}\), all of modulus \(\tfrac1{\sqrt2}\).
Against the cross basis, \(\braket{+|c_1}=\tfrac{1+i}2\),
\(\braket{-|c_1}=\tfrac{1-i}2\), \(\braket{+|c_2}=\tfrac{1+i}2\) and
\(\braket{-|c_2}=\tfrac{i-1}2\), and each of \(1\pm i\) and \(i-1\) has
modulus \(\sqrt2\), so each overlap has modulus \(\tfrac1{\sqrt2}\).
\(\;\blacksquare\) See \texttt{Code/Ch12/TwoSettings.py}.

The protocol will use only the plain and cross settings. The circular
basis, which Bennett and Brassard mention and do not use, returns in
§12.4 as a test that tells a transpose from a conjugate.

\vskip 0pt plus 4\baselineskip\penalty-100\vskip 0pt plus
-4\baselineskip

\begin{quote}
\textbf{Dirac's Blackboard.}\nopagebreak

\emph{Why the third basis needs \(i\).} A real unit vector in the plane
is \((\cos\theta,\sin\theta)\). Its overlap with \(\ket0\) has squared
modulus \(\cos^2\theta\), and with \(\ket+\) it is
\(\tfrac12(\cos\theta+\sin\theta)^2=\tfrac12(1+\sin2\theta)\). To be
conjugate to the plain basis we need \(\cos^2\theta=\tfrac12\), so
\(\theta\) is an odd multiple of \(\tfrac\pi4\). To be conjugate to the
cross basis we need \(\sin2\theta=0\), so \(\theta\) is a multiple of
\(\tfrac\pi2\). No real \(\theta\) does both. With complex entries the
obstacle goes: \(\ket{c_1}=\tfrac1{\sqrt2}(1,i)\) has overlap
\(\tfrac12\) with everything in sight, as Result 12.1 shows. This is
what the BB84 paper means when it says that, ``owing to the complex
nature of its coefficients'', the space of a single photon admits a
third basis conjugate to the other two. Three questions, each a fair
coin to the other two, fit on one qubit only because amplitudes are
complex numbers.
\end{quote}

\textbf{In the sketchbook.} Recall from Chapter 11 (§11.1) that a
normalized plain state is a red point with \(\tfrac1{\sqrt2}\) written
in front: the red point with phase \(a\pi\) is \(\sqrt2\,\ket a\). In
the same way, a green effect with phase \(b\pi\) is
\(\bra0+e^{ib\pi}\bra1\), which is \(\sqrt2\,\bra+\) for \(b=0\) and
\(\sqrt2\,\bra-\) for \(b=1\). A plain state read with a cross effect is
the closed picture of the two joined by one wire, and a closed picture
is a number (§2.2):

\[\tfrac1{\sqrt2}\cdot\tfrac1{\sqrt2}\;
\begin{tikzpicture}[sd, x=1.4em, y=1.4em,
  zsp/.style={circle,draw,line width=0.9pt,fill=green!55,minimum size=1.1em,inner sep=0pt},
  xsp/.style={circle,draw,line width=0.9pt,fill=red!75,minimum size=1.1em,inner sep=0pt},
  ph/.style={font=\scriptsize,inner sep=1pt}]
  \draw[wire] (0,0) -- (2.2,0);
  \node[xsp,label={[ph]above:{$a\pi$}}] at (0,0) {};
  \node[zsp,label={[ph]above:{$b\pi$}}] at (2.2,0) {};
\end{tikzpicture}
\;=\;\tfrac12\cdot\sqrt2\,e^{iab\pi}\;=\;\frac{(-1)^{ab}}{\sqrt2}.\]

For all four choices of \(a\) and \(b\) the number has the same modulus,
\(\tfrac1{\sqrt2}\). Only its sign depends on the phases. That is Result
12.1 as a picture: a red point meeting a green effect gives the same
odds whichever red point and whichever green effect.

A reading, though, is more than an amplitude. It is a process that ends
in a record, and since Chapter 9 we draw such processes on doubled
wires. The \textbf{reader} at setting \(t\) is Chapter 9's read-spider
(row 68): a green dot that takes a quantum system as two lines, its wire
\(q\) and its mirror copy \(\bar q\), and emits one classical wire
carrying the outcome. At the cross setting a yellow Hadamard box sits on
each of the two quantum legs; the mirror copy of \(H\) is \(H\) again,
because \(H\) is real. At the plain setting there are no boxes. The
\textbf{sender} at setting \(r\) is the same green dot run the other
way: one classical wire comes in carrying Alice's bit, and the doubled
wire goes out, with the same yellow boxes on both quantum legs at the
cross setting. Chapter 11 met this dot once already, in §11.3, as the
green dot that ``sets a stone to whatever the record says''. Here it
gets a name and a row of its own (row 89). Drawn at the cross setting:

\[E_1\;=\;
\begin{tikzpicture}[sd, x=1.2em, y=1.2em,
  zsp/.style={circle,draw,line width=0.9pt,fill=green!55,minimum size=1.2em,inner sep=0pt},
  hbox/.style={draw,line width=0.9pt,fill=yellow!85,minimum size=0.85em,inner sep=1pt}]
  \draw[cwire] (-2,0) -- (0,0);
  \draw[wire] (0,0) -- (1,0.9) -- (4,0.9);
  \draw[wire] (0,0) -- (1,-0.9) -- (4,-0.9);
  \node[hbox] at (2.3,0.9) {};
  \node[hbox] at (2.3,-0.9) {};
  \node[zsp] at (0,0) {};
  \node[sdlabel,anchor=east] at (-2.1,0) {$x$};
  \node[sdlabel,anchor=west] at (4.1,0.9) {$q$};
  \node[sdlabel,anchor=west] at (4.1,-0.9) {$\bar q$};
\end{tikzpicture}
\qquad\qquad
R_1\;=\;
\begin{tikzpicture}[sd, x=1.2em, y=1.2em,
  zsp/.style={circle,draw,line width=0.9pt,fill=green!55,minimum size=1.2em,inner sep=0pt},
  hbox/.style={draw,line width=0.9pt,fill=yellow!85,minimum size=0.85em,inner sep=1pt}]
  \draw[wire] (0,0.9) -- (3,0.9) -- (4,0);
  \draw[wire] (0,-0.9) -- (3,-0.9) -- (4,0);
  \draw[cwire] (4,0) -- (6,0);
  \node[hbox] at (1.7,0.9) {};
  \node[hbox] at (1.7,-0.9) {};
  \node[zsp] at (4,0) {};
  \node[sdlabel,anchor=east] at (-0.1,0.9) {$q$};
  \node[sdlabel,anchor=east] at (-0.1,-0.9) {$\bar q$};
  \node[sdlabel,anchor=west] at (6.1,0) {$y$};
\end{tikzpicture}\]

A yellow box on a leg is the same thing as drawing that leg dashed,
Chapter 8's Hadamard edge (row 56). We draw the boxes, because in a
moment two of them will meet on one line.

On the blackboard the two dots are, with
\(\delta=\ket{00}\bra0+\ket{11}\bra1\) the copier of Chapter 11,

\[E_r=(H^r\otimes H^r)\,\delta,\qquad R_t=\delta^{\dagger}\,(H^t\otimes H^t).\]

The sender sends \(\ket x\) to \(H^r\ket x\otimes H^r\ket x\). That is
the doubled form of the pure state \(H^r\ket x\): the state on the wire
and its conjugate on the mirror copy, which is the same vector because
\(H^r\ket x\) is real. The reader takes a doubled state
\(\ket\psi\otimes\ket{\psi^{*}}\) to
\(\sum_y\lvert\bra yH^t\ket\psi\rvert^2\,\ket y\), the vector of
probabilities of the reading at setting \(t\) (Result 9.1, with \(H^t\)
in front).

\textbf{\emph{Why this definition.}} Why draw the sender and reader on
doubled wires, when the blackboard managed with one amplitude? Because
the round ends in a record, and a record is a probability, not an
amplitude. The doubled picture carries probabilities directly: sender,
then reader, is a process from Alice's classical bit to Bob's classical
bit, a table of odds. And the next two sections will put things on the
line, such as a spy whose copy is later forgotten, that are not pure
processes at all and can only be drawn doubled.

Now the headline. For a statement that something happens exactly when
two bits agree we need a symbol. Write \([\,y=x\,]\) for the number that
is \(1\) if \(y=x\) and \(0\) otherwise.

\begin{quote}
\textbf{Result 12.2 (BB84 correctness).} Let Alice send the bit \(x\) at
setting \(r\) and let Bob read at setting \(t\). Then Bob reads \(y\)
with probability
\[P(y\mid x,r,t)=\begin{cases}[\,y=x\,] & r=t,\\[2pt] \tfrac12 & r\neq t.\end{cases}\]
As doubled processes on the classical wires, exactly,
\[R_t\,E_r=\begin{cases}I & r=t,\\[2pt] \tfrac12\,J & r\neq t,\end{cases}\]
where \(J\) is the \(2\times2\) matrix with every entry \(1\). In
pictures: \textbf{matching settings fuse into a wire; mismatched
settings disconnect into noise.}
\end{quote}

\emph{Proof idea.} When the settings match, the reading undoes the
sending: two Hadamards in a row cancel, \(H^2=I\), and Bob's question is
the one Alice's state answers with certainty. When they differ, one
Hadamard is left over, and by Result 12.1 it spreads every plain state
evenly over both outcomes. In the picture, the matching round has a box
cancelling a box on each line, and the two green dots melt into a bare
wire. The mismatched round has a green dot and a red dot joined by two
lines, which Chapter 8's Hopf law says is no join at all: Bob's record
is left tied to nothing Alice did. The proof is given twice, in full, in
§12.6. See \texttt{Code/Ch12/TwoSettings.py}, which checks all sixteen
amplitudes and builds both rounds as pictures.

\vskip 0pt plus 4\baselineskip\penalty-100\vskip 0pt plus
-4\baselineskip

\begin{quote}
\textbf{Feynman's Notebook.}\nopagebreak

\emph{A conference in Bangalore.} The paper that named the protocol is
C. H. Bennett and G. Brassard, ``Quantum cryptography: Public key
distribution and coin tossing'', in the proceedings of the International
Conference on Computers, Systems \& Signal Processing, Bangalore, India,
December 10--12, 1984, pages 175--179. It gives the authors' addresses
as IBM Research, Yorktown Heights, and the department IRO of the
Université de Montréal; the authors have since posted a reprint on the
arXiv (2003.06557). Gisin, Ribordy, Tittel and Zbinden, reviewing the
field, remark that it appeared ``in a conference in India, totally
unknown to physicists''. The word \emph{conjugate} is older. Bennett and
Brassard take it from Stephen Wiesner's ``Conjugate Coding'', which
their reference list gives as a manuscript of about 1970, published in
\emph{SIGACT News} in 1983. Gisin et al.~place Wiesner at Columbia
University in the 1970s, and say his paper ``appeared only a decade
later''. The BB84 paper credits it with two applications: money that is
impossible to counterfeit, and two or three messages multiplexed so that
reading one destroys the others. To polarize light, Bennett and Brassard
suggest ``a Polaroid filter or calcite crystal''.
\end{quote}

\begin{center}\rule{0.5\linewidth}{0.5pt}\end{center}

\hypertarget{say-only-the-setting}{%
\section{§12.2 Say Only the Setting}\label{say-only-the-setting}}

Result 12.2 is a lock with two keys, and a key only works if both ends
use the same one. Alice and Bob live in different towns. If they agree
their settings in advance, anyone who overhears the agreement can read
the line at the right setting and learn everything. If they do not agree
them, half their rounds are coin tosses. The way out is to agree on the
settings \emph{afterwards}.

\textbf{On the blackboard.} Here is the protocol, as Bennett and
Brassard describe it and as Gisin and his co-authors retell it.

\begin{enumerate}
\def\labelenumi{\arabic{enumi}.}
\tightlist
\item
  For each round, Alice picks a bit \(x\) and a setting \(r\) at random,
  and sends \(H^r\ket x\).
\item
  Bob picks a setting \(t\) at random, independently of Alice, reads,
  and writes down his outcome \(y\).
\item
  Over an ordinary public channel, Bob announces his settings, one per
  round, and Alice says in which rounds hers was the same. Neither ever
  announces a bit.
\item
  They keep the rounds where \(r=t\) and strike the rest. The kept bits
  are the \textbf{sifted key}, Gisin et al.'s name for it.
\end{enumerate}

The public channel has to satisfy one condition, and it is worth saying
plainly. It need not be secret, but it must be \textbf{authentic}:
anyone may listen, nobody may alter or insert a message. Bennett and
Brassard assume a channel ``susceptible to eavesdropping but not to the
injection or alteration of messages''; Gisin et al.~put it as ``does not
have to be confidential, but has to be authentic''. Without that, a spy
could simply pose as Bob to Alice and as Alice to Bob. Gisin et al.~warn
that without some authentication Alice ``could actually be communicating
directly with Eve!''

A short run makes it concrete. The two coin-toss readings in rounds 2
and 4 could have come out either way.

\begingroup\renewcommand{\arraystretch}{1.3}

\begin{longtable}[]{@{}lcccccc@{}}
\toprule\noalign{}
round & 1 & 2 & 3 & 4 & 5 & 6 \\
\midrule\noalign{}
\endhead
\bottomrule\noalign{}
\endlastfoot
Alice's bit \(x\) & \(0\) & \(1\) & \(1\) & \(0\) & \(0\) & \(1\) \\
Alice's setting \(r\) & plain & plain & cross & cross & plain & cross \\
Bob's setting \(t\) & plain & cross & cross & plain & plain & cross \\
Bob's reading \(y\) & \(0\) & \(0\) & \(1\) & \(1\) & \(0\) & \(1\) \\
announced aloud & plain & cross & cross & plain & plain & cross \\
kept? & yes & no & yes & no & yes & yes \\
\end{longtable}

\endgroup

The kept rounds are \(1,3,5,6\), and on each of them Bob's reading is
Alice's bit: the sifted key is \(0,1,0,1\) at both ends. The two struck
rounds are the ones where a coin was tossed. A listener heard ``plain,
cross, cross, plain, plain, cross'' and ``yes, no, yes, no, yes, yes'',
and nothing else. Notice also who chose the key. Alice chose her bits,
but which of them survive was decided by both sets of random settings
together. Gisin et al.~make the same point: ``neither Alice nor Bob can
decide which key results from the protocol''.

\begin{quote}
\textbf{Result 12.3 (sifting).} Let \(x\), \(r\) and \(t\) be
independent fair coins in every round. Then, exactly: (i) the settings
match, \(r=t\), with probability \(\tfrac12\); (ii) on every kept round,
\(y=x\); (iii) over all rounds, before sifting, \(y\neq x\) with
probability \(\tfrac14\); (iv) the announcement tells a listener nothing
about the kept bit: \(P(x=0\mid r,t,\ r=t)=\tfrac12\) for each announced
pair.
\end{quote}

\emph{Proof.} (i) Of the four equally likely pairs \((r,t)\), two match.
(ii) Result 12.2 with \(r=t\). (iii) By Result 12.2, \(y\neq x\) happens
only when \(r\neq t\), and then with probability \(\tfrac12\); so
\(P(y\neq x)=\tfrac12\cdot\tfrac12=\tfrac14\). (iv) The announcement
consists of \(r\) and \(t\), and \(x\) was chosen independently of both,
so \(x\) is still a fair coin given any value of \((r,t)\). On a kept
round \(y=x\), so the same holds for \(y\). \(\;\blacksquare\) See
\texttt{Code/Ch12/SiftedKey.py}, which enumerates all eight triples
\((x,r,t)\) exactly and also runs a seeded sequence of rounds.

Part (iii) is Gisin et al.'s remark that before the settings are
compared Bob holds ``a string of bits with 25\% errors, called the raw
key''. Sifting removes all of them at the cost of half the rounds.

\textbf{\emph{Why this definition.}} Why announce the settings and not,
say, a few of the bits? Because the settings are exactly the information
that decides which rounds are useful, and, by (iv), no information at
all about what the useful rounds contain. A bit announced is a bit lost
to the key. A setting announced costs nothing, because it was chosen
independently of the bit it accompanies.

\textbf{In the sketchbook.} By Result 12.2, a kept round is a bare
classical wire and a struck round falls apart into \(\tfrac12\) times a
green effect on Alice's wire and a green point on Bob's. Now feed the
kept round a fair coin, and let Alice keep a copy of her bit. A fair
coin is the green point with \(\tfrac12\) in front:
\(\tfrac12(\ket0+\ket1)\) is the vector of probabilities
\((\tfrac12,\tfrac12)\). Copying a classical bit needs no permission.
The classical wire carries a settled value, and the green copier of
Chapter 11 copies exactly the settled values of its own basis (Result
11.2). The coin, the copier and the kept round are all green, so fusion
(Theorem 8.1) melts them into one dot:

\[\tfrac12\;
\begin{tikzpicture}[sd, x=1.2em, y=1.2em,
  zsp/.style={circle,draw,line width=0.9pt,fill=green!55,minimum size=1.2em,inner sep=0pt}]
  \draw[cwire] (0,0) -- (1.6,0);
  \draw[cwire] (1.6,0) -- (2.6,1) -- (9.4,1);
  \draw[cwire] (1.6,0) -- (2.6,-1) -- (3.8,-1);
  \draw[wire] (3.8,-1) -- (4.8,-0.2) -- (6.2,-0.2) -- (7.2,-1);
  \draw[wire] (3.8,-1) -- (4.8,-1.8) -- (6.2,-1.8) -- (7.2,-1);
  \draw[cwire] (7.2,-1) -- (9.4,-1);
  \node[zsp] at (0,0) {};
  \node[zsp] at (1.6,0) {};
  \node[zsp] at (3.8,-1) {};
  \node[zsp] at (7.2,-1) {};
  \node[sdlabel,anchor=west] at (9.6,1) {$k_A$};
  \node[sdlabel,anchor=west] at (9.6,-1) {$k_B$};
\end{tikzpicture}
\;\;\overset{\text{fusion}}{=}\;\;
\tfrac12\;
\begin{tikzpicture}[sd, x=1.2em, y=1.2em,
  zsp/.style={circle,draw,line width=0.9pt,fill=green!55,minimum size=1.2em,inner sep=0pt}]
  \draw[cwire] (0,0) -- (1.2,1) -- (3,1);
  \draw[cwire] (0,0) -- (1.2,-1) -- (3,-1);
  \node[zsp] at (0,0) {};
  \node[sdlabel,anchor=west] at (3.2,1) {$k_A$};
  \node[sdlabel,anchor=west] at (3.2,-1) {$k_B$};
\end{tikzpicture}\]

The round is drawn at the plain setting; at the cross setting it carries
two boxes on each line, which cancel in pairs (§12.6). Here \(k_A\) is
Alice's record of her bit and \(k_B\) is Bob's reading. The fusion
leaves the doubled lines between sender and reader closing one loop on
the fused dot, and a self-loop on a spider is erased with no number (van
de Wetering §4.1, as in §11.2). The right-hand side is a green dot with
two classical legs, \(\tfrac12\bigl(\ket{00}+\ket{11}\bigr)\): the joint
distribution \(p(k_A,k_B)=\tfrac12[\,k_A=k_B\,]\). It is a bend on the
two records, a shared fair coin. That is what a key \emph{is}. Section
12.4 will find the same bend with no coin tossed by anyone.

\begin{center}\rule{0.5\linewidth}{0.5pt}\end{center}

\hypertarget{eve-is-a-box-on-the-wire}{%
\section{§12.3 Eve Is a Box on the
Wire}\label{eve-is-a-box-on-the-wire}}

Now put someone on the line. Gisin et al.~describe what a spy would like
to do: send each qubit on to Bob ``in its original state, keeping a copy
for herself''. Chapter 11 showed that she cannot. What she can do is
what Eve did at the carrier's office in Chapter 11: press each passing
qubit with the green copier \(\delta\), keep one output, and send the
other on. By Result 11.4 that is a look at the plain setting. The
difference now is that Eve need not read her kept output at once. She
can hold it until Alice and Bob have announced their settings, and then
read it however she likes; Gisin et al.~allow the spy exactly this, and
more, in their account of the ideal attack (§VI.1).

\textbf{On the blackboard.} Take a kept round, so Alice and Bob use the
same setting, and follow the two cases.

\emph{A plain round.} Alice sends \(\ket x\). The press makes
\(\delta\ket x=\ket x\ket x\). Bob reads his copy plain and gets \(x\);
Eve's copy is \(\ket x\) too. Bob's key bit is right, and Eve holds it.

\emph{A cross round.} Alice sends \(H\ket x\), which is \(\ket+\) or
\(\ket-\). The press makes \(\delta\ket\pm=\ket{\Phi^{\pm}}\) (Result
11.2): not two copies of a cross state but a Bell pair, one half with
Bob and one half with Eve. Bob, reading his half at the cross setting,
reads the half of a Bell pair, which on its own is \(\tfrac12I\) (Result
9.3): a fair coin. He agrees with Alice only half the time. And Eve? If
she reads her half plain, she gets \(0\) or \(1\) with probability
\(\tfrac12\) each from both \(\ket{\Phi^+}\) and \(\ket{\Phi^-}\), so
her record says nothing about \(x\). If she waits, hears ``cross'', and
reads her half at the cross setting, she does no better. Rewriting the
two pairs at the cross setting,

\[\ket{\Phi^+}=\tfrac1{\sqrt2}\bigl(\ket{++}+\ket{--}\bigr),\qquad \ket{\Phi^-}=\tfrac1{\sqrt2}\bigl(\ket{+-}+\ket{-+}\bigr),\]

so her cross reading \(k_E\) and Bob's cross reading \(k_B\) satisfy
\(k_B\oplus k_E=x\), where \(\oplus\) is addition of bits mod \(2\).
That tells her \(x\) only if she also knows Bob's reading, and she never
sees it. Her own record, alone, is a fair coin whatever \(x\) was.

\begin{quote}
\textbf{Result 12.4 (the press on the wire).} Let Eve's press \(\delta\)
sit on the line in a kept round. Let \(k_B\) be Bob's reading and
\(k_E\) Eve's reading of her copy. Then, exactly: (i) at the plain
setting, \(p(k_B,k_E\mid x)=[\,k_B=x\,]\,[\,k_E=x\,]\): Bob is right and
Eve holds the key bit; (ii) at the cross setting, with Eve reading
plain, \(p(k_B,k_E\mid x)=\tfrac14\) for all four pairs; with Eve
reading cross, \(p(k_B,k_E\mid x)=\tfrac12\,[\,k_B\oplus k_E=x\,]\).
Either way Bob disagrees with Alice with probability \(\tfrac12\), and
Eve's record alone is independent of \(x\); (iii) with Alice's setting a
fair coin, Eve learns on average half a bit about each key bit, and Bob
disagrees with Alice on a quarter of the kept rounds.
\end{quote}

\emph{Proof.} (i) and (ii) are the calculations above; for Eve reading
plain on a cross round, write \(\ket{\Phi^{\pm}}\) with Bob's half at
the cross setting and Eve's at the plain one,
\(\ket{\Phi^+}=\tfrac12\bigl(\ket+(\ket0+\ket1)+\ket-(\ket0-\ket1)\bigr)\),
and likewise for \(\ket{\Phi^-}\): all four amplitudes have modulus
\(\tfrac12\). (iii) On plain rounds Eve's record determines \(x\), which
is one bit; on cross rounds it is independent of \(x\), which is none.
The two kinds of round are equally likely, so the average is
\(\tfrac12\) bit. The disagreement rate is \(0\) on plain rounds and
\(\tfrac12\) on cross rounds, which averages to \(\tfrac14\).
\(\;\blacksquare\) See \texttt{Code/Ch12/EvesBox.py}.

\vskip 0pt plus 4\baselineskip\penalty-100\vskip 0pt plus
-4\baselineskip

\begin{quote}
\textbf{Dirac's Blackboard.}\nopagebreak

\emph{Half a bit, counted.} ``Half a bit'' has a precise meaning. Let
\(h(p)=-p\log_2p-(1-p)\log_2(1-p)\) be the binary entropy, the
uncertainty in bits of a coin that lands \(0\) with probability \(p\); a
fair coin has \(h(\tfrac12)=1\). The information Eve's record carries
about Alice's bit is the uncertainty she had before, minus what remains
once she has read it:
\[\mathcal I(x;k_E)=h\bigl(\tfrac12\bigr)-\sum_k p(k_E{=}k)\,h\bigl(p(x{=}0\mid k_E{=}k)\bigr).\]
On a plain round \(x=k_E\), every \(p(x{=}0\mid k_E{=}k)\) is \(0\) or
\(1\), every \(h\) vanishes, and \(\mathcal I=1\). On a cross round
\(x\) is a fair coin given either record, every \(h\) is \(1\), and
\(\mathcal I=1-1=0\). Averaged over the two settings,
\(\tfrac12(1+0)=\tfrac12\) bit. The disagreement is simpler: \(0\) and
\(\tfrac12\), averaging \(\tfrac14\). Eve cannot improve on the plain
setting by choosing cross, because she does not know which rounds are
which until the settings are announced, and by then the qubit has gone
on to Bob.
\end{quote}

Is that a good trade for Eve? Bennett and Brassard answered in general,
and we state their answer without proving it. No measurement on a photon
in transit, by an eavesdropper who learns the photon's basis only after
measuring, ``can yield more than 1/2 expected bits of information about
the key bit encoded by that photon'', and any such measurement yielding
\(b\) bits of expected information ``must induce a disagreement with
probability at least \(b/2\)'' when the photon is later measured in its
original basis. The press gets \(b=\tfrac12\) and causes disagreement
\(\tfrac14=b/2\): it sits exactly on the bound. These are also the
numbers of the example Bennett and Brassard give of this ``optimum
tradeoff'', a different strategy in which the eavesdropper measures
every photon at the plain setting and sends on a fresh one, ``learning
the correct polarizations of half the photons and inducing disagreements
in 1/4'' of those later measured in the original basis. Gisin et
al.~call that strategy intercept--resend and give the same figures, 50\%
information and 25\% errors in the sifted key. They add that a spy who
applies it to only a tenth of the rounds causes about 2.5\% errors for
about 5\% information.

\textbf{\emph{Why this definition.}} Why take Chapter 11's copier as the
model spy, and not the most general attack quantum mechanics allows?
Because the copier is the look, drawn: Result 11.4 showed that pressing
a qubit and keeping one output is a plain reading whose outcome has not
yet been read. It is the simplest spy that learns anything, and it
already sits on Bennett and Brassard's bound. The most general spy is
the business of the security theorems of §12.8, and no picture in this
chapter reaches it.

\textbf{In the sketchbook.} Draw the press on the doubled line. The
copier acts on the wire, and its mirror copy acts on the mirror copy of
the wire; \(\delta\) is real, so the mirror copier is the green copier
again. Each copier has one leg in and two out. Its through-leg, \(b'\)
on the wire and \(\bar b'\) on the mirror, goes on to Bob's reader; its
kept leg, \(e\) and \(\bar e\), goes to Eve's reader. To bring each pair
of legs to its own reader, \(e\) and \(\bar b'\) cross once; a crossing
neither joins nor breaks a line. Here is a cross round, with Alice and
Bob at the cross setting and Eve reading plain:

\begin{center}
\begin{tikzpicture}[sd, x=1.2em, y=1.1em,
  zsp/.style={circle,draw,line width=0.9pt,fill=green!55,minimum size=1.2em,inner sep=0pt},
  hbox/.style={draw,line width=0.9pt,fill=yellow!85,minimum size=0.85em,inner sep=1pt}]
  \draw[cwire] (-2,0) -- (0,0);
  \draw[wire] (0,0) -- (1,1.6) -- (3.2,1.6);
  \draw[wire] (0,0) -- (1,-1.6) -- (3.2,-1.6);
  \node[hbox] at (1.9,1.6) {};
  \node[hbox] at (1.9,-1.6) {};
  \draw[wire] (3.2,1.6) -- (4.2,2.4) -- (10,2.4) -- (11,1.6);
  \draw[wire] (3.2,1.6) -- (4.2,0.8) -- (5,0.8) to[out=0,in=180] (7.6,-0.8) -- (10,-0.8) -- (11,-1.6);
  \draw[wire] (3.2,-1.6) -- (4.2,-0.8) -- (5,-0.8) to[out=0,in=180] (7.6,0.8) -- (10,0.8) -- (11,1.6);
  \draw[wire] (3.2,-1.6) -- (4.2,-2.4) -- (10,-2.4) -- (11,-1.6);
  \node[hbox] at (9,2.4) {};
  \node[hbox] at (9,0.8) {};
  \draw[cwire] (11,1.6) -- (13,1.6);
  \draw[cwire] (11,-1.6) -- (13,-1.6);
  \node[zsp] at (0,0) {};
  \node[zsp] at (3.2,1.6) {};
  \node[zsp] at (3.2,-1.6) {};
  \node[zsp] at (11,1.6) {};
  \node[zsp] at (11,-1.6) {};
  \node[sdlabel,anchor=east] at (-2.1,0) {$x$};
  \node[sdlabel,anchor=south] at (1.25,1.65) {$b$};
  \node[sdlabel,anchor=north] at (1.25,-1.65) {$\bar b$};
  \node[sdlabel,anchor=south] at (6.5,2.45) {$b'$};
  \node[sdlabel,anchor=north] at (4.6,0.75) {$e$};
  \node[sdlabel,anchor=south] at (4.6,-0.75) {$\bar b'$};
  \node[sdlabel,anchor=north] at (6.5,-2.45) {$\bar e$};
  \node[sdlabel,anchor=west] at (13.2,1.6) {$k_B$};
  \node[sdlabel,anchor=west] at (13.2,-1.6) {$k_E$};
\end{tikzpicture}
\end{center}

From left to right: Alice's sender at the cross setting; the press on
\(b\) and its mirror on \(\bar b\); Bob's reader at the cross setting on
\(b',\bar b'\); Eve's reader at the plain setting on \(e,\bar e\). This
is Eve as a box on the wire: a picture with a side road.

At the \textbf{plain} setting there are no boxes anywhere, every dot is
green and phase-free, and fusion (Theorem 8.1) makes the whole picture
one green dot with three classical legs, \(x\), \(k_B\) and \(k_E\). Six
lines join five dots, two more than a tree needs, so two self-loops are
removed on the way (van de Wetering §4.1). A green dot with legs \(x\)
in and \(k_B,k_E\) out is \(\sum_i\ket{ii}\bra i\): both records equal
Alice's bit, which is Result 12.4(i).

At the \textbf{cross} setting, ask what Eve learns while Bob's record is
forgotten. Forgetting a record is the green effect \(\sum_j\bra j\) on
its classical wire, and a reading followed by forgetting is a discard, a
cap on the two copies (Result 9.5). A unitary followed by a discard is a
discard (Result 9.4), so Bob's two yellow boxes go too, and his side of
the press is capped. A copier with one output capped onto its mirror
copy fuses with it into Chapter 11's four-legged green dot (row 85), and
Eve's reader, green and without boxes, fuses into that as well. What is
left is Alice's sender, its two boxes, and one green dot \(D\) with
Eve's record as its only classical leg. Then four steps:

\[\begin{tikzpicture}[sd, x=1.2em, y=1.2em,
  zsp/.style={circle,draw,line width=0.9pt,fill=green!55,minimum size=1.2em,inner sep=0pt},
  hbox/.style={draw,line width=0.9pt,fill=yellow!85,minimum size=0.85em,inner sep=1pt}]
  \draw[cwire] (-1.8,0) -- (0,0);
  \draw[wire] (0,0) to[out=60,in=120] (3.6,0);
  \draw[wire] (0,0) to[out=-60,in=-120] (3.6,0);
  \node[hbox] at (1.8,0.95) {};
  \node[hbox] at (1.8,-0.95) {};
  \draw[cwire] (3.6,0) -- (5.4,0);
  \node[zsp] at (0,0) {};
  \node[zsp] at (3.6,0) {};
  \node[sdlabel,anchor=east] at (-1.9,0) {$x$};
  \node[sdlabel,anchor=west] at (5.5,0) {$k_E$};
\end{tikzpicture}
\;\overset{\text{colour}}{=}\;
\begin{tikzpicture}[sd, x=1.2em, y=1.2em,
  zsp/.style={circle,draw,line width=0.9pt,fill=green!55,minimum size=1.2em,inner sep=0pt},
  xsp/.style={circle,draw,line width=0.9pt,fill=red!75,minimum size=1.2em,inner sep=0pt},
  hbox/.style={draw,line width=0.9pt,fill=yellow!85,minimum size=0.85em,inner sep=1pt}]
  \draw[cwire] (-2.4,0) -- (0,0);
  \node[hbox] at (-1.1,0) {};
  \draw[wire] (0,0) to[out=60,in=120] (3,0);
  \draw[wire] (0,0) to[out=-60,in=-120] (3,0);
  \draw[cwire] (3,0) -- (4.8,0);
  \node[xsp] at (0,0) {};
  \node[zsp] at (3,0) {};
\end{tikzpicture}\]

\[\;\overset{\text{Hopf}}{=}\;
\tfrac12\;
\begin{tikzpicture}[sd, x=1.2em, y=1.2em,
  zsp/.style={circle,draw,line width=0.9pt,fill=green!55,minimum size=1.2em,inner sep=0pt},
  xsp/.style={circle,draw,line width=0.9pt,fill=red!75,minimum size=1.2em,inner sep=0pt},
  hbox/.style={draw,line width=0.9pt,fill=yellow!85,minimum size=0.85em,inner sep=1pt}]
  \draw[cwire] (-2.4,0) -- (0,0);
  \node[hbox] at (-1.1,0) {};
  \draw[cwire] (1.4,0) -- (3.2,0);
  \node[xsp] at (0,0) {};
  \node[zsp] at (1.4,0) {};
\end{tikzpicture}
\;\overset{\text{colour}}{=}\;
\tfrac12\;
\begin{tikzpicture}[sd, x=1.2em, y=1.2em,
  zsp/.style={circle,draw,line width=0.9pt,fill=green!55,minimum size=1.2em,inner sep=0pt}]
  \draw[cwire] (-1.8,0) -- (0,0);
  \draw[cwire] (1.4,0) -- (3.2,0);
  \node[zsp] at (0,0) {};
  \node[zsp] at (1.4,0) {};
  \draw[sdcut] (0.7,1.2) -- (0.7,-1.2);
\end{tikzpicture}\]

\textbf{Colour change} (Theorem 8.2): put two yellow boxes on Alice's
classical leg, which changes nothing since \(H^2=I\); now the sender
wears a box on every leg, so it is a red dot, with one box left on
\(x\). \textbf{Hopf} (Theorem 8.6): a red dot and a green dot joined by
two lines fall apart, with the factor \(\tfrac12\). \textbf{Colour
change} again: a red effect after a yellow box is a green effect,
\(\sqrt2\bra0H=\bra0+\bra1\). The result is \(\tfrac12\) times a green
effect on \(x\) beside a green point on \(k_E\):
\(p(k_E\mid x)=\tfrac12\) for every \(x\) and \(k_E\). The dotted grey
curve is the cut of Chapter 11 (row 87), and it crosses nothing. Alice's
bit and Eve's record are disconnected. Every step is exact.

The other half of the story is Bob's disagreement, and it comes from the
same rules. Forget Eve's record instead. The press with one output
discarded is the four-legged green dot (row 85), a look nobody read,
sitting on the line between Alice's cross sender and Bob's cross reader.
Colour-change both of them to red. The four-legged dot is now joined to
a red dot by two lines on each side. Hopf applies twice, with
\(\tfrac12\) each time, and leaves the green dot with no legs at all,
whose value is \(1+1=2\) (§8.1). So the round is
\(\tfrac12\cdot\tfrac12\cdot2=\tfrac12\) times a green effect on Alice's
wire and a green point on Bob's: \(\tfrac12J\), exactly the mismatched
round of Result 12.2, although the settings matched. Eve's look has
turned a kept round into a struck one. That is the mark her look leaves.

In the dialogue, just after the tell-glass is set in the sleeve's place
and its reflection appears in the tray, Achilles glances at \emph{Three
Worlds}: a pond: leaves on its surface, trees reflected in it, a fish
beneath.

\begin{center}\rule{0.5\linewidth}{0.5pt}\end{center}

\hypertarget{backwards-to-the-source}{%
\section{§12.4 Backwards to the Source}\label{backwards-to-the-source}}

So far Alice has sent and Bob has read. There is a second way to grow
the same key, in which neither of them sends anything. A source between
them emits a pair, one half to each, and both of them read. This is the
idea behind the protocol of Ekert (1991), and in the form used here the
two ways turn out to be one.

\textbf{On the blackboard.} The source emits the Bell state of Chapter
6,

\[\ket{\Phi^+}=\tfrac1{\sqrt2}\bigl(\ket{00}+\ket{11}\bigr)=\tfrac1{\sqrt2}\,\text{cup},\]

where, as always in this book, the cup is unnormalized: cup
\(=\ket{00}+\ket{11}\), and cap \(=\bra{00}+\bra{11}\). Gisin et
al.~write this state as their Eq. (9), with spins up and down in place
of \(0\) and \(1\). Alice reads her half at setting \(r\), Bob reads his
at setting \(t\), and each writes down the outcome, \(k_A\) and \(k_B\).

Try the four pairs of settings. Both plain: the outcomes are \(00\) or
\(11\), each with probability \(\tfrac12\), so Alice and Bob agree. Both
cross: the same state rewritten at the cross setting is
\(\tfrac1{\sqrt2}(\ket{++}+\ket{--})\), as we saw in §12.3, and again
they agree. Plain against cross, in either order: all four outcomes come
with probability \(\tfrac14\). So the pair behaves exactly like a sent
qubit whose bit is a fair coin. Alice's reading plays the part of her
choice of \(x\), and Bob's reading is correct when the settings match
and a coin toss when they do not. Why should a reading at one end act
like a sending from that end?

The reason is the transpose trick of Chapter 6. Gisin et al.~state it as
their Eq. (10): for all unitaries \(U_1\) and \(U_2\),

\[(U_1\otimes U_2)\,\Phi^{(+)}=(\mathbb{1}\otimes U_2U_1^{t})\,\Phi^{(+)},\]

where \(U_1^{t}\) is the transpose and \(\mathbb 1\) the identity. A
process on Alice's half can be moved to Bob's half, where it arrives
\emph{transposed}. Reading is a process too, and the same move carries
Alice's reading to Bob's end. There it arrives as a \emph{preparation},
and its effect arrives conjugated: for every state \(\ket\varphi\),

\[(\bra\varphi\otimes I)\,\ket{\Phi^+}=\tfrac1{\sqrt2}\,\ket{\varphi^{*}},\]

where \(\ket{\varphi^{*}}\) is \(\ket\varphi\) with every amplitude
complex-conjugated, Chapter 6's \(\ket{b^{*}}\) (§6.4). When Alice's
reading at setting \(r\) gives \(k_A\), Bob's half is left in the
conjugate of \(H^r\ket{k_A}\). For the plain and cross settings that
conjugate is \(H^r\ket{k_A}\) itself, because those states are real. So
Alice's reading prepares for Bob exactly the state she would have sent
in BB84.

\begin{quote}
\textbf{Result 12.5 (E91 is BB84 bent).} Exactly: (i) for all
\(2\times2\) matrices \(U_1,U_2\),
\(\;(U_1\otimes U_2)\ket{\Phi^+}=(I\otimes U_2U_1^{T})\ket{\Phi^+}\);
(ii) for every state \(\ket\varphi\),
\(\;(\bra\varphi\otimes I)\ket{\Phi^+}=\tfrac1{\sqrt2}\ket{\varphi^{*}}\);
(iii) for Alice at the plain or cross setting and Bob at the plain,
cross or circular setting, the pair gives
\[p_{\text{pair}}(k_A,k_B\mid r,t)=\tfrac12\,P_{\text{BB84}}(y{=}k_B\mid x{=}k_A,r,t),\]
the BB84 round of Result 12.2 fed a fair coin; (iv) with both at the
circular setting,
\(p_{\text{pair}}(k_A,k_B)=\tfrac12\,[\,k_A\neq k_B\,]\): the two
readings \textbf{always disagree}, where the BB84 round at the circular
setting always agrees.
\end{quote}

Here \(P_{\text{BB84}}\) at the circular setting means Alice sends
\(\ket{c_1}\) for \(x=0\) and \(\ket{c_2}\) for \(x=1\), and Bob's
effects are \(\bra{c_1}\) and \(\bra{c_2}\).

\emph{Proof.} (i) Theorem 6.4 slides \(U_1\) round the cup:
\((U_1\otimes I)\,\text{cup}=(I\otimes U_1^{T})\,\text{cup}\). Apply
\(I\otimes U_2\) to both sides, and divide by \(\sqrt2\). (ii)
\((\bra\varphi\otimes I)\,\text{cup}=\sum_i\braket{\varphi|i}\ket i=\sum_i\overline{\varphi_i}\ket i=\ket{\varphi^{*}}\);
divide by \(\sqrt2\). (iii) Alice's effect at setting \(r\) with outcome
\(k_A\) is \(\bra{k_A}H^r\), which is real. By (ii) she leaves Bob's
half in \(\tfrac1{\sqrt2}H^r\ket{k_A}\), which is exactly what BB84's
sender emits for \(x=k_A\), scaled by \(\tfrac1{\sqrt2}\). Bob's reading
at any setting then gives the BB84 amplitude times \(\tfrac1{\sqrt2}\),
and squaring gives the factor \(\tfrac12\). (iv) Alice's effect
\(\bra{c_1}\) leaves Bob's half in
\(\tfrac1{\sqrt2}\ket{c_1^{*}}=\tfrac1{\sqrt2}\cdot\tfrac1{\sqrt2}(1,-i)\),
and \((1,-i)=-i\,(i,1)\), so \(\ket{c_1^{*}}=-i\ket{c_2}\). Likewise
\(\ket{c_2^{*}}=-i\ket{c_1}\). Bob, reading at the circular setting,
gets the other answer every time. In BB84, \(\ket{c_1}\) is read as
\(c_1\) every time. \(\;\blacksquare\) See
\texttt{Code/Ch12/BentKey.py}, which checks (i) on random complex
non-symmetric matrices, checks (ii) on random non-real states, and
confirms that the dagger in place of the transpose fails.

The circular test is the whole point of (iv). For the plain and cross
settings, a transpose and a conjugate look the same, because every
matrix and every state involved is real. Only a non-real setting can
tell them apart, and at the circular setting the pair and the sent qubit
give opposite answers. The BB84 paper's own two-photon state is a
different one, the singlet \(0.7071(r_1r_2-r_2r_1)\) in its notation,
whose halves are ``always found to have opposite polarization''. Bennett
and Brassard note that the same state vector ``is demonstrably equal
to'' \(0.7071(d_1d_2-d_2d_1)\) and to \(0.7071(c_1c_2-c_2c_1)\); with
the paper's own diagonal vectors, the first holds as a vector only up to
an overall sign, which no reading can see. So the singlet is
anticorrelated at every setting, the circular one included. The pair of
this chapter, \(\ket{\Phi^+}\), is not the singlet. It agrees at the
plain and cross settings and disagrees at the circular one, and that
difference is the conjugation.

\vskip 0pt plus 4\baselineskip\penalty-100\vskip 0pt plus
-4\baselineskip

\begin{quote}
\textbf{Dirac's Blackboard.}\nopagebreak

\emph{A transpose, not a dagger.} Eq. (10) moves \(U_1\) across the pair
as \(U_1^{T}\). It is tempting to write \(U_1^{\dagger}\), since the
dagger is the operation this book uses most for turning things round.
Check it on one entry. Take
\(U_1=\bigl(\begin{smallmatrix}1&0\\0&i\end{smallmatrix}\bigr)\), a
diagonal matrix, so \(U_1^{T}=U_1\) while
\(U_1^{\dagger}=\bigl(\begin{smallmatrix}1&0\\0&-i\end{smallmatrix}\bigr)\).
Then
\[(U_1\otimes I)\,\text{cup}=\ket{00}+i\ket{11}=(I\otimes U_1^{T})\,\text{cup},\]
\[(I\otimes U_1^{\dagger})\,\text{cup}=\ket{00}-i\ket{11}.\] The dagger
gives the wrong sign on \(\ket{11}\). In general,
\(\bigl(\bra{jk}(U_1\otimes I)\,\text{cup}\bigr)=(U_1)_{jk}\) and
\(\bigl(\bra{jk}(I\otimes M)\,\text{cup}\bigr)=M_{kj}\), so the two
sides agree exactly when \(M_{kj}=(U_1)_{jk}\), which is the transpose.
The dagger would also conjugate each entry. That conjugation is not
lost, though. It turns up in (ii), where an \emph{effect} slid round the
bend arrives as a conjugated \emph{state}, because an effect is already
the conjugate transpose of a state.
\end{quote}

\textbf{In the sketchbook.} Draw the pair doubled, as in Chapters 10 and
11: the cup on \((a,b)\) and a second cup on the mirror copies
\((\bar a,\bar b)\), with \(\tfrac12\) beside them, since
\(\ket{\Phi^+}\bra{\Phi^+}=\tfrac12\,\text{cup}\otimes\text{cup}\) in
doubled form. Put Alice's reader on \(a\) and \(\bar a\), at the cross
setting. Now move her two yellow boxes round their cups. Theorem 6.4
moves a box round a bend as its transpose, and \(H^{T}=H\), so each box
arrives on the other leg unchanged. A cup is a phase-free green dot with
two legs (§8.1), so the two cups and Alice's green reader are green dots
joined in a chain, and fusion (Theorem 8.1) makes them one:

\[\tfrac12\;
\begin{tikzpicture}[sd, x=1.2em, y=1.05em,
  zsp/.style={circle,draw,line width=0.9pt,fill=green!55,minimum size=1.2em,inner sep=0pt},
  hbox/.style={draw,line width=0.9pt,fill=yellow!85,minimum size=0.85em,inner sep=1pt}]
  \sdcup{0}{2.4}{-2.4}
  \sdcup{0}{0.8}{-0.8}
  \draw[wire] (0,2.4) -- (2.8,2.4) -- (3.8,1.6);
  \draw[wire] (0,0.8) -- (2.8,0.8) -- (3.8,1.6);
  \node[hbox] at (1.5,2.4) {};
  \node[hbox] at (1.5,0.8) {};
  \draw[cwire] (3.8,1.6) -- (5.6,1.6);
  \draw[wire] (0,-0.8) -- (5.6,-0.8);
  \draw[wire] (0,-2.4) -- (5.6,-2.4);
  \node[zsp] at (3.8,1.6) {};
  \node[sdlabel,anchor=south] at (0.6,2.45) {$\bar a$};
  \node[sdlabel,anchor=south] at (0.6,0.85) {$a$};
  \node[sdlabel,anchor=west] at (5.7,1.6) {$k_A$};
  \node[sdlabel,anchor=west] at (5.7,-0.8) {$b$};
  \node[sdlabel,anchor=west] at (5.7,-2.4) {$\bar b$};
\end{tikzpicture}
\;\overset{\text{slide}}{=}\;
\tfrac12\;
\begin{tikzpicture}[sd, x=1.2em, y=1.05em,
  zsp/.style={circle,draw,line width=0.9pt,fill=green!55,minimum size=1.2em,inner sep=0pt},
  hbox/.style={draw,line width=0.9pt,fill=yellow!85,minimum size=0.85em,inner sep=1pt}]
  \sdcup{0}{2.4}{-2.4}
  \sdcup{0}{0.8}{-0.8}
  \draw[wire] (0,2.4) -- (2.8,2.4) -- (3.8,1.6);
  \draw[wire] (0,0.8) -- (2.8,0.8) -- (3.8,1.6);
  \draw[cwire] (3.8,1.6) -- (5.6,1.6);
  \draw[wire] (0,-0.8) -- (5.6,-0.8);
  \draw[wire] (0,-2.4) -- (5.6,-2.4);
  \node[hbox] at (3,-0.8) {};
  \node[hbox] at (3,-2.4) {};
  \node[zsp] at (3.8,1.6) {};
\end{tikzpicture}\]

\[\;\overset{\text{fusion}}{=}\;
\tfrac12\;
\begin{tikzpicture}[sd, x=1.2em, y=1.05em,
  zsp/.style={circle,draw,line width=0.9pt,fill=green!55,minimum size=1.2em,inner sep=0pt},
  hbox/.style={draw,line width=0.9pt,fill=yellow!85,minimum size=0.85em,inner sep=1pt}]
  \draw[cwire] (0,-0.8) -- (1,1.6) -- (3.6,1.6);
  \draw[wire] (0,-0.8) -- (3.6,-0.8);
  \draw[wire] (0,-0.8) -- (1,-2.4) -- (3.6,-2.4);
  \node[hbox] at (2.2,-0.8) {};
  \node[hbox] at (2.2,-2.4) {};
  \node[zsp] at (0,-0.8) {};
  \node[sdlabel,anchor=west] at (3.7,1.6) {$k_A$};
  \node[sdlabel,anchor=west] at (3.7,-0.8) {$b$};
  \node[sdlabel,anchor=west] at (3.7,-2.4) {$\bar b$};
\end{tikzpicture}\]

The last picture is Alice's sender at the cross setting (row 89), with
its classical leg turned to point out of the picture instead of in, and
\(\tfrac12\) in front. It is the fair coin, Alice's copy of it and her
sender of §12.2, fused into one dot. That is Result 12.5(iii) as a
picture: \textbf{the cup and Alice's reader fuse into her sender}. The
reading has run backwards to the source and forwards again to Bob. Gisin
et al.~give the same image in words: in two-photon schemes, ``everything
is as if Alice's photon propagated backwards in time from Alice to the
source and then forwards from the source to Bob.'' The blackboard needed
Eq. (10) and a remark about conjugates to reach this. The picture needed
one slide and one fusion.

The circular setting shows the conjugation in the picture too. A
circular state is a cross state given a quarter-turn: \(\ket{c_1}\) is
\(Z(\tfrac\pi2)H\ket0\), a green dot of phase \(\tfrac\pi2\) applied
after the yellow box, and \(Z(\tfrac\pi2)H\ket1=-i\ket{c_2}\). So a
circular \emph{sender} puts a box and then a green \(+\tfrac\pi2\) dot
on the wire. A circular \emph{reader} undoes that: a green
\(-\tfrac\pi2\) dot, then the box, then the read-spider. On the mirror
copy the phases are negated, because the mirror conjugates, so the
circular reader carries \(+\tfrac\pi2\) on \(\bar q\) (row 93). Now
slide Alice's circular reader round the bend. A phase dot on one leg of
a cup can slide to the other leg unchanged, since a green dot is its own
transpose (only the connectivity matters, §8.5). So on Bob's wire \(b\)
the slid reader arrives as a box followed by a \(-\tfrac\pi2\) dot: a
circular sender with the \emph{wrong} quarter-turn. Bob's circular
reader adds its own \(-\tfrac\pi2\). Follow the wire \(b\) from the
source to Bob's read-spider:

\[\begin{tikzpicture}[sd, x=1.2em, y=1.2em,
  zsp/.style={circle,draw,line width=0.9pt,fill=green!55,minimum size=1.1em,inner sep=0pt},
  hbox/.style={draw,line width=0.9pt,fill=yellow!85,minimum size=0.85em,inner sep=1pt},
  ph/.style={font=\scriptsize,inner sep=1pt}]
  \draw[wire] (0,0) -- (8,0);
  \node[hbox] at (1.2,0) {};
  \node[zsp,label={[ph]above:{$-\tfrac\pi2$}}] at (3.2,0) {};
  \node[zsp,label={[ph]above:{$-\tfrac\pi2$}}] at (5.0,0) {};
  \node[hbox] at (6.8,0) {};
\end{tikzpicture}
\;\overset{\text{fusion}}{=}\;
\begin{tikzpicture}[sd, x=1.2em, y=1.2em,
  zsp/.style={circle,draw,line width=0.9pt,fill=green!55,minimum size=1.1em,inner sep=0pt},
  hbox/.style={draw,line width=0.9pt,fill=yellow!85,minimum size=0.85em,inner sep=1pt},
  ph/.style={font=\scriptsize,inner sep=1pt}]
  \draw[wire] (0,0) -- (5.6,0);
  \node[hbox] at (1.2,0) {};
  \node[zsp,label={[ph]above:{$\pi$}}] at (2.8,0) {};
  \node[hbox] at (4.4,0) {};
\end{tikzpicture}
\;\overset{\text{colour}}{=}\;
\begin{tikzpicture}[sd, x=1.2em, y=1.2em,
  xsp/.style={circle,draw,line width=0.9pt,fill=red!75,minimum size=1.1em,inner sep=0pt},
  ph/.style={font=\scriptsize,inner sep=1pt}]
  \draw[wire] (0,0) -- (2.8,0);
  \node[xsp,label={[ph]above:{$\pi$}}] at (1.4,0) {};
\end{tikzpicture}\]

The two quarter-turns fuse into a half-turn, \(-\pi\), which is the same
phase as \(\pi\) (a whole turn is nothing); the green half-turn between
two boxes is a red half-turn by colour change; and a red half-turn is
\(X\), a flip (§8.2). On the mirror copy the two \(+\tfrac\pi2\) dots
make the same flip. So Bob always reads the opposite of Alice: Result
12.5(iv), seen. In the BB84 round, Alice's circular sender carries
\(+\tfrac\pi2\), Bob's reader \(-\tfrac\pi2\), and they cancel to a bare
wire.

Two remarks keep the history straight. First, Ekert's own protocol used
a third setting and a test of Bell's inequality, as Gisin et al.~and
Horodecki et al.~both describe; this chapter does not run that test, and
uses the pair only at BB84's two settings. Second, that simplified use
is not a novelty of ours. Horodecki et al.~describe the protocol of
Bennett, Brassard and Mermin, in which Alice and Bob check their pairs
at the plain and cross settings only, and call this simplified form of
Ekert's protocol ``formally equivalent to the BB84 protocol''.

\textbf{\emph{Why this definition.}} Why bother with the pair at all, if
it gives the same key as a sent qubit? Because it moves the source out
of Alice's hands. Gisin et al.~point out that a source which simply sent
both parties the same BB84 state would have to be trusted, and ``could
be in Eve's hand''. A source of \(\ket{\Phi^+}\) can be tested by the
two ends themselves, and §12.7 shows what the test finds.

\vskip 0pt plus 4\baselineskip\penalty-100\vskip 0pt plus
-4\baselineskip

\begin{quote}
\textbf{Feynman's Notebook.}\nopagebreak

\emph{A second discovery.} Gisin et al.~credit the pair version to Artur
Ekert, then at Oxford University, who in 1991, ``while elaborating on a
suggestion of David Deutsch (1985), discovered QC independently of the
BB84 paper''. His paper, ``Quantum cryptography based on Bell's
theorem'' (\emph{Physical Review Letters}, 1991), we know only through
two reviews. Ekert proposed to found the security on Bell's inequality.
He gave Alice and Bob a third choice of basis, so that while they built
a key they also gathered the data to test the inequality; Gisin et
al.~note that this lowers the chance of matching bases from \(\tfrac12\)
to \(\tfrac29\). Horodecki et al.~add, in a footnote, that Ekert used
the fact that the singlet state cannot be correlated with any
environment, which is the monogamy of §12.7. The following year, Gisin
et al.~report, Bennett, Brassard and Mermin argued that the Bell
violation is not needed for security and stressed how close Ekert's
scheme is to BB84. Gisin et al.~also credit them with first noting the
close analogy between the one-photon and two-photon schemes: when Alice
detects her photon, she ``effectively prepares Bob's photon in a given
state''.
\end{quote}

\begin{center}\rule{0.5\linewidth}{0.5pt}\end{center}

\hypertarget{the-rosetta-table-10}{%
\section{§12.5 The Rosetta Table}\label{the-rosetta-table-10}}

The rows this chapter adds to the running dictionary, numbered on from
Chapter 11's eighty-eight. There is no new picture element in them.
Every row is a new \emph{use} of something we already had: the
read-spider run backwards as a sender (89), the red dot laid on
classical wires as a parity check (96), a green quarter-turn for the
circular setting (93), the copier on the line (92), the half-turn rings
on a bend (94) and the cut (91, 95).

\begingroup\renewcommand{\arraystretch}{1.5}

\begin{longtable}[]{@{}
  >{\raggedright\arraybackslash}p{(\columnwidth - 4\tabcolsep) * \real{0.2727}}
  >{\raggedright\arraybackslash}p{(\columnwidth - 4\tabcolsep) * \real{0.3636}}
  >{\raggedright\arraybackslash}p{(\columnwidth - 4\tabcolsep) * \real{0.3636}}@{}}
\toprule\noalign{}
\begin{minipage}[b]{\linewidth}\raggedright
\#
\end{minipage} & \begin{minipage}[b]{\linewidth}\raggedright
Algebra (Dirac)
\end{minipage} & \begin{minipage}[b]{\linewidth}\raggedright
Picture (Coecke)
\end{minipage} \\
\midrule\noalign{}
\endhead
\bottomrule\noalign{}
\endlastfoot
89 & two settings: plain \(\{\ket0,\ket1\}\), cross \(\{\ket+,\ket-\}\);
conjugate bases, \(\lvert\bra iH\ket j\rvert^2=\tfrac12\); the sender
\(E_r=(H^r\otimes H^r)\,\delta\) and the reader
\(R_t=\delta^{\dagger}(H^t\otimes H^t)\) & the \textbf{sender}: a green
dot with a classical wire in and a doubled wire out, a yellow box on
each quantum leg at the cross setting (a dashed edge, row 56); the
\textbf{reader}: Chapter 9's read-spider with the same boxes \\
90 & matching settings: \(R_rE_r=\delta^{\dagger}\delta=I\) &
\textbf{matching settings fuse into a wire}: the boxes cancel in pairs,
the two dots fuse along two lines, the self-loop goes; a bare classical
wire \\
91 & mismatched settings: \(R_tE_r=\tfrac12J\): Bob's reading a fair
coin, Alice's bit forgotten & \textbf{mismatched settings disconnect
into noise}: colour change, then Hopf removes both lines; a green effect
on Alice's wire and a green point on Bob's, \(\tfrac12\) beside them; a
cut crosses nothing \\
92 & Eve's press \(\delta\) on the line; plain: \(k_B=k_E=x\); cross:
Bob errs half the time and Eve's record is independent of \(x\);
averages \(\tfrac12\) bit and \(\tfrac14\) & \textbf{Eve is a box on the
wire}: the copier and its mirror on the doubled line, one leg to her
reader; with that leg forgotten, the four-legged dephasing dot (row
85) \\
93 & \((U_1\otimes U_2)\ket{\Phi^+}=(I\otimes U_2U_1^{T})\ket{\Phi^+}\);
\((\bra\varphi\otimes I)\ket{\Phi^+}=\tfrac1{\sqrt2}\ket{\varphi^{*}}\);
the pair \(=\tfrac12\,\)BB84 fed a fair coin & \textbf{E91 is BB84
bent}: the cup and Alice's reader fuse into her sender; a reading slid
round the bend arrives conjugated: unchanged at plain and cross; the
circular reader's \(-\tfrac\pi2\) dot keeps its phase round the bend,
where a circular sender carries \(+\tfrac\pi2\), so it arrives as the
other circular state \\
94 & plain disagreement \(=\) \(\ket{\Psi^{\pm}}\) (bit flips); cross
disagreement \(=\) \(\ket{\Phi^-},\ket{\Psi^-}\) (phase flips)
(Shor--Preskill) & on the bend: a \textbf{red} half-turn ring for a
plain disagreement, a \textbf{green} one for a cross disagreement (row
73); passing both checks leaves the bare bend \\
95 & monogamy, extreme form: the pass projector
\(\tfrac14(I+ZZ)(I+XX)=\ket{\Phi^+}\bra{\Phi^+}=\tfrac12\,\)cup\(\,\)cap;
a state that always passes is
\(\ket{\Phi^+}\bra{\Phi^+}\otimes\rho_{\mathrm{Eve}}\) & the pass
projector is a cap then a cup, \(\times\tfrac12\): \textbf{disconnected
in the middle}; a cut through empty paper separates the pair from
anything else \\
96 & parity \(w=k_A\oplus k_B\), with
\(X(2\to1)=\tfrac1{\sqrt2}\,[\,w=k_A\oplus k_B\,]\); \textbf{the Alarm}:
\(p(k_E,w)=\tfrac12[\,w=0\,]\) (plain), \(\tfrac12[\,k_E=w\,]\) (cross)
& a red dot on two classical wires (Chapter 8's parity, now on records);
\textbf{the master}: at plain Eve's record and the alarm fall apart, at
cross one bend joins them \\
\end{longtable}

\endgroup

Every row is exact, scalar included; none needs Chapter 8's ``up to a
non-zero scalar''. Two things behind these rows are cited and not
proved: Bennett and Brassard's general bound behind row 92, and the
quantitative form of monogamy of which row 95 is the extreme case. No
row says anything about security.

\begin{center}\rule{0.5\linewidth}{0.5pt}\end{center}

\hypertarget{the-theorem-twice-11}{%
\section{§12.6 The Theorem, Twice}\label{the-theorem-twice-11}}

\begin{quote}
\textbf{Result 12.2 (BB84 correctness), restated.} Alice sends \(x\) at
setting \(r\) and Bob reads at setting \(t\). Then
\(P(y\mid x,r,t)=[\,y=x\,]\) when \(r=t\) and \(\tfrac12\) when
\(r\neq t\); as doubled processes, \(R_tE_r=I\) when \(r=t\) and
\(\tfrac12J\) when \(r\neq t\), exactly.
\end{quote}

\hypertarget{diracs-blackboard-11}{%
\subsection{Dirac's Blackboard}\label{diracs-blackboard-11}}

\textbf{Step 1 --- the amplitude.} Alice's state is \(H^r\ket x\) and
Bob's effect for outcome \(y\) is \(\bra yH^t\) (§12.1). The amplitude
is \(\bra y\,H^tH^r\,\ket x\), the \((y,x)\) entry of the matrix
\(M_{rt}=H^tH^r\).

\textbf{Step 2 --- the four matrices.} Square the Hadamard:
\[H^2=\tfrac12\begin{pmatrix}1&1\\1&-1\end{pmatrix}\begin{pmatrix}1&1\\1&-1\end{pmatrix}=\tfrac12\begin{pmatrix}2&0\\0&2\end{pmatrix}=I .\]
So the four settings give
\[M_{00}=I,\qquad M_{11}=H^2=I,\qquad M_{01}=H,\qquad M_{10}=H .\] These
four matrices hold all sixteen amplitudes. When \(r=t\) the amplitude is
\(\braket{y|x}=[\,y=x\,]\). When \(r\neq t\) it is
\(\bra yH\ket x=(-1)^{xy}/\sqrt2\): that is \(+\tfrac1{\sqrt2}\) for
\((x,y)=(0,0),(0,1),(1,0)\) and \(-\tfrac1{\sqrt2}\) for \((1,1)\).

\textbf{Step 3 --- the probabilities.} Square the moduli. When \(r=t\),
\([\,y=x\,]^2=[\,y=x\,]\). When \(r\neq t\),
\(\bigl((-1)^{xy}/\sqrt2\bigr)^2=\tfrac12\) for all four pairs: the sign
\((-1)^{xy}\) is squared away. That is the first form of the result.

\textbf{Step 4 --- the doubled process.} Now the second form. With
\(\delta\ket x=\ket{xx}\) and \(\delta^{\dagger}=\sum_y\ket y\bra{yy}\),
and \(M=M_{rt}\) real,
\[\begin{aligned}R_tE_r&=\delta^{\dagger}(H^t\otimes H^t)(H^r\otimes H^r)\,\delta=\delta^{\dagger}(M\otimes M)\,\delta\\&=\sum_{x,y}\bra{yy}(M\otimes M)\ket{xx}\,\ket y\bra x=\sum_{x,y}M_{yx}^2\,\ket y\bra x .\end{aligned}\]
The entry in row \(y\), column \(x\) is the probability of Step 3. For
\(M=I\) this is \(I\). For \(M=H\) every entry is \(\tfrac12\), so it is
\(\tfrac12J\). \(\;\blacksquare\)

\hypertarget{coeckes-sketchbook-6}{%
\subsection{Coecke's Sketchbook}\label{coeckes-sketchbook-6}}

\textbf{Matching settings}, drawn at \(r=t=1\) (at \(r=t=0\) there are
no boxes and the chain starts at its second picture). Alice's sender on
the left, Bob's reader on the right, a box on each end of each line:

\[\begin{tikzpicture}[sd, x=1.2em, y=1.2em,
  zsp/.style={circle,draw,line width=0.9pt,fill=green!55,minimum size=1.2em,inner sep=0pt},
  hbox/.style={draw,line width=0.9pt,fill=yellow!85,minimum size=0.85em,inner sep=1pt}]
  \draw[cwire] (-1.8,0) -- (0,0);
  \draw[wire] (0,0) -- (0.9,0.9) -- (5.1,0.9) -- (6,0);
  \draw[wire] (0,0) -- (0.9,-0.9) -- (5.1,-0.9) -- (6,0);
  \node[hbox] at (2.1,0.9) {};
  \node[hbox] at (2.1,-0.9) {};
  \node[hbox] at (3.9,0.9) {};
  \node[hbox] at (3.9,-0.9) {};
  \draw[cwire] (6,0) -- (7.8,0);
  \node[zsp] at (0,0) {};
  \node[zsp] at (6,0) {};
  \node[sdlabel,anchor=east] at (-1.9,0) {$x$};
  \node[sdlabel,anchor=west] at (7.9,0) {$y$};
\end{tikzpicture}
\;\overset{\text{cancel}}{=}\;
\begin{tikzpicture}[sd, x=1.2em, y=1.2em,
  zsp/.style={circle,draw,line width=0.9pt,fill=green!55,minimum size=1.2em,inner sep=0pt}]
  \draw[cwire] (-1.8,0) -- (0,0);
  \draw[wire] (0,0) to[out=50,in=130] (3.4,0);
  \draw[wire] (0,0) to[out=-50,in=-130] (3.4,0);
  \draw[cwire] (3.4,0) -- (5.2,0);
  \node[zsp] at (0,0) {};
  \node[zsp] at (3.4,0) {};
\end{tikzpicture}\]

\[\;\overset{\text{fusion}}{=}\;
\begin{tikzpicture}[sd, x=1.2em, y=1.2em,
  zsp/.style={circle,draw,line width=0.9pt,fill=green!55,minimum size=1.2em,inner sep=0pt}]
  \draw[cwire] (-1.8,0) -- (0,0);
  \draw[cwire] (0,0) -- (1.8,0);
  \draw[wire] (0,0) .. controls (-1.1,1.9) and (1.1,1.9) .. (0,0);
  \node[zsp] at (0,0) {};
\end{tikzpicture}
\;\overset{\text{loop}}{=}\;
\begin{tikzpicture}[sd, x=1.2em, y=1.2em,
  zsp/.style={circle,draw,line width=0.9pt,fill=green!55,minimum size=1.2em,inner sep=0pt}]
  \draw[cwire] (-1.8,0) -- (1.8,0);
  \node[zsp] at (0,0) {};
\end{tikzpicture}
\;\overset{\text{identity}}{=}\;
\begin{tikzpicture}[sd, x=1.2em, y=1.2em]
  \draw[cwire] (0,0) -- (3,0);
  \node[sdlabel,anchor=east] at (-0.1,0) {$x$};
  \node[sdlabel,anchor=west] at (3.1,0) {$y$};
\end{tikzpicture}\]

\textbf{Cancel} (\(H^2=I\), Chapter 3; row 56): two boxes in a row on
one line cancel, on each of the two lines. \textbf{Fusion} (Theorem
8.1): two green dots that share a line are one green dot; they share
two, so the second becomes a loop from the fused dot to itself.
\textbf{Loop removal}: a self-loop on a phase-free spider is erased,
with no number (van de Wetering §4.1). \textbf{Identity} (§8.1): a
phase-free green dot with two legs is a plain wire. What is left is a
bare classical wire: \(R_rE_r=I\), exactly, with no scalar anywhere on
the way.

\textbf{Mismatched settings}, drawn at \(r=0\), \(t=1\): a plain sender,
so no boxes on Alice's side, and a cross reader.

\[\begin{tikzpicture}[sd, x=1.2em, y=1.2em,
  zsp/.style={circle,draw,line width=0.9pt,fill=green!55,minimum size=1.2em,inner sep=0pt},
  hbox/.style={draw,line width=0.9pt,fill=yellow!85,minimum size=0.85em,inner sep=1pt}]
  \draw[cwire] (-1.8,0) -- (0,0);
  \draw[wire] (0,0) -- (0.9,0.9) -- (3.3,0.9) -- (4.2,0);
  \draw[wire] (0,0) -- (0.9,-0.9) -- (3.3,-0.9) -- (4.2,0);
  \node[hbox] at (2.6,0.9) {};
  \node[hbox] at (2.6,-0.9) {};
  \draw[cwire] (4.2,0) -- (6,0);
  \node[zsp] at (0,0) {};
  \node[zsp] at (4.2,0) {};
  \node[sdlabel,anchor=east] at (-1.9,0) {$x$};
  \node[sdlabel,anchor=west] at (6.1,0) {$y$};
\end{tikzpicture}
\;\overset{\text{colour}}{=}\;
\begin{tikzpicture}[sd, x=1.2em, y=1.2em,
  zsp/.style={circle,draw,line width=0.9pt,fill=green!55,minimum size=1.2em,inner sep=0pt},
  xsp/.style={circle,draw,line width=0.9pt,fill=red!75,minimum size=1.2em,inner sep=0pt},
  hbox/.style={draw,line width=0.9pt,fill=yellow!85,minimum size=0.85em,inner sep=1pt}]
  \draw[cwire] (-1.8,0) -- (0,0);
  \draw[wire] (0,0) to[out=60,in=120] (3,0);
  \draw[wire] (0,0) to[out=-60,in=-120] (3,0);
  \draw[cwire] (3,0) -- (5.4,0);
  \node[hbox] at (4.1,0) {};
  \node[zsp] at (0,0) {};
  \node[xsp] at (3,0) {};
\end{tikzpicture}\]

\[\;\overset{\text{Hopf}}{=}\;
\tfrac12\;
\begin{tikzpicture}[sd, x=1.2em, y=1.2em,
  zsp/.style={circle,draw,line width=0.9pt,fill=green!55,minimum size=1.2em,inner sep=0pt},
  xsp/.style={circle,draw,line width=0.9pt,fill=red!75,minimum size=1.2em,inner sep=0pt},
  hbox/.style={draw,line width=0.9pt,fill=yellow!85,minimum size=0.85em,inner sep=1pt}]
  \draw[cwire] (-1.8,0) -- (0,0);
  \draw[cwire] (1.4,0) -- (3.8,0);
  \node[hbox] at (2.5,0) {};
  \node[zsp] at (0,0) {};
  \node[xsp] at (1.4,0) {};
\end{tikzpicture}
\;\overset{\text{colour}}{=}\;
\tfrac12\;
\begin{tikzpicture}[sd, x=1.2em, y=1.2em,
  zsp/.style={circle,draw,line width=0.9pt,fill=green!55,minimum size=1.2em,inner sep=0pt}]
  \draw[cwire] (-1.8,0) -- (0,0);
  \draw[cwire] (1.4,0) -- (3.2,0);
  \node[zsp] at (0,0) {};
  \node[zsp] at (1.4,0) {};
  \draw[sdcut] (0.7,1.2) -- (0.7,-1.2);
  \node[sdlabel,anchor=east] at (-1.9,0) {$x$};
  \node[sdlabel,anchor=west] at (3.3,0) {$y$};
\end{tikzpicture}\]

\textbf{Colour change} (Theorem 8.2): put two boxes on Bob's classical
leg, which changes nothing since \(H^2=I\); now his reader wears a box
on every leg, so it is a red dot, with one box left over on \(y\).
\textbf{Hopf} (Theorem 8.6): a green dot and a red dot joined by two
lines fall apart, with the factor \(\tfrac12\). \textbf{Colour change}
again: a red point followed by a box is a green point,
\(H\sqrt2\ket0=\ket0+\ket1\). The result is \(\tfrac12\) times a green
effect on \(x\) beside a green point on \(y\), which is
\(\tfrac12\sum_{x,y}\ket y\bra x=\tfrac12J\), exactly. The cut (row 87)
crosses nothing: Bob's reading is tied to nothing Alice did. The case
\(r=1\), \(t=0\) is the same chain reflected, with the colour change
made on Alice's sender and its leftover box on \(x\); it is the chain of
§12.3 with Eve's dot in Bob's place. \(\;\blacksquare\) See
\texttt{Code/Ch12/TwoSettings.py}, which builds all four rounds as
pictures and compares them, scalar included, with the matrices of the
blackboard.

\hypertarget{the-verdict-11}{%
\subsection{The verdict}\label{the-verdict-11}}

The blackboard is shorter. Its whole proof is one amplitude,
\(\bra yH^tH^r\ket x\), and the fact \(H^2=I\); the rest is squaring.
The sketchbook needs four rules for the matching round and three for the
mismatched one, and it has to be told what ``a green effect beside a
green point'' means before it can be read as \(\tfrac12J\). The
blackboard also keeps something the picture never had. The doubled
picture computes probabilities from the start, so the sign \((-1)^{xy}\)
of the mismatched amplitude, which the blackboard sees and then squares
away, never appears in it.

What the picture shows, and the blackboard does not, is \emph{why} the
mismatched round is useless. On the blackboard, \(\tfrac12J\) is a
matrix of equal numbers, and the reader must notice that equal numbers
mean independence. In the picture the round has come apart: there is no
line from Alice's bit to Bob's record, and a cut passes through empty
paper. And the same rules, with no new idea, turned the shared pair of
§12.4 into the sent qubit with one slide and one fusion, where the
blackboard needed a transpose identity and a remark about conjugates.
None of this is a claim about security. The two proofs establish the
same equality of processes, and nothing more.

\vskip 0pt plus 4\baselineskip\penalty-100\vskip 0pt plus
-4\baselineskip

\begin{quote}
\textbf{Coecke's Sketchbook.}\nopagebreak

\emph{Independence is disconnection.} The Hopf law did the work in every
mismatched round of this chapter, and it is worth saying what it means.
A green dot and a red dot joined by two lines come apart. Van de
Wetering's survey reads the rule as a statement about the two bases
(§4.5): ``having maximal information about the Z observable gives
minimal information about the X observable''. Here the Z observable is
our plain setting and the X observable our cross one, and the rule is
Result 12.1, drawn. In the same place the survey reads the bialgebra law
as a stronger form of the same idea, \emph{strong} complementarity;
§12.7 uses it. The picture adds one thing to the words. ``Minimal
information'' is a quantity, but ``disconnected'' is a shape, and a
shape can be checked at a glance. Whenever a round of this chapter came
apart, a cut through empty paper was available. Whenever a line
survived, as between Alice's bit and Eve's record at the plain setting,
it carried the information straight across.
\end{quote}

\begin{center}\rule{0.5\linewidth}{0.5pt}\end{center}

\hypertarget{the-alarm}{%
\section{§12.7 The Alarm}\label{the-alarm}}

The last step turns Eve's mark into an alarm. Alice and Bob have to look
for the mark on purpose. Bennett and Brassard have them compare, over
the public channel, some of the bits on which they think they should
agree. The compared bits are then lost to the key, and they suggest a
random subset, ``(say one third)'' of them, so that a spy who touched
more than a few photons would be unlikely to escape. In the pair form,
Shor and Preskill have Alice choose some of the pairs as \textbf{check
bits}, read at both settings. A check that finds a disagreement where
there should be none rings the alarm.

\textbf{On the blackboard.} What exactly does a check see? Write the
four Bell states of Chapter 10 as
\(\ket{B_{m_1m_2}}=(Z^{m_1}X^{m_2}\otimes I)\ket{\Phi^+}\), so that
\(\ket{B_{00}}=\ket{\Phi^+}\), \(\ket{B_{10}}=\ket{\Phi^-}\),
\(\ket{B_{01}}=\ket{\Psi^+}\) and \(\ket{B_{11}}=\ket{\Psi^-}\), where
\(\ket{\Psi^{\pm}}=\tfrac1{\sqrt2}(\ket{01}\pm\ket{10})\). The bit
\(m_2\) records whether the pair has suffered a flip, \(X\), and \(m_1\)
whether it has suffered a sign, \(Z\). For example,
\(\ket{\Psi^+}=\tfrac1{\sqrt2}(\ket{01}+\ket{10})\) read plain at both
ends always gives \(01\) or \(10\): the two readings always differ.
Rewritten at the cross setting it is
\(\tfrac1{\sqrt2}(\ket{++}-\ket{--})\), so read crosswise at both ends
the two readings always agree. A flip shows up at the plain check and
not at the cross one.

\begin{quote}
\textbf{Result 12.6 (what the checks see).} Exactly:
\[\ket{\Psi^+}\bra{\Psi^+}+\ket{\Psi^-}\bra{\Psi^-}=\ket{01}\bra{01}+\ket{10}\bra{10},\]
\[\ket{\Phi^-}\bra{\Phi^-}+\ket{\Psi^-}\bra{\Psi^-}=\ket{+-}\bra{+-}+\ket{-+}\bra{-+}.\]
So a plain check disagrees exactly on the Bell states with a flip
(\(m_2=1\)), and a cross check disagrees exactly on those with a sign
(\(m_1=1\)).
\end{quote}

\emph{Proof.} On the left of the first identity the cross terms
\(\pm\ket{01}\bra{10}\) and \(\pm\ket{10}\bra{01}\) come in with
opposite signs and cancel, leaving the two diagonal terms. For the
second, write both states at the cross setting. Substituting
\(\ket0=\tfrac1{\sqrt2}(\ket++\ket-)\) and
\(\ket1=\tfrac1{\sqrt2}(\ket+-\ket-)\) gives
\(\ket{\Phi^-}=\tfrac1{\sqrt2}(\ket{+-}+\ket{-+})\) and
\(\ket{\Psi^-}=\tfrac1{\sqrt2}(\ket{-+}-\ket{+-})\), and the same
cancellation happens. The right-hand sides are the projectors onto ``the
two plain readings differ'' and ``the two cross readings differ''.
\(\;\blacksquare\) See \texttt{Code/Ch12/Alarm.py}.

Shor and Preskill put the consequence plainly: the rates of bit flip
errors and of phase flip errors that Alice and Bob estimate from their
check bits ``are the same as they would have estimated using the Bell
basis measurement''. The checks at the two settings read off the two
bits \(m_2\) and \(m_1\), and nothing else.

\textbf{In the sketchbook.} Chapter 10 drew a Bell state as a bend
carrying two half-turn rings, a red ring \(m_2\pi\) and then a green
ring \(m_1\pi\), on its top leg (rows 73, 75). The cup is unnormalized,
cup \(=\ket{00}+\ket{11}\), hence the \(\tfrac1{\sqrt2}\):

\[\ket{B_{m_1m_2}}\;=\;\tfrac1{\sqrt2}\;
\begin{tikzpicture}[sd, x=1.4em, y=1.4em,
  zsp/.style={circle,draw,line width=0.9pt,fill=green!55,minimum size=1.1em,inner sep=0pt},
  xsp/.style={circle,draw,line width=0.9pt,fill=red!75,minimum size=1.1em,inner sep=0pt},
  ph/.style={font=\scriptsize,inner sep=1pt}]
  \sdcup{0}{1.5}{0}
  \draw[wire] (0,1.5) -- (4.6,1.5);
  \draw[wire] (0,0) -- (4.6,0);
  \node[xsp,label={[ph]above:{$m_2\pi$}}] at (1.4,1.5) {};
  \node[zsp,label={[ph]above:{$m_1\pi$}}] at (3.0,1.5) {};
\end{tikzpicture}\]

A plain check asks whether the two legs have the same plain value. The
bare bend, \(\ket{00}+\ket{11}\), always does. A red half-turn is a flip
of the plain value (§8.2), so the red ring makes the legs differ; a
green half-turn only changes a sign, which no plain reading can see. A
cross check asks the same about the cross values, and there colour
change swaps the roles: the green ring flips a cross value and the red
ring only signs it. So the plain check sees the red ring and the cross
check sees the green one. A pair that passes both checks every time has
neither ring, and what is left is the bare bend.

That last sentence is about Bell states. What about everything else?
Here is the spy of §12.3 again, now on a pair. Eve's press on Bob's half
turns the pair into
\[(I\otimes\delta)\ket{\Phi^+}=\tfrac1{\sqrt2}\bigl(\ket{000}+\ket{111}\bigr),\]
with the qubits in the order Alice, Bob, Eve. At the plain setting all
three readings agree: Alice and Bob pass the plain check every time, and
Eve holds the key. At the cross setting they fail half the time (Result
12.4). Could a cleverer Eve prepare some three-part state, entangled
with her own system in some cleverer way, that passes \emph{both} checks
every time and still leaves her a record tied to the key? With classical
bits she could: a coin copied three ways agrees with itself at every
comparison anyone cares to make. Horodecki et al.~note that classically,
if two bits are perfectly correlated, ``there is no constraints on
correlations'' between either of them and a third. With qubits, the
answer is no.

\begin{quote}
\textbf{Result 12.7 (monogamy, extreme form).} Let \(\rho\) be a state
of Alice's qubit, Bob's qubit and any finite system of Eve's. If Alice's
and Bob's plain readings agree with probability \(1\), and their cross
readings agree with probability \(1\), then
\[\rho=\ket{\Phi^+}\bra{\Phi^+}\otimes\rho_{\mathrm{Eve}}\] for some
state \(\rho_{\mathrm{Eve}}\) of Eve's system. Whatever Eve then reads
is independent of everything Alice and Bob read.
\end{quote}

\emph{Proof idea.} Each check is a projector, and a state that passes a
check with certainty lives inside its range. The two ranges meet in a
single line, the line of \(\ket{\Phi^+}\). A state on Alice and Bob that
is pinned to one pure state has nothing left over to share with anyone.

\emph{Proof.} The plain check passes on the range of
\(P_{\mathrm{plain}}=\ket{00}\bra{00}+\ket{11}\bra{11}=\tfrac12(I+Z\otimes Z)\),
and the cross check on the range of
\(P_{\mathrm{cross}}=\ket{++}\bra{++}+\ket{--}\bra{--}=\tfrac12(I+X\otimes X)\).
The first range is spanned by \(\ket{\Phi^+}\) and \(\ket{\Phi^-}\). The
second is spanned by \(\ket{++}+\ket{--}=\sqrt2\ket{\Phi^+}\) and
\(\ket{++}-\ket{--}=\sqrt2\ket{\Psi^+}\). They meet in the line of
\(\ket{\Phi^+}\) alone, so the projector onto both,
\[P=\tfrac14(I+Z\otimes Z)(I+X\otimes X)=\ket{\Phi^+}\bra{\Phi^+},\] has
rank \(1\). (The two factors commute, since
\((Z\otimes Z)(X\otimes X)=ZX\otimes ZX=XZ\otimes XZ\), the two minus
signs from \(ZX=-XZ\) cancelling.) Take first a pure state \(\ket\psi\)
of all three parts. If
\(\bra\psi(P_{\mathrm{plain}}\otimes I)\ket\psi=1\), the projection of
the unit vector \(\ket\psi\) has length \(1\), so \(\ket\psi\) lies in
the range of \(P_{\mathrm{plain}}\otimes I\); likewise for the cross
check. So
\(\ket\psi=(P\otimes I)\ket\psi=\ket{\Phi^+}\otimes\bigl((\bra{\Phi^+}\otimes I)\ket\psi\bigr)\),
a product with \(\ket{\Phi^+}\). For a mixed state, write
\(\rho=\sum_kp_k\ket{\psi_k}\bra{\psi_k}\) with every \(p_k>0\). The
passing probability is
\(\sum_kp_k\bra{\psi_k}(P_{\mathrm{plain}}\otimes I)\ket{\psi_k}\), an
average of numbers at most \(1\). It equals \(1\) only if every term
does, so each \(\ket{\psi_k}\) is \(\ket{\Phi^+}\otimes\ket{\eta_k}\),
and
\(\rho=\ket{\Phi^+}\bra{\Phi^+}\otimes\sum_kp_k\ket{\eta_k}\bra{\eta_k}\).
\(\;\blacksquare\) See \texttt{Code/Ch12/Alarm.py}, which checks the
projector identity and projects random three-part states.

In the sketchbook the projector is already a picture.
\(\ket{\Phi^+}\bra{\Phi^+}=\tfrac12\,\text{cup}\,\text{cap}\), with the
cup unnormalized as always, cup \(=\ket{00}+\ket{11}\): a cap, then a
cup, and \(\tfrac12\):

\[P\;=\;\tfrac12\;
\begin{tikzpicture}[sd, x=1.3em, y=1.3em]
  \draw[wire] (0,1) -- (1.2,1);
  \draw[wire] (0,-1) -- (1.2,-1);
  \sdcap{1.2}{1}{-1}
  \sdcup{4.6}{1}{-1}
  \draw[wire] (4.6,1) -- (5.8,1);
  \draw[wire] (4.6,-1) -- (5.8,-1);
  \draw[sdcut] (2.9,1.6) -- (2.9,-1.6);
\end{tikzpicture}\]

The picture is \textbf{disconnected in the middle}. Whatever comes in
from the left, including any tie to a system of Eve's, is closed off by
the cap, and what goes out on the right is a fresh bend tied to nothing.
The cut passes through empty paper. Horodecki et al.~state the extreme
case in one sentence: ``If two qubits A and B are maximally quantumly
correlated they cannot be correlated at all with third qubit C.'' They
go on to a trade-off between the two correlations when neither is
perfect. That quantitative form is beyond this chapter, and nothing here
proves it.

\textbf{\emph{Why this definition.}} Why check at both settings, when
either one alone catches the press half the time? Because each check
alone can be fooled. The state \(\tfrac1{\sqrt2}(\ket{000}+\ket{111})\)
passes the plain check every time and gives Eve the key; a state with
Eve's copy taken at the cross setting passes the cross check every time
and gives her the cross key. Only the pair of checks pins the state down
to a single line, and only a single line has nothing left over.

\hypertarget{a-check-round-with-a-spy}{%
\subsection{A check round with a spy}\label{a-check-round-with-a-spy}}

Now put everything in one picture: a check round with a spy. The source
emits the pair, doubled. Alice reads her half at the announced setting
\(S\) and records \(k_A\). Eve's press sits on Bob's half, with its
mirror on the mirror copy. Bob reads the press's through-legs at \(S\)
and records \(k_B\). Eve holds her kept legs until the setting has been
announced, which is her best case: Gisin et al.'s ideal eavesdropper may
keep her system for as long as she likes and choose her measurement
``after listening to all the public discussion between Alice and Bob''.
Then she reads at \(S\) and records \(k_E\). Finally, the alarm. A red
dot on two classical wires is Chapter 8's parity, now laid on records
(row 96): it takes \(k_A\) and \(k_B\) and emits \(w=k_A\oplus k_B\). On
the blackboard it is
\(X(2\to1)=\tfrac1{\sqrt2}\,[\,w=k_A\oplus k_B\,]\), so the parity
itself is \(\sqrt2\) times the red dot. The alarm sounds when \(w=1\).

The lines of the picture are the doubled lines \(a\), \(\bar a\)
(Alice's half and its mirror), \(b\), \(\bar b\) (Bob's half, into the
press), \(b'\), \(\bar b'\) (the press's through-legs, to Bob) and
\(e\), \(\bar e\) (its kept legs, to Eve), and the classical wires
\(k_A\), \(k_B\), \(k_E\) and \(w\). The cups are unnormalized, cup
\(=\ket{00}+\ket{11}\), so the doubled pair carries \(\tfrac12\); with
the \(\sqrt2\) of the parity, the picture carries \(\tfrac1{\sqrt2}\).
It is drawn at the cross setting; at the plain setting the yellow boxes
are absent.

\begin{center}
\begin{tikzpicture}[sd, x=1.25em, y=1.1em,
  zsp/.style={circle,draw,line width=0.9pt,fill=green!55,minimum size=1.2em,inner sep=0pt},
  xsp/.style={circle,draw,line width=0.9pt,fill=red!75,minimum size=1.2em,inner sep=0pt},
  hbox/.style={draw,line width=0.9pt,fill=yellow!85,minimum size=0.85em,inner sep=1pt}]
  \node[sdlabel, font=\normalsize, anchor=east] at (-3.9,8) {$\tfrac1{\sqrt2}$};
  \sdcup{0}{12}{4}
  \sdcup{0}{10}{8}
  \draw[wire] (0,12) -- (5,12) -- (6.5,11);
  \draw[wire] (0,10) -- (5,10) -- (6.5,11);
  \node[hbox] at (3,12) {};
  \node[hbox] at (3,10) {};
  \draw[cwire] (6.5,11) -- (15,11) -- (16.5,9.5);
  \draw[wire] (0,8) -- (3.5,8);
  \draw[wire] (0,4) -- (3.5,4);
  \draw[wire] (3.5,8) -- (4.5,8.8) -- (11,8.8) -- (12,8);
  \draw[wire] (3.5,8) -- (4.5,7.2) -- (5.5,7.2) to[out=0,in=180] (8.5,4.8) -- (11,4.8) -- (12,4);
  \draw[wire] (3.5,4) -- (4.5,4.8) -- (5.5,4.8) to[out=0,in=180] (8.5,7.2) -- (11,7.2) -- (12,8);
  \draw[wire] (3.5,4) -- (4.5,3.2) -- (11,3.2) -- (12,4);
  \node[hbox] at (10,8.8) {};
  \node[hbox] at (10,7.2) {};
  \node[hbox] at (10,4.8) {};
  \node[hbox] at (10,3.2) {};
  \draw[cwire] (12,8) -- (15,8) -- (16.5,9.5);
  \draw[cwire] (16.5,9.5) -- (19,9.5);
  \draw[cwire] (12,4) -- (19,4);
  \node[zsp] at (6.5,11) {};
  \node[zsp] at (3.5,8) {};
  \node[zsp] at (3.5,4) {};
  \node[zsp] at (12,8) {};
  \node[zsp] at (12,4) {};
  \node[xsp] at (16.5,9.5) {};
  \node[sdlabel,anchor=south] at (1.5,12.05) {$\bar a$};
  \node[sdlabel,anchor=south] at (1.5,10.05) {$a$};
  \node[sdlabel,anchor=south] at (1.5,8.05) {$b$};
  \node[sdlabel,anchor=north] at (1.5,3.95) {$\bar b$};
  \node[sdlabel,anchor=south] at (7,8.85) {$b'$};
  \node[sdlabel,anchor=north] at (5.1,7.15) {$e$};
  \node[sdlabel,anchor=south] at (5.1,4.85) {$\bar b'$};
  \node[sdlabel,anchor=north] at (7,3.15) {$\bar e$};
  \node[sdlabel,anchor=south] at (10,11.1) {$k_A$};
  \node[sdlabel,anchor=north] at (14,7.9) {$k_B$};
  \node[sdlabel,anchor=west] at (19.2,9.5) {$w$};
  \node[sdlabel,anchor=west] at (19.2,4) {$k_E$};
\end{tikzpicture}
\end{center}

This is \textbf{The Alarm}. From left to right: the doubled pair, two
nested cups; Alice's reader on \(a,\bar a\), top; Eve's press and its
mirror on \(b\) and \(\bar b\); the lines \(e\) and \(\bar b'\) crossing
once, so that each reader gets a wire and its mirror (a crossing neither
joins nor breaks a line); Bob's reader and Eve's reader; and, far right,
the red parity dot on \(k_A\) and \(k_B\).

\begin{quote}
\textbf{Result 12.8 (the Alarm).} In a check round at announced setting
\(S\), with Eve's press on Bob's half and Eve reading at \(S\) after the
announcement, Eve's record and the alarm are distributed as
\[p(k_E,w)=\begin{cases}\tfrac12\,[\,w=0\,] & S=\text{plain},\\[2pt] \tfrac12\,[\,k_E=w\,] & S=\text{cross}.\end{cases}\]
At the plain setting the alarm never sounds and Eve's record is the key.
At the cross setting Eve's record \emph{is} the alarm bit: she learns
whether she was caught, and nothing of the key, since \(k_A\) and
\(k_E\) are then independent fair coins.
\end{quote}

\emph{Proof, on the blackboard.} The press turns the pair into
\(\tfrac1{\sqrt2}(\ket{000}+\ket{111})\) on \((a,b',e)\), as above, and
all three are read at \(S\). At the plain setting the three readings are
\(000\) or \(111\), each with probability \(\tfrac12\). So
\(k_A=k_B=k_E\), \(w=0\) always, and \(k_E\) is a fair coin. At the
cross setting apply \(H\) to each qubit:
\[\begin{aligned}(H\otimes H\otimes H)\tfrac1{\sqrt2}\bigl(\ket{000}+\ket{111}\bigr)&=\tfrac1{\sqrt2}\bigl(\ket{+++}+\ket{---}\bigr)\\&=\tfrac12\!\!\sum_{k_A\oplus k_B\oplus k_E=0}\!\!\ket{k_Ak_Bk_E},\end{aligned}\]
since each \(\ket{k_Ak_Bk_E}\) appears in \(\ket{+++}\) with amplitude
\(\tfrac1{2\sqrt2}\) and in \(\ket{---}\) with amplitude
\(\tfrac{(-1)^{k_A+k_B+k_E}}{2\sqrt2}\). So each of the four triples of
even parity has probability \(\tfrac14\) and the others none. Then
\(w=k_A\oplus k_B=k_E\), and summing over \(k_A\) gives
\(p(k_E,w)=2\cdot\tfrac14\,[\,k_E=w\,]=\tfrac12[\,k_E=w\,]\). The pair
\((k_A,k_E)\) runs over all four values with probability \(\tfrac14\)
each. \(\;\blacksquare\)

\emph{Proof, in the sketchbook.} At the \textbf{plain} setting there are
no boxes. The two cups, the two copiers and the three readers are all
phase-free green dots, and fusion (Theorem 8.1) makes them one green
dot. It has the three classical legs \(k_A\), \(k_B\), \(k_E\), and two
self-loops, which are removed (van de Wetering §4.1): the quantum lines
joined those dots with two more lines than a tree on them needs. The
green dot and the red parity dot now share the two lines \(k_A\) and
\(k_B\), and Hopf (Theorem 8.6) takes them apart with the factor
\(\tfrac12\):

\[\tfrac1{\sqrt2}\;
\begin{tikzpicture}[sd, x=1.2em, y=1.2em,
  zsp/.style={circle,draw,line width=0.9pt,fill=green!55,minimum size=1.2em,inner sep=0pt},
  xsp/.style={circle,draw,line width=0.9pt,fill=red!75,minimum size=1.2em,inner sep=0pt}]
  \draw[cwire] (0,0) to[out=50,in=130] (3.4,0);
  \draw[cwire] (0,0) to[out=-50,in=-130] (3.4,0);
  \draw[cwire] (3.4,0) -- (5.2,0);
  \draw[cwire] (0,0) -- (0,-1.8) -- (5.2,-1.8);
  \node[zsp] at (0,0) {};
  \node[xsp] at (3.4,0) {};
  \node[sdlabel,anchor=west] at (5.3,0) {$w$};
  \node[sdlabel,anchor=west] at (5.3,-1.8) {$k_E$};
\end{tikzpicture}
\;\overset{\text{Hopf}}{=}\;
\tfrac1{2\sqrt2}\;
\begin{tikzpicture}[sd, x=1.2em, y=1.2em,
  zsp/.style={circle,draw,line width=0.9pt,fill=green!55,minimum size=1.2em,inner sep=0pt},
  xsp/.style={circle,draw,line width=0.9pt,fill=red!75,minimum size=1.2em,inner sep=0pt}]
  \draw[cwire] (0,0) -- (1.8,0);
  \draw[cwire] (0,-1.8) -- (1.8,-1.8);
  \node[xsp] at (0,0) {};
  \node[zsp] at (0,-1.8) {};
  \draw[sdcut] (-0.9,-0.9) -- (2.4,-0.9);
  \node[sdlabel,anchor=west] at (1.9,0) {$w$};
  \node[sdlabel,anchor=west] at (1.9,-1.8) {$k_E$};
\end{tikzpicture}\]

A red point is \(\sqrt2\ket0\) and a green point is \(\ket0+\ket1\), so
the right-hand side is
\(\tfrac1{2\sqrt2}\cdot\sqrt2\,[\,w=0\,]=\tfrac12[\,w=0\,]\) for either
\(k_E\). Eve's record and the alarm are disconnected, and the alarm is
pinned at silence.

At the \textbf{cross} setting, six steps.

\begin{enumerate}
\def\labelenumi{\arabic{enumi}.}
\tightlist
\item
  \textbf{Fusion} (Theorem 8.1). The cup on \((a,b)\) fuses with the
  copier on \(b\) into a green dot \(G_1\) with legs \(a\), \(b'\),
  \(e\); likewise the mirror cup and the mirror copier make \(G_2\) with
  legs \(\bar a\), \(\bar b'\), \(\bar e\).
\item
  \textbf{Colour change} (Theorem 8.2). Each reader has a box on both
  quantum legs. Put two boxes on its classical leg (\(H^2=I\)) and it
  becomes a red dot, with one box left on its record.
\item
  \textbf{Bialgebra} (Theorem 8.5). The green \(G_1,G_2\) and the red
  readers of Alice and Bob are joined by all four lines
  \(a,\bar a,b',\bar b'\): a complete square. The bialgebra law replaces
  it with one red dot, holding the green dots' other legs \(e,\bar e\),
  and one green dot, holding the red dots' other legs \(k_A,k_B\) (each
  with its box), joined by one link. The factor is \(\tfrac1{\sqrt2}\),
  which with the \(\tfrac1{\sqrt2}\) in front makes \(\tfrac12\):
\end{enumerate}

\[\tfrac12\;
\begin{tikzpicture}[sd, x=1.2em, y=1.1em,
  zsp/.style={circle,draw,line width=0.9pt,fill=green!55,minimum size=1.2em,inner sep=0pt},
  xsp/.style={circle,draw,line width=0.9pt,fill=red!75,minimum size=1.2em,inner sep=0pt},
  hbox/.style={draw,line width=0.9pt,fill=yellow!85,minimum size=0.85em,inner sep=1pt}]
  \draw[wire] (0,-1.6) -- (0,1.6);
  \draw[cwire] (0,1.6) -- (1,2.6) -- (5,2.6) -- (6,1.6);
  \draw[cwire] (0,1.6) -- (1,0.6) -- (5,0.6) -- (6,1.6);
  \node[hbox] at (2.2,2.6) {};
  \node[hbox] at (2.2,0.6) {};
  \draw[cwire] (6,1.6) -- (8,1.6);
  \draw[wire] (0,-1.6) to[out=50,in=130] (3.4,-1.6);
  \draw[wire] (0,-1.6) to[out=-50,in=-130] (3.4,-1.6);
  \draw[cwire] (3.4,-1.6) -- (8,-1.6);
  \node[hbox] at (4.8,-1.6) {};
  \node[zsp] at (0,1.6) {};
  \node[xsp] at (0,-1.6) {};
  \node[xsp] at (3.4,-1.6) {};
  \node[xsp] at (6,1.6) {};
  \node[sdlabel,anchor=north] at (1.7,-2.35) {$e,\bar e$};
  \node[sdlabel,anchor=south] at (3.6,2.65) {$k_A$};
  \node[sdlabel,anchor=north] at (3.6,0.55) {$k_B$};
  \node[sdlabel,anchor=west] at (8.1,1.6) {$w$};
  \node[sdlabel,anchor=west] at (8.1,-1.6) {$k_E$};
\end{tikzpicture}\]

\begin{enumerate}
\def\labelenumi{\arabic{enumi}.}
\setcounter{enumi}{3}
\tightlist
\item
  \textbf{Colour change} on the parity dot. It becomes green with a box
  on each of its legs. The boxes on \(k_A\) and \(k_B\) meet the boxes
  already there and cancel (\(H^2=I\)); one box stays on \(w\).
\item
  \textbf{Fusion}, twice. The green dot on top and the green parity dot
  share \(k_A\) and \(k_B\): they fuse into one green dot, and the
  second shared line is a self-loop, removed. The red dot below and
  Eve's red reader share \(e\) and \(\bar e\): they fuse likewise, with
  a self-loop removed.
\item
  \textbf{Identity} (§8.1). Each of the two surviving dots has two legs,
  so each is a plain wire. What is left is one wire from \(k_E\) through
  a box, through the link, through a box to \(w\), and the two boxes
  cancel.
\end{enumerate}

\[\text{The Alarm at cross}\;=\;\tfrac12\;
\begin{tikzpicture}[sd, x=1.2em, y=1.2em,
  zsp/.style={circle,draw,line width=0.9pt,fill=green!55,minimum size=1.2em,inner sep=0pt}]
  \draw[cwire] (0,0) -- (1,0.9) -- (2.8,0.9);
  \draw[cwire] (0,0) -- (1,-0.9) -- (2.8,-0.9);
  \node[zsp] at (0,0) {};
  \node[sdlabel,anchor=west] at (2.9,0.9) {$w$};
  \node[sdlabel,anchor=west] at (2.9,-0.9) {$k_E$};
\end{tikzpicture}
\;=\;\tfrac12\bigl(\ket{00}+\ket{11}\bigr)=\tfrac12\,[\,k_E=w\,].\]

A green dot with two legs is a bend (§8.1): \textbf{one bend joins Eve's
record to the alarm}. Every step is exact. \(\;\blacksquare\) See
\texttt{Code/Ch12/Alarm.py}, which builds The Alarm at both settings,
checks it against the blackboard, and asks the automatic simplifier for
the two shapes: nothing joins \(k_E\) to \(w\) at the plain setting, and
a path does at the cross setting.

Average over the two settings, which are equally likely. The alarm never
sounds on plain check rounds and sounds on half of the cross ones, so it
sounds on a quarter of all checked rounds. Eve's record is the key bit
on plain rounds and independent of it on cross rounds, so she holds half
a bit per key bit. These are the numbers of Result 12.4, and of Bennett
and Brassard's optimum example. The difference is that now they are
heard: the mark Eve's look leaves has become the sound of the bell.

The Alarm also bends back into the form of §12.3. Slide Alice's boxes
round the bend and fuse, as in §12.4: the cup and Alice's reader become
her sender, with \(k_A\) as its classical leg, and the picture becomes
Eve's box on the line of a sent qubit, with Alice's bit copied to the
alarm.

\begin{center}\rule{0.5\linewidth}{0.5pt}\end{center}

\hypertarget{what-the-lock-does-not-prove}{%
\section{§12.8 What the Lock Does Not
Prove}\label{what-the-lock-does-not-prove}}

Take stock of what has been proved. A round at matching settings
delivers Alice's bit to Bob every time, and a round at mismatched
settings delivers a coin (Result 12.2). Announcing the settings costs
nothing and keeps half the rounds (Result 12.3). A spy who presses the
line learns half a bit per key bit and pays for it with errors on a
quarter of the kept rounds (Result 12.4), and the checks hear those
errors as an alarm (Result 12.8). A pair that passes both checks every
time has nothing left over for anyone (Result 12.7). Each of these is an
exact equality, and each has been proved.

None of them says the key is secret. The press is one spy among
infinitely many. A real eavesdropper may entangle each passing qubit
with a probe of her own design, keep all her probes, and measure them
together after every word of the public discussion. And real checks
never come out perfect: some errors come from the line and the
detectors, and Alice and Bob cannot tell those from a spy's. Result 12.7
covers only the case in which the checks never fail. What is needed is a
theorem that turns a \emph{small} measured error rate into a
\emph{small} bound on what any eavesdropper can know. That theorem
exists.

\begin{quote}
\textbf{Result 12.9 (Shor and Preskill; stated, not proved).} The BB84
protocol, with its sifting, a random sample of check bits at both
settings, and suitable classical error correction and privacy
amplification, is secure against an eavesdropper able to perform any
operation quantum mechanics permits. Precisely: the probability is
exponentially small that Alice and Bob agree on a key about which Eve
can obtain more than an exponentially small amount of information. This
holds as long as the bit- and phase-error rates measured on the check
bits are below \(11\%\), the point at which the rate
\(1-2h(\mathcal D)\) reaches \(0\), where \(\mathcal D\) is the error
rate and \(h\) the binary entropy of §12.3. The proof assumes that
Alice's source emits perfect single photons.
\end{quote}

Here is the shape of their argument, which we do not reproduce. They
first give a protocol in which Alice and Bob share many noisy pairs, and
purify them. The pairs are purified into fewer, nearly perfect copies of
\(\ket{\Phi^+}\) using quantum error-correcting codes of a family known
as CSS codes. That protocol they prove secure with methods from a proof
by Lo and Chau. Then they show that Alice and Bob can do all of it with
single qubits sent and read at two settings, so that the pair protocol
collapses into BB84. The last step works because this family of codes
corrects bit errors and phase errors separately. The checks enter
through Result 12.6: the bit-flip rate and the phase-flip rate are what
the plain and cross checks measure. What the check bits estimate is what
the codes must correct.

Gisin et al.~give the thresholds in context. Against \emph{individual}
attacks, where Eve attaches an independent probe to each qubit and
measures her probes one after the other, BB84 is secure up to an error
rate \(\mathcal D_0=\tfrac12\bigl(1-\tfrac1{\sqrt2}\bigr)\approx15\%\).
They note that this is also the error rate up to which Alice and Bob
could still see a violation of Bell's inequality in the CHSH form of
Chapter 7. Against \emph{coherent} attacks, where Eve may process
several qubits together, the bound is \(11\%\), ``precisely that
obtained in Mayers proof (after improvement by P. Shor and J. Preskill
(2000))''. Between \(11\%\) and \(15\%\), Gisin et al.~wrote, it was not
known what happens.

Two classical steps finish the protocol, and neither is quantum at all.
\emph{Error correction} removes the disagreements that noise, or a spy,
left in the sifted key. \emph{Privacy amplification} shrinks what Eve
knows. Gisin et al.~give its simplest form. Alice picks two bits of the
key, announces only \emph{which} two (say bit \(103\) and bit \(537\)),
and both replace them by their sum mod \(2\). If Eve knows the first bit
and nothing of the second, she knows nothing of the sum. If she knows
each with probability \(60\%\), she guesses the sum right only with
probability \(0.6^2+0.4^2=52\%\). Repeated, and done in larger blocks,
this drives her information as low as Alice and Bob wish, at the cost of
key length.

We should be exact about the scope of the pictures. The rewrites of this
chapter prove \emph{correctness}, Result 12.2; the numbers of one spy,
Results 12.4 and 12.8; the identities behind the checks, Result 12.6;
and monogamy in its extreme case, Result 12.7. They prove no security
theorem, no key rate and no threshold. The \(11\%\) and \(15\%\) are
stated and cited here, and proved in neither of our two notations. The
blackboard does not prove them either, in this chapter. Result 12.9 is a
theorem about many rounds at once, about error-correcting codes and
about bounds on information, and neither of our strands was built for
that.

\vskip 0pt plus 4\baselineskip\penalty-100\vskip 0pt plus
-4\baselineskip

\begin{quote}
\textbf{Coecke's Sketchbook.}\nopagebreak

\emph{What the rewrites do not prove.} A chain of rewrites proves that
two pictures are the same process. It is tempting to read more into a
picture that falls apart, and to say that a disconnected Eve ``cannot
know'' the key. The picture licenses something smaller. In §12.6, Bob's
record is disconnected from Alice's bit in \emph{that} round, with
\emph{that} reader. In §12.7, Eve's record is disconnected from the
alarm for \emph{that} spy, the copier, at the plain setting. Change the
spy and the picture changes with it, and there are infinitely many
spies. Monogamy is the nearest thing here to a statement about every
spy, and even it is only the extreme case: a pair that never fails a
check. A key that fails a few checks in a hundred needs Result 12.9, and
Result 12.9 needs a theory of errors and codes that is not in the
sketchbook. The pictures show what one protocol does with one spy. They
are not a security proof.
\end{quote}

\vskip 0pt plus 4\baselineskip\penalty-100\vskip 0pt plus
-4\baselineskip

\begin{quote}
\textbf{Feynman's Notebook.}\nopagebreak

\emph{A simple proof, sixteen years on.} Peter W. Shor, of AT\&T Labs
Research in Florham Park, New Jersey, and John Preskill, of the
Lauritsen Laboratory of High Energy Physics at the California Institute
of Technology, posted their proof on the arXiv as quant-ph/0003004; part
of the work was done while Shor was visiting Caltech. Their opening sets
the scene: BB84 was proposed in 1984, but ``for many years, it was not
rigorously proven secure against an adversary able to perform any
physical operation permitted by quantum mechanics''. They count three
recent proofs, none entirely satisfactory. Lo and Chau's is easy to
understand but needs a protocol that requires a quantum computer. The
other two, one of them by Dominic Mayers, cover a protocol based on BB84
but are ``quite complicated''. Shor and Preskill's proof, they write,
was inspired by noticing that the codes they use ``are hidden in the
inner workings'' of Mayers's proof. They end on its weakness. Like the
two earlier proofs for BB84, it needs perfect single-photon sources,
while most experiments use weak coherent sources, which ``no currently
known proof covers''.
\end{quote}

The lock is built. Two settings that are each a coin to the other gave a
key that arrives whole when the settings match. Announcing the settings
afterwards kept the key and lost nothing. A spy on the line became a
green box with a side road, and the price of her look could be written
down. The shared pair bent back into the sent qubit with one yank. And a
check round with a spy in it, The Alarm, came apart at the plain setting
and joined Eve's record to the bell at the cross setting. The part the
pictures could not reach, security against every spy, is there to be
cited, and we have said whose it is.

This is the end of Part III. The pictures have become proofs: of copying
and its impossibility, of teleportation, and now of a lock made of
light. Part IV asks something different of them. Until now, but for
Chapter 5's single oracle, every box in the book has been ours, built
from gates we chose. In Chapter 13 someone else hands us a box and will
not say what is inside it. We may ask it a question, once. The box turns
out to be a web of CNOTs, and yellow boxes turn its colours over. Its
secret can be read straight off the picture, if we ask in the right way.

\hypertarget{minuet-for-one-wish}{%
\chapter{Minuet for One Wish}\label{minuet-for-one-wish}}

\begin{center}\rule{0.5\linewidth}{0.5pt}\end{center}

\emph{霜の夜や}\\
\emph{見えぬ峰より}\\
\emph{返る声}

\emph{shimo no yo ya}\\
\emph{mienu mine yori}\\
\emph{kaeru koe}

\emph{A night of hard frost ---}\\
\emph{from the peaks I cannot see}\\
\emph{my own voice comes back.}

\begin{center}\rule{0.5\linewidth}{0.5pt}\end{center}

\emph{{[}The Tortoise's parlour, towards eight on the same frosty
evening. The second pot stands empty between the two cups, and the stove
has burned down low.{]}}

\emph{{[}It has settled to a dull red behind its mica, and its iron
ticks now and then as it cools. On the lower pane the two ferns stand
where they met at the middle; above them the glass is black, and past it
the lane lies pale under the rime, the curtained windows opposite still
lit. Achilles has come back from the window to his chair by the fender
with his stockinged feet to the stove, and the Tortoise sits across from
him with her empty cup in her lap. Low as the stove is, the parlour is
still the warmest room in the house.{]}}

\textbf{Achilles:} \emph{{[}tilting the pot, and finding nothing in
it{]}} That's the second pot gone, and I haven't the face to ask for a
third.

\textbf{Tortoise:} You'd only have to carry it up the stairs.

\textbf{Achilles:} Up the stairs? \emph{{[}He looks at the window.{]}}
The ferns have met in the middle of the pane. I was about to say that
means it's time I went home.

\textbf{Tortoise:} It means the frost has finished with the parlour.
There's something upstairs I've meant to show you since before the
frost, and tonight is cold enough and late enough for it.

\textbf{Achilles:} In all the winters I've been coming here, I've never
once been up your stairs. I'd assumed they only led to more stairs.

\textbf{Tortoise:} They lead to one small room, and nobody has set foot
in it since the frost came. \emph{{[}She takes the candlestick from the
mantel and lights it at the stove.{]}} This is all the light the top
room gets. Keep close behind it on the stairs: one of the treads is
loose, and I never remember which.

\textbf{Achilles:} \emph{{[}picking up her work-bag{]}} I'll carry this,
then. It's heavier than knitting.

\textbf{Tortoise:} My cloth, and some string, and a few things you
haven't seen yet. Mind the tread.

\emph{{[}They go up, the candle first. Halfway, a tread gives under
Achilles's foot with a groan, and he stops short; the Tortoise, above
him, does not turn round.{]}}

\textbf{Tortoise:} That one.

\emph{{[}The stair is narrow, and makes one turn at a small landing.
Their shadows go up the wall beside them, huge and unsteady, stretching
tall when the candle dips and shrinking when it rises, sliding round the
turn of the landing and back; not one of them keeps its place for two
steps together. The flame leans in the cold that comes down the stairs
to meet them, and every shadow leans with it.{]}}

\textbf{Achilles:} \emph{{[}at the head of the stairs, his breath
showing{]}} It's colder with every step. You could keep butter up here.

\textbf{Tortoise:} I do, in summer. \emph{{[}She sets the candlestick on
the small table in the middle of the room, takes her drawing-cloth from
the work-bag, and spreads it beside the candle.{]}} Sit. Yours is the
chair without the wobble.

\emph{{[}The top room: a small table and two chairs, one window over the
lane, and on the wall beside the window a shelf. On the shelf stand an
old brass lamp and a lidded casket.{]}}

\emph{{[}The top room is small and very cold, tucked under the slope of
the roof, and it has one of everything. One door, which had stuck in its
frame with the frost and gave all at once when the Tortoise put her
shoulder to it, and which stands open now on the dark of the stairs. One
window over the lane, its pane furred white with frost, the ferns grown
well past the middle, so that the lane shows through nowhere. One shelf,
on the wall beside the window. One light: the candle on the table. On
the floor under the shelf the Tortoise's summer larder is put by for the
winter, a row of empty earthenware butter-crocks and, among them, one
tall glass jar. The lamp on the shelf is tarnished nearly brown, with no
oil in it, and the casket beside it is of dark old wood, the dust lying
undisturbed along its lid; neither has been touched since the frost
came, nor for a long while before. On the wall facing the window, not
above the shelf, hangs a print in a plain dark frame: Escher's
\textup{Belvedere}, a pillared lookout whose upper storey is set
crosswise on the lower; a ladder stands inside it and leans against the
outside. The candle just reaches it across the table, and it takes what
little light there is, as it has taken it every winter since long before
the Tortoise came to the house. Then the Tortoise lowers herself into
the other chair, with her back to that wall, and her shadow goes up over
half the print. With the candle still at last, the flame stands up, and
every shadow in the room settles where it belongs, thrown straight back
from it: the chair-backs up the walls, the lamp and the casket together
on the slope of the ceiling, and Achilles's, still standing, huge across
the frosted pane.{]}}

\textbf{Achilles:} \emph{{[}sitting, and looking up at the shelf{]}} Is
that what you've brought me up to see? A lamp and a jewel box?

\textbf{Tortoise:} Neither of them is mine. They came with the house,
and they aren't mine to lend.

\begin{center}\rule{0.5\linewidth}{0.5pt}\end{center}

\hypertarget{first-challenge-one-wish-one-answer}{%
\section{\texorpdfstring{First Challenge --- \emph{One Wish, One
Answer}}{First Challenge --- One Wish, One Answer}}\label{first-challenge-one-wish-one-answer}}

\textbf{Tortoise:} \emph{{[}holding the candlestick up to the shelf{]}}
This is the Genie's casket. It has been sealed as long as I've had the
house, and it opens once, after it has done its one piece of work.

\textbf{Achilles:} There are knobs down this side. And little windows
down the other.

\textbf{Tortoise:} Down the left face, the knobs: each turns down for
dark or up for lit, and tonight they all stand down. The lowest, set a
little apart from the rest, is the flag knob; the ones above it are
wish-knobs.

\textbf{Tortoise:} Down the right face, a window level with each knob,
and the flag window lowest. A window shows nothing until the casket
passes; then it shows dark, a plain pane, or lit, a small gleam behind
the glass, and it keeps what it shows.

\textbf{Achilles:} And these rings round every knob and every window?

\textbf{Tortoise:} Collars, like the collars on my keys downstairs, but
with only two stops: the upright stroke and the cross. Every one of them
stands upright tonight. Setting the knobs and turning the collars is
free, all night if you like; only the passage costs.

\textbf{Achilles:} \emph{{[}standing, and bending to the lid{]}} There's
a little enamel plate on the top. \emph{{[}reading{]}} ``The Terms. One
wish, one passage. Behind each wish-knob there is a mark, or none. With
every collar upright, each wish-window shows its own knob as it was set,
and the flag window shows the flag knob as it was set, turned over once
for each marked knob that was set lit. What the collars do crosswise is
the wisher's business.''

\textbf{Achilles:} So the secret is which knobs are marked. And it says
nothing about how the marks are made.

\textbf{Tortoise:} Nothing at all. The Terms are short, and they keep to
them.

\textbf{Achilles:} \emph{{[}his hand going to the lamp{]}} And the lamp
---

\textbf{Tortoise:} Is how you wish. Rub it once and the Genie comes; he
grants one wish, and the casket passes once. Not yet: we should know
what to wish before you rub it.

\textbf{Achilles:} \emph{{[}taking his hand back, and sitting down
again{]}} One wish. Then we'd better spend it on the right knobs.

\textbf{Tortoise:} \emph{{[}setting the candlestick back on the
table{]}} Then we'll practise on a box we can spend all night on.
\emph{{[}She tears a strip from the back of her notebook and rules it
across into a row of little squares.{]}} One square for each wish-knob.
This is your slip: ink some of the squares where I can't see, and fold
it, and the inked squares are your marks.

\emph{{[}With the candlestick back on the table, the shadows that swung
round the walls when she held it up to the shelf have come back to their
places, each where it lay before: the chair-backs, the lamp and the
casket on the slope of the ceiling, her own over the print. Achilles
pulls his chair in close and holds his hands to the flame a moment,
working the cold out of his fingers.{]}}

\emph{{[}Achilles turns his shoulder to her, inks some of the squares,
and folds the slip along its length so that the squares are inside.{]}}

\textbf{Tortoise:} \emph{{[}bringing a small counter out of the
work-bag{]}} And this is our flag, a disc. It's gilt on one face, for
lit, and black on the other, for dark. I name a wish and hand you the
disc with whichever face up I choose; you look at your slip, turn the
disc over once for each inked square among the places I named lit, and
hand it back.

\textbf{Achilles:} Just as the Terms say the flag window goes.

\textbf{Tortoise:} Just so. The first and the third lit, the rest dark.
\emph{{[}She puts the disc in his hand, black face up.{]}}

\emph{{[}The disc lies on Achilles's palm, black face up: dark. He opens
the slip behind his other hand. Of the two places she named lit, the
first is inked and the third is bare. He turns the disc over once
between finger and thumb: the black face goes down against his palm, and
the gilt comes up. It is the same small counter in the same hand; the
slip is folded again, and the Tortoise has not seen it; the wish she
named has not changed. Only the face that shows has changed: dark has
become lit.{]}}

\textbf{Achilles:} \emph{{[}handing it back, gilt up{]}} There.

\textbf{Tortoise:} Lit. One face, one answer.

\textbf{Achilles:} Only because you asked so meanly: two places, one
disc. Let me choose the wish. Light every place at once, and the disc
will have to account for every mark I made; with more places lit, it has
more to say.

\textbf{Tortoise:} Choose, then. \emph{{[}She hands him the disc, black
up.{]}}

\emph{{[}Every place lit. Achilles runs his thumb along the whole slip
under his hand and turns the disc once for each inked square he passes:
over to gilt, and then over again to black. It comes to rest on his palm
black up, exactly as she handed it to him. The slip and its marks are as
they were; the disc shows the same face as before the wish, as though
nothing had been asked.{]}}

\textbf{Achilles:} Black. As if I'd made no marks at all. \emph{{[}He
frowns at the slip.{]}} Or two, or four. Every mark turns it over, and a
disc turned twice is back where it started. It has only two faces.
Whatever I light, all it can ever say is whether the count came out odd
or even.

\textbf{Tortoise:} One face, one answer, however many places you light.
Now my first wish again, the first and the third lit, but I'll hand you
the disc gilt up.

\emph{{[}The disc goes to him gilt up. The first place is inked, as
before; he turns the disc once, and it goes back to her black up. Last
time it went to him black and came back gilt; this time it went to him
gilt and came back black. The same one mark has turned it the same once.
What it was handed is what has changed, and so what comes back.{]}}

\textbf{Achilles:} Turned the other way from before. Of course. The one
mark turns over whatever face it's handed.

\textbf{Tortoise:} Then the same wish twice running. \emph{{[}She hands
it to him black up, the first and the third lit.{]}}

\emph{{[}He turns it once, gilt up, and hands it back. She hands it
straight back to him, gilt up, naming the same wish; the same inked
square turns it the same once more, and it lies on her palm black up,
the face she first gave. The slip has not moved, nor has the wish. The
disc has been turned twice, and the two turnings have undone each
other.{]}}

\textbf{Achilles:} Back as you gave it. The wish, asked twice, undoes
itself.

\emph{{[}The Tortoise leans back from the table, and her chair, which
has its wobble as promised, rocks over onto its other short leg with a
small clack. She leans in again, and it rocks back onto the first with
another. It has two ways to stand, and no third.{]}}

\textbf{Tortoise:} Now I'll read your slip. The first place alone lit.
\emph{{[}The disc goes to him black up.{]}}

\textbf{Achilles:} \emph{{[}turning it once{]}} Gilt.

\textbf{Tortoise:} Marked. The second alone. \emph{{[}Black up.{]}}

\textbf{Achilles:} \emph{{[}handing it straight back{]}} Black.

\textbf{Tortoise:} Bare.

\emph{{[}So on down the slip, a single place lit each time and the disc
handed over black each time: it comes back black at the third, black at
the fourth, gilt at the fifth, black at the sixth.{]}}

\textbf{Tortoise:} The first and the fifth.

\textbf{Achilles:} \emph{{[}unfolding the slip: the first square inked,
and the fifth{]}} The first and the fifth. A wish for every knob, and
one place read by each.

\textbf{Tortoise:} Could you do it with fewer? Change places.
\emph{{[}She tears another strip, rules it, turns her shoulder, inks it,
and folds it.{]}} My slip, my disc; you're the wisher. Read my marks
with fewer wishes than there are knobs.

\textbf{Achilles:} In pairs, then; you can't hide from pairs. The first
two lit.

\emph{{[}She takes the disc black up, looks at her slip, turns it once,
and hands it back gilt.{]}}

\textbf{Achilles:} The third and the fourth.

\emph{{[}Black comes back black.{]}}

\textbf{Achilles:} The fifth and the sixth.

\emph{{[}Black side up to her again; her slip turns it once, and it
comes back gilt.{]}}

\textbf{Achilles:} And the first, the third and the fifth.

\emph{{[}Black side up once more; it comes back black, as it went.{]}}

\textbf{Achilles:} One of the first two is marked, and one of the last
two. The third and fourth go together, both or neither. And the first,
third and fifth come out even. So I say the second and the sixth, and
nothing in the middle.

\textbf{Tortoise:} Perhaps. Or perhaps this. \emph{{[}She tears a third
strip, rules it, inks it in front of him, and lays it open beside her
folded one.{]}}

\emph{{[}The new slip lies open on the cloth: the first square inked,
the third, the fourth and the sixth. Achilles goes down his wishes
against it one at a time, a finger counting the inked squares among the
places he named lit. The first two: one inked square, the first; odd,
and the disc would have turned. The third and fourth: both inked; even,
and the disc would have come back as given. The fifth and sixth: one,
the sixth; odd. The first, third and fifth: two; even. Wish by wish, the
open slip would have turned the disc exactly as her folded one did.{]}}

\textbf{Achilles:} It answers every one of my wishes just as yours did.
And it isn't your slip?

\textbf{Tortoise:} \emph{{[}unfolding hers: the second square inked, and
the sixth{]}} It isn't. You guessed right, and your wishes could never
have told you so. Whatever fewer wishes you choose, I can ink you a
second slip that answers them all alike; and the wisher never holds the
slip he is wishing against.

\textbf{Achilles:} \emph{{[}sitting back{]}} What does one wish buy?

\textbf{Tortoise:} One answer, and only the one.

\textbf{Achilles:} Then the rest must be bought one at a time.

\textbf{Tortoise:} \emph{{[}drawing the work-bag to her{]}} Then let me
draw what your slip is. My drawing-cloth: whatever I lay on it is a
picture of a thing, not the thing, and no wish is ever spent on it.

\textbf{Tortoise:} \emph{{[}bringing out string and a little box of
rings, green and red{]}} My string, and my rings. A green ring threaded
on one cord and tied by a short length of string to a red ring threaded
on another is a mark: whenever the green ring's cord is lit, the red
ring turns the other cord over.

\emph{{[}The cloth lies bare between them. She lays a cord across it
from left to right for each square of Achilles's slip, one under another
like the knobs down the casket's face, and below them, set a little
apart, a longer cord for the flag. On the cord for the first square she
threads a green ring, and on the flag cord a red ring beneath it, and
ties the two together with a short length that runs down across the
cords between without touching them. On the cord for the fifth square
she threads another green ring, tied down to a red ring of its own
further along the flag cord. The other cords she leaves straight and
bare. What was hidden in a folded slip lies open now: the marks have
become ties, and the unmarked places are cords with nothing on them at
all.{]}}

\textbf{Tortoise:} That is all a box of marks is.

\textbf{Achilles:} A tie from every mark down to the flag. Light a
marked cord and its tie turns the flag over; light a bare one and
nothing happens. \emph{{[}He runs a finger along the flag cord, past one
red ring and then the other.{]}} My disc, drawn. Then the casket will
take a wish for every knob, and the Genie grants one.

\textbf{Tortoise:} Unless the one wish is set some other way. What do
you suppose the collars are for?

\begin{center}\rule{0.5\linewidth}{0.5pt}\end{center}

\hypertarget{second-challenge-the-mark-on-the-cords}{%
\section{\texorpdfstring{Second Challenge --- \emph{The Mark on the
Cords}}{Second Challenge --- The Mark on the Cords}}\label{second-challenge-the-mark-on-the-cords}}

\textbf{Achilles:} The Terms say that crosswise is the wisher's
business. So it's ours.

\emph{{[}They both draw in to the cloth, and the Tortoise's chair clacks
over as she comes. Bent over the table from either side, their two heads
throw two great shadows up the slope of the ceiling, which lean together
and apart as they talk, cross, slide over each other, and never once
keep still. The candle has begun to run: a drop of wax slips down its
side and stops there, gone white in the cold.{]}}

\textbf{Tortoise:} Then we'll mind it. First, a knob on the cloth.
\emph{{[}She takes a handful of small rings from the box, each with a
bead set in it.{]}} A ring at a cord's loose end, with that one cord and
nothing else, is a knob's setting. Red with its bead at rest is a dark
knob; red with its bead set halfway round, like the half-turn rings
you've met before, is a lit one.

\emph{{[}She threads a red end-ring on the loose left end of every cord
of the web, the flag cord too, each bead at rest. Every cord now begins
at a red end-ring: every knob dark, as the knobs stand on the
casket.{]}}

\textbf{Tortoise:} \emph{{[}shaking a new packet out onto the cloth:
short yellow tubes, each no longer than a fingernail{]}} And these are
new. Each is a yellow tube that slides along a cord, and a tube on a
cord is a collar turned crosswise.

\textbf{Achilles:} Mr Crab's flip-slot, sold by the packet.

\textbf{Tortoise:} The same trick, and much cheaper. Two tubes that meet
on a cord slide off together, since a collar turned twice is no turn at
all; and a tube slid out onto an end-ring changes its colour.

\textbf{Tortoise:} Every wish-knob crosswise and dark, then, and the
flag crosswise and dark as well.

\emph{{[}Each wish-cord begins at a red end-ring, bead at rest. She
threads a tube on near the end and slides it outward until it runs onto
the ring; the tube goes into the ring and is gone, and the ring is
green, its bead still at rest, the cord behind it just as it was. Down
the web, one after another, every wish-cord's end-ring goes from red to
green. The ties, the green rings on the first and fifth squares' cords
and the red rings on the flag cord are untouched. Last she does the same
to the flag's own end-ring: green, at rest.{]}}

\textbf{Tortoise:} A window reads the far end of its cord by the same
code. Read crosswise, a green end at rest shows dark, and a green end
with its bead halfway shows lit.

\textbf{Tortoise:} Two more things my rings do. Two rings of one colour
on one cord slide together into one, and their beads' turns add up. And
an end-ring pushed into a ring of the other colour vanishes into it: a
copy of it comes out on every other length of cord that leads from that
ring, and that ring is gone.

\textbf{Achilles:} Mr Crab's rules, from the evening of his koekpan.

\textbf{Tortoise:} The same rules; my rings keep them too. Now watch the
flag.

\emph{{[}The flag's green end-ring, bead at rest, sits at the start of
the flag cord. Along the cord lies the red ring where the first square's
tie comes down, and beyond it the red ring for the fifth square's tie.
She slides the green end-ring along the flag cord into the first red
ring. It goes in and does not come out. The red ring is gone from the
cord, and in its place two green end-rings, both at rest: one at the
foot of the tie, one on the flag cord just past where the red ring was.
The one at the foot of the tie she slides up it to the green ring on the
first square's cord; green meets green, and they run together into that
ring, whose bead stays exactly where it was. The flag's copy she slides
on into the red ring for the fifth square, and the same happens there:
that red ring gone, a green copy up the tie into the fifth square's
ring, adding nothing, and one last green copy left on the flag cord, at
rest.{]}}

\emph{{[}What lies on the cloth now: not a tie anywhere, and not a red
ring. On the first and fifth squares' cords the tied ring and the
end-ring have slid together into one green end-ring, at rest; every
other wish-cord begins at its green end-ring, at rest, as she set it.
Every wish-cord ends as it began. The flag cord ends as it began too,
green and at rest.{]}}

\textbf{Achilles:} Then the web is empty. There are no marks.

\textbf{Tortoise:} The slip has marks; you inked them.

\textbf{Achilles:} Then the flag gave nothing to anyone.

\textbf{Tortoise:} It had nothing to give. Now the flag crosswise and
lit.

\emph{{[}She lays the web again as it was: the ties down from the first
and fifth squares' green rings to their red rings on the flag cord,
every wish-cord beginning at a green end-ring at rest. At the start of
the flag cord she sets a green end-ring as before; but before she lets
go of it she turns its bead halfway round.{]}}

\emph{{[}She slides the flag's end-ring, its bead halfway round, into
the first red ring. The red ring is gone, and out come two copies, each
green, each with its bead halfway round: one at the foot of the tie, one
on the flag cord beyond. She slides the first up the tie into the green
ring on the first square's cord. Green runs into green, and the ring's
bead, which was at rest, now stands halfway round. The flag's copy she
slides on into the red ring for the fifth square: gone, and two copies
again, halfway round, one climbing that tie into the fifth square's
green ring, one left on the flag cord. Nothing has come to the bare
cords, and their end-rings sit at rest. What has changed: the first and
fifth squares' cords each carry a half-turn they did not have, exactly
where the ties came up. What has held: the flag cord, which ends as it
began, green with its bead halfway round. It has given its turn away
twice and still has it.{]}}

\textbf{Achilles:} The flag's turn has gone to the first and the fifth,
my marks, and the flag still has its own. It hasn't given anything away.

\textbf{Achilles:} \ldots Oh. It's your token, wedged up on its edge on
the porch in the spring. The turn had nothing on the flag to spend
itself on, so it jumped back up the tie onto the cord that asked.

\emph{{[}Up the lane the church clock begins to strike eight. At the
first stroke the tall glass jar under the shelf hums back, faint and
clear, and holds its hum a moment after the stroke has gone; the crocks
on either side of it stay dumb. So at every stroke, to the eighth, and
each stroke from the tower comes as full as the one before. Achilles
glances down at the jar, and back to the cloth.{]}}

\textbf{Tortoise:} And read crosswise, which windows gleam?

\textbf{Achilles:} The first and the fifth. Every mark comes back on its
own knob's cord.

\textbf{Tortoise:} One more box. \emph{{[}She unties both ties, lifts
the green and red rings off the cords, and resets the end-rings: green
at rest on every wish-cord, green with its bead halfway on the flag.{]}}
A box with no marks at all.

\emph{{[}She slides the flag's end-ring along the flag cord. There is no
red ring on it now to go into; it runs the whole length and comes to
rest at the far end, halfway round, as it set out. Nothing climbs to any
wish-cord, for there is no tie to climb. Every wish-cord lies as she set
it, green and at rest, and every bead on the cloth is where it was
before she pushed.{]}}

\textbf{Tortoise:} If every wish-window stays dark, the box has no
marks. If any wish-window gleams, it has.

\textbf{Achilles:} And the flag window?

\textbf{Tortoise:} Set crosswise and lit, it comes back lit every time,
marks or none. Leave it out of the test.

\textbf{Achilles:} Then set the casket as you set the cloth, every
collar crosswise, the wish-knobs dark and the flag lit, and the one wish
brings back every mark at once. \emph{{[}He looks up at the shelf.{]}}
But would the Genie count that as one wish?

\begin{center}\rule{0.5\linewidth}{0.5pt}\end{center}

\hypertarget{third-challenge-the-turned-web}{%
\section{\texorpdfstring{Third Challenge --- \emph{The Turned
Web}}{Third Challenge --- The Turned Web}}\label{third-challenge-the-turned-web}}

\textbf{Tortoise:} There's only one way to find out what the Genie
counts.

\emph{{[}Achilles reaches up and takes the lamp down from the shelf. It
sits in his two hands as it has sat on the shelf all evening: tarnished
brass, cold, its wick dry, never lit. He rubs its side once with his
thumb. The lamp does not warm or brighten, and nothing in it catches.
From its spout a thread of something paler than smoke rises, straight as
a plumb line in the still cold air, thickens, and stands; and where it
stands there is all at once a Genie, upright beside the table, with the
air of someone who has come about a form. The lamp in Achilles's hands
is exactly as it was, cold and dull and no lighter. What has changed is
that there is a third at the table. Achilles sets the lamp down beside
the candlestick.{]}}

\emph{{[}While the pale thread stood, the candle flame leaned over
towards the lamp as though something were drawing on the air, and every
shadow in the room leaned with it. Now it stands up straight again,
stiller than it has stood all evening, and not a shadow on the walls
stirs. On the wall by the shelf there is a third shadow beside the other
two, thinner than theirs, as though the candle can only half find
him.{]}}

\textbf{Genie:} The Office of Wishes. The lamp has been rubbed once,
which is correct.

\textbf{Genie:} One wish is one passage of the casket, whatever its
knobs and collars are set to. A wish with every collar crosswise is
still one wish.

\textbf{Achilles:} You've been listening.

\textbf{Genie:} The Office anticipates. A plain wish, every collar
upright, returns the knobs' settings in the wish-windows and one reading
in the flag window: the flag knob's setting, turned once for each marked
knob set lit. That is one answer.

\textbf{Achilles:} And setting the knobs ---

\textbf{Genie:} Is not a wish. Setting is free; the Office charges for
passage. There is one wish.

\textbf{Achilles:} Then I wish to know which knobs are marked.

\textbf{Genie:} Referred. \emph{{[}He inclines his head as though
listening up a stairwell, and straightens at once.{]}} The Meta-Genie
regrets. The marks are the casket's to tell, and the casket tells only
by passage. Nothing has been spent.

\emph{{[}A draught comes up the dark stairwell and in at the open door,
lifts the candle flame tall and thin, and lets it fall back to its own
height, all in the one instant. Nothing else in the room has had time to
move.{]}}

\textbf{Achilles:} That was quick.

\textbf{Genie:} Any request other than one passage of the casket as it
is now set is a meta-wish, and it goes up. The Meta-Genie referred yours
to the one above, who referred it higher, and so on all the way; the
answer came down all the way at once. The Office has no stairs.

\textbf{Achilles:} Then I wish for more wishes.

\textbf{Genie:} Referred. \emph{{[}The same tilt of the head, no longer
than a blink.{]}} The reply has come down: there is one wish. The
Meta-Genie adds that a wish for wishes is a wish about wishes, and is
not charged.

\textbf{Achilles:} Then at least let us look inside first.

\textbf{Genie:} Referred. \emph{{[}A blink.{]}} The lid lifts on a spent
casket and on no other. Your wish is untouched.

\emph{{[}The Tortoise laughs under her breath and tucks her claws into
her sleeves. Achilles, a little pink about the ears, blows on his
fingers and holds them to the candle.{]}}

\textbf{Tortoise:} Achilles, you'll talk your wish into the ground. Let
me show you the setting on the cloth first.

\textbf{Genie:} The Office is in no hurry. The Office is never in a
hurry.

\textbf{Achilles:} May I at least bring the casket down to the table? Or
is that a meta-wish too?

\textbf{Genie:} That is not a wish. That is furniture. Carry it level.

\emph{{[}Achilles lifts the casket down from the shelf in both hands and
sets it on the table beside the cloth, its knobs down the left, its
windows down the right, as the cords run on the cloth. It stands as it
has stood all evening: every knob down, every collar upright, every
window a plain pane, the seal whole.{]}}

\textbf{Tortoise:} \emph{{[}laying the slip's web out once more on the
cloth: a cord for each square, the flag cord below, the ties down from
the first and fifth squares' green rings to their red rings{]}} Your
slip's web again, set as the casket will be set. Every wish-knob dark,
\emph{{[}a red end-ring, bead at rest, at the start of each
wish-cord{]}} and the flag lit. \emph{{[}At the start of the flag cord,
a red end-ring with its bead halfway round.{]}} Now every collar
crosswise.

\emph{{[}She threads tubes onto the web. On every cord that runs from a
knob's end to a window's, one tube goes on just inside its end-ring, and
one near its far end, where the window would be. On each tie she threads
a pair of tubes, and a pair on the flag cord between its two red
rings.{]}}

\textbf{Tortoise:} A pair is no turn at all. It only puts the tubes
where I need them; and a ring with a tube on every cord that leads into
it changes colour, as Mr Crab's did when every leg went through his
slot.

\emph{{[}She slides the tubes along toward the rings. On the bare cords
the two tubes run in from either end, meet in the middle, and slide off
together; those cords lie bare again, their end-rings red and at rest,
as before. On the first square's cord one tube comes in from the knob's
end and one from the window's end, and the tie's pair parts, one tube
running up to the green ring and one down to the red; the same on the
fifth square's cord and its tie; on the flag cord, the tube from the
knob's end runs to the first red ring, the middle pair parts toward the
two red rings, and the tube from the far end runs to the second. Now
every ring in the web has a tube on every length of cord that leads into
it. Every ring changes colour at once, and its tubes come off it: the
green rings on the first and fifth squares' cords are red now, and the
red rings on the flag cord are green. Not a tube is left on the cloth.
The end-rings have not changed: red at rest on the wish-cords, red and
halfway round on the flag. Each tie still joins the same two cords it
joined before; but a tie runs from a green ring to a red one, and now
the green is on the flag.{]}}

\textbf{Achilles:} The ties run the other way. They came down from my
marks to the flag; now they go up from the flag to my marks.

\textbf{Achilles:} So what you're saying is that the collars have moved
the ties.

\textbf{Tortoise:} Nothing has moved. Every tie joins the cords it
always joined; only the rings have turned.

\textbf{Achilles:} And last time you slid the tubes outward, onto the
end-rings.

\textbf{Tortoise:} Outward or inward, it's the same collar, and the two
ways agree.

\emph{{[}Achilles looks up from the cloth, past the candle, to the print
on the wall facing the window: Escher's \textup{Belvedere}.{]}}

\textbf{Achilles:} A pillared lookout whose upper storey is set
crosswise on the lower; a ladder stands inside it and leans against the
outside. It has been there all along.

\emph{{[}The Tortoise turns in her chair to look where he is looking,
and her shadow, which has lain over that wall since the candle came up,
slides aside off the frame. For the first time this evening the print
hangs clear in the candlelight, and seems no further off than the
candle.{]}}

\textbf{Tortoise:} Push the flag in.

\emph{{[}The flag's end-ring, red with its bead halfway round, sits at
the start of the flag cord. Along the cord lies the green ring whose tie
climbs to the first square's cord, and beyond it the green ring for the
fifth. Achilles slides the end-ring along into the first green ring. The
ring is gone, and out come two copies, each red with its bead halfway
round: one at the foot of the tie, one on the flag cord beyond. He sends
the first up the tie into the red ring on the first square's cord. Red
runs into red: the cord's end-ring at rest, the ring, and the copy all
slide together into one red end-ring with its bead halfway round. The
flag's copy he slides on into the next green ring, and the same happens
there, up the tie to the fifth square's cord. The last copy comes to
rest at the flag cord's far end, red, halfway round. The bare cords have
not been touched: red and at rest from end to end.{]}}

\textbf{Tortoise:} No tubes left, so read them plainly: red at rest is
dark, and red halfway round is lit.

\textbf{Achilles:} Lit at the first and the fifth, my marks, and the
flag lit. One push, and every mark read at once.

\textbf{Achilles:} \emph{{[}turning to the casket{]}} Now the casket.
Every collar crosswise.

\emph{{[}The casket stands as it has all evening: every knob down, every
collar at the upright stroke, every window plain. Achilles goes down the
left face, turning each wish-knob's collar till he feels it stop at the
cross; then down the right face, each wish-window's collar; then the
flag window's.{]}}

\textbf{Achilles:} \emph{{[}his fingers on the flag knob's collar{]}} Or
leave this one upright, as the Terms describe it?

\textbf{Tortoise:} That would gamble the only wish.

\emph{{[}He turns the flag knob's collar to the cross, and then turns
the flag knob up, to lit. Every wish-knob he leaves down. What has
changed: every collar, now at the cross, and the flag knob, now up. What
has held: the wish-knobs, dark; the plain windows; the seal; the
lid.{]}}

\textbf{Achilles:} Is that one wish?

\textbf{Genie:} Every collar crosswise, every knob as set: one passage.
One wish.

\textbf{Achilles:} Then I wish it.

\textbf{Genie:} Granted.

\emph{{[}The casket as it stands: every collar at the cross, the flag
knob up, the wish-knobs down, every window a plain pane, the seal whole.
At the word, nothing on it moves: no knob, no collar, not the lid.
Behind the second wish-window a small gleam comes up and stays. Behind
the fourth, another. Behind the flag window, one more. Behind the first,
the third, the fifth and the sixth wish-windows the panes stay plain.
Achilles touches the flag knob, and it will not turn: every knob is
locked where he set it. The seal on the lid has parted of itself. What
has changed: the gleams, the lock, the seal. What has held: every
setting, every collar, and the lid, still down.{]}}

\textbf{Achilles:} The second and the fourth, and the flag. Not my marks
at all. The casket's own.

\emph{{[}Achilles's eye goes from the casket to the window. Low on the
pane, where the frost lies thickest, a pair of initials stands out white
that nobody could have made out on the bare glass: scratched into it
long ago with a diamond, and picked out by the frost, the whole of them
at once.{]}}

\textbf{Tortoise:} Your slip was only the rehearsal. The casket had its
own secret all along.

\textbf{Genie:} The casket is spent. A secret read is no longer a
secret, and the lid may be lifted.

\emph{{[}The casket as it stands: spent, its knobs locked, the gleams
behind the second and fourth wish-windows and the flag window, the seal
parted. Achilles lifts the lid on its hinge. The knobs do not move, nor
the collars, nor the gleams; only the lid swings back, and the candle
shines down into the casket.{]}}

\emph{{[}Inside, from the top of the column down. The cord from the
first knob runs straight across to the first window, knotted at the knob
and knotted at the window, with nothing on it. The cord from the second
knob runs in only as far as a green ring and is knotted there; from that
ring a cord runs on across to the second window, knotted at the ring and
at the window; and from the same ring a tie drops away down the casket,
knotted to the ring above and, below, to the upper red ring, low in the
casket near the flag knob, passing behind every cord beneath it without
touching one. The cord from the third knob runs straight across to the
third window. The fourth knob is like the second: its cord runs to a
green ring and is knotted there; a cord runs from that ring to the
fourth window; and a tie runs down from the ring, behind the cords below
it, to the lower red ring, beneath the upper. The cord from the fifth
knob runs straight across to the fifth window, and the cord from the
sixth knob straight across to the sixth. From the flag knob a cord runs
to the upper red ring and is knotted there; another cord, knotted at
both its ends, runs from the upper red ring to the lower; and from the
lower red ring a last cord runs to the flag window. Every length is
knotted at both its ends, and no cord passes through a ring. What has
changed is only that the inside can be seen. What has held is everything
the windows say, and every gleam can now be followed back, length by
length, to the knob it came from.{]}}

\emph{{[}A smell comes up out of the open casket, of old wood and
something like cloves, shut in since long before the house was the
Tortoise's. Achilles and the Tortoise lean in over it together; the
Genie, standing a little back, does not trouble to look.{]}}

\textbf{Achilles:} Green rings at the second and the fourth knobs, red
rings low down, where the flag's lengths meet, and the ties coming down
from the marks. It's your web as you first laid it, not turned at all.
So the turning was never inside.

\textbf{Tortoise:} It was all in the collars.

\textbf{Genie:} The wish has been spent in accordance with the Terms.
The Office does not collect spent boxes.

\emph{{[}The lamp lies where Achilles set it down, beside the
candlestick, cold, as it has been since he rubbed it. The Genie thins
from the feet upward into the pale thread he came as; the thread leans,
runs down into the lamp's spout, and is gone. The lamp does not stir or
warm. The casket stays open on the table, its windows still showing what
they showed, and nothing else in the room has moved. There are two at
the table again.{]}}

\textbf{Tortoise:} \emph{{[}taking up the lamp and setting it back on
its shelf{]}} Back where it lives.

\emph{{[}Up on its shelf, out of the candle's reach, the lamp is only
brass again, one dull shape beside the clean patch in the dust where the
casket stood. The thin third shadow has gone from the wall by the shelf;
there are two again, and the Tortoise's chair clacks as she sits back
down in it.{]}}

\textbf{Achilles:} \emph{{[}looking into the open casket{]}} Every box
of marks, read in one wish. Is there any secret a wish can't read like
that?

\begin{center}\rule{0.5\linewidth}{0.5pt}\end{center}

\hypertarget{fourth-challenge-the-pairs-box}{%
\section{\texorpdfstring{Fourth Challenge --- \emph{The Pairs
Box}}{Fourth Challenge --- The Pairs Box}}\label{fourth-challenge-the-pairs-box}}

\textbf{Tortoise:} Some. Here is one. \emph{{[}She clears the web from
the cloth.{]}} This is a pairs box, drawn. Its settings come in pairs
that answer alike, and what it hides is where the two of a pair differ;
in place of a flag it has a pair of answer-cords.

\emph{{[}Across the top of the cloth she lays a short row of wish-cords,
each beginning at a red end-ring at rest, and below them, where the flag
cord would lie, a pair of answer-cords, each beginning at a red end-ring
at rest. On the cords for the first two places she threads a green ring
each, and ties both down to the upper answer-cord, each tie to a red
ring of its own; on the third place's cord a green ring, tied down to a
red ring on the lower answer-cord. Past its last red ring, each
answer-cord turns down and runs to the cloth's near edge.{]}}

\textbf{Tortoise:} Which settings answer alike?

\textbf{Achilles:} \emph{{[}touching the ties as he goes{]}} Light the
first place and the upper answer turns over; light the second instead
and it turns over just the same. Light both and it turns twice, and
comes back as if I'd lit neither. The third place turns the lower
answer, and nothing else touches that. So two settings answer alike
exactly when they differ at the first two places together, and nowhere
else.

\emph{{[}Achilles sits back and tucks his hands under his arms; his
breath still shows in the candlelight, and the cold has got into his
fingers again. The Tortoise's chair clacks as she leans in over the
cloth.{]}}

\textbf{Tortoise:} Then the secret is the first two places. You can see
the ties; a wisher couldn't. A crosswise wish on this box would light
its wish-windows in some pattern, and we'll call the pattern a clue.

\textbf{Tortoise:} Nobody reads the answer-cords: a wisher who wants a
clue leaves the answer windows alone. So I'll fold them away. Whatever
is under the fold is never looked at.

\emph{{[}The two answer-cords run from their end-rings, through their
red rings, and down to the cloth's near edge, where their far ends lie
loose. The Tortoise takes the near edge of the cloth and folds it up and
back over them, a hand's breadth deep. The loose ends are gone under the
fold; everything above it lies where it lay. Nothing has been cut, and
nothing under the fold will be read: the answer-cords now run on under
the folded edge and stop where nobody looks.{]}}

\emph{{[}On the upper answer-cord, the dark end-ring, the red ring for
the first place's tie and the red ring for the second's all lie on one
cord, all red. She slides them together, and they run into one red ring,
its bead at rest. The green rings on the first two places' cords are
tied now, both of them, to that one red ring, and its own cord runs on
down under the fold. On the lower answer-cord, the end-ring and the red
ring run together in the same way: one red ring, tied to the third
place's green ring, its cord running on under the fold. No ring is
worked past the fold. What has changed: each answer-cord's start and its
red rings have gone into one ring. What has held: the ties, the
wish-cords, and the fold.{]}}

\textbf{Tortoise:} A red ring counts the ties that come into it, odd or
even, as it always has. With its own cord under the fold, all it can do
is hold every clue to a check: a clue lights the secret's places an even
number of times, and here that means the first two places both lit or
neither. The third place's ring has one tie and nobody reading it, so
the third place is free.

\textbf{Achilles:} Then let me rub the lamp and fetch a real clue.

\textbf{Tortoise:} The Genie has gone, and so has our wish. We'll
suppose the clues instead. Suppose one wish had lit only the third
window, and another the first two.

\textbf{Achilles:} The first lit the third window alone, so the third
can't be a secret place; they would share one lit place, and one is odd.
The second lit the first two, so the secret holds both of them or
neither. Neither is no secret at all. Both, then: the first two places,
just as the ties said.

\emph{{[}Up the lane the church clock gives a single stroke. It gives
the same one stroke at every half hour, and says the half, and not of
which hour; but anyone in the house who heard it strike eight needs no
telling.{]}}

\textbf{Tortoise:} And a wish that lit nothing at all?

\textbf{Achilles:} Would pass every check there is, and tell me nothing.

\textbf{Tortoise:} Two clues, then: two wishes, had the Genie granted
them.

\textbf{Achilles:} What does one wish buy?

\textbf{Tortoise:} One answer, and only the one.

\textbf{Achilles:} Then the rest must be bought one at a time.

\emph{{[}They have both leaned in over the cloth, and on the slope of
the ceiling their two shadows, which came up the stairs apart and have
leaned and crossed all evening, have run together over the table into
one, and hold there. Every shadow in the room stands a little higher on
the walls than it did when the candle was first set down.{]}}

\textbf{Achilles:} \emph{{[}looking from the cloth to the open
casket{]}} So that's the evening. A plain wish buys one answer. A
crosswise wish on a box of marks buys the whole secret, read straight
off its ties. On a pairs box it buys one clue: a line through the
secret, not the secret. And no turning of the collars buys what the box
doesn't hold.

\textbf{Tortoise:} And some questions are no quicker for any wish,
however it's set; that was proved by people who counted wishes, not by
drawing.

\textbf{Achilles:} \emph{{[}looking at the candle{]}} The candle's lower
than I thought.

\textbf{Tortoise:} It always burns faster up here. The cold makes it
gutter.

\textbf{Achilles:} \emph{{[}going to the window{]}} And the frost is
thicker on this pane than on yours downstairs. It's come right over the
middle.

\textbf{Tortoise:} Nobody breathes on it up here.

\textbf{Achilles:} \emph{{[}coming back to the table, and looking into
the open casket{]}} Its windows are still showing. The second, the
fourth, and the flag. Will they keep?

\textbf{Tortoise:} Till the house falls down, I expect. It's spent; it
has nothing left to do but show them.

\textbf{Achilles:} \emph{{[}unwinding his scarf and putting it round the
Tortoise's shoulders{]}} You'll catch cold up here.

\textbf{Tortoise:} I never catch anything. \emph{{[}She settles the
scarf.{]}} Thank you.

\textbf{Achilles:} Shall we go down?

\textbf{Tortoise:} In a little while. \emph{{[}She draws her chair
nearer the candle.{]}}

\textbf{Achilles:} If someone hands you a box and won't say what is
inside it, how much can one question tell you?

\textbf{Tortoise:} \emph{{[}taking up the candlestick{]}} That depends
on the box, and on how the question is set. Mind the loose tread going
down; I still don't remember which it is.

\emph{{[}She does not get up at once. The candle stands low in its stick
between her two claws, guttered down to a small flame, and Achilles
leans on the table beside her with his hands cupped to it. Every shadow
in the room has climbed as the candle burned down, and they all stand
tall and still now: the chair-backs high on the walls, the two of theirs
together on the slope of the ceiling. The casket stands open on the
table where Achilles set it. On the shelf by the window the lamp is only
brass, out of the light. On the wall facing the window, past the
Tortoise's shoulder, the print hangs half in her shadow again at the
edge of the light, a print on a cold wall. Then she stands, and the
candle goes up with her, and every shadow in the room swings round and
shortens and follows it out of the door. The stair takes them down its
one turn; halfway, a tread gives under Achilles's foot, as it did coming
up, and neither of them says which. Behind them the top room is dark and
still: the lamp on its shelf, the casket open on the table, the print on
the wall that nobody can see now, and the frost going on with its work
over the whole of the pane. Up the stairwell, into the dark they have
left, climbs a little of the parlour's warmth.{]}}

\hypertarget{chapter-13-asking-the-oracle}{%
\chapter{Chapter 13: Asking the
Oracle}\label{chapter-13-asking-the-oracle}}

\emph{--- Minuet for One Wish ---}

\emph{Part IV: Algorithms}

\emph{Escher: ``Belvedere'' --- a pillared lookout whose upper storey is
set crosswise on the lower; a ladder stands inside it and leans against
the outside}

\begin{center}\rule{0.5\linewidth}{0.5pt}\end{center}

If someone hands you a sealed box and will not say what is inside it,
how much can one question tell you?

Everyone has played the game. A friend thinks of something, and you may
ask questions that are answered yes or no. Each answer is one bit of
news, and a careful player learns to make every bit count: a question
that splits the possibilities in half is worth asking, and a question
whose answer you could have guessed is wasted. A secret of six yes-or-no
facts needs six good questions, and no cleverness in the wording gets it
done in five. That feels less like a rule of the game than a fact about
the world. A question is a thing you put in, an answer is a thing you
get back, and one answer can carry only so much.

Now change the game in one respect. The box is not a friend but a
device, and it will do exactly one thing, as often as you like: take in
a question and hand back an answer, with no opinion about what you are
doing. You pay for each use. What you would like to know is not the
answer to any single question but something about the box as a whole:
whether it treats every question alike, say, or which hidden setting
inside it decides its answers. Asked the ordinary way, the device gives
one answer per use, and the counting of the last paragraph applies in
full.

Chapter 5 found a crack in that counting. A single use of a quantum box,
asked in a superposition and read after one more turn of the wires, told
whether a function gave the same answer at both of its inputs, a fact
that no single ordinary question can settle. This chapter follows the
crack as far as it goes. We will meet a box whose whole secret, all its
hidden bits at once, comes out of one use. We will draw that box and
find that it is nothing but a web of simple gates, and that the secret
is written in the web's wiring for anyone who looks at the picture the
right way round. Then we will meet a box that one use cannot empty, and
see exactly what one use does buy from it.

We will also be careful about what we have not shown. A picture can
prove that a certain question returns a certain answer. Whether some
other method, classical or quantum, could have done as well with fewer
questions is a different kind of claim, a claim about all possible
methods at once, and no picture in this chapter proves one. Where such a
claim is needed we prove it by counting on the blackboard, or we state
it and say whose it is.

\begin{center}\rule{0.5\linewidth}{0.5pt}\end{center}

\hypertarget{one-wish-one-answer}{%
\section{§13.1 One Wish, One Answer}\label{one-wish-one-answer}}

Chapter 12 ended with a promise: ``In Chapter 13 someone else hands us a
box and will not say what is inside it. We may ask it a question, once.
The box turns out to be a web of CNOTs, and yellow boxes turn its
colours over. Its secret can be read straight off the picture, if we ask
in the right way.'' Before any of that we have to say what a question
is, and what it costs.

\textbf{On the blackboard.} The box computes a function \(f\) that takes
a string of \(n\) bits to one bit, \(f:\{0,1\}^n\to\{0,1\}\). We cannot
look inside it. What we can do is Chapter 5's oracle: the reversible box
\(U_f\) that acts on \(n\) query wires and one answer wire by

\[U_f\,\ket x\ket y\;=\;\ket x\ket{y\oplus f(x)},\]

where \(x\) is an \(n\)-bit string, \(y\) is a bit, and \(\oplus\) is
addition mod \(2\). In this chapter we call the \(n\) query wires the
\textbf{rods} and the answer wire the \textbf{flag}. Ket order is top
wire first, so the rods come first and the flag last. One use of \(U_f\)
is one \textbf{query}. This is the operation Cleve, Ekert, Macchiavello
and Mosca call the \(f\)-controlled-NOT, their Eq. (2.1) for one bit and
Eq. (3.3) for \(n\) bits. Bennett, Bernstein, Brassard and Vazirani (we
will call them BBBV) write the same thing for a quantum Turing machine:
a query changes the contents of a query tape from \(\ket{x\circ b}\) to
\(\ket{x\circ(b\oplus A(x))}\), where \(\circ\) joins two strings and
\(A\) is their name for the box.

A query costs one unit, and it costs one unit whatever is on the rods.
BBBV describe an oracle as ``a special subroutine call whose invocation
only costs unit time'', and they allow the query tape to hold an
arbitrary superposition of strings when the call is made. So a query on
a superposition of all \(2^n\) strings is still one query. This is the
only resource the chapter counts.

Here is a small box to make it concrete. Take \(n=2\) and let \(f\)
return the first bit, \(f(x_1x_2)=x_1\). Then
\(U_f\ket{10}\ket0=\ket{10}\ket1\): the first rod is \(1\), so the flag
is flipped. Ask again with the same rods, and
\(U_f\ket{10}\ket1=\ket{10}\ket0\): the flag is flipped back. With the
rods at \(01\) nothing happens either time. BBBV make the same
observation in general: if the answer bit starts at \(0\) it ends
holding \(A(x)\), and ``if the target bit is already in state
\(\ket{A(x)}\), calling the oracle will reset it to \(\ket0\)''. They
say this ability to uncompute ``will often prove essential'' for
interference.

\begin{quote}
\textbf{Result 13.1 (a query undoes itself).} For every
\(f:\{0,1\}^n\to\{0,1\}\), \(U_f\) permutes the \(2^{n+1}\) basis
states, so it is unitary, and \(U_f^2=I\).
\end{quote}

\emph{Proof.} Apply \(U_f\) twice to a basis state:
\(\ket x\ket y\mapsto\ket x\ket{y\oplus f(x)}\mapsto\ket x\ket{y\oplus f(x)\oplus f(x)}=\ket x\ket y\),
since \(f(x)\oplus f(x)=0\). So \(U_f^2=I\) on a basis, hence
everywhere. A map that is its own inverse on the basis states is a
bijection of them, so the matrix of \(U_f\) has exactly one \(1\) in
each row and each column and \(0\) elsewhere. Its columns are distinct
basis vectors, which are orthonormal, so \(U_f\) is unitary.
\(\;\blacksquare\) BBBV note the same fact for their oracles: the
transformation that evaluates a classical Boolean function ``is an
involution''. See \texttt{Code/Ch13/PlainWishes.py}.

\textbf{In the sketchbook.} The tall box \(U_f\) of Chapter 5 (§5.3, row
37) is laid across the rods and the flag, drawn wide so that it is never
mistaken for a wire. Drawn here for three rods, Result 13.1 says that
two tall boxes in a row are bare wires:

\[\begin{tikzpicture}[sd, x=1em, y=1em]
  \foreach \y in {4.5,3,1.5,0} { \draw[wire] (0,\y) -- (7,\y); }
  \node[sdbox, minimum height=6.2em] at (2,2.25) {$U_f$};
  \node[sdbox, minimum height=6.2em] at (5,2.25) {$U_f$};
\end{tikzpicture}
\;=\;
\begin{tikzpicture}[sd, x=1em, y=1em]
  \foreach \y in {4.5,3,1.5,0} { \draw[wire] (0,\y) -- (3,\y); }
\end{tikzpicture}\]

exactly. The picture says nothing more, because the tall box hides
everything inside it. That is the point of an oracle, and it is about to
change.

\textbf{\emph{Why this definition.}} Why count uses of the box, and not
time, or bits of output? Because the box is sealed. Whatever happens
inside it is not ours to count, and the one thing we control is how
often we consult it. And why the form \(y\oplus f(x)\), rather than
simply writing \(f(x)\) on the flag? Because writing over \(y\) would
destroy it, and a process that destroys information cannot be unitary.
Adding mod \(2\) keeps the flag's old value recoverable, which is
exactly Result 13.1.

\hypertarget{a-box-of-marks}{%
\subsection{A box of marks}\label{a-box-of-marks}}

A single definite query, \(\ket x\ket y\) with \(x\) a fixed string,
returns one bit of news: the flag comes back as \(y\oplus f(x)\), and
the rods come back unchanged. Which boxes have secrets worth reading?
The simplest interesting ones are the \textbf{linear} functions. For a
string \(a\in\{0,1\}^n\) and a bit \(b\), let

\[f(x)\;=\;(a\cdot x)\oplus b,\qquad a\cdot x\;=\;a_1x_1\oplus a_2x_2\oplus\cdots\oplus a_nx_n .\]

The product \(a\cdot x\) is the dot product \textbf{taken mod \(2\)}:
multiply the bits place by place and add the results mod \(2\). This is
Cleve et al.'s Eq. (3.7), with their dot product (3.2). The string \(a\)
is the box's \textbf{secret}, and \(b\) is a constant bit. Call rod
\(i\) \textbf{marked} if \(a_i=1\), and write \(w\) for the number of
marks. So \(f(x)\) is the parity of \(x\) on the marked rods, flipped
once more if \(b=1\).

Take \(n=3\), \(a=101\), \(b=0\). Then \(f(x)=x_1\oplus x_3\), and
\(U_f\ket{x_1x_2x_3}\ket y=\ket{x_1x_2x_3}\ket{y\oplus x_1\oplus x_3}\).
Read that as two operations: flip the flag if rod \(1\) is \(1\), then
flip it if rod \(3\) is \(1\). Each is a CNOT (Chapter 4) with its
control on a marked rod and its target on the flag. Rod \(2\) is not
touched at all.

\begin{quote}
\textbf{Result 13.2 (the web lemma).} If \(f(x)=(a\cdot x)\oplus b\),
then
\[U_f\;=\;X_{\text{flag}}^{\,b}\prod_{i:\,a_i=1}\mathrm{CNOT}_{i\to\text{flag}},\]
where \(\mathrm{CNOT}_{i\to\text{flag}}\) flips the flag when rod \(i\)
is \(1\) and \(X_{\text{flag}}\) flips the flag. The factors commute, so
their order does not matter. In pictures, \(U_f\) is \(2^{w/2}\) times a
\textbf{web of ties}: a green dot on each marked rod, tied to its own
red dot on the flag, the unmarked rods straight, and a red \(\pi\) dot
on the flag if \(b=1\). Two webs in a row are \(2^{-w}\) times bare
wires.
\end{quote}

\emph{Proof.} On a basis state, each \(\mathrm{CNOT}_{i\to\text{flag}}\)
adds \(x_i\) to the flag, so together they add
\(\bigoplus_{a_i=1}x_i=a\cdot x\), and \(X_{\text{flag}}^{\,b}\) adds
\(b\). The flag ends as \(y\oplus(a\cdot x)\oplus b=y\oplus f(x)\) and
the rods are untouched, which is \(U_f\ket x\ket y\). Two linear maps
that agree on a basis are equal. Each of these operations only adds
something to the flag, and additions mod \(2\) commute. For the picture:
a bare green--red pair is \(\mathrm{CNOT}/\sqrt2\) (row 64), and a red
\(\pi\) dot with one leg in and one out is \(X\) exactly (its matrix,
§8.1), so the web is \((\tfrac1{\sqrt2})^{w}\,U_f\). \(\;\blacksquare\)
See \texttt{Code/Ch13/PlainWishes.py}.

\textbf{In the sketchbook.} Here is the web of \(a=101\), \(b=0\): rods
\(1\), \(2\), \(3\) from the top, then the flag.

\[U_f\;=\;2\;\,
\begin{tikzpicture}[sd, x=1.2em, y=1.1em,
  zsp/.style={circle,draw,line width=0.9pt,fill=green!55,minimum size=1.0em,inner sep=0pt},
  xsp/.style={circle,draw,line width=0.9pt,fill=red!75,minimum size=1.0em,inner sep=0pt}]
  \draw[wire] (1.5,4.2) -- (1.5,0);
  \draw[wire] (3.2,1.4) -- (3.2,0);
  \foreach \y in {4.2,2.8,1.4,0} { \draw[wire] (0,\y) -- (4.7,\y); }
  \node[zsp] at (1.5,4.2) {};
  \node[zsp] at (3.2,1.4) {};
  \node[xsp] at (1.5,0) {};
  \node[xsp] at (3.2,0) {};
  \node[sdlabel,anchor=west] at (4.8,4.2) {$1$};
  \node[sdlabel,anchor=west] at (4.8,2.8) {$2$};
  \node[sdlabel,anchor=west] at (4.8,1.4) {$3$};
  \node[sdlabel,anchor=west] at (4.8,0) {flag};
\end{tikzpicture}\]

The tie from rod \(1\) runs down past rods \(2\) and \(3\) without
touching them: a crossing neither joins nor breaks a line. The factor
\(2=2^{w/2}\) with \(w=2\) is drawn because, as in Chapters 9 to 12,
every picture equation in this chapter is exact, with its scalar written
in.

Now ask twice. Put two webs in a row. On each marked rod two green dots
share a wire, and fusion (Theorem 8.1, row 57) makes them one. On the
flag all the red dots share the flag's wire, and fusion makes them one
red dot. Each marked rod's green dot is now tied to that red dot by two
separate ties, and the Hopf law (Theorem 8.6, row 62) removes a double
tie at the price of \(\tfrac12\). What is left on every wire is a
phase-free dot with one leg in and one out, which is a plain wire
(§8.1):

\[\begin{tikzpicture}[sd, x=1.2em, y=1.1em,
  zsp/.style={circle,draw,line width=0.9pt,fill=green!55,minimum size=1.0em,inner sep=0pt},
  xsp/.style={circle,draw,line width=0.9pt,fill=red!75,minimum size=1.0em,inner sep=0pt}]
  \draw[wire] (1.2,4.2) -- (1.2,0);
  \draw[wire] (2.6,1.4) -- (2.6,0);
  \draw[wire] (4.0,4.2) -- (4.0,0);
  \draw[wire] (5.4,1.4) -- (5.4,0);
  \foreach \y in {4.2,2.8,1.4,0} { \draw[wire] (0,\y) -- (6.6,\y); }
  \node[zsp] at (1.2,4.2) {}; \node[zsp] at (2.6,1.4) {};
  \node[xsp] at (1.2,0) {}; \node[xsp] at (2.6,0) {};
  \node[zsp] at (4.0,4.2) {}; \node[zsp] at (5.4,1.4) {};
  \node[xsp] at (4.0,0) {}; \node[xsp] at (5.4,0) {};
\end{tikzpicture}
\;\overset{\text{fusion}}{=}\;
\begin{tikzpicture}[sd, x=1.2em, y=1.1em,
  zsp/.style={circle,draw,line width=0.9pt,fill=green!55,minimum size=1.0em,inner sep=0pt},
  xsp/.style={circle,draw,line width=0.9pt,fill=red!75,minimum size=1.0em,inner sep=0pt}]
  \draw[wire] (1.2,4.2) to[bend right=12] (2.4,0);
  \draw[wire] (1.2,4.2) to[bend left=22] (2.4,0);
  \draw[wire] (3.6,1.4) to[bend right=25] (2.4,0);
  \draw[wire] (3.6,1.4) to[bend left=25] (2.4,0);
  \foreach \y in {4.2,2.8,1.4,0} { \draw[wire] (0,\y) -- (4.8,\y); }
  \node[zsp] at (1.2,4.2) {}; \node[zsp] at (3.6,1.4) {};
  \node[xsp] at (2.4,0) {};
\end{tikzpicture}
\;\overset{\text{Hopf}}{=}\;\tfrac14\;
\begin{tikzpicture}[sd, x=1.2em, y=1.1em]
  \foreach \y in {4.2,2.8,1.4,0} { \draw[wire] (0,\y) -- (2.4,\y); }
\end{tikzpicture}\]

With \(U_f=2\times\)web, the left side is \(\tfrac14U_f^2\), so the
picture has proved \(U_f^2=I\) again, for this box. For \(b=1\) each web
carries a red \(\pi\) dot on the flag. They fuse into the flag's red dot
as well, and their phases add to \(2\pi\), which is no turn at all. See
\texttt{Code/Ch13/PlainWishes.py}, which checks the web against the
matrix \(U_f\) and derives the bare wires.

\textbf{\emph{Why this definition.}} Why draw a CNOT for every mark,
when one tall box would do? Because the tall box hides the secret and
the web displays it. In the web, the marks \emph{are} the ties. Every
rule of Chapter 8 acts on dots and ties, so once the box is drawn as a
web the rules can reach its secret. Most boxes are not webs. A web
always computes some \((a\cdot x)\oplus b\), and there are only
\(2^{n+1}\) of those among the \(2^{2^n}\) functions on \(n\) bits: for
\(n=2\), eight of the sixteen. For every other box the tall box is all
we have.

\hypertarget{how-many-plain-questions}{%
\subsection{How many plain questions?}\label{how-many-plain-questions}}

With definite queries, how many are needed to learn \(a\)? Suppose first
that \(b\) is known; say \(b=0\). Asking at the unit string \(e_i\),
which has a single \(1\) in place \(i\), returns \(a\cdot e_i=a_i\). So
\(n\) queries read the whole secret, one rod at a time. Can a cleverer
choice do it in fewer? Take \(n=3\) and try the two queries \(110\) and
\(011\). They return \(a_1\oplus a_2\) and \(a_2\oplus a_3\). Now
compare two secrets that differ by \(v=111\), for instance \(a=100\) and
\(a\oplus v=011\). The first gives \(100\cdot110=1\) and
\(100\cdot011=0\); the second gives \(011\cdot110=1\) and
\(011\cdot011=0\), again \(1\) and \(0\). The two secrets answer both
questions alike, because \(v\cdot110=0\) and \(v\cdot011=0\). Two
questions cannot tell them apart.

\begin{quote}
\textbf{Result 13.3 (the classical baseline).} Use only definite
queries, chosen one after another in any way that may depend on the
earlier answers, and suppose the method must name \(a\) with certainty.
(i) If \(b\) is known, every such method uses at least \(n\) queries,
and \(n\) suffice. (ii) If \(b\) is unknown, every such method uses at
least \(n+1\) queries, and \(n+1\) suffice.
\end{quote}

\emph{Proof.} (i) Run the method against the secret \(a\), and let
\(q_1,\dots,q_k\) be the query strings it uses, with \(k<n\). The \(k\)
equations \(v\cdot q_j=0\) in the \(n\) unknown bits of \(v\) have a
solution \(v\neq0\) (fewer equations than unknowns; see the Dirac's
Blackboard below). Now run the same method against \(a\oplus v\). Its
first query is \(q_1\), since nothing has been answered yet, and the
answer is \((a\oplus v)\cdot q_1\oplus b=(a\cdot q_1)\oplus b\), the
same as before. So its second query is \(q_2\), with the same answer,
and so on to the end. The method sees exactly the same answers for \(a\)
and for \(a\oplus v\neq a\), so it names the same secret for both and is
wrong for one of them. The unit strings show that \(n\) suffice. (ii)
Now suppose \(k\le n\). The \(k\) equations \(v\cdot q_j\oplus c=0\) in
the \(n+1\) unknowns \((v,c)\) have a nonzero solution. If \(v\) were
\(0\), then \(c\) would be \(1\) and the equations would read \(1=0\);
so \(v\neq0\). The secrets \((a,b)\) and \((a\oplus v,\,b\oplus c)\)
give the same answer to every query, and the argument of (i) finishes.
Asking first at \(x=0^n\), the all-zero string, returns \(b\), and then
the \(n\) unit strings return every \(a_i\oplus b\): \(n+1\) suffice.
\(\;\blacksquare\)

A method that tosses coins but must still be certain is certain on every
run, and each run is a fixed method of the kind above, so the same
bounds hold for it. Cleve et al.~give the reason for (i) in one
sentence: \(a\) ``contains \(n\) bits of information and each classical
evaluation of \(f\) yields a single bit''.
\texttt{Code/Ch13/PlainWishes.py} checks both parts exhaustively for
\(n=4\). There is no picture on this page. No rewrite proves a lower
bound, and this one is proved by counting.

\vskip 0pt plus 4\baselineskip\penalty-100\vskip 0pt plus
-4\baselineskip

\begin{quote}
\textbf{Dirac's Blackboard.}\nopagebreak

\emph{Fewer equations than unknowns, mod 2.} Why must \(k<n\) equations
\(v\cdot q_j=0\) have a solution \(v\neq0\)? The same way as over the
real numbers. Take the queries \(110\) and \(011\): the equations are
\(v_1\oplus v_2=0\) and \(v_2\oplus v_3=0\). The first says \(v_1=v_2\),
the second \(v_3=v_2\), and \(v_2\) is free. Setting it to \(1\) gives
\(v=111\). In general, use each equation in turn to express one unknown
through the others, adding it mod \(2\) to the later equations to remove
that unknown from them; an equation that has become \(0=0\) is dropped.
Each equation removes at most one unknown. With \(k<n\) equations at
least one unknown is never removed, and we may set it to \(1\) and read
the others off. Addition mod \(2\) is its own subtraction, so nothing in
the elimination needs anything but \(\oplus\).
\end{quote}

\begin{center}\rule{0.5\linewidth}{0.5pt}\end{center}

\hypertarget{the-mark-on-the-rods}{%
\section{§13.2 The Mark on the Rods}\label{the-mark-on-the-rods}}

A definite query moves a mark onto the flag: the flag is flipped if
\(f(x)=1\). Chapter 5 showed the other way round. With the flag held in
the difference state \(\ket-=\tfrac1{\sqrt2}(\ket0-\ket1)\), the flag
does not change, and the sign of \(f(x)\) lands on the rods instead:

\[U_f\,\ket x\ket-\;=\;(-1)^{f(x)}\,\ket x\ket-\]

(row 37; Cleve et al.'s Eq. (2.2), which they write without the
normalization). For a box of marks there is more to see, because the box
is a web, and each tie can be watched on its own.

\textbf{On the blackboard.} Start with the trap. Put the flag in the
other crosswise state, \(\ket+=\tfrac1{\sqrt2}(\ket0+\ket1)\), and one
tie: a CNOT from a rod in any state \(\ket\psi\) to the flag. A CNOT
applies \(X\) to the flag when the rod is \(1\), and \(X\ket+=\ket+\).
So \(\mathrm{CNOT}\,(\ket\psi\ket+)=\ket\psi\ket+\), and nothing at all
has happened. Now the flag at \(\ket-\). Here \(X\ket-=-\ket-\), so on a
basis state \(\mathrm{CNOT}\,\ket x\ket-=(-1)^x\ket x\ket-\), and by
linearity

\[\mathrm{CNOT}\,(\ket\psi\ket-)\;=\;(Z\ket\psi)\ket-,\]

where \(Z=\bigl(\begin{smallmatrix}1&0\\0&-1\end{smallmatrix}\bigr)\) is
the phase flip. The flag is untouched, and the tie has left a \(Z\) on
the rod.

\textbf{In the sketchbook.} Recall the points of Chapters 11 and 12: a
red point is \(\sqrt2\,\ket0\) and a red \(\pi\) point is
\(\sqrt2\,\ket1\); a green point is \(\ket0+\ket1=\sqrt2\,\ket+\) and a
green \(\pi\) point is \(\ket0-\ket1=\sqrt2\,\ket-\). So the flag's
\(\ket-\) is \(\tfrac1{\sqrt2}\) times a green \(\pi\) point. Feed it
into the red end of a tie. A phase-free red dot copies green points, by
the copy rule (Theorem 8.3, row 59), with a factor \(\tfrac1{\sqrt2}\).
The green \(\pi\) point goes in and comes out twice: once on the flag,
which leaves the tie as it came, and once up the tie. Up the tie it
meets the green control, and fusion (Theorem 8.1) adds its \(\pi\) to
the control's phase. A green dot with phase \(\pi\), one leg in and one
out, is \(Z\) (its matrix, §8.1):

\[\begin{tikzpicture}[sd, x=1.2em, y=1.2em,
  zsp/.style={circle,draw,line width=0.9pt,fill=green!55,minimum size=1.0em,inner sep=0pt},
  xsp/.style={circle,draw,line width=0.9pt,fill=red!75,minimum size=1.0em,inner sep=0pt},
  ph/.style={font=\scriptsize,inner sep=1pt}]
  \draw[wire] (0,1.8) -- (3.6,1.8);
  \draw[wire] (1.8,1.8) -- (1.8,0);
  \draw[wire] (0,0) -- (3.6,0);
  \node[zsp] at (1.8,1.8) {};
  \node[xsp] at (1.8,0) {};
  \node[zsp,label={[ph]below:{$\pi$}}] at (0,0) {};
\end{tikzpicture}
\;\overset{\text{copy}}{=}\;\tfrac1{\sqrt2}\;
\begin{tikzpicture}[sd, x=1.2em, y=1.2em,
  zsp/.style={circle,draw,line width=0.9pt,fill=green!55,minimum size=1.0em,inner sep=0pt},
  ph/.style={font=\scriptsize,inner sep=1pt}]
  \draw[wire] (0,1.8) -- (3.8,1.8);
  \draw[wire] (1.4,1.8) -- (1.4,0.3);
  \node[zsp] at (1.4,1.8) {};
  \node[zsp,label={[ph]left:{$\pi$}}] at (1.4,0.3) {};
  \draw[wire] (2.8,-0.3) -- (3.8,-0.3);
  \node[zsp,label={[ph]below:{$\pi$}}] at (2.8,-0.3) {};
\end{tikzpicture}
\;\overset{\text{fusion}}{=}\;\tfrac1{\sqrt2}\;
\begin{tikzpicture}[sd, x=1.2em, y=1.2em,
  zsp/.style={circle,draw,line width=0.9pt,fill=green!55,minimum size=1.0em,inner sep=0pt},
  ph/.style={font=\scriptsize,inner sep=1pt}]
  \draw[wire] (0,1.8) -- (3.2,1.8);
  \node[zsp,label={[ph]above:{$\pi$}}] at (1.6,1.8) {};
  \draw[wire] (1.8,0) -- (3.2,0);
  \node[zsp,label={[ph]below:{$\pi$}}] at (1.8,0) {};
\end{tikzpicture}\]

With a phase-free green point in place of the green \(\pi\) point, the
same two steps copy a phase-free point up the tie, fusion adds nothing,
and the control is left as a green dot with no phase and two legs: a
plain wire. Nothing comes back.

\begin{quote}
\textbf{Result 13.4 (kickback as a copy).} Exactly: (i) a green \(\pi\)
point fed into the red end of a tie is copied onto the rod as a green
\(\pi\) dot, a \(Z\), and the flag keeps its green \(\pi\) point, with
the factor \(\tfrac1{\sqrt2}\) drawn above; on the blackboard,
\(\mathrm{CNOT}\,(\ket\psi\ket-)=(Z\ket\psi)\ket-\) for every rod state
\(\ket\psi\); (ii) a phase-free green point is copied as nothing; on the
blackboard, \(\mathrm{CNOT}\,(\ket\psi\ket+)=\ket\psi\ket+\); and for
every \(f\), marked or not,
\((H^{\otimes n}\otimes I)\,U_f\,(H^{\otimes n}\otimes I)\ket{0}^{\otimes n}\ket+=\ket0^{\otimes n}\ket+\).
\end{quote}

\emph{Proof.} The pictures above, with copy and fusion; and the two
blackboard lines before them. The scalar checks: the left side is
\(\tfrac1{\sqrt2}\mathrm{CNOT}\) applied to a rod and \(\sqrt2\ket-\),
which is \(Z\otimes\ket-\); the right side is
\(\tfrac1{\sqrt2}\,Z\otimes\sqrt2\ket-\), the same. For the last claim
of (ii), \(U_f\ket x\ket+=\ket x\ket+\) for every \(x\), since
\(X\ket+=\ket+\), so \(U_f\) acts on \(\ket\psi\ket+\) as the identity
for every \(f\), and the two layers of Hadamards cancel.
\(\;\blacksquare\) See \texttt{Code/Ch13/KickedMark.py}, which tests (i)
on a random complex rod state and the kickback of row 37 on a random
non-linear \(f\).

The picture shows why kickback is a copy. The red dot copies the flag's
green point, one copy stays on the flag and one goes up the tie, and a
green point is exactly the kind of point that a green control absorbs by
fusion. The flag gives nothing away; it is copied. And the only green
points that survive the copy with a trace are those with phase \(\pi\).

\textbf{\emph{Why the difference state.}} Chapter 5 answered this on the
blackboard: \(\ket-\) is the eigenvector of \(X\) with eigenvalue
\(-1\), and the \(-1\) is what lands on the rod (row 37). The picture
gives a second answer. The flag's point must be green to be copied at
all, since a red dot copies only green points; and of the two green
points, the phase-free one fuses into the control as nothing. Only the
green \(\pi\) point leaves a mark.

\hypertarget{the-one-question-test-read-off-the-web}{%
\subsection{The one-question test, read off the
web}\label{the-one-question-test-read-off-the-web}}

Chapter 5 proved Deutsch and Jozsa's test for any width (§5.5, row 38).
Put \(n\) rods at \(\ket0\), a Hadamard on each before and after the
box, the flag at \(\ket-\), and one query. The amplitude of finding
every rod at \(0\) is

\[\frac1{2^n}\sum_{x\in\{0,1\}^n}(-1)^{f(x)},\]

which is \(\pm1\) when \(f\) is constant and \(0\) when \(f\) is
balanced (zero on exactly half the strings). We do not prove that again.
What is new is what the number is for a box of marks. It needs one sum,
which the next section needs as well. For a statement that holds exactly
when two strings are equal we use the bracket of Chapter 12:
\([\,c=0\,]\) is the number that is \(1\) if \(c\) is the all-zero
string and \(0\) otherwise.

A small case first. For \(n=2\) and \(c=10\), the four signs
\((-1)^{c\cdot x}\) at \(x=00,01,10,11\) are \(+1,+1,-1,-1\), and they
cancel. For \(c=00\) they are all \(+1\), and the sum is \(4\).

\begin{quote}
\textbf{Result 13.5 (the sign sum).} For every \(c\in\{0,1\}^n\),
\[\sum_{x\in\{0,1\}^n}(-1)^{c\cdot x}\;=\;2^n\,[\,c=0\,].\]
\end{quote}

\emph{Proof.} Since \(c\cdot x\) is a sum mod \(2\) and
\((-1)^{u\oplus v}=(-1)^u(-1)^v\), the sign is a product,
\((-1)^{c\cdot x}=\prod_i(-1)^{c_ix_i}\), with one factor for each bit
of \(x\). A sum over all strings of a product of one-bit factors is the
product of the one-bit sums:
\[\sum_{x}(-1)^{c\cdot x}=\prod_{i=1}^{n}\;\sum_{x_i\in\{0,1\}}(-1)^{c_ix_i}=\prod_{i=1}^{n}\bigl(1+(-1)^{c_i}\bigr).\]
Each factor is \(2\) if \(c_i=0\) and \(0\) if \(c_i=1\). So the product
is \(2^n\) if \(c=0\) and \(0\) otherwise. \(\;\blacksquare\) See
\texttt{Code/Ch13/KickedMark.py} and \texttt{Code/Ch13/TurnedWeb.py}.

\begin{quote}
\textbf{Result 13.6 (Deutsch and Jozsa, read off the web).} If
\(f(x)=(a\cdot x)\oplus b\), the amplitude of row 38 is
\((-1)^b\,[\,a=0\,]\). So a box of marks is constant exactly when it has
no marks, and balanced otherwise; the one-query test reads ``no ties''
or ``some tie''.
\end{quote}

\emph{Proof.}
\(2^{-n}\sum_x(-1)^{(a\cdot x)\oplus b}=(-1)^b\,2^{-n}\sum_x(-1)^{a\cdot x}=(-1)^b[\,a=0\,]\)
by Result 13.5. If \(a=0\), then \(f=b\) is constant. If \(a\neq0\) the
sum \(\sum_x(-1)^{a\cdot x}\) is \(0\), so \(f\) takes each value on
exactly half the strings. \(\;\blacksquare\) See
\texttt{Code/Ch13/KickedMark.py}.

\textbf{In the sketchbook.} The test needs no sum when the box is a web.
Take one marked rod: a red point, a yellow box, the green control, a
yellow box, then the output; and the flag as a green \(\pi\) point into
the tie's red end. Kickback as a copy (Result 13.4) puts a green \(\pi\)
on the rod between its two yellow boxes. A green dot with a yellow box
on each leg is a red dot of the same phase, by colour change (Theorem
8.2, row 58), and the red \(\pi\) fuses with the red point before it:

\[\begin{tikzpicture}[sd, x=1.2em, y=1.2em,
  zsp/.style={circle,draw,line width=0.9pt,fill=green!55,minimum size=1.0em,inner sep=0pt},
  xsp/.style={circle,draw,line width=0.9pt,fill=red!75,minimum size=1.0em,inner sep=0pt},
  hbox/.style={draw,line width=0.9pt,fill=yellow!85,minimum size=0.8em,inner sep=1pt},
  ph/.style={font=\scriptsize,inner sep=1pt}]
  \draw[wire] (0,1.8) -- (5,1.8);
  \draw[wire] (2.5,1.8) -- (2.5,0);
  \draw[wire] (0,0) -- (5,0);
  \node[xsp] at (0,1.8) {};
  \node[hbox] at (1.25,1.8) {};
  \node[zsp] at (2.5,1.8) {};
  \node[hbox] at (3.75,1.8) {};
  \node[xsp] at (2.5,0) {};
  \node[zsp,label={[ph]below:{$\pi$}}] at (0,0) {};
\end{tikzpicture}
\;\overset{\text{copy, fusion}}{=}\;\tfrac1{\sqrt2}\;
\begin{tikzpicture}[sd, x=1.2em, y=1.2em,
  zsp/.style={circle,draw,line width=0.9pt,fill=green!55,minimum size=1.0em,inner sep=0pt},
  xsp/.style={circle,draw,line width=0.9pt,fill=red!75,minimum size=1.0em,inner sep=0pt},
  hbox/.style={draw,line width=0.9pt,fill=yellow!85,minimum size=0.8em,inner sep=1pt},
  ph/.style={font=\scriptsize,inner sep=1pt}]
  \draw[wire] (0,1.8) -- (5,1.8);
  \node[xsp] at (0,1.8) {};
  \node[hbox] at (1.25,1.8) {};
  \node[zsp,label={[ph]above:{$\pi$}}] at (2.5,1.8) {};
  \node[hbox] at (3.75,1.8) {};
  \draw[wire] (2,0) -- (5,0);
  \node[zsp,label={[ph]below:{$\pi$}}] at (2,0) {};
\end{tikzpicture}
\;\overset{\text{colour, fusion}}{=}\;\tfrac1{\sqrt2}\;
\begin{tikzpicture}[sd, x=1.2em, y=1.2em,
  xsp/.style={circle,draw,line width=0.9pt,fill=red!75,minimum size=1.0em,inner sep=0pt},
  zsp/.style={circle,draw,line width=0.9pt,fill=green!55,minimum size=1.0em,inner sep=0pt},
  ph/.style={font=\scriptsize,inner sep=1pt}]
  \draw[wire] (0,1.8) -- (2,1.8);
  \node[xsp,label={[ph]above:{$\pi$}}] at (0,1.8) {};
  \draw[wire] (0,0) -- (2,0);
  \node[zsp,label={[ph]below:{$\pi$}}] at (0,0) {};
\end{tikzpicture}\]

A red \(\pi\) point is \(\sqrt2\,\ket1\): the marked rod comes back
reading \(1\). An unmarked rod has no tie, its two yellow boxes meet and
cancel (\(H^2=I\)), and it comes back as its red point, reading \(0\).
So every rod reads \(0\) exactly when the web has no ties. That is
Result 13.6, seen rod by rod. The next section does the same for the
whole web at once.

Two honest remarks bound this test. First, it reads ``no ties'' only for
a box of marks. Cleve et al.~are explicit that ``a balanced function
need not be of this form'', and a general balanced \(f\) has no web; for
it, only Chapter 5's sum of signs decides, and there the blackboard is
the only strand that works. Second, the test beats the classical worst
case, \(2^{n-1}+1\) evaluations, only because it is certain. BBBV,
describing this line of results, say at once that ``these computational
problems are in BPP relative to the same oracle'': a classical machine
that tosses coins, and may err with small probability, solves them
quickly. Such a machine can query \(k\) random strings and answer
``constant'' if all \(k\) answers agree; when \(f\) is balanced, each
answer is a fair coin and the \(k\) agree with probability \(2^{1-k}\).

\vskip 0pt plus 4\baselineskip\penalty-100\vskip 0pt plus
-4\baselineskip

\begin{quote}
\textbf{Feynman's Notebook.}\nopagebreak

\emph{From one bit to many.} Cleve and his co-authors trace the line.
David Deutsch proposed the first problem of this kind in 1985, for a
function on one bit. Deutsch and Jozsa generalised it in 1992, in
``Rapid solution of problems by quantum computation'' (\emph{Proc. R.
Soc. Lond. A} 439, 553--558), to a function on \(n\) bits promised to be
constant or balanced. Any classical algorithm, Cleve et al.~note, needs
\(2^{n-1}+1\) evaluations in the worst case to answer with certainty.
The network of this chapter, with a single \(f\)-controlled-NOT, is
theirs: they call it ``a slight improvement'' of Deutsch and Jozsa's
original algorithm, ``which involves two \(f\)-controlled-NOT
operations''. Chapter 5 told the same story for one bit. BBBV place
Deutsch and Jozsa at the start of the field's evidence that quantum
machines can be stronger: their work ``showed how to exploit some
inherently quantum mechanical features''. Then comes the honest
qualification about BPP, quoted above.
\end{quote}

\begin{center}\rule{0.5\linewidth}{0.5pt}\end{center}

\hypertarget{one-wish-asked-crosswise}{%
\section{§13.3 One Wish, Asked
Crosswise}\label{one-wish-asked-crosswise}}

The test of §13.2 asked only whether the web has ties. It did so by
bringing a mark back onto every marked rod, and then it asked only
whether any rod was marked. Why not read every rod?

Try it on \(a=101\), \(n=3\), \(b=0\). Put each rod at \(\ket+=H\ket0\)
and the flag at \(\ket-\). By kickback, the query multiplies each string
\(\ket x\) by \((-1)^{f(x)}=(-1)^{x_1\oplus x_3}=(-1)^{x_1}(-1)^{x_3}\),
and a sign that is a product of one-rod signs gives a state that is a
product of one-rod states:

\[U_f\;\ket+\ket+\ket+\ket-\;=\;\frac{\ket0-\ket1}{\sqrt2}\;\frac{\ket0+\ket1}{\sqrt2}\;\frac{\ket0-\ket1}{\sqrt2}\;\ket-\;=\;\ket-\ket+\ket-\ket- .\]

Now a Hadamard on each rod turns \(\ket-\) into \(\ket1\) and \(\ket+\)
into \(\ket0\): the rods read \(101\). One query has returned the whole
secret, where Result 13.3 says that plain queries need three. The rest
of this section says why it always works, and what the picture sees.

\hypertarget{the-turned-tie}{%
\subsection{The turned tie}\label{the-turned-tie}}

First a small surprise about a single CNOT. Put a Hadamard on both of
its wires, before and after. Feed in \(\ket{01}\): the Hadamards make
\(\ket+\ket-\); the CNOT, by kickback, makes \(\ket-\ket-\); the
Hadamards make \(\ket{11}\). The bottom bit was \(1\), and the
\emph{top} bit has been flipped. Feed in \(\ket{10}\): \(\ket-\ket+\),
then nothing (since \(X\ket+=\ket+\)), then \(\ket{10}\). The top bit
was \(1\), and nothing happened to the bottom. The control has moved to
the bottom wire.

\begin{quote}
\textbf{Result 13.7 (the turned web).} Exactly,
\[(H\otimes H)\,\mathrm{CNOT}_{1\to2}\,(H\otimes H)\;=\;\mathrm{CNOT}_{2\to1}.\]
Hence, for \(f(x)=(a\cdot x)\oplus b\),
\[H^{\otimes(n+1)}\,U_f\,H^{\otimes(n+1)}\;=\;Z_{\text{flag}}^{\,b}\prod_{i:\,a_i=1}\mathrm{CNOT}_{\text{flag}\to i}:\]
the flag controls, and every marked rod is a target. In pictures: yellow
boxes on every outer leg of the web turn every dot's colour, and the
ties now run \emph{from} the flag \emph{to} the marked rods, with no
scalar.
\end{quote}

\emph{Proof.} Write
\(\mathrm{CNOT}_{1\to2}=\ket0\bra0\otimes I+\ket1\bra1\otimes X\).
Conjugating by \(H\otimes H\) gives
\(\ket+\bra+\otimes I+\ket-\bra-\otimes Z\), since \(H\ket0=\ket+\),
\(H\ket1=\ket-\) and \(HXH=Z\). With \(\ket\pm\bra\pm=\tfrac12(I\pm X)\)
this is
\[\begin{aligned}\tfrac12(I+X)\otimes I+\tfrac12(I-X)\otimes Z\;&=\;I\otimes\tfrac12(I+Z)+X\otimes\tfrac12(I-Z)\\&=\;I\otimes\ket0\bra0+X\otimes\ket1\bra1,\end{aligned}\]
which is \(\mathrm{CNOT}_{2\to1}\). For the web, \(H^{\otimes(n+1)}\) is
its own inverse, so conjugating the product of Result 13.2 conjugates
each factor. Each \(\mathrm{CNOT}_{i\to\text{flag}}\) becomes
\(\mathrm{CNOT}_{\text{flag}\to i}\) (the Hadamards on the other wires
cancel in pairs), and \(X_{\text{flag}}\) becomes
\(HXH=Z_{\text{flag}}\). \(\;\blacksquare\) See
\texttt{Code/Ch13/TurnedWeb.py}.

\textbf{In the sketchbook.} Take one tie with a yellow box on each of
its four outer legs. The green control has boxes on two of its three
legs, and the red target likewise; the tie itself is bare. Put a
cancelling pair of yellow boxes on the tie, which changes nothing since
\(H^2=I\) (Chapter 3; row 56). Now each dot wears a yellow box on every
leg it holds, so colour change (Theorem 8.2) turns the green dot red and
the red dot green, and the boxes come off. The pair on the tie is split,
one box absorbed at each end:

\[\begin{tikzpicture}[sd, x=1.2em, y=1.2em,
  zsp/.style={circle,draw,line width=0.9pt,fill=green!55,minimum size=1.0em,inner sep=0pt},
  xsp/.style={circle,draw,line width=0.9pt,fill=red!75,minimum size=1.0em,inner sep=0pt},
  hbox/.style={draw,line width=0.9pt,fill=yellow!85,minimum size=0.8em,inner sep=1pt}]
  \draw[wire] (0,3) -- (4,3);
  \draw[wire] (0,0) -- (4,0);
  \draw[wire] (2,3) -- (2,0);
  \node[hbox] at (0.8,3) {}; \node[hbox] at (3.2,3) {};
  \node[hbox] at (0.8,0) {}; \node[hbox] at (3.2,0) {};
  \node[zsp] at (2,3) {}; \node[xsp] at (2,0) {};
\end{tikzpicture}
\;\overset{H^2=I}{=}\;
\begin{tikzpicture}[sd, x=1.2em, y=1.2em,
  zsp/.style={circle,draw,line width=0.9pt,fill=green!55,minimum size=1.0em,inner sep=0pt},
  xsp/.style={circle,draw,line width=0.9pt,fill=red!75,minimum size=1.0em,inner sep=0pt},
  hbox/.style={draw,line width=0.9pt,fill=yellow!85,minimum size=0.8em,inner sep=1pt}]
  \draw[wire] (0,3) -- (4,3);
  \draw[wire] (0,0) -- (4,0);
  \draw[wire] (2,3) -- (2,0);
  \node[hbox] at (0.8,3) {}; \node[hbox] at (3.2,3) {};
  \node[hbox] at (0.8,0) {}; \node[hbox] at (3.2,0) {};
  \node[hbox] at (2,2.05) {}; \node[hbox] at (2,0.95) {};
  \node[zsp] at (2,3) {}; \node[xsp] at (2,0) {};
\end{tikzpicture}
\;\overset{\text{colour}}{=}\;
\begin{tikzpicture}[sd, x=1.2em, y=1.2em,
  zsp/.style={circle,draw,line width=0.9pt,fill=green!55,minimum size=1.0em,inner sep=0pt},
  xsp/.style={circle,draw,line width=0.9pt,fill=red!75,minimum size=1.0em,inner sep=0pt}]
  \draw[wire] (0,3) -- (3,3);
  \draw[wire] (0,0) -- (3,0);
  \draw[wire] (1.5,3) -- (1.5,0);
  \node[xsp] at (1.5,3) {}; \node[zsp] at (1.5,0) {};
\end{tikzpicture}\]

A red dot on top tied to a green dot below is a CNOT with its control on
the \emph{bottom} wire (row 64, read upside down), and both bare pairs
carry the same \(\tfrac1{\sqrt2}\), so the equation holds with no
scalar. For a whole web, the same two moves run on every tie at once:
one cancelling pair on each tie, one on each stretch of the flag between
two of its red dots, and then colour change on every dot. That is Result
13.7 in pictures, \textbf{the turned web}.

Now the headline. Its statement is the same network as Chapter 5's test,
with nothing changed but the question we ask of the output.

\begin{quote}
\textbf{Result 13.8 (Bernstein--Vazirani, in Cleve et al.'s one-query
form).} Let \(a\in\{0,1\}^n\), \(b\in\{0,1\}\) and
\(f(x)=(a\cdot x)\oplus b\), and let
\(\ket-=\tfrac1{\sqrt2}(\ket0-\ket1)\). Then, for every \(n\), \(a\) and
\(b\), exactly,
\[(H^{\otimes n}\otimes I)\,U_f\,(H^{\otimes n}\otimes I)\;\ket{0}^{\otimes n}\ket-\;=\;(-1)^b\,\ket{a}\ket- .\]
One query, and reading the rods returns \(a\) with certainty.
\end{quote}

The attribution needs care, because two different things are credited
here. The \emph{problem} is Ethan Bernstein and Umesh Vazirani's, from
1993: given a box that computes \((a\cdot x)\oplus b\), find \(a\). They
also gave the first quantum algorithm for it, which, as Cleve et
al.~record, ``employs two \(f\)-controlled-NOT operations instead of
one''. The one-query network stated above is Cleve et al.'s
presentation, the same network as their one-query form of Deutsch and
Jozsa's test; their Figure 4 is drawn once for both. Cleve et al.~also
point out why the two problems share a network: a function
\((a\cdot x)\oplus b\) ``is constant if \(a=00\cdots0\) and balanced
otherwise'', so every Bernstein--Vazirani box satisfies Deutsch and
Jozsa's promise, though not every balanced function is of this form.

\emph{Proof idea.} On the blackboard, kickback writes the sign
\((-1)^{a\cdot x}\) on every string, the sign factors into one sign per
rod, and the second layer of Hadamards reads each rod's sign as a bit.
In the sketchbook, the Hadamards turn the web round (Result 13.7), so
the ties run from the flag to the marked rods, and the flag's \(1\) is
copied up every tie. Both proofs are given in full in §13.6. See
\texttt{Code/Ch13/TurnedWeb.py}, which checks the equation, sign
included, for every \(a\) and both \(b\) at \(n=4\), and for random
secrets at \(n=6\).

\textbf{\emph{Why this definition.}} Why must the flag be crosswise too?
Because with the flag left at \(\ket0\) the query entangles the rods
with it: the state after the query is
\(2^{-n/2}\sum_x\ket x\ket{(a\cdot x)\oplus b}\). Take \(b=0\), split
the state by the flag's value, and apply the second layer to the rods;
write \(z\) for an \(n\)-bit string read from the rods. Using Result
13.5 on
\(\sum_{x:\,a\cdot x=0}(-1)^{x\cdot z}=\tfrac12\sum_x\bigl(1+(-1)^{a\cdot x}\bigr)(-1)^{x\cdot z}\),
one finds, for \(a\neq0\),
\[\tfrac12\bigl(\ket{0^n}+\ket a\bigr)\ket0+\tfrac12\bigl(\ket{0^n}-\ket a\bigr)\ket1 .\]
The rods read \(0^n\) or \(a\), each with probability \(\tfrac12\), and
the flag is a fair coin independent of them. With the flag at \(\ket1\),
or with \(b=1\), the two branches change places. So a query with the
flag upright gambles the one question. Only the difference state keeps
the flag out of the answer.

In the dialogue, just after the Tortoise's tubes have turned every ring
on her cloth web, Achilles looks up at \emph{Belvedere}: a pillared
lookout whose upper storey is set crosswise on the lower; a ladder
stands inside it and leans against the outside.

\vskip 0pt plus 4\baselineskip\penalty-100\vskip 0pt plus
-4\baselineskip

\begin{quote}
\textbf{Dirac's Blackboard.}\nopagebreak

\emph{Why the second layer returns \(a\).} After the query the rods hold
\(2^{-n/2}(-1)^b\sum_x(-1)^{a\cdot x}\ket x\). The second layer sends
\(\ket x\) to \(2^{-n/2}\sum_z(-1)^{x\cdot z}\ket z\), so the amplitude
of \(\ket z\) is
\[(-1)^b\,2^{-n}\sum_x(-1)^{x\cdot(a\oplus z)}=(-1)^b\prod_{i=1}^{n}\tfrac12\bigl(1+(-1)^{a_i\oplus z_i}\bigr),\]
the sum factored bit by bit as in Result 13.5. Each factor is \(1\) when
\(z_i=a_i\) and \(0\) when \(z_i\neq a_i\). Take \(n=2\), \(a=10\). For
\(z=10\) both factors are \(1\). For \(z=00\) the first factor is
\(\tfrac12(1-1)=0\); for \(z=11\) the second is; for \(z=01\) both are.
The second layer never compares whole strings. It asks each rod
separately whether its sign was \(+\) or \(-\), and \(2^n\) roads
collapse to \(n\) one-rod questions.
\end{quote}

\vskip 0pt plus 4\baselineskip\penalty-100\vskip 0pt plus
-4\baselineskip

\begin{quote}
\textbf{Coecke's Sketchbook.}\nopagebreak

\emph{The Hadamards turn the web.} Colour change is usually applied to
one dot. Here it is applied to a whole web at once, and what makes that
possible is a single bookkeeping move: a cancelling pair of yellow boxes
on every tie and on every stretch of the flag between its red dots.
After that, every dot is dressed in yellow on every leg, every dot
changes colour together, and every tie reverses its direction. The
picture never forms the sum over \(2^n\) strings that the blackboard
needs. It works with the ties, and there are only \(w\) of them. This is
not a trick that works on any oracle. It works because this oracle
\emph{is} a web: the box is a handful of CNOTs, and the secret is which
rods they touch. A general \(f\) is a tall box with nothing inside to
turn, and for it the blackboard's sum is the only proof we have.
\end{quote}

\vskip 0pt plus 4\baselineskip\penalty-100\vskip 0pt plus
-4\baselineskip

\begin{quote}
\textbf{Feynman's Notebook.}\nopagebreak

\emph{An interferometer with many arms.} Ethan Bernstein and Umesh
Vazirani posed the problem of this section in ``Quantum complexity
theory'', in the proceedings of the 25th ACM Symposium on Theory of
Computing, 1993, pages 11--20; a journal version appeared in the
\emph{SIAM Journal on Computing} in 1997. Cleve et al.~describe their
network as an experiment. In a Mach--Zehnder interferometer a photon
meets a half-silvered mirror, travels two paths through phase shifters,
and meets a second mirror before the detectors. Cleve et al.~identify
each half-silvered mirror with a Hadamard, and remark that the second
mirror ``effectively erases all information about the path taken'',
which is what lets the paths interfere. Their summary of the known
algorithms is a three-part pattern: ``a Fourier transform, followed by
an \(f\)-controlled-\(U\), followed by another Fourier transform'', and
they insist that the last transform ``is essential''. This chapter's
Hadamard layers are the simplest Fourier transform, and the query sits
between them.
\end{quote}

\begin{center}\rule{0.5\linewidth}{0.5pt}\end{center}

\hypertarget{the-pairs-box}{%
\section{§13.4 The Pairs Box}\label{the-pairs-box}}

The Bernstein--Vazirani box has its whole secret in its wiring, and a
crosswise query reads the wiring whole. Here is a box with a different
kind of promise, whose secret no single query has been shown to return.

\textbf{On the blackboard.} A \textbf{pairs box} computes a function
\(f:\{0,1\}^n\to\{0,1\}^m\) with \(m\) output bits, with this promise:
there is a secret string \(s\neq0\) such that
\[f(x)=f(x')\quad\text{exactly when}\quad x'\in\{x,\;x\oplus s\}.\] So
the inputs come in pairs \(\{x,x\oplus s\}\), each pair shares an
answer, and different pairs have different answers. The task is to find
\(s\). This is Simon's problem. The oracle is
\(U_f\ket x\ket y=\ket x\ket{y\oplus f(x)}\) as before, where now \(y\)
is an \(m\)-bit string on \(m\) \textbf{answer wires} that take the
flag's place, and \(\oplus\) acts bit by bit.

Here is the smallest pairs box we will draw. Take \(n=3\), \(m=2\), and
let \(f(x)=Mx\) for the matrix
\[M=\begin{pmatrix}1&1&0\\0&0&1\end{pmatrix},\qquad f(x_1x_2x_3)=(x_1\oplus x_2,\;x_3),\]
with the arithmetic mod \(2\). Then \(f(000)=f(110)=00\),
\(f(001)=f(111)=01\), \(f(010)=f(100)=10\) and \(f(011)=f(101)=11\). The
pairs differ by \(s=110\). This box happens to be linear, so by Result
13.2 it is a web: rods \(1\) and \(2\) are tied to the first answer wire
and rod \(3\) to the second. Linearity is not part of Simon's promise,
though at this size it comes free: with three rods and two answer wires
every box with the promise is a web, perhaps with some answer bits
flipped. With more answer wires there are boxes with the promise that
are not webs. What follows about one round holds for every \(f\) with
the promise, web or not.

One \textbf{round} of Simon's method is the network we have used all
chapter, with the answer wires starting at \(\ket0\): rods at \(\ket0\),
a Hadamard on each rod, one query, a Hadamard on each rod, and then the
rods are read. The answer wires are never looked at: they are
\textbf{discarded}, in the sense of Chapter 9. The string read from the
rods we call the \textbf{clue}, \(z\).

\begin{quote}
\textbf{Result 13.9 (one round of Simon's problem).} For every \(f\)
with the promise, one round leaves the rods in the state
\[\rho\;=\;2^{-n}\bigl(I+Z^{s}\bigr)\;=\;2^{1-n}P_s,\] where
\(Z^{s}=Z^{s_1}\otimes\cdots\otimes Z^{s_n}\) applies \(Z\) to the rods
where \(s\) has a \(1\), and \(P_s=\tfrac12(I+Z^{s})\). The clue \(z\)
appears with probability \(2^{1-n}\,[\,z\cdot s=0\,]\): it is equally
likely to be any of the \(2^{n-1}\) strings with \(z\cdot s=0\), and
never anything else.
\end{quote}

\emph{Proof.} After the first layer and the query the state is
\(\ket\Psi=2^{-n/2}\sum_x\ket x\ket{f(x)}\). The second layer acts only
on the rods, and discarding acts only on the answer wires, so we may
discard first. Discarding is the partial trace (Result 9.2):
\[\operatorname{tr}_{\text{ans}}\ket\Psi\bra\Psi=2^{-n}\sum_{x,x'}\braket{f(x')|f(x)}\,\ket x\bra{x'}=2^{-n}\sum_{x,x'}[\,f(x)=f(x')\,]\,\ket x\bra{x'} .\]
By the promise, \(f(x)=f(x')\) exactly when \(x'=x\) or
\(x'=x\oplus s\), so this is
\(2^{-n}\sum_x\bigl(\ket x\bra x+\ket x\bra{x\oplus s}\bigr)=2^{-n}\bigl(I+X^{s}\bigr)\),
where \(X^{s}\) flips the rods where \(s\) has a \(1\) (it sends
\(\ket{x\oplus s}\) to \(\ket x\)). The second layer gives
\(\rho=2^{-n}\bigl(I+H^{\otimes n}X^{s}H^{\otimes n}\bigr)=2^{-n}\bigl(I+Z^{s}\bigr)\),
since \(HXH=Z\) and \(HIH=I\) on each rod. Finally
\(Z^{s}\ket z=(-1)^{s\cdot z}\ket z\), so the diagonal entry at \(z\) is
\(2^{-n}\bigl(1+(-1)^{s\cdot z}\bigr)=2^{1-n}[\,z\cdot s=0\,]\).
\(\;\blacksquare\) See \texttt{Code/Ch13/PairsBox.py}, which checks the
pairs box and a six-rod instance with \(s=100010\).

For the pairs box, the clue is one of \(000\), \(001\), \(110\),
\(111\), each with probability \(\tfrac14\). One clue does not name
\(s\). Each clue is one parity check that the secret passes,
\(z\cdot s=0\): a line through the secret, not the point. How many are
needed? Suppose two rounds give the clues \(001\) and \(110\). The first
says \(s_3=0\); the second says \(s_1\oplus s_2=0\). Since \(s\neq0\),
the only candidate left is \(s=110\).

\begin{quote}
\textbf{Result 13.10 (clues fix the secret).} If \(z_1,\dots,z_{n-1}\)
are clues that are linearly independent mod \(2\) (no nonempty set of
them adds up to \(0\)), then \(s\) is the only string \(v\neq0\) with
\(v\cdot z_j=0\) for all \(j\). On average, at most \(2(n-1)\) rounds
are needed to collect \(n-1\) independent clues.
\end{quote}

\emph{Proof.} Every clue satisfies \(z_j\cdot s=0\), so \(s\) is a
solution. The map \(v\mapsto(v\cdot z_1,\dots,v\cdot z_{n-1})\) has
\(n-1\) independent rows, so its image has \(2^{n-1}\) elements, and
each image value is taken by \(2^n/2^{n-1}=2\) strings \(v\). In
particular exactly two strings solve all the equations, \(0\) and \(s\).
For the rounds: suppose \(k<n-1\) independent clues are in hand. The
strings they add up to, \(2^k\) of them, all satisfy \(z\cdot s=0\), and
a new clue is equally likely to be any of the \(2^{n-1}\) such strings,
so it falls among the old ones with probability
\(2^k/2^{n-1}\le\tfrac12\). A new independent clue therefore takes at
most \(2\) rounds on average, and \(n-1\) of them at most \(2(n-1)\).
\(\;\blacksquare\) See \texttt{Code/Ch13/PairsBox.py}, which checks
every independent pair of clues for the pairs box.

The all-zero clue always falls among the old ones: it tells nothing.

\textbf{In the sketchbook.} Draw the round doubled, as Chapter 9 taught:
each wire as two lines, the wire and its mirror copy, with the mirror
copy of each gate on the mirror lines (row 66). The mirror conjugates,
but every gate and point here is real, so the mirror copy of the circuit
is the same circuit, drawn below it with its rows in reverse order. To
discard an answer wire, join it to its own mirror copy by a cap (row
67); as always the cap is unnormalized, cap \(=\bra{00}+\bra{11}\). The
answer wires start at red points, \(\sqrt2\ket0\), like the rods. Rods
\(1\), \(2\), \(3\) are at the top and their mirrors \(\bar1\),
\(\bar2\), \(\bar3\) at the bottom; the four answer lines in the middle
are joined in pairs by the two nested caps:

\begin{center}
\begin{tikzpicture}[sd, x=1.2em, y=1.3em,
  zsp/.style={circle,draw,line width=0.9pt,fill=green!55,minimum size=0.85em,inner sep=0pt},
  xsp/.style={circle,draw,line width=0.9pt,fill=red!75,minimum size=0.85em,inner sep=0pt},
  hbox/.style={draw,line width=0.9pt,fill=yellow!85,minimum size=0.7em,inner sep=1pt}]
  \draw[wire] (3,9) -- (3,6);
  \draw[wire] (4.2,8) -- (4.2,6);
  \draw[wire] (5.4,7) -- (5.4,5);
  \draw[wire] (3,0) -- (3,3);
  \draw[wire] (4.2,1) -- (4.2,3);
  \draw[wire] (5.4,2) -- (5.4,4);
  \foreach \y in {9,8,7,2,1,0} {
    \draw[wire] (0,\y) -- (8.4,\y);
    \node[hbox] at (1.5,\y) {};
    \node[hbox] at (6.9,\y) {};
    \node[xsp] at (0,\y) {};
  }
  \foreach \y in {6,5,4,3} { \draw[wire] (0,\y) -- (6.6,\y); \node[xsp] at (0,\y) {}; }
  \sdcap{6.6}{5}{4}
  \sdcap{6.6}{6}{3}
  \node[zsp] at (3,9) {}; \node[zsp] at (4.2,8) {}; \node[zsp] at (5.4,7) {};
  \node[xsp] at (3,6) {}; \node[xsp] at (4.2,6) {}; \node[xsp] at (5.4,5) {};
  \node[zsp] at (3,0) {}; \node[zsp] at (4.2,1) {}; \node[zsp] at (5.4,2) {};
  \node[xsp] at (3,3) {}; \node[xsp] at (4.2,3) {}; \node[xsp] at (5.4,4) {};
  \node[sdlabel,anchor=west] at (8.5,9) {$1$};
  \node[sdlabel,anchor=west] at (8.5,8) {$2$};
  \node[sdlabel,anchor=west] at (8.5,7) {$3$};
  \node[sdlabel,anchor=west] at (8.5,2) {$\bar3$};
  \node[sdlabel,anchor=west] at (8.5,1) {$\bar2$};
  \node[sdlabel,anchor=west] at (8.5,0) {$\bar1$};
\end{tikzpicture}
\end{center}

Read as a doubled state on the six open legs, this picture is \(4\rho\):
the ten points carry \(2^{5}\) and the six bare ties \(2^{-3}\). Four
rewrites finish it, and none of them costs a scalar.

\begin{enumerate}
\def\labelenumi{\arabic{enumi}.}
\tightlist
\item
  \textbf{Colour change} (Theorem 8.2). On each rod, a red point
  followed by a yellow box is a green point:
  \(H\,\sqrt2\ket0=\ket0+\ket1\).
\item
  \textbf{Fusion} (Theorem 8.1). That green point fuses into the rod's
  green control. What remains is a phase-free green dot with two legs,
  the tie and the rest of the rod, which is a plain wire (§8.1). Each
  tie now runs straight through to its rod's last yellow box.
\item
  \textbf{Fusion}, across the caps. On each answer line the red point
  fuses into the red targets. A cap is a plain wire, so the red dots of
  an answer wire and of its mirror copy fuse through it into one red
  dot. The first answer wire and its mirror give one red dot tied to
  rods \(1\), \(2\), \(\bar1\), \(\bar2\); the second gives one tied to
  rods \(3\) and \(\bar3\).
\item
  \textbf{Colour change.} Each of those red dots has a yellow box on
  every leg, the rods' last boxes, so each turns green, and the boxes
  come off.
\end{enumerate}

What is left is a green dot with legs on \(1,2,\bar1,\bar2\), and a
green dot with legs on \(3,\bar3\), which is a cup:

\[4\rho\;=\;
\begin{tikzpicture}[sd, x=1.2em, y=1.3em,
  zsp/.style={circle,draw,line width=0.9pt,fill=green!55,minimum size=0.95em,inner sep=0pt}]
  \draw[wire] (0,2.5) -- (1.6,5) -- (3.4,5);
  \draw[wire] (0,2.5) -- (1.6,4) -- (3.4,4);
  \draw[wire] (0,2.5) -- (1.6,1) -- (3.4,1);
  \draw[wire] (0,2.5) -- (1.6,0) -- (3.4,0);
  \node[zsp] at (0,2.5) {};
  \draw[wire] (2.4,3) -- (3.4,3);
  \draw[wire] (2.4,2) -- (3.4,2);
  \sdcup{2.4}{3}{2}
  \node[sdlabel,anchor=west] at (3.5,5) {$1$};
  \node[sdlabel,anchor=west] at (3.5,4) {$2$};
  \node[sdlabel,anchor=west] at (3.5,3) {$3$};
  \node[sdlabel,anchor=west] at (3.5,2) {$\bar3$};
  \node[sdlabel,anchor=west] at (3.5,1) {$\bar2$};
  \node[sdlabel,anchor=west] at (3.5,0) {$\bar1$};
\end{tikzpicture}
\qquad\text{so}\qquad
\rho\;=\;\tfrac14\;
\begin{tikzpicture}[sd, x=1.2em, y=1.3em,
  zsp/.style={circle,draw,line width=0.9pt,fill=green!55,minimum size=0.95em,inner sep=0pt},
  xsp/.style={circle,draw,line width=0.9pt,fill=red!75,minimum size=0.95em,inner sep=0pt}]
  \draw[wire] (0,2.4) -- (4,2.4);
  \draw[wire] (0,1.2) -- (4,1.2);
  \draw[wire] (0,0) -- (4,0);
  \draw[wire] (1.2,2.4) -- (2,-1.2);
  \draw[wire] (2.8,1.2) -- (2,-1.2);
  \node[zsp] at (1.2,2.4) {};
  \node[zsp] at (2.8,1.2) {};
  \node[xsp] at (2,-1.2) {};
  \node[sdlabel,anchor=west] at (4.1,2.4) {$1$};
  \node[sdlabel,anchor=west] at (4.1,1.2) {$2$};
  \node[sdlabel,anchor=west] at (4.1,0) {$3$};
\end{tikzpicture}\]

Every step was exact. The last equality reads the doubled picture as a
matrix, in Chapter 9's way: a doubled state on the legs \(q\) and
\(\bar q\) lists the entries \(\rho_{z,z'}\) of \(\rho\), with \(z\) on
the wires and \(z'\) on the mirrors. The cup on \(3,\bar3\) is the
identity on rod \(3\) (row 69, without its \(\tfrac12\)). The green dot
on \(1,2,\bar1,\bar2\) is \(\ket{0000}+\ket{1111}\), whose only entries
are at \((00,00)\) and \((11,11)\): the matrix
\(\ket{00}\bra{00}+\ket{11}\bra{11}\) on rods \(1\) and \(2\). Drawn as
a process with legs in and out, that matrix is a green dot on rod \(1\)
and a green dot on rod \(2\), tied to one phase-free red dot. The red
dot has two legs, so it is a plain wire (§8.1), and the two green dots
fuse through it into one; the drawing keeps the red dot because of what
it does in general.

Call that drawing the \textbf{parity gadget} of \(s\): a phase-free
green dot on each rod where \(s\) has a \(1\), each tied to one
phase-free red dot with \(|s|\) legs, where \(|s|\) is the number of
\(1\)s in \(s\), and the other rods bare. The red dot reads a parity, as
it did in Chapter 8's Octet and on Chapter 12's records (row 96).

\begin{quote}
\textbf{Result 13.11 (the parity gadget).} For every \(s\neq0\),
exactly, \[P_s\;=\;2^{|s|/2-1}\times\text{(the parity gadget of }s).\]
So one round leaves \(\rho=2^{1-n}P_s=2^{|s|/2-n}\times\) gadget. For
the pairs box, \(|s|=2\) and \(n=3\): \(P_s\) is the gadget itself and
\(\rho=\tfrac14\times\) gadget, as drawn.
\end{quote}

\emph{Proof.} A phase-free green dot with legs in, out and tie sends
\(\ket z\) to \(\ket z\ket z\) (Chapter 11's copier, row 81), so the
green dots put a copy of each \(z_i\) with \(s_i=1\) on its tie. A
phase-free red dot with \(k\) legs, all of them ties coming in, is the
effect \(\bra+^{\otimes k}+\bra-^{\otimes k}\) (row 54). On a string
\(t\) of \(k\) bits it gives
\(2^{-k/2}\bigl(1+(-1)^{t_1\oplus\cdots\oplus t_k}\bigr)\), which is
\(2^{1-k/2}\) if the parity of \(t\) is even and \(0\) if it is odd.
With \(k=|s|\) the ties carry the bits of \(z\) at the places of \(s\),
whose parity is \(s\cdot z\). So the gadget sends \(\ket z\) to
\(2^{1-|s|/2}[\,s\cdot z=0\,]\ket z=2^{1-|s|/2}P_s\ket z\).
\(\;\blacksquare\) See \texttt{Code/Ch13/PairsBox.py}, which also checks
a secret of weight three, where the factor is not \(1\).

The scalar matters. For \(|s|\neq2\) the gadget is not \(P_s\), and
\(\rho\) is not \(2^{1-n}\) times the gadget.

\textbf{\emph{Why this definition.}} Why discard the answer wires,
rather than read them? Reading them would tell us \(f\) of some random
string, which says nothing about \(s\) that one plain query would not.
And nothing we do to the answer wires after the query can change the
statistics of the rods: a process on one part followed by discarding it
is the same as discarding it, for every trace-preserving process
(causality, Result 9.4). So the discard is the honest summary of what
the rods alone can say. In the picture it is also what makes the round
simple: the caps are what let the answer wires' red dots fuse with their
mirrors.

How much better is this than asking plainly? Bennett, Bernstein,
Brassard and Vazirani report the answer in their introduction, and we
state it without proof.

\begin{quote}
\textbf{Result 13.12 (Simon; stated, not proved).} Simon proved ``the
existence of an oracle relative to which BQP cannot even be simulated by
probabilistic machines allowed to run for \(2^{n/2}\) steps'' (BBBV §1).
Here BQP is the class of problems a quantum machine solves in polynomial
time with error probability at most \(\tfrac13\), and a probabilistic
machine is a classical one that may toss coins.
\end{quote}

BBBV add that Simon's paper ``also introduced an important new technique
which was one of the ingredients'' in Shor's factoring algorithm. That
is where Chapter 14 begins. The separation is a statement about all
classical machines, and neither the pictures nor the blackboard of this
chapter prove it. What we have proved is what one round does: it buys
one clue, uniformly chosen among the parity checks the secret passes.

\begin{center}\rule{0.5\linewidth}{0.5pt}\end{center}

\hypertarget{the-rosetta-table-11}{%
\section{§13.5 The Rosetta Table}\label{the-rosetta-table-11}}

The rows this chapter adds to the running dictionary, numbered on from
Chapter 12's ninety-six. There is no new picture element in them. Every
row is a new use of something Chapters 5, 8 and 9 already drew: the tall
box (97), the green--red CNOT (98, 99), the copy rule (100, 101), and
the doubled picture with its caps (102).

\begingroup\renewcommand{\arraystretch}{1.5}

\begin{longtable}[]{@{}
  >{\raggedright\arraybackslash}p{(\columnwidth - 4\tabcolsep) * \real{0.1250}}
  >{\raggedright\arraybackslash}p{(\columnwidth - 4\tabcolsep) * \real{0.4375}}
  >{\raggedright\arraybackslash}p{(\columnwidth - 4\tabcolsep) * \real{0.4375}}@{}}
\toprule\noalign{}
\begin{minipage}[b]{\linewidth}\raggedright
\#
\end{minipage} & \begin{minipage}[b]{\linewidth}\raggedright
Algebra (Dirac)
\end{minipage} & \begin{minipage}[b]{\linewidth}\raggedright
Picture (Coecke)
\end{minipage} \\
\midrule\noalign{}
\endhead
\bottomrule\noalign{}
\endlastfoot
97 & a \textbf{query} is one use of
\(U_f\ket x\ket y=\ket x\ket{y\oplus f(x)}\); \(U_f\) is unitary and
\(U_f^2=I\); a superposed query is one query & Chapter 5's tall box
\(U_f\); two in a row are bare wires \\
98 & linear \(f=(a\cdot x)\oplus b\):
\(U_f=X^{b}_{\text{flag}}\prod_{a_i=1}\mathrm{CNOT}_{i\to\text{flag}}\)
& \textbf{the oracle is a web of ties}: a green dot on each marked rod
tied to its own red dot on the flag; bare, each tie is
\(\mathrm{CNOT}/\sqrt2\) (row 64); asked twice, Hopf leaves bare
wires \\
99 &
\((H\otimes H)\,\mathrm{CNOT}_{1\to2}\,(H\otimes H)=\mathrm{CNOT}_{2\to1}\);
\(H^{\otimes(n+1)}U_fH^{\otimes(n+1)}=\) fan-out from the flag onto the
marked rods (\(b=0\)) & \textbf{the turned web} (drawn for \(b=0\); if
\(b=1\) the flag's red \(\pi\) comes through the turn as a green \(\pi\)
and only signs the result): yellow boxes on every leg turn every dot's
colour; the ties now run from the flag to the rods (exact, no scalar) \\
100 & kickback, \(\mathrm{CNOT}\,(\ket\psi\ket-)=(Z\ket\psi)\ket-\) (row
37 in spiders) & \textbf{kickback as a copy}: a green \(\pi\) point into
a red target is copied onto the control's wire as a green \(\pi\); a
phase-free green point copies as nothing \\
101 & \textbf{Bernstein--Vazirani}:
\((H^{\otimes n}\otimes I)\,U_f\,(H^{\otimes n}\otimes I)\ket0^{\otimes n}\ket-=(-1)^b\ket a\ket-\)
& \textbf{the secret is read off the wiring}: after the turn, the flag's
\(1\) is copied up every tie, and exactly the marked rods end at \(1\)
(the master, \emph{The Turned Web}) \\
102 & one Simon round, answer wires discarded:
\(\rho=2^{-n}(I+Z^{s})=2^{1-n}P_s\); clues \(z\) with \(z\cdot s=0\),
each equally likely; \(P_s=2^{\lvert s\rvert/2-1}\times\) gadget & the
discarded answer wires are caps (row 67); the round is
\(2^{\lvert s\rvert/2-n}\times\) \textbf{the parity gadget} (green dots
on the \(s\)-rods, one red dot with \(\lvert s\rvert\) legs); for the
pairs box \(\rho=\tfrac14\times\) gadget: a line through the secret \\
\end{longtable}

\endgroup

Every row is exact, scalar included, as in Chapters 9 to 12. Nothing in
these rows is a statement about how many queries a task needs: rows 97
to 102 are equalities of processes, and the bounds of this chapter are
proved by counting (Result 13.3) or cited (Results 13.12 and 13.14).

\begin{center}\rule{0.5\linewidth}{0.5pt}\end{center}

\hypertarget{the-theorem-twice-12}{%
\section{§13.6 The Theorem, Twice}\label{the-theorem-twice-12}}

\begin{quote}
\textbf{Result 13.8 (Bernstein--Vazirani, in Cleve et al.'s one-query
form), restated.} For \(f(x)=(a\cdot x)\oplus b\) and
\(\ket-=\tfrac1{\sqrt2}(\ket0-\ket1)\), for every \(n\), \(a\) and
\(b\), exactly,
\((H^{\otimes n}\otimes I)\,U_f\,(H^{\otimes n}\otimes I)\ket0^{\otimes n}\ket-=(-1)^b\ket a\ket-\).
\end{quote}

\hypertarget{diracs-blackboard-12}{%
\subsection{Dirac's Blackboard}\label{diracs-blackboard-12}}

These are Cleve et al.'s lines, their Eqs. (3.4), (3.5), (3.8) and
(3.9), with the normalization they omit written back in.

\textbf{Step 1: the first layer.} One Hadamard on each rod sends
\(\ket0\) to \(\tfrac1{\sqrt2}(\ket0+\ket1)\), and the product over
\(n\) rods is the equal sum of all strings:
\[(H^{\otimes n}\otimes I)\,\ket0^{\otimes n}\ket-\;=\;2^{-n/2}\sum_{x\in\{0,1\}^n}\ket x\ket- .\]

\textbf{Step 2: the query.} By kickback (row 37),
\(U_f\ket x\ket-=(-1)^{f(x)}\ket x\ket-\), one road at a time:
\[\begin{aligned}U_f\;2^{-n/2}\sum_x\ket x\ket-\;&=\;2^{-n/2}\sum_x(-1)^{(a\cdot x)\oplus b}\ket x\ket-\\&=\;(-1)^b\,2^{-n/2}\sum_x(-1)^{a\cdot x}\ket x\ket- .\end{aligned}\]

\textbf{Step 3: the second layer.} On one rod,
\(H\ket{x_i}=\tfrac1{\sqrt2}\sum_{z_i}(-1)^{x_iz_i}\ket{z_i}\); on all
of them, \(H^{\otimes n}\ket x=2^{-n/2}\sum_z(-1)^{x\cdot z}\ket z\). So
the state becomes
\[\begin{aligned}&(-1)^b\,2^{-n}\sum_{z}\Bigl(\sum_x(-1)^{a\cdot x}(-1)^{x\cdot z}\Bigr)\ket z\ket-\\&\qquad=\;(-1)^b\,2^{-n}\sum_{z}\Bigl(\sum_x(-1)^{x\cdot(a\oplus z)}\Bigr)\ket z\ket-,\end{aligned}\]
since
\((a\cdot x)\oplus(x\cdot z)=\bigoplus_ix_i(a_i\oplus z_i)=x\cdot(a\oplus z)\).

\textbf{Step 4: the sum.} By Result 13.5 with \(c=a\oplus z\), the inner
sum is \(2^n[\,a\oplus z=0\,]=2^n[\,z=a\,]\). Only \(z=a\) survives,
with coefficient \((-1)^b\,2^{-n}\,2^n=(-1)^b\):
\[(H^{\otimes n}\otimes I)\,U_f\,(H^{\otimes n}\otimes I)\,\ket0^{\otimes n}\ket-\;=\;(-1)^b\,\ket a\ket- .\qquad\blacksquare\]

\hypertarget{coeckes-sketchbook-7}{%
\subsection{Coecke's Sketchbook}\label{coeckes-sketchbook-7}}

Draw the smallest web that has everything in it: three rods with
\(a=101\), so rods \(1\) and \(3\) are marked, and rod \(1\)'s tie
passes rod \(2\) on its way to the flag; take \(b=0\). Each rod is a red
point, a yellow box, the control if it has one, a yellow box. The flag
is \(\ket-\), which is \(\tfrac1{\sqrt2}\) times a green \(\pi\) point,
followed by the two red targets. The picture carries \(2^{3/2}\) from
the rods' points, \(\sqrt2\) from the flag's point, and \(\tfrac12\)
from the two bare ties, \(2\) in all, so the left side of Result 13.8 is
\(\tfrac12\) times it.

\textbf{Step 1: dress every dot.} Colour change (Theorem 8.2) read
backwards on a single point: a green \(\pi\) point is a red \(\pi\)
point with a yellow box on its leg, since \(H\sqrt2\ket1=\sqrt2\ket-\).
Then insert a cancelling pair of yellow boxes (\(H^2=I\); row 56) on
each tie, on the flag between its two red dots, and on the flag's
output. Now every dot has a yellow box on every leg it holds:

\[\tfrac12\;
\begin{tikzpicture}[sd, x=1.1em, y=1.1em,
  zsp/.style={circle,draw,line width=0.9pt,fill=green!55,minimum size=0.95em,inner sep=0pt},
  xsp/.style={circle,draw,line width=0.9pt,fill=red!75,minimum size=0.95em,inner sep=0pt},
  hbox/.style={draw,line width=0.9pt,fill=yellow!85,minimum size=0.7em,inner sep=1pt},
  ph/.style={font=\scriptsize,inner sep=1pt}]
  \draw[wire] (2.8,6.6) -- (2.8,0);
  \draw[wire] (5.8,3) -- (5.8,0);
  \foreach \y in {6.6,4.8,3,0} { \draw[wire] (0,\y) -- (8.8,\y); }
  \foreach \y in {6.6,4.8,3} { \node[xsp] at (0,\y) {}; \node[hbox] at (1.3,\y) {}; \node[hbox] at (7.4,\y) {}; }
  \node[zsp] at (2.8,6.6) {}; \node[zsp] at (5.8,3) {};
  \node[xsp] at (2.8,0) {}; \node[xsp] at (5.8,0) {};
  \node[zsp,label={[ph]below:{$\pi$}}] at (0,0) {};
\end{tikzpicture}
\;\overset{\text{colour},\,H^2=I}{=}\;\tfrac12\;
\begin{tikzpicture}[sd, x=1.1em, y=1.1em,
  zsp/.style={circle,draw,line width=0.9pt,fill=green!55,minimum size=0.95em,inner sep=0pt},
  xsp/.style={circle,draw,line width=0.9pt,fill=red!75,minimum size=0.95em,inner sep=0pt},
  hbox/.style={draw,line width=0.9pt,fill=yellow!85,minimum size=0.7em,inner sep=1pt},
  ph/.style={font=\scriptsize,inner sep=1pt}]
  \draw[wire] (2.8,6.6) -- (2.8,0);
  \draw[wire] (5.8,3) -- (5.8,0);
  \foreach \y in {6.6,4.8,3,0} { \draw[wire] (0,\y) -- (8.8,\y); }
  \foreach \y in {6.6,4.8,3} { \node[xsp] at (0,\y) {}; \node[hbox] at (1.3,\y) {}; \node[hbox] at (7.4,\y) {}; }
  \node[hbox] at (1.3,0) {};
  \node[hbox] at (2.8,5.65) {}; \node[hbox] at (2.8,0.95) {};
  \node[hbox] at (5.8,2.05) {}; \node[hbox] at (5.8,0.95) {};
  \node[hbox] at (3.85,0) {}; \node[hbox] at (4.75,0) {};
  \node[hbox] at (6.95,0) {}; \node[hbox] at (7.9,0) {};
  \node[zsp] at (2.8,6.6) {}; \node[zsp] at (5.8,3) {};
  \node[xsp] at (2.8,0) {}; \node[xsp] at (5.8,0) {};
  \node[xsp,label={[ph]below:{$\pi$}}] at (0,0) {};
\end{tikzpicture}\]

\textbf{Step 2: turn the web.} Colour change (Theorem 8.2) on every dot
at once: the two controls turn red, the two targets turn green, and
every box on a dot's leg is absorbed. The pairs on the ties and on the
flag's middle are split, one box absorbed at each end. On rod \(2\),
which has no dot, its two boxes meet and cancel (\(H^2=I\)). One box is
left, the outer half of the pair on the flag's output. The ties now run
from the flag up to the marked rods (Result 13.7).

\textbf{Step 3: copy and fuse, first tie.} The flag's red \(\pi\) point
enters the first green dot, and the copy rule (Theorem 8.3) sends a red
\(\pi\) point up the tie and another on along the flag, at the price of
\(\tfrac1{\sqrt2}\). Up the tie, it fuses (Theorem 8.1) with rod \(1\)'s
red dot and with that rod's red point: rod \(1\) is left holding a
single red \(\pi\) point.

\[\overset{\text{colour}}{=}\;\tfrac12\;
\begin{tikzpicture}[sd, x=1.2em, y=1.1em,
  zsp/.style={circle,draw,line width=0.9pt,fill=green!55,minimum size=1.0em,inner sep=0pt},
  xsp/.style={circle,draw,line width=0.9pt,fill=red!75,minimum size=1.0em,inner sep=0pt},
  hbox/.style={draw,line width=0.9pt,fill=yellow!85,minimum size=0.8em,inner sep=1pt},
  ph/.style={font=\scriptsize,inner sep=1pt}]
  \draw[wire] (2,6) -- (2,0);
  \draw[wire] (4.6,3) -- (4.6,0);
  \foreach \y in {6,4.5,3,0} { \draw[wire] (0,\y) -- (7,\y); }
  \foreach \y in {6,4.5,3} { \node[xsp] at (0,\y) {}; }
  \node[xsp] at (2,6) {}; \node[xsp] at (4.6,3) {};
  \node[zsp] at (2,0) {}; \node[zsp] at (4.6,0) {};
  \node[hbox] at (5.9,0) {};
  \node[xsp,label={[ph]below:{$\pi$}}] at (0,0) {};
\end{tikzpicture}
\;\overset{\text{copy, fusion}}{=}\;\tfrac12\cdot\tfrac1{\sqrt2}\;
\begin{tikzpicture}[sd, x=1.2em, y=1.1em,
  zsp/.style={circle,draw,line width=0.9pt,fill=green!55,minimum size=1.0em,inner sep=0pt},
  xsp/.style={circle,draw,line width=0.9pt,fill=red!75,minimum size=1.0em,inner sep=0pt},
  hbox/.style={draw,line width=0.9pt,fill=yellow!85,minimum size=0.8em,inner sep=1pt},
  ph/.style={font=\scriptsize,inner sep=1pt}]
  \draw[wire] (4.6,3) -- (4.6,0);
  \foreach \y in {6,4.5,3} { \draw[wire] (0,\y) -- (7,\y); }
  \draw[wire] (2.4,0) -- (7,0);
  \node[xsp,label={[ph]above:{$\pi$}}] at (0,6) {};
  \node[xsp] at (0,4.5) {}; \node[xsp] at (0,3) {};
  \node[xsp] at (4.6,3) {};
  \node[zsp] at (4.6,0) {};
  \node[hbox] at (5.9,0) {};
  \node[xsp,label={[ph]below:{$\pi$}}] at (2.4,0) {};
\end{tikzpicture}\]

\textbf{Step 4: copy and fuse, second tie; colour change on the flag.}
The same at the second green dot: a red \(\pi\) point climbs to rod
\(3\) and fuses there, another goes on along the flag, and a second
\(\tfrac1{\sqrt2}\) is paid. On the flag a red \(\pi\) point meets the
last yellow box, and colour change makes it a green \(\pi\) point:

\[\overset{\text{copy, fusion, colour}}{=}\;\tfrac12\cdot\tfrac12\;
\begin{tikzpicture}[sd, x=1.2em, y=1.1em,
  zsp/.style={circle,draw,line width=0.9pt,fill=green!55,minimum size=1.0em,inner sep=0pt},
  xsp/.style={circle,draw,line width=0.9pt,fill=red!75,minimum size=1.0em,inner sep=0pt},
  ph/.style={font=\scriptsize,inner sep=1pt}]
  \foreach \y in {6,4.5,3,0} { \draw[wire] (0,\y) -- (2.4,\y); }
  \node[xsp,label={[ph]above:{$\pi$}}] at (0,6) {};
  \node[xsp] at (0,4.5) {};
  \node[xsp,label={[ph]below:{$\pi$}}] at (0,3) {};
  \node[zsp,label={[ph]below:{$\pi$}}] at (0,0) {};
\end{tikzpicture}\]

\[=\;\tfrac14\,\bigl(\sqrt2\ket1\bigr)\bigl(\sqrt2\ket0\bigr)\bigl(\sqrt2\ket1\bigr)\bigl(\sqrt2\ket-\bigr)\;=\;\ket{101}\ket- .\]

Nothing in the chain depended on how many rods there were or which of
them were marked. Step 1 dresses each dot separately, Step 2 turns each
tie separately, and Steps 3 and 4 climb one tie at a time, paying
\(\tfrac1{\sqrt2}\) for each. So for any \(n\) and \(a\), with \(w\)
marks, the same moves leave a red \(\pi\) point on every marked rod, a
red point on every unmarked rod and a green \(\pi\) point on the flag.
The scalars agree: the left side of Result 13.8 is \(2^{(w-n-1)/2}\)
times the picture, the copies contribute \(2^{-w/2}\), and the \(n+1\)
points are worth \(2^{(n+1)/2}\) times \(\ket a\ket-\), a product of
\(1\). For \(b=1\) the web carries a red \(\pi\) dot on the flag.
Dressed and turned with the rest, it becomes a green \(\pi\) dot, a
\(Z\), and when the flag's red \(\pi\) point passes it the result is
\(Z\sqrt2\ket1=-\sqrt2\ket1\): the sign \((-1)^b\), and nothing else.
\(\;\blacksquare\) See \texttt{Code/Ch13/TurnedWeb.py}, which builds
this form for random secrets at \(n=6\) and reduces it by the rules.

There is a second route, in Cleve et al.'s own shape: leave the web
unturned, copy the green \(\pi\) point through each red target onto its
control (Result 13.4), and let the rods' last yellow boxes turn each
green \(\pi\) between its two boxes into a red \(\pi\), as in §13.2. It
reaches the same points by the same rules in a different order.

\hypertarget{the-verdict-12}{%
\subsection{The verdict}\label{the-verdict-12}}

On length it is close, and the picture wins narrowly. The blackboard's
four steps are short, but the last one rests on the sign sum, a sum over
\(2^n\) strings that needed its own proof. The sketchbook needs four
pictures and five named moves, and forms no sum at all. The blackboard
also has a shortcut that is easy to miss: the sign \((-1)^{a\cdot x}\)
is a product of one-rod signs, so the state after the query is a product
of \(\ket+\)s and \(\ket-\)s, as in the example of §13.3. That makes the
algebra almost as quick as the picture, provided one notices it.

On showing, the picture wins outright. In the sketchbook the secret is
not computed at all: it is the wiring. The marked rods are the ones with
ties, the Hadamards turn the ties round, and the flag's \(1\) climbs
exactly those ties. The blackboard proves that the amplitude of every
wrong string is zero, and it does not show \emph{where} the right one
came from. The picture's advantage has a price, and the price is honesty
about its reach. It works because this oracle is a web. For Deutsch and
Jozsa's general balanced function there are no ties to turn, and only
the blackboard's sum decides. Neither proof says anything about how many
queries a classical method needs; both prove one equality of processes.

\begin{center}\rule{0.5\linewidth}{0.5pt}\end{center}

\hypertarget{the-turned-web}{%
\section{§13.7 The Turned Web}\label{the-turned-web}}

Here is the chapter's master diagram, one instance of the headline drawn
whole. There are six rods and a flag, and the secret is \(a=010100\),
with \(b=0\): rods \(2\) and \(4\) are marked. It is drawn in the
symmetric form, with the flag entering at \(\ket1\) and a yellow box on
the flag on both sides, like every rod. That form agrees with Cleve et
al.'s, since \(H\ket1=\ket-\) and \(H\ket-=\ket1\), and it gives every
wire the same dress: a point, a yellow box, the web, a yellow box.

From left to right: a red point on each rod (\(\ket0\)) and a red
\(\pi\) point on the flag (\(\ket1\)); a yellow box on every wire; the
web, a green dot on rods \(2\) and \(4\), each tied to its own red dot
on the flag; a yellow box on every wire; open outputs. The two red dots
on the flag are drawn unfused, one for each tie, which is what a
gate-by-gate translation gives: one green--red pair per CNOT (row 64).
Fusing them is a legal rewrite, but it gives a different picture of the
same process, and this is the picture we draw. The ties run down past
the rods between without touching them. The picture is called
\textbf{The Turned Web}.

\begin{center}
\begin{tikzpicture}[sd, x=1.2em, y=1.1em,
  zsp/.style={circle,draw,line width=0.9pt,fill=green!55,minimum size=1.0em,inner sep=0pt},
  xsp/.style={circle,draw,line width=0.9pt,fill=red!75,minimum size=1.0em,inner sep=0pt},
  hbox/.style={draw,line width=0.9pt,fill=yellow!85,minimum size=0.8em,inner sep=1pt},
  ph/.style={font=\scriptsize,inner sep=1pt}]
  \draw[wire] (3,8.5) -- (3,0);
  \draw[wire] (5,5.1) -- (5,0);
  \foreach \y in {10.2,8.5,6.8,5.1,3.4,1.7,0} {
    \draw[wire] (0,\y) -- (8,\y);
    \node[hbox] at (1.4,\y) {};
    \node[hbox] at (6.6,\y) {};
  }
  \foreach \y in {10.2,8.5,6.8,5.1,3.4,1.7} { \node[xsp] at (0,\y) {}; }
  \node[xsp,label={[ph]below:{$\pi$}}] at (0,0) {};
  \node[zsp] at (3,8.5) {};
  \node[zsp] at (5,5.1) {};
  \node[xsp] at (3,0) {};
  \node[xsp] at (5,0) {};
  \node[sdlabel,anchor=west] at (8.1,10.2) {$1$};
  \node[sdlabel,anchor=west] at (8.1,8.5) {$2$};
  \node[sdlabel,anchor=west] at (8.1,6.8) {$3$};
  \node[sdlabel,anchor=west] at (8.1,5.1) {$4$};
  \node[sdlabel,anchor=west] at (8.1,3.4) {$5$};
  \node[sdlabel,anchor=west] at (8.1,1.7) {$6$};
  \node[sdlabel,anchor=west] at (8.1,0) {flag};
\end{tikzpicture}
\end{center}

\begin{quote}
\textbf{Result 13.13 (The Turned Web).} Exactly,
\[2^{-5/2}\times\text{The Turned Web}\;=\;\ket{010100}\ket1,\] that is,
\(H^{\otimes7}\,U_f\,H^{\otimes7}\,\ket{000000}\ket1=\ket{010100}\ket1\)
for \(f(x)=010100\cdot x\). Exactly the marked rods read \(1\), and the
flag comes back at \(1\).
\end{quote}

\emph{Proof, on the blackboard.} By Result 13.7 with \(b=0\),
\(H^{\otimes7}U_fH^{\otimes7}=\mathrm{CNOT}_{\text{flag}\to2}\,\mathrm{CNOT}_{\text{flag}\to4}\).
On \(\ket{000000}\ket1\) the flag is \(1\), so both flip their rods:
\(\ket{010100}\ket1\). For the scalar, the seven points give \(2^{7/2}\)
and the two bare ties \(2^{-1}\), so the picture is
\(2^{5/2}\,\ket{010100}\ket1\).

\emph{Proof, in the sketchbook.} The chain of §13.6, with nothing to
prepare on the flag, since it already wears its yellow boxes. Insert a
cancelling pair of yellow boxes on each tie and on the flag's middle
line, between its two red dots. Colour change turns the controls red and
the targets green; on rods \(1\), \(3\), \(5\) and \(6\) the two boxes
cancel. Then the flag's red \(\pi\) point is copied through the first
green dot, climbs the tie and fuses into rod \(2\), while its copy runs
on along the flag's middle line into the second green dot; there it is
copied again, climbs into rod \(4\), and runs on to the flag's output,
whose box has already been absorbed:

\[2^{-5/2}\;
\begin{tikzpicture}[sd, x=1.2em, y=1.1em,
  zsp/.style={circle,draw,line width=0.9pt,fill=green!55,minimum size=1.0em,inner sep=0pt},
  xsp/.style={circle,draw,line width=0.9pt,fill=red!75,minimum size=1.0em,inner sep=0pt},
  ph/.style={font=\scriptsize,inner sep=1pt}]
  \draw[wire] (2.4,8.5) -- (2.4,0);
  \draw[wire] (4.2,5.1) -- (4.2,0);
  \foreach \y in {10.2,8.5,6.8,5.1,3.4,1.7,0} { \draw[wire] (0,\y) -- (6,\y); }
  \foreach \y in {10.2,8.5,6.8,5.1,3.4,1.7} { \node[xsp] at (0,\y) {}; }
  \node[xsp,label={[ph]below:{$\pi$}}] at (0,0) {};
  \node[xsp] at (2.4,8.5) {};
  \node[xsp] at (4.2,5.1) {};
  \node[zsp] at (2.4,0) {};
  \node[zsp] at (4.2,0) {};
\end{tikzpicture}
\;\overset{\text{copy, fusion}}{=}\;2^{-5/2}\cdot\tfrac12\;
\begin{tikzpicture}[sd, x=1.2em, y=1.1em,
  xsp/.style={circle,draw,line width=0.9pt,fill=red!75,minimum size=1.0em,inner sep=0pt},
  ph/.style={font=\scriptsize,inner sep=1pt}]
  \foreach \y in {10.2,8.5,6.8,5.1,3.4,1.7,0} { \draw[wire] (0,\y) -- (2.2,\y); }
  \foreach \y in {10.2,6.8,3.4,1.7} { \node[xsp] at (0,\y) {}; }
  \node[xsp,label={[ph]above right:{$\pi$}}] at (0,8.5) {};
  \node[xsp,label={[ph]above right:{$\pi$}}] at (0,5.1) {};
  \node[xsp,label={[ph]below:{$\pi$}}] at (0,0) {};
\end{tikzpicture}
\;=\;\ket{010100}\ket1 .\]

The first picture is the web turned: the ties now run from the flag. The
last is one point per wire, red \(\pi\) on exactly rods \(2\) and \(4\)
and the flag, worth \(2^{7/2}\ket{010100}\ket1\), and
\(2^{-5/2}\cdot\tfrac12\cdot2^{7/2}=1\). \(\;\blacksquare\) See
\texttt{Code/Ch13/TurnedWeb.py}, which builds The Turned Web, evaluates
it with the scalar, and reduces it by the rules to the same seven
points.

The secret was chosen not to read the same in both directions. Read from
the bottom rod up, the outputs would give \(001010\), the wrong secret;
the book's convention, top wire first, gives \(010100\), and the code
checks that it is this one.

Look at the picture once more with the rewrite in mind. Before the turn
the ties hang from the marked rods down to the flag, as a box of marks
is built: each marked rod adds its bit to the flag. After the turn they
rise from the flag to the marked rods, and the flag's \(1\) is handed up
each of them. What was inside the box below is outside it above. The
rods that no tie touches come back as they went in, and the rods that a
tie touches come back turned, and the secret is read in the return.

\begin{center}\rule{0.5\linewidth}{0.5pt}\end{center}

\hypertarget{what-a-wish-cannot-buy}{%
\section{§13.8 What a Wish Cannot Buy}\label{what-a-wish-cannot-buy}}

Take stock. A query is one use of a reversible box that undoes itself
(Result 13.1). A box that computes \((a\cdot x)\oplus b\) is a web of
CNOTs, one tie for each mark (Result 13.2), and plain questions need
\(n\) of them, or \(n+1\) if \(b\) is unknown, to read its secret
(Result 13.3). With the flag at \(\ket-\), each tie copies a mark back
onto its rod (Result 13.4), which reads Deutsch and Jozsa's test off the
web (Result 13.6). With the whole network crosswise the web turns round,
and one query returns the entire secret (Results 13.7 and 13.8, and The
Turned Web). A pairs box yields to one round of Simon's method only a
single clue, a parity check on the secret (Result 13.9), and on average
at most \(2(n-1)\) rounds find it (Result 13.10).

Now the honest part. Every rewrite in this chapter proves that two
processes are equal: that this network, with this box, returns this
state. The Turned Web proves that one query returns \(a\). It proves
nothing about how many queries any \emph{other} method needs. The claim
that plain questions need at least \(n\) was proved on the blackboard by
counting, in Result 13.3, and no picture took part. It is a claim about
every classical method at once, and a picture shows one method.

The gap it measures is modest. For Bernstein and Vazirani's problem as
posed here, the classical cost is \(n\) queries, or \(n+1\), against
one: a gap that grows in proportion to \(n\), not exponentially.
Bernstein and Vazirani's own separation used a harder, recursive
problem, which BBBV call recursive Fourier sampling, and showed, in
BBBV's account, that relative to an oracle there are problems in BQP
that cannot be solved with small error probability by probabilistic
machines restricted to \(n^{o(\log n)}\) steps. We state that and do not
prove it.

And some questions are no quicker for any quantum method at all. The
sharpest statement of this in the chapter's sources is BBBV's.

\begin{quote}
\textbf{Result 13.14 (Bennett, Bernstein, Brassard and Vazirani, Theorem
3.5; stated, not proved).} ``For any \(T(n)\) which is \(o(2^{n/2})\),
relative to a random oracle, with probability 1, \(\text{BQTime}(T(n))\)
does not contain NP.''
\end{quote}

Here \(\text{BQTime}(T(n))\) is the class of problems accepted with
probability at least \(\tfrac23\) by an oracle quantum machine whose
running time is bounded by \(T(n)\) (BBBV §2); \(T(n)\) is
\(o(2^{n/2})\) when \(T(n)/2^{n/2}\) tends to \(0\); and a random oracle
answers every string with a fair coin, fixed once and for all. In words:
relative to such a box, no quantum machine running in much less than
\(2^{n/2}\) steps solves every problem in NP, the problems whose
solutions can be checked quickly. BBBV explain what the bound means. The
proof, they write, shows that no quantum algorithm ``can gain more than
this quadratic factor in success probability compared to classical
algorithms, when attempting to answer NP-type questions formulated
relative to a random oracle''. It does not show that NP is out of reach
of quantum computers. What it establishes, in their words, is ``that
there is no black-box approach to solving NP--complete problems by using
some uniquely quantum-mechanical features''. The bound is tight: they
note that Grover's search, the subject of Chapter 15, meets it.

Every bound in this chapter, then, is of one of two kinds. Result 13.3
is proved on the blackboard by counting, and Results 13.12 and 13.14 are
stated and cited. None is drawn.

\vskip 0pt plus 4\baselineskip\penalty-100\vskip 0pt plus
-4\baselineskip

\begin{quote}
\textbf{Coecke's Sketchbook.}\nopagebreak

\emph{What a rewrite does not count.} It is tempting to look at The
Turned Web and say that the picture \emph{shows} why one query suffices
where \(n\) were needed. It shows something smaller. It shows that this
network, with this box, returns \(a\). The number \(n\) belongs to a
different question: what every classical method must do. A picture is
one method, drawn. To show that no method does better you would have to
range over all of them, and the rules of Chapter 8 have nothing to say
about methods that are not drawn. There is a subtler temptation too. The
web is small, \(w\) ties and a flag, and one might read its smallness as
the reason the quantum method is cheap. But the classical method has the
same web to work with and still needs \(n\) questions. The size of a
picture is not the cost of a computation. The rewrite is an equality,
and that is all it is.
\end{quote}

\vskip 0pt plus 4\baselineskip\penalty-100\vskip 0pt plus
-4\baselineskip

\begin{quote}
\textbf{Feynman's Notebook.}\nopagebreak

\emph{Four authors and a limit.} ``Strengths and weaknesses of quantum
computing'' is by Charles H. Bennett of IBM Research, Ethan Bernstein of
Microsoft Corporation, Gilles Brassard of the Université de Montréal and
Umesh Vazirani of UC Berkeley. The preprint is dated 12 December 1996
and carries the note ``To appear in SIAM Journal on Computing (special
issue on quantum computing)''. Its introduction is a short history of
the evidence that quantum machines are stronger: Deutsch's quantum
Turing machine, Deutsch and Jozsa, Bernstein and Vazirani's universal
machine and their proof that BPP \(\subseteq\) BQP \(\subseteq\) PSPACE,
then Simon's oracle, and Shor. Simon's paper, ``On the power of quantum
computation'', appears in their references in the proceedings of the
35th IEEE Symposium on Foundations of Computer Science, 1994, pages
116--123; a journal version was published in the \emph{SIAM Journal on
Computing} in 1997. Against that list of strengths the paper sets its
weakness: the random-oracle bound, and the remark that Grover's
algorithm makes it tight.
\end{quote}

\begin{center}\rule{0.5\linewidth}{0.5pt}\end{center}

The chapter's question has an answer now, or rather three. If the box is
a web of marks and we ask in the right way, one question tells us
everything the box holds. If the box keeps its secret in pairs, one
question tells us one parity check the secret passes, and a handful of
questions tell us the rest. And if the box is a table filled by coin
tosses and the question is a search of the kind NP asks, no quantum
method gains more than a square root over the ordinary count; that was
proved by people who counted queries, and we have only stated it. Where
one question did buy something, its pattern was the same each time: a
layer of Hadamards, one query, a layer of Hadamards. Every wire the
query did not touch came back as it went in, and every wire it touched
came back turned.

Chapter 14 keeps the pattern and changes the layers. Simon's round, with
the second layer of Hadamards replaced by a Fourier transform over a
larger group, is where Shor began: BBBV name Simon's technique as one of
the ingredients of Shor's result. The box will hide not a string but a
period, the query will be a controlled multiplication, and the pictures
will again be kickback diagrams, the marks copied back along their ties.
What will change is how much of the work the pictures can do. The number
theory that turns a period into a factor stays on the blackboard, and we
will say so.

\hypertarget{appendix-references}{%
\chapter{Appendix: References}\label{appendix-references}}

The works each chapter leans on, chapter by chapter, with the sections
where they are cited. Where a work is freely available its arXiv number
is given. A work cited in several chapters is listed under each.

\begingroup\raggedright

\hypertarget{chapter-1-not-just-a-bit-1}{%
\section{Chapter 1: Not Just a Bit}\label{chapter-1-not-just-a-bit-1}}

\begin{itemize}
\tightlist
\item
  S. Abramsky and B. Coecke, ``A categorical semantics of quantum
  protocols'', 2004. arXiv:quant-ph/0402130. \emph{Cited in §1.3.}
\item
  R. P. Feynman, ``Simulating Physics with Computers'',
  \emph{International Journal of Theoretical Physics}, 1982.
  doi:10.1007/BF02650179. \emph{Cited in §1.5.}
\item
  M. A. Nielsen and I. L. Chuang, \emph{Quantum Computation and Quantum
  Information}, Cambridge University Press, 2000.
  ISBN~978-1-107-00217-3. \emph{Cited in §1.1 and §1.5.}
\item
  R. Penrose, ``Applications of negative dimensional tensors'', in D. J.
  A. Welsh (ed.), \emph{Combinatorial Mathematics and its Applications},
  Academic Press, 1971, pp.~221--244. \emph{Cited in §1.3.}
\item
  P. Selinger, ``A survey of graphical languages for monoidal
  categories'', 2009. arXiv:0908.3347. \emph{Cited in §1.3.}
\end{itemize}

\hypertarget{chapter-2-amplitudes-1}{%
\section{Chapter 2: Amplitudes}\label{chapter-2-amplitudes-1}}

\begin{itemize}
\tightlist
\item
  S. Abramsky and B. Coecke, ``A categorical semantics of quantum
  protocols'', 2004. arXiv:quant-ph/0402130. \emph{Cited in §2.1 and
  §2.9.}
\item
  P. A. M. Dirac, ``A new notation for quantum mechanics'',
  \emph{Proceedings of the Cambridge Philosophical Society}, 1939.
  doi:10.1017/S0305004100021162. \emph{Cited in §2.1.}
\item
  M. A. Nielsen and I. L. Chuang, \emph{Quantum Computation and Quantum
  Information}, Cambridge University Press, 2000.
  ISBN~978-1-107-00217-3. \emph{Cited in §2.1 and §2.3.}
\end{itemize}

\hypertarget{chapter-3-turning-the-sphere-1}{%
\section{Chapter 3: Turning the
Sphere}\label{chapter-3-turning-the-sphere-1}}

\begin{itemize}
\tightlist
\item
  S. Abramsky and B. Coecke, ``A categorical semantics of quantum
  protocols'', 2004. arXiv:quant-ph/0402130. \emph{Cited in §3.8.}
\item
  M. A. Nielsen and I. L. Chuang, \emph{Quantum Computation and Quantum
  Information}, Cambridge University Press, 2000.
  ISBN~978-1-107-00217-3. \emph{Cited in §3.4, §3.5 and §3.6.}
\item
  P. Selinger, ``A survey of graphical languages for monoidal
  categories'', 2009. arXiv:0908.3347. \emph{Cited in §3.8.}
\item
  J. van de Wetering, ``ZX-calculus for the working quantum computer
  scientist'', 2020. arXiv:2012.13966. \emph{Cited in §3.4 and §3.6.}
\end{itemize}

\hypertarget{chapter-4-uniting-wires-1}{%
\section{Chapter 4: Uniting Wires}\label{chapter-4-uniting-wires-1}}

\begin{itemize}
\tightlist
\item
  B. Coecke and A. Kissinger, \emph{Picturing Quantum Processes},
  Cambridge University Press, 2017. ISBN~978-1-316-21931-7. \emph{Cited
  in §4.9.}
\item
  P. Selinger, ``A survey of graphical languages for monoidal
  categories'', 2009. arXiv:0908.3347. \emph{Cited in §4.6 and §4.9.}
\end{itemize}

\hypertarget{chapter-5-roads-that-cancel-1}{%
\section{Chapter 5: Roads That
Cancel}\label{chapter-5-roads-that-cancel-1}}

\begin{itemize}
\tightlist
\item
  R. Cleve, A. Ekert, C. Macchiavello and M. Mosca, ``Quantum algorithms
  revisited'', 1997. arXiv:quant-ph/9708016. \emph{Cited in §5.1, §5.3,
  §5.4 and §5.5.}
\item
  D. Deutsch, ``Quantum theory, the Church-Turing principle and the
  universal quantum computer'', \emph{Proceedings of the Royal Society
  of London A}, 1985. doi:10.1098/rspa.1985.0070. \emph{Cited in §5.4.}
\item
  D. Deutsch and R. Jozsa, ``Rapid solution of problems by quantum
  computation'', \emph{Proceedings of the Royal Society of London A},
  1992. doi:10.1098/rspa.1992.0167. \emph{Cited in §5.5.}
\end{itemize}

\hypertarget{chapter-6-entanglement-1}{%
\section{Chapter 6: Entanglement}\label{chapter-6-entanglement-1}}

\begin{itemize}
\tightlist
\item
  A. Einstein, B. Podolsky and N. Rosen, ``Can Quantum-Mechanical
  Description of Physical Reality Be Considered Complete?'',
  \emph{Physical Review}, 1935. doi:10.1103/PhysRev.47.777. \emph{Cited
  in §6.1.}
\item
  R. Horodecki, P. Horodecki, M. Horodecki and K. Horodecki, ``Quantum
  entanglement'', 2007. arXiv:quant-ph/0702225. \emph{Source for the
  history in §6.1 and §6.8.}
\item
  R. Penrose, ``Applications of negative dimensional tensors'', in D. J.
  A. Welsh (ed.), \emph{Combinatorial Mathematics and its Applications},
  Academic Press, 1971, pp.~221--244. \emph{Cited in §6.3.}
\end{itemize}

\hypertarget{chapter-7-inequalities-nature-breaks-1}{%
\section{Chapter 7: Inequalities Nature
Breaks}\label{chapter-7-inequalities-nature-breaks-1}}

\begin{itemize}
\tightlist
\item
  J. S. Bell, ``On the Einstein Podolsky Rosen paradox'',
  \emph{Physics}, 1964. CERN CDS record 111654. \emph{Cited in §7.1 and
  §7.2.}
\item
  N. Brunner, D. Cavalcanti, S. Pironio, V. Scarani and S. Wehner,
  ``Bell nonlocality'', 2013. arXiv:1303.2849. \emph{Source for the game
  form in §7.2, for Tsirelson's bound in §7.3, and for the history and
  attributions in §7.2, §7.3 and §7.7.}
\item
  J. F. Clauser, M. A. Horne, A. Shimony and R. A. Holt, ``Proposed
  experiment to test local hidden-variable theories'', \emph{Physical
  Review Letters}, 1969. doi:10.1103/PhysRevLett.23.880. \emph{Cited in
  §7.2.}
\item
  A. Einstein, B. Podolsky and N. Rosen, ``Can Quantum-Mechanical
  Description of Physical Reality Be Considered Complete?'',
  \emph{Physical Review}, 1935. doi:10.1103/PhysRev.47.777. \emph{Cited
  in §7.1, §7.2 and §7.4.}
\item
  D. M. Greenberger, M. A. Horne and A. Zeilinger, ``Going beyond Bell's
  theorem'', 2007. arXiv:0712.0921. \emph{Cited in §7.4, §7.6 and §7.7.}
\item
  B. Hensen, H. Bernien, A. E. Dréau et al., ``Loophole-free Bell
  inequality violation using electron spins separated by 1.3 km'', 2015.
  arXiv:1508.05949. \emph{Cited in §7.7.}
\end{itemize}

\hypertarget{chapter-8-spiders-red-and-green-1}{%
\section{Chapter 8: Spiders, Red and
Green}\label{chapter-8-spiders-red-and-green-1}}

\begin{itemize}
\tightlist
\item
  S. Abramsky and B. Coecke, ``A categorical semantics of quantum
  protocols'', \emph{Proceedings of the 19th IEEE Symposium on Logic in
  Computer Science (LiCS'04)}, 2004. arXiv:quant-ph/0402130. \emph{Cited
  in §8.7.}
\item
  M. Backens, ``The ZX-calculus is complete for stabilizer quantum
  mechanics'', 2013. arXiv:1307.7025. \emph{Cited in §8.5 and §8.7.}
\item
  B. Coecke and R. Duncan, ``Interacting quantum observables:
  categorical algebra and diagrammatics'', 2009. arXiv:0906.4725.
  \emph{Cited in §8.5 and §8.7.}
\item
  R. Penrose, ``Applications of negative dimensional tensors'', in D. J.
  A. Welsh (ed.), \emph{Combinatorial Mathematics and its Applications},
  Academic Press, 1971, pp.~221--244. \emph{Cited in §8.7.}
\item
  P. Selinger, ``A survey of graphical languages for monoidal
  categories'', 2009. arXiv:0908.3347. \emph{Cited in §8.5.}
\end{itemize}

\hypertarget{chapter-9-now-you-look-1}{%
\section{Chapter 9: Now You Look}\label{chapter-9-now-you-look-1}}

\begin{itemize}
\tightlist
\item
  B. Coecke and S. Perdrix, ``Environment and classical channels in
  categorical quantum mechanics'', 2010. arXiv:1004.1598. \emph{Cited in
  §9.3 and §9.6.}
\item
  M. A. Nielsen and I. L. Chuang, \emph{Quantum Computation and Quantum
  Information}, Cambridge University Press, 2000.
  ISBN~978-1-107-00217-3. \emph{Cited in §9.6.}
\item
  W. H. Zurek, ``Decoherence, einselection, and the quantum origins of
  the classical'', 2003. arXiv:quant-ph/0105127. \emph{Source for the
  history in §9.1.}
\end{itemize}

\hypertarget{chapter-10-teleportation-1}{%
\section{Chapter 10: Teleportation}\label{chapter-10-teleportation-1}}

\begin{itemize}
\tightlist
\item
  S. Abramsky and B. Coecke, ``A categorical semantics of quantum
  protocols'', \emph{Proceedings of the 19th IEEE Symposium on Logic in
  Computer Science (LiCS'04)}, 2004. arXiv:quant-ph/0402130. \emph{Cited
  in §10.1, §10.3, §10.7 and §10.8.}
\item
  C. H. Bennett, G. Brassard, C. Crépeau, R. Jozsa, A. Peres and W. K.
  Wootters, ``Teleporting an unknown quantum state via dual classical
  and Einstein-Podolsky-Rosen channels'', \emph{Physical Review
  Letters}, 1993. doi:10.1103/PhysRevLett.70.1895. \emph{Cited in
  §10.3.}
\item
  C. H. Bennett and S. J. Wiesner, ``Communication via one- and
  two-particle operators on Einstein-Podolsky-Rosen states'',
  \emph{Physical Review Letters}, 1992. doi:10.1103/PhysRevLett.69.2881.
  \emph{Cited in §10.4.}
\item
  D. Gottesman and I. L. Chuang, ``Quantum teleportation is a universal
  computational primitive'', 1999 (published in \emph{Nature} as
  ``Demonstrating the viability of universal quantum computation using
  teleportation and single-qubit operations''). arXiv:quant-ph/9908010.
  \emph{Cited in §10.7.}
\item
  J. van de Wetering, ``ZX-calculus for the working quantum computer
  scientist'', 2020. arXiv:2012.13966. \emph{Cited in §10.1 and §10.6.}
\end{itemize}

\hypertarget{chapter-11-copying-is-forbidden-1}{%
\section{Chapter 11: Copying Is
Forbidden}\label{chapter-11-copying-is-forbidden-1}}

\begin{itemize}
\tightlist
\item
  H. Barnum, C. M. Caves, C. A. Fuchs, R. Jozsa and B. Schumacher,
  ``Noncommuting mixed states cannot be broadcast'', 1995.
  arXiv:quant-ph/9511010. \emph{Cited in §11.2, §11.3, §11.5 and §11.8.}
\item
  B. Coecke and R. Duncan, ``Interacting quantum observables:
  categorical algebra and diagrammatics'', 2009. arXiv:0906.4725.
  \emph{Cited in §11.1.}
\item
  D. Dieks, ``Communication by EPR devices'', \emph{Physics Letters A}
  92, 271--272, 1982. doi:10.1016/0375-9601(82)90084-6. \emph{Cited in
  §11.2.}
\item
  A. K. Pati and S. L. Braunstein, ``Impossibility of deleting an
  unknown quantum state'', \emph{Nature} 404, 164, 2000.
  arXiv:quant-ph/9911090. \emph{Cited in §11.2, §11.4, §11.5, §11.6 and
  §11.7.}
\item
  J. van de Wetering, ``ZX-calculus for the working quantum computer
  scientist'', 2020. arXiv:2012.13966. \emph{Cited in §11.1, §11.2,
  §11.5 and §11.7.}
\item
  W. K. Wootters and W. H. Zurek, ``A single quantum cannot be cloned'',
  \emph{Nature} 299, 802--803, 1982. doi:10.1038/299802a0. \emph{Cited
  in §11.2.}
\end{itemize}

\hypertarget{chapter-12-locks-made-of-light-1}{%
\section{Chapter 12: Locks Made of
Light}\label{chapter-12-locks-made-of-light-1}}

\begin{itemize}
\tightlist
\item
  C. H. Bennett and G. Brassard, ``Quantum cryptography: public key
  distribution and coin tossing'', \emph{Proceedings of the
  International Conference on Computers, Systems \& Signal Processing},
  Bangalore, 175--179, 1984. arXiv:2003.06557. \emph{Cited in §12.1,
  §12.2, §12.3, §12.4, §12.5 and §12.7.}
\item
  C. H. Bennett, G. Brassard and N. D. Mermin, ``Quantum cryptography
  without Bell's theorem'', \emph{Physical Review Letters} 68, 557--559,
  1992. doi:10.1103/PhysRevLett.68.557. \emph{Cited in §12.4.}
\item
  D. Deutsch, ``Quantum theory, the Church-Turing principle and the
  universal quantum computer'', \emph{Proceedings of the Royal Society
  of London A} 400, 97--105, 1985. doi:10.1098/rspa.1985.0070.
  \emph{Cited in §12.4.}
\item
  A. K. Ekert, ``Quantum cryptography based on Bell's theorem'',
  \emph{Physical Review Letters} 67, 661--663, 1991.
  doi:10.1103/PhysRevLett.67.661. \emph{Cited in §12.4.}
\item
  N. Gisin, G. Ribordy, W. Tittel and H. Zbinden, ``Quantum
  cryptography'', 2001. arXiv:quant-ph/0101098. \emph{Cited in §12.1,
  §12.2, §12.3, §12.4, §12.7 and §12.8.}
\item
  R. Horodecki, P. Horodecki, M. Horodecki and K. Horodecki, ``Quantum
  entanglement'', 2007. arXiv:quant-ph/0702225. \emph{Cited in §12.4 and
  §12.7.}
\item
  H.-K. Lo and H. F. Chau, ``Unconditional security of quantum key
  distribution over arbitrarily long distances'', \emph{Science} 283,
  2050--2056, 1999. arXiv:quant-ph/9803006. \emph{Cited in §12.8.}
\item
  D. Mayers, ``Unconditional security in quantum cryptography'', 1998.
  arXiv:quant-ph/9802025. \emph{Cited in §12.8.}
\item
  P. W. Shor and J. Preskill, ``Simple proof of security of the BB84
  quantum key distribution protocol'', 2000. arXiv:quant-ph/0003004.
  \emph{Cited in §12.5, §12.7 and §12.8.}
\item
  J. van de Wetering, ``ZX-calculus for the working quantum computer
  scientist'', 2020. arXiv:2012.13966. \emph{Cited in §12.2, §12.3,
  §12.6 and §12.7.}
\item
  S. Wiesner, ``Conjugate coding'', \emph{ACM SIGACT News} 15:1, 78--88,
  1983. doi:10.1145/1008908.1008920. \emph{Cited in §12.1.}
\end{itemize}

\hypertarget{chapter-13-asking-the-oracle-1}{%
\section{Chapter 13: Asking the
Oracle}\label{chapter-13-asking-the-oracle-1}}

\begin{itemize}
\tightlist
\item
  C. H. Bennett, E. Bernstein, G. Brassard and U. Vazirani, ``Strengths
  and weaknesses of quantum computing'', preprint dated 12 December
  1996. arXiv:quant-ph/9701001. \emph{Cited in §13.1, §13.2, §13.4 and
  §13.8.}
\item
  E. Bernstein and U. Vazirani, ``Quantum complexity theory'',
  \emph{Proceedings of the 25th Annual ACM Symposium on Theory of
  Computing}, 11--20, 1993. doi:10.1145/167088.167097. \emph{Cited in
  §13.3, §13.4, §13.5, §13.6 and §13.8.}
\item
  R. Cleve, A. Ekert, C. Macchiavello and M. Mosca, ``Quantum algorithms
  revisited'', 1997. arXiv:quant-ph/9708016. \emph{Cited in §13.1,
  §13.2, §13.3, §13.6 and §13.7.}
\item
  D. Deutsch, ``Quantum theory, the Church-Turing principle and the
  universal quantum computer'', \emph{Proceedings of the Royal Society
  of London A} 400, 97--105, 1985. doi:10.1098/rspa.1985.0070.
  \emph{Cited in §13.2 and §13.8.}
\item
  D. Deutsch and R. Jozsa, ``Rapid solution of problems by quantum
  computation'', \emph{Proceedings of the Royal Society of London A}
  439, 553--558, 1992. doi:10.1098/rspa.1992.0167. \emph{Cited in §13.2,
  §13.3, §13.6 and §13.8.}
\item
  D. R. Simon, ``On the power of quantum computation'', \emph{SIAM
  Journal on Computing} 26, 1474--1483, 1997.
  doi:10.1137/S0097539796298637. \emph{Cited in §13.4, §13.5 and §13.8.}
\end{itemize}

\endgroup

\backmatter
\end{document}
